% 日本語 LuaLaTeX。2026-11-25 / v4.70
% エンジン、出力先、SyncTeX は repo ルートの .latexmkrc に従う。
\RequirePackage{fix-cm}
\documentclass[a4paper,11pt]{ltjsarticle}
\usepackage{amsmath,amssymb,amsthm,mathtools,mathrsfs}
\DeclareFontFamily{U}{rsfs}{\skewchar\font127}
\DeclareFontShape{U}{rsfs}{m}{n}{<-6> rsfs5 <6-8> rsfs7 <8-> rsfs10}{}
\usepackage[margin=25mm]{geometry}
\usepackage{booktabs,tabularx,array,enumitem}
\usepackage[unicode,pdfusetitle,colorlinks=true,
  hypertexnames=false,linkcolor=blue,citecolor=blue,urlcolor=blue]{hyperref}
\hypersetup{pdftitle={非可換二次曲面のavatarとAS正則Z代数}}
\numberwithin{equation}{section}
\theoremstyle{plain}
\newtheorem{theorem}{定理}[section]
\newtheorem{proposition}[theorem]{命題}
\newtheorem{lemma}[theorem]{補題}
\newtheorem{corollary}[theorem]{系}
\theoremstyle{definition}
\newtheorem{definition}[theorem]{定義}
\newtheorem{assumption}[theorem]{仮定}
\newtheorem{example}[theorem]{例}
\newtheorem{problem}[theorem]{問題}
\theoremstyle{remark}
\newtheorem{remark}[theorem]{注意}
\renewcommand{\proofname}{証明}
\newcommand{\kk}{k}
\newcommand{\OO}{\mathcal O}
\newcommand{\PP}{\mathbb P}
\newcommand{\Gm}{\mathbb G_m}
\newcommand{\HH}{\mathcal H_2}
\newcommand{\HHstack}{\mathfrak H_2}
\DeclareMathOperator{\Hom}{Hom}
\DeclareMathOperator{\RHom}{RHom}
\DeclareMathOperator{\End}{End}
\DeclareMathOperator{\Ext}{Ext}
\DeclareMathOperator{\Aut}{Aut}
\DeclareMathOperator{\GL}{GL}
\DeclareMathOperator{\PGL}{PGL}
\DeclareMathOperator{\SO}{SO}
\DeclareMathOperator{\SL}{SL}
\DeclareMathOperator{\Spin}{Spin}
\DeclareMathOperator{\Gr}{Gr}
\DeclareMathOperator{\Spec}{Spec}
\DeclareMathOperator{\rank}{rank}
\DeclareMathOperator{\Pic}{Pic}
\DeclareMathOperator{\Br}{Br}
\DeclareMathOperator{\Jac}{Jac}
\DeclareMathOperator{\coh}{coh}
\DeclareMathOperator{\qgr}{qgr}
\DeclareMathOperator{\Qgr}{Qgr}
\DeclareMathOperator{\Perf}{Perf}
\DeclareMathOperator{\Sym}{Sym}
\DeclareMathOperator{\coker}{coker}
\DeclareMathOperator{\im}{im}
\DeclareMathOperator{\Supp}{Supp}
\title{非可換二次曲面のavatarとAS正則$\mathbb Z$-代数：\\
  八・十・十二本矢の分類と内在的忘却}
\author{}
\date{2026年11月25日\quad v4.70（全指定相対スタックの比較と適用範囲）}
\begin{document}
\maketitle
\begin{abstract}
非可換二次曲面の三つの表示、すなわち最小射影分解型を指定した
AS 正則 $\mathbb Z$-代数、反標準曲線と線束による幾何データ、
および八・十・十二本矢の重み付きポテンシャルを相対的に比較する。
成分が局所自由な指定 AS 族は、有限な四頂点の窓と全整数添字の双対線束から
同型を含めて再構成できる。四周期標識の存在はこれらの線束に依存する。
指定表示の適合性を有限窓で判定し、余分な自己同型と、
各標準頂点を構造対象に選んだときの適合開を記述する。

構造対象付き曲面のスタックでは、二つの ruling の順序を選ぶ操作が
全域の有限エタール二重被覆を与える。cubic の稠密開ではこれを
Okawa--Ueda のポテンシャル表示からの四重射と比較できる。
二つの乗法がともに階数三となる noncubic 境界まで含めて、
三本の半不変式から全スタックの準射影的良モジュライ空間を構成する。
\end{abstract}
\tableofcontents

\section{Introduction}
\label{sec:qa-introduction}

本論文の第一の主定理は、
\emph{avatarと、最小射影分解型を指定したAS正則$\mathbb Z$-代数}
の同型類の全単射である。
その証明は第~\ref{sec:qa-as-avatar-classification}節に、
十本矢の独立ASクラスとQP表示の正確な関係、
rank-three壁、mutation後のNoether性・商圏同値は
第~\ref{sec:qa-as10-mutation-classification}節に置く。
同節では二重壁のscheme-theoretic横断性と
非零二捩れの相対有限étale三次も記述する。
二つのmutation壁の交点と、cubic型への到達不能性の完全判定は
第~\ref{sec:qa-as-cubic-obstruction}節で証明する。
最小分解からの有理Hilbert級数の導出と
AS--Gorenstein性の証明を別の論理段階として扱う。
また noncubic 二形式に自然な点曲線$E_f$と
spectral 曲線$C_q$が2-isogenousだが一般に非同型であることを
第~\ref{sec:qa-point-spectral-isogeny}節で説明し、
全noncubic幾何点では$E_f$の二conicファイバー積と
二pencilを用いて、mutation後のAS heartとの
pointed認識を証明する。
第~\ref{sec:qa-point-spectral-isogeny}節では
MCM行列kernelから二pencilを回収し、
四頂点Endの全Homと積を点曲線への制限で同定する。
標準頂点$H_0$でなくshift後の$H_1$に構造対象を対応させる。
順序付き分解の普遍族では、有限graded End代数の比較を
任意の非被約基底変更へ延長する。
第~\ref{sec:qa-root-descent-brauer}節では
根対内の順序交換の厳密な降下と、
因子のスカラー変換も忘れた一般二次形式では
conicのBrauer類が非零となりpencilの線束降下を妨げること、
分岐根選択による次数四の\'{e}tale分裂被覆、
およびBrauer gerbe上の相対avatarと根対相互交換の
ruling対合を証明する。
さらに第~\ref{sec:qa-twisted-morita-weights}節で
determinant補正後の四項のBrauer重みを計算し、
Azumaya双加群の有向Morita圏を係数空間へ構成し、
階数54の通常の有限代数としても明示する。
一般点のruling交換は頂点固定Moritaでは実現できない。

この障害を越えるため、第~\ref{sec:qa-symmetric-morita-descent}節では
$(\mathcal O,X_c,X_d,\mathcal O(1))$への相対mutationを構成する。
二形式交換が中間二頂点を交換する対称表示を使い、
共通座標表現による増幅で全商$[P/G]$へ階数$45$の通常代数を降下させる。
実際のspinor generatorとの比較により、
構造対象を零頂点に固定したcompact Morita同値を得る。
第~\ref{sec:qa-relative-serre-completion}節は有限separable対角bar解消を用いて
全tensor次数の相対Serre完成と任意基底変更を扱い、
零頂点の反Serre角によって元の相対pointed heartを回収する。
ここで二形式交換の符号指標$\varepsilon$を含むcanonical Serre作用を計算し、
論文Bの正規化表示との次数線束捩れを明示する。

第~\ref{sec:qa-spinor-pointing-obstruction}節では全point schemeを
曲線成分と二つの孤立点の非交和として直接計算し，spinorの生成障害を得る．
第~\ref{sec:qa-general-vertex-az}節は，頂点を周期的に消した商の有限性を
AZ豊富性と同定し，二つのmutation壁と論文Aのcubic逆構成を用いて
全標準頂点のH2所属を決定する．この体上判定には非被約曲線も含む．
楕円base locus外の点加群やfat pointが現れること自体，および
MCM加群の生成性をtrace商で検出する方法には既知の類例がある
\cite[\S10.5]{SVdB2013}\cite[Lemma~6.6 and proof of Proposition~6.10]{RogalskiSierraStafford2020}．
本稿では中央二重平方表示の二周期二点と，指定標準頂点の全反Serre軌道を扱う．
これらの具体的判定の包括的新規性比較は未完了である．

第~\ref{sec:qa-spectral-pointed-atkin-lehner}節では、
spectral曲線のavatarに対して同じ指定AS実現を行うと、
一般に元と異なるpointed対象が得られることを証明する。
二同種の商公式\cite[III, Example~4.5]{SilvermanAEC2009}と
点曲線型の逆構成を組み合わせると、smooth幾何同型類上の対応は
$\chi\mapsto\chi/(\chi-1)$となる。
自己双対点は正確に$\chi=0,2$であり、
曲線の$j$値だけではpointed同値を判定できない例も示す。
さらに共有根と一致根対の双方への退化で像の$\chi$に極が現れ、
noncubic境界の閉性により、H2全体へも族として延長不能となる。
これは固定spectral曲線上のPicard線束の延長障害より強い主張である。
二同種の双対化そのものは古典的であり、
本稿の追加部分は指定pointingとの同定と内在的境界への適用である。
この適用の新規性について網羅的な文献調査は未完了である。

相対構成は第~\ref{sec:qa-spectral-relative-nonfaithful-map}節で与える。
二つのinvolutionの商から順序付き二conic被覆を作ると、
その逆構成のfiberは二捩れPicard類を実際の二次線束へ持ち上げる
$\mu_2$-gerbeである。順序忘却を合成した幾何fiberは
$B\mu_2\amalg B\mu_2$となる。
一般慣性の中心の二次分裂体を比較し、
粗対合を覆ういかなる射も一般点で表現可能にはならないことも示す。
さらに両境界に対する平坦な代数表示を構成し、
周期解消から特殊fiberの全高次Extを計算する。
この計算は二重根または零形式を持つ全中央二重平方表示へ適用され、
その有界coherent導来圏の右Serre関手の不存在を与える。


第~\ref{sec:qa-half-period-ruling}節ではsmooth locus上で
$2t=P$の次数四有限étale torsorを作り、平行移動による
二pencilと有限Endの幾何比較に加え、
Brauer差のindex四と点曲線のperiod二・index四、
半周期torsorが二conicを同型化する最小次数四の分裂被覆であることを証明する。
分岐根を選ぶもう一つの四重被覆とは同型でないことも、
全根が分裂した三変数族の二同種模型で直接示す。
同じ模型で独立平方根の符号変更と根対内部交換を合成し、
半周期torsorの一般$D_4$ monodromyを確定する。
OU側の一般$S_4$ monodromyは非零二捩れを選ぶ
次数三の標識被覆で$D_4$へ減少するが、
さらに二乗固有値の二対分割を持つ共通incidence基底上に
OUのspectral曲線とnoncubicの$C_q$を同定する。
そこでは両四重射のmonodromyは共に$D_4$だが、
全固有値体への引戻しでOU側のみ分裂するため非同型と証明する。
さらに二torsion局所系$J(C_q)[2]$と$J(E_f)[2]$も
その体上で非同型である。
自然な二同種のKummer写像
$\phi[2]_*:H^1(J(C_q)[2])\to H^1(J(E_f)[2])$で
OUの四重射を押し出しても、
noncubicの半周期Kummer類とは一致しない。
点曲線側の相対mutation・pointed heartへの降下は上記の構成で得られるが、
任意のavatar/AS/QP表示stackの同型層を含む認識は別問題である。
ここでは有限directed foundation、relation-dual QP、
underlying AS代数、O-marked内在的heartの四水準を混同しない。

本論文の第二の問題は、\textbf{O-marked quadric、potential、avatar}を
単に隣接する記述として並べるのではなく、
具体的な同型類の全単射と、それを誘導するモジュライ間の射で結ぶことである。
特に、potentialの頂点と四項foundation、avatarの二ruling、
内在的な構造対象$\OO$をそれぞれ区別して忘却を行う。
滑らかな楕円cubicの三線束辞書と$r_1$商、
共有根・根対一致を含む全noncubic境界の二捩れavatarと二根対の
幾何的全単射、およびrank-one torsion-free spectral中間射と
可逆Picard延長障害を本稿の具体的な比較定理とする。
幾何的基礎は第~\ref{sec:qa-self-contained-foundations}節と
第~\ref{sec:qa-self-contained-helix}節で自己完結に証明し、
同型類・相対spectral射の証明は
第~\ref{sec:qa-avatar-charts}節に置く。
共有根の最小解消と奇偶障害、
さらに二次積のscheme-theoretic欠損と三型の
正規化極限・固定曲線上の慣性・
二次積付きtorsion-free相対groupoidへの射は
第~\ref{sec:qa-nodal-resolution-parity}節で証明する。
$\Sigma_2$の最深境界・非被約因子・古典零切断・全相対AS認識は
結論に含めず、必要な追加構成を明記する。


三次元 AS 正則 $\mathbb Z$-代数の四周期表示では、同じ幾何的構成に対して
八本、十本、十二本の矢を持つ箙が現れる。
これらを別々の分類集合として並べるだけでは、
生成元と関係が同時に消去される変形や、異なる foundation の比較を捉えられない。
本稿は、まず共通の有限代数のモジュライを構成し、その上の二つの開条件と
二つの階数低下条件として矢数を整理する。
さらにgeometric cubic部分の隣接四重tensorの非退化性は
四次4式・六次16式の明示的多項式判定へ書き換える。
加えて第~\ref{subsec:qa-two-mass-finite-gerbe}節では、
raw八・十・十二矢potentialの非最小表示を同じ有限foundationに対応させ、
全二mass（nilpotentを含む）で余剰ゲージ慣性を明示する。
有限表示の相対同型層と、全次数AS完成や心の一致は別の問題として扱う。
次に反標準 avatar の幾何的 chart を比較し、
第~\ref{subsec:qa-two-mass-rigidification-ruling-glue}節では
明示$H_2$開上の全QP表示をmass/頂点scalarsで剛体化し、
Serre完成を経て二rulingを相対mutationで貼る。
通常のright-equivalenceを全heartの同値と混同しない。
この全相対認識での左右Noether性には非対称な注意が必要である．
有限squareの反対代数は四交差$H_2$開を保つが，
逆順にした\emph{normal}四項の零番頂点は
一重rank-three壁では強い$H_2$ pointingとならない．
その具体的障害を命題~
\ref{prop:qa-q12-opposite-normal-rank-three-obstruction}で示す．
左Noether性は，反対squareの正しいdual heartに
逆順双対normal四項を置き，
その\emph{第二}の頂点が構造対象$Q_0$であることを使って
Artin--Zhangのample条件を証明する
（補題~\ref{lem:qa-q12-opposite-normal-dual-heart-az}）．
相対AZ2の二段presentationとNoetherianな一般End環の仮定は
命題~\ref{prop:qa-relative-flat-foundation-az}，
補題~\ref{lem:qa-relative-az-endring-finiteness}で確認する．
AZ再構成の原理は既知入力
\cite[Theorem 4.5(1)]{AZ1994}
\cite[Theorem 7.3]{OACU2026}であり，
本稿の主張は適合pointingの選択・相対Extと境界を含む具体的適用である．
最後に構造対象付き Abel 圏への忘却に必要な条件を分離し、
Okawa--Ueda の quintuple/potential GIT と比較する。
特に滑らかな楕円型の稠密開では二つのnormalized rulingが
有限étale二重忘却の二sheetになり、そのstack商として
内在的モジュライを記述できる。
さらに二捩れ特別座では、内在的慣性の法表現を求め、
OU不変量による粗分岐式と交点の横断$A_1$局所型まで計算する。
これらを既知のpotential/quintupleの次数四射や
Picard根の四選択とは区別して扱う。
複素乗法で$j=0,1728$となっても二cubic開では追加の交換はなく、
一cubic境界ではfiberのsheetを一つ失うことも区別する。

\subsection{三つの表示と内在的モジュライの比較}

本稿で扱う問題には、次の三つの水準がある。
とくに既存GIT商を内在的モジュライと結び付ける相対比較を本稿Cの
中心研究課題とする（第~\ref{sec:qa-git-comparison}節）。
\begin{enumerate}
\item \emph{有限表示のモジュライ。}
  四頂点間の Hom 次元と結合律を固定した代数スタックは有限型 Artin である。
  隣接する二つの乗法が両方とも同型である部分、一方が同型である二つの部分、
  双方に一次元の欠損を許す部分を比較する。
  これが八本・十本・十二本の開部分の正確な有限代数版である。
\item \emph{AS 代数と avatar。}
  四周期の全 AS 条件と幾何的 avatar の対応には、有限 foundation の分類に加えて
  全単純分解、noetherian 性、相対平坦性、降下が必要である。
  体上の既存原稿の分類結果と、今回の有限代数・幾何的 chart の証明を区別する。
\item \emph{内在的な対象。}
  構造対象 $\OO$ だけを指定した $\HHstack$ は、選んだ箙、potential、
  foundation、平方根、ambient を忘れる。
  しかしそれらの存在条件を失うわけではない。
  特に AS3 正則性や $\Sigma_2$ という有限表現モジュライだけから、
  $\HHstack$ の全相対条件は導けない。
\end{enumerate}

有限代数版では、二種類の十本矢 chart の交わりが八本矢 chart である。
二つの十本矢 chart の和集合の補集合が最小十二本矢の層となる。
これに対し、各 chart の内部で正確に十本が最小生成矢として残る部分は
八本矢部分の境界である。
このため、本稿では上限を許す chart と最小本数が一定の層を同じ記号で表さない。

直積型 collection
\[
 (\OO,\OO(1,0),\OO(0,1),\OO(1,1))
\]
そのものの有限箙は八本、四本の最小関係の双対を加えた relation QP は十二本である。
固定した一頂点 mutation により、八本と一方の十本を含む chart が得られる。
従って直積型と十本矢 chart の対応は mutation を介した対応であり、
元の箙の矢数を十本と数えるという意味ではない。

\subsection{先行研究と依存関係}

有限箙の関係式を moduli stack として扱う一般的枠組みは
Abdelgadir--Okawa--Ueda \cite[Definitions 3.5, 4.1; Propositions 3.6, 4.2]{AOU2014}
による。本稿は四頂点の Hom 次元を固定して結合律を具体的な行列方程式で書き、
二つの乗法の階数が最小箙に与える影響を計算する。
これは一般論の新しい証明を主張するものではない。
Derksen--Weyman--Zelevinsky \cite[Theorem 4.6; Section 5]{DWZ2008}
は完備道代数の potential の splitting と mutation を構成した。
ここで必要な消去は内部次数四のため有限回の三角変換で済み、
通常の道代数にも任意基底にも同じ式を適用できる。
ただし pair の消去と quotient stack の慣性の消去は別の操作である。

AS 正則 quadric の cubic 型については Van den Bergh \cite{VdB2011}、
その幾何的データと potential のモジュライについては
Okawa--Ueda \cite[Theorem 1.1]{OU2014} が基礎となる。
同定理 (4) の有限射の向きと、本稿で扱う二段の忘却を次に示す。
上段の $M$ は複素数体上のコンパクト GIT 商、下段は標数零の代数閉体上で
共通の稠密開に制限した有効表示スタックと構造対象付きスタックである。
\[
 \begin{aligned}
 M_{\mathrm{pot}}=\mathbb P(1,2,3,4)
   &\xrightarrow{\ 4\ } M_{\mathrm{qui}}=\mathbb P(2,4,4,6),\\
 \mathfrak P^{\circ}
   &\xrightarrow{\ 4\ } \mathfrak Q^{\circ}
    \xrightarrow{\ 2\ } \mathfrak H_2^{\circ}.
 \end{aligned}
\]
下段で一つの構造対象付き曲面の上にある二点は二つの normal ruling の選択、
それぞれの上の四点はその quintuple を与える potential の選択である。
全域で有限エタールに延びる ruling の被覆と、既存 GIT 商への大域射の
区別は noncubic 曲線に現れる。後者の二射はこの曲線をそれぞれ一点に送る
（定理~\ref{thm:qa-global-two-ou-morphisms}）。

表示側の四次元スカラー慣性の核を命題~\ref{prop:qa-ou-scalar-gerbe}、
剛体化後にも残る一般有限慣性を命題
~\ref{prop:qa-ou-generic-quintuple-finite-inertia}で計算する。
OUの一般安定化群という表現は単位連結成分と全慣性を区別して読む。
また、potential$\to$quintupleの既知GIT次数四射を
定理~\ref{thm:qa-ou-rigid-potential-quintuple-generic-etale}で
一般の表現可能有限étaleスタック射へ強める。
さらに命題~\ref{prop:qa-ou-generic-inertia-faithful-to-h2}では
点曲線の点対象からpointed忘却の一般慣性単射を証明する。
論文Bの普遍smoothnessと合わせた一般étale性・KS同型は
定理~\ref{thm:qa-quintuple-h2-generic-etale}、
一般次数二とpotential合成の次数八は
系~\ref{cor:qa-quintuple-h2-generic-degree-two}と
系~\ref{cor:qa-potential-h2-generic-degree-eight}に示す。
さらに稠密開でのスタック的な商同型は
定理~\ref{thm:qa-h2-generic-ruling-quotient}で与える。
Artin形式変形上の二normal foundationの一意持上げは
命題~\ref{prop:qa-ou-formal-two-ruling-lifts}に示す。
さらに全H2のnormal ruling選択の有限étale二重被覆は
定理~\ref{thm:qa-global-two-ruling-etale-cover}、
finite foundationの実現可能開への開埋込と
全相対$S_2$商同型は
定理~\ref{thm:qa-intrinsic-h2-open-in-foundation-stack}に与える。
有限多項式判定の存在は系~\ref{cor:qa-realizable-foundation-finite-open-criterion}に、
全square表示上での四組・計二十式以内の具体的判定は
定理~\ref{thm:qa-square-explicit-h2-open}と
系~\ref{cor:qa-square-h2-twenty-polynomials}に示す。
被覆の可換gerbe上での非自明性は系
~\ref{cor:qa-ruling-cover-nontrivial-commutative-gerbe}で確認する。
二捩れ特別座のspectral方程式と粗局所分岐は
命題~\ref{prop:qa-ou-spectral-torsion-divisors}、
定理~\ref{thm:qa-ou-special-reflection-coarse-branch}、
系~\ref{cor:qa-ou-coarse-branch-invariant-equations}に、
交点の慣性・横断$A_1$モデルは
定理~\ref{thm:qa-ou-double-torsion-intersection-local-model}に示す。
片側二捩れ特別座の全慣性群型は
命題~\ref{prop:qa-elliptic-special-inertia-group-types}に示す。
特殊ruling交換の正確な二捩れ条件は
命題~\ref{prop:qa-elliptic-ruling-swap-two-torsion}に証明する。
特別点の2-fiberは命題
~\ref{prop:qa-special-ruling-two-fiber-no-stack-ramification}に示す。
square・rotated square の反標準 avatar と十本矢 mutation については、
大川--植田の日本語草稿 v7 \cite{AvatarVSeven2026}も参照されたい。
本稿では、指定十二矢表示から幾何avatarへの体上の逆復元と
四層の分類を定理~\ref{thm:qa-avatar-presented-field-dictionary}と
系~\ref{cor:qa-avatar-four-stratum-dictionary}で証明する。
相対再埋め込みと並進は命題~\ref{prop:qa-pushforward}および
定理~\ref{thm:qa-translation}で構成する。
全mass配置の相対比較に用いる階数、消去、有限窓の再構成も本文で証明する。

構造対象付き $\HH$ の体上の分類は \cite{QuadricIntrinsic2026}、
相対スタック $\HHstack$ の代数性と境界の幾何は \cite{QuadricModuli2026}
を用いる。本稿ではこれらの結果を表示からの忘却射の標的に適用する。
全 $\HHstack$ の代数性と既存の相対開埋込は論文Bの入力であり、
本稿は有限square商とその明示開を通じて全群oidを比較する。
任意に別途与えたAS/QP族の全相対認識、および
公刊に向けた包括的な新規性確認は未完了である。

全体の良モジュライについては、定理~\ref{thm:qa-global-h2-quasiprojective-good-moduli}が
三本の同じ指標の半不変式による構成を与える。
四次行列$M$の左右kernel条件を
$P(M)=\det M$と$Q(M)=e_3(JM^{\mathrm{t}}JM)$でまとめ、
二つの交差縮約から同一指標の三式を作る。
各非消滅開はアフィンで、二式の比は不変な正則関数になるので、
局所商の重なりは飽和している。
線型簡約群のアフィン商と飽和開の貼合せ自体は
\cite[Proposition 7.9, Theorems 13.2 and 13.6]{AlperGoodModuli2013}の既知の原理である。
本稿はこの原理を全$H_2$開へ適用するための半不変式と共通指標を明示する。
この被覆に用いる三式と、共通指標に関する半不変式環の生成元は役割が異なる。
定理~\ref{thm:qa-global-h2-full-semiinvariant-ring} は、後者を ruling 順序付きで
五生成一関係式、順序を忘れると重み $(2,2,3,4)$ の四変数多項式環と決定する。

\paragraph{固定表示のモジュライとの比較}
四量子ビットの不変式環と重み付き射影商は
\cite[\S3, (3.9)--(3.16)]{OU2014}の既知結果である．
Mori--Okawa--Ueda の関係式スタックと elliptic quadruple の双有理比較でも，
四因子 tensor と三次元の自明 scalar 核が明示されている
\cite[Theorem 1.1, \S5.3, (5.24)--(5.26)]{MOU2019}．
また固定 helix 型の marked AS surface は，very strong helix を標識に持つ
potential stack の AS-regular 開として扱われ，elliptic tuple との比較は双有理である
\cite[Definition 1.1, Theorem 6.4]{OU2025MarkedAS}．
従って，固定表示の商を作ること自体や，その AS 開の存在を本稿の成果とはしない．
本稿の対象は構造対象を保つ完全同値を射とする全相対 strong $H_2$ であり，
二つの積がともに階数三となる noncubic 部分も含む．
前記の良モジュライ定理はこの全開を覆い，cubic chart では既存商との
比較を命題~\ref{prop:qa-global-cubic-four-qubit-good-moduli}で与える．
さらに命題~\ref{prop:qa-cubic-adjugate-explicit-invariants}と
系~\ref{cor:qa-cubic-geometric-two-complement-curves}では，
$e_i=e_i(JT^{\mathrm{t}}JT)$，$d=\det T$ により二半不変式と被約補曲線を明示する．
命題~\ref{prop:qa-cubic-weighted-inversion-gluing}と
系~\ref{cor:qa-both-invertible-good-moduli-affine-open}は，
この開の二つのコピーを貼り合わせ，両積可逆な非順序商を
三次元 affine 空間の直線の補として記述する．
定理~\ref{thm:qa-global-two-ou-morphisms}は全域の既存 GIT 商への二射を与えるが，
noncubic 曲線を収縮するため，全良モジュライの埋込みではない．
この具体的な相対範囲の比較と，包括的な新規性の確認は区別する．

\paragraph{既知の降下原理と本稿の具体化}
局所対象を Hom の双対で増幅してスカラーの曖昧さを消す操作は，
twisted sheaf における重みの相殺と重み零の降下という既知の原理に属する
\cite[Lemma 2.1.1.13]{LieblichTwisted2007}
\cite[Lemma 3.1.1.7]{LieblichPeriodIndex2008}．
増幅した対象の対角 End が Azumaya 代数となり得ることも，
同じ理論の基本的な現象である
\cite[Proposition 3.1.2.1 and its proof]{LieblichPeriodIndex2008}．
また古典的な $\mathbb P^1\times\mathbb P^1$ の twisted form については，
二つの中間 ruling を一つの block にまとめた分解の降下と，
順序付き square への精密化を妨げる二重被覆が
Antieau--Elmanto により記述されている
\cite[Section 5.3, Proposition 5.7]{AntieauElmanto2021}．
本稿で行うのは，これらの原理を $\mathfrak H_2$ の内在的 square に適用し，
指定対象 $\mathcal O$，ruling 交換および全許容基底変更と整合する
\[
 T^{\mathrm{can}}
   =\bigoplus_{i=0}^3
       E_i\otimes\operatorname{Hom}(\mathcal O,E_i)^\vee
\]
を古典型も含む全 $\mathfrak H_2$ 上で具体的に構成することである．
局所重複度 $(1,2,2,4)$ は階数 $81$ の通常の有限代数を与え，
その canonical Serre 完成の各非負次数の局所自由性と，
指定 $\mathcal O$ を保つ actual tails の認識までを証明する
（命題~\ref{prop:qa-universal-h2-square-morita81}，
定理~\ref{thm:qa-universal-h2-rank81-completion}）．
一般の降下原理そのものや，階数 $81$ の最小性を新規な主張とはしない．


\subsection{規約}

$\kk$ は標数零の代数閉体とする。
族は $\kk$ 上局所有限型スキームの fppf site 上で考え、非被約基底を含む。
矢の積は合成順で、$ab$ は先に $b$、次に $a$ を通る道である。
有限 foundation は四つの順序付き頂点を保持し、
右加群規約に対応する導来同値と標準 heart の exact 同値を区別する。
基底上の係数を消去に使う条件は可逆性であり、単に零でないことではない。

境界比較の具体的な構成は，第~\ref{sec:qa-spectral-boundary-models}節の
共通正則模型，pencilの多重次数，四捩れ補正の全torsorを経て，
第~\ref{sec:qa-global-fricke-boundary}節の大域的Fricke射へ進む．
blowup，形式関数定理，最小正則模型の一意性は既知の一般結果
\cite{StacksBlowups,StacksFormalFunctions,StacksMinimalModels}を用いる．
本稿で直接証明する部分は，指定された二模型の比較，入力二捩れの全有限分岐後の障害，
全表示群に対する出力作用とその降下，および線束境界上の全fiberである．
これらの組合せについて包括的な新規性の確定は行っていない．

一般square曲線族の相対比較には，二つの乗法行列式を用いる．
定理~\ref{thm:qa-relative-vertex-determinant-open} は全標準頂点の幾何適合開を
非被約底変更まで記述し，命題~\ref{prop:qa-relative-universal-cartier-walls} は
普遍楕円Picard上の二壁のscheme構造を決定する．
さらに定理~\ref{thm:qa-relative-square-e0-e3-h2} は，実際のconeによる
rank十六からrank二十四へのnormal窓を通じて，零番・三番頂点の
相対$\HH$族とpointed compact Morita同値を構成する．
さらに定理~\ref{thm:qa-relative-all-vertex-h2-open} は，中間二頂点も同じ相対heart内で認識し，
命題~\ref{prop:qa-relative-cubic-from-avatar-window} は二つの適合開で実際の全次数cubic族を回復する．
別に定義された全相対square AS $\qgr$ heartとの比較は，なお個別に確認する．

% !TEX root = ../quiver-avatar-moduli.tex
% 2026-11-21 / v1.0
\section{矢の本数、安定化表示、および開部分スタック}
\label{sec:qa-presentations}

本節では、有限 foundation の最小矢、relation-dual の戻り矢を加えた
箙、重み付き箙の非最小表示を区別する。
この区別のもとで、八本矢のモジュライが十本矢のモジュライの開部分となる
という主張を、任意の底上で正確に定式化する。
基礎体を $\kk$ とし、積 $ba$ は $a$ の後に $b$ を合成する規約とする。
本節の代数的な消去には標数の制限を要しない。

\subsection{直積型 collection の最小表示}

\begin{proposition}[直積型 collection の矢と関係]
\label{prop:qa-product-count}
$X=\PP^1\times\PP^1$ とし、
\[
 (E_0,E_1,E_2,E_3)
   =(\OO,\OO(1,0),\OO(0,1),\OO(1,1))
\]
と置く。有限代数 $B=\End_X(\bigoplus_{i=0}^3E_i)$ の最小箙は
\[
 a_i:0\longrightarrow1,\quad b_j:0\longrightarrow2,\quad
 u_j:1\longrightarrow3,\quad v_i:2\longrightarrow3
 \qquad(i,j=0,1)
\]
の八本の矢を持ち、最小関係は
\begin{equation}
\label{eq:qa-square-relations}
 r_{ij}=u_ja_i-v_ib_j\qquad(i,j=0,1)
\end{equation}
の四本である。従って relation-dual の戻り矢
$c_{ij}:3\to0$ を加えた relation QP は十二本矢であり、potential は
\begin{equation}
\label{eq:qa-square-potential}
 W_{\mathrm{sq}}=\sum_{i,j=0}^1c_{ij}(u_ja_i-v_ib_j)
\end{equation}
である。これは道長に関して reduced である。
\end{proposition}

\begin{proof}
$V=H^0(X,\OO(1,0))$、$U=H^0(X,\OO(0,1))$ はともに二次元である。
非恒等 Hom のうち、$\Hom(E_0,E_3)$ 以外の非零成分は
$V,U,U,V$ の四つである。
$\Hom(E_1,E_2)=H^0(X,\OO(-1,1))=0$ であり、
$\Hom(E_0,E_3)=V\otimes U$ への二つの道の合成はともに同型である。
従って二段階の道空間 $(U\otimes V)\oplus(V\otimes U)$ から
$V\otimes U$ への核は四次元で、\eqref{eq:qa-square-relations} が基底となる。
それより長い道は存在せず、これらが関係の全てである。
関係に一次部分はないので八本の矢は最小であり、
関係を他の矢で左右から乗じて新たな関係を得る余地もない。
四本の双対矢を加えると十二本となる。
\eqref{eq:qa-square-potential} は道長三なので、その二次部分は零である。
\end{proof}

\begin{remark}[十二本という数だけでは箙を指定できない]
\label{rem:qa-two-twelve-quivers}
命題~\ref{prop:qa-product-count} の箙を $Q_{\mathrm{sq}}$ と書く。
後述する $Q_{12}$ は四辺に各二本の矢を持つ有向四角形と四本の対角矢からなる。
$Q_{\mathrm{sq}}$ は $3\to0$ に四本の矢を持ち、$1\to2$ の矢を持たない。
従って両者は同じ十二本矢の箙ではない。
また、直積型 collection 自身の relation QP を十本矢と数えることはできない。
十本矢との対応を主張する際には、mutation と collection の変更を明示する必要がある。
この数え方は \cite[Remark 38.2]{AvatarVSeven2026} と一致する。
\end{remark}

\subsection{二つの可縮 pair を持つ共通表示}

四頂点を $\mathbb Z/4\mathbb Z$ で標識し、$Q_8$ を
各 $i\to i+1$ に二本の矢を持つ有向四角形とする。
その矢の重みを一とする。
$Q_{10}^{\mathrm{odd}}$ はこれに
\[
 x:1\longrightarrow3,\qquad y:3\longrightarrow1
\]
を加えた箙で、追加矢の重みは二とする。
$Q_{10}^{\mathrm{even}}$ は代わりに
\[
 u:0\longrightarrow2,\qquad v:2\longrightarrow0
\]
を加えたもの、$Q_{12}$ は両方の pair を加えたものとする。
以下の potential は重み四の巡回同値類であり、potential の定数倍は
同一視しない。

\begin{lemma}[共通 potential の形]
\label{lem:qa-universal-potential}
$Q_{12}$ の重み四の potential は一意に
\begin{equation}
\label{eq:qa-universal-potential}
 W=H+yq+px+vq'+p'u+\lambda yx+\mu vu
\end{equation}
と書ける。ここで $H$ は $Q_8$ の道長四の potential、
$q:1\to3$、$p:3\to1$、$q':0\to2$、$p':2\to0$ は
$Q_8$ の道長二の線形結合であり、$\lambda,\mu$ は係数である。
従って基底を固定した potential 空間は
\[
 P_8\simeq\mathbb A^{16},\qquad
 P_{10}^{\mathrm{odd}}\simeq P_{10}^{\mathrm{even}}\simeq\mathbb A^{25},
 \qquad P_{12}\simeq\mathbb A^{34}
\]
である。これらの記述は任意の $\kk$-scheme 上で成立する。
\end{lemma}

\begin{proof}
重み四の cycle が重み二の矢を持たなければ、四辺を一周するので
$2^4=16$ 個の道が巡回同値類の基底を与える。
重み二の矢を一つ持つ場合は、残りの二本が $Q_8$ の道長二の道となり、
四本の対角矢それぞれについて四個の可能性がある。
重み二の矢を二つ持つ cycle は $yx$ と $vu$ の二つだけである。
この三種類は巡回回転で混ざらないので、直和分解と次元の式を得る。
同じ道の基底を任意の可換係数環上で用いれば相対的な主張も従う。
\end{proof}

\begin{proposition}[任意の底上の有限な消去]
\label{prop:qa-relative-elimination}
$T$ 上で \eqref{eq:qa-universal-potential} が与えられているとする。
$\lambda\in\Gamma(T,\OO_T)^\times$ なら、
\[
 x'=x+\lambda^{-1}q,\qquad y'=y+\lambda^{-1}p
\]
によって
\begin{equation}
\label{eq:qa-one-pair-elimination}
 W=\lambda y'x'+H-\lambda^{-1}pq+vq'+p'u+\mu vu
\end{equation}
となる。従って通常の重み付き道代数で定義した Jacobian 代数について
\[
 \Jac(Q_{12},W)
 \simeq
 \Jac(Q_{10}^{\mathrm{even}},H-\lambda^{-1}pq+vq'+p'u+\mu vu)
\]
が成り立つ。$\mu$ が可逆な場合も同様で、両方が可逆なら
\begin{equation}
\label{eq:qa-two-pair-elimination}
 \Jac(Q_{12},W)
 \simeq\Jac(Q_8,H-\lambda^{-1}pq-\mu^{-1}p'q').
\end{equation}
これらの構成は任意の基底変換と両立する。
\end{proposition}

\begin{proof}
$p,q$ は $x,y,u,v$ を含まないので、表示した変換は逆も多項式である
重み付き道代数の自己同型である。式は積を展開して確認できる。
変換後の cyclic derivative は $\lambda x'$ と $\lambda y'$ を含むため、
Jacobian 商では $x'=y'=0$ である。
残りの cyclic derivative は残余 potential のものに一致する。
二つの pair の変換は互いの矢を変更せず、順序を問わず実行できる。
用いた演算は係数の加減乗除と、表示した可逆元の逆元だけである。
従って同型は任意の基底変換後にも同じ式で成立する。
\end{proof}

\begin{remark}[一般の splitting theorem との関係]
一般の QP の splitting theorem は完備道代数上の
\cite[Theorem 4.6]{DWZ2008} による。
命題~\ref{prop:qa-relative-elimination} では混合項が追加矢を含まないため、
有限回の三角変換だけで分離できる。従って完備化も極限操作も不要である。
この相対計算は AS 正則性や各次数成分の平坦性を主張していない。
それらは別の仮定または定理を必要とする。
\end{remark}

幾何点上の最小矢数は次の四つの場合に分かれる。
表の $\lambda=0$、$\mu=0$ は、その幾何点の剰余体での等式である。
\begin{center}
\begin{tabular}{c|c|c}
係数 & 残る最小対角矢 & 最小矢数 \\
\hline
$\lambda\ne0,\ \mu\ne0$ & なし & $8$ \\
$\lambda=0,\ \mu\ne0$ & $x,y$ & $10$（odd）\\
$\lambda\ne0,\ \mu=0$ & $u,v$ & $10$（even）\\
$\lambda=0,\ \mu=0$ & $x,y,u,v$ & $12$
\end{tabular}
\end{center}
消去後の potential は道長三以上なので、Jacobian ideal に一次関係はない。
これが表の最小性の根拠である。
非被約な底上では、全幾何点で係数が零でも係数自身が零とは限らない。
例えば $\lambda=\epsilon$ on $\Spec\kk[\epsilon]/(\epsilon^2)$ は可逆でなく、
最後の pair を消去できない。

\subsection{通常の QP スタックに残る慣性}

ここでは頂点を固定する全ての重み付き道代数自己同型を射とする。
$G_8=(\GL_2)^4$ と置くと、重み一の矢の像は基底変更で尽くされる。
十本矢では、さらに各追加矢の非零スカラー倍と、
同じ始点・終点を持つ道長二の道の付加が許される。従って
\[
 G_{10}=(G_8\times\Gm^2)\ltimes\mathbb G_a^8,
 \qquad
 G_{12}=(G_8\times\Gm^4)\ltimes\mathbb G_a^{16}.
\]
群の作用は potential の変数変換であり、
$\mathcal P_r=[P_r/G_r]$ を通常の重み付き QP 商スタックとする。
この節では頂点スカラーの rigidification は行わない。

\begin{proposition}[可縮 pair の残余慣性]
\label{prop:qa-stabilized-inertia}
上の群と頂点標識を用いると、
\begin{align}
 [\{\lambda\ne0\}\subset P_{10}^{\mathrm{odd}}/G_{10}]
   &\simeq \mathcal P_8\times B\Gm,\label{eq:qa-qp-gerbe-one}\\
 [\{\lambda\mu\ne0\}\subset P_{12}/G_{12}]
   &\simeq \mathcal P_8\times B(\Gm^2).
   \label{eq:qa-qp-gerbe-two}
\end{align}
特に、消去によってこれらの通常の QP スタックを
$\mathcal P_8$ と同定することはできない。
\end{proposition}

\begin{proof}
十本矢の場合、命題~\ref{prop:qa-relative-elimination} により
$p=q=0$ とでき、さらに一方の矢のスケールを変えて $\lambda=1$ とできる。
これは係数環上の可逆元の逆を用いる代数的な正規化である。
正規化した切片は $H+yx$ の形の potential の空間 $P_8$ である。

この切片の二点 $H+yx$ と $\widetilde H+yx$ を結ぶ自己同型を調べる。
重み一部分の作用を $g\in G_8$、対角矢の作用を
$x\mapsto\alpha x+q_0$、$y\mapsto\beta y+p_0$ と書く。
変換後に $yq_0$ と $p_0x$ の混合項が消えるには、
$\alpha,\beta$ が可逆なので $p_0=q_0=0$ が必要である。
$yx$ の係数から $\alpha\beta=1$ を得る。
従って切片間の全ての射は
\[
 (g,t):\quad H\longmapsto gH,\quad
 x\longmapsto tx,\quad y\longmapsto t^{-1}y
\]
であり、その群は $G_8\times\Gm$ である。
後半の $\Gm$ は $P_8$ に自明に作用する。
各対象が fppf 局所にこの切片へ移り、切片間の射がちょうど上の群で
記述されるので、商スタックの同値~\eqref{eq:qa-qp-gerbe-one} を得る。
二つの pair について同じ正規化を独立に行えば
\eqref{eq:qa-qp-gerbe-two} を得る。
ここでは切片とその安定化群を明示したので、単に gerbe と述べるだけでなく
直積との同値まで得られている。
\end{proof}

\begin{remark}[表示を忘れる操作]
消去される pair の変換 $x\mapsto tx$、$y\mapsto t^{-1}y$ は
Jacobian 商では恒等である。
従って通常の QP groupoid から代数 groupoid への忘却は忠実でない。
同じ障害は \cite[Proposition 52.3]{AvatarVSeven2026} で指摘されている。
一般に関係式の表示の自己同型と商代数の自己同型を区別すべきことは
\cite[Proposition 3.3]{AOU2014} にも現れる。
有限有向箙の関係式スタックの有限型性
\cite[Definition 3.5, Proposition 3.6]{AOU2014} は有用であるが、
その定理だけから、ここで扱う無限次元 graded 代数の全次数平坦な族や
内在的モジュライの代数性は従わない。
\end{remark}

\subsection{代数の射を用いた正確な開埋込み}

\begin{definition}[代数の fppf スタック]
\label{def:qa-algebra-stack}
$\mathfrak M_{12}^{\mathrm{alg}}$ の $T$ 上の対象は、
$A_0=\OO_T^4$ と四つの標識付き冪等元を持つ非負次数付き代数
$A=\bigoplus_{d\ge0}A_d$ で、各 $A_d$ が有限階数局所自由、かつ
fppf 局所に $\Jac(Q_{12},W)$ の形の表示を持つものとする。
射は頂点と次数を保つ全ての代数同型とし、potential、矢の基底、
表示からの全射はデータに含めない。
$\mathfrak M_{10}^{\mathrm{odd,alg}}$、
$\mathfrak M_{10}^{\mathrm{even,alg}}$、$\mathfrak M_8^{\mathrm{alg}}$
もそれぞれの箙を用いて同様に定義する。
これらは fppf 降下を持つ代数のスタックの充満部分スタックである。
この定義に AS 正則性は含めない。
\end{definition}

各対象について、添字を $\mathbb Z/4\mathbb Z$ で読み、有限表示層
\begin{equation}
\label{eq:qa-indecomposables}
 C_i(A)=\operatorname{coker}\left(
 e_{i+2}A_1e_{i+1}\otimes_{\OO_T}e_{i+1}A_1e_i
 \longrightarrow e_{i+2}A_2e_i\right)
\end{equation}
を定める。これは重み二における新しい生成元の層であり、
定義は表示に依存しない。

\begin{theorem}[矢の少ない代数スタックの開性]
\label{thm:qa-open-algebra-stacks}
可縮 pair を追加することによって、充満な開埋込み
\begin{equation}
\label{eq:qa-open-diagram}
 \begin{aligned}
 \mathfrak M_8^{\mathrm{alg}}&\lhook\joinrel\longrightarrow
 \mathfrak M_{10}^{\mathrm{odd,alg}}
 \lhook\joinrel\longrightarrow\mathfrak M_{12}^{\mathrm{alg}},\\
 \mathfrak M_8^{\mathrm{alg}}&\lhook\joinrel\longrightarrow
 \mathfrak M_{10}^{\mathrm{even,alg}}
 \lhook\joinrel\longrightarrow\mathfrak M_{12}^{\mathrm{alg}}.
 \end{aligned}
\end{equation}
を得る。
両十本矢開部分の交わりは $\mathfrak M_8^{\mathrm{alg}}$ である。
正確には、それぞれの像は
\begin{align*}
 \mathfrak M_{10}^{\mathrm{odd,alg}}&=\{C_0=C_2=0\},\\
 \mathfrak M_{10}^{\mathrm{even,alg}}&=\{C_1=C_3=0\},\\
 \mathfrak M_8^{\mathrm{alg}}&=\{C_0=C_1=C_2=C_3=0\}
\end{align*}
である。開埋込みという主張は任意の scheme からの基底変換について成立し、
各スタックの代数性を仮定しない。
\end{theorem}

\begin{proof}
十本矢の表示に $vu$ を追加しても Jacobian 代数は変わらず、
八本矢からは $yx$ と $vu$ を追加できる。
従って同じ代数を対象として送る関手が存在する。
射は両辺とも全ての次数付き頂点固定代数同型なので、関手は充満忠実である。

局所表示~\eqref{eq:qa-universal-potential} で
\eqref{eq:qa-indecomposables} を計算する。
重み二の関係は四本の重み二の矢による cyclic derivative だけである。
道長二の積を割ると、これらの関係は
$\lambda x=\lambda y=\mu u=\mu v=0$ となる。従って局所的に
\[
 C_1\simeq C_3\simeq\OO_T/(\lambda),\qquad
 C_0\simeq C_2\simeq\OO_T/(\mu).
\]
特に $C_0=C_2=0$ と $\mu$ の可逆性は同値であり、
命題~\ref{prop:qa-relative-elimination} により
$Q_{10}^{\mathrm{odd}}$ の表示を得る。他の二つの場合も同様である。
逆に小さい箙の表示では、対応する生成元層は明らかに零である。
これにより像の表示を得る。

有限表示層 $C_i$ の消滅 locus は、その support の補集合という開集合である。
また cokernel は任意の基底変換と両立する。
従って任意の $T\to\mathfrak M_{12}^{\mathrm{alg}}$ による各関手の
基底変換は、対応する $C_i$ の消滅で定まる $T$ の開部分そのものである。
これは開埋込みの主張と交わりの式を同時に証明する。
\end{proof}

\begin{remark}[最小型、相対型、および追加条件]
\label{rem:qa-open-limits}
十本矢の開部分に八本矢の対象が入るのは、
定義~\ref{def:qa-algebra-stack} が非最小の十本矢表示も許すためである。
幾何点で最小矢がちょうど十本の対象だけを取れば、八本矢の対象は含まれない。
相対的に固定した最小型を扱うときは、残る $C_i$ が階数一の局所自由層で
あることを課す。この条件は局所表示では対応する mass 係数が底上で
零となることを要求し、全幾何点で零という条件より強い。

本節の開性は、同型と基底変換に不変な追加条件を課した共通の
部分スタックへ制限しても保たれる。
ただし AS 条件を満たす部分がどの部分であるか、その部分の代数性、
avatar からの全域被覆、さらに構造対象を保つ内在的圏への射は
この定理の結論に含まれない。
また $\mathfrak M_r^{\mathrm{alg}}$ では四頂点と周期的な graded 表示を
保持している。これを忘れた $\mathcal O$-marked の exact 同値を、
次数付き代数同型と同一視してはならない。
\end{remark}

% !TEX root = ../quiver-avatar-moduli.tex
% 2026-11-21
\section{有限 foundation の二つの十本矢開部分}
\label{sec:qa-finite-foundations}

本節では、矢の本数による開部分を実際の代数スタックとして構成する。
対象は四頂点の有限 directed algebra であり、AS 正則性や標準 heart は
この構成の条件に含めない。最後に数える矢は、有限代数そのものの矢と
最小関係に双対な戻り矢を合わせた relation completion の矢である。
従って、ここで得る開部分の定理と、AS 代数・avatar の相対分類定理とは
異なる主張である。

\subsection{全ての積を保持する有限表示}

\(U,V,W\) を二次元、\(E,F\) を四次元、\(N\) を六次元の
\(\kk\)-ベクトル空間とし、頂点成分を
\[
 A_{01}=U,\quad A_{12}=V,\quad A_{23}=W,\qquad
 A_{02}=E,\quad A_{13}=F,\quad A_{03}=N
\]
と固定する。対角成分は \(A_{ii}=\kk\)、逆向き成分は零とする。
非自明な積は四写像
\[
 \begin{aligned}
 a&:V\otimes U\longrightarrow E,&
 b&:W\otimes V\longrightarrow F,\\
 c&:W\otimes E\longrightarrow N,&
 d&:F\otimes U\longrightarrow N
 \end{aligned}
\]
であり、結合律はただ一つの写像等式
\begin{equation}
 c(1_W\otimes a)=d(b\otimes1_U)
 \quad\text{in }\Hom(W\otimes V\otimes U,N)
 \label{eq:qa-associativity}
\end{equation}
である。この等式の定める閉アフィンスキームを \(Z\) と書く。
\(Z\subset\mathbb A^{128}\) は有限型である。
全ての積を保持するため、後に \(a\) または \(b\) の階数が落ちても
有限代数のベクトル空間としての階数二十四は変わらない。

\begin{definition}[foundation の許容開部分]
\label{def:qa-foundation-open}
\(Z\) の開部分 \(Z^{\mathrm{nd}}\) を次の条件で定める。
\begin{enumerate}
\item \((c,d): (W\otimes E)\oplus(F\otimes U)\to N\) は全射である。
\item 全ての幾何ファイバーで、\(\ker a\) と \(\ker b\) は
非零の分解可能テンソルを含まない。
\end{enumerate}
ここで \(Z^{\mathrm{nd}}\) の記号は上記の有限条件だけを意味し、
AS 正則性を意味しない。
\end{definition}

\begin{lemma}
\label{lem:qa-foundation-open-ranks}
定義~\ref{def:qa-foundation-open} は基底変更と両立する開条件である。
各幾何点で \(\rank a,\rank b\in\{3,4\}\) である。
\end{lemma}
\begin{proof}
全射性は最大 minor の非消滅条件である。
\(a\) について、\(Z\times\PP(V)\times\PP(U)\) 上で
\(a(v\otimes u)=0\) となる閉 incidence を取る。
射影空間の積からの射影は proper なので、その \(Z\) への像は閉である。
その補集合が第二条件を表す。\(b\) についても同じである。

二次元空間のテンソル積の射影化 \(\PP^3\) において、
分解可能テンソルは Segre 二次曲面をなす。
代数閉体上の射影直線はこの二次曲面と交わる。
従って二次元以上の核は非零の分解可能テンソルを含むので、
許される核の次元は零または一である。
\end{proof}

\[
 G=\GL(U)\times\GL(V)\times\GL(W)
       \times\GL(E)\times\GL(F)\times\GL(N)
\]
を各空間の基底変更で作用させる。
頂点の対角単元 \((\lambda_0,\lambda_1,\lambda_2,\lambda_3)\)
が \(A_{ij}\) 上で \(\lambda_j\lambda_i^{-1}\) 倍する作用の像を
\(K\subset G\) と書く。
共通スカラーを除けば \(K\simeq\Gm^3\) であり、
\(K\) は \(Z\) 上自明に作用する中心部分群である。

\begin{proposition}[有限代数の代数スタック]
\label{prop:qa-foundation-stack}
\[
 \mathfrak F^{\mathrm{lin}}_{\le12}=[Z^{\mathrm{nd}}/G],
 \qquad
 \mathfrak F_{\le12}=[Z^{\mathrm{nd}}/(G/K)]
\]
は有限型 Artin スタックである。
前者は上記の階数の頂点成分と積を持つ有限代数のスタックであり、
後者は頂点ごとの比較同型のスカラーによる中心 \(K\) を剛体化したものである。
この命題ではスタックの滑らかさも AS locus の開性も主張しない。
\end{proposition}
\begin{proof}
各頂点成分を局所自由層とする同じデータは、fppf 局所に基底を選ぶと
\(Z^{\mathrm{nd}}\) の点で表される。
代数の同型はこれらの空間の基底変更と正確に一致する。
従って前者の商は指定した moduli を表す。
\(G\) と \(G/K\) は滑らかな有限型代数群なので、
両商は有限型 Artin スタックである。
\(K\) は各対象に中心自己同型として含まれるため、
\([Z^{\mathrm{nd}}/G]\to[Z^{\mathrm{nd}}/(G/K)]\) はその剛体化である。
\end{proof}

\subsection{最小生成元と最小関係の数}

\begin{proposition}[八・十・十二本の有限代数的判定]
\label{prop:qa-foundation-arrow-count}
\(Z^{\mathrm{nd}}\) の幾何点に対応する代数を \(A\) とし、
\[
 \delta_a=4-\rank a,\qquad \delta_b=4-\rank b
 \quad\in\{0,1\}
\]
と置く。\(A\) の最小 directed quiver は
\(6+\delta_a+\delta_b\) 本の矢を持ち、最小関係は
\(2+\delta_a+\delta_b\) 本である。
従って、最小関係ごとに一本の戻り矢を加えた relation completion は
\begin{equation}
 8+2(\delta_a+\delta_b)
 \label{eq:qa-foundation-arrow-count}
\end{equation}
本の矢を持つ。
\end{proposition}
\begin{proof}
隣接頂点間には各二本、合計六本の矢がある。
\(E/\operatorname{im}a\) と \(F/\operatorname{im}b\) の基底を取り、
それぞれ \(0\to2\)、\(1\to3\) の追加生成矢とする。
\((c,d)\) の全射性により、\(0\to3\) の全ての元は
\(0\to2\to3\) または \(0\to1\to3\) を経由する積で生成されるので、
\(0\to3\) の追加生成矢はない。
従って生成矢の数は \(6+\delta_a+\delta_b\) である。

\(0\to2\)、\(1\to3\) の最小関係は、それぞれ
\(K_a=\ker a\)、\(K_b=\ker b\) であり、その次元は
\(\delta_a,\delta_b\) である。
これらの関係を課す前の \(0\to3\) の道の空間は
\[
 (W\otimes V\otimes U)
 \oplus\bigl(W\otimes(E/\operatorname{im}a)\bigr)
 \oplus\bigl((F/\operatorname{im}b)\otimes U\bigr)
\]
という、次元 \(8+2\delta_a+2\delta_b\) の空間である。
短い関係の生成する部分空間は最初の直和成分内の
\[
 (W\otimes K_a)+(K_b\otimes U)
\]
であり、この和は直和である。
実際、両核が一次元で共通の非零テンソルがあれば、
それは \(w\otimes\kappa_a=\kappa_b\otimes u\) と書ける。
\(w\) に対して非零となる線形形式で縮約すると、
\(\kappa_a\) が分解可能になり、定義に反する。
一方の核が零の場合は明らかである。

従って短い関係を課した後の \(0\to3\) 成分は次元八である。
結合律~\eqref{eq:qa-associativity} により \(N\) への写像が定まり、
これは \((c,d)\) の全射性によって全射である。
残る核は二次元であり、ちょうど二本の新たな最小関係を与える。
添字次数は三であるが、追加矢があるときは道長二と三の項を含み得る。
既に短い関係の生成する部分で割っているため、これらは最小関係であり、
四頂点の有向性から、それ以外の始点・終点に新たな関係はない。
これで関係数と \eqref{eq:qa-foundation-arrow-count} を得る。
\end{proof}

添字の差による重みを有限代数の矢 \(i\to j\) に与え、
同じ始点・終点を持つ関係の戻り矢には重み \(4-(j-i)\) を与える。
関係 potential の各項は重み四となる。
隣接矢六本と \(0\to3\) の二つの関係の戻り矢は、
四周期の各隣接頂点を結ぶ二本ずつの重み一の矢を与える。
\(\delta_a=1\) なら \(0\to2\)、\(2\to0\) の重み二の矢を一組、
\(\delta_b=1\) なら \(1\to3\)、\(3\to1\) の重み二の矢を一組加える。
従って重み付きの箙自体は、八本矢、二種類の十本矢、十二本矢の
それぞれの配置に正確に一致する。
この配置の一致は、potential の復元や AS 正則性、
QP の同型と有限代数の同型の比較まで含む主張ではない。

\subsection{二つの開部分とその共通部分}

\begin{theorem}[矢数を上限とする開部分]
\label{thm:qa-foundation-open-diagram}
\(Z^{\mathrm{nd}}\) 上で \(a\)、\(b\) が可逆となる開部分を
それぞれ \(Z_a,Z_b\) と書く。
\[
 \begin{aligned}
 \mathfrak F_{\le10}^{(a)}&=[Z_b/(G/K)],&
 \mathfrak F_{\le10}^{(b)}&=[Z_a/(G/K)],\\
 \mathfrak F_8&=[(Z_a\cap Z_b)/(G/K)]
 \end{aligned}
\]
と置く。上付き \((a)\)、\((b)\) は、階数低下を許す写像を示す。
これらは \(\mathfrak F_{\le12}\) の開部分スタックであり、
\begin{equation}
 \mathfrak F_8
 =\mathfrak F_{\le10}^{(a)}
       \cap\mathfrak F_{\le10}^{(b)}.
 \label{eq:qa-foundation-chart-intersection}
\end{equation}
両十本矢開部分の合併は、最小 relation completion の矢数が
高々十である locus である。
その補集合は、幾何点上では \(\rank a=\rank b=3\) の
最小十二本矢 locus である。
\end{theorem}
\begin{proof}
可逆性は行列式の非消滅であり、\(G/K\) の作用で不変である。
不変開部分からの商の射は representable な open immersion である。
交わりと合併は同じ基底上で計算できる。
矢数との対応は命題~\ref{prop:qa-foundation-arrow-count} による。
\end{proof}

\begin{remark}[最小本数の層と開部分の区別]
\label{rem:qa-exact-arrow-strata}
例えば \(\mathfrak F_{\le10}^{(a)}\) は八本矢の点も含む。
最小十本矢の点だけを取った層は
\(\mathfrak F_{\le10}^{(a)}\setminus\mathfrak F_8\) であり、
八本矢の moduli を開部分として含むのはこの層ではない。
非被約基底について「階数三の層」と呼ぶときは、単なる幾何点条件でなく、
対応する大きな minor の消滅と三次 minor の可逆性を課した
scheme-theoretic な階数 locus を用いる。
本節の開部分は可逆性で定義されており、この区別による曖昧さを持たない。
また、選んだ最小関係とその戻り矢から作る QP の自己同型群を、
有限代数の自己同型群と同一視したわけではない。
QP から有限代数への回復と射の比較には別の証明が必要である。
\end{remark}

\subsection{cubic表示の幾何条件に対する次数六以下の明示的検査}
\label{subsec:qa-explicit-geometric-fourfold-test}

既存のnormal foundation chartで最初の積が同型となる場合、
論文B~\cite[\texttt{prop:cubic-foundation-open}]{QuadricModuli2026}
は二次元空間$V_0,V_1,V_2,V_3$の
四重テンソル$w\in V_0\otimes\cdots\otimes V_3$の
geometric性を使ってcubic逆構成を行う。
以下はこのgeometric開を
\emph{完全に明示的な次数四・六の多項式条件}に書き直す。

\begin{lemma}[四重テンソルのgeometric性のadjugate判定]
\label{lem:qa-explicit-cubic-geometric-adjugate}
$k$を代数閉体とし、各$V_i$を二次元として、
$w\in V_0\otimes V_1\otimes V_2\otimes V_3$
を取る。$i\in\mathbb Z/4\mathbb Z$について、
隣接二因子の縮約行列を
\[
 M_i:(V_i\otimes V_{i+1})^\vee
       \longrightarrow V_{i+2}\otimes V_{i+3}
\]
とする。基底を選んで$M_i$を$4\times4$行列とみなす。
$A_i=\operatorname{adj}(M_i)$をその古典的adjugateとし、
その第$j$列を二因子の$2\times2$テンソルに
reshapeした行列を$C_{ij}$（$1\le j\le4$）とする。
\[
 d_i=\det M_i,\qquad
 h_{ij}=\det C_{ij}
 \tag{Geo-poly}\label{eq:qa-cubic-geometric-minors}
\]
と置く。$d_i$は$w$の成分の次数四、
$h_{ij}$は次数六の多項式である。

Van den Berghのgeometric条件
\[
 (\phi_i\otimes\phi_{i+1})\mathbin{\lrcorner}w\ne0
 \quad\text{for all }i,\quad
 0\ne\phi_i\in V_i^\vee,\ \
 0\ne\phi_{i+1}\in V_{i+1}^\vee
\]
と次の有限条件は同値である。
\[
 \boxed{\quad
  \text{各$i=0,1,2,3$について、}
  \quad d_i\ne0
  \ \text{または}\ 
  h_{i1},h_{i2},h_{i3},h_{i4}
  \text{の少なくとも一つが非零}.\quad}
 \tag{Geo-test}\label{eq:qa-cubic-geometric-finite-test}
\]
従ってgeometric四重テンソル開は、
\emph{四つの次数四式と十六の次数六式だけ}
を使って定義される。
各$M_i$の行・列の基底変更に対し
この非消滅条件は不変なので、
有限foundationの相対chartへ降下する。

これは前段の一般存在定理
~\ref{cor:qa-realizable-foundation-finite-open-criterion}
で未計算だった\emph{全実現可能開}の方程式ではなく、
$F$が可逆であるcubic部分における
geometric条件の具体的な判定式である。
\end{lemma}
\begin{proof}
固定した$i$について、$K_i=\ker M_i$を
四次元空間$(V_i\otimes V_{i+1})^\vee$の部分空間とする。
geometric条件は、$K_i$に非零の
\emph{分解可能}二因子テンソルがないことと同値。
これは
$\PP(K_i)\cap\operatorname{Seg}(\PP^1\times\PP^1)=\varnothing$
という条件である。
Segre像は$\PP^3$の滑らかな二次曲面で、
分解可能な$2\times2$行列は行列式零によって特徴付けられる。

$\operatorname{rank}M_i=4$なら$K_i=0$なので条件は成立し、
$d_i\ne0$と同値である。
$\operatorname{rank}M_i\le2$なら
$\dim K_i\ge2$なので$\PP(K_i)$は少なくとも直線を含む。
代数閉体上で任意の射影直線は二次曲面と交わる
（線上で行列式を制限した二次形式に零点がある、
恒等零の場合も交わる）。
従ってgeometric条件は失敗する。
このとき全ての$3\times3$小行列式が零なので
$A_i=0$となり、全ての$h_{ij}$は零である。

残る$\operatorname{rank}M_i=3$の場合、
$A_i$は階数一の非零行列で、像が
一次元の$K_i$と一致する。
その非零列は$K_i$を張るので、
$K_i$が分解可能テンソルを含まないことは、
いずれかの列をreshapeした$2\times2$行列の
行列式が非零となることと同値である。
これがいずれかの$h_{ij}\ne0$という条件になる。

基底変更で$K_i$は
$(\operatorname{GL}(V_i)\otimes\operatorname{GL}(V_{i+1}))^\vee$
により移り、Segreの分解可能座も不変だから、
判定条件は基底によらない。
四つの$i$について適用すれば結論を得る。
\end{proof}

\subsection{既存の normal foundation chart の位置}

\begin{proposition}[一方の積を可逆にした表示]
\label{prop:qa-existing-foundation-chart}
\(b\) が可逆な開部分では \(F\simeq W\otimes V\) と同定でき、
結合律は \(d\) を \(a,c\) から一意に決める。
従って、この開部分の線形データは既存の normal foundation 表示
\[
 F_{\mathrm{old}}:P\otimes W_1\to W_2,
 \qquad G_{\mathrm{old}}:Q\otimes W_2\to W_3
\]
において
\((P,W_1,W_2,Q,W_3)=(V,U,E,W,N)\)、
\((F_{\mathrm{old}},G_{\mathrm{old}})=(a,c)\) としたものに一致する。
この同定の下で \((c,d)\) の全射性は \(c\) の全射性に等しい。
既存の表示の全開部分と比較する際には、
本節の分解可能な核を持たない条件を追加する。
\end{proposition}
\begin{proof}
\eqref{eq:qa-associativity} の右辺で \(b\otimes1_U\) が同型なので、
\(d=c(1_W\otimes a)(b^{-1}\otimes1_U)\) である。
残る空間と二写像は表示したものに一致する。
\(d\) は \(c\) を経由するので、両者の像の和は \(c\) の像である。
\end{proof}

この比較により、\cite[\texttt{prop:foundation-stack}]{QuadricModuli2026} に現れる一方の積を恒等化した chart と、
本節の二つの積がともに階数三になる境界とを区別できる。
古い二写像表示で \(\rank F_{\mathrm{old}}=2\) として得られる十二本矢表示は、
本節の許容開部分には入らない。
その核は二次元であり分解可能テンソルを含むからである。
また、その十二本矢は一方の対角位置に二組の追加矢を持つのに対し、
本節の最小十二本矢は両対角位置に一組ずつの追加矢を持つ。
矢の総数だけから、これらを同じ型と扱ってはならない。

% !TeX root = ../nc-quadric-avatar.tex
% 2026-11-23 / v4.25。可逆cubicとmass壁の相対QP群oid。

\subsection{八矢QPから有限cubic foundationへの全相対群oid}
\label{subsec:qa-cubic-qp-foundation-groupoid}

本項では前記の四頂点とweightを保持し，標数零の体$k$上で扱う．
以下の有限Grassmann商の同型自体は任意の基礎体上で成立する．
八矢箙$Q_8$の隣接矢束を
\[
 U:0\to1,\qquad V:1\to2,\qquad W:2\to3,\qquad T:3\to0
\]
（いずれも階数二）とし，
\(\mathcal V=W\otimes V\otimes U\) と置く．
頂点で始点を固定した巡回道を用いれば，
weight四のpotentialは正確に
\[
 H\in\mathcal V\otimes T,\qquad
 \partial_T H:\ T^\vee\longrightarrow\mathcal V
 \tag{QP8}\label{eq:qa-cubic-qp-flattening}
\]
に対応する．ここでは巡回同値を取っても余分な係数商は生じない．
任意のスキーム上で \(\partial_T H\) が階数二の局所自由直和因子
である開を \(P_8^{(2)}\) と書く．
これはrank-two flatteningの全minorの単位イデアル条件である．
\[
 G_8=\GL(U)\times\GL(V)\times\GL(W)\times\GL(T),
 \qquad G_3=\GL(U)\times\GL(V)\times\GL(W)
\]
を重み付きright-equivalence群とする．$Q_8$では任意の
重み一の矢はその始点終点の二矢の線型結合であり，
同じ始点終点の重み一の長い道は存在しない．従って
重みを保つ右同値の全群は正に$G_8$である．
非重み付き完備道代数の全自己同型群と同一視しない．

\begin{theorem}[八矢QPと階数二関係束の相対groupoid同値]
\label{thm:qa-cubic-qp-grassmann-groupoid}
次のfppfスタックの自然な同値がある：
\[
 \boxed{\quad
 [P_8^{(2)}/G_8]\ \simeq\
 [\operatorname{Gr}(2,\mathcal V)/G_3]\quad}.
 \tag{QG}\label{eq:qa-cubic-qp-grassmann-stack}
\]
左辺はweight四の四次potentialと重み付きright-equivalenceを
すべての同型として持つ．右辺は三つの階数二束と
$\mathcal V$内の階数二直和因子$R$を同型とともに持つ．
同値は$H$を
$R=\operatorname{im}(\partial_T H)$へ送る．
任意のスキームへの基底変更，非被約基底，Isom-sheafおよび
慣性群スキームと整合する．
\end{theorem}
\begin{proof}
$G_3$をいったん忘れ，各矢束を固定する．
$\operatorname{Gr}(2,\mathcal V)$の普遍階数二部分束を
$\mathcal R\subset\mathcal V$とする．
$\partial_T H$の像が$\mathcal R$であるという条件は，
$T^\vee\xrightarrow{\sim}\mathcal R$という同型を選ぶことに等しい．
従って
\[
 P_8^{(2)}
 =\operatorname{Isom}_{\operatorname{Gr}(2,\mathcal V)}
       (T^\vee,\mathcal R)
\]
は普遍部分束の$\GL(T)$-torsorであり，Zariski局所に自明である．
$\GL(T)$は$P_8^{(2)}$上自由に作用し，
その幾何商およびfppf商は正確にGrassmannスキームとなる．
このtorsorは$G_3$-同変なので，残りの群による商を取れば
表示した二商スタックの同値が得られる．
具体的な逆関手では，$R$に対して$T=R^\vee$と置き，
$\mathcal V\otimes R^\vee$内の恒等テンソルをpotentialとする．
基底変更に伴う第四矢束の自動的な同型は全right-equivalenceを回収する．
単なる幾何点の全単射でなくtorsor同値なので，
相対同型層と慣性も回収する．
\end{proof}

\begin{proposition}[四頂点finite foundationの全積を保つ同型]
\label{prop:qa-cubic-qp-finite-foundation-groupoid}
有限代数スタック
\(\mathfrak F_8^{\rm lin}=[(Z_a\cap Z_b)/G]\)
（命題~\ref{thm:qa-foundation-open-diagram}のunrigidified版）
には自然な同値
\[
 \boxed{\quad
 \mathfrak F_8^{\rm lin}
 \ \simeq\
 [\operatorname{Gr}(2,\mathcal V)/G_3]
 \ \simeq\ [P_8^{(2)}/G_8]\quad}.
 \tag{FG}\label{eq:qa-cubic-qp-finite-groupoid}
\]
ここでは有限代数の階数二十四の全成分・全積を保存する．
QPからの関手は、最後の矢群$T$に沿うcyclic derivativesが張る
三段道の関係束を取り，四頂点のdirected windowの有限End代数を
回収する．Jacobian商の全次数AS正則性は主張しない．
\end{proposition}
\begin{proof}
有限foundationの三つの隣接Hom束を$U,V,W$とする．
$Z_a\cap Z_b$では二つの乗法
\[
 a:V\otimes U\xrightarrow{\sim}E,\qquad
 b:W\otimes V\xrightarrow{\sim}F
\]
が可逆である．
従って$E=V\otimes U$，$F=W\otimes V$へ基底変更で
厳密に規格化できる．
結合律$c(1_W\otimes a)=d(b\otimes1_U)$により
$W\otimes E$と$F\otimes U$の両経路の合成は，同一の写像
\[
 m:\mathcal V=W\otimes V\otimes U\longrightarrow N
\]
となる．元の$(c,d)$の合併全射性は$m$の全射性と同値である．
$\operatorname{rank}\mathcal V=8$，$\operatorname{rank}N=6$だから，
$R=\ker m$は階数二の直和因子である．逆に任意の$R$から
\[
 E=V\otimes U,\quad F=W\otimes V,\quad
 N=\mathcal V/R
\]
および二つの標準的な商写像$c,d$を取れば，全積と結合律を
満たす有限directed代数が得られる．この構成は互いに逆で，
$E,F,N$のbasis torsorも含めて任意の基底変更と可換する．
したがって有限代数の全isomorphism groupoidは
$[\operatorname{Gr}(2,\mathcal V)/G_3]$に等しい．
定理~\ref{thm:qa-cubic-qp-grassmann-groupoid}と合わせれば結論が従う．

さらに八矢QPの微分$\partial_T H$は二つの長さ三の関係を与える．
四頂点のdirected windowでは頂点$3\to0$の矢を通る道は存在しない．
従って他のcyclic derivativesに由来する余分な関係はなく，
最初の四頂点の有限Endの長さ三の商は正確に$\mathcal V/R$である．
二段道には関係がないので，有限代数の全積が一致する．
\end{proof}

\begin{corollary}[頂点スカラーの慣性とcubic $H_2$の相対認識]
\label{cor:qa-cubic-qp-pointed-h2-gerbe}
前命題の$G_3$の三つのscalar部分群
\(K_{\rm vert}\simeq\Gm^3\)はGrassmann上自明に作用する．
その剛体化は有限foundationの
\(\mathfrak F_8=[(Z_a\cap Z_b)/(G/K)]\) と一致し，
\[
 \mathfrak F_8\simeq
 [\operatorname{Gr}(2,\mathcal V)/\operatorname{PGL}_2^3].
 \tag{Fr}\label{eq:qa-cubic-qp-rigid-foundation}
\]
八矢側では，四つの矢群のscalarの積が一となる
\[
 K_{\rm QP}=
 \{(\alpha,\beta,\gamma,\delta)\in\Gm^4:
                 \alpha\beta\gamma\delta=1\}
 \simeq\Gm^3
\]
がpotentialを点ごとに固定する．
$G_8/K_{\rm QP}$への変更が正に上の剛体化を与え，
幾何点の安定化群にはfppf群スキームの完全列
\[
 1\longrightarrow K_{\rm QP}\longrightarrow
 \operatorname{Aut}_{[P_8^{(2)}/G_8]}(H)
 \longrightarrow\operatorname{Aut}_{\mathfrak F_8}(R)
 \longrightarrow1
 \tag{St}\label{eq:qa-cubic-qp-stabilizer-exact}
\]
がある．

四重テンソルのVan den Bergh geometric開$\mathscr G$を
指定すると、この同値はその不変開へ制限される．
幾何開では$\partial_T H$が階数二である．
既存の cubic 表示の相対同型
\(\mathfrak U_f\simeq[\mathscr G/G_f]\)
（命題~\ref{prop:qa-global-cubic-four-qubit-good-moduli}）
により，矢群のdualとcyclic順序を合わせて
\[
 \boxed{\quad
 [\mathscr G/G_8]\ \longrightarrow\
 [\mathscr G/(G_8/K_{\rm QP})]\simeq\mathfrak U_f
 \quad}
 \tag{HC}\label{eq:qa-cubic-qp-h2-cubic-gerbe}
\]
はbanded $\Gm^3$-gerbeである．
同値は実際のcubic表示の相対groupoidを含むが，
選んだ正規窓・rulingの忘却や任意AS族のheartとの同一視は
追加で検証を要する．
\end{corollary}
\begin{proof}
四頂点のscalar $(\lambda_0,\ldots,\lambda_3)$は
隣接Homに
$(\lambda_1/\lambda_0,\lambda_2/\lambda_1,
\lambda_3/\lambda_2)$で作用する．この三つの値は独立で，
三重tensorの関係部分空間には共通scalarとして作用するから
Grassmann上自明である．
第四矢群のscalarはその逆積となり，表示した
$K_{\rm QP}$に一致する．中心部分群の作用が自明な
商スタックの剛体化について，二つの表示を比べれば
群スキームの完全列が従う．
$K_{\rm vert}$の群作用を割った有効群は
$\operatorname{PGL}_2^3$である．

幾何条件は四重テンソルの不変開条件である．
もし第四矢群に関するflatteningの階数が一以下なら，
四因子テンソルを二因子対二因子の$4\times4$行列として見たとき
その階数は高々二となる．
しかしgeometric条件はその行列と転置のkernelがSegreを避ける
ことを含み，階数二以下を許さない
（補題~\ref{lem:qa-cubic-adjugate-two-invariants}）．
従って$\mathscr G\subset P_8^{(2)}$である．
$f$可逆な$cubic$ H2 chartにおける残余群$G_f$は，
前の四重テンソル表示の四つの矢群の有効群であり，
双対に使う因子についてscalarを逆に読めば
$G_f\simeq G_8/K_{\rm QP}$となる．
命題~\ref{prop:qa-global-cubic-four-qubit-good-moduli}は
この商を実際の相対cubic H2 chartに同定するので，
結論を得る．
\end{proof}

\begin{corollary}[二cubic正規化を貼る操作は一般には右同値でない]
\label{cor:qa-cubic-qp-inversion-not-right-equivalent}
両cubic chartの交わり$\Omega^\times$上で、
$f$可逆な正規化と$g$可逆な正規化から得る二つの
四重テンソル$T,T^{-1}$は、一般点で頂点を固定する
八矢weight-preserving right-equivalenceでは結ばれない。
同値となり得る点の粗good moduli像は
\[
 e_3=de_1\quad\text{または}\quad e_3=-de_1
 \tag{Fix}\label{eq:qa-cubic-qp-inversion-right-equivalence-locus}
\]
という二つの固有閉部分の和に含まれる。
従って二cubic chartの全域比較には、
八矢QPの単一のright-equivalence商だけでなく、
実際のmutation・window変更に対応する貼合せが必要である。
\end{corollary}
\begin{proof}
命題~\ref{prop:qa-cubic-weighted-inversion-gluing}によって、
二つのchartの四重テンソル不変式は
$[e_1:e_2:d:e_3]$と
$\iota[e_1:e_2:d:e_3]=[e_3:d^2e_2:d^3:d^4e_1]$
である。
weight-preserving right-equivalenceなら四因子の
$\GL_2^4$作用による同じorbitであるため、その四量子ビットGIT
同次不変式も一致する。
命題~\ref{prop:qa-cubic-inversion-weight-two-fixed-locus}により、
この等式の固定部分は正に表示した二つの閉部分である。
それらが真であることは、例えばCartanの固有値を
$(1,2,3,4)$とした可逆行列の不変式
\[
 (e_1,e_2,d,e_3)=(30,273,24,820)
\]
からも従う。この行列では
$p_{\rm adj}=9625$、$q_{\rm adj}=6510100$が非零なので
$\Omega^\times$に属し、一方$820\ne\pm24\cdot30$である。
従って一般点で右同値にはならない。
\end{proof}

\begin{lemma}[nilpotent mass上の対角慣性は非平坦になり得る]
\label{lem:qa-cubic-qp-nilpotent-mass-scaling}
一組の重み二の可縮pair $x,y$を八矢potential $H$に加えて
\[
 W_\lambda=H+\lambda yx
\]
とする。$x\mapsto tx,\ y\mapsto sy$で与えられる
二つの矢の対角スカラー変換のうち、
potentialを固定する部分群スキームは
\[
 \mathcal T_\lambda=
 \operatorname{Spec}
 \frac{R[t^{\pm1},s^{\pm1}]}{(\lambda(ts-1))}
 \ \subset\ (\mathbb G_{m,R})^2.
 \tag{Mass}\label{eq:qa-qp-mass-scaling-group}
\]
$\lambda$が可逆なら$\mathcal T_\lambda\simeq\Gm_R$、
$\lambda=0$なら$\mathcal T_\lambda\simeq(\Gm_R)^2$である。
特に$R=k[\epsilon]/(\epsilon^2)$、$\lambda=\epsilon$では
$\mathcal T_\lambda$は$R$上平坦でない。

この補題は通常のright-equivalence商内の明示的な部分慣性の
計算であり、$W_\lambda$が全次数flatなAS正則族をなすこと、
あるいは全慣性群スキームが$\mathcal T_\lambda$に等しいことは
主張しない。
\end{lemma}
\begin{proof}
$H$は$x,y$を含まず、変換によって
$W_\lambda$は$H+\lambda st\,yx$へ移る。
二種類の巡回道は独立だから固定条件は
$\lambda(st-1)=0$に等しい。
これは任意可換$R$-代数への基底変更後も同じ等式である。
$\lambda$が単元なら$st=1$、
$\lambda=0$なら条件は空である。

双対数環での非平坦性を証明する。
表示した環を$A$とし、$z=st-1\in A$と置く。
$\epsilon z=0$だが、$z$は$A/\epsilon A
=k[t^{\pm1},s^{\pm1}]$で非零である。
もし$A$が$R$上flatなら、$R$の完全列
$0\to(\epsilon)\to R\to k\to0$
を$A$でtensorすることで
$\ker(\epsilon:A\to A)=\epsilon A$となる。
これは$z\notin\epsilon A$と矛盾する。
\end{proof}


\subsection{非最小十矢のmassゲージと相対群oid}
\label{subsec:qa-ten-arrow-mass-gerbe}

次は八矢領域から階数三のscheme壁への移行を扱う．
四辺のrank-two矢束を$U,V,W,T$とし，偶数対角の
$x:0\to2,\ y:2\to0$を追加する．
$Q_8$の矢だけで作る道
$q:0\to2$，$p:2\to0$を用いて，
\[
 \mathcal W_\lambda=H+yq+px+\lambda yx
 \tag{MG}\label{eq:qa-mass-gauge-potential}
\]
とする．ここで$H$は四辺を一周するquartic potentialである．
$q\in V\otimes U$，$p\in T\otimes W$であり，これらを
$x,y$に依存しないものとして保つ．
$\lambda$は任意の可換係数環$R$の元で，零因子でもよい．

\begin{theorem}[massを横断する滑らかなゲージ群]
\label{thm:qa-mass-gauge-smooth-group}
\[
 \mathcal U_\lambda
 =\Spec R[t,(1+\lambda t)^{-1}],\qquad
 t\star s=t+s+\lambda ts
 \tag{UG}\label{eq:qa-mass-gauge-group-law}
\]
は$R$上smoothかつflatな相対次元一の可換群スキームである．
その$\mathcal W_\lambda$への重みを保つ右同値作用は
$a=1+\lambda t$として
\[
 x\longmapsto ax+tq,\qquad
 y\longmapsto a^{-1}(y-tp),
 \qquad U,V,W,T\text{の各矢は固定}
 \tag{GA}\label{eq:qa-mass-gauge-action}
\]
で与えられる．この右同値はpotentialを巡回同値の意味で
厳密に固定し，$\Jac(Q_{10},\mathcal W_\lambda)$には
恒等自己同型を誘導する．
$R/(\lambda)$上では$\mathcal U_0\simeq\mathbb G_a$，
$\lambda$が単元の開では
$t\mapsto1+\lambda t$で$\mathcal U_\lambda\simeq\Gm$となる．
$\lambda=\epsilon$ on $R=k[\epsilon]/(\epsilon^2)$でも
$\mathcal U_\epsilon$はsmooth/flatである．
$(\lambda,q)$が係数環上単位idealを生成するとき，この作用は
right-equivalenceの安定化群への群スキーム単射である．
\end{theorem}
\begin{proof}
$1+\lambda t$を可逆化した$\mathbb A^1_R$の主開なので
smooth/flatは直ちに従う．積の式に対し
$1+\lambda(t\star s)=(1+\lambda t)(1+\lambda s)$，
逆元は$-t/(1+\lambda t)$である．
右同値の逆写像も同じ形式で与えられる．

式~\eqref{eq:qa-mass-gauge-action}をpotentialに代入する．
$x'=ax+tq$，$y'=a^{-1}y-ta^{-1}p$と置くと，
$y'q,px',\lambda y'x'$の展開で，
$yx,yq,px,pq$の巡回係数はそれぞれ
\[
 \lambda,\qquad
 a^{-1}(1+\lambda t)=1,\qquad
 a(1-\lambda t/a)=1,\qquad
 t-t/a-\lambda t^2/a=0 .
\]
$pq$と$qp$は巡回同値であり，ゆえに
$\mathcal W_\lambda(x',y')\sim_{\rm cyc}
\mathcal W_\lambda(x,y)$である．
さらに
\[
 x'-x=t(q+\lambda x)=t\,\partial_y\mathcal W_\lambda,
 \qquad
 y'-y=-\frac{t}{a}(p+\lambda y)
      =-\frac{t}{a}\,\partial_x\mathcal W_\lambda .
\]
従ってJacobian商では各生成矢の像が等しい．
$\lambda$が単元なら$a$が任意の可逆元を走り，
$\lambda=0$なら加法群となる．
もし$\lambda$が単元でない点で$q$の一成分が単元なら，
$x$に加えた$q$の係数から$t$が復元される．
$\lambda$が単元なら$x$のスカラー$a$から$t$が復元される．
単位ideal条件によりこれらがZariski局所に適用でき，
作用は単射である．
\end{proof}

\begin{theorem}[一積可逆なscheme壁の十矢QP群oid]
\label{thm:qa-ten-arrow-relative-finite-gerbe}
$Z^{\rm nd}$と$G$を有限foundationの
定義~\ref{def:qa-foundation-open}・命題~\ref{prop:qa-foundation-stack}
に従うものとする．$b:W\otimes V\to F$が可逆な開の
unrigidified stackを
\[
 \mathfrak F^{\rm lin}_{b\text{-inv}}
   =[(Z^{\rm nd}\cap D(\det b))/G]
\]
と置く．十矢potential空間のうち
式~\eqref{eq:qa-mass-gauge-potential}の係数が次の
二条件を満たす開groupoidを$\mathfrak P^{\rm good}_{10,e}$
とする：
\begin{enumerate}
 \item 全幾何点で，$\lambda\ne0$または$q$が非分解可能
       （すなわち$D(\lambda)\cup D(\det q)$）．
 \item $\mathcal D=(V\otimes U)\oplus\mathcal O x$，
       $r=q+\lambda x$とし，
       $\partial_T\mathcal W_\lambda:T^\vee\to W\otimes\mathcal D$
       を取ると，次の写像の像が階数四の直和因子である：
       \[
       (W\otimes\mathcal O r)\oplus T^\vee
          \longrightarrow W\otimes\mathcal D,\qquad
       (w\otimes r,\tau)\longmapsto
             w\otimes r+\partial_T\mathcal W_\lambda(\tau).
       \tag{FR}\label{eq:qa-mass-gauge-relation-bundle}
       \]
\end{enumerate}
この開はweighted right-equivalence groupoidの制限として
任意の基底変更に安定である．
その各対象から
\[
 E=\mathcal D/(\mathcal O r),\quad
 F=W\otimes V,\quad
 N=(W\otimes\mathcal D)/
       (W\otimes r+\operatorname{im}\partial_T\mathcal W_\lambda)
 \tag{FD}\label{eq:qa-mass-gauge-finite-reconstruction}
\]
を構成すると，有限foundationの全積を保存する射
\[
 \Phi:\mathfrak P^{\rm good}_{10,e}
          \longrightarrow\mathfrak F^{\rm lin}_{b\text{-inv}}
\]
が得られる．$\Phi$はfppf局所にessentially surjectiveで，
全Isom sheafの射もfppf局所に全射である．
各対象における自己同型の核は
定理~\ref{thm:qa-mass-gauge-smooth-group}の
$\mathcal U_\lambda$に正確に等しい．
すなわち$\Phi$はfppf局所に$\mathcal U_\lambda$を
相対慣性とするgerbeであり，幾何的rank三の壁上では
$\mathbb G_a$，可逆cubic開では$\Gm$となる．
$\lambda$がnilpotentでも群スキームはsmooth/flatである．
この定理は有限代数のgroupoidについての主張であり，
任意のpotentialのAS正則性を含まない．
\end{theorem}
\begin{proof}
第一条件は幾何点での条件であり，
$\lambda=0$では$q$の非分解性は二因子の基底変更と
可逆scalar倍で保存される．
$x$への二段道の付加は$q$を$\lambda$倍の道だけ変えるので，
この条件は右同値に不変である．第二条件も，関係束とその
直和因子性の不変性から保たれる．

$r$は各幾何点で非零のprimitive元であり，
$\mathcal O r\subset\mathcal D$はrank-one直和因子である．
第二条件によって$E,N$はそれぞれrank四・六の束となる．
二段の積$a:V\otimes U\to E$は商写像，$b$は恒等，
三段の積$c:W\otimes E\to N$は商写像，
$d:(W\otimes V)\otimes U\to N$はその合成である．
従って結合律と$(c,d)$の全射性を満たす．
$\lambda\ne0$なら$a$は可逆であり，
$\lambda=0$では$\ker a=kq$がSegreを避ける．
ゆえに全ての幾何点で$Z^{\rm nd}$の条件を満たす．
$x,y,T$のright-equivalenceはrelationの枠やliftを変更する
だけなので，この構成は相対groupoid射となる．

逆方向を局所的に構成する．
$b$が可逆なfinite foundationでは結合律により
$d$は$c$と$a$から決まり，$(c,d)$の全射性は$c$の全射性と同値である．
$a$は各幾何点でrank三または四である．
ある三次minorを可逆化し，源$V\otimes U$の第四基底を$q$，
標的$E$の第四補元を$x$と選べば
$a\sim\operatorname{diag}(I_3,-\lambda)$となり，
$E\simeq((V\otimes U)\oplus\mathcal O x)/(q+\lambda x)$を得る．
rank三の点では$q$は既存の非分解可能kernel線の生成元なので，
$\det q$は単元に取れる．
$c$の核$K\subset W\otimes E$はrank二の直和因子である．
その局所基底を$T^\vee$で標識し，
$W\otimes\mathcal D\to W\otimes E$へのliftを選べば
$\partial_T\mathcal W_\lambda$の二本を得る．
各liftを三段道成分と$x$成分に分けると
$H+px$が復元し，$y(q+\lambda x)$を追加して
十矢potentialを得る．構成~\eqref{eq:qa-mass-gauge-finite-reconstruction}
は元の全積を回収する．
これらの枠・splitting・liftはfppf局所に取れ，
任意の有限foundation同型も同じ選択を合わせることで
十矢right-equivalenceへ局所に持ち上がる．

最後に自己同型核を計算する．
有限foundationを恒等に保つ右同値は$U,V,W$上で恒等であり，
$x\mapsto\alpha x+z$について
$\alpha x+z-x$が$E$で零である．
$r=q+\lambda x$がprimitiveだから，これは唯一の$t$について
\[
 \alpha=1+\lambda t,\qquad z=tq
\]
と同値である．$\alpha$は単元なので$t\in\mathcal U_\lambda$．
また$\partial_T\mathcal W_\lambda$の$W\otimes E$への像は
rank二直和因子であり，$N$を恒等に固定する
potential自己同型は$T$の基底を変更できない．
残る$y\mapsto\beta y+h$をpotentialの係数と比較する．
$r$がprimitiveであることから
\[
 \beta=(1+\lambda t)^{-1},
 \qquad h=-t(1+\lambda t)^{-1}p
\]
が一意に強制される．これは
式~\eqref{eq:qa-mass-gauge-action}そのものである．
従って核の群スキームは$\mathcal U_\lambda$に尽き，
任意基底変更でも同じ計算が成立する．
\end{proof}

\begin{corollary}[実際のmutation AS族への適用範囲]
\label{cor:qa-mass-gauge-actual-as-window}
定理~\ref{thm:qa-relative-mutation-nonminimal-as-resolutions}
の局所十矢表示
$\mathcal W_\delta=\delta yx+yq+px+H$は
上記$\Phi$の対象である．特に
$\delta=\epsilon$ on $k[\epsilon]/(\epsilon^2)$でも
finite foundationの全成分と全積は平坦に回収され，
raw QP自己同型からfinite foundation自己同型への核には
滑らかな$\mathcal U_\epsilon$が現れる．
補題~\ref{lem:qa-cubic-qp-nilpotent-mass-scaling}の
非平坦な対角scalar部分慣性$\mathcal T_\epsilon$は，
別のraw potential $H+\epsilon yx$ に対する模型である．
実際のmutation AS potential $H+yq+px+\epsilon yx$ に
そのまま同じ部分群があるとは主張しない．
これに対し本定理の$\mathcal U_\epsilon$は
実際のmutation AS族に作用する滑らかなゲージ核である．
全相対heartと選択したmutation窓の一致は
定理~\ref{thm:qa-relative-rot-mut-completion-heart}による．
ここから選択した窓を忘れた全AS/QP stackと
内在的$H_2$の同値を推論しない．
\end{corollary}
\begin{proof}
補題~\ref{lem:qa-relative-mutation-seven-three-presentation}は
$q$のprimitive性と$\delta$の任意の冪零値を許し，
命題~\ref{prop:qa-relative-mutation-ordinary-jacobian}は
同じpotentialが実際のflatな相対有限代数・全次数完成を
与えることを示す．前定理を適用すれば結論が従う．
\end{proof}

\begin{remark}[十二矢QPの余剰慣性と残る比較]
\label{rem:qa-cubic-qp-remaining-inertia}
命題~\ref{prop:qa-stabilized-inertia}で既証明の
\(
 \mathcal P_{12}^{\lambda\mu\ne0}
 \simeq\mathcal P_8\times B(\Gm^2)
\)
をgeometric開に制限すれば、十二矢QPから上の有限
foundationへの忘却には、通常の四頂点scalarの$\Gm^3$に加え
二つの可縮pairの$\Gm^2$がある．
従ってraw十二矢QPの全right-equivalence群oidを
$\mathfrak U_f$へ忠実に同定してはならない．
本項は八矢の相対同型層を可逆cubic開で決着させるものであり，
rank-three壁・noncubic・nilpotent massを持つ十／十二矢での
相対QP群oidや同じheartを保つ比較までは証明しない．
特に\(\operatorname{coker}\mu_G=\OO/(\delta)\)が
非平坦となる一次変形
（命題~\ref{prop:qa-relative-as10-minimal-arrow-wall}）を
幾何fiberごとの最小矢数だけから復元することはできない．
\end{remark}

% !TeX root = ../quiver-avatar-moduli.tex
% 2026-11-22: square / rotated-square の幾何的・有限代数的基礎を本文内で証明
\section{Square と rotated square の自己完結な幾何的基礎}
\label{sec:qa-self-contained-foundations}

この節は，未公刊草稿 \cite{AvatarVSeven2026} の第I--II部から
本稿Cの比較に必要な部分を取り出して，仮定と証明を与え直す．
とくに \emph{反標準曲線上の幾何的 avatar，同型類の相互変換，
有限 directed foundation の復元} は以下の証明だけで読める．
これと，無限周期の AS--Gorenstein 分解，Noether 性，
内在的 Hilbert 関手の認識とは別の定理である．
後者の全証明を有限次元の計算から自動的に推論しない．

\subsection{定義と非被約曲線を含む消滅補題}

$k$ は標数零の代数閉体とする．
$Y$ は，$h^2=2,\ hf=1,\ f^2=0,\ -K_Y=2h$ を満たす
標識付き $\Sigma_0=\PP^1\times\PP^1$
（$h=f+g$）または
$\Sigma_2=\PP_{\PP^1}(\OO\oplus\OO(2))$
（負断面 $S=h-2f$）とする．
$D\in|2h|$ を有効 Cartier 因子とし，
$Y=\Sigma_2$ では $D\cap S=\varnothing$ を課す．
以下，$F=\OO_D(f)$，$H=\OO_D(h)$，
$C=\OO_D(h-2f)$ とする．
$\Sigma_2$ の場合には $C\simeq\OO_D$ である．
曲線は非被約・可約でもよい．

\begin{lemma}[部分曲線と Picard 消滅]
\label{lem:qa-core-curve-vanishing}
上の $D$ について
$\omega_D\simeq\OO_D$，
$h^0(\OO_D)=h^1(\OO_D)=1$ が成り立つ．
真の非零な有効部分曲線 $0<Z<D$ には
$Z(D-Z)\ge2$，したがって $p_a(Z)\le0$ が成り立つ．
さらに $R\in\Pic^0(D)$ とすると，
\begin{enumerate}
\item $R\not\simeq\OO_D$ なら $H^0(D,R)=H^1(D,R)=0$;
\item $FR$ と $(H F^{-1})R$ は大域生成な次数二の
      pencil で，$h^0=2$，$H^1=0$;
\item $HR$ は very ample で，$h^0=4$，$H^1=0$;
\item $CR\not\simeq\OO_D$ なら $H^0(D,CR)=H^1(D,CR)=0$.
\end{enumerate}
この主張は $D$ の滑らかさも被約性も仮定しない．
\end{lemma}
\begin{proof}
随伴公式と
$0\to\omega_Y\to\OO_Y\to\OO_D\to0$
から最初の主張が従う．
$\Sigma_0$ で $Z=af+bg$（$0\le a,b\le2$）なら
$Z(D-Z)=2a+2b-2ab$ であり，$0,D$ 以外は二以上である．
$\Sigma_2$ では $D\cap S=\varnothing$ より
$Z\cdot S=0$ であるから，$Z=aS+bf$ と書けば $b=2a$．
$0<Z<D$ のとき $Z\equiv h$ であり，
$Z(D-Z)=h^2=2$ となる．
随伴公式
$2p_a(Z)-2=-Z(D-Z)$ も得られる．

ここで非被約曲線にも使える二つの観察を証明しておく．
線束 $N$ の切断が generic に非零となる最大の部分曲線を
$Z$ とすると，その切断は
$N|_Z(-(D-Z))$ の generic に単射な切断を誘導する．
ゆえに
\[
 \deg_Z N\ge Z(D-Z).                                      \tag{*}
\]
また長さ $\ell$ の零次元部分 scheme $\xi\subset D$ と
線束 $M$ に対して $H^1(I_\xi M)\ne0$ なら，
Serre 双対性から $I_\xi M\to\omega_D$ の非零射を得る．
像の generic annihilator が定める $Z$ 上に射を因数分解し，
$Z$ の双対化層と余核の長さを比較すると
\[
 \deg_Z M\le 2p_a(Z)-2+\ell .                             \tag{**}
\]
ここで $\omega_Z=\mathcal{H}om_D(\OO_Z,\omega_D)$
という随伴を用いる．両議論は局所的な annihilator と
Cohen--Macaulay 曲線の双対性だけに依存するので，
nilpotent layer を持つ部分曲線にも適用できる．

$R\in\Pic^0(D)$ の各成分次数は零である．
$R$ に切断があれば，真部分 $Z$ への
generic support は $(*)$ に反し，全体への切断なら
零点を持たず $R\simeq\OO_D$ となる．
$\omega_D\simeq\OO_D$ と双対性により (1) が従う．
$F$ と $HF^{-1}$ は各部分曲線上非負次数を持つ．
$\Sigma_2$ では $HF^{-1}\simeq F$ なので同じ主張である．
$(**)$ を $\ell=0,1$ に適用すると
$H^1$ の消滅と長さ一の点の生成が従い，
Riemann--Roch から切断数は二である．
$H$ は真部分曲線上正次数を持つ
（$\Sigma_2$ の唯一の例外 $S$ は $D$ と交わらない）．
$(**)$ を $\ell=0,1,2$ に適用すると $HR$ は
全ての長さ二の cluster を分離するので very ample，
かつ $h^0(HR)=4$ となる．
最後に $\Sigma_2$ では $C=\OO_D$ なので (1) に帰着する．
$\Sigma_0$ では $C=g-f$ であり，
真部分曲線 $Z=af+bg$ に対する
$C\cdot Z=a-b\ge 2a+2b-2ab$ は成立しない．
$(*)$ は $CR$ の切断を排除し，全体への切断は
$CR\simeq\OO_D$ を強制する．
$C^{-1}R^{-1}$ に同じ議論を適用し，双対性で
$H^1(CR)=0$ を得る．
\end{proof}

\subsection{二つの pencil と $\Sigma_2$ への再埋め込み}

\begin{lemma}[独立な pencil と一致する pencil]
\label{lem:qa-core-two-pencils}
$A,B\in\Pic(D)$ を大域生成の次数二線束とし，
$h^0(A)=h^0(B)=2$，
$AB$ を very ample で $h^0(AB)=4$ とする．
\begin{enumerate}
\item $H^*(BA^{-1})=0$ なら，完全 pencil の積は
$D\hookrightarrow\PP H^0(A)^\vee\times\PP H^0(B)^\vee$
という閉埋め込みで，像は $(2,2)$ Cartier 因子である．
\item $B\simeq A$ かつ $p:D\to\PP^1$ が有限なら，
$A^2$ の完全線形系から
$D\hookrightarrow\Sigma_2$ を得る．
像は負断面を避ける $|2h|$ の Cartier 因子であり，
$\OO_D(f)=A,\ \OO_D(h)=A^2$ である．
この埋め込みは標識付き曲面の同型を除き一意である．
\end{enumerate}
\end{lemma}
\begin{proof}
(1) $A$ の評価列を $B$ で tensor した
\[
 0\to BA^{-1}\to H^0(A)\otimes B\to AB\to0
\]
に大域切断を取る．消滅仮定によって
$H^0(A)\otimes H^0(B)\xrightarrow{\sim}H^0(AB)$
である．従って Segre 埋め込みとの合成は
very ample な $|AB|$ であり，もとの積も閉埋め込みである．
Cohen--Macaulay 曲線の滑らかな曲面内への余次元一閉埋め込みは
Cartier であり，各 ruling との交点数二から類は $(2,2)$ となる．

(2) $H^0(A^2)$ は四次元で，
$\Sym^2H^0(A)$ は三次元部分空間となる．
$p$ が有限ゆえ非定数であり，
$\Sym^2H^0(A)\to H^0(A^2)$ は単射である．
補元を $z$ とすると，$A^2$ への商
$\OO_D\oplus p^*\OO_{\PP^1}(2)\to A^2$ は
$(z,1)$ で与えられる．
これは quotient 規約の
$\Sigma_2=\PP_{\PP^1}(\OO\oplus\OO(2))$
への射を定める．第二成分が可逆なので負断面を避ける．
$|h|$ との合成は完全な $|A^2|$ であり，閉埋め込みである．
$fD=2,\ hD=4$ から像は $2h$ である．
$z$ の変更は三角形の束自己同型で吸収される．
\end{proof}

\begin{theorem}[square と rotated square の avatar の相互変換]
\label{thm:qa-core-avatar-square-rotated}
次の二種類の標識付き曲線型 avatar の同型類は自然に全単射する：
\[
\begin{array}{ll}
\text{\rm square:}&
(D\subset\PP^1\times\PP^1;\ F,G,P),\
P\in\Pic^0(D)\setminus\{\OO_D\},\\
\text{\rm rotated:}&
(D\subset Y;\ F,H,P),\
Y=\Sigma_0\text{ または }\Sigma_2,\
P\not\simeq\OO_D,\ C,\
D\cap S=\varnothing\text{ if }Y=\Sigma_2.
\end{array}
\]
ここで square の $F,G$ は順序付き ruling，
rotated の $F,H$ は $f,h$ の制限，$C=HF^{-2}$ である．
同型類の対応は
\begin{equation}
\begin{aligned}
\mathrm{rotated}\longmapsto\mathrm{square}:&
  \quad F_s=F,\quad G_s=HF^{-1}P^{-1},\\
\mathrm{square}\longmapsto\mathrm{rotated}:&
  \quad F_r=F_s,\quad H_r=F_sG_sP .
\end{aligned}
\label{eq:qa-core-rsq-sq}
\end{equation}
さらに rotated の Hilbert 曲面が $\Sigma_2$ となるのは
$G_sP\simeq F_s$ の場合に限る．
\emph{square の許容条件に}
$P\ne F_sG_s^{-1}$ \emph{や}
$P\ne G_sF_s^{-1}$ \emph{を追加してはならない．}
\end{theorem}
\begin{proof}
rotated からは補題~\ref{lem:qa-core-curve-vanishing}より
$F$ と $HF^{-1}P^{-1}$ は二つの大域生成 pencil で，
積 $HP^{-1}$ は very ample である．
その比 $HF^{-2}P^{-1}=CP^{-1}$ は
仮定により非自明で，補題の (4) から acyclic である．
補題~\ref{lem:qa-core-two-pencils}(1) により
$(2,2)$ Cartier 因子へ埋め込める．

逆に square から $F_s$ と $G_sP$ を取る．
これらは大域生成次数二 pencil で，
積 $F_sG_sP$ は very ample である．
$G_sP\not\simeq F_s$ なら，比
$G_sF_s^{-1}P$ は非自明なので
補題~\ref{lem:qa-core-curve-vanishing}(4) から acyclic であり，
$\Sigma_0$ 上に再埋め込みができる．
$G_sP\simeq F_s$ のとき
$F_s^2\simeq F_sG_sP$ は very ample である．
$F_s$ は $G_s$ とともに各成分上正次数を持ち
（$P$ の各成分次数は零），
対応する $D\to\PP^1$ は有限である．
補題~\ref{lem:qa-core-two-pencils}(2) で $\Sigma_2$ を得る．
この場合 $H_r=F_s^2$ なので $C_r=\OO_D$ である．
一般の場合 $C_r=G_sF_s^{-1}P$ であり，
$P=C_r$ なら $F_s=G_s$ となる．
元の $(2,2)$ 埋め込みでは二つの ruling は異なり
$H^*(G_sF_s^{-1})=0$ なので，これは起こらない．
最後に式~\eqref{eq:qa-core-rsq-sq}の合成は
$F_s,G_s$，または $F_r,H_r$ を厳密に回収する．
完全 pencil と補元の一意性により曲面上の同型類も回収する．
\end{proof}

\subsection{有限 directed foundation と反標準行列式}

以下 square の場合には $H_s=F_sG_s$ と書く．
四頂点の directed algebra は，$i<j$ に
$\Hom_D(L_i,L_j)$，対角に $k$，逆向きに零を置き，
積を切断の乗法で定めたものとする．
この定義は $\Perf(D)$ の full dg endomorphism algebra
を意味しない．

\begin{proposition}[二種類の有限 foundation と点曲線の復元]
\label{prop:qa-core-directed-foundations}
曲線型 avatar に対して
\[
\begin{aligned}
B_{\rm sq}&=\End_{\rm dir}(P,F_s,G_s,F_sG_s),\\
B_{\rm rsq}&=\End_{\rm dir}
 (P,C_r,H_rF_r^{-1},H_r^2F_r^{-1})
\end{aligned}
\]
を考える．square の Hom 次元上三角行列は
$\left(\begin{smallmatrix}
1&2&2&4\\0&1&0&2\\0&0&1&2\\0&0&0&1
\end{smallmatrix}\right)$，
rotated のそれは
$\left(\begin{smallmatrix}
1&0&2&6\\0&1&2&6\\0&0&1&4\\0&0&0&1
\end{smallmatrix}\right)$
である．いずれも対応する四頂点箙とその最小関係を与え，
四頂点の点表現を三頂点へ忘れる射は，
関係の普遍評価行列の余核の相対射影化として
\emph{もとの} $D\hookrightarrow Y$ を回収する．
特に square の行列は rank 四から三への Cartier 退化で，
rotated は rank 二から一への退化である．
\end{proposition}
\begin{proof}
線束の差の $H^0$ を補題~\ref{lem:qa-core-curve-vanishing}
で計算すれば表示の Hom 次元を得る．
square で非自明な矢は四つの二次元空間
$H^0(F_sP^{-1}),H^0(G_sP^{-1}),H^0(F_s),H^0(G_s)$
であり，関係は $0\to3$ の四次元空間である．
rotated では $0\to2$ と $1\to2$ の二次元空間，
$2\to3$ の四次元空間が生成矢で，
$0\to3$ と $1\to3$ に各二次元の関係がある．
後者の二つの関係は pencil の評価列から直接得られる．

square での四つの関係を同時に扱うには，
$\PP^1$ の対角の Koszul 分解を二つ外積して
$\PP^1\times\PP^1$ の対角分解を作る．
それを $\iota_*(H_sP)$ に作用させて大域切断を取ると，
補題の消滅と $h^0=2,2,4$ から
\[
 0\longrightarrow\OO_Y(-f)^{\oplus2}
       \oplus\OO_Y(-g)^{\oplus2}
 \longrightarrow\OO_Y^{\oplus4}
 \longrightarrow\iota_*(H_sP)\longrightarrow0
\]
となる．曲面上で双対化し，Cartier 双対性を使うと
\begin{equation}
0\to\OO_Y^{\oplus4}
 \xrightarrow{\kappa_{\rm sq}}
 \OO_Y(g)^{\oplus2}\oplus\OO_Y(f)^{\oplus2}
 \to\iota_*(H_sP^{-1})\to0 .
\label{eq:qa-core-square-matrix}
\end{equation}
大域切断と $(-f),(-g)$ への tensor は
元の四矢空間と四関係を与えるので，
$\kappa_{\rm sq}$ は普遍評価行列である．
余核は局所的に $\OO_Y/(q)$ の自由 rank-one
加群であり，局所行基本変形により
$\kappa_{\rm sq}\sim\operatorname{diag}(1,1,1,q)$．
ゆえに $\PP_Y(\coker\kappa_{\rm sq})=V(q)=D$ である．

rotated では $N=H_rF_r^{-1}P^{-1}$ とおく．
$\Sigma_0$ では同じ積の対角分解と
$H^*(P^{-1})=H^*(C_rP^{-1})=0$ を用い，
$\Sigma_2$ では有限平坦な二重射
$p:D\to\PP^1$ の $p_*N\simeq\OO_{\PP^1}^{\oplus2}$
を用いる．いずれも
\[
0\to\OO_Y(-h)^{\oplus2}\to\OO_Y^{\oplus2}
  \to\iota_*N\to0
\]
となる．$\OO_Y(h)$ で tensor すると
\begin{equation}
0\to\OO_Y^{\oplus2}
 \xrightarrow{\kappa_{\rm rsq}}
 \OO_Y(h)^{\oplus2}
 \to\iota_*(H_r^2F_r^{-1}P^{-1})\to0 .
\label{eq:qa-core-rotated-matrix}
\end{equation}
元の二関係と八道から大域切断で
$\kappa_{\rm rsq}$ を回収する．
局所標準形は $\operatorname{diag}(1,q)$ であり，
その相対射影化は再び $D$ である．
最後の普遍線束の商を
$H_s$ または $H_r^2F_r^{-1}$ で正規化すると
$P$ が回収される．
これは単なる幾何点の行列式の一致ではなく，
余核による scheme-theoretic な復元である．
\end{proof}

\begin{remark}[この節から AS 分類へ進むための残りの入力]
\label{rem:qa-core-as-scope}
上の構成は \emph{任意の}曲線型 avatar に対して成立する．
有限 foundation から全整数添字の AS 正則
$\mathbb Z$-代数を得るには，さらに四つの canonical
kernel--cokernel 列を全周期で構成し，
mutation と Serre 関手を比較する必要がある．
その際に非可換圏の実際の最小射影分解と
AS--Gorenstein 双対性を点検しなければならない．
本稿の他の節が有限 foundation と AS 代数を
区別するのはこのためである．
\end{remark}

% !TeX root = ../nc-quadric-avatar.tex
% 2026-11-22: 幾何的avatarの四つのcanonical列
\section{反標準 avatar の canonical complexes と周期}
\label{sec:qa-self-contained-helix}

この節は前節の幾何的補題から四項の canonical complex を
構成する．これは有限 foundation と全周期の関係を明示する
ためのものであり，任意の相対基底上の AS 正則性と
内在的な心の同定を無条件に主張するものではない．

\begin{proposition}[rotated square の四つの完全列]
\label{prop:qa-core-rotated-complexes}
定理~\ref{thm:qa-core-avatar-square-rotated} の
rotated square 曲線型 avatar を固定し，
\[
\begin{gathered}
L_0=P,\quad L_1=HF^{-2},\quad
L_2=HF^{-1},\quad L_3=H^2F^{-1},\\
L_4=H^2P,\quad L_5=H^3F^{-2},\quad
L_6=H^3F^{-1}P^2,\quad L_7=H^4F^{-1}P^6 .
\end{gathered}
\]
これらの線束の間には directed foundation の評価写像と
最小関係から作られる次の四つの局所分裂完全列がある：
\begin{align}
0&\to L_0\to L_2^{\oplus2}\to L_3^{\oplus2}
     \to L_4\to0,\label{eq:qa-core-rsq-complex0}\\
0&\to L_1\to L_2^{\oplus2}\to L_3^{\oplus2}
     \to L_5\to0,\label{eq:qa-core-rsq-complex1}\\
0&\to L_2\to L_3^{\oplus4}
     \to L_4^{\oplus2}\oplus L_5^{\oplus2}
     \to L_6\to0,\label{eq:qa-core-rsq-complex2}\\
0&\to L_3\to L_4^{\oplus2}\oplus L_5^{\oplus2}
     \to L_6^{\oplus4}\to L_7\to0.
     \label{eq:qa-core-rsq-complex3}
\end{align}
任意の $q\in\mathbb Z$ に対して，$i=0,1,2,3$ とし
\[
\begin{array}{c|cccc}
i&0&1&2&3\\ \hline
G_i&0&h-2f&h-f&2h-f\\
d_i&0&0&2&6\\
\varepsilon_i&1&0&0&0
\end{array}
\]
と書くと，
\begin{equation}
L_{4q+i}=\OO_D(G_i+2qh)
           P^{\,4q(q-1)+d_iq+\varepsilon_i}
\label{eq:qa-core-rsq-periodic-L}
\end{equation}
は同じ四列の全整数平行移動を持つ．
\end{proposition}
\begin{proof}
$A=L_2L_0^{-1}=HF^{-1}P^{-1}$ は
補題~\ref{lem:qa-core-curve-vanishing} により二切断の
大域生成 pencil である．評価列
$0\to A^{-1}\to\OO_D^{\oplus2}\to A\to0$
を $L_2$ で tensor すると
\[
0\to L_0\to L_2^{\oplus2}\to L_2^2L_0^{-1}\to0 .
\]
他方 $L_3(L_2^2L_0^{-1})^{-1}=FP$ も pencil なので，
その評価列を $L_3$ で tensor した
\[
0\to L_2^2L_0^{-1}\to L_3^{\oplus2}\to L_4\to0
\]
と splice すれば第一列となる．
第二列は $L_2L_1^{-1}=F$ と
$L_3H^{-1}=HF^{-1}$ という二 pencil の評価列
\[
0\to L_1\to L_2^{\oplus2}\to H\to0,\qquad
0\to H\to L_3^{\oplus2}\to L_5\to0
\]
を splice して得られる．
ここで $HF^{-1}$ は二切断の大域生成 pencil である．

第三列には，二つの pencil $a_+=FP$，$b_+=HF^{-1}$
を用いる．比 $b_+a_+^{-1}=HF^{-2}P^{-1}$ は
補題~\ref{lem:qa-core-curve-vanishing} より acyclic，
積 $a_+b_+=HP$ は very ample である．
よって補題~\ref{lem:qa-core-two-pencils} で得る
別の $(2,2)$ 埋め込み上に，
命題~\ref{prop:qa-core-directed-foundations} の
四行四列の対角分解を適用する．
その曲線への導来制限では，
Cartier 因子の標準 Tor 項
$\operatorname{Tor}_1(\iota_*Q,\OO_D)=Q(-D)$
が左端に現れる．
$Q=HP^2$ を代入し，全体を $L_3$ で tensor すると
第三列の四項と係数 $4,2+2,1$ が得られる．
これは局所的に長さ一の Cartier 解消を制限して
得た完全列なので，単に Euler 特性の等号ではない．

第四列も，二つの pencil
$a_0=HF^{-1}P$，$b_0=FP^2$ を使う．
その比 $a_0b_0^{-1}=HF^{-2}P^{-1}$ は acyclic，
積は $HP^3$ で very ample である．
同じ四行四列分解を $Q=HP^2$ に適用し，
導来制限と双対化を行った後 $L_6$ で tensor すると
第四列を得る．
第三，第四列の中間写像は対角分解の評価から得られ，
大域切断に戻すと四本の矢と二本ずつの最小関係に一致する．
全列は線束の局所分裂完全列である．

周期については，各連続四項を共通線束で
tensor して最初の項を $P$ に正規化すると，
二つの ruling の役割を担う線束は
$F_q=FP^{2q}$ と $H_q=HP^{4q}$ となる．
$H_qF_q^{-2}=HF^{-2}$ は不変なので，これらは
同じ非自明性と消滅条件を満たす．
前節の二 pencil の再埋め込みを各 $q\in\mathbb Z$
に適用でき，上の四列を繰り返して構成できる．
線束の差分を帰納的に追うと
式~\eqref{eq:qa-core-rsq-periodic-L} を得る．
\end{proof}

\begin{corollary}[square の四つの完全列]
\label{cor:qa-core-square-complexes}
square avatar の $L_0=P,L_1=F,L_2=G,L_3=FG$ に対し
\[
(L_4,L_5,L_6,L_7)
=(F^2G^2P,\ F^3G^2P^2,\ F^2G^3P^2,\ F^3G^3P^4)
\]
とおくと，
\begin{align*}
0&\to L_0\to L_1^{\oplus2}\oplus L_2^{\oplus2}
       \to L_3^{\oplus4}\to L_4\to0,\\
0&\to L_1\to L_3^{\oplus2}\to L_4^{\oplus2}
       \to L_5\to0,\\
0&\to L_2\to L_3^{\oplus2}\to L_4^{\oplus2}
       \to L_6\to0,\\
0&\to L_3\to L_4^{\oplus4}
       \to L_5^{\oplus2}\oplus L_6^{\oplus2}
       \to L_7\to0
\end{align*}
が canonical な局所分裂完全列となる．
\end{corollary}
\begin{proof}
定理~\ref{thm:qa-core-avatar-square-rotated} により，
square の一つ後方または前方の foundation を
rotated の四項と比較する．
正確には rotated の $(L_{-1},L_0,L_1,L_2)$ を
$L_{-1}=F^{-1}P^2$ から出発して
共通線束 $FP^{-1}$ で tensor すると
$(P,F,G,FG)$ となる．
前命題の列を一つ添字移動し，
全項にこの共通線束を tensor すれば四列を得る．
この操作は曲線上の線束の標準化であって，
非可換圏に $(-)\otimes FP^{-1}$ という自己同値を
仮定したものではない．
\end{proof}

\begin{proposition}[全整数 helix の前向き Ext 消滅]
\label{prop:qa-core-forward-ext-vanishing}
上記の square および rotated square の曲線上の線束列について
\[
 \Ext_D^p(L_i,L_j)=0\qquad(p>0,\ i<j)
\]
が成り立つ．曲線が可約・非被約であっても同様である．
これは曲線上の forward acyclicity の主張であり，
この段階で例外 collection の fullness や
非可換 AS--Gorenstein 性を主張するものではない．
\end{proposition}
\begin{proof}
Gorenstein 曲線 $D$ 上で
$\Ext^p_D(L_i,L_j)=H^p(D,L_jL_i^{-1})$ であり，
$p\ge2$ は自動的に零である．
rotated の連続する四商は，$q$ 番目の区間で
\[
 C P^{-1},\quad FP^{2q},\quad HP^{4q},
       \quad FP^{2q+1},
 \qquad C=HF^{-2}.
\]
唯一の総次数零の商 $CP^{-1}$ は
avatar の除外条件 $P\ne C$ と
補題~\ref{lem:qa-core-curve-vanishing}(4) で acyclic である．
その他の商は各成分上非負次数の線束であり，
少なくとも一つ含めば全体上の次数は正である．
複数の連続商の tensor 積は，単独の
$CP^{-1}$ 以外は各部分曲線上非負次数を持ち，
$D$ 全体上正次数を持つ．
ここで $C$ 自体の成分次数は零ではない場合があるが，
$C$ の両側には $F$ または $H$ の商が隣接し，
それらとの積は各部分曲線に非負次数を持つ．
補題~\ref{lem:qa-core-curve-vanishing}の証明で得た
cluster 消滅の不等式 $(**)$ を $\ell=0$ に適用すれば
$H^1=0$ である．

square では四つの連続商は
\[
 FP^{2q-1},\quad GF^{-1},\quad FP^{2q},
       \quad FGP^{4q+1}
\]
である．$GF^{-1}$ の単独の acyclicity は
$\Sigma_0$ 上の Cartier 制限列，
または補題~\ref{lem:qa-core-curve-vanishing}(4)
による．それ以外は上と同様に
nef な正次数の商を少なくとも一つ含み，
$GF^{-1}$ を隣接商と組にすると
各部分曲線上の次数は非負になる．
同じ cluster 判定から $H^1=0$ が従う．
\end{proof}

\begin{remark}[AS 正則性への論理的な境界]
\label{rem:qa-core-helix-not-as}
前命題と系は，両端と中間項を持つ
\emph{実際に完全な}曲線上の複体を与える．
非可換な全 $\mathbb Z$-代数の AS 正則性には，
これらを canonical exceptional helix の
最小射影分解として実現し，
反対向きの $\Ext^3$ による Gorenstein 性を示す
別の導来圏比較が必要である．
この比較のためには acyclic spherical helix と
有限 directed algebra の Serre 関手を用いる．
その\emph{体上の全AS正則$\mathbb Z$-代数への橋}を
第~\ref{sec:qa-as-avatar-classification}節で
最小分解の必要条件・逆構成・相互逆性まで証明する．
ただし相対基底変更と元の O-marked heart の直接同値は
この四列だけからは結論しない．
\end{remark}

% !TeX root = ../nc-quadric-surface-avatar.tex
% 2026-11-22 / AS/avatar classification, source v7 Parts I-II
\section{AS正則\texorpdfstring{$\mathbb Z$}{Z}-代数とavatarの相互逆な分類}
\label{sec:qa-as-avatar-classification}

本節の目的は有限directed foundationの記述を
\emph{全整数頂点のAS正則$\mathbb Z$-代数の同型類}
との全単射へ昇格させることである。
元の証明は\cite[第2--8節，第11--12節]{AvatarVSeven2026}にある。
前二節で証明した曲線上の消滅・評価行列・canonical複体を入力とし、
残っていたAS--Gorenstein性、全helixの一意性と
必要条件の逆向きもここで示す。
一般的なspherical helix実現定理だけは
\cite[Theorem 5]{OkawaUedaSpherical2020}を外部入力とし、
本件でどのようにAS分解を得るかを明記する。

\subsection{最小分解で指定する二型}
\label{subsec:qa-as-minimal-types}

$A=\bigoplus_{j\ge i}A_{ji}$は局所有限、
$A_{ii}=k$、$A_{ji}=0$ ($j<i$)の
positively directed $\mathbb Z$-代数とする。
$P_i=e_iA$、$S_i$を右射影・右単純加群とする。
頂点をすべて保つ代数同型で分類し、
矢の基底や四周期同型の選択は追加データとしない。

\begin{definition}[二つのAS型]
\label{def:qa-as-square-rotated-minimal}
$I=\mathrm{rsq},\mathrm{sq}$に対し
$A\in\mathrm{AS}_{I}$とは、
すべての$i$で最小右射影分解
\begin{equation}
0\to P_{i-4}\to Q_{2,i}\to Q_{1,i}
 \to P_i\to S_i\to0
\label{eq:qa-as-two-minimal-resolutions}
\end{equation}
を持ち、さらに
\begin{equation}
\Ext_A^q(S_i,A)=0\ (q\ne3),\qquad
\Ext_A^3(S_i,A)\simeq {}_\ell S_{i-4}
\label{eq:qa-as-two-gorenstein}
\end{equation}
を満たすものとする。
$Q_{2,i},Q_{1,i}$は次の四周期表で指定する。
\[
\begin{array}{c|c|l|l}
I&i\bmod4&Q_{2,i}&Q_{1,i}\\ \hline
\mathrm{rsq}&0&P_{i-2}^{\oplus2}&P_{i-1}^{\oplus2}\\
&1&P_{i-3}^{\oplus2}&P_{i-2}^{\oplus2}\\
&2&P_{i-3}^{\oplus4}&P_{i-2}^{\oplus2}\oplus P_{i-1}^{\oplus2}\\
&3&P_{i-3}^{\oplus2}\oplus P_{i-2}^{\oplus2}&P_{i-1}^{\oplus4}\\ \hline
\mathrm{sq}&0&P_{i-3}^{\oplus2}\oplus P_{i-2}^{\oplus2}&P_{i-1}^{\oplus4}\\
&1&P_{i-2}^{\oplus2}&P_{i-1}^{\oplus2}\\
&2&P_{i-3}^{\oplus2}&P_{i-2}^{\oplus2}\\
&3&P_{i-3}^{\oplus4}&P_{i-2}^{\oplus2}\oplus P_{i-1}^{\oplus2}
\end{array}
\]
左右Noether性・幾何性・周期のtrivializationは仮定しない。
\end{definition}

\begin{lemma}[局所双対性と全$\mathbb Z$-代数の回収]
\label{lem:qa-as-serre-unique-recovery}
$A\in\mathrm{AS}_I$、$E_i=\pi P_i\in\Qgr A$とする。
$\mathcal D_{\rm lin}=\operatorname{thick}\{E_i\mid i\in\mathbb Z\}$
において連続する四対象はstrong exceptional tilting generatorである。
そのEnd代数$B$は大域次元二で、
Serre関手$S_B$により
\begin{equation}
E_{i+4}\simeq S_B^{-1}(E_i)[2],\qquad
A_{ji}=\Hom(E_i,E_j)\quad(j\ge i)
\label{eq:qa-as-serre-completion}
\end{equation}
を得る。全代数は初期四頂点のEnd代数から
頂点を保つ同型を除いて一意に決まる。
その数値Euler行列は
\[
G_{\rm rsq}=\begin{pmatrix}
1&0&2&6\\0&1&2&6\\0&0&1&4\\0&0&0&1
\end{pmatrix},\quad
G_{\rm sq}=\begin{pmatrix}
1&2&2&4\\0&1&0&2\\0&0&1&2\\0&0&0&1
\end{pmatrix},
\]
かつ$\dim A_{i+4m,i}=(2m+1)^2$である。
\end{lemma}
\begin{proof}
最小分解を各頂点で評価したEuler漸化式は、
負の成分が零で対角が$k$という初期条件から
全Hom次元を一意に定める。
可換基準collection
$(\OO,\OO(1,0),\OO(0,1),\OO(1,1))$と
$(\OO,\OO(-1,1),\OO(0,1),\OO(1,2))$
にも同じ漸化式が成立し、上の行列を得る。
四周期のHilbert関数は
$h^0(\PP^1\times\PP^1,\OO(2m,2m))=(2m+1)^2$。
これを可換基準モデルとの類比だけに依存させず、
最小分解表から有理母関数として直接検算する。
$r=i\bmod4$、$h_{r,n}=\dim A_{i,i-n}$、
$H_r(z)=\sum_{n\ge0}h_{r,n}z^n$とおき、
$Q_{1,i},Q_{2,i}$中の$P_{i-k}$の重複度を
それぞれ$a_{r,k},b_{r,k}$とする。
表のEuler漸化式は
\[
 (1-z^4)H_r(z)
 =1+\sum_{k=1}^3
 (a_{r,k}-b_{r,k})z^kH_{r-k\bmod4}(z).
 \tag{Euler4}\label{eq:qa-as-four-phase-euler}
\]
これら四本の一次方程式の有理関数解について、
$\zeta^4=1$を満たす原始四乗根$\zeta$を取ると、
各$r$で
\begin{equation}
 \frac14\sum_{j=0}^3H_r(\zeta^jz)
 =\frac{1+6z^4+z^8}{(1-z^4)^3}
 =\sum_{m\ge0}(2m+1)^2z^{4m}.
 \label{eq:qa-as-four-phase-rational-proof}
\end{equation}
左辺は四の倍数次数への射影である。
式の有理関数恒等性は

\noindent\texttt{computations/verify\_as\_avatar\_euler.py}で
四元連立の厳密な逆行列計算として再現できる。
ここで確認したのは\emph{最小分解が完全である場合の
Euler数値帰結}であって、
分解の完全性やAS--Gorenstein性を
母関数だけで証明したのではない。

まず有限区間で局所双対の消滅を確認する。
$P_j$ を頂点 $b$ 未満の部分加群で割った有限長商を
$U_{j,b}$ とする。頂点によるフィルトレーションの
因子は $S_\ell^{\oplus\dim A_{j\ell}}$（$b\le\ell\le j$）である。
\eqref{eq:qa-as-two-gorenstein} を射影 $P_i$ の成分で読むと、
$\Ext_A^q(S_\ell,P_i)$ は $(q,\ell)=(3,i+4)$ の場合だけ
$k$ となる。従ってフィルトレーションの長完全列から
$\Ext_A^q(U_{j,b},P_i)$ は $q\ne3$ では零であり、
$q=3$ の次元は $b\le i+4\le j$ なら $\dim A_{j,i+4}$、
それ以外は零である。
さらに切断を下げるとき、追加される因子の
$\Ext_A^3(-,P_i)$ は零なので、$b\le i+4$ 以後の
遷移写像は同型となる。
従ってこの計算はフィルトレーションの Euler 特性だけでなく、
導来 torsion を計算する直極限の各頂点での消滅を与える。

加群としての同定には、有限長加群 $N$ を
頂点 $i+4$ 未満の部分と $i+4$ 以下の部分で切る。
その二つの外側の部分の $\Ext_A^\bullet(-,P_i)$ は上の計算で零である。
従って包含と商の誘導する同型により
$\Ext_A^3(N,P_i)\simeq\Hom_k(N_{i+4},k)$ を得る。
この同型は $\Ext_A^3(S_{i+4},P_i)\simeq k$ の同定を固定すれば
$N$ に自然である。
特に $U_{j,b}$ に適用した同型は $j$ を変える道の作用と可換であり、
直極限で $R\tau(P_i)\simeq D(Ae_{i+4})[-3]$ となる。
この計算を tails に移すため、$J_{i,b}$ を
$P_i$ の頂点 $b$ 未満の部分加群とする。
$P_i/J_{i,b}$ は有限長で、各単純の長さ三の解消から
有限生成射影による長さ高々三の解消を持つ。
$J_{i,b}$ も有限生成射影による長さ高々二の解消を持つ。
さらに上の添字規約では、$J_{i,b}$ の射影解消を
頂点 $b$ 未満の射影の直和で選べる。
従って下に有界な torsion 加群 $N$ に対して、
十分小さい $b$ で $\Ext_A^p(J_{i,b},N)=0$ が全 $p$ で成り立つ。
一般の torsion 加群はそのような部分加群の filtered colimit
であり、有限射影解消により同じ極限消滅が成り立つ。
従って Gabriel 商の section 関手には
論文Bの \cite[式 \texttt{eq:cubic-tails-section-colimit} と
\texttt{eq:cubic-tails-derived-localization}]{QuadricModuli2026}
と同じ局所化三角を適用できる。
ここで用いたのは有限区間の射影解消と torsion 上の消滅である。

この三角を $P_j$ に適用して頂点 $i$ で評価すると、
\[
 \Hom(E_i,E_j)=A_{ji},\qquad
 \Ext^2(E_i,E_j)=D(A_{i,j+4}),\qquad
 \Ext^q(E_i,E_j)=0\quad(q>0,\ q\ne2)
\]
を得る。特に長さ四の窓の全対で高次 Ext が消える。
最小分解を$\pi$で移した四項完全列を
正負双方へ繰り返すと連続四対象は全$E_j$を生成する。
有限quiverの最小関係はsourceからsinkへ向かい、
関係同士のoverlapを持たない。
よってその有限End代数の単純分解の長さは高々二。
$E_j$ を末項とする四項窓の直和を $T$ とする。
$B=\End(T)$ は有限次元かつ有限大域次元なので、
導来 Morita 同値で Serre 関手は $DB$ との導来テンソルに対応する。
この窓の各 $E_i$ について $\Hom(E_i,E_{j-4})=0$ であり、
先の局所双対から
\[
 \RHom(E_i,E_{j-4}[2])\simeq D\Hom(E_j,E_i)[0]
\]
となる。この式は $i$ の窓内の射に自然である。
実際、$\Ext^2$ の同定は有限切断の $i$ 番射影への
道の作用を双対化したものである。
従って $\RHom(T,E_{j-4}[2])$ は、ただ一つのコホモロジーとして
右 $B$ 加群 $D\Hom(E_j,T)$ を持つ。
これは $S_B(E_j)$ に対応する加群であり、
Morita 同値により $S_B(E_j)\simeq E_{j-4}[2]$ を得る。
窓を変えても生成する圏は同じなので、その Serre 関手は同一である。
よって全 $j$ で $S_B^{-1}[2](E_j)\simeq E_{j+4}$ となる。
有限窓の導来圏でこの自己同値を反復し、合成を通常の射の合成に
取れば、全整数の代数の復元 \eqref{eq:qa-as-serre-completion} を得る。
これは周期同型の選択を分類データへ加える主張ではない。
\end{proof}

\subsection{AS条件からavatarを復元する}
\label{subsec:qa-as-to-avatar}

\begin{lemma}[AS条件が強制する階数]
\label{lem:qa-as-required-ranks}
rotated型の有限quiverを
$0\overset{2}{\longrightarrow}2
 \overset{4}{\longrightarrow}3$、
$1\overset{2}{\longrightarrow}2$と書き、
$R_0,R_1$を二つの二次元関係空間とする。
$V=H^0(2\to3)$とし、
二つの関係評価行列を$M_0(v),M_1(v)$、
随伴関係の$4\times4$行列を
$T(c,d),\widetilde T(a,b)$とすると
\begin{enumerate}
\item $R_0,R_1\subset V\otimes k^2$に
 非零の分解可能テンソルはない。
\item $\det M_0(v)=\det M_1(v)=0$、
 $v\ne0$なら両行列の階数は一。
\item $c,d$がともに非零なら
 $\operatorname{rk}T(c,d)\ge3$、
 $a,b$がともに非零なら
 $\operatorname{rk}\widetilde T(a,b)\ge3$。
\end{enumerate}
\end{lemma}
\begin{proof}
まず無限の関係を有限窓から回収する。
有限窓の代数 $B$ は大域次元二であり、$R_0,R_1$ は
有向道代数の関係 ideal の最小生成である。
未変形の Keller の relation completion
\cite[Theorem 6.10]{KellerDeformedCY2011} を用いると、
$T_B\RHom_{B^e}(B,B^e)[2]$ の零次コホモロジーは
有限窓の関係に双対戻り矢を加えた通常の Jacobian 代数である。
ここでは元の矢のテンソル次数を零、戻り矢を一とする。
この次数は一窓をまたぐ回数であり、関係の巡回微分は
両辺で同じ回数を保つ。
逆 Serre 双加群 $\RHom_{B^e}(B,B^e)[2]$ は
$S_B^{-1}[2]$ を表すので、その $m$ 次テンソル成分の
頂点対 $(i,j)$ は前補題により $\Hom(E_i,E_{j+4m})$ になる。
積は双加群のテンソル積、すなわち添字を移した射の合成である。
従ってこの Jacobian 代数を四頂点ずつ展開すると元の $A$ を得る。
以下の周期商で使う関係はこの通常の道 ideal の全生成である。

四周期の有限quiverのsourceとsinkを連結し、
各ブロック内の元の関係とその随伴微分関係を
すべて課した商を考える。
一ブロックのboundary cornerを$W$とすると、
$m$ブロックのcornerは$W^{\otimes m}$である。
境界間の道は一意に分割される。
以下では各場合ごとに境界を選び、
境界をまたぐ関係の像が零であることを確認する。
そのとき残る関係は各ブロックの一因子だけに作用するので、
ブロックごとの商のテンソル積が全境界 corner となる。
$\dim W\ge2$なら成分次元は少なくとも$2^m$となり、
前補題の$(2m+1)^2$に反する。

使う関係を先に書く。
$v_0,\dots,v_3$を$V$の基底とし、
$r_{0,\alpha}=\sum_{\ell,i}A_{\ell i\alpha}v_\ell a_i$、
$r_{1,\beta}=\sum_{\ell,j}B_{\ell j\beta}v_\ell b_j$
（$\alpha,\beta=0,1$）とする。
次周期の戻り矢$c_\alpha,d_\beta$を加えると、
各$r_{0,\alpha},r_{1,\beta}$に加えて
\[
\begin{aligned}
 \sum_{\ell,\alpha}A_{\ell i\alpha}c_\alpha v_\ell&=0
 &&(i=0,1),\\
 \sum_{\ell,\beta}B_{\ell j\beta}d_\beta v_\ell&=0
 &&(j=0,1),\\
 \sum_{i,\alpha}A_{\ell i\alpha}a_i'c_\alpha+
 \sum_{j,\beta}B_{\ell j\beta}b_j'd_\beta&=0
 &&(\ell=0,1,2,3)
\end{aligned}
\tag{Adj-rel}\label{eq:qa-as-forcing-adjoint-relations}
\]
および周期移動を持つ。以下は消した矢で切った
\emph{全随伴関係の像}を含む商である。

(1) 分解可能な関係を$r_{0,0}=v_0a_0$へ正規化する。
他方の枝の頂点 $1$ とその周期移動を零にし、
$v_0=a_0=c_1=0$を課す。境界は各周期の頂点 $3$ とする。
式~\eqref{eq:qa-as-forcing-adjoint-relations}の
$c_0$係数は$v_0$か$a_0$を含み、
$c_1$係数は零となるため、境界をまたぐ $c v$ 型の関係は消える。
残る境界から次の境界への道は三本で、
$r_{0,1}'c_0$による関係は高々一本。
従って$\dim W\ge3-1=2$である。
$R_1$の場合は二枝を交換する。

(2) $M_0(v)=0$かつ$\det M_1(v)=0$とする。
$v=v_0$とし、$M_1(v_0)$の左右kernelに
沿って$b_0,d_0$を選ぶ。
$v_1=v_2=v_3=b_1=d_1=0$を課し、境界を各周期の頂点 $2$ とする。
境界をまたぐ $v a,v b$ 型の $r_0,r_1$ はこの切断で消える。
$c v,d v$ 型の関係も、$M_0(v_0)=0$ と
$d_0$ の kernel 条件により消える。
$a'c,b'd$ 型のうち $\ell=0$ の式は $b_0$ の kernel 条件から零となる。
四本の$a_i'c_\alpha$と一本の$b_0'd_0$、
合計五本の境界道に、
$\ell=1,2,3$の高々三関係だけが残る。
よって$\dim W\ge5-3=2$である。
$M_1(v)=0$の場合も同様。

(3) $\operatorname{rk}T(c,d)\le2$なら
$c=d=(1,0)$に正規化し、
$A_{\ell i0},B_{\ell j0}$の四係数列が
$\langle v_0,v_1\rangle$に入るようにできる。
境界を各周期の頂点 $3$ とし、$v_0=v_1=c_1=d_1=0$を課す。
境界をまたぐ $c v,d v$ 型の関係は係数列が
$\langle v_0,v_1\rangle$ に入るため消える。
残った$a_0'c_0,a_1'c_0,b_0'd_0,b_1'd_0$を
高々二関係で割る空間$U'$は
$\dim U'\ge2$である。
次の境界までの道空間は
\[
 W=(\langle v_2',v_3'\rangle\otimes U')/
       \langle r_{0,1}'c_0,r_{1,1}'d_0\rangle
\]
であるから$\dim W\ge2\dim U'-2\ge2$。
$c_0,d_0$を可逆元や恒等射と扱っていないので、
この次元計算は四周期の正当な次数付き商に対するもの。
各場合で$\dim W\ge2$は指数増大に矛盾する。
左最小分解に同じ議論を適用すると
$\widetilde T$の階数条件を得る。
\end{proof}

\begin{proposition}[一点表示の幾何]
\label{prop:qa-as-extraction}
任意の$A\in\mathrm{AS}_{\rm rsq}$から
\[
Y=\PP_{\PP^1}(\mathcal E)\simeq\Sigma_0\ \text{または}\ \Sigma_2,
\quad
D\hookrightarrow Y,\quad
P\in\Pic^0(D)
\]
を回収できる。曲線型では
$D\in|2h|$、$\Sigma_2$なら$D\cap S=\varnothing$、
かつ$P\ne\OO_D,\OO_D(h-2f)$となる。
初期四項は
\[
(P,\OO_D(h-2f),\OO_D(h-f),\OO_D(2h-f)).
\]
$D=Y$の場合には$\Sigma_0$の古典型一点を得る。
\end{proposition}
\begin{proof}
保持三頂点の安定表現は
$\{([b],[v])\in\PP^1\times\PP^3\mid M_1(v)b=0\}$
である。$[b]$ を固定すると、関係は $V$ の二つの一次式となる。
この二式が従属なら、$R_1$ の非零元の $b$ による縮約が零となり、
その元は $V\otimes\ker b$ に属する。
$\ker b$ は一次元なので、これは分解可能テンソルを与え、
前補題に反する。
よって普遍関係写像
$\OO(-1)^2\to V\otimes\OO_{\PP^1}$
は各点で階数二であり、単元 minor により局所分裂単射である。
その余核 $\mathcal E$ は rank 二、次数二の大域生成束であり、
$M_1(v)b=0$ の incidence は $\PP(\mathcal E)$ と
試験スキーム上の商線束の関手として同型である。
$\mathcal E$ の分裂次数は大域生成性から非負、和は二なので
$\OO(1)^2$ または $\OO\oplus\OO(2)$ となり、
$Y$ の二型を得る。
忘れた頂点を加える評価は
$\kappa:\OO_Y^2\to\OO_Y(h)^2$である。
零点では $\det M_0(v)=\det M_1(v)=0$ であり、
前補題により $M_0(v)$ の階数は一なので、
局所的に$\kappa\sim\operatorname{diag}(1,q)$。
従って$\PP_Y(\coker\kappa)=V(q)=D$は
scheme-theoretic閉埋め込みであり、
$\det\kappa$の類は$2h$。
$\Sigma_2$ の負断面 $S$ は
$\mathcal E=\OO\oplus\OO(2)\twoheadrightarrow\OO$ の商に対応する。
この商では $\OO_Y(h)|_S\simeq\OO_S$ であり、
普遍評価 $V\otimes\OO_S\to\OO_S$ は一つの固定商 $[v_*]$ となる。
従って incidence 式 $M_1(v_*)b=0$ はすべての $[b]\in\PP^1$ で成り立つ。
二つの独立な $b$ を代入すると $M_1(v_*)=0$ を得る。
もし $\det M_0(v_*)=0$ なら、両行列式が零なのに
$M_1$ が階数零となり、前補題の第二条件に反する。
従って $\det M_0$ は $S$ に沿って単元であり、
スキームとして $D\cap S=\varnothing$ である。
この議論は $\Sigma_2$ 上で $\det\kappa$ が恒等的に零となる場合も排除する。
普遍右表現の双対と余核$Q$を用いて
$P=Q^{-1}\OO_D(2h-f)$と正規化する。

点表示の移動を随伴関係
\eqref{eq:qa-as-forcing-adjoint-relations} で記述する。
既知の評価を $v,a,b$ とすると、$c,d$ は
$M_0(v),M_1(v)$ の関係空間側の核線から得る。
いずれも階数一なので核線は局所直和因子であり、
単元 minor による記述は基底変更と可換である。
$c,d$ を代入した $T(c,d)$ では、既知の $v$ が左核に入るため
行列式は零となり、前補題の第三条件によって階数は三である。
従って右核から次の組 $(a',b')$ が一意に定まる。
その一方が零なら、他方と $c$ または $d$ によって
非零の関係テンソルの縮約が零となる。
この関係は二次元矢空間の一次元部分に支持されるので
分解可能となり、前補題の第一条件に反する。
従って $a',b'$ はともに非零である。

次に $\widetilde T(a',b')$ には既知の $(c,d)$ が核を与える。
同じ階数三の議論により、新しい $v'$ が一意に定まる。
随伴関係は $M_0(v')a'=M_1(v')b'=0$ を含むので、
$(v',a',b')$ は再び元の四頂点の点表示を与える。
逆向きにも同じ核線を用いると、階数一と階数三での
一意性により元の表示に戻る。
すべての写像は単元 minor 上の核線で構成されるので
逆性はスキーム上で成立する。
この移動の自己同型を $\sigma$ とし、
矢を $L_i\to L_j$ の評価と書くために普遍右表現を双対化する。
初期の線束を $L_0=P,L_1=HF^{-2},L_2=HF^{-1},L_3=H^2F^{-1}$
と正規化すると、二つの行列の余核は $Q$ および $HF$ である。
後者は $b:B\otimes\OO\twoheadrightarrow F$ の核に $H$ を捩った
部分束が $M_1(v)$ の像になることから従う。
したがって二つの階数一行列の核は、
行列式を取るとそれぞれ $QH^{-2}$ と $FH^{-1}$ である。
その双対が $L_4L_3^{-1}$ と $L_5L_3^{-1}$ を与えるので、
$L_4=H^2P$、$L_5=H^3F^{-2}$ を得る。

続く二つの階数三の写像では、
核と余核の線束を含む完全列の型は
\[
\begin{aligned}
0&\to L_6^{-1}\to
 (L_4^{-1})^{\oplus2}\oplus (L_5^{-1})^{\oplus2}
 \xrightarrow{T(c,d)} (L_3^{-1})^{\oplus4}
 \to L_2^{-1}\to0,\\
0&\to L_7^{-1}\to (L_6^{-1})^{\oplus4}
 \to (L_4^{-1})^{\oplus2}\oplus (L_5^{-1})^{\oplus2}
 \to L_3^{-1}\to0 .
\end{aligned}
\]
最初の余核は既知の $v$、次の余核は既知の $(c,d)$ の評価である。
二次元・四次元の定数係数空間の行列式を除けば、
$L_6=L_4^2L_5^2L_2L_3^{-4}=H^3F^{-1}P^2$ と
$L_7=L_6^4L_3L_4^{-2}L_5^{-2}=H^4F^{-1}P^6$ となる。
よって移動後の偏極は
$\sigma^*F=L_6L_5^{-1}=FP^2$ と
$\sigma^*H=L_7L_6^{-1}=HP^4$ である。
曲線型では、二つの式から $\sigma^*(HF^{-2})=HF^{-2}$ となる。
$\Sigma_0$ 上で成分の類を $af+bg$ と書くと、
$HF^{-2}$ のその成分上の次数は $a-b$ である。
この値が同じで類が異なる二成分が共存すると仮定すると、
一方の二係数は他方よりともに少なくとも一大きい。
小さい方も非零の有効類なので、二成分の和のいずれかの
係数が三以上となり、$D\in|2f+2g|$ に反する。
従って $\sigma$ は各成分の数値類を保ち、
特に $F$ の多重次数を保つ。
$\Sigma_2$ では各成分が負断面を避けるため、その類は
正整数倍の $h$ であり、類の異なる二成分は $2h$ 内に共存しない。
この場合も $F$ の多重次数は保たれる。
従って $\sigma^*F=FP^2$ を各成分で比較して $\deg P=0$ を得る。
ここで次数は被約既約成分への制限で測る。
多重次数零が $\Pic^0(D)$ を与える理由は、
まず $D_{\mathrm{red}}$ の正規化に移すと、
各正規化成分の Picard 群の単位成分が次数零部分だからである。
正規化による Picard 群の核は、導手の有限スキーム上の
単元群の商として記述され、連結である。
実際、代数閉体上の有限局所代数の単元群は
剰余体の乗法群と冪零イデアルに由来する連結な群からなる。
さらに $D$ の冪零イデアル $\mathcal J$ に対し
$1+\mathcal J$ を冪で濾過すると、逐次商は加法的な準連接層
$\mathcal J^r/\mathcal J^{r+1}$ である。
曲線上でこれらの $H^2$ は零で、$H^1$ はベクトル群なので、
被約化への Picard 制限は全射で核は連結である。
従って $P$ は非被約の場合も $\Pic^0(D)$ に属する。
曲線型では $\kappa$ の行列式は非零で、
曲面上の二項解消の余核は $\iota_*Q$ である。
これを $-2h+f$ と $-h-f$ で捩ると、余核はそれぞれ
$\iota_*P^{-1}$ と $\iota_*(HF^{-2}P^{-1})$ となる。
$\Sigma_0$ では $h=f+g$ なので、一つ目の捩れの
二つの曲面項はそれぞれ $-f-2g,-g$、
二つ目は $-2f-g,-f$ であり、Künneth によってすべて acyclic である。
$\Sigma_2$ では、相対次数 $-1$ の項の順像は全次数で零であり、
$-f$ の順像は $\OO_{\PP^1}(-1)$ である。
残る $-2h+f$ については相対双対性と
$\omega_\pi=\OO_Y(-2h+2f)$ により
$R^1\pi_*\OO_Y(-2h+f)=\OO_{\PP^1}(-1)$、$\pi_*=0$ となる。
従って両曲面で
$H^\bullet(P^{-1})=H^\bullet(HF^{-2}P^{-1})=0$ を得る。
自明線束には単位切断があるので、二つの除外条件が従う。
$\det\kappa\equiv0$ なら、負断面での可逆性により
$\Sigma_2$ は起こらない。
$Y=\Sigma_0$ 上で $\kappa$ は至るところ階数一なので、
余核 $Q$ は $Y$ 上の線束であり、点表示のスキームは $Y$ 全体である。
この場合の四周期の自己同型 $\sigma$ は
$\omega_Y^{-1}\simeq H^2$ を保つ。
$\Pic(Y)\simeq\mathbb Z^2$ には二捩れがないので、
$H$ は反標準束の一意な平方根として $\sigma$ に保たれる。
従って四周期の式 $\sigma^*H=HP^4$ から $P^4=\OO_Y$ を得て、
同じ Picard 群の計算から $P=\OO_Y$ となる。
この場合の線形代数には曲線型の単射解消を使わない。
\end{proof}


\subsection{Avatarからの全AS代数の逆構成}
\label{subsec:qa-as-avatar-inverse}

\begin{lemma}[canonical truncation]
\label{lem:qa-as-canonical-truncation}
rotated square avatarから前節の$L_i$を全整数に作る。
任意の$i$でcanonical四項列を
初期四頂点の共変Yonedaへ制限すると、
頂点単純加群の最小射影分解になる。
特に三回の右mutationは
\begin{equation}
T^{-1}_{L_{i+3}}T^{-1}_{L_{i+2}}T^{-1}_{L_{i+1}}(L_i)
 \simeq L_{i+4}[-2]
\label{eq:qa-as-canonical-mutation}
\end{equation}
を満たす。
\end{lemma}
\begin{proof}
$pencil$の評価列で作った
\eqref{eq:qa-core-rsq-complex0}と
\eqref{eq:qa-core-rsq-complex1}は
Yonedaで最初の二種の最小関係と一次syzygyを与える。
後二列は四行四列の対角分解の
Cartier曲線への導来制限と双対化から来るので
Tor項を含めて完全である。
各位相で最初に現れる非零な矢空間への評価は、
pencil または四次線束の全切断空間を使うので全射である。
第一・第二位相の残る頂点では、第一 pencil の評価列を
その頂点の線束で双対化して用いる。
評価の余核は pencil trick により次数二の線束の $H^1$ に
含まれるので零となる。第三位相の後二頂点では、四次線束 $H$
に対して $a_+$ または $b_+$ の評価列を用いる。
乗法の余核は $H a_+^{-1}$ または $H b_+^{-1}$ の $H^1$ に
含まれ、これらも次数二の pencil なので零である。
第四位相の最後の頂点では、第四列を作る前の曲面上の
$0\to\OO^{\oplus4}\to\OO(a)^{\oplus2}\oplus\OO(b)^{\oplus2}
 \to\iota_*Q\to0$
の大域切断が、$H^1(\OO)=0$ により所要の全射を与える。
中間の完全性もEuler特性だけには依存しない。
canonical複体を一時的に
$0\to L_i\to M_1\to M_2\to L_{i+4}\to0$と書く。
共変Yonedaの頂点$b$における
$\Hom(M_1,L_b)\to\Hom(L_i,L_b)$のkernelの次元を
計算すると
\[
\begin{array}{c|c|ccc}
i\bmod4&\text{相対次数}&b=1&b=2&b=3\\ \hline
0&(0,0,2,6)&0&0&2\\
1&(0,2,6,8)&0&2&4\\
2&(0,4,6,6)&0&2&2\\
3&(0,2,2,4)&0&0&4
\end{array}
\]
となる。ここで $L_b$ は窓の第 $b$ 頂点の略記である。
窓内の逆向きの通常の Hom は、
命題~\ref{prop:qa-core-forward-ext-vanishing}と
$\omega_D\simeq\OO_D$ による Serre 双対で零である。
従って通常の Hom を使ったこの計算は directed Yoneda の評価と一致する。
同じ双対性により、窓の外の端についても
$\Hom_D(L_{i+4},L_b)=\Ext_D^1(L_b,L_{i+4})^\vee=0$ となる。
特に総次数が等しい対でも、前向き消滅の次数零の場合がこれを保証する。

短完全列 $0\to\operatorname{coker}(L_i\to M_1)\to
M_2\to L_{i+4}\to0$ に $\Hom_D(-,L_b)$ を適用すると、
$\Hom_D(M_2,L_b)$ から $\Hom_D(M_1,L_b)$ への写像は単射である。
像は表の kernel に含まれ、その次元は表の数と一致する。
従って窓の各頂点で中間は完全である。
頂点 $i$ では対角の $k$ が augmentation として残り、
すべての微分は異なる頂点間の成分からなるので解消は最小である。
三回のmutationは頂点1,2,3を順次消す
Yoneda truncationなので、
$D$上に実現すると最後の$L_{i+4}[-2]$のみが残る。
\end{proof}

\begin{proposition}[spherical helixによるAS正則性]
\label{prop:qa-as-helix-reconstruction}
任意の曲線型rotated square avatarから
$B=\End_{\rm dir}(L_0,L_1,L_2,L_3)$を作る。
その標準射影対象を$E_0,\ldots,E_3$とし、
$E_{i+4}=S_B^{-1}(E_i)[2]$で延長すると
\[
 A(D,P)_{ji}=\Hom(E_i,E_j)\quad(j\ge i)
\]
は$\mathrm{AS}_{\mathrm{rsq}}$に属する。
\end{proposition}
\begin{proof}
$\omega_D=\OO_D$かつ
$H^0(\OO_D)=H^1(\OO_D)=k$だから
各$L_i$は一次元spherical objectである。
全整数のcanonical完全列と
命題~\ref{prop:qa-core-forward-ext-vanishing}、
前補題はacyclic spherical helixを与える。
このacyclicityは同文献Definition~4の全$i<j$に対する
次数零以外のHom消滅である。
これを入力として、対応するexceptional helixの
前向きHomが次数零に集中することに
\cite[Theorem 5]{OkawaUedaSpherical2020}を用いる。
指定した最小分解の頂点・重複度と双対の同定には、
以下のcanonical複体とSerre traceの計算を使う。

$C$を四線束のfull dg Endとする。
$B$は前向き次数零Homと対角の恒等射を含み、
$B\to C$ は dg の実現写像である。
$H^0(C)$ の後ろ向き成分は前向き Ext 消滅の双対により零で、
$H^0(B)\to H^0(C)$ は同型となる。
その cone の唯一のコホモロジーは $H^1(C)$ である。
Serre trace による pairing は合成に関して自然であるから、
$H^1(C)=DB$ は $B$ 双加群としての同型である。
単一のコホモロジーをもつ cone は双加群の導来圏で
$DB[-1]$ に同型なので、三角形
\begin{equation}
 DB[-2]\longrightarrow B\longrightarrow C
 \longrightarrow DB[-1]
 \label{eq:qa-as-dg-triangle}
\end{equation}
がある。
有限代数$B$のSerre関手を用いると
realizationとその右随伴の間に
$S_B[-2]\to\mathrm{id}\to\iota^*\iota_*$を得る。
前補題とこの三角形から
$E_i$のrealizationは全整数で$L_i$となり、
\begin{equation}
\RHom(E_i,E_{j-4})\longrightarrow
\RHom(E_i,E_j)\longrightarrow
\RHom_D(L_i,L_j)
\label{eq:qa-as-rhom-comparison}
\end{equation}
が成立する。
$0<j-i<4$ では $j-4<i$ であり、
$i$ と $j-4$ は同じ窓に入るので最初の項は零である。
従って実現は各窓の前向き Hom 上で合成を保つ同型を与え、
前補題の窓内の計算を例外対象の窓に移せる。
$j-i=4$ の最初の項は対角の $k$ である。
$j-i>4$ では最初の項に距離についての帰納法を使うと、
最後の項も前向き Ext 消滅により次数零に集中するため、
中央の前向き Hom も同じ性質を持つ。
この比較の向きは \cite[式 (9), (30)--(32)]{OkawaUedaSpherical2020}
の $E_{j-4}=S_B(E_j)[-2]$ と一致する。

四項複体を得るため、まず窓 $[i-4,i-1]$ を取る。
前補題の共変 Yoneda 解消の三つの射影項をこの窓に戻すと、
$E_{i-4}$ から始まる逐次右 mutation の複体になる。
実際、三項の 畳み込み は後三対象への導来 Hom が零であり、
$E_{i-4}$ への射の cone は後三対象で生成される。
従って窓の半直交分解によって一意に定まる右 mutation である。
窓の Serre 関手による同定から、その最後の項を $E_i$ とした
四項複体 $K_i$ の全複体は零である。
$K_i$ は次数 $-3,-2,-1,0$ に置き、両端は
$E_{i-4},E_i$ とする。中間項の重複度は前補題の解消から
定義~
ef{def:qa-as-square-rotated-minimal} の rotated 型の表となる。

$RHom(E_j,K_i)$ の項別スペクトル系列で次数零の行を取ると、
右射影複体の頂点 $j$ における評価になる。
$j<i$ では前向き消滅と幅四の窓の半直交性により
他の行は零なので、この行は完全である。
$j=i$ では唯一の例外が左端の
$Ext^2(E_i,E_{i-4})=k$ である。実際、
$S_B(E_i)=E_{i-4}[2]$ なのでこれは Serre 双対から従う。
全複体の消滅により、この $k$ から次数零の行の右端の
コホモロジーへの $d_3$ が同型である。
$j>i$ での評価は directed の定義により零である。
従って右射影複体の唯一のコホモロジーは次数零の $S_i$ である。

双対複体を左頂点 $j$ で評価するには $RHom(K_i,E_j)$ を用いる。
$j>i-4$ では同じ前向き消滅と半直交性により次数零の行が完全であり、
$j<i-4$ では directed の定義により評価は零である。
$j=i-4$ では、次数 $0,1,2,3$ のうち最初の項から来る
$Ext^2(E_i,E_{i-4})=k$ だけが次数零の行の外に残る。
全複体が零なので $d_3$ はこれを次数零の行の
次数三のコホモロジーに同定する。
従って $Ext^3_A(S_i,A)={}_ell S_{i-4}$ であり、
他の次数の Ext は零である。
異なる頂点間の成分しか現れないので右解消は最小である。
\end{proof}


\begin{proposition}[有限Endと全代数の相互逆性]
\label{prop:qa-as-inverse-foundations}
$A\in\mathrm{AS}_{\rm rsq}$から取り出した
$(Y,D,P)$の初期四頂点Endは
\[
A|_{\{0,1,2,3\}}\simeq
\End_{\rm dir}
(P,\OO_D(h-2f),\OO_D(h-f),\OO_D(2h-f))
\]
であり、古典型は$Y$上の可換基準collectionである。
逆にavatarの有限Endから構成した代数の
点表現と余核線束は元の$(Y,D,P)$を回収する。
\end{proposition}
\begin{proof}
後者は命題~\ref{prop:qa-core-directed-foundations}の
二行二列評価行列の局所標準形
$\operatorname{diag}(1,q)$から
$\PP_Y(\coker\kappa)=D$を回収し、
忘れられた線束の正規化で$P$を得ることによる。
前者は、AS代数から作った
\[
0\to\OO_Y^{\oplus2}\xrightarrow{\kappa}
\OO_Y(h)^{\oplus2}\to\iota_*Q\to0
\]
を $\OO_Y(-h)$ で捩ると、$\OO_Y(-h)$ の全コホモロジー消滅により
$H^0(QH^{-1})$ は $0\to2$ の二矢空間と同型になる。
元の列の大域切断は
\[
 0\longrightarrow R_0\longrightarrow
 H^0(Y,\OO_Y(h))\otimes H^0(D,QH^{-1})
 \longrightarrow H^0(D,Q)\longrightarrow0
\]
である。最初の射は普遍評価行列の係数なので、
元の二関係の包含そのものである。
最後の射は可逆層の切断の乗法であり、
$0\to3$ の合成を同定する。

保持した $1,2,3$ の三頂点では、$f,h,h+f$ から $2h$ を引いた
$f-2h,-h,f-h$ のすべてが acyclic なので、
それぞれの切断空間は $Y$ から $D$ への制限で変わらない。
$Y$ 上の $f$ と $h$ の乗法は、
$\PP^1$ 上で $\mathcal E$ の評価列を $\OO(1)$ で捩って得る
全射 $V\otimes H^0(\OO(1))\to H^0(\mathcal E(1))$ であり、
核は $R_1$ である。
従って $1\to3$ の合成も元の二関係で割った積と一致する。
残る
$0\to1$のHomの消滅は
$H^\bullet(\OO_D(h-2f)P^{-1})=0$である。
古典型では $P=\OO_Y$ なので $Q=\OO_Y(2h-f)$ である。
一定階数一の $\kappa$ の像を $I$ とすると、
$0\to I\to\OO_Y(h)^2\to Q\to0$ の行列式から
$I=\OO_Y(f)$、さらに $\ker\kappa=\OO_Y(-f)$ となる。
$\OO_Y^2\twoheadrightarrow I$ の二切断は、
$\OO_Y(f)$ の全生成 pencil の基底をなす。
また $I\hookrightarrow\OO_Y(h)^2$ は、
$\OO_Y(h-f)$ の全生成 pencil の基底を双対化して得る包含である。
いずれも二切断が線形従属なら、正次数の線束を
一切断で生成することになるため不可能である。
この分解で $H^0(\OO_Y^2)\to H^0(I)$ は同型であり、
$H^1(I)=0$ なので、$0\to3$ の積の核は元の $R_0$ となる。
$1\to3$ の積は曲面上の $\mathcal E$ の評価列から同様に $R_1$ を回収する。
従って古典型でも全合成を含めて可換基準 collection を得る。
全$\mathbb Z$-代数は
補題~\ref{lem:qa-as-serre-unique-recovery}が決めるから
有限Endの同型で止まらない。
\end{proof}

\begin{theorem}[squareとrotated squareのavatar分類]
\label{thm:qa-as-avatar-bijection}
定義~\ref{def:qa-as-square-rotated-minimal}の
最小分解型を持つAS正則$\mathbb Z$-代数について
\begin{equation}
\mathrm{AS}_{\rm rsq}/\cong\simeq
 \mathrm{Av}_{\rm rsq}/\cong,\qquad
\mathrm{AS}_{\rm sq}/\cong\simeq
 \mathrm{Av}_{\rm sq}/\cong .
\label{eq:qa-as-avatar-two-bijections}
\end{equation}
$\mathrm{Av}_{\rm rsq}$は
$D\in|2h|$（$\Sigma_2$では$D\cap S=\varnothing$）と
$P\in\Pic^0(D)\setminus\{\OO_D,\OO_D(h-2f)\}$、
$\mathrm{Av}_{\rm sq}$は順序付きrulingを持つ
$D\subset\PP^1\times\PP^1$の$(2,2)$反標準因子と
$P\in\Pic^0(D)\setminus\{\OO_D\}$からなる。
各々に一つの古典型を含める。
曲線の滑らかさ・被約性・既約性は課さない。
逆構成の有限foundationは
\[
B_{\rm rsq}=\End_{\rm dir}(P,HF^{-2},HF^{-1},H^2F^{-1}),
\qquad
B_{\rm sq}=\End_{\rm dir}(P,F,G,FG).
\]
\end{theorem}
\begin{proof}
rotated型では命題~\ref{prop:qa-as-extraction}が必要性、
命題~\ref{prop:qa-as-helix-reconstruction}が
任意avatarのAS実現、
命題~\ref{prop:qa-as-inverse-foundations}と
補題~\ref{lem:qa-as-serre-unique-recovery}が
両構成の相互逆性を与える。
代数同型は点表現・普遍線束を保ち、
avatar同型はHomと合成を保つので、
頂点を保つ代数同型類の全単射となる。

square型代数を一つ進めた
$A^{[1]}_{ji}=A_{j+1,i+1}$は
表からrotated型であり逆shiftも存在する。
一方定理~\ref{thm:qa-core-avatar-square-rotated}は
二つのavatarの同型類を相互変換する。
rotatedの一つ前の四項を曲線上の共通線束
$FP^{-1}$で正規化すると
$(L_{-1},L_0,L_1,L_2)\otimes FP^{-1}
=(P,F,G,FG)$となる。
これは曲線上の操作であって、
非可換圏上にtensor自己同値を仮定しない。
全代数のSerre復元もshiftと可換なので
square側の全単射が従う。
とくに$P=FG^{-1}$や$P=GF^{-1}$を
squareの許容条件から排除してはならない。
この共鳴はrotatedの$\Sigma_2$への移行を
記録し得るがAS正則性を破らない。
\end{proof}

\begin{proposition}[左右Noether性]
\label{prop:qa-as-avatar-noetherian}
上記二型の全代数は左右Noetherianである。
曲線型には四次の正規正則family
$g_i\in A_{i+4,i}$があり、その商は
$D$のtwisted section algebraを与える。
\end{proposition}
\begin{proof}
式~\eqref{eq:qa-as-rhom-comparison}の次数零から
$j>i$について
\[
0\to A_{j-4,i}\xrightarrow{\,g\,}A_{ji}
\to H^0(D,L_jL_i^{-1})\to0
\]
を得る。$t=S_B^{-1}[2]$ とし、三角形の最初の自然変換を
$\eta:t^{-1}\to\operatorname{id}$ と書く。
$E_i\to E_j$ の射 $a$ に対して自然性は
\[
 a\,\eta_{E_i}=\eta_{E_j}\,t^{-1}(a)
\]
を与える。$t^{-1}$ は各 Hom で同型なので、
$\eta$ を前から掛けた像と後ろから掛けた像は一致する。
上の短完全列から後合成による乗法は単射であり、
この式を使えば前合成による乗法も単射となる。
また自然性を $\eta_{E_j}$ 自身に適用すると
$\eta_{E_j}t^{-1}(\eta_{E_j})=
\eta_{E_j}\eta_{E_{j-4}}$ となる。
後合成の単射性により $t^{-1}(\eta_{E_j})=\eta_{E_{j-4}}$
である。従って $g_i=\eta_{E_{i+4}}$ は四周期の
折り畳みと整合する両側正規正則の族となる。
偏極の移動公式に加えて、正規化した初期窓では
$P=(L_0L_1^{-1})HF^{-2}$ である。
この式を次窓に適用すると
$\sigma^*P=(L_4L_5^{-1})(\sigma^*H)(\sigma^*F)^{-2}=P$
を得る。したがって周期式~\eqref{eq:qa-core-rsq-periodic-L}に
$\sigma^*F=FP^2$、$\sigma^*H=HP^4$、$\sigma^*P=P$ を代入すれば、
全整数で $L_{i+4}=\sigma^*L_i\otimes\OO_D(2h)$ となる。
$2h$ は $\Sigma_2$ でも $D\cap S=\varnothing$ により
各成分上正次数なので $D$ 上 ample である。
曲線上の数値類は各既約成分の次数で測られ、
$\sigma$ はこれらを有限置換する。
従って数値類への作用は準冪単であり、
\cite{Keeler2003,Keeler2018Corrigendum} の判定で
この偏極は $\sigma$-ample かつ $\sigma^{-1}$-ample となる。
twisted section algebra は左右 Noetherian である。
四周期をまとめた代数を $J$ とし、$g$ を四位相における
$g_i$ の和とする。$J/gJ$ の対角四隅は
$\OO_D(2h)$ のねじれ切断環である。
異なる位相の隅は、それぞれの対角隅上で
可逆層の切断加群となる。両方向の ampleness により
これらは左右で有限生成なので、$J/gJ$ は対角環の直積上で
左右有限加群となり、左右 Noetherian である。

正規正則性により、$g$-adic 随伴次数環は
$J/gJ$ 上の skew polynomial 環となる。
実際、各 $g^mJ/g^{m+1}J$ は $J/gJ$ と同定でき、
乗法は $g$ を移す正規自己同型でねじれる。
従ってこの随伴次数環は左右 Noetherian である。
内部次数で斉次な右 ideal について、その $g$-adic 初項の
ideal の有限斉次生成系を持ち上げる。
元の斉次元から初項を消すたびに $g$-adic 位数は増え、
$g$ の内部次数が正なので、その元の次数を固定すれば
有限回で停止する。よってこれらの持ち上げは斉次 ideal を生成する。
一般の右 ideal には内部次数の最高次項を使い、
次数を下げる帰納法で有限生成性を得る。左 ideal も同様である。
最後に $A$ の各射影行・列の斉次部分加群は、
折り畳み後に位相の冪等元を保つ斉次部分加群として扱える。
斉次生成元を元の添字へ戻せるので、$A$ も左右 Noetherian である。
古典型は通常の可換曲面のNoether性に帰着する。
\end{proof}

\begin{proposition}[Noether性からcoherentな導来同値へ]
\label{prop:qa-as-avatar-coherent-morita}
前定理のsquare／rotated square型の代数$A$について
$X=\Qgr A$、$\coh X=\qgr A$と置き、
$T=E_0\oplus E_1\oplus E_2\oplus E_3$、
$B=\End(T)$とする。
すると
\[
 \RHom_X(T,-):
 D^b(\coh X)\xrightarrow{\ \sim\ }D^b(\mathrm{mod}\,B)
\]
は三角圏同値である。
さらに$\coh X$の完全自己同値$\nu$があり、
$\nu E_i\simeq E_{i-4}$、
$D^b(\coh X)$のSerre関手は
$\nu[2]$と同型になる。
これは\emph{導来}Morita同値であり、
$\coh X\simeq\mathrm{mod}\,B$という
Abel圏同値を主張しない。
\end{proposition}
\begin{proof}
命題~\ref{prop:qa-as-avatar-noetherian}より
$A$は左右Noetherianである。
有限生成graded加群の有限表示のsyzygyに、
左右の最小単純分解と$\operatorname{Tor}$の検出を繰り返すと、
射影分解は長さ三以内で停止する。
実際、さらに非零のsyzygyがあれば、
左単純加群との$\operatorname{Tor}_{>3}$に非零の項を持つが、
左単純の長さ三の射影分解に矛盾する。
従って各連接対象は有限個の$E_i$のconeから生成される。
補題~\ref{lem:qa-as-serre-unique-recovery}の
四周期完全列により、これらの$E_i$は
$T$が生成する厚部分圏に入る。
$T$はstrong exceptional tiltingであり、
そのEndは有限大域次元二の$B$である。
これで導来Morita同値が$\coh X$の
有界導来圏全体へ延長される。

四周期のSerre同型
$E_{i+4}\simeq S_B^{-1}(E_i)[2]$を
全整数で自然化した$\mathbb Z$-代数のtransportは、
graded加群のexactな四頂点ずらしと可換で、
torsionを保つ。
従って商圏の完全自己同値$\nu$を与え、
$\nu E_i=E_{i-4}$となる。
有限tilting generatorとその全Endにおいて
$\nu[2]$は$B$のSerre関手と一致する。
生成性からSerre関手として同型となる。
\end{proof}

\begin{remark}[一般定理への依存と未証明の比較]
\label{rem:qa-as-classification-scope}
本節では元PDFの体上分類の
必要性・逆構成・全代数の相互逆性を移管した。
spherical helix実現の一般定理
\cite[Theorem 5]{OkawaUedaSpherical2020}と
twisted section algebraのNoether性は外部入力である。
これはAS代数の同型類の対応であり、
別に定義された内在的O-marked$\HHstack$との
pointed同値や任意の非被約基底上のAS相対分類を
意味しない。
\end{remark}

% !TeX root = ../nc-quadric-avatar.tex
% 2026-11-22 / minimal ten-arrow AS type, mutation and twelve-arrow status
\section{最小十本矢型の独立AS分類と十二本矢表示}
\label{sec:qa-as10-mutation-classification}

前節のsquare AS代数から一頂点mutationを行うと、
重み四relation potentialのreductionにより
八本矢cubic型または最小十本矢型が現れる。
ここで矢の本数だけからAS正則性を結論してはならない。
元草稿\cite[第VIII部，第X部]{AvatarVSeven2026}の
Ginzburg次数零集中・独立逆構成・評価/余評価の証明を
本稿のavatar分類に接続する。

\subsection{Mutation像によらない最小十本矢AS条件}
\label{subsec:qa-as10-minimal}

\begin{definition}[最小十本矢AS型]
\label{def:qa-as10-minimal}
$i$が奇数なら$\varepsilon_i=1$、
偶数なら$\varepsilon_i=0$とする。
positively directedで局所有限、
$A_{ii}=k$の$\mathbb Z$-代数が
$\mathrm{AS}_{10}^{\mathrm{min}}$に属するとは、
全$i$の右単純$S_i$に対して
\begin{equation}
0\to P_{i-4}\to
 P_{i-3}^{\oplus2}\oplus P_{i-2}^{\oplus\varepsilon_i}
 \to P_{i-1}^{\oplus2}\oplus P_{i-2}^{\oplus\varepsilon_i}
 \to P_i\to S_i\to0
\label{eq:qa-as10-minimal-resolution}
\end{equation}
が最小であり、さらに
$\Ext_A^q(S_i,A)=0$ ($q\ne3$)、
$\Ext_A^3(S_i,A)\simeq {}_\ell S_{i-4}$
を満たすことをいう。
同型は全整数頂点を保つ。
QP、mutation元、周期同型と分解の基底は
定義の標識ではない。
\end{definition}

\begin{lemma}[Hilbert関数と有限quiver]
\label{lem:qa-as10-hilbert-foundation}
$A\in\mathrm{AS}_{10}^{\mathrm{min}}$なら
\[
\dim A_{i+2r,i}=(r+1)^2,\qquad
\dim A_{i+2r+1,i}=(r+1)(r+2).
\]
四頂点のfinite directed algebraは大域次元二で、
最小quiverに
$0\to1\to2\to3$の各二矢と
追加の$x:1\to3$を持ち、
最小関係は$1\to3$の一二次関係と
$0\to3$の二三次関係である。
この有限foundationは全$\mathbb Z$-代数を
Serre関手で一意に回収する。
\end{lemma}
\begin{proof}
\eqref{eq:qa-as10-minimal-resolution}の
二つの$P_{i-2}$項はEuler漸化式で相殺され、
$h_n-2h_{n-1}+2h_{n-3}-h_{n-4}=0$、
$h_0=1$、$h_n=0$ ($n<0$)を得る。
生成関数は
$1/((1-t)^2(1-t^2))$であり、表示の次元が従う。
最小分解の前二段を四頂点へ制限すると、
列挙した矢と三つの関係が現れ、
関係同士の次のoverlapはないため
大域次元二である。
左右AS双対性による局所コホモロジーの証明は
補題~\ref{lem:qa-as-serre-unique-recovery}と同じで、
連続四項のSerre関手による
$E_{i+4}=S_B^{-1}(E_i)[2]$が全代数を回収する。
\end{proof}

\subsection{Ginzburg次数零集中によるAS条件}
\label{subsec:qa-as10-ginzburg}

\begin{lemma}[relation QPの次数零集中]
\label{lem:qa-as10-ginzburg-concentration}
square型、cubic型または
$\mathrm{AS}_{10}^{\mathrm{min}}$の有限foundationから
最小関係の双対戻り矢で作るrelation QPを$(Q,W)$とする。
非完備およびarrow-adic完備Ginzburg dg代数について
\[
 H^j\Gamma(Q,W)=H^j\widehat\Gamma(Q,W)=0
 \qquad(j\ne0)
\]
が成立する。
\end{lemma}
\begin{proof}
有限foundationの大域次元は二なので、
未変形の3-Calabi--Yau completion
$\Pi_3(B)=T_B\RHom_{B^e}(B,B^e)[2]$を考える。
\cite[Theorem 6.10]{KellerDeformedCY2011}の
未変形の場合はrelation Ginzburg dg代数と
$\Pi_3(B)$の準同型を与える。
その$m$番目のtensor成分は
$S_B^{-m}[2m]$の全周期Homとして計算できる。
$m=0$では有限foundationのtilting性、
$m\ge1$では
前節~\ref{lem:qa-as-serre-unique-recovery}または
前補題の前向きExt消滅によって
次数零だけに集中する。
従って直和も次数零集中を満たす。
ここで用いたのは未変形completionであり、
任意のdeformation parameterのCY性ではない。

元の矢の内部次数は$1,2$、
反対矢は$4-d$、頂点loopは$4$である。
全dg生成元の内部次数は正で、
各内部次数成分は有限次元複体。
完備化は内部次数の直積であり、
体上の直積は完全なので
各次数の集中は完備版でも保たれる。
逆向きも各内部次数成分を取れば成立する。
\end{proof}

\begin{lemma}[mutation後の最小AS分解]
\label{lem:qa-as10-keller-yang}
前補題のQPを、loopとoriented $2$-cycleを通らない
指定のmutation頂点で変換すると、
reduced QPのGinzburg dg代数も次数零に集中する。
その正次数付きJacobian algebraの頂点単純加群は
長さ三の完全な最小自己双対分解を持ち、
AS--Gorenstein条件を満たす。
\end{lemma}
\begin{proof}
\cite[Theorem 6.2]{KellerYang2011}を用いる。
必要なのはGinzburgの次数零集中、
mutation頂点にloopがないこと、および
局所的な2-cycle除外条件である。
squareの指定中間頂点にloop・2-cycleはなく、
最小十本矢型の2-cycleは別の二頂点間だけにあり、
逆mutation頂点を通らない。
前補題が第一の入力を保証する。
mutation後の完備Ginzburgの次数零集中が得られ、
正の内部次数付けで非完備版へ降りる。
Ginzburgの頂点単純の標準semifree解消を
$H^0\Gamma$に移すと、
恒等射、入る矢、最小関係、四次最上項からなる
Jacobian complexを得る。
次数零集中によりこれは\emph{実際に完全}である。
双対複体には最上次数の左単純だけが残る。
reduced QPでは微分はすべて正次数だから
最小性も成り立つ。
単なるdg導来同値から通常代数のAS正則性を
推論したのではない。
\end{proof}


% !TeX root = ../nc-quadric-avatar.tex
% 2026-11-22 / cubic reachability and derived obstruction
\section{二重壁でのcubic型への到達不能性}
\label{sec:qa-as-cubic-obstruction}

二つの乗法がともに階数三となることだけでは、
他のmutation列や別のexceptional foundationによる
cubic型への到達を排除できない。
本節ではこの強い否定方向を
\cite[第19--21節]{AvatarVSeven2026}から
必要な数値・層論・平滑化の証明を保って統合する。
分類対象は引き続き全整数頂点を保つAS代数、
比較対象は構造対象を指定しない商Abel圏・導来圏である。

\subsection{有限foundationでのcubic候補の条件}
\label{subsec:qa-as-cubic-candidate}

$A\in\mathrm{AS}_{\mathrm{sq}}$の有限foundationを
$B=\End(E_0,E_1,E_2,E_3)$、
$\mathcal D=\operatorname{per}B\simeq D^b(\qgr A)$とする。
そのEuler行列は
\[
G_{\mathrm{sq}}=
\begin{pmatrix}
1&2&2&4\\0&1&0&2\\0&0&1&2\\0&0&0&1
\end{pmatrix}.
\]
\begin{definition}[cubic候補foundation]
\label{def:qa-as-cubic-candidate}
$\mathcal D$の四対象$(Q_0,Q_1,Q_2,Q_3)$が
cubic候補foundationであるとは、
strong exceptionalであり、
その$K_0(\mathcal D)$の類が整数基底をなし、
Euler行列が
\[
G_c=
\begin{pmatrix}
1&2&4&6\\0&1&2&4\\0&0&1&2\\0&0&0&1
\end{pmatrix}
\]
で、二つの隣接合成
\begin{equation}
\Hom(Q_{i+1},Q_{i+2})\otimes\Hom(Q_i,Q_{i+1})
\xrightarrow{\ \sim\ }\Hom(Q_i,Q_{i+2})
\quad(i=0,1)
\label{eq:qa-as-cubic-two-compositions}
\end{equation}
が同型であることをいう。
fullnessは別途仮定しない。
任意のcubic AS代数の標準四対象はこの条件を満たす。
\end{definition}

\begin{lemma}[cubic候補のrankと数値線束型]
\label{lem:qa-as-cubic-candidate-numerical}
cubic候補$(Q_0,\ldots,Q_3)$が存在すれば
共通の$[1]$ shiftの後、全$Q_i$の数値rankは一、
かつ適当な$a,b\in\mathbb Z$に対して
\[
c_1(Q_i)=(a,b),\ (a+1,b),\
(a+1,b+1),\ (a+2,b+1)
\]
または二つのrulingを交換した順序を取る。
数値Chern座標を$(r,a,b,s)$と書けば
各$Q_i$は$s=ab$を満たす。
\end{lemma}
\begin{proof}
squareの標準基底に対応する数値座標を
\[
[E_0]=(1,0,0,0),\quad [E_1]=(1,1,0,0),\quad
[E_2]=(1,0,1,0),\quad [E_3]=(1,1,1,1)
\]
とする。$x^2=y^2=0,\ xy=1$の形式的Chern記号で
Serre作用は
\[
S_*(r,a,b,s)=(r,a-2r,b-2r,s-2a-2b+4r)
\]
であり、
\[
(S_*-1)^2v=8r(v)\delta,\qquad
\delta=(0,0,0,1).
\]
cubicのEuler行列$G_c$から計算した
同じ作用の二乗は、四つの候補基底ベクトルを
同一のprimitiveな交代結合
$[Q_0]-[Q_1]-[Q_2]+[Q_3]$へ送る。
従って各rankは共通整数$r$であり、
この結合が$K_0$のprimitive元なので$r=\pm1$。
共通shiftで$r=1$とする。
自己Euler値は$1+2s-2ab=1$なので$s=ab$。
二つのrank一数値類のEuler値は
$(1+a_j-a_i)(1+b_j-b_i)$。
隣接方向のEuler値は二、逆向きは零だから、
差は$(1,0)$または$(0,1)$。
二段隔てたEuler値四により差は$(1,1)$となるので、
隣接差は交互となり主張を得る。
\end{proof}

\begin{lemma}[曲線への制限とdeterminant]
\label{lem:qa-as-cubic-restriction-determinant}
$(D,F,G,P)$をsquare avatarとし、
$\rho:\operatorname{per}B\to\Perf(D)$を
四線束$(P,F,G,FG)$へのrealizationとする。
任意の$X,Y\in\mathcal D$に自然な三角形
\begin{equation}
\RHom(Y,X)^\vee[-2]\longrightarrow
\RHom(X,Y)\longrightarrow
\RHom_D(\rho X,\rho Y)
\label{eq:qa-as-cubic-spherical-restriction}
\end{equation}
がある。
数値類$[X]=(r,a,b,s)$に対し
\begin{equation}
\det(\rho X)=F^aG^bP^{\,r-a-b+s}.
\label{eq:qa-as-cubic-determinant}
\end{equation}
$D$が滑らかで$X$がrank一のexceptional対象なら
$\rho X$は次数零のshiftを除き線束である。
\end{lemma}
\begin{proof}
前節のdg双加群三角形
\eqref{eq:qa-as-dg-triangle}とSerre双対性を
$\RHom(X,-)$へ適用すると
\eqref{eq:qa-as-cubic-spherical-restriction}を得る。
$X=Y$がexceptionalなら、
右辺の自己$\RHom$は$k\oplus k[-1]$となる。
$D$が滑らかなら$\coh D$は遺伝的で、
有界複体はコホモロジー層のshiftの直和になる。
Endが$k$でrankが一だから、
制限は一つのrank-one indecomposable vector bundleの
shift、すなわち線束のshiftである。
（torsionとvector bundleの拡大は分裂する。）
rank一の符号を共通shiftで正に揃える。

数値基底の座標を解くと
\[
[X]=(r-a-b+s)[E_0]+(a-s)[E_1]
 +(b-s)[E_2]+s[E_3].
\]
perfect複体のdeterminantの加法性と
$\rho(E_0,E_1,E_2,E_3)=(P,F,G,FG)$から
\eqref{eq:qa-as-cubic-determinant}を得る。
\end{proof}

\begin{proposition}[滑らかな二重壁では任意のcubic候補が不存在]
\label{prop:qa-as-cubic-smooth-no-candidate}
$D$を滑らかな楕円曲線とし、
$P^{\otimes2}\simeq\OO_D$、
$P\not\simeq\OO_D$、$G\simeq FP$とする。
$\mathcal D=\operatorname{per}B(D,F,G,P)$には
cubic候補foundationは存在しない。
\end{proposition}
\begin{proof}
候補を仮定する。
補題~\ref{lem:qa-as-cubic-candidate-numerical}により
全rankを一とし、補題
~\ref{lem:qa-as-cubic-restriction-determinant}により
$\rho Q_i=N_i[2m_i]$と書く。
隣接する$N_i$の次数差は二で、
$\Ext$の三角形と正次数線束の$H^1=0$から
前向きHomが次数零に集中するため、
各$m_i$は共通である。
共通shiftを除き、$N_i$を四線束とする。

\eqref{eq:qa-as-cubic-determinant}で
$G=FP$、$s=ab$、$P^2=\OO_D$を用いると
$\det(\rho Q)=F^{a+b}P^{1-a+ab}$となる。
数値差が$(1,0),(0,1),(1,0)$の場合、
連続する三つの商線束は
\[
(N_1N_0^{-1},N_2N_1^{-1},N_3N_2^{-1})
=(FP^{b-1},FP^{a+1},FP^b).
\]
逆順のrulingなら
$(FP^a,FP^b,FP^{a+1})$となる。
$P$が位数二なので商は$F,FP$の二種類だけで、
最初と最後が異なる。
従って三つのうち隣接する二つは同型である。
その共通次数二線束を$L$とすると、
合成写像
\[
 H^0(D,L)\otimes H^0(D,L)\longrightarrow H^0(D,L^2)
\]
のkernelは
$\bigwedge^2H^0(D,L)$という一次元空間である。
他方\eqref{eq:qa-as-cubic-spherical-restriction}によって
候補のforward Homと合成は
この線束上のHomと合成に一致する。
これは\eqref{eq:qa-as-cubic-two-compositions}の
同型性に矛盾する。
特定のmutation列に依存せず、
全てのcubic候補を排除した。
\end{proof}

% !TeX root = ../nc-quadric-surface-avatar.tex
\section{反標準 avatar の族と二つの十本矢 chart}
\label{sec:qa-avatar-charts}

本節では幾何的な族を直接定義し、その間の開埋め込みを証明する。
代数の同型類から幾何的 avatar への対応は
\cite[定義15.2・定理15.3]{Q12Avatar2026}に記載されているが、
その同型類分類を相対 AS 代数の stack 同値として使用しない。
族上の再埋め込みは\cite[第4--5節]{Q12Weyl2026}との比較も踏まえ、
本稿の補題~\ref{lem:qa-core-curve-vanishing}、
\ref{lem:qa-core-two-pencils}と定理~\ref{thm:qa-core-avatar-square-rotated}で
必要な幾何を自足的に証明し、
開部分の判定に必要なコホモロジー、非被約因子、垂直成分および降下を以下で確認する。
統合ノート\cite[第7.11--7.12項・第9節・第13.3項]{Synthesis2026}は
これらの出典と主張の水準を整理した資料であり、独立の検証の代用とはしない。

\paragraph{AS 側の辞書の対象.}
入力する体上分類の対象は
\cite[第15.1項・定理15.3]{Q12Avatar2026} の $A_Q^{\mathrm{pres}}$ である。
すなわち、局所有限な連結 positively directed $\mathbb Z$-代数で、
大域次元三、Gorenstein parameter 四の AS 正則性を満たし、
\eqref{eq:qa-universal-potential} の十二本矢 potential
（可縮な二次 pair を許す）の展開として表示できるものを、
全整数頂点を固定する代数同型で分類する。
最小化した四型では、右射影 $\mathsf P_i=e_iA$ と右単純 $\mathsf S_i$ に対して
\[
\begin{aligned}
0\longrightarrow\mathsf P_{i-4}
 &\longrightarrow\mathsf P_{i-3}^{\oplus2}\oplus\mathsf P_{i-2}^{\oplus\epsilon_i}\\
 &\longrightarrow\mathsf P_{i-1}^{\oplus2}\oplus\mathsf P_{i-2}^{\oplus\epsilon_i}
 \longrightarrow\mathsf P_i\longrightarrow\mathsf S_i\longrightarrow0
\end{aligned}
\]
という実際に完全な最小分解がある。
$\epsilon_i=0$、奇数 $i$ のみで $\epsilon_i=1$、
偶数 $i$ のみで $\epsilon_i=1$、全 $i$ で $\epsilon_i=1$ が、
順に八本、奇数対角十本、偶数対角十本、十二本に対応する。
右側の AS--Gorenstein 条件は
$\Ext_A^q(\mathsf S_i,\mathsf P_j)=0$ for $(q,j)\ne(3,i-4)$、
$\Ext_A^3(\mathsf S_i,\mathsf P_{i-4})\simeq\kk$ である。
分解の Euler 特性だけを課す条件ではない。
Noether 性はこの定義の追加仮定ではなく、引用する分類の結論である。
potential 表示、trivial pair の係数、周期同型、分解の基底は選択データにしない
\cite[定義2.1・第15.1項]{Q12Avatar2026}\cite[定義48.1]{AvatarVSeven2026}。
この整数頂点固定の同型は、四頂点 graded 代数の同型を射とした
$\mathfrak M_r^{\mathrm{alg}}$ と同じ相対同型層であるとは仮定しない。
本節の幾何的 stack の同型との比較も、体上の同型類辞書と区別する。

\subsection{分類する幾何データ}

基礎体 $\kk$ は標数零の代数閉体とする。
射影束は一次元商を分類する規約を用いる。
次の二つの標識付き曲面を考える。
\[
\begin{array}{c|c|c}
Y&h&f\\\hline
\Sigma _0=\PP^1\times\PP^1&\OO(1,1)&\OO(1,0)\\
\Sigma _2=\PP_{\PP^1}(\OO\oplus\OO(2))&\OO_{\PP(\OO\oplus\OO(2))}(1)&\pi^*\OO_{\PP^1}(1)
\end{array}
\]
いずれも $h^2=2$, $hf=1$, $f^2=0$, $-K_Y=2h$ である。
$\Sigma _2$ の負断面は $S=h-2f$ である。
本節の $\Pic^0$ は Picard 群の単位成分を表す。

\begin{definition}[曲線型の相対 avatar]
\label{def:qa-avatar}
局所 Noetherian な $\kk$-scheme $T$ 上の
$\mathcal A_{12}^{\times}$ の対象を
\[
 (Y\xrightarrow{\pi}B\longrightarrow T;\ h,f;\ D\hookrightarrow Y;\ P)
\]
とする。ここで $B/T$ は $\PP^1$-fibration、$Y/B$ は
$T$ 上 fppf 局所で $\PP_B(\mathcal E)$ と表示される射影直線族であり、
各幾何 fiber の $(Y,h,f)$ は上表のいずれかである。
$h,f$ は相対 Picard の類で、$f$ は指定した ruling に対応する。
$D$ は $T$ 上平坦な相対有効 Cartier 因子で、相対因子類は $2h$ である。
$P$ は $\Pic^0_{D/T}$ の section で、各幾何 fiber 上で非自明とする。
射は ruling、標識、$D$ および $P$ の類を保つ $T$ 上の同型である。
線束の scalar trivialization はデータに加えない。
\end{definition}

この定義は $\Sigma _0$ と $\Sigma _2$ の間で型が跳ぶ族を許す。
実際の線束を用いる式は fppf 局所で解釈し、基底からの線束の tensor は区別しない。
以下で
\[
 p=\pi|_D:D\longrightarrow B,\qquad F=\OO_Y(f)|_D,\qquad H=\OO_Y(h)|_D
\]
と書く。特に $H$ は線束であり、内在的モジュライの記号とは無関係である。
曲面、因子および相対 Picard 類の降下により、定義\ref{def:qa-avatar}は
fppf stack を定める。本節の「開部分 stack」は任意の基底への引戻しが
開部分 scheme となる意味で用いる。

\begin{remark}[有限表示と scalar の扱い]
fppf 局所で $B=\PP^1_T$ とし、$\mathcal E=\pi_*\OO_Y(h)$ を取ると、
各幾何 fiber の分裂型は $\OO(1)^{\oplus2}$ または $\OO\oplus\OO(2)$ である。
評価列は
\[
 0\longrightarrow \mathcal R\otimes\OO_B(-1)
 \longrightarrow \mathcal V\otimes\OO_B
 \longrightarrow\mathcal E\longrightarrow0,
 \qquad\rank\mathcal R=2,\quad\rank\mathcal V=4
\]
となる。したがって線形行列、その余核の射影束上の相対反標準因子、
および相対 Picard 類によって局所的な parameter を記述できる。
この表示の基底変更をすべて新たな自己同型として保持すると、
定義\ref{def:qa-avatar}と異なる stack になる。
本節の射はあくまで標識付き曲面と因子の同型である。
\end{remark}

\subsection{非被約因子を含むコホモロジー}

\begin{lemma}[支持に対する判定]
\label{lem:qa-support}
体上で $D\in|2h|$ とする。このとき
$\omega_D\simeq\OO_D$、$h^0(\OO_D)=h^1(\OO_D)=1$ である。
真の非零部分因子 $0<Z<D$ について $Z(D-Z)\geq2$ が成り立つ。
線束 $N$ の非零切断が正確に $Z\leq D$ に支持されるなら
\[
 \deg_ZN\ \geq\ Z(D-Z).
\]
従って、線束 $L$ がすべての $0<Z\leq D$ について
\begin{equation}
 \deg_ZL+Z(D-Z)>0
 \label{eq:qa-support-vanishing}
\end{equation}
を満たせば $H^1(D,L)=0$ である。
\end{lemma}

\begin{proof}
随伴と $0\to\omega_Y\to\OO_Y\to\OO_D\to0$ により
双対化層と構造層のコホモロジーを得る。
$\Sigma _0$ では $D=2f+2g$、$Z=af+bg$ と書けて
\[
 Z(D-Z)=2(a+b-ab),\qquad 0\leq a,b\leq2.
\]
$Z=0,D$ を除いたすべての場合で右辺は二以上である。
$\Sigma _2$ では $D=2S+4f$、$Z=aS+bf$ と書くと
\[
 Z(D-Z)=2a^2+2b(1-a),\qquad0\leq a\leq2,\quad0\leq b\leq4.
\]
$a=0$ では $b\geq1$、$a=1$ では右辺が二、
$a=2$ では真部分因子なら $b\leq3$ なので、同じ結論を得る。

切断の generic annihilator によって、その正確な支持である部分因子 $Z$ を取る。
各既約成分の generic な Artin 局所環で切断の消滅次数を取り除いて $Z$ の重複度を定めると、
切断は $N|_Z(-(D-Z))$ の、各 generic point で可逆な切断になる。
$Z$ は Cohen--Macaulay なので、この切断が定める
$\OO_Z\to N|_Z(-(D-Z))$ は単射で、余核は有限長である。
Euler 特性を比較すれば不等式を得る。
これは部分因子を重複度込みで取る議論であり、$D$ の被約性を使わない。
最後は $\omega_D\simeq\OO_D$ による Serre 双対を用い、
$L^{-1}$ の切断にこの不等式を適用すればよい。
\end{proof}

\begin{lemma}[再埋め込みに必要な消滅]
\label{lem:qa-vanishing}
$R\in\Pic^0(D)$ とする。このとき
\[
 H^1(D,HR)=H^1(D,HF^{-1}R)=0,\qquad
 \chi(HR)=4,\quad\chi(HF^{-1}R)=2.
\]
また $R\not\simeq\OO_D$ なら $H^*(D,R)=0$ である。
\end{lemma}

\begin{proof}
$R$ は各部分因子上で次数零である。
$h$ は nef で全次数が四なので、補題\ref{lem:qa-support}を適用できる。
$\Sigma _0$ では $h-f$ も nef で全次数が二である。
$\Sigma _2$ では $Z=aS+bf$ 上で
\[
 \deg_Z(h-f)+Z(D-Z)=2a^2-a+3b-2ab.
\]
$a=0,1,2$ に対する右辺はそれぞれ $3b$, $1+b$, $6-b$ であり、
$Z\neq0$ の許される範囲で正である。Riemann--Roch が Euler 特性を与える。
$R$ の非零切断は真部分因子には支持されない。
全体に支持される切断の余核は次数零なので零であり、$R$ は自明になる。
$R^{-1}$ にも同じ議論を適用し、Serre 双対を使う。
\end{proof}

\subsection{Picard twist による相対的な再埋め込み}

\begin{proposition}[垂直成分がある場合も含む直像]
\label{prop:qa-pushforward}
定義\ref{def:qa-avatar}の族で、$R\in\Pic^0_{D/T}(T)$ を取る。
fppf 局所で線束によって表示すると
\[
 Rp_*(HR)\simeq\mathcal E_R:=p_*(HR)
\]
は $B$ 上の階数二の局所自由層であり、任意の基底変換と両立する。
各幾何 fiber $B_t\simeq\PP^1$ 上で次数二かつ大域生成である。
評価写像は閉埋め込み
\[
 j_R:D\hookrightarrow Y_R:=\PP_B(\mathcal E_R)
\]
を定め、像は相対反標準 Cartier 因子である。
新しい標識の制限は $F_R=F$、$H_R=HR$ である。
\end{proposition}

\begin{proof}
まず体上で $b\in B$、$G_b=\pi^{-1}(b)$ とする。
$G_b$ が $D$ の成分でなければ $D\cap G_b$ は長さ二の scheme であり、
導来 fiber の切断は二次元で次数零に集中する。
$G_b\leq D$ なら、局所 parameter $t$ と $D$ の方程式 $t^m a$、$t\nmid a$ を用いると
\[
 \operatorname{ann}_{\OO_Y/(t^ma)}(t)=(t^{m-1}a)/(t^ma).
\]
したがって $L=HR$ の $B$ 方向の導来制限は
\[
 \mathcal H^{-1}=L|_{G_b}\otimes\OO_{G_b}(-D+G_b)\simeq\OO_{\PP^1}(-1),
 \qquad
 \mathcal H^0=L|_{G_b}\simeq\OO_{\PP^1}(1)
\]
を持つ。ここで $\deg_{G_b}L=1$、$D G_b=2$、$G_b^2=0$ を使った。
次数 $-1$ の層は acyclic であるから、この場合も直像の導来 fiber は $\kk(b)^2[0]$ となる。
垂直成分を持つ $p$ は有限でも平坦でもないが、この計算にはその仮定を要しない。

$i_*L$ は Cartier immersion $i:D\hookrightarrow Y$ の線束の直像なので perfect である。
滑らかな射影的射 $\pi$ による perfect pushforward と導来基底変換を用いると、
$R\pi_*i_*L$ は perfect である。
すべての導来 fiber が同じ階数二の次数零複体なので、局所的に可逆 minor で
可縮項を消去することにより、階数二の局所自由層となる。
族の場合、$i_*L$ は $T$ 上平坦であり、まず $T$ の幾何 fiber、次に $B_t$ の一点へ
制限すれば同じ計算になる。このため非被約な基底も含めて結論が成り立つ。

補題\ref{lem:qa-vanishing}と Riemann--Roch により
\[
 \deg\mathcal E_R=\chi(HR)-2=2,\qquad
 H^1(B,\mathcal E_R(-1))=H^1(D,HF^{-1}R)=0.
\]
$\PP^1$ 上の分裂定理より、分裂型は $\OO(1)^{\oplus2}$ または $\OO\oplus\OO(2)$ である。

各 $B$-fiber での評価は、長さ二の scheme 上の完全切断の評価、
または $\OO_{\PP^1}(1)$ の評価である。従って全射で、誘導される fiber の射は閉埋め込みである。
Nakayama の補題から全評価も全射となる。$j_R$ は proper かつ quasi-finite なので有限であり、
有限射の基底変換と再度の Nakayama により
$\OO_{Y_R}\to(j_R)_*\OO_D$ は全射である。
体上では像は滑らかな曲面内の Cohen--Macaulay な余次元一部分 scheme なので Cartier である。
像の類を $a h_R+b f_R$ とすると、ruling との交点から $a=2$、
$h_R$ との交点 $\deg_D(HR)=4$ から $b=0$ を得る。
族では $D/T$ の平坦性と各 fiber の Cartier 性によって相対 Cartier 性が従う。
各 fiber で自明な差の線束は、$H^1(\OO_{Y_t})=0$ と基底変換により基底から来るので、
相対 Picard 類も $2h_R$ である。
\end{proof}

\begin{theorem}[族上の幾何的並進]
\label{thm:qa-translation}
前命題において $R=P^n$ とした構成は $\mathcal A_{12}^{\times}$ の自己同値
$\mathsf T^n$ を定める。自然な同型
\[
 \mathsf T^m\mathsf T^n\simeq\mathsf T^{m+n},\qquad
 \mathsf T^0\simeq\operatorname{id}
\]
があり、逆操作は $P^{-n}$ を用いる同じ構成である。
\end{theorem}

\begin{proof}
抽象的な $D$、ruling $p:D\to B$ および $P$ は変更せず、$H$ を $HP^n$ に変える。
二回適用した直像は $p_*(HP^nP^m)$ なので、tensor の結合則が合成則を与える。
$R=\OO_D$ の場合、制限列
\[
 0\longrightarrow\OO_Y(-h)\longrightarrow\OO_Y(h)\longrightarrow i_*H\longrightarrow0
\]
と $R\pi_*\OO_Y(-h)=0$ が元の $\mathcal E=\pi_*\OO_Y(h)$ を回収する。
相対的には左項に基底線束の tensor が入ることがあるが、同じ消滅が成り立つ。
評価の自然性により射影束だけでなく $D$ の埋め込みも回収する。
線束の代表を基底線束で tensor すると直像も同じ線束で tensor されるため、
射影束および相対標識は変わらない。
$q:D\to T$ について、各幾何 fiber で $h^0(\OO_D)=1$ であるから、
コホモロジーと基底変換により $q_*\OO_D=\OO_T$ が任意の基底変換後も成立する。
従って線束代表の自己同型は基底からの scalar であり、射影化で消える。
従って構成と合成同型は fppf 降下し、射と基底変換に整合する。
\end{proof}

この定理は\cite[定理5.4]{Q12Weyl2026}の操作である。
ここで証明したのは幾何的 avatar の stack 上の並進であり、
無限二面体群のすべての関係を stack 上へ持ち上げる主張ではない。

\subsection{二つの開 chart と四つの層}

\begin{definition}[幾何的な八本矢・十本矢 chart]
\label{def:qa-open-charts}
$\mathcal U\subset\mathcal A_{12}^{\times}$ を、全幾何 fiber が $\Sigma _0$ である部分とする。
\[
 \mathcal A_{10,\mathrm e}^{\times}:=\mathcal U,
 \qquad
 \mathcal A_{10,\mathrm o}^{\times}:=\mathsf T^{-1}(\mathcal U),
 \qquad
 \mathcal A_8^{\times}:=\mathcal U\cap\mathsf T^{-1}(\mathcal U)
\]
と置く。添字 $\mathrm e,\mathrm o$ は、後述する標準表示で残る偶数対角・奇数対角に対応する。
ここでは定義自体を幾何データだけで与えている。
\end{definition}

\begin{theorem}[avatar の開部分 stack]
\label{thm:qa-open-charts}
次の射は開埋め込みである。
\[
 \mathcal A_8^{\times}\hookrightarrow
 \mathcal A_{10,\mathrm o}^{\times}\hookrightarrow\mathcal A_{12}^{\times},
 \qquad
 \mathcal A_8^{\times}\hookrightarrow
 \mathcal A_{10,\mathrm e}^{\times}\hookrightarrow\mathcal A_{12}^{\times}.
\]
幾何点での層は次の表で与えられる。
\begin{center}
\begin{tabular}{c|c|c}
曲面&曲線型の条件&属する層\\\hline
$\Sigma _0$&$P\not\simeq F^2H^{-1}$&$\mathcal A_8^{\times}$\\
$\Sigma _0$&$P\simeq F^2H^{-1}$&$\mathcal A_{10,\mathrm e}^{\times}\setminus\mathcal A_8^{\times}$\\
$\Sigma _2$&$D\cap S=\varnothing$&$\mathcal A_{10,\mathrm o}^{\times}\setminus\mathcal A_8^{\times}$\\
$\Sigma _2$&$S\leq D$&$\mathcal A_{12}^{\times}\setminus
(\mathcal A_{10,\mathrm e}^{\times}\cup\mathcal A_{10,\mathrm o}^{\times})$
\end{tabular}
\end{center}
最後の行は幾何点集合として閉 locus である。
全 fiber がこの型となる族の充満部分関手に、閉埋め込みを自動的に与える主張ではない。
\end{theorem}

\begin{proof}
$\mathcal E=\pi_*\OO_Y(h)$ の二つの分裂型は
\[
 h^0(B_t,\mathcal E_t(-2))=
 \begin{cases}0&Y_t=\Sigma _0,\\1&Y_t=\Sigma _2\end{cases}
\]
によって区別される。コホモロジーの上半連続性から $\mathcal U$ は開である。
自己同値 $\mathsf T$ による逆像と交わりも開なので、すべての開埋め込みを得る。
残るのは幾何点での条件である。

$\Sigma _0$ で $H=FG$ と書くと
\[
 h^0(B,\mathcal E_P(-2))=h^0(D,GF^{-1}P).
\]
$Z=af+bg$ に対し、$GF^{-1}P$ の次数は $a-b$ である。
真部分因子では
\[
 Z(D-Z)-\deg_Z(GF^{-1}P)=a+3b-2ab>0.
\]
よって補題\ref{lem:qa-support}により、この線束の非零切断は全体に支持される。
全次数は零なので、非零切断があることは線束が自明であることと同値である。
従って $\mathsf T(Y)=\Sigma _2$ となる条件はちょうど $P\simeq F^2H^{-1}$ である。

$\Sigma _2$ で $D\cap S=\varnothing$ なら $H\simeq F^2$ だから
$h^0(D,HF^{-2}P)=h^0(D,P)=0$ であり、$\mathsf T(Y)=\Sigma _0$ となる。
$S\leq D$ なら $C=D-S$ と置く。$C\sim h+2f$ であり、
$\OO_Y(-C)$ は $\pi$ に沿って次数 $-1$ なので acyclic である。
従って $h^0(\OO_C)=1$, $h^1(\OO_C)=0$、標数零より $\Pic^0(C)=0$ を得る。
残余因子の完全列は
\[
 0\longrightarrow P|_C\longrightarrow\OO_D(S)P
 \longrightarrow\OO_S(-2)\longrightarrow0.
\]
$P|_C\simeq\OO_C$ なので $h^0(D,HF^{-2}P)=1$ である。
従って $\mathsf T(Y)=\Sigma _2$ のままである。
最後に $D S=0$ であるから、有効因子 $D$ は $S$ を含むか $S$ を避けるかのいずれかである。
これですべての場合を得る。閉性は二つの開部分の和の補集合であることによる。
\end{proof}

\begin{remark}[矢数の名称と体上の辞書]
幾何的な開埋め込みと表の分解は、上の証明の通り
射影束の分裂型、Picard 類および垂直成分条件から従う。
各層を $8,10,10,12$ 本の最小矢数と対応させる体上の解釈は、
$A_Q^{\mathrm{pres}}$ を対象とする
\cite[定理15.3]{Q12Avatar2026}の同型類分類を用いる。
同じ矢数の相対 QP 表示を比較する際の階数条件は、
上の幾何学的な層の名称からは
推論せず、各代数表示の積と関係について検証する。
\end{remark}


\begin{remark}[「十本矢全体」の二通りの意味]
最小十本矢の二つの層の和は $\mathcal A_8^{\times}$ を含まない。
これに対し上の $\mathcal A_{10,\mathrm o}^{\times}$ と
$\mathcal A_{10,\mathrm e}^{\times}$ は、可縮な pair を許す十本矢表示の
幾何的 chart に対応し、八本矢部分を開部分として含む。
二つの十本矢 chart の和は、最深の十二本矢層の補集合である。
この区別なしに「矢の少ないモジュライが矢の多いモジュライの開部分である」
と述べることはできない。
\end{remark}

\subsection{直積型 collection と全 square avatar}

$\mathcal A_{\mathrm{sq}}^{\times}$ を、順序付きの二つの ruling を持つ
$\Sigma _0$-族、その相対 $(2,2)$ Cartier 因子 $D$、および fiberwise 非自明な
$P\in\Pic^0_{D/T}$ の stack とする。
曲線上の ruling の制限を $F_s,G_s$ と書く。
これは $h=f+g$ と置くことによって $\mathcal U$ と同じ幾何的 stack である。

\begin{corollary}[全 square chart の相対的な対応]
\label{cor:qa-square-chart}
幾何的 stack の同値
\[
 \mathcal A_{\mathrm{sq}}^{\times}
 \xrightarrow[\ \mathsf T^{-1}\ ]{\ \sim\ }
 \mathcal A_{10,\mathrm o}^{\times}
\]
があり、制限線束は
\begin{equation}
 F=F_s,\qquad H=F_sG_sP^{-1},\qquad P=P
 \label{eq:qa-square-dictionary}
\end{equation}
となる。逆には $\mathsf T$ によって得られる $\Sigma _0$ 上で
\[
 F_s=F,\qquad G_s=HPF^{-1}
\]
を取る。square 座標で八本矢部分となる条件は
$P\not\simeq G_sF_s^{-1}$ であり、壁 $P\simeq G_sF_s^{-1}$ は
負断面を避ける $\Sigma _2$ 上の奇数対角十本矢層に移る。
\end{corollary}

\begin{proof}
定義から $\mathsf T^{-1}:\mathcal U\to\mathsf T^{-1}(\mathcal U)$ は同値であり、
その線束の式が\eqref{eq:qa-square-dictionary}である。
逆の式は $H_s=HP=F_sG_s$ による。
元の $\Sigma _0$ で $\mathsf T^{-1}$ が $\Sigma _2$ を与える条件は、
定理\ref{thm:qa-open-charts}の計算を $P^{-1}$ に対して行って
$P^{-1}=F_sG_s^{-1}$、すなわち $P=G_sF_s^{-1}$ である。
その像は $\mathcal A_{10,\mathrm o}^{\times}$ にあるので負断面を避ける。
\end{proof}

この系では反標準因子を滑らか、被約または既約と仮定していない。
従って「十本でちょうど全ての quadric の divisor と $\Pic^0$ が出る」
という記述は、\emph{奇数対角の非最小十本矢 chart と全 square avatar の対応}
として正確に定式化できる。最小十本矢だけの主張ではなく、$P=\OO_D$ も除いている。

代数側の square foundation は
\[
 B_{\mathrm{sq}}=\End_{\to}(P,F_s,G_s,F_sG_s)
\]
であり、可換な基準 collection は
$(\OO,\OO(1,0),\OO(0,1),\OO(1,1))$ である
\cite[定義1.2・定理1.3]{AvatarVSeven2026}。
最初の二項の左 mutation は曲線上で
\[
 (P^2F_s^{-1},P,G_s,F_sG_s)
\]
を与える。共通に $F_sP^{-1}$ を tensor すると
\[
 (P,F_s,F_sG_sP^{-1},F_s^2G_sP^{-1})=(P,F,H,HF)
\]
となり、上の幾何的辞書と一致する。
mutationによる指定型AS代数の対応は定理~
\ref{thm:qa-as10-global-bijection}と
\ref{thm:qa-as10-relative-families}で証明する。
元の体上および仮定付き相対対応は
\cite[定理51.4・系51.5、定理53.2]{AvatarVSeven2026}にある。
幾何的な開埋め込みは独立に証明した。

\subsection{滑らかな楕円型における三者の明示的辞書}
\label{subsec:qa-square-elliptic-pointed-dictionary}

\begin{theorem}[VdB楕円四つ組とsquare avatarの同型類の全単射]
\label{thm:qa-elliptic-vdb-square-avatar-bijection}
$k$を標数零の代数閉体とする。
滑らかな楕円曲線$E$上のVan den Bergh許容四つ組
$(E,L_0,L_1,L_2)$で、各$L_i$は次数二、
\[
 L_1L_0^{-1}\ne\OO_E,\qquad
 L_2L_0^{-1}\ne\OO_E,\qquad
 L_1L_2^{-1}\ne\OO_E
 \tag{VdB-adm}\label{eq:qa-vdb-square-admissibility}
\]
を満たすものの、曲線・三線束を保つ同型類を考える。
これは、\emph{順序付き}二rulingを持つ
滑らかなsquare avatar
$(Y=\PP^1\times\PP^1,D,P)$で
\[
 D\in|\OO_Y(2,2)|,\qquad D\text{ は滑らか},\qquad
 P\in\Pic^0(D)\setminus\{\OO_D\},\qquad
 P\not\simeq G_sF_s^{-1}
 \tag{Sq-cub}\label{eq:qa-square-elliptic-cubic-open}
\]
を満たすものの同型類と自然に全単射する。
ここで$F_s=\OO_D(1,0)$、
$G_s=\OO_D(0,1)$とし、
二rulingの順序は同型の中で保持する。

全単射の\emph{両方向の式}は
\[
 \boxed{
 \begin{aligned}
 (E;L_0,L_1,L_2)&\longmapsto
  (D\hookrightarrow\PP H^0(F_s)^\vee
           \times\PP H^0(G_s)^\vee;\ P),\\
 F_s&=L_2,\qquad
 G_s=L_1L_2L_0^{-1},\qquad
 P=L_2L_0^{-1},\\
 (D;F_s,G_s,P)&\longmapsto
  (E=D;\ L_0=F_sP^{-1},\
   L_1=G_sP^{-1},\ L_2=F_s).
 \end{aligned}}
 \tag{Sq-VdB}\label{eq:qa-square-vdb-explicit}
\]
で与えられる。
この同型類の比較は、周期同型や根を
追加標識にしないVdBのcubic$\mathbb Z$-代数
の幾何的分類を通じて、
\emph{選んだcubic normal foundation}の
相対的な$k$点の辞書になる。

差パラメータ
$\eta=L_1L_0^{-1}$、
$\delta=L_2L_0^{-1}$を用いれば、
\[
 \eta=G_sF_s^{-1},\qquad \delta=P.
 \tag{Sq-diff}\label{eq:qa-square-vdb-differences}
\]
\end{theorem}
\begin{proof}
まず$L_0,L_1,L_2$から
$P=L_2L_0^{-1}$、
$F_s=L_2$、
$G_s=L_1L_2L_0^{-1}$を作ると、
$F_s,G_s$はいずれも次数二で、
$G_sF_s^{-1}=L_1L_0^{-1}$である。
式~\eqref{eq:qa-vdb-square-admissibility}より
$F_s\not\simeq G_s$である。
滑らかな楕円曲線上の相異なる次数二線束の
完全線形系は、それぞれ二重被覆
$E\to\PP^1$を与え、その積は
$\PP^1\times\PP^1$への閉埋込みである。
像は$(2,2)$因子で、反標準因子となる。
また$P\ne\OO_E$、
$P\not\simeq G_sF_s^{-1}$である。
これは曲線同型を保ち、線形系の基底の選択を
$\PGL_2^2$で割ればavatarの同型類を定める。

逆にsquare avatarから
$L_0=F_sP^{-1}$、
$L_1=G_sP^{-1}$、
$L_2=F_s$を作ると、
式~\eqref{eq:qa-square-elliptic-cubic-open}により
$\eta=G_sF_s^{-1}\ne\OO_E$、
$\delta=P\ne\OO_E$、
$\eta\ne\delta$となる。
従ってVdBの許容条件が成立する。
二つの式を順次適用すると三線束と
二つのrulingの完全線形系をそのまま回収するので、
互いに逆である。
VdBの幾何的quintuple復元
\cite[Theorem 5.5.10]{VdB2011}を適用すれば
最後の解釈も従う。
\end{proof}

\begin{theorem}[O-markingを保つruling反射と楕円型の商全単射]
\label{thm:qa-elliptic-avatar-pointed-quotient}
前定理のsquare avatarについてさらに
\[
 G_sF_s^{-1}\not\simeq P^2
 \tag{two-cub}\label{eq:qa-square-two-cubic-open}
\]
を課した開を$\mathcal A_{\mathrm{sq,ell}}^{2c}$とする。
この開上では
\[
 \boxed{\quad
 r_1(F_s,G_s,P)
  =(G_sP^{-1},F_sP,P)
 \quad}
 \tag{Sq-R}\label{eq:qa-square-pointed-reflection}
\]
は、右辺の二つの次数二線束が定める
新しい$(2,2)$埋込みを用いて、
幾何的avatarの同型類の対合を定める。
前定理の式~\eqref{eq:qa-square-vdb-explicit}の下で、
これは論文Bの
\cite[\texttt{lem:elliptic-two-reflections}]{QuadricModuli2026}
の\emph{構造対象$\OO$を保つ}ruling交換
\[
 (L_0,L_1,L_2)\longmapsto
 (L_0L_1L_2^{-1},L_2,L_1)
\]
に一致する。
従って、$y=(\mathcal C,\OO)$が滑らかな
楕円型かつ二つのnormal foundationがともに
cubic型となる\emph{幾何的なO-marked同型類集合}は
\[
 \boxed{\quad
 \bigl|\mathcal A_{\mathrm{sq,ell}}^{2c}(k)\bigr|
    /\langle r_1\rangle
 \ \xrightarrow{\ \sim\ }\
 \bigl|\mathfrak H_{2,\mathrm{ell}}^{2c}(k)\bigr|
 \quad}
 \tag{Sq-H2}\label{eq:qa-square-pointed-bijection}
\]
と全単射する。
左辺はavatar同型類の集合をとってから
$r_1$の軌道で割ることを意味する。
ここでは$S_2$の商の全相対stack同値を
この幾何的全単射だけから推論しない。

一方、helixの一段の移動に対応する
\[
 r_0(F_s,G_s,P)=(G_s,F_sP,P)
 \tag{Sq-T}\label{eq:qa-square-unpointed-shift}
\]
はVdB四つ組の
$(L_0,L_1,L_2)\mapsto(L_1,L_2,L_3)$に
対応するが、一般に$\OO$を変更する。
従って\emph{O-marked}同型類を得るために
$r_0$まで同一視してはならない。

この商の記述は、一般の二つのavatarの
$r_1$-軌道が二元であることを要求しない。
特殊な二捩れ慣性点では二つが同型になる。
\end{theorem}
\begin{proof}
式~\eqref{eq:qa-square-vdb-explicit}を
$r_1$の右辺へ適用すると
\[
 \begin{aligned}
 L'_0&=G_sP^{-2}
       =L_0L_1L_2^{-1},\\
 L'_1&=F_s=L_2,\qquad
 L'_2=G_sP^{-1}=L_1
 \end{aligned}
\]
を得る。新しい二ruling線束が相異なる条件は
$G_sP^{-1}\not\simeq F_sP$、
すなわち式~\eqref{eq:qa-square-two-cubic-open}
に他ならない。
もう一度適用すると$(F_s,G_s,P)$を回収するため
対合である。

論文Bの二つの内在的foundationの系
\cite[\texttt{cor:two-intrinsic-foundations}]{QuadricModuli2026}と、
pointed groupoid比較の系
\cite[\texttt{cor:h2-point-groupoid-foundation}]{QuadricModuli2026}は、指定$\OO$から始まるnormal foundationは
exact同値の下でちょうど二択であり、
二つともcubic型なら他方は上記$R$で得られること、
またfoundation同型はpointed exact同型を
完全に検出することを示す。
Van den Berghの四つ組復元を合わせれば、
二cubic楕円型のpointed同型類に関する
核と像はちょうど$r_1$の同型類軌道となる。
これが式~\eqref{eq:qa-square-pointed-bijection}である。

同様に$r_0$の右辺から得る
$(L'_0,L'_1,L'_2)$は
\[
 (G_sP^{-1},F_s,G_s)=(L_1,L_2,L_3),
 \quad L_3=L_0^{-1}L_1L_2
\]
である。論文Bの同じ補題では、この添字移動は
underlying Abel圏の同値類は保つが、
一般に構造対象$\OO$を保たない。
従って$r_0$を割るのは別の
unpointed moduli問題である。
\end{proof}

\begin{corollary}[OU potentialからsquare avatar、O-marked対象への一般次数]
\label{cor:qa-ou-potential-square-avatar-pointed-degrees}
論文C
~\ref{thm:qa-ou-rigid-potential-quintuple-generic-etale}
の正則半単純な滑らかな楕円開を、
前定理の二cubic開~\eqref{eq:qa-square-two-cubic-open}
および一般慣性同型の稠密開へ制限する。
この開ではpotentialの幾何的四重表示が
VdB四つ組$(E,L_0,L_1,L_2)$を与え、
式~\eqref{eq:qa-square-vdb-explicit}により
\emph{square avatarへの相対的な対応}がある。
全ての線束・曲線が相対的に構成される
étale chart上では、これは基底変更と可換な
代数スタックの射となり、
相対的な曲線同型を通じて降下する。
一方、対応の幾何的fiberの一般次数は
\[
 \boxed{
 \begin{array}{ccc}
 \text{有効potential表示}
   &\xrightarrow{\;4\;}&
  \text{順序付きsquare avatar}\\
 &&\big\downarrow\vcenter{\rlap{$\scriptstyle 2$}}\\
 &&\text{O-marked }\mathfrak H_2
 \end{array}}
 \tag{Pot-Sq-H2}\label{eq:qa-potential-square-pointed-degrees}
\]
であり、potentialからO-marked対象への
一般次数は$4\cdot2=8$である。
ここで上の$4$はpotential/quintupleの
\emph{表示間}の一般四重射、
下の$2$は二つのnormal foundationの忘却であり、
$\Pic^0$の四つの平方根の選択を
追加標識として数えた次数ではない。
\end{corollary}
\begin{proof}
Cの
~\ref{thm:qa-ou-rigid-potential-quintuple-generic-etale}
による一般有限étale四重射は、
指定された幾何quintupleの相対四重テンソルを与える。
滑らかな正則半単純開では、C
~\ref{cor:qa-ou-relative-regular-semisimple-potential}
の点曲線と二rulingの相対構成が、
基底変更と可換な三つの次数二線束を与える。
それらのテンソル積と二つの完全線形系で
式~\eqref{eq:qa-square-vdb-explicit}の
square avatarを相対的に作る。
曲線・線束の同型を保つので基底変更と降下に
関して自然である。
VdBの四つ組の幾何的同型類とsquare avatarの
全単射は定理
~\ref{thm:qa-elliptic-vdb-square-avatar-bijection}で証明した。
これとCの
~\ref{cor:qa-quintuple-h2-generic-degree-two}の
二つのfoundationの一般次数を合成すれば
表示された幾何一般次数が従う。
\end{proof}

\begin{remark}[一cubic境界と古典点での辞書の射程]
\label{rem:qa-square-avatar-boundary-cautions}
$G_sF_s^{-1}=P^2$では、$r_1$後の
$F'_s=G_sP^{-1}$と$G'_s=F_sP$が一致し、
square$(2,2)$埋込みを構成できない。
これは論文C
~\ref{cor:qa-ou-one-cubic-fiber-and-nonfiniteness}
の一cubic境界であり、
二つ目のnormal foundationはnoncubicとして
内在的な選択被覆
~\ref{thm:qa-global-two-ruling-etale-cover}
に残る。

また$P=\OO_D$を任意の$(2,2)$因子に
そのまま許すと、これらの曲線から得る
可換な四項foundation代数は$D$に依存せず、
同じ可換quadricを重複して数える。
従って全同型類の全単射には古典点を
別途一点として追加し、元のavatarの
$P=0$座を単に追加するだけでは足りない。
$\Sigma_2$の古典型は別のweighted
stackを生み、無条件には$\HHstack$に入らない。
$\Sigma_2$曲線型でのO-marked適合性、
非被約反標準因子を含む全対応、
およびpotentialの安定化・自己同型群の比較は
式~\eqref{eq:qa-square-pointed-bijection}
の証明には含まれない。
\end{remark}

\begin{proposition}[二つのwallは異なり、noncubic avatar壁には追加認識が必要]
\label{prop:qa-avatar-walls-and-noncubic-dimension}
滑らかな楕円型square avatarについて、
$\eta=G_sF_s^{-1}$、$\delta=P$と書く。
次の二つのwallを区別する。
\begin{enumerate}
\item $\eta=2\delta$、$\delta\notin E[2]$：
式~\eqref{eq:qa-square-vdb-explicit}の四つ組は
geometric cubic型であるが、
$r_1$後の第二foundationはnoncubicになる。
したがって元のO-marked点は
\emph{cubic locusに属する}。
\item $\eta=\delta\ne0$：
式~\eqref{eq:qa-square-vdb-explicit}の
$L_1=L_2$なので元のfoundationが
geometric cubic型ではない。
このwallはsquareの八本矢開の補集合であり、
曲線型avatarの十本矢・$\Sigma_2$ chartへ
移行する。
\end{enumerate}
後者の全smooth楕円avatar同型類の
パラメータ空間は、楕円曲線の$j$と
非零点$\delta\in E$を局所座標として
\emph{二次元}である。
一方、論文Bの内在noncubic境界
$\mathfrak N\subset\HHstack$は
\emph{純一次元}のArtinスタックである。
従ってこのwall\emph{全体}を
内在noncubic境界と表現可能quasi-finiteな
モジュライ同型で結び付けることはできない。
O-marked適合性はこのwall上で
さらに制限されるか、あるいは
同じpointed対象へ正の相対次元を持って
潰れるかのいずれかである。

特に非零二捩れ$P$の条件$P^2=\OO_E$は
楕円曲線のmoduli上で一次元を定め、
次に検討すべき\emph{候補}ではあるが、
それが内在noncubic境界をちょうど尽くすことは
ここでは証明していない。
\end{proposition}
\begin{proof}
最初のwallでは$L_0=F_sP^{-1}$、
$L_1=G_sP^{-1}$、$L_2=F_s$が
互いに相異なり、VdB許容条件を満たす。
$r_1$による新しい二線束は
$G_sP^{-1}=F_sP$と$F_sP$で一致するので、
第二はnoncubicとなる。
これはC
~\ref{cor:qa-ou-one-cubic-fiber-and-nonfiniteness}
の一cubic境界に一致する。
次のwallでは$G_s=F_sP$だから
$L_1=G_sP^{-1}=F_s=L_2$となる。
$\mathcal A_8^\times$の条件
~\eqref{eq:qa-square-elliptic-cubic-open}
に反することも直接に分かる。
$\Sigma_2$との対応は
系~\ref{cor:qa-square-chart}に従う。

最後の次元を計算する。
楕円曲線$E$のmoduliは一次元で、
$F_s\in\Pic^2(E)$の選択は曲線の
translationによって動かせる。
標数零では$[2]:E\to E$が全射なので、
このtranslationの軌道は$\Pic^2(E)$全体である。
$\eta=\delta$の下では$G_s=F_sP$だから
残るパラメータは
$E$と$0\ne P\in\Pic^0(E)$で、
そのmoduliの一般次元は$1+1=2$。
一方、論文B
\cite[\texttt{thm:noncubic-boundary-smooth-codim-two}]{QuadricModuli2026}
は全$\HHstack$が純次元三で
$\mathfrak N$が滑らかな余次元二、
すなわち純次元一であることを証明している。
有限型quasi-finite射の像は源の各成分と同じ
次元を持つので、この二次元wall全体を
一次元の境界にquasi-finiteに写すことはできない。
最後の二捩れ条件の次元も
有限étaleな$E[2]\setminus0$から従う。
\end{proof}

\begin{theorem}[二捩れsquare壁とnoncubic二根対の幾何同型類の全単射]
\label{thm:qa-avatar-two-torsion-root-pair-bijection}
$k$を標数零の代数閉体とする。
次の二集合を考える。
\begin{enumerate}
\item 滑らかな楕円曲線$E$、次数二線束$F_s$、
非自明な$P\in\Pic^0(E)[2]$の三つ組
$(E,F_s,P)$の曲線・線束・Picard類を
保つ同型類集合。$G_s=F_sP$とすると
$F_s\not\simeq G_s$だから、二完全線形系は
$(2,2)$埋込みを定め、対応するsquare avatarは
壁$P=G_sF_s^{-1}$上にある。
\item 射影直線$B\simeq\PP^1$上の、
互いに交わらない\emph{被約}二点因子の
非順序対$\{B_1,B_2\}$で、
$\deg B_i=2$かつ$B_1+B_2$は
四つの相異なる点からなるものの、
$\PGL_2$同型類集合。
\end{enumerate}
この二集合の間に自然な全単射がある。
対応は、$F_s$の二重被覆
\[
 \pi_F:E\longrightarrow\PP H^0(E,F_s)^\vee
\]
の四つの分岐点に対し、
$P$に対応する非零二捩れ平行移動が
誘導する二組の分割を取ることで与えられる。
逆方向には、四分岐点$B_1+B_2$の
二重被覆$\pi_F:E\to B$を取り、
その分割に対応する唯一の
非零$E[2]$の元を$P$とする。
このとき$F_s=\pi_F^*\OO_B(1)$である。

論文B
\cite[\texttt{thm:noncubic-intrinsic-presentation-equivalence}]{QuadricModuli2026}
の非cubic根対表示によれば、
従ってこの二捩れsquare壁の幾何的avatar同型類は、
\emph{四根が相異なるnoncubic O-marked H2境界}
の同型類とも全単射する。

最後の全単射は\emph{論文Bの根対分類を介した}
同型類の対応である。
この幾何辞書だけから
avatar側のAS$\mathbb Z$-代数の
$\qgr$とそのnoncubic O-marked heartの
同値、または対応するモジュライstackの射は
まだ結論しない。
\end{theorem}
\begin{proof}
次数二の線束$F_s$は楕円曲線から$\PP^1$への
二重被覆を与える。
その被覆対合を$\iota_F$と書くと、
その固定点集合$R_F$は四点からなり、
像$B_F$が四つの相異なる分岐点となる。
$E$は標識なしでも$\Pic^1(E)$のtorsorで、
$P\in\Pic^0(E)[2]\setminus0$は
$E$上の固定点のない位数二平行移動$t_P$を
定める。
次数二線束に対する平行移動の作用は
$t_P^*F_s\simeq F_sP^{\pm2}\simeq F_s$なので、
$t_P$は$\iota_F$と可換であり、
$R_F$を保つ。
その四点への作用は自由であり、
ちょうど二つの長さ二の軌道に分ける。
二つの軌道の$\pi_F$による像が
$B_F=B_1+B_2$の二分割となる。
これは線束の基底・楕円原点の選択に依存せず、
曲線と線束の同型に自然である。

逆に四つの相異なる分岐点$B_1+B_2$から、
$B$上の二重被覆は、
$\mathcal O_B(2)$とその四点を零点に持つ
切断を用いて、被覆同型を除き一意に構成される。
$F_s=\pi_F^*\OO_B(1)$は次数二である。
四つの分岐点は$E$上の被覆対合の固定点と
一対一に対応する。任意の一つを楕円原点に
選べば、その四点が$E[2]$となる。
$E[2]\simeq(\mathbb Z/2)^2$の各非零元の
平行移動は四点を二つの組に分割し、
三つの非零元と三つの二分割との間に
自然な全単射を与える。
従って指定の$\{B_1,B_2\}$から
一意の非零二捩れ$P$を回収する。
原点を変えても$\Pic^0(E)$における$P$の
同型類は変わらない。
二方向の構成が互いに逆であることも
この四点の分割から直ちに従う。

具体的に、適当な射影座標で
$B_1=\{0,\infty\}$、$B_2=\{1,\lambda\}$、
$\lambda\notin\{0,1,\infty\}$とすれば、
$E$はLegendre曲線
\[
 y^2=x(x-1)(x-\lambda)
\]
で与えられる。$B_1,B_2$をそれぞれ
集合として保ち、その二点を交換する
$\PP^1$の対合は$x\mapsto\lambda/x$である。
その固定点を持たないliftは
\[
 (x,y)\longmapsto
 \left(\frac{\lambda}{x},
             -\frac{\lambda y}{x^2}\right)
 \tag{Root-2}\label{eq:qa-avatar-root-pair-translation}
\]
であり、$E$上では非零二捩れ点
$(0,0)$によるtranslationである。
この式からも二分割と二捩れ点の
相互復元を直接確認できる。

最後に論文Bの非cubic被約境界は、
二つの二根集合の非順序対を
$\PGL_2$で割った群oidとして記述されており、
四根相異なる開ではまさに第二集合である。
従ってBの分類との合成が最後の
同型類全単射を与える。
これはBのO-marked対象への
逆構成による対応であり、
元のavatar AS代数の直接認識とは区別する。
\end{proof}

\subsection{二根対のquartic resolventと退化avatarの非可逆Picardパラメータ}
\label{subsec:qa-rootpair-resolvent-singular-avatar}

\begin{theorem}[noncubic根対の三重quartic resolventと三つの退化型]
\label{thm:qa-noncubic-rootpair-quartic-resolvent}
$k$を標数零の代数閉体、$W$を二次元ベクトル空間とし、
$\mathcal D_2=\PP(\operatorname{Sym}^2W^\vee)$、
$\mathcal D_4=\PP(\operatorname{Sym}^4W^\vee)$と置く。
二つの二次形式をその順序を忘れて掛ける射
\[
 \mu:\operatorname{Sym}^2(\mathcal D_2)\longrightarrow
             \mathcal D_4,\qquad
 \{[f_1],[f_2]\}\longmapsto[f_1f_2]
 \tag{Res-4}\label{eq:qa-rootpair-resolvent-map}
\]
は全射有限射であり、一般次数は$3$である。
根が相異なる四次形式$q$の上の三点は
四根の非順序二根対への三分割を表す。
アフィンmonic chartで
\[
 q(x)=x^4+a x^3+b x^2+c x+d
\]
とすると、四根$\alpha_1,\ldots,\alpha_4$の
二組分割から得る
$\alpha_i\alpha_j+\alpha_k\alpha_\ell$の
三つの値は
\[
 R_q(z)=z^3-bz^2+(ac-4d)z
             -(a^2d+c^2-4bd)
 \tag{Res-cub}\label{eq:qa-rootpair-resolvent-cubic}
\]
の三根に一致する。
この三重射は古典的なquartic resolventの
幾何的な表示である
\cite[Section 4.3]{Cremona1999}。

$\mathcal D_2^{\mathrm{sf}}$を二つの相異なる
根を持つ二次形式の開部分とする。
論文Bのnoncubic根対表示の
\[
 \mathcal R_{\mathrm{sf}}=
 \operatorname{Sym}^2(\mathcal D_2^{\mathrm{sf}})
\]
上では$q=f_1f_2$の幾何的根配置は
次の三型に限られ、それぞれ
\[
 C_q:\quad y^2=f_1f_2
 \quad\text{in }\operatorname{Tot}\OO_{\PP^1}(2)
 \tag{Res-curve}\label{eq:qa-rootpair-spectral-curve}
\]
の射影的二重被覆が次の形となる。
\[
 \begin{array}{c|c|c}
 \text{根配置}&q\text{の重複度}&C_q\\\hline
 B_1\cap B_2=\varnothing&(1,1,1,1)&\text{smooth genus }1\\
 |B_1\cap B_2|=1&(2,1,1)&\text{既約nodal genus }1\\
 B_1=B_2&(2,2)&\text{二つの}\PP^1\text{が二節点で交わる}
 \end{array}
 \tag{Res-strata}\label{eq:qa-rootpair-three-strata}
\]
三型とも$C_q$は連結なGorenstein曲線で
算術種数$1$である。
この構成は$\operatorname{PGL}(W)$と
根対交換の作用に自然であり、論文Bの
根対スタック
$\mathfrak R=[((\mathcal D_2^{\mathrm{sf}})^2)/
(\operatorname{PGL}(W)\times S_2)]$
からquartic二重被覆のスタックへの射を与える。
ただし全$\operatorname{Sym}^2(\mathcal D_2)$の
有限射と、平方因子を除いた
$\mathcal R_{\mathrm{sf}}$上の制限射は区別する。
\end{theorem}
\begin{proof}
四次形式の任意の因数分解は
その根の重複度付き四点集合から
二点ずつを選ぶ有限個の操作なので、
$\mu$の幾何fiberは有限である。
$\operatorname{Sym}^2(\mathcal D_2)$は射影的、
従って$\mu$はproperかつquasi-finiteで
有限である。代数閉体上ではすべての四次形式が
二次形式の積となるので全射。
四根が相異なれば分割は三つあり、
一般次数は三となる。
式~\eqref{eq:qa-rootpair-resolvent-cubic}は
四根の基本対称式を展開すれば得られる。
例えば三値を
$u_{12|34}=\alpha_1\alpha_2+\alpha_3\alpha_4$
などと置くと、
$\sum u=b$、
$\sum_{r<s}u_ru_s=ac-4d$、
$\prod u=a^2d+c^2-4bd$である。

$f_1,f_2$の根が各々相異なるので、
積の各根の重複度は高々二である。
共通根の個数が$0,1,2$の三場合が
表の各行に一致する。
四単根なら二重被覆はsmoothでRiemann--Hurwitzにより
genus一。一つの二重根の近くでは
$y^2=t^2u(t)$、$u(0)\ne0$なので、
étale局所で$y'^2=t^2$という普通節点になる。
残り二単根から二重被覆は既約で、
正規化はgenus零。
二つの二重根なら$f_1=f_2$（射影類として）で、
$y^2=f_1^2$は二つのsection
$y=\pm f_1$の和である。
各sectionは$\PP^1$で、$f_1$の二単根で横断する。
二重被覆の代数
$\OO_{\PP^1}\oplus\OO_{\PP^1}(-2)y$
から$\chi(\OO_{C_q})=1+(-1)=0$であり、
全型で$p_a(C_q)=1$。
局所的に平面曲線の節点またはsmoothなので
Gorensteinでもある。
群の作用については積と二重被覆の
構成が自然であることから従う。
\end{proof}

\begin{theorem}[根対からの相対torsion-free spectral avatar]
\label{thm:qa-noncubic-rootpair-torsionfree-spectral-module}
論文Bの二形式表示
\[
 (f_1,f_2)\in
 \bigl(\operatorname{Sym}^2W^\vee\bigr)_{\mathrm{nd}}^2
\]
上で、$B=\PP(W)$と置き、
$\widetilde p:\widetilde C\to B$、
$p:C\to B$を以下の有限射、
$r:\widetilde C\to C$をその間の自然な射とする。
\[
 \widetilde C=
 \operatorname{Spec}_{\PP(W)}
  \frac{\OO_{\PP(W)}[c,d]}
       {(c^2-f_1,\ d^2-f_2)},\qquad
 C=\operatorname{Spec}_{\PP(W)}
  \frac{\OO_{\PP(W)}[y]}{(y^2-f_1f_2)}
\]
を取る。ただし$c,d$は局所的に
$\OO_{\PP(W)}(-1)$の生成元、$y=cd$は
$\OO_{\PP(W)}(-2)$の生成元と解釈する。
$\widetilde C$から$C$への有限射は
$(c,d)\mapsto(-c,-d)$の商であり、
$\OO_C$-加群としての反不変部分と
その$B$への直像はそれぞれ
\[
 \mathcal M=(r_*\OO_{\widetilde C})^-,
 \qquad
 p_*\mathcal M
 \cong\OO_B(-1)c\oplus\OO_B(-1)d
 \tag{Res-M}\label{eq:qa-rootpair-spectral-module}
\]
であり、$\mathcal M$は
\[
 y c=f_1 d,\qquad y d=f_2 c
 \tag{Res-MF}\label{eq:qa-rootpair-spectral-matrix-factorization}
\]
を満たす。全二形式族と任意の基底変更にわたり、
$C$は有限平坦二重被覆、$\mathcal M$は
\emph{相対平坦なrank-one torsion-free sheaf}
（各幾何fiber上で次数零）となる。
この構成は二形式の基底変更・スカラー変更・
二形式交換に関して自然なので、
論文Bのnoncubic表示スタック
$\mathfrak N^{\mathrm{pres}}=[P/G]$から
「種数一の有限二重被覆とrank-one
torsion-free sheaf」の相対群oidへの射を与える。

二根対が交わらない滑らかなfiberでは、
$\widetilde C\to C$はétale $\mu_2$-torsor、
$\mathcal M$は非自明な
$\Pic^0(C)[2]$の線束である。
これは定理
~\ref{thm:qa-avatar-two-torsion-root-pair-bijection}
で定めた分割の2-torsionと一致する。
他方、根対が一根を共有するとき
$\mathcal M$はその節点で局所自由でなく、
二根対が一致するときは二つの節点で
局所自由でない。
従って\emph{全noncubic根対配置にわたって}
自然に存在するのは通常の$\Pic^0$線束ではなく、
この相対torsion-free spectral moduleである。
\end{theorem}
\begin{proof}
二つの二次平方根の代数は、
$\OO_B$上の基底$1,c,d,cd$を持つ
階数四の局所自由代数である。
$\widetilde p_*\OO_{\widetilde C}$の
同時符号変換の不変部分は$1,cd$、
反不変部分は$c,d$なので
式~\eqref{eq:qa-rootpair-spectral-module}を得る。
$y=cd$による積は
式~\eqref{eq:qa-rootpair-spectral-matrix-factorization}
で与えられ、形式的には
\[
 \begin{pmatrix}0&f_2\\f_1&0\end{pmatrix}^2
      =f_1f_2 I_2
\]
というmatrix factorizationである。
$p_*\mathcal M$は$B$上で
$\OO(-1)^{\oplus2}$だから、
二形式パラメータとその任意基底変更に対して平坦。
各幾何fiberでは二次拡大のgeneric点でrank一で、
局所的に二次平方根代数の直和因子であるから
torsion-freeである。
Euler特性は
\[
 \chi(C,\mathcal M)
 =\chi(\PP^1,\OO(-1)^{\oplus2})=0
 =\chi(C,\OO_C)
\]
なので次数零である。

$f_1,f_2$に共通零点がなければ、
$c,d$が同時に零となる点がないため
同時符号対合は自由で、
$\widetilde C\to C$は有限étaleの二重torsor。
このとき反不変部分$\mathcal M$は
invertibleかつ$\mathcal M^2\simeq\OO_C$。
$H^0(C,\mathcal M)=H^0(B,p_*\mathcal M)
 =H^0(\PP^1,\OO(-1)^{\oplus2})=0$
なので非自明である。
$\widetilde C$の三つの二次部分被覆の
分岐分割から、これは前定理の2-torsionに一致する。

共有根の局所では単元の平方根をétale局所に取り、
\[
 \OO_{C,o}^{\wedge}
  \simeq k[[x,y]]/(y^2-x^2),\qquad
 \mathcal M_o\cong
   \langle c,d\mid yc=xd,\ yd=xc\rangle .
\]
$\mathcal M_o/(x,y)\mathcal M_o$の次元は二なので、
rank一局所自由ではない。
一致根対では$C$は二つのsectionの和で、
その二つの交点の局所方程式も同じになり、
$\mathcal M$は各節点で非局所自由である。
二形式の交換は$c,d$の交換、
スカラー変更は平方根のスカラー変更として
作用し、構成は元の表示群の作用に整合する。
\end{proof}

\begin{corollary}[noncubic中央二重平方環の可換点曲線からの幾何射]
\label{cor:qa-noncubic-central-square-point-curve-spectral-avatar}
論文A・Bのnoncubic中央二重平方環
\[
 B_f=k[u,v]\langle c,d\rangle/
          (c^2-f_1(u,v),\,d^2-f_2(u,v))
\]
に対し、通常の可換化を施した環と
その射影曲線は
\[
 B_f^{\mathrm{ab}}=
 k[u,v,c,d]/(c^2-f_1,\,d^2-f_2),\qquad
 E_f=\operatorname{Proj}(B_f^{\mathrm{ab}})
 \subset\PP^3
\]
である。これは定理
~\ref{thm:qa-noncubic-rootpair-torsionfree-spectral-module}
の四重被覆$\widetilde C$と自然に同型であり、
\[
 \boxed{\quad
 E_f\xrightarrow{/\,\mu_2}C_{f_1f_2}
 \xrightarrow{\;2:1\;}\PP^1,
 \qquad
 \mathcal M=(r_*\OO_{E_f})^- .
 \quad}
 \tag{NC-point-avatar}\label{eq:qa-noncubic-pointcurve-spectral-factorization}
\]
$\mu_2$は$(c,d)\mapsto(-c,-d)$として作用する。
この対合の固定点は$f_1,f_2$の
共通根の上に限られるので、二根対が
非交差なら第一射はétale、
一根共有または一致根対では
固定点に対応する節点でétaleでなくなる。
この場合、有限平坦な二重torsorと
みなしてはならない。

この射はnoncubic線形代数$B_f$から
rank-one torsion-free spectral avatarへの
\emph{具体的な幾何構成}である。
非可換な$B_f$の$\qgr$と
$\OO$-marked avatar側のheartが同値である
と主張するものではない。
\end{corollary}
\begin{proof}
$u,v$は中央にあるので、可換化において
追加される関係は$cd=dc$だけである。
そのProJの$\PP^1=\operatorname{Proj}k[u,v]$上の
相対座標環は$1,c,d,cd$を基底とする
二重平方根代数となる。
これは前定理の$\widetilde C$そのものである。
同時符号の不変部分は$1,cd$、
反不変部分は$c,d$なので、
商と$\mathcal M$の表示も一致する。
非自明な同時符号の固定点は、
射影座標$(u,v)$が非零である限り
$c=d=0$を満たす点に限られ、
これは$f_1=f_2=0$と同値である。
$(u,v)=(0,0)$では二つの平方関係から
$c=d=0$となり射影点は存在しない。
従って固定点の条件が従う。
\end{proof}

\begin{proposition}[共有根退化での可逆Picard延長の局所障害]
\label{prop:qa-rootpair-nodal-picard-nonextension}
$R=k[[s]]$上の二根対の退化
\[
 f_1=x(x-1),\qquad f_2=(x-s)(x-2)
\]
を取り、$C_R:y^2=f_1f_2$とする。
その一般fiberは滑らかな楕円曲線であり、
特殊fiberは既約な一節点曲線となる。
一般fiberの自然な非零2-torsion線束
$\mathcal M_K$は、
\emph{$C_R$全体の可逆線束には延長できない}。
従って二根対から
$C:y^2=f_1f_2$と
\emph{可逆な}$P\in\Pic^0(C)$を
割り当てる滑らかなavatar辞書は、
この族に沿ってそのまま相対射としては延長できない。

この障害は、$C_R$の節点における
表面$A_1$特異点と、
そこでの非自明なrank-one reflexive module
によって検出される。
$C_R$以外の曲線へのavatar再埋込み、
$\Sigma_2$へのelementary transform、
またはcompactified Picard型データ
\cite{AltmanKleiman1980}による
延長の可能性は排除しない。
ただし同じ一径数基底の上で全空間の$A_1$特異点を
最小解消するだけでは不十分である。
命題~\ref{prop:qa-nodal-resolution-odd-exceptional-degree}の
反不変moduleの可逆化は例外曲線上の次数が奇数であり、
任意の垂直Cartier因子でtwistしても
相対的な多重次数零線束にはならない。
第~\ref{sec:qa-spectral-boundary-models}節の二traitでは，
任意の有限DVR基底変更後の最小正則模型でも二乗自明な線束への延長は存在しない
（定理~\ref{thm:qa-boundary-input-two-torsion-line-obstruction}）．
線束自体への延長と，非可逆TF層・二次積を保持する標的は区別する．
\end{proposition}
\begin{proof}
$s=0$では$f_1f_2=x^2(x-1)(x-2)$であり、
特殊曲線は既約nodal genus一である。
共通根$x=s=0$の完成局所環では、
単元の平方根で座標を正規化すると
\[
 A=k[[x,s,y]]/(y^2-x(x-s)).
\]
二形式を別々に平方根化する被覆の完成環は
$B=k[[c,d]]$で、
\[
 x=c^2,\quad x-s=d^2,\quad y=cd,\quad
 A\simeq B^{\{\pm1\}}
\]
となる。$B$上の同時符号の反不変部分
$M=Ac+Ad$は、商$A$上の
rank-one maximal Cohen--Macaulay、従って
rank-one reflexive moduleである。
その最小生成元は$c,d$の二つなので非自由であり、
$A$は表面$A_1$特異点を持つ。
これは前定理の$\mathcal M$の完成stalkである。

相対曲線の全空間$C_R$はnormalであり、
その唯一の特異点で上記$A_1$局所型を持つ。
$\mathcal M$はnormal曲面上の
rank-one reflexive sheafで、
一般fiberでは$\mathcal M_K\in\Pic^0(C_K)[2]$。
もし可逆線束$L$が全空間上に存在して
$L_K\simeq\mathcal M_K$となれば、
normal曲面上のrank-one reflexive sheafの
Weil因子記述から、$L$と$\mathcal M$の差は
特殊fiberに支持されるWeil因子となる。
この特殊fiberは\emph{既約}かつ
パラメータ$s$で切られたCartier因子なので、
その整数倍は主因子である。
従って$L\simeq\mathcal M$となり、
節点での非局所自由性に矛盾する。
最後の留保は$C_R$を固定したときの
可逆延長障害と、別のavatarモデルの存在を
混同しないためである。
\end{proof}

\begin{theorem}[全noncubic根対と二捩れsquare avatarの体上同型類全単射]
\label{thm:qa-noncubic-rootpair-all-geometric-square-avatar-bijection}
$k$を標数零の代数閉体とする。
論文Bのnoncubic根対の全幾何点、すなわち
$\PP^1$上の\emph{各々は被約}な二点因子の
非順序対$\{B_1,B_2\}$の
$\PGL_2$-同型類集合を考える。
これは、次の条件を持つsquare avatarの
\emph{同型類集合}と自然に全単射する。
\[
 \begin{gathered}
 (Y=\PP^1\times\PP^1;\
  D\in|\OO_Y(2,2)|;\ F_s,G_s,P),\\
 P=G_sF_s^{-1}\ne\OO_D,\qquad
 P^{\otimes2}\simeq\OO_D,\\
 D\text{ はsmooth genus }1,\quad
 \text{既約一節点、または二}\PP^1
 \text{が二節点で交わる曲線}.
 \end{gathered}
 \tag{Sq-all-nc}\label{eq:qa-all-noncubic-square-avatar-bijection}
\]
ここで$F_s=\OO_D(1,0)$、
$G_s=\OO_D(0,1)$であり、
二rulingの順序は保持する。
この全単射は、論文Bの
\cite[\texttt{thm:noncubic-intrinsic-presentation-equivalence}]{QuadricModuli2026}
を介して、\emph{全noncubic O-marked H2幾何同型類}の
集合論的なavatar分類を与える。

構成は次の通りである。
まず$q=f_1f_2$から
式~\eqref{eq:qa-rootpair-spectral-curve}の
二重被覆$p:C_q\to\PP^1$と
$F_s=p^*\OO_{\PP^1}(1)$を取る。
\begin{enumerate}
\item 四根相異なる場合は、
定理~\ref{thm:qa-avatar-two-torsion-root-pair-bijection}
の分割に対応する非零
$P\in\Pic^0(C_q)[2]$を取る。
\item 一根共有または二根対一致の場合は
$\Pic^0(C_q)\simeq\mathbb G_m$なので、
その\emph{唯一の}非自明二捩れ点$P=-1$
を取る。
\end{enumerate}
いずれの場合も$G_s=F_sP$は次数二で
大域生成され、二つの完全線形系の積
\[
 C_q\longrightarrow
 \PP H^0(F_s)^\vee\times\PP H^0(G_s)^\vee
\]
は$(2,2)$反標準Cartier因子への閉埋込みとなる。
この逆操作は、square avatarの
第一rulingの二重被覆$D\to\PP^1$の
分岐四次因子と$P$から、
四根相異なる場合には三分割の一つを、
一根共有・一致根対では唯一の
被約二根対分割を復元することである。

\emph{重要：} これは\emph{体上の同型類の全単射}である。
上の二場合で特別に選んだ
$P=-1\in\mathbb G_m$は、
必ずしも滑らかな族の2-torsion線束の極限でない。
実際、命題
~\ref{prop:qa-rootpair-nodal-picard-nonextension}
の一根共有族では延長できない。
従ってこの全単射を、そのまま
通常のPicard avatarスタックへの
相対モジュライ射、または元のavatar
AS代数の$\qgr$のpointed同値と
読み替えてはならない。
\end{theorem}
\begin{proof}
smoothな場合は定理
~\ref{thm:qa-avatar-two-torsion-root-pair-bijection}
による。以下、特異な二場合を調べる。

\emph{既約一節点の場合。}
$C_q$の正規化は$\nu:\PP^1\to C_q$で、
節点の二逆像を$a,b$とする。
$\Pic^0(C_q)$の線束は、
正規化上の自明線束の
$a,b$での貼合せスカラー
$\lambda\in k^\times$で分類される。
従って$\Pic^0(C_q)\simeq\mathbb G_m$で、
非自明二捩れは貼合せ比$-1$だけである。

$F_s=p^*\OO(1)$は正規化上次数二で、
節点での貼合せ比を$1$に規格化できる。
$G_s=F_sP$は正規化上の$\OO(2)$で、
貼合せ比は$-1$となる。
$H^0(C_q,G_s)$は$H^0(\PP^1,\OO(2))$の
$a,b$における値が符号反対になる
二次元部分空間である。
これはbasepoint-freeである。
実際、$a,b$の二評価汎関数の
非自明な線形結合は、
二次Veroneseの双対円錐上で
単一点の評価汎関数には一致しない。
（異なる二点を結ぶ割線が二次曲線と
交わるのはその両端のみである。）
従って$|G_s|$も有限な次数二射を与える。

$F_s\not\simeq G_s$なので、
この二射の積は総次数二ではあり得ず、
$C_q$上の一般点で単射である。
像の二rulingへの次数は各二だから
$(2,2)$因子で、算術種数は一。
元の$C_q$も算術種数一なので、
有限双有理射の構造層の余核はEuler特性零、
従って零であり、積射は閉埋込みとなる。

\emph{二根対一致の場合。}
$C_q=C_+\cup C_-$は二つの$\PP^1$が
相異なる二節点で交わる曲線である。
正規化上の各成分の次数零線束は自明であり、
二つの節点での貼合せ比
$(\lambda_1,\lambda_2)$を
成分ごとの定数倍で同一視すると、
$\lambda_1/\lambda_2\in\mathbb G_m$
だけが残る。
よって$\Pic^0(C_q)\simeq\mathbb G_m$、
非自明二捩れはこの比が$-1$のものに限る。
$F_s$は各成分で$\OO_{\PP^1}(1)$であり、
$G_s=F_sP$も各成分で次数一、
二節点での値の二条件から
$H^0(C_q,G_s)$は二次元で大域生成される。
二つの成分の像は、
第一$\PP^1$から第二$\PP^1$への
二つの$(1,1)$グラフとなる。
貼合せ比$-1$ゆえ二グラフは異なり、
ちょうど二つの節点で横断する。
従って像は元の$C_q$と同型な
$(2,2)$Cartier因子である。

逆向きは平方freeな二つの二次因子の
分割を数えることで得られる。
四単根なら三通りで非零$E[2]$の三点に対応する。
$(2,1,1)$なら
$(a,b)\mid(a,c)$の一通り、
$(2,2)$なら
$(a,b)\mid(a,b)$の一通りしか
二因子ともsquarefreeにする分割はない。
同じ各曲線の$\Pic^0[2]\setminus0$の
点数が$3,1,1$なので、逆は一意である。
元の第二rulingは
$G_s=F_sP$で復元されるから、
曲線・線束の同型はavatarの同型となる。
最後にBのnoncubic根対分類と合成すれば
O-marked同型類への全単射が得られる。
相対射に関する留保は直前命題の
非延長例により必要である。
\end{proof}

\begin{theorem}[全noncubic O-marked同型類の$\Sigma_2$反標準avatar分類]
\label{thm:qa-all-noncubic-sigma2-anticanonical-avatar-bijection}
$k$を標数零の代数閉体とする。
$\Sigma_2=\PP_{\PP^1}(\OO\oplus\OO(2))$の
負断面を$S=h-2f$とし、
$D\in|2h|=|-K_{\Sigma_2}|$を
$D\cap S=\varnothing$を満たす
相対反標準Cartier因子とする。
このとき、次の二集合は自然に全単射する。
\[
 \boxed{
 \begin{aligned}
 &\{\text{全noncubic O-marked quadricの
             幾何exact同型類}\}\\
 &\quad\longleftrightarrow
 \left\{
 \begin{array}{l}
 (D\subset\Sigma_2,P),\quad
 D\cap S=\varnothing,\\
 0\ne P\in\Pic^0(D)[2]
 \end{array}
 \right\}
/\operatorname{Aut}(\Sigma_2;h,f).
 \end{aligned}}
 \tag{NC-Sigma2}\label{eq:qa-all-noncubic-sigma2-avatar}
\]
右辺で$D$はあらかじめ滑らか、既約、被約と
仮定しない。しかし$P$の非自明二捩れ性から、
実際には$D$はsmooth genus一、
既約一節点、または二つの$\PP^1$が二節点で
交わる型のいずれかになる。
この三型は論文Bの二根対の
非交差／一根共有／一致に正確に対応する。

具体的な同型類の対応は
定理~\ref{thm:qa-noncubic-rootpair-all-geometric-square-avatar-bijection}
のsquare avatar
$(D_s;F_s,G_s,P)$に、
系~\ref{cor:qa-square-chart}の
幾何的変換$\mathsf T^{-1}$を適用したものである。
$\Sigma_2$側の制限線束を
$F=\OO_{\Sigma_2}(f)|_D$、
$H=\OO_{\Sigma_2}(h)|_D$とすると、
\[
 H\simeq F^2,\qquad
 F_s=F,\quad G_s=FP,\quad
 H_s=F_sG_s=F^2P
 \tag{NC-Sigma2-dict}\label{eq:qa-noncubic-sigma2-square-dictionary}
\]
が全対応の明示公式となる。

これは\emph{体上の集合論的全単射}であり、
$\Sigma_2$ avatarから元のAS代数の
pointed heartを直接回復する定理でも、
Picard line bundleを使う相対スタック同値でもない。
smooth域でこのsquare換算を指定AS10 mutationと一番頂点の
pointingへ接続した対象は、
定理~\ref{thm:qa-spectral-avatar-pointed-involution}により
元の対象の二同種双対に対応する。
元とpointed同値となるのは$\chi=0,2$だけであり、
この対応先は二つのnoncubic境界でH2族にも延長できない
（系~\ref{cor:qa-spectral-involution-no-h2-extension}）。
命題~\ref{prop:qa-rootpair-nodal-picard-nonextension}
の共有根退化では、固定されたspectral曲線族上の
可逆$P$への相対延長が不可能なことに注意する。
\end{theorem}
\begin{proof}
$D\cap S=\varnothing$のとき、
$S=h-2f$の定義から
$\OO_D(S)\simeq\OO_D$となり、
$H\simeq F^2$が従う。
また$D$はどの$\PP^1$-fiberも成分として含まない。
実際、fiberは$S$と交わるので、
そのような成分は$D\cap S=\varnothing$に反する。
従って$p:D\to\PP^1$は有限平坦次数二である。
そのtrace-zero直和分解を取ると
\[
 p_*\OO_D\simeq\OO_{\PP^1}\oplus
                \OO_{\PP^1}(-2)
\]
となる。後項の次数は
$\chi(\OO_D)=0$と$p_a(D)=1$から決まる。
よって$D$は四次式$q$を用いた
$y^2=q$型の二重被覆になる。
$q$が四単根、一つの二重根、二つの二重根なら
前定理の三型になる。
それ以外の四次根配置は
$(3,1)$のcuspidal型、
$(4)$の二枝接触型、
または$q=0$の非被約double sectionである。
これらの$\Pic^0$は標数零で
加法型$\mathbb G_a$となり、
非自明二捩れ点を持たない。
この主張は正規化と特異点での単元の貼合せ、
非被約double sectionでは
$1+\OO_{\PP^1}(-2)$の加法的拡大から得られる。
従って$P\ne\OO_D,\ P^2=\OO_D$は
三つのsmooth/nodal/I2型を強制する。

幾何的変換については
系~\ref{cor:qa-square-chart}が、
$\Sigma_2$の負断面を避けるavatarを
square chartへ送る操作$\mathsf T$を
任意基底変更と可換な同値として構成している。
そこで$H=F^2$を代入すると
square側の線束は
\[
 F_s=F,\qquad G_s=HPF^{-1}=FP
\]
となり、$P=G_sF_s^{-1}$かつ$P^2=\OO_D$である。
逆にこの二捩れsquare wallのデータに
$\mathsf T^{-1}$を適用すれば
$H=F_sG_sP^{-1}=F_s^2$となり、
系~\ref{cor:qa-square-chart}と
定理~\ref{thm:qa-open-charts}により
負断面を避ける$\Sigma_2$に入る。
両操作は互いに逆であり、
曲面・ruling・因子・Picard類の同型を保つ。
これと前定理の全noncubic二根対の
幾何同型類全単射、論文Bの
\cite[\texttt{thm:noncubic-intrinsic-presentation-equivalence}]{QuadricModuli2026}
を合成して結論を得る。
相対射に関する留保は直前の
$A_1$局所非延長の結果による。
\end{proof}

\begin{corollary}[noncubic良moduliの同じ点を区別する二つの$\Sigma_2$ avatar]
\label{cor:qa-sigma2-avatar-separates-special-noncubic-classes}
論文Bのnoncubic境界の良moduli
$\chi:\mathfrak N\to\mathbb A^1$において、
一根共有と二根対一致はともに$\chi=1$へ写るが、
これらは異なるpointed exact同型類である。
定理~\ref{thm:qa-all-noncubic-sigma2-anticanonical-avatar-bijection}
の$\Sigma_2$ avatarによれば、
前者は\emph{既約nodal}反標準因子$D$と
唯一の非零二捩れ$P$、
後者は\emph{二$\PP^1$が二節点で交わる}
反標準因子$D$と唯一の非零二捩れ$P$に対応する。
従って曲線$D$の既約成分構成が、
良moduliでは同一視される二つの
O-marked同型類を区別する。
\end{corollary}
\begin{proof}
良moduliの特殊fiberの二幾何同型類と
$\chi=1$への写像は論文B
\cite[\texttt{prop:noncubic-good-moduli-fibers}]{QuadricModuli2026}
による。
根対の共通根の個数に応じた
二重被覆$C_q$の特異型の分類は
定理~\ref{thm:qa-noncubic-rootpair-quartic-resolvent}
にあり、$\Sigma_2$での再埋込みは
前定理で証明した。
既約性・成分数が異なるので
二avatarは同型になり得ない。
\end{proof}

\begin{remark}[節点型avatarの二つの明示$(2,2)$模型]
\label{rem:qa-singular-square-avatar-explicit-models}
定理~\ref{thm:qa-noncubic-rootpair-all-geometric-square-avatar-bijection}
の特異な二型は、いずれも
$\PP^1_X\times\PP^1_Z$上で明示できる。
既約一節点型では、正規化$\PP^1_t$の
$t=1,-1$を貼り合わせ、非自明二捩れの
貼合せ比を$-1$とする。
$F_s$の二切断は$1,t^2$、
$G_s$の二切断は$t,1-t^2$であり、
\[
 X=t^2,\qquad Z=\frac{1-t^2}{t}
 \quad\Longrightarrow\quad
 XZ^2=(1-X)^2 .
 \tag{Sq-node}\label{eq:qa-nodal-square-avatar-model}
\]
これは$(X,Z)=(1,0)$に普通節点を持つ
既約$(2,2)$曲線である。
式は各$\PP^1$で斉次化するものとする。

一致根対型では、$C_+\cup C_-$を二つの
$\PP^1$の和として、その二交点を
$X=0,\infty$に取る。$F_s$は両成分で
$X$座標、非自明二捩れ$P$による
$G_s$は第一成分で$Z=X$、
第二成分で$Z=-X$となる。
従って
\[
 (Z-X)(Z+X)=0
 \tag{Sq-I2}\label{eq:qa-i2-square-avatar-model}
\]
は二つの$(1,1)$曲線が
$(0,0),(\infty,\infty)$で横断する
反標準$(2,2)$因子を与える。
この二模型は、特異な場合の
$P=-1\in\Pic^0(D)[2]$が本当に
square avatarの二rulingを作ることを
直接に示している。
\end{remark}

\begin{proposition}[Picard零切断の全体を古典pointとして追加できない理由]
\label{prop:qa-avatar-zero-picard-collapse}
$Y=\PP^1\times\PP^1$、
$D\in|\OO_Y(2,2)|$を任意の反標準Cartier因子とし、
$P=\OO_D$と置く。
このとき曲線上の四線束
\[
 (\OO_D,\ F_s,\ G_s,\ F_sG_s)
\]
の有向End代数（各積を含む）は、
$D$の選択に依存せず、
$Y$上の可換quadricの標準四項collection
\[
 (\OO_Y,\OO_Y(1,0),\OO_Y(0,1),\OO_Y(1,1))
\]
の有向End代数と同型である。
従って$\Pic^0$の零切断をそのまま
曲線型avatarに追加すると、
異なる$D$が同一の可換foundation代数を与え、
O-marked分類への単射性が失われる。
古典的な可換quadricは、曲線型とは別に
一点として追加する必要がある。
\end{proposition}
\begin{proof}
相異なる四項の間のHomに現れる
$\OO_Y(a,b)$は、
$a,b\in\{-1,0,1\}$の範囲にある。
各々についてCartier制限列
\[
 0\to\OO_Y(a-2,b-2)\to
 \OO_Y(a,b)\to\OO_D(a,b)\to0
\]
を考えると、$a-2\le-1$、$b-2\le-1$であり、
Künnethの公式から
$H^0(Y,\OO(a-2,b-2))=
H^1(Y,\OO(a-2,b-2))=0$である。
従って四項間の全Homの制限は同型となる。
構成は線束の合成に関して自然なので、
有向End代数の積も保存する。
右辺は$D$に依存しない。
\end{proof}

\begin{remark}[二つの $\Sigma _2$ への移行]
\label{rem:qa-rotated-square}
square helix の foundation を一つ進めた rotated square の偏極は
\[
 F_r=F_s,\qquad H_r=F_sG_sP
\]
である（本稿の幾何的証明は定理~\ref{thm:qa-core-avatar-square-rotated}；
原草稿\cite[命題11.1・定理15.1]{AvatarVSeven2026}とも整合する）。
一方、\eqref{eq:qa-square-dictionary}の偏極は $H=F_sG_sP^{-1}$ である。
前者の $\Sigma _2$ 条件は $P=F_sG_s^{-1}$、後者の条件は $P=G_sF_s^{-1}$ となる。
これらは一般には別の壁である。
foundation の移動や mutation とともに選ばれる構造対象も変わるため、
固定された構造対象の内在的一点 Hilbert 関手を同じものとして扱わない。
\end{remark}

\subsection{古典点と最深境界の扱い}

本節の stack は一貫して $P$ が各 fiber で非自明な曲線型を扱った。
体上の分類では、$P=\OO_D$ とした有限 directed algebra は $D$ の選択に依存せず、
\[
 \mathrm{Cl}_0:\operatorname{coh}(\PP^1\times\PP^1),\qquad
 \mathrm{Cl}_2:\operatorname{coh}(\PP(1,1,2))
\]
という二つの古典型に対応する
\cite[注意11.3・定理15.3]{Q12Avatar2026}。
後者の $\PP(1,1,2)$ は重み付き射影 stack である。
これは同型類の説明であり、相対 Picard の零切断を二点へ自動的に縮約する
stack の構成ではない。零切断では $D$ とその自己同型が残るため、
相対的に古典点を加えるには別途の収縮または同一視を定義する必要がある。

最深の十二本矢層についても、三頂点の有限安定表現のモジュライが $\Sigma _2$ になることと、
標準 heart における内在的一点 Hilbert 曲面が $\Sigma _2$ になることは異なる。
この区別のための内在的な検査として，
例~\ref{ex:qa-weighted-not-h2}では重み付き古典模型の全標準頂点で，
例~\ref{ex:qa-hirzebruch-not-h2}では $\coh\Sigma_2$ 自身の全可逆層で，
反Serre軌道の生成条件が破れることを示した。
古典型の負断面上には、一つの foundation では安定であっても、
次の foundation では最後の頂点から生成されなくなる表現が存在する
\cite[命題14.1]{Q12Avatar2026}。
したがって本節の開部分 stack だけから、全 $\mathcal A_{12}^{\times}$ の
内在的 quadric モジュライへの射、あるいはその射の定義域が開であることを結論しない。
そこでは構造対象、標準 heart、豊富性、Serre 平方根および相対 AS 持上げを
別に検証しなければならない。

% !TeX root = ../nc-quadric-avatar.tex
% 2026-11-22 / two genuinely distinct genus-one curves in noncubic comparison
\section{Noncubic の点曲線と spectral 曲線の二重 isogeny}
\label{sec:qa-point-spectral-isogeny}

noncubic の中央二重平方環から得る可換点曲線と，
二根対の四次 resolvent から得る spectral 曲線は，
滑らかな場合でさえ一般に同型ではない．
両者の関係を明示することは，avatar の幾何的分類を
内在的 O-marked heart の同値と取り違えないために必要である．

\begin{theorem}[点曲線とspectral曲線の二重 isogeny]
\label{thm:qa-point-spectral-dual-two-isogeny}
$k$ を標数 $0$ の代数閉体，
$f_1,f_2\in k[u,v]_2$ を squarefree かつ互いに
根を共有しない二次形式とする．
\[
 \begin{split}
 E_f&=\{c^2=f_1(u,v),\ d^2=f_2(u,v)\}
                   \ \subset\PP^3_{[u:v:c:d]},\\
 C_q&=\{y^2=f_1(u,v)f_2(u,v)\}
             \ \longrightarrow\PP^1_{[u:v]}
 \end{split}
\]
と置く．$E_f$ は
定理~\cite[\texttt{thm:ore-anticommuting-central}]{QuadricIntrinsic2026}
に現れる中央二重平方環の可換点曲線であり，
$C_q$ は
式~\eqref{eq:qa-rootpair-spectral-curve}の
spectral 曲線である．
次が成り立つ．
\begin{enumerate}
\item $E_f$ と $C_q$ は滑らかな連結種数一曲線である．
 同時符号対合
 $\tau(u,v,c,d)=(u,v,-c,-d)$ は $E_f$ に自由に作用し，
 \[
   r:E_f\longrightarrow C_q,\qquad
   [u:v:c:d]\longmapsto([u:v],y=cd)
   \tag{Point--spectral}\label{eq:qa-point-spectral-cover}
 \]
 はその有限 \'{e}tale 二重商である．
 原点を選べば $r$ は次数二の isogeny となる．
\item $\alpha(c,d)=(-c,d)$ と $\beta(c,d)=(c,-d)$ とする．
 これらの商は滑らかな conic
 \[
    Q_d=\{d^2=f_2(u,v)\},\qquad
    Q_c=\{c^2=f_1(u,v)\}
 \]
 であり，商射を $p_d:E_f\to Q_d$，
 $p_c:E_f\to Q_c$ と書く．
 $Q_d,Q_c\simeq\PP^1$ 上で各々唯一の次数一線束
 $N_d,N_c$ を取り，
 $F_d=p_d^*N_d,\ F_c=p_c^*N_c$ とすると
 \[
   \deg F_d=\deg F_c=2,\quad
   F_d^{\otimes2}\simeq F_c^{\otimes2}
       \simeq\OO_{E_f}(1),\quad
   0\ne P_E:=F_cF_d^{-1}\in\Pic^0(E_f)[2].
   \tag{Point-avatar}\label{eq:qa-point-curve-double-wall}
 \]
 従って $(E_f,F_d,F_c,P_E)$ 自身が
 二重壁上の square avatar を与える．
 また $\tau$ は $P_E$ による非自明二捩れ平行移動である．
\item
 $M_C=(r_*\OO_{E_f})^-$ は
 定理~\ref{thm:qa-noncubic-rootpair-torsionfree-spectral-module}
 の非自明二捩れ線束である．
 その $\Pic^0(C_q)[2]$ の点は，
 $r$ の\emph{双対 isogeny}の核を生成する：
 \[
     0\longrightarrow
       \{\OO_{C_q},M_C\}
       \longrightarrow\Pic^0(C_q)
       \xrightarrow{\ r^*\ }\Pic^0(E_f)
       \longrightarrow0.
       \tag{Dual-2}\label{eq:qa-point-spectral-dual-isogeny}
 \]
 特に原点の選択の下で
 $E_f\simeq C_q/\langle M_C\rangle$ である．
\end{enumerate}
\end{theorem}
\begin{proof}
$E_f\to\PP^1$ は $c,d$ の符号を独立に変える
biquadratic な次数四被覆である．
$f_1,f_2$ の四分岐点が互いに異なるため，
各分岐点では一方の平方根だけが分岐し，
$E_f$ は滑らかで連結である．
Riemann--Hurwitz により
\[
 2g(E_f)-2=4(-2)+4\cdot2=0.
\]
同時符号対合の固定点では$c=d=0$を要するが，
これは$f_1=f_2=0$を意味して仮定に反する．
この対合の不変関数は
$u,v$ と $cd=y$ で，$y^2=f_1f_2$であるから，
商は$C_q$である．
商は非分岐二重被覆で，$C_q$も種数一となる．

$\alpha$ と $\beta$ の不変式から，それぞれの
商が $Q_d,Q_c$ であることが直ちに従う．
squarefree 性によって conic は滑らかであり，
$k$ は代数閉だから $\PP^1$ に同型である．
Conic の超平面線束は $\OO_{\PP^1}(2)$なので，
その平方根となる次数一線束は一意に存在する．
$E_f$ の埋込の線形座標
$(u,v,d)$，$(u,v,c)$ はそれぞれ
$Q_d$，$Q_c$ の超平面切断の引戻しだから
$F_d^2=F_c^2=\OO_{E_f}(1)$となる．
それぞれ次数二で大域生成されており，
相異なる二つの pencil は
補題~\ref{lem:qa-core-two-pencils}により
$E_f$を$\PP^1\times\PP^1$の$(2,2)$因子として
閉埋め込みする．
もし $F_d\simeq F_c$ なら二つの次数二線形系は
基底の射影変換を除いて同じ射を定め，
その deck involution も一致するはずである．
しかし $\alpha\ne\beta$ なので非自明性を得る．

滑らかな種数一曲線の次数二線束 $L$ は，
原点を選んで$x\mapsto \ell_L-x$という
involution を定める．
従って二線束$F_d,F_c$の involution
$\alpha,\beta$の積は，その差
$F_cF_d^{-1}=P_E$による平行移動である
（符号の違いは二捩れゆえ消える）．
これで第二項が従う．

最後に$E_f\to C_q$は \'{e}tale な
$\mu_2$-torsorであり，
$r_*\OO_{E_f}=\OO_{C_q}\oplus M_C$，
$M_C^{\otimes2}\simeq\OO_{C_q}$である．
$H^0(C_q,M_C)=0$は
式~\eqref{eq:qa-rootpair-spectral-module}から従う．
従って $M_C$ は非自明である．
有限 \'{e}tale 同種写像の双対性，
または線束の Galois descent から，
$\ker r^*=\{\OO_{C_q},M_C\}$となる．
次数二の isogeny の双対も次数二なので
$r^*$ は全射となり，最後の完全列を得る．
\end{proof}



\begin{proposition}[非自由MCM二型の行列核は点曲線の二pencilを回収する]
\label{prop:qa-point-mcm-matrix-pencils}
前定理の仮定の下で，$f_1=-\ell_1\ell_2$，
$f_2=-m_1m_2$と因数分解し，
$\ell_1,\ell_2$および$m_1,m_2$を
それぞれ独立な一次形式とする．
中央二重平方環 $B_f$ で
\[
 M_c=\begin{pmatrix}c&\ell_1\\\ell_2&-c\end{pmatrix},
 \qquad
 M_d=\begin{pmatrix}d&m_1\\m_2&-d\end{pmatrix}
 \tag{MCM-matrices}\label{eq:qa-point-mcm-matrices}
\]
と置くと，$M_c^2=M_d^2=0$であり，
それぞれは元のnoncubic normal helixに使う
二つの線形matrix factorizationの表示となる．
$E_f$上では両行列が全点で階数一となり，
その右核の直線から定まる射
\[
 \kappa_c:E_f\longrightarrow\PP^1,\qquad
 \kappa_d:E_f\longrightarrow\PP^1
\]
は，それぞれ前定理の
$p_c:E_f\to Q_c$，$p_d:E_f\to Q_d$
と conic の同型 $Q_c,Q_d\simeq\PP^1$
との合成に一致する（終域の射影変換を除く）．
従って
\[
 \kappa_c^*\OO_{\PP^1}(1)\simeq F_c,\qquad
 \kappa_d^*\OO_{\PP^1}(1)\simeq F_d .
 \tag{MCM-pencils}\label{eq:qa-point-mcm-pencil-classes}
\]
これは\emph{matrix factorizationの点上のkernel}と
幾何的な二pencilの一致であり，
ASのfinite directed End代数全体の同型や
$\pi B_f$を保つpointed心の同型を
まだ結論しない．
\end{proposition}
\begin{proof}
$u,v$ は $B_f$ の中心にあり，
$c^2=f_1=-\ell_1\ell_2$，
$d^2=f_2=-m_1m_2$だから
行列の平方はそれぞれ零である．
これは
\cite[\texttt{thm:ore-anticommuting-central}]{QuadricIntrinsic2026}
の二つの非自由MCM型に用いる
線形行列分解の表示である．
$E_f$の点において$M_c$が零行列となるなら
$c=\ell_1=\ell_2=0$が必要だが，
$\ell_1,\ell_2$は$\PP^1_{[u:v]}$上で
共通零点を持たない．
従って行列の階数は一である．
同じ議論が$M_d$にも適用できる．

$M_c$の成分は$u,v,c$だけに依存するので，
$\kappa_c$は conic $Q_c$ を経由する．
$\ell_1,\ell_2$を $s,t$ とする座標で
$Q_c$の式は $c^2+st=0$である．
その標準パラメータ表示
\[
 [s:t:c]=[a^2:-b^2:ab]
\]
の下で
\[
 M_c=
 \begin{pmatrix}a\\-b\end{pmatrix}
 \begin{pmatrix}b&a\end{pmatrix}
\]
となるから，右核は直線$[a:-b]$である．
これは conicのパラメータ$[a:b]$への
射影線形自己同型であり，
$\kappa_c$は$Q_c\simeq\PP^1$を通る次数二射となる．
よって$\kappa_c^*\OO(1)=F_c$を得る．
$m_1,m_2,d$を用いた同じ計算で
$F_d$の主張も従う．
\end{proof}

\begin{proposition}[MCMの隣接積と二pencilのVeronese積の一致]
\label{prop:qa-point-mcm-composition-veronese}
前命題の仮定の下で，中央二重平方環の
一次部分を
\[
 V=(B_f)_1=\langle u,v,c,d\rangle_k
     \xrightarrow{\ \sim\ }
     H^0(E_f,\OO_{E_f}(1))
 \tag{Degree-one}\label{eq:qa-point-degree-one-restriction}
\]
と同定する．
二つの非自由MCM型 $X_c,X_d$ に関する
noncubic normal helix の隣接合成
\[
 \begin{aligned}
 \mu_c:\;&
 \Hom(\pi X_c,\pi B_f(1))
 \otimes\Hom(\pi B_f,\pi X_c)
 \longrightarrow\Hom(\pi B_f,\pi B_f(1)),\\
 \mu_d:\;&
 \Hom(\pi X_d,\pi B_f(1))
 \otimes\Hom(\pi B_f,\pi X_d)
 \longrightarrow\Hom(\pi B_f,\pi B_f(1))
 \end{aligned}
\]
の像をそれぞれ $U_c,U_d\subset V$ とする．
一方，$E_f$ 上の二つの次数二pencilに対する
通常の対称積
\[
 m_c:\Sym^2H^0(E_f,F_c)\to H^0(E_f,F_c^2),
 \qquad
 m_d:\Sym^2H^0(E_f,F_d)\to H^0(E_f,F_d^2)
\]
の像を前式の$V$へ写す．すると
\[
 \boxed{
 \im\mu_c=\im m_c=U_c=\langle u,v,c\rangle,\qquad
 \im\mu_d=\im m_d=U_d=\langle u,v,d\rangle.}
 \tag{Two-rank3}\label{eq:qa-point-two-rank-three-subspaces}
\]
特に
\[
 \dim U_c=\dim U_d=3,\qquad
 U_c\cap U_d=\langle u,v\rangle,\qquad U_c+U_d=V.
\]
よってnoncubic normal helixの二つのMCM型に伴う
\emph{最小十本矢の隣接合成像}は，点曲線$E_f$の
二つのconic商pencilのVeronese像と一致する．
ただし，この命題は二つの像部分空間と
MCMのkernel写像を同定するもので，
四頂点のfinite directed End代数の
\emph{すべての積}の同型を意味しない．
\end{proposition}
\begin{proof}
$E_f\subset\PP^3$は二quadricの完全交叉で，
一次の斉次関係はない．
よって$V=H^0(E_f,\OO_{E_f}(1))$となる．
定理
~\cite[\texttt{thm:normal-helix} および
 \texttt{cor:normal-cubic-defect}]{QuadricIntrinsic2026}
における非自由MCMの隣接合成は，
線形matrix factorizationの
$2\times2$成分を一次切断とみなしたものだから，
式~\eqref{eq:qa-point-mcm-matrices}より
\[
 \im\mu_c=\langle c,\ell_1,\ell_2\rangle
          =\langle u,v,c\rangle,\qquad
 \im\mu_d=\langle d,m_1,m_2\rangle
          =\langle u,v,d\rangle
\]
である．ここで各一次形式対は独立である．

一方，前命題から
$F_c=p_c^*N_c$，$N_c\simeq\OO_{\PP^1}(1)$であり，
$H^0(E_f,F_c)=H^0(Q_c,N_c)$である．
実際，$p_{c*}\OO_{E_f}$の不変部分は$\OO_{Q_c}$で，
反不変部分の負次数の線束には切断がない．
従って$m_c$は
\[
 \Sym^2H^0(Q_c,N_c)\longrightarrow
 H^0(Q_c,N_c^2)=H^0(Q_c,\OO_{Q_c}(1))
\]
の通常のVeronese積の引戻しである．
この積は三次元切断空間への同型であり，
$Q_c\subset\PP^2_{[u:v:c]}$の
超平面切断の像は
$\langle u,v,c\rangle$である．
同様に$m_d$の像は$\langle u,v,d\rangle$となる．
二つの部分空間の共通部分と和は明らかである．
\end{proof}

\begin{lemma}[noncubic 点曲線への制限と四頂点Endの同型]
\label{lem:qa-pointed-abelianization-four-end}
標数零の代数閉体上、$f_1,f_2$がそれぞれ相異なる二根を持ち、
互いに根を共有しないとする。
$B=B_f$、$R=B/(cd-dc)$、$E=\operatorname{Proj}R=E_f$とする。
$f_1=-\ell_1\ell_2$, $f_2=-m_1m_2$に対する
二つの右$B$-加群$X_c=\operatorname{coker}M_c$,
$X_d=\operatorname{coker}M_d$を前命題のように取る。
$F=F_c$, $G=F_d$と書けば、次の自然な制限は
頂点と全合成を保つ有限次元代数の同型である：
\begin{equation}
\End_{\rm dir}^{\qgr B}
 (\pi B,\pi X_c,\pi B(1),\pi X_d(1))
 \ \xrightarrow{\ \sim\ }\
\End_{\rm dir}^{E}
 (\OO_E,F,F^2,F^2G).
\label{eq:qa-pointed-four-end-abelianization}
\end{equation}
ここで$F^2=G^2=\OO_E(1)$であり、
左辺では$B$の通常シフト$(1)$を記すが、
これ自体を適合Serre平方根とは主張しない。
\end{lemma}
\begin{proof}
$R_c=k[u,v,c]/(c^2-f_1)$とおく。
$R=R_c\oplus dR_c(-1)$は$R_c$上自由である。
滑らかなconic $Q_c=\operatorname{Proj}R_c\simeq\PP^1$
上で$M_c$は$s=\ell_1,t=\ell_2$により
\[
c^2+st=0,\qquad
[s:t:c]=[a^2:-b^2:ab],\qquad
M_c=\binom{a}{-b}(b\ \ a)
\]
となる。従って$\operatorname{coker}M_c$を
$R_c$へ制限した次数付き加群は
$\bigoplus_{n\ge0}H^0(Q_c,\OO_{\PP^1}(2n+1))$である。
$R_c\to R$への自由な拡大と
$p_{c*}\OO_E\simeq\OO_{\PP^1}\oplus\OO_{\PP^1}(-2)$から
\[
X_c\otimes_B R
 \simeq\bigoplus_{n\ge0}H^0(E,F\OO_E(n)),
\qquad
X_d\otimes_B R
 \simeq\bigoplus_{n\ge0}H^0(E,G\OO_E(n)).
\tag{Spinor--sections}\label{eq:qa-pointed-spinor-sections}
\]
両式は単にsheafificationの同型ではなく、
非負次数のsection moduleの同型である。
特に右辺は有限長torsionを持たない。
$B\to R$の核は次数二の$cd-dc$で生成されるので、
$B_0,B_1$ならびに$(X_c)_0,(X_c)_1,
(X_d)_0,(X_d)_1$は、それぞれ$R$への制限で不変である。

Aの適合表示をここで明確にする。
分類論文Aの
\cite[\texttt{thm:ore-anticommuting-central}]{QuadricIntrinsic2026}
で構成するZhang自己同型$\sigma$は、
線形座標変換を伴って$c,d$を交換し、
$B^\sigma$の通常シフトを適合根とする。
右Zhang同値でそのnormal helixの四連続対象を$B$側へ
戻すと、自由加群のshiftの作用を$\sigma$で輸送し、
二つのspinor型も交換するので、
添字を選んで
$(\pi B,\pi X_c,\pi B(1),\pi X_d(1))$
を得る。この輸送は次数付き加群の単なる同一視ではなく、
右作用の$\sigma$-twistを含む。
従って
\cite[\texttt{thm:normal-helix},
\texttt{cor:normal-cubic-defect}]{QuadricIntrinsic2026}
のHom計算と、深さ三によるgraded/tails比較が適用できる。
四頂点の上三角Hom次元は
\[
\begin{pmatrix}
1&2&4&6\\0&1&2&4\\0&0&1&2\\0&0&0&1
\end{pmatrix}.
\tag{Four-End-dims}\label{eq:qa-pointed-four-end-dimensions}
\]
右辺の曲線上の線束も次数$0,2,4,6$であり、
$E$が滑らかな種数一曲線なので同じ次元を持つ。

左辺の各Homは生成元の像で決まり、
$B_0,B_1,(X_d)_0,(X_d)_1$のいずれかに値を持つ。
特に$\Hom(X_c,X_d(1))$の射は
$X_c$の次数零の二生成元の$(X_d)_1$内の像で決まる。
上記の低次数の制限単射性と
\eqref{eq:qa-pointed-spinor-sections}により、
どのHomも右辺のsheaf Homへ単射となる。
両辺の次元が一致するため全て同型である。
$-\otimes_B R$は右完全で、個々の加群射の合成を保つため、
上の六つの同型は同時に全積を保つ。

最初の二段階合成はさらに直接確認できる。
$X_c$の次数零生成元を$e_1,e_2$とし、
$\Hom_B(X_c,B(1))$の基底を
$(c,\ell_1),(-\ell_2,c)$と取ると、
合成の$2\times2$行列は
\[
\begin{pmatrix}
c&\ell_1\\-\ell_2&c
\end{pmatrix}.
\]
従って唯一の二次関係は二つの対角道の差であり、
像は$\langle c,\ell_1,\ell_2\rangle$である。
上のconic表示で行列は
$\binom{a}{b}(b\ \ a)$となり、
$H^0(F)\otimes H^0(F)\to H^0(F^2)$の
反対称核と一致する。この計算は像の一致だけでなく、
二つの二次元Homの\emph{双線形積そのもの}を同定する。
\end{proof}

\begin{theorem}[滑らかなnoncubicの位相を合わせたpointed認識]
\label{thm:qa-point-smooth-noncubic-pointed}
\label{prop:qa-point-mcm-pointed-finite-end-reduction}
前補題の仮定の下で、
square avatar $(E,F,G,P)$を
$F=F_c$, $G=F_d$, $P=GF^{-1}\ne\OO_E$で定める。
$P^2=\OO_E$かつ$G=FP$である。
このsquare型AS$\mathbb Z$-代数を論文Cの指定mutationで
最小十本矢型$A_{\rm av}$に移し、
mutation後のhelixを$(H_i)_{i\in\mathbb Z}$とする。
このとき
\begin{equation}
 \End_{\rm dir}^{\qgr B_f}(E_0,E_1,E_2,E_3)
 \ \simeq\
 \End_{\rm dir}^{\qgr A_{\rm av}}(H_1,H_2,H_3,H_4),
\qquad E_0=\pi B_f,
\label{eq:qa-pointed-shifted-four-end}
\end{equation}
という頂点と全合成を保つ同型が存在する。
論文Aのnormal-helixと本稿のAS10 Serre完成・Noether性・AZ再構成を
適用すると、
\begin{equation}
(\qgr B_f,\pi B_f)\simeq
(\qgr A_{\rm av},\,\pi e_1A_{\rm av})
\tag{Smooth-pointed}\label{eq:qa-pointed-smooth-equivalence}
\end{equation}
というpointed exact同値を得る。
ここで構造対象の像はmutation\emph{後の一番頂点}であり、
ゼロ番頂点ではない。
\end{theorem}
\begin{proof}
square側の曲線上の初期四項は$(P,F,G,FG)$である。
補題~\ref{lem:qa-as10-all-index-mutual-inverse}で扱うevaluation mutationでは
\[
 H_0=K_0,\quad H_1=P,\quad H_2=G,\quad H_3=FG,
\quad
 K_0=\ker\bigl(H^0(FP^{-1})\otimes P\to F\bigr)
      =P^2F^{-1}.
\]
次のsquare項は
系~\ref{cor:qa-core-square-complexes}の式から
$L_4=F^2G^2P,\ L_5=F^3G^2P^2$である。
ゆえに一周期後のevaluation kernelは
\[
 H_4=K_1=L_4^2L_5^{-1}=FG^2.
\]
$P^2=\OO_E$, $G=FP$を代入し、
四項$(H_1,H_2,H_3,H_4)=(P,FP,F^2P,F^3)$を
共通線束$P^{-1}$でtensorすると
\[
(\OO_E,F,F^2,F^3P^{-1})
   =(\OO_E,F,F^2,F^2G).
\]
共通線束のtensorは曲線上の全Homとその積を保つ。
前補題と本稿の全整数mutationのEnd同定を合わせれば
\eqref{eq:qa-pointed-shifted-four-end}を得る。
これは線束上の計算であり、任意のnc heartに
線束tensor自己同値を仮定するものではない。

Aのnormal helixはfullでample、
Cのmutation後のhelixもfull very strongで、
Noether性とAZ amplenessを持つ。
有限Endの同型はtilting Morita同値を与え、
Serre関手と自然に可換する。
従って四項をSerre完成して得られる全整数の
前向きHomと積は一致する。
補題~\ref{lem:qa-as-serre-unique-recovery}と
補題~\ref{lem:qa-as10-all-index-mutual-inverse}および
\cite[\texttt{thm:normal-helix}]{QuadricIntrinsic2026}
の再構成を適用して
\eqref{eq:qa-pointed-smooth-equivalence}を得る。
\end{proof}

\begin{remark}[位相の取り違えと相対境界]
\label{rem:qa-pointed-phase-obstruction}
同じnoncubic normal四項をmutation後の
$(H_0,H_1,H_2,H_3)$に\emph{頂点固定で}
同定することはできない。
前者の$0\to1\to2$合成の核は一次元だが、
後者は最小十本矢型の定義により
$0\to2$に最小関係を持たず、核は零である。
従って一添字shiftは単なる記法でなく必要条件である。
なお定理は四根相異なるsmooth noncubicの
\emph{体上のpointed同値}を述べる。
一根共有・二根対一致における同じ式の相対化、
非被約底上のgroupoid、$C_q$側の通常Picardへの延長、
およびpotential全体の相対忘却は主張しない。
\end{remark}

\begin{corollary}[非零二捩れ付き楕円曲線上の双対二重商]
\label{cor:qa-point-spectral-atkin-lehner}
滑らかな楕円曲線と非零二捩れの組
$(T,L)$ の同型類に対し
\[
 \mathcal W_2(T,L)=
 \bigl(T/\langle L\rangle,\ L^\vee\bigr)
 \tag{Dual-wall}\label{eq:qa-point-spectral-dual-wall}
\]
と置く．ここで $L$ を原点選択によって
位数二の点とみなし，$L^\vee$は
商 isogeny の双対核を指定する非零二捩れである．
この対応は同型類上の involution：
$\mathcal W_2^2=\mathrm{id}$ を定める．
前定理の根対から生じる二つの組は
\[
 \mathcal W_2(C_q,M_C)=(E_f,P_E),\qquad
 \mathcal W_2(E_f,P_E)=(C_q,M_C)
\]
を満たす．
これは楕円曲線と二捩れ\emph{だけ}の対応であり，
$F_d,F_c$と$p^*\OO_{\PP^1}(1)$という
異なる degree-two pencil まで
自動的に同一視するものではない．
\end{corollary}
\begin{proof}
有限 \'{e}tale 二重 isogeny
$T\to T/\langle L\rangle$に対し，
双対 isogeny の核は位数二である．
さらに双対を取れば，楕円曲線の主偏極による
二重双対の自然同型から元の isogeny に戻る．
従って商と双対核の操作は二回で元に戻る．
前定理で$r:E_f\to C_q$の核は$P_E$であり，
$r^\vee:C_q\to E_f$の核は$M_C$と示したので，
二つの表示が得られる．
\end{proof}

\begin{example}[二曲線は一般に同型でない]
\label{ex:qa-point-spectral-j-distinction}
アフィン座標$x=u/v$で
$f_1=x(x-1)$，$f_2=(x-2)(x-3)$とする．
$C_q$の射影模型は
\[
 y^2=x(x-1)(x-2)(x-3).
\]
一方 $E_f$の関数体で
$t=d/(x-3)$と置けば
\[
 x=\frac{3t^2-2}{t^2-1},\qquad
 z=(t^2-1)c,\qquad
 z^2=(3t^2-2)(2t^2-1)=6t^4-7t^2+2.
\]
四次式 $az^4+bz^3+cz^2+dz+e$ の
不変式を $I=12ae-3bd+c^2$，
$J=72ace+9bcd-27ad^2-27b^2e-2c^3$
とすると
$j=1728\cdot4I^3/(4I^3-J^2)$である．
$C_q$の四次式には $(I,J)=(13,-70)$，
$E_f$の上の四次式には $(I,J)=(193,-5362)$を得るので，
\[
 \boxed{j(C_q)=\frac{35152}{9},\qquad
        j(E_f)=\frac{28756228}{3}.}
 \tag{Distinct-j}\label{eq:qa-point-spectral-distinct-j}
\]
従って両曲線は同型でない．
この有理数の計算は
\texttt{computations/verify\_rootpair\_resolvent\_spectral.py}
にも厳密算術の検算を置く．
\end{example}

\begin{remark}[pointed 認識への帰結と退化境界]
\label{rem:qa-point-spectral-pointed-warning}
定理~\ref{thm:qa-noncubic-rootpair-all-geometric-square-avatar-bijection}
で使った square avatar の曲線は \emph{$C_q$} である．
一方，noncubic 中央二重平方環そのものの
可換点曲線は \emph{$E_f$} であり，
第二項は $E_f$ 上にも自然な
二捩れ double-wall avatar があることを示す．
式~\eqref{eq:qa-point-spectral-distinct-j}により，
二つの構成を同一曲線とみなして
avatar由来AS代数と$\qgr B_f$の
pointed同値を結論することはできない．
\emph{どちらの曲線が特定のnormal helixの
AS avatarとして抽出されるか}，
およびそのexceptional foundationの対応は
まだ直接同定していない．これは次の認識問題である．

$f_1,f_2$ が一根を共有すると，
$\tau$には固定点が生じ，
$r$は有限 \'{e}tale 二重isogenyではなくなる．
その際に正しい相対対象は
式~\eqref{eq:qa-rootpair-spectral-module}の
torsion-free $M_C$とその二次積であり，
双対isogenyの短完全列を特異fiberへ
そのまま延長してはならない．
\end{remark}

% !TeX root = ../nc-quadric-avatar.tex
% 2026-11-22 / root-factor monodromy and generic Brauer obstruction

\section{根の内部交換と相対avatarのBrauer障害}
\label{sec:qa-root-descent-brauer}

本節では有限directed End代数の相対比較と、その上の線束を
無標識の二次形式モジュライへ降下させる問題を区別する。
根対を\emph{内部で}交換する操作には厳密な行列表示があるが、
因子の逆数倍も忘れると一般conicが非分裂になる。
とくに$E_f$と$C_q$の相違とは独立な障害が現れる。

\begin{proposition}[二つの根対の内部交換による厳密な有限End降下]
\label{prop:qa-root-swaps-linearization}
命題~\ref{prop:qa-point-relative-split-root-end}の
順序付き一次形式空間$S$上で
$W=(\mathbb Z/2)^2$を、$\ell_1,\ell_2$の交換と
$m_1,m_2$の交換で作用させる。
$W$は$S$に自由に作用する。
$\iota\in k$を$\iota^2=-1$となるように一つ選ぶと、
相対spinor加群$\mathcal X_c,\mathcal X_d$および
四項$(\mathcal B,\mathcal X_c,\mathcal B(1),
\mathcal X_d(1))$には、$W$の\emph{厳密な}
線形化を与えられる。
従って相対有限directed Endの同型
\[
 \mathfrak B_S\xrightarrow{\sim}\mathfrak G_S
\]
は$W$-同変にでき、
有限\'{e}tale商$S/W$上の有限局所自由代数の同型へ降下する。
降下後の比較も$S/W$の任意の基底変更と可換する。
根対の内部交換と、二つの二次形式の\emph{相互交換}は
異なる作用であり、後者についてはここでは述べない。
\end{proposition}
\begin{proof}
$\ell_1=s,\ell_2=t$と記し、右加群の
二項行列表示を
\[
 M(s,t)=\begin{pmatrix}c&s\\t&-c\end{pmatrix},
 \quad
 J=\begin{pmatrix}0&1\\-1&0\end{pmatrix},
 \quad J^2=-I_2
\]
とする。行列の各成分$s,t,c$は相互に可換するので
直接の乗法から
\[
 M(t,s)=JM(s,t)J,\qquad
 (\iota J)M(s,t)=M(t,s)(-\iota J).
 \tag{Swap-matrix}\label{eq:qa-root-swap-matrix}
\]
を得る。
ここで$(\iota J)^2=(-\iota J)^2=I_2$である。
従って二項表示の自由加群に対し、
余核側で$\iota J$、関係側で$-\iota J$を使えば、
二回の根交換が恒等である鎖同型が定まる。
その余核への作用が$\mathcal X_c$の厳密な
$\mathbb Z/2$-線形化である。
$\ell$の交換は$\mathcal X_d$と$\mathcal B$には恒等で作用する。
同様に$m$の交換には$d,m_1,m_2$を使い、
$\mathcal X_d$に同じ行列で線形化する。
二つの作用は異なる加群成分にのみ非自明に働くので可換し、
四項と全Hom・合成に$W$-作用を与える。

$f_1=-\ell_1\ell_2$、$f_2=-m_1m_2$は
それぞれの内部交換で不変だから、
$\mathcal B,\mathcal R=\mathcal B/(cd-dc)$には
基底変換と両立する自明な係数の輸送がある。
$\mathcal X_i\otimes_{\mathcal B}\mathcal R$の
sheafificationに上の線形化を輸送し、
その線束を$\mathcal F,\mathcal G$と同定する。
四項の後二項
$L_2\simeq\mathcal F^2$、$L_3\simeq\mathcal F^2\mathcal G$
には、それぞれ$\mathcal B(1)$と$\mathcal X_d(1)$の
制限から線形化を輸送する。
ここで$\mathcal F^2\simeq\OO_{\mathcal E}(1)$の
選択された同型が、$\mathcal F$の線形化の
tensor平方と自動的に両立するとは仮定しない。
必要なのは四つの\emph{対象ごと}に整合する線形化であり、
それにより命題~\ref{prop:qa-point-relative-split-root-end}の
有限End射は自然性から$W$-同変である。
二つの一次形式が独立なので、それぞれの交換が固定する
$S$の幾何点は存在しない。
$W$は有限定数群で標数零だから
$S\to S/W$は有限\'{e}tale torsorである。
有限局所自由加群・その積と代数同型の
faithfully flat descentが適用でき、主張が従う。

$\iota$を$-\iota$へ変えると、各交換の
線形化が符号指標だけ変わる。
従ってこの命題は線形化の\emph{存在}を述べ、
選択に依らない正準線形化とは主張しない。
\end{proof}

\begin{proposition}[普遍二次形式のconicは一般に非分裂]
\label{prop:qa-universal-conic-brauer-obstruction}
$k$を標数零の代数閉体とし、
\[
 U=\operatorname{Spec}k[a,b,d,(b^2-4ad)^{-1}],
 \quad K=k(a,b,d),\quad \Delta=b^2-4ad
\]
と置く。$K$上の滑らかなconic
\begin{equation}
 Q:\quad c^2=au^2+buv+dv^2\subset\PP^2_K
 \label{eq:qa-generic-conic}
\end{equation}
は$K$-有理点を持たない。
そのSeveri--Brauer類は
\[
 \alpha(Q)=(a,\Delta)\ne0\quad\text{in}\quad
 \operatorname{Br}(K)[2]
 \tag{Brauer-class}\label{eq:qa-conic-brauer-class}
\]
であり、$\Delta$を一様化元とする離散付値の
留数は$[a]\ne1$である。
従って$Q$には幾何次数一の線束が存在せず、
$\bar K$上の同型$Q_{\bar K}\simeq\PP^1$で定まる
次数一束$\OO_{\PP^1}(1)$は$K$上へ降下しない。
ここで射影平面から制限した$\OO_Q(1)$は
幾何次数\emph{二}である。

さらに二つの独立な一般二次形式の係数体を$K_{12}$とし、
$Q_c,Q_d$のファイバー積である滑らかな
$E=Q_c\times_{\PP^1}Q_d$を考える。
幾何学的な次数二pencil
$F_c=p_c^*\OO_{(Q_c)_{\bar K_{12}}\simeq\PP^1}(1)$と
$F_d=p_d^*\OO_{(Q_d)_{\bar K_{12}}\simeq\PP^1}(1)$は
Galois不変なPicard\emph{類}ではあるが、
一般には$E$上の実際の線束として降下しない。
そのHochschild--Serre障害はそれぞれ
\begin{equation}
 \partial_E[F_c]=\alpha(Q_c),\qquad
 \partial_E[F_d]=\alpha(Q_d),\qquad
 \partial_E[F_dF_c^{-1}]=\alpha(Q_d)-\alpha(Q_c).
 \label{eq:qa-point-pencils-brauer-obstruction}
\end{equation}
であり、三つとも$K_{12}$のBrauer群で非零である。
\end{proposition}
\begin{proof}
$K=k(a,b,\Delta)$と見なせる。
$a\ne0$の開上で
$z=u+(b/2a)v$と置くと、
\eqref{eq:qa-generic-conic}は
\[
 c^2-a z^2=-\frac{\Delta}{4a}v^2
 \tag{Diagonal-conic}\label{eq:qa-diagonal-conic}
\]
となる。右辺は$-\Delta/(4a)$を係数とする
ノルム方程式で、対応するquaternion類は
$(a,-\Delta/(4a))=(a,\Delta)$である
（$-1$と$4$は$k$で平方）。
$\Delta$の付値に対し$a$は単元で、
剰余体$k(a,b)$内の$a$は平方ではない。
従ってquaternion類の留数は非自明な$[a]$となる。

非分裂性はBrauer類の一般論を用いずに次のようにも示せる。
$\Delta$-進DVRで仮に\eqref{eq:qa-diagonal-conic}に
射影有理点があるとする。座標$(c,z,v)$をprimitiveに規格化し、
$\Delta$で還元すると
$\bar c^{\,2}=a\bar z^{\,2}$となる。
$a$が剰余体で非平方なので$\bar c=\bar z=0$である。
従って$\bar v\ne0$だが、
左辺の付値は二以上、右辺の付値は一となり矛盾する。
滑らかな種数零曲線上の幾何次数一線束が
$K$上に存在すればRiemann--Rochにより
その非零切断の零点から次数一の$K$-有効因子、
すなわち$K$-有理点が得られるので、存在しない。

二つの独立係数のときも、
第一の判別式$\Delta_c$で取った留数は
$\alpha(Q_c)$に対し$[a_c]\ne1$、
$\alpha(Q_d)$に対し零である。
同様に$\Delta_d$では役割が交換する。
よって二つのquaternion類はそれぞれ非零で、
和または差も非零である。

最後の主張は、滑らかな射影幾何的整曲線に対する
Galois descentのPicard--Brauer境界
\[
 0\longrightarrow\Pic(E)
 \longrightarrow\Pic(E_{\bar K_{12}})^{G_{K_{12}}}
 \xrightarrow{\ \partial_E\ }\operatorname{Br}(K_{12})
\]
とそのpullbackに関する自然性による。
滑らかなconicの幾何次数一Picard類の障害は
そのSeveri--Brauer類$\alpha(Q_i)$であり、
$p_i:E\to Q_i$で引き戻した類の障害も同じである。
Picard群上の差を取れば第三の式を得る。
これは$\Pic^0$の\emph{類}に相当する
Galois不変点があることと、その類の
実際の線束による代表があることの違いである。
\end{proof}

\begin{corollary}[二つの分岐根を選ぶ次数四の\'{e}tale分裂被覆]
\label{cor:qa-root-choice-etale-splitting}
$U_2$を、標数零の代数閉体上の二つの
非退化二元二次形式$(f_1,f_2)$の全パラメータ空間とする。
二つの根対の相互共有は許す。
$Z_i\subset U_2\times\PP^1$を$f_i=0$の
相対零点とし、
\[
 \widetilde U_2=Z_1\times_{U_2}Z_2
       \longrightarrow U_2
 \tag{Root-choice-cover}\label{eq:qa-root-choice-four-cover}
\]
と置く。この射は、共有根および一致根対の境界を含めて
有限\'{e}tale次数四である。
$\widetilde U_2$上の二つのconic
$Q_c,Q_d$には、それぞれ自然な分岐点section
$\sigma_c,\sigma_d$が存在し、
\[
 N_c=\OO_{Q_c}(\sigma_c),\qquad
 N_d=\OO_{Q_d}(\sigma_d)
\]
は相対次数一の実際の線束である。
従って$E=Q_c\times_{\PP^1}Q_d$上には
$F=p_c^*N_c,\ G=p_d^*N_d$という
大域的な二つの次数二pencilが存在し、
各幾何fiber上で
\[
 F^2=G^2=\OO_E(1),\qquad
 P=GF^{-1}\in\Pic^0(E)[2]\setminus\{0\}
\]
となる。これらの等式はまず各幾何fiberのPicard類の等式であり、
基底上の線束捩れを固定する正準の
平方同型を主張しない。
この被覆は一般conicのBrauer障害を明示的に消すが、
相対AS heartのpointed同値や
$\widetilde U_2\to U_2$についての
線束の降下を自動的には与えない。
\end{corollary}
\begin{proof}
各$f_i$の判別式が可逆なので
$Z_i\to U_2$は有限\'{e}tale次数二である。
二つを独立に基底変更しても有限\'{e}tale性が保たれ、
fiberの次数は$2\cdot2=4$となる。
根の共有は$f_1,f_2$の相互関係であり、
個々の$Z_i\to U_2$が\'{e}taleであることを妨げない。

選ばれた根$p_i$では$f_i(p_i)=0$なので
$Q_i:c_i^2=f_i$の上に$(p_i,c_i=0)$という
sectionが定まる。$Q_i$は相対的に滑らかなconicで、
このsectionは有効Cartier因子を与える。
従って$\OO_{Q_i}(\sigma_i)$は相対次数一。
$E$への引戻しは相対次数二である。
各幾何fiberでは$Q_i\simeq\PP^1$なので
$N_i^{\otimes2}$はconic超平面束になり、
補題~\ref{lem:qa-pointed-all-rootpair-geometry}から
$P$が非自明な多重次数零の二捩れとなる。
各式の相対的な捩れは、線束の基底からの引戻しを
標識から除く既存のavatar規約に従う。
\end{proof}

\begin{remark}[相対有限Endと無標識avatarの正確な境界]
\label{rem:qa-brauer-scope}
命題~\ref{prop:qa-root-swaps-linearization}は
\emph{根対内の有限な順序交換}を消す降下である。
二つの一次形式を
$(\ell_1,\ell_2)\mapsto(t\ell_1,t^{-1}\ell_2)$
と変える$\Gm$作用も同時に割った
二次形式の係数空間への線束の降下ではない。
実際、
\[
 D(t)=\operatorname{diag}(t,1),\qquad
 M(t\ell_1,t^{-1}\ell_2)=D(t)M(\ell_1,\ell_2)D(t)^{-1}
\]
だが、根交換の行列$J$について
\[
 J D(t)J^{-1}
   =\operatorname{diag}(1,t)
   =tD(t^{-1}) .
\]
右辺の余分なスカラーは、
線束の表示が単なる$PGL_2$作用から
自動的に$GL_2$作用へ持ち上がらないことの
有限行列的な兆候である。
実際に線束が降下しないことは
命題~\ref{prop:qa-universal-conic-brauer-obstruction}で
基礎体のBrauer類から証明した。
この障害から、\emph{別の}Morita表示、twisted sheaf、
または内在的$\mathcal H_2$の相対stack自体が
存在しないとは結論しない。
従って次の課題はBrauer類を殺す適切な
gerbe/相対AS表示の構成と、二つの根対の相互交換の
相対mutation・pointingの降下である。
\end{remark}

% !TeX root = ../nc-quadric-avatar.tex
% 2026-11-22: relative Azumaya vertex classes and finite directed Morita model
\section{Brauer重みを持つ四頂点Morita表示}
\label{sec:qa-twisted-morita-weights}
前節のgerbe上の幾何avatarの四項Endは、慣性の重みを
記録してAzumaya頂点・双加群矢により係数空間へ降下する。
ここでは有限EndのMorita表示のみを扱い、
相対AS完成およびpointed \qgr の同値は主張しない。
gerbeの重み分解およびtwisted sheafのMorita理論は
\cite[Theorems~4.7 and 5.4]{BerghSchnurer2021}
および\cite[Theorem~6.2 and Remark~6.3]{BerghSchnurer2021}
の一般論を使う。

\begin{lemma}[conic pencilのdeterminant補正]
\label{lem:qa-morita-determinant-normalization}
前節の$U_2$と$\mathscr G=\mathscr G_c\times_{U_2}\mathscr G_d$を取る。
$q_i:Q_i\to\mathscr G$を普遍smooth conic、
$N_i$を普遍相対幾何次数一束とし、
$V_i=q_{i*}N_i$（$i=c,d$）と置く。
$V_i$は階数二、全基底変更可換で、
一次元商の規約で$Q_i\simeq\PP_{\mathscr G}(V_i)$となる。
conic超平面束を$\OO_{Q_i}(1)$とすれば
\begin{equation}
 N_i^{\otimes2}\simeq
 \OO_{Q_i}(1)\otimes q_i^*\det V_i .
 \label{eq:qa-morita-determinant-square}
\end{equation}
$E=Q_c\times_{\PP^1}Q_d$、
$p_i:E\to Q_i$、$\pi:E\to\mathscr G$、
$F=p_c^*N_c$、$G=p_d^*N_d$、
$H=\OO_E(1)$と書くと
\begin{equation}
 F^2\simeq H\pi^*\det V_c,\quad
 G^2\simeq H\pi^*\det V_d,\quad
 (GF^{-1})^2\simeq\pi^*(\det V_d(\det V_c)^{-1}).
 \label{eq:qa-morita-fg-square}
\end{equation}
従って$GF^{-1}$が相対Picardの二捩れsectionであることと、
gerbe上でその二乗線束が実際に自明であることは異なる。
\end{lemma}
\begin{proof}
各幾何conicは$\PP^1$であり、$N_i$の相対次数は一だから
$V_i$は階数二で$R^1q_{i*}N_i=0$、
相対評価は全射である。評価が定める
$Q_i\to\PP(V_i)$は各fiberで次数一の完全線形系なので同型。
商規約の相対Euler列は
$\omega_{Q_i/\mathscr G}\simeq N_i^{-2}\otimes q_i^*\det V_i$
を与える。一方、二次Cartier超曲面
$Q_i\subset\PP^2_{U_2}$の随伴から
$\omega_{Q_i/\mathscr G}\simeq\OO_{Q_i}(-1)$である。
これで最初の式が得られる。
両conicの超平面束は、共通の$\PP^1_{[u:v]}$の
$\OO(1)$の引戻しだから、$E$上で同じ$H$になる。
二式の引戻しと商で残りの式が従う。
$\det V_c,\det V_d$の慣性重みは$(2,0),(0,2)$であり、
最後の右辺の自明化を仮定してはならない。
\end{proof}

\begin{theorem}[四頂点のAzumaya双加群による有限Morita表示]
\label{thm:qa-twisted-four-vertex-morita}
前補題の記号で、gerbe上の四線束を
\begin{equation}
 (L_0,L_1,L_2,L_3)=(\OO_E,F,H,HG)
 \label{eq:qa-morita-normalized-four}
\end{equation}
と正規化する。
$T_{ij}=\pi_*(L_i^{-1}L_j)$（$i\le j$）は
有限局所自由で、任意の基底変更と可換する。
非対角成分の階数と慣性$\Gm^2$の指標は
\begin{equation}
\begin{array}{c|rrrrrr}
(i,j)&01&02&03&12&13&23\\\hline
\operatorname{rk}T_{ij}&2&4&6&2&4&2\\
\operatorname{wt}T_{ij}&(1,0)&(0,0)&(0,1)&(-1,0)&(-1,1)&(0,1)
\end{array}
\label{eq:qa-morita-hom-weights}
\end{equation}
であり、$T_{ii}=\OO_{\mathscr G}$である。
四項の重みは
$w_0=(0,0)$、$w_1=(1,0)$、$w_2=(0,0)$、$w_3=(0,1)$。
$W_0=W_2=\OO_{\mathscr G}$、
$W_1=V_c$、$W_3=V_d$とする。
重み零の降下を$\operatorname{Desc}$で記し、
\[
 A_i=\operatorname{Desc}\End(W_i),\qquad
 {\mathsf M}_{ij}
 =\operatorname{Desc}(T_{ij}\otimes W_i\otimes W_j^\vee)
 \quad(i\le j)
\]
と定義する。$A_i$は$U_2$上のAzumaya代数であり
\begin{equation}
 \operatorname{rk}(A_0,A_1,A_2,A_3)=(1,4,1,4),\qquad
 ([A_0],[A_1],[A_2],[A_3])=(0,\alpha_c,0,\alpha_d).
 \label{eq:qa-morita-vertex-classes}
\end{equation}
（oppositeの規約による符号は二捩れのため影響しない。）
$\mathsf M_{ij}$は有限局所自由な$(A_i,A_j)$双加群、
$\mathsf M_{ii}=A_i$であり、
\begin{equation}
 {\mathsf M}_{ij}\otimes_{A_j}{\mathsf M}_{jk}
 \longrightarrow{\mathsf M}_{ik}
 \qquad(i\le j\le k)
 \label{eq:qa-morita-bimodule-product}
\end{equation}
は結合的で単位を持つ。
この四頂点Azumaya双加群有向圏は
$U_2$の任意の基底変更と可換する。
分岐根を選ぶ有限étale分裂被覆の上で
$V_i$と$\det V_c$をさらに局所自明化すると、
命題~\ref{prop:qa-point-relative-split-root-end}の
spinor有限End代数の通常のMorita増幅に一致する。
\emph{相対AS完成と相対pointed $\qgr$の同値は結論しない。}
\end{theorem}
\begin{proof}
$\pi:E\to\mathscr G$は$(2,2)$ Cartier族である。
$i<j$なら$L_jL_i^{-1}$の各既約成分上の次数は正である。
相対$(2,2)$分解とcohomology-and-base-changeによって
$R^1\pi_*(L_jL_i^{-1})=0$となり、
$T_{ij}$は指定階数の有限局所自由束で基底変更と可換する。
対角は$\pi_*\OO_E=\OO_{\mathscr G}$である
（$R^1\pi_*\OO_E$の消滅は用いない）。

gerbeの慣性$(\lambda_c,\lambda_d)$は
$N_c,N_d$を各々の文字でスカラー倍する。
$F,G$の重みは$(1,0),(0,1)$、
二次形式の係数に由来する$H$の重みは$(0,0)$である。
$T_{ij}$の重みは$w_j-w_i$となり、
表の全項が従う。

$W_i$の重みは$w_i$である。
$\End(W_i)$は重み零、
$T_{ij}\otimes W_i\otimes W_j^\vee$の重みは
$(w_j-w_i)+w_i-w_j=0$である。
対角化可能群をbandとするgerbeの重み零sheafは
基底へ一意に降下する
\cite[Theorem~4.7]{BerghSchnurer2021}。
そのため$A_i$と$\mathsf M_{ij}$が存在する。
$V_c,V_d$はそれぞれのBrauer類のrank-two twisted bundleなので、
$\End(V_c),\End(V_d)$の降下は階数四Azumayaである。

$W_i$への左$\End(W_i)$作用と
$W_j^\vee$への右$\End(W_j)$作用で
$\mathsf M_{ij}$を双加群にする。
合成には線束の乗法$T_{ij}T_{jk}\to T_{ik}$と
\[
 (W_i\otimes W_j^\vee)
 \otimes_{\End W_j}(W_j\otimes W_k^\vee)
 \simeq W_i\otimes W_k^\vee
\]
というMorita evaluationを用いる。
結合律と単位はこの二つの構成のそれから従う。
重み零なので降下後も各等式は成り立つ。

分岐根選択の被覆上では両conicが$\PP^1$-bundleに分裂し、
各$V_i$を局所自明化できる。
すると$A_i$は行列代数、$\mathsf M_{ij}$は
$T_{ij}$のMorita増幅になる。
式~\eqref{eq:qa-morita-fg-square}で
$\det V_c$を自明化すると四項は
$(\OO_E,F,F^2,F^2G)$に一致する。
命題~\ref{prop:qa-point-relative-split-root-end}を適用して
局所spinor有限EndとのMorita同型が得られる。
一般の全整数AS完成の降下は使っていない。
\end{proof}

\begin{corollary}[ruling交換の頂点固定Morita降下は不可能]
\label{cor:qa-morita-ruling-exchange-phase}
$U_2$の二つの二次形式の交換$\epsilon$は
上記の頂点Brauer類を
\[
 (0,\alpha_c,0,\alpha_d)
   \longmapsto(0,\alpha_d,0,\alpha_c)
\]
と交換する。
一般点では$\alpha_c,\alpha_d$が独立なので、
この交換を全四頂点\emph{固定}のMorita自己同値として
降下させることは不可能である。
特にO-pointingを固定した相対比較でも、
残りのcollectionの位相変更・mutationを省略して
頂点固定同型を主張してはならない。
適切なmutationを伴う同値の存在・不存在は
この結論に含まない。
\end{corollary}
\begin{proof}
頂点固定のMorita同値が存在すれば、
特に一番頂点のAzumaya代数$A_1$と
$\epsilon^*A_1$はMorita同値となる。
そのBrauer類はそれぞれ$\alpha_c,\alpha_d$なので、
$\alpha_c=\alpha_d$が必要である。
しかし命題~\ref{prop:qa-universal-conic-brauer-obstruction}は
一般係数体上で$\alpha_c-\alpha_d\ne0$を証明した。
この矛盾は頂点固定の同値だけを排除し、
異なる四項の導来mutationには適用されない。
\end{proof}

% !TeX root = ../nc-quadric-avatar.tex
% 2026-11-22 / symmetric foundation and equivariant pointed Morita descent
\section{対称foundationと無標識商へのMorita降下}
\label{sec:qa-symmetric-morita-descent}
本節では有限圏の射を$\Hom(T_i,T_j)$と書き、通常の合成を持つ
$\Gamma=\End(\bigoplus_iT_i)$の右Yoneda加群$e_i\Gamma$を使う。
従って$\Hom(e_i\Gamma,e_j\Gamma)=e_j\Gamma e_i$である。
前節の$\mathsf M_{ij}$を上三角位置$e_iAe_j$に置く規約とは
反対環で移り合う。以下のMorita比較にはこの向きの変更を含める。

\subsection{二つのrulingを対称に置く相対foundation}

\begin{proposition}[対称squareへの相対mutation]
\label{prop:qa-symmetric-square-relative-mutation}
命題~\ref{prop:qa-point-relative-split-root-end}の順序付き分解普遍底を
$S$とし、$T$を任意の可換$S$-代数とする。
四頂点の有限代数$\mathfrak B_T$の標準射影対象を
$(U_0,U_1,U_2,U_3)$と書き、元の四項
$(\mathcal B,\mathcal X_c,\mathcal B(1),\mathcal X_d(1))$
と同じ順に標識する。
\[
 V=\Hom(U_2,U_3),\qquad
 K=\operatorname{Cone}(V\otimes_TU_2\xrightarrow{\rm ev}U_3)[-1]
\]
と置く。このとき
\begin{equation}
 (U_0,U_1,K,U_2)
 \tag{Sym-square}\label{eq:qa-symmetric-square-foundation}
\end{equation}
は相対full strong exceptional collectionである。
その有限End代数$\Gamma_{\rm sq,T}$は階数$16$の有限局所自由代数で、
全Hom、積、上のmutation、および
\[
 \Perf\mathfrak B_T\simeq\Perf\Gamma_{\rm sq,T}
\]
は任意の$T$-基底変更と可換する。
この同値は$U_0$を保つ。
また$E_T$上の四線束を用いる自然な同定
\begin{equation}
 \Gamma_{\rm sq,T}\simeq
 \End_{\rm dir}^{E_T}(\OO,F,G,H),\qquad H=F^2=G^2
 \label{eq:qa-symmetric-square-curve-end}
\end{equation}
を持つ。そのHom階数行列は
\[
 \begin{pmatrix}
 1&2&2&4\\0&1&0&2\\0&0&1&2\\0&0&0&1
 \end{pmatrix}.
\]
ここで曲線上の\emph{導来}Endとの同一視はしていない。
対角の$\Ext^1_E(L,L)=H^1(E,\OO_E)$は消えない。
\end{proposition}
\begin{proof}
$M_d^2=0$の周期的線形分解は、全ての係数環上で
\begin{equation}
 0\longrightarrow X_d\longrightarrow B(1)^{\oplus2}
 \longrightarrow X_d(1)\longrightarrow0
 \label{eq:qa-symmetric-square-spinor-evaluation}
\end{equation}
を与える。完全性は$B$の交互語PBW表示、または中心
$T[u,v,h]$（$h=cd+dc$）上の四元自由表示から確認できる。
従って全幾何fiberで$K$は$X_d$に一致する。
二型は根対一致の場合にも相異なる。
実際、同型なら点曲線への制限で$F\simeq G$となり、
補題~\ref{lem:qa-pointed-all-rootpair-geometry}の
$H^*(E,GF^{-1})=0$に反する。
strongnessを有限Homから確認する。
評価三角に$\RHom(U_i,-)$を適用した三つの写像は
\[
 V\otimes T_{02}\longrightarrow T_{03}\quad(8\to6),\qquad
 V\otimes T_{12}\longrightarrow T_{13}\quad(4\to4),\qquad
 V\otimes T_{22}\longrightarrow T_{23}\quad(2\to2)
\]
である。最初は曲線の評価列から全射で核は$H^0(G)=V$、
第二は同列に$F^{-1}$を掛けた$H^*(GF^{-1})=0$から同型、
第三は恒等写像である。逆向きには元の有向性から
$\RHom(K,U_0)=\RHom(K,U_1)=0$、
$\RHom(K,U_2)=V^\vee[0]$を得る。
また$\RHom(K,U_3)=(\End V/\OO)[0]$なので、
評価三角へ$\RHom(K,-)$を適用すると
$\End V\to\End V/\OO$の核から$\RHom(K,K)=\OO[0]$となる。
これで新四項の全非零次数Ext消滅と表示したHom階数が確認できる。
生成性は同じ評価三角から従う。
特に二つのspinorの相互Ext消滅を省略して推論してはいない。

まず被約な普遍底$S$上で議論する。
$U_i$は有限代数のprojective、$K$はprojectiveの二項複体だから、
上の四項間の各$\RHom$は$S$上perfectである。
全幾何fiberで次数零に集中するので、perfect fiber判定により
それ自体が表示した階数の局所自由束となる。
この同定は非平坦なものを含む任意の基底変更に保たれる。
従って$T$上の主張も成立する。
また$U_3$は$K$と$U_2$から評価三角で回収されるので、
四項は元の全projectiveを生成し、fullnessも相対的に成立する。
tilting Moritaの構成はこの評価と導来tensorから得られ、
$U_0$を動かさず任意の基底変更と整合する。

曲線上ではpencil $G$の評価列に$H=G^2$を掛けると
$0\to G\to H^{\oplus2}\to HG\to0$となる。
元の四頂点Endの同定の下で、この列は上のmutationに対応する。
二つの中間項の間は$H^*(E,GF^{-1})=0$、
他の前向きHomは正次数線束の切断なので
$(2,2)$ Cartier分解から全基底変更可換な局所自由束となる。
制限は全幾何fiberで同型であり、両辺が局所自由なので
全Homと積について式~\eqref{eq:qa-symmetric-square-curve-end}を得る。
階数の和は$4+(2+2+4+2+2)=16$である。
\end{proof}

\begin{corollary}[square avatarそのものの零頂点pointing]
\label{cor:qa-symmetric-square-field-pointing}
標数零の代数閉体上、二形式がそれぞれsquarefreeなら、
共有根と根対一致を含め、点曲線側square avatar
\[
 (E_f,F_d,F_c,P_E),\qquad P_E=F_dF_c^{-1}
\]
のsquare型AS $\mathbb Z$-代数を$A_{\rm sq}$とすると
\[
 (\qgr B_f,\pi B_f)\simeq(\qgr A_{\rm sq},\pi e_0A_{\rm sq})
\]
というpointed exact同値がある。
従来の最小十本矢型の一添字shiftと矛盾せず、
ここでは末尾のruling evaluationで対称squareへ戻した。
\end{corollary}
\begin{proof}
$F=F_c,G=F_d,P=GF^{-1}$と置く。
square avatar $(E,G,F,P)$の初期曲線四項は
$(P,G,F,GF)$である。
$P^2=\OO$、$F^2=G^2=H$を使い全項を$P^{-1}$でtensorすると
\[
 (\OO,F,G,H)
\]
になる。共通tensorは全Homと積を保つ。
前命題により、その有限Endは
$(\pi B_f,\pi X_c,\pi X_d,\pi B_f(1))$のEndに一致する。
両側はfull strong foundationで、Serre完成は有限Endから一意に定まる。
補題~\ref{lem:qa-as-serre-unique-recovery}、
命題~\ref{prop:qa-as-avatar-coherent-morita}と
分類論文Aのnormal helixおよびAZ再構成を適用する。
初期零頂点が$\pi B_f$なので同値はpointedである。
\end{proof}

\begin{proposition}[determinantを保つ五十四次元表示から三十次元表示へのmutation]
\label{prop:qa-symmetric-morita-rank30-mutation}
定理~\ref{thm:qa-twisted-four-vertex-morita}のgerbe上で、
$F,G,H,V_c,V_d$を同じ記号とする。
四項$(\OO,F,H,HG)$に付随する階数$54$の上三角代数の反対環$\mathsf A_{\mathrm{Mor}}^{\mathrm{op}}$から、
四項$(\OO,F,G,H)$と係数束
\[
 (W_0,W_1,W_2,W_3)=(\OO,V_c,V_d,\OO)
\]
に付随する有限Morita代数$\mathsf A_{\rm sq}$へ、
零頂点を保つ基底変更可換な相対tilting mutationが存在する。
$\mathsf A_{\rm sq}$は係数空間$U_2$上の階数$30$の有限局所自由代数で、
頂点のBrauer類は$(0,\alpha_c,\alpha_d,0)$である。
二形式交換は中間二頂点を交換し、零頂点を固定する。
\end{proposition}
\begin{proof}
gerbe上のpencil評価は
$0\to G^{-1}\det V_d\to\pi^*V_d\to G\to0$である。
式~\eqref{eq:qa-morita-fg-square}の
$G^2=H\pi^*\det V_d$を代入すると
\begin{equation}
 0\longrightarrow G\longrightarrow H\pi^*V_d
 \longrightarrow HG\longrightarrow0
 \label{eq:qa-symmetric-gerbe-evaluation}
\end{equation}
を得る。$\pi^*V_d^\vee$でtensorすれば
\begin{equation}
 0\longrightarrow G\pi^*V_d^\vee
 \longrightarrow H\pi^*\End(V_d)
 \longrightarrow HG\pi^*V_d^\vee\longrightarrow0.
 \label{eq:qa-symmetric-amplified-evaluation}
\end{equation}
全項は慣性重み零である。
従ってこの列と、その写像から作るfinite-categoryのmutationは
$U_2$へ降下する。
局所に$V_d$を自明化すれば前命題のmutationのMorita増幅に一致し、
評価写像の自然性から重なり上の比較はcocycleを満たす。
この議論は$\det V_d$を自明化しない。

新四項の重みは$(0,e_c,e_d,0)$であり、
各Homに$W_i\otimes W_j^\vee$を掛けると重み零となる。
対角階数は$(1,4,4,1)$、四辺は各$4$、
零頂点から終頂点への成分は$4$なので、総階数は
$1+4+4+1+4+4+4+4+4=30$である。
二つの中間項間のHomは零。
二形式交換は$F,V_c$と$G,V_d$を同時に交換し、
$\OO,H$を固定するため、全合成と零頂点を保つ。
逆mutationと対合のcocycleも評価・余評価の自然性から従う。
\end{proof}

\begin{remark}[相対heartへの比較に必要な追加論証]
有限代数間のtilting同値だけからheartの同値は従わない。
本件の点曲線側では、論文Bの
\cite[\texttt{lem:nc-relative-normal-tilting},
\texttt{prop:nc-root-ample-orbit}]{QuadricModuli2026}
が、元の二形式族の相対compact tiltingと
$t=S^{-1}[2]=(-)_\iota(2)$のheart上のexactな作用を与える。
同じ零頂点と相対Serreを保つ同値から
$\bigoplus_{m\ge0}\Hom(\OO,t^m\OO)$の全積付き同型を得た後に、
\cite[\texttt{prop:relative-h2-antiserre-reconstruction}]{QuadricModuli2026}
を適用する必要がある。
その適用範囲は元のrelative H2のsite（局所有限型、非被約を許す）である。
全可換環上の素朴なAbel公理や全QP stackの同値はこの議論に含まれない。
\end{remark}


\subsection{対称な有限Morita表示の全表示群への降下}
\label{subsec:qa-symmetric-full-group-descent}

二つのrulingを同じ位相に置く対称四項
$(\OO_E,F,G,H)$では、根対交換は中間二頂点だけを交換する。
この表示は、非対称四項
$(\OO_E,F,H,HG)$についての頂点固定Brauer障害と両立する。
ただし表示群の共通次数スカラーを除くには、
$H$の線型化の重みも処理しなければならない。
本項では二次元の共通座標表現による増幅を使い、
非被約基底変更にも適合する通常の有限代数として降下させる。

\begin{lemma}[conicの次数一束の選択に依らない階数二束]
\label{lem:qa-conic-canonical-rank-two-bundle}
$Q\to T$をsmooth conic族とし、$\mathscr G_Q\to T$を
相対幾何次数一の可逆層の$\Gm$-gerbeとする。
普遍束$N$と$V=q_*N$に対し、
$N\otimes q^*V^\vee$は$Q\times_T\mathscr G_Q$から
$Q$上の階数二ベクトル束$U_Q$へ降下する。
この構成はconic族の同型と任意の基底変更に関して自然であり、
次数一束の代表・その同型・その基底線束捩れの選択に依らない。
特にconicを恒等射として作用させる群は$U_Q$に自明に作用する。
\end{lemma}
\begin{proof}
$N$と$V$は同じ慣性指標を持つので
$N\otimes q^*V^\vee$の慣性重みは零である。
重み零sheafの降下
\cite[Theorem~4.7]{BerghSchnurer2021}を適用する。
選択独立性は局所表示でも直接見える。
$N'=N\otimes q^*M$とすれば、基底変更と射影公式により
$V'=V\otimes M$であるから、自然な縮約が
\[
 N'\otimes q^*(V')^\vee
 \simeq N\otimes q^*V^\vee
\]
を与える。二つの代表の同型をスカラー倍しても
そのスカラーは$N$と$V^\vee$で打ち消し合う。
この縮約は三つ以上の基底線束のtensor積について結合的であり、
Cechの三重重なり上のcocycleは厳密に一致する。
conicの同型によるpullbackと$q_*N$の基底変更同型も
この縮約と可換するため、自然性と最後の主張を得る。
\end{proof}

二形式の係数空間を$P=U_2$とし、論文Bの表示群を
\[
 \widetilde G=(\operatorname{GL}(W)\times\Gm^2)\rtimes S_2,
 \qquad
 D=\{(\lambda I_W,\lambda,\lambda)\},
 \qquad G=\widetilde G/D
\]
と書く。$W=ku\oplus kv$は共通二次元座標表現である。
$\widetilde G$の作用は
$(A,a_c,a_d)\cdot(f_c,f_d)
 =(a_c^{-2}\operatorname{Sym}^2(A)f_c,
   a_d^{-2}\operatorname{Sym}^2(A)f_d)$。
対応する次数付き座標輸送は
$W\mapsto A(W)$、$c\mapsto a_c c$、$d\mapsto a_d d$である。
このため$D$は$E\subset\PP^3_P$と両conicには恒等に作用し、
自然線型化を持つ$H=\OO_E(1)$と$W$にはいずれも重み一で作用する。

$p_c:E\to Q_c$、$p_d:E\to Q_d$に対し
\[
 U_c=p_c^*U_{Q_c},\qquad U_d=p_d^*U_{Q_d},\qquad
 K=H\otimes\pi^*W^\vee
\]
と置く。ここで$W$は$P$上の定数束に標準$\widetilde G$線型化を入れたもの。
$U_c,U_d$は個々には根対を順序付ける部分群上の対象だが、
$U_c\oplus U_d$は$S_2$の交換を含む全群に自然に同変である。

\begin{theorem}[無標識二形式商上の階数四十五のpointed有限代数]
\label{thm:qa-unordered-rank45-morita-descent}
$P$上の族$E$に
\[
 \mathcal V_\square
 =\OO_E\oplus U_c\oplus U_d\oplus K
 \qquad (\operatorname{rk}=1+2+2+2=7)
\]
を取る。この束の$\widetilde G$線型化では$D$が自明に作用し、
従って$[E/G]\to[P/G]=\mathfrak N^{\mathrm{pres}}$上へ降下する。
右加群のYoneda規約に合わせ、通常の合成を積とする
\begin{equation}
 \mathcal A_\square^{45}
   =\pi_*\mathcal End_E(\mathcal V_\square)
 \label{eq:qa-rank45-ordinary-end}
\end{equation}
と置く。従って$e_j\mathcal A_\square^{45}e_i$が
$i$番から$j$番へのHomであり、右射影加群$e_i\mathcal A_\square^{45}$の
間のHomもこの向きである。これは$\mathfrak N^{\mathrm{pres}}$上の
単位付き有限局所自由な通常の代数で、階数は$45$である。
全非被約基底変更を含む任意の基底変更と可換する。

根対を順序付ける二重被覆上では四頂点に分かれ、
各頂点の隅代数の階数は$(1,4,4,4)$、
非零非対角成分の階数は
\[
 \begin{array}{c|ccccc}
 (i,j)&01&02&03&13&23\\\hline
 \operatorname{rk}&4&4&8&8&8
 \end{array}.
\]
従って$45=(1+4+4+4)+(4+4+8+8+8)$である。
無標識商上では$e_0,e_3,e_1+e_2$が大域的な冪等元で、
中間二頂点はruling交換で置換される。
特に$e_0\mathcal A_\square^{45}e_0=\OO$であり、
すべての降下同型はこの指定頂点とその恒等元を保つ。

各代数閉幾何点では、この代数は
$(\OO_E,F,G,H)$の基本有限directed End
（次元$16$）のMorita増幅である。
係数空間$P$への引戻しは、同四項に階数
$(1,2,2,1)$の増幅を施した階数$30$の代数とMorita同値であり、
このMorita同値も零番の指定対象を保つ。
この定理は有限代数・perfect Morita表示の降下であり、
単独で相対AS完成や相対$\qgr$ heartの一致を結論しない。
\end{theorem}
\begin{proof}
補題~\ref{lem:qa-conic-canonical-rank-two-bundle}により、
$U_c,U_d$にはconic族の全同型について自然な輸送がある。
$\operatorname{GL}(W)$と二つの係数スカラーはconicの同型を誘導し、
$S_2$は二つの束を厳密に交換する。
$D$はconicには恒等なので$U_c,U_d$に自明に作用する。
一方$K=H\otimes W^\vee$上の$D$重みは$1-1=0$である。
$S_2$は$W$を保ち、$H$には$c,d$の座標交換として作用する。
従って$\mathcal V_\square$には全$\widetilde G$線型化があり、
$D$を通じて自明なので$G$線型化へ降下する。
合成同型は幾何的pullbackとtensor縮約の自然同型から作るため、
二回の交換も三重重なりのcocycleも恒等に一致する。
$\OO_E$の直和因子には単位sectionを固定する標準線型化を使う。

有限局所自由性と階数はgerbe上で計算できる。
この基底では
\[
 U_c=F\otimes V_c^\vee,
 \quad U_d=G\otimes V_d^\vee,
 \quad K=H\otimes W^\vee.
\]
全幾何fiberで$GF^{-1}$は非自明な多重次数零線束である。
smooth/nodal/$I_2$のいずれでも
$H^0(GF^{-1})=H^1(GF^{-1})=0$であり、逆の束にも同じ消滅がある。
実際、次数零の非零sectionは全成分上で零点を持たず、
線束の自明化となるので非自明性に反する。
$H^1$には$\omega_E\simeq\OO_E$とSerre双対を使う。
相対コホモロジー・基底変更により、この消滅は
任意の非被約引戻し後にも成立する。

残る前向きHomの基本階数は
\[
 \dim H^0(F)=\dim H^0(G)
 =\dim H^0(HF^{-1})=\dim H^0(HG^{-1})=2,
 \qquad\dim H^0(H)=4.
\]
これらは各成分上正次数で$H^1=0$である。
後向きの負次数Homは零、対角は$\pi_*\OO_E=\OO_P$である。
ここで対角$R^1\pi_*\OO_E$の消滅は仮定しない。
全Homの基底変更可換性は、既存の$(2,2)$ Cartier族の計算、
あるいは$E\to\PP^1_P$の
$\OO\oplus\OO(-1)^{\oplus2}\oplus\OO(-2)$という直像表示から従う。
従って上記の五つの基本Homに
$r_ir_j$、$r=(1,2,2,2)$を掛けた階数が表の値となる。
対角は$r_i^2$であり、全階数は$45$。
積は実際の束の射の合成から来るため結合的で、
基底変更および群の輸送と可換する。

二次形式の交換は$U_c,U_d$の直和因子だけを入れ替え、
$\OO_E$と$K$を保つ。
従って主張した冪等元は無標識商へ降下する。
零番には増幅を施していないので、その隅は$\OO$のままである。
各幾何点では$V_c,V_d,W$は二次元の通常ベクトル空間であり、
元の四つの射影generatorをそれぞれ$1,2,2,2$個に増幅した代数になる。
全ての増幅階数が正なので基本代数とMorita同値である。
係数空間上では$W$は通常の二次元自明束であるから、
最後の因子を二回重ねることだけを除けば階数$30$の増幅を得る。
この加法generatorの変更は零番因子を一個のまま保つ。
\end{proof}

\begin{remark}[頂点Brauer類と増幅の役割]
\label{rem:qa-rank45-polarization-brauer}
順序付き二重被覆上で、階数四の中間頂点のBrauer類は
conicの類$\alpha_c,\alpha_d$である。
最後の隅は$\operatorname{End}(W^\vee)$の$G$同変降下であり、
基底の$\PP(W)$に対応する階数四Azumaya代数である。
これは$\operatorname{End}(W)$の降下の反対代数であるが、
Brauer類は二捩れなので同じ類を表す。
係数空間上では行列代数だが、無標識商上のその類を
追加の論証なしに零と扱わない。

順序付き部分群では、$a_c$の一次元指標を用いて
$H\otimes a_c^{-1}$を作るだけでも共通次数スカラーを消せる。
しかし根対交換はこの選択を$H\otimes a_d^{-1}$に変えるため、
対称な厳密線型化を自動的には与えない。
$H\otimes W^\vee$という増幅は、
二次元座標表現が根対交換で保たれることを利用して
この非対称な選択を避ける。
これは基本線束の下降を証明する操作ではなく、
その障害をAzumaya隅に保持した有限Morita表示である。
またruling交換は中間二頂点を実際に置換するため、
系~\ref{cor:qa-morita-ruling-exchange-phase}の
\emph{四頂点固定}の不可能性と矛盾しない。
\end{remark}


\subsection{実際のspinor加群の降下と内在的compact圏}
\label{subsec:qa-symmetric-spinor-gluing}

以下の構成は、有限End代数を曲線上で降下させるだけでなく、
内在的な相対noncubic heartの中のgeneratorを直接与える。
通常のEnd代数と、前節の上三角双加群代数の積の向きを区別する。
本項では$\Gamma=\operatorname{End}(\bigoplus T_i)$に対し
$P_i=\operatorname{Hom}(\bigoplus T_j,T_i)=e_i\Gamma$を右加群とし、
$\operatorname{Hom}_\Gamma(P_i,P_j)=e_j\Gamma e_i$
というYoneda規約を一貫して用いる。
旧表示の$\mathsf M_{ij}$を$e_i\mathsf A_{\mathrm{Mor}}e_j$に置く
上三角積は、同じ束の通常Endの\emph{反対環}である。
その表示を本項の右Yonedaに移す際は
$\mathsf A_{\mathrm{Mor}}^{\mathrm{op}}$を用いる。
階数と二捩れBrauer類は変わらないが、
この反対環の操作を省略して右加群圏を同一視しない。

\begin{lemma}[conic section moduleからのspinorの自然構成]
\label{lem:qa-intrinsic-spinor-induced-module}
$P=U_2$上で
\[
 C=\OO_P[u,v],\quad
 R_i=C[c_i]/(c_i^2-f_i),\quad
 \mathcal B=C\langle c,d\rangle/(c^2-f_c,d^2-f_d),
 \qquad (c_c=c,\ c_d=d)
\]
と置く。$\mathscr G_i$上で
\[
 S_i(N_i)=\bigoplus_{n\ge0}
      q_{i*}(N_i\otimes\OO_{Q_i}(n)),\qquad
 X_i(N_i)=S_i(N_i)\otimes_{R_i}\mathcal B
\]
を次数付き右加群として定義する。
この構成は任意の基底変更と可換し、
根が分解する局所座標では既存の行列分解のcokernel $X_i$に一致する。
$\mathcal B$は左$R_i$加群として自由である。
各次数の$X_i(N_i)$は有限局所自由で、階数は$(n+1)(n+2)$である。

$V_i=q_{i*}N_i$とすれば
\begin{equation}
 Y_i=\operatorname{Desc}_{\mathscr G_i/P}
                  (X_i(N_i)\otimes V_i^\vee)
 \label{eq:qa-spinor-amplified-descent}
\end{equation}
は$P$上の実際の有限生成次数付き右$\mathcal B$加群である。
その次数$n$成分の階数は$2(n+1)(n+2)$。
構成はconicの同型、係数スカラー、共通$\operatorname{GL}(W)$作用に自然で、
二形式交換は$Y_c$と$Y_d$を厳密に交換する。
\end{lemma}
\begin{proof}
$\mathcal B$の中央係数$C$上のPBW基底は$c,d$の交互語である。
左$R_c$加群としては、空語と$d$から始まる交互語
\[
 1,d,dc,dcd,dcdc,\ldots
\]
を基底に取れる。各$C$-交互語はこの基底の一つ、または
その左に$c$を付けたものに一意に分かれるためである。
$d$側も同様であり、左$R_i$上の自由性が従う。

smooth conic上では$N_i\otimes\OO_{Q_i}(n)$の幾何次数は$2n+1$で、
$n\ge0$なら一次コホモロジーは零。
その直像は階数$2n+2$のベクトル束で任意の基底変更と可換する。
左$R_i$自由基底は各非負次数に一個なので、
$X_i(N_i)_n$の階数は
$\sum_{j=0}^n(2j+2)=(n+1)(n+2)$。
基底変更と可換する同じ直和記述が平坦性も与える。

根を分解して$Q_i\simeq\PP^1$、$N_i\simeq\OO(1)$とすると、
$S_i(N_i)$は二次Veronese環の奇数section moduleである。
$s=a^2,t=-b^2,c_i=ab$と置けば、
\[
 \begin{pmatrix}c_i&s\\t&-c_i\end{pmatrix}
  =\binom{a}{-b}(b\ \ a)
\]
であり、この行列のcokernelは
$\bigoplus_{n\ge0}H^0(\PP^1,\OO(2n+1))$である。
左$R_i$自由な$\mathcal B$へtensorしても表示は完全なので、
既存のspinor $X_i$を得る。特に有限生成性がfppf局所に成り立ち降下する。

$N_i$を$N_i\otimes q_i^*M$へ替えると
$S_i(N_i)$、$X_i(N_i)$、$V_i$はすべて$M$をtensorする。
従って式~\eqref{eq:qa-spinor-amplified-descent}のtensor縮約は
$M$を正確に消し、慣性指標も零になる。
補題~\ref{lem:qa-conic-canonical-rank-two-bundle}と同じcocycle計算で
$Y_i$は厳密に降下し、全同型と基底変更に自然である。
\end{proof}

\begin{proposition}[対称spinor generatorの相対pointed Morita表示]
\label{prop:qa-symmetric-spinor-relative-tilting}
上の記号を用い、locally finite typeな基底上の相対noncubic heartで
\begin{equation}
 \mathcal T_\square
 =\pi\mathcal B\ \oplus\ \pi Y_c\ \oplus\ \pi Y_d\
       \oplus\bigl(\pi\mathcal B(1)\otimes W^\vee\bigr)
 \label{eq:qa-rank45-actual-spinor-generator}
\end{equation}
を取る。これは相対tilting compact generatorである。
この対象とその構造対象因子は$\mathfrak N^{\mathrm{pres}}=[P/G]$上へ降下する。
自然な可換化による有限Homの比較は
\begin{equation}
 \operatorname{End}(\mathcal T_\square)
       \simeq\mathcal A_\square^{45}
 \label{eq:qa-rank45-actual-spinor-end}
\end{equation}
となり、任意のNoetherian基底変更と可換する。
従ってcanonical compact圏には、構造対象を零番射影対象へ送る
相対的な同値
\begin{equation}
 D(\operatorname{Ind}\mathcal C)^c
       \simeq\operatorname{Perf}(\mathcal A_\square^{45})
 \label{eq:qa-rank45-intrinsic-perfect-equivalence}
\end{equation}
を得る。ここで右辺は右加群のperfect圏であり、
$\mathcal O$の像は$e_0\mathcal A_\square^{45}$である。
このperfect圏の同値と、任意に選んだAS表示の標準heartが一致することは区別する。
\end{proposition}
\begin{proof}
最初に相対tilting性を根が分解するfppf被覆上で確認する。
$Y_i$はspinor $X_i$の階数二の増幅なので、
増幅を外した四項
$(\pi\mathcal B,\pi X_c,\pi X_d,\pi\mathcal B(1))$
を考えればよい。
$G=p_d^*N_d$の評価列と
$G^2=H\det V_d$から、conic上およびそのsection moduleに
\[
 0\longrightarrow N_d\longrightarrow
 V_d\otimes\OO_{Q_d}(1)\longrightarrow
 N_d\otimes\OO_{Q_d}(1)\longrightarrow0
\]
を得る。この列を各非負次数で直像してから
左$R_d$自由な$\mathcal B$へtensorすると、実際の完全列
\begin{equation}
 0\longrightarrow X_d\longrightarrow
 V_d\otimes\mathcal B(1)\longrightarrow X_d(1)
       \longrightarrow0
 \label{eq:qa-spinor-relative-euler-sequence}
\end{equation}
を得る。すべて任意の基底変更と可換する。

論文Bの相対smooth proper性、compactness、相対normal tiltingの議論
\cite[\texttt{prop:nc-relative-smooth},
\texttt{prop:nc-compact-generator},
\texttt{lem:nc-relative-normal-tilting}]{QuadricModuli2026}
を普遍二形式族へ適用すると、
$(\pi\mathcal B,\pi X_c,\pi\mathcal B(1),\pi X_d(1))$
は相対full strong foundationである。
ここで特定の正規化された係数形に範囲を制限しない。
根分解被覆上で両spinorは中心$\OO[u,v,cd+dc]$上
$Z^2\oplus Z(-1)^2$となる同じ行列計算を持ち、
全幾何fiberでAのnormal foundationに一致する。
このためBの証明で用いる相対perfectness、derived Homの基底変更、
fiberwise full性の検出が同じまま適用できる。

式~\eqref{eq:qa-spinor-relative-euler-sequence}は末尾二項のevaluation mutationである。
命題~\ref{prop:qa-symmetric-square-relative-mutation}の有限Hom計算により、
これを$(\pi X_d,\pi\mathcal B(1))$に替えた四項も相対full strongとなる。
その計算は$H^*(GF^{-1})=0$と二つのpencilの評価列を使い、
幾何点だけのmutationから相対性を推論するものではない。
正の階数の局所自由束による増幅は生成する加法閉包を変えないので、
$\mathcal T_\square$もtilting compact generatorである。

$\mathcal B\to\mathcal B^{\mathrm{ab}}$による制限は
$Y_i$を$U_i$に、$\mathcal B(1)\otimes W^\vee$を$K$に送る。
順序付き根分解被覆上では既存の相対有限End比較
（命題~\ref{prop:qa-point-relative-split-root-end}）と
式~\eqref{eq:qa-spinor-relative-euler-sequence}のmutation比較から、
この制限は対称四項のすべての有限Homと積の同型を与える。
対象・制限射・評価射の三つがconic section moduleから自然に作られるため、
これらの同型は重なりで一致して式~\eqref{eq:qa-rank45-actual-spinor-end}へ降下する。

最後に共通次数スカラー$\lambda\in D$を追う。
$Y_i$の次数$n$には$\lambda^n$が作用し、
$\mathcal B(1)\otimes W^\vee$にも
$\lambda^{n+1}\lambda^{-1}=\lambda^n$が作用する。
$\mathcal B$の次数$n$も同じ重みである。
次数スカラーによる右加群の輸送には
\[
 m_n\longmapsto\lambda^{-n}m_n
\]
という自然同型があり、この同型が四因子すべてに共通に働く。
特に構造対象の次数零の単位では恒等である。
従って$\widetilde G$の作用を$G$へ剛体化する降下データと
$\mathcal T_\square$の線型化は完全に整合する。
二形式交換は中間二因子を交換し、残る二因子とpointingを保つ。
相対tilting Morita
\cite[Sections~8.2 and 9.2]{Keller1994}
を適用すると式~\eqref{eq:qa-rank45-intrinsic-perfect-equivalence}を得る。
end代数・tilting対象・その評価はすべてbase changeに可換なので、
この同値も任意のNoetherian基底変更に整合する。
\end{proof}


% !TeX root = ../nc-quadric-avatar.tex
% 2026-11-22 / all-degree relative Serre completion and equivariant heart
\section{相対Serre完成と構造対象付きheart}
\label{sec:qa-relative-serre-completion}
\subsection{有限Azumaya表示の相対Serre tensor代数}
\label{subsec:qa-relative-morita-serre}

以下でいうcompletionは，内部次数に関する\emph{直和}
$T_A(\Theta)=\bigoplus_{m\ge0}\Theta^{\otimes_A^{\mathbf L}m}$
である．無限直積によるarrow-adic completionの基底変更は
この主張には含めない．また，この構成だけで
AS正則性，Noether性，または標準$\qgr$心の一致を結論しない．

\begin{lemma}[Azumaya頂点を持つ有向代数の有限対角解消]
\label{lem:qa-directed-azumaya-bar}
$S$をスキームとし，$A$を有限局所自由な
$\OO_S$-代数とする．（必要なら順を反転して）大域的な直交冪等元
$e_0,\ldots,e_{r-1}$の和を$1$とし，
$e_iAe_j=0$ ($i>j$)，$D_i=e_iAe_i$は
$\OO_S$上Azumaya，$M_{ij}=e_iAe_j$は
$\OO_S$上有限局所自由と仮定する．
$D=\bigoplus_iD_i$，$J=\bigoplus_{i<j}M_{ij}$と置く．
すると$J^{\otimes_Dr}=0$であり，
\begin{equation}
0\longrightarrow A\otimes_DJ^{\otimes_D(r-1)}\otimes_DA
\longrightarrow\cdots\longrightarrow
A\otimes_DJ\otimes_DA
\longrightarrow A\otimes_DA\longrightarrow A\longrightarrow0
\label{eq:qa-directed-finite-bar}
\end{equation}
は$A$の相対enveloping代数
$A^e=A\otimes_{\OO_S}A^{\mathrm{op}}$上の
有限局所射影解消である．この解消は任意の基底変更と可換する．
特に$A$は$S$上properかつhomologically smoothである．
同じ結論は，$D\subset A$と増大冪零ideal $J$，
$A=D\oplus J$が大域的に与えられ，$D$が
$\OO_S$上有限局所自由separableで，
あるfppf被覆上で上記のAzumaya頂点表示を持つ場合にも成り立つ．
頂点そのものが大域的に番号付けできる必要はない．
\end{lemma}
\begin{proof}
$J$のtensor成分は
$M_{i_0i_1}\otimes_{D_{i_1}}\cdots
\otimes_{D_{i_{n-1}}}M_{i_{n-1}i_n}$，
$i_0<\cdots<i_n$，に分かれるので$n\ge r$では零である．
$D_i\otimes D_j^{\mathrm{op}}$はAzumayaであり，
その上の$\OO_S$-有限局所自由加群は
代数上も有限局所射影である．これはAzumaya代数を
行列代数にするfppf被覆上でMorita同値を用いて確認でき，
有限局所射影性の忠実平坦降下によって元の底へ戻る．
同じ議論とMorita contractionから，
$J^{\otimes_Dn}$は$D^e$上有限局所射影となる．
従って各誘導加群
$A\otimes_DJ^{\otimes_Dn}\otimes_DA$は
$A^e$上有限局所射影である．

$A=D\oplus J$という$D$-双加群分解により，
式\eqref{eq:qa-directed-finite-bar}は相対bar解消の
normalized部分である．微分は隣接する二因子の積の
交代和で，最後の写像は乗法である．
bar解消の単位挿入による収縮とdegenerate部分の除去により
増大複体は完全となる．この議論は基礎環が体であることを用いない．
分解$A=D\oplus J$，tensor成分，微分と収縮は
すべて任意の基底変更に保たれる．
したがって非平坦な基底変更の後にも同じ解消が得られる．
$A$自体は$\OO_S$上有限局所自由なのでproperでもある．
最後の拡張では，$D^e$上の有限局所射影性と
$J^{\otimes_Dr}=0$を与えられた被覆上で確認する．
bar複体は大域的な$D,J$だけから定義されるので，
完全性と有限局所射影性の忠実平坦降下で同じ結論を得る．
\end{proof}

\begin{proposition}[相対inverse dualizing complexと全tensor次数]
\label{prop:qa-relative-serre-tensor-base-change}
前補題の仮定の下で
\[
\Theta_A=\RHom_{A^e}(A,A^e)[2],\qquad
K_m(A)=\Theta_A^{\otimes_A^{\mathbf L}m}\quad(m\ge0),
\qquad K_0(A)=A
\]
とする．各$K_m(A)$は$\OO_S$上perfectであり，
任意の$f:S'\to S$について自然な同型
\begin{equation}
Lf^*K_m(A)\simeq K_m(f^*A)
\label{eq:qa-serre-tensor-all-base-change}
\end{equation}
が積と整合する．従って
$\Pi_3(A)=\bigoplus_{m\ge0}K_m(A)$は
任意基底変更と可換な内部次数付きdg代数を定める．

さらに，全幾何点$\bar s\to S$，すべての$m\ge0$について
\begin{equation}
H^pK_m(A_{\bar s})=0\qquad(p\ne0)
\tag{Fwd}\label{eq:qa-relative-geometric-forward}
\end{equation}
を仮定する．この場合$K_m(A)$は
有限局所自由加群$C_m(A)=H^0K_m(A)$に準同型で，
\[
C(A)=\bigoplus_{m\ge0}C_m(A)
\]
は通常の内部次数付き代数となる．各次数と積は
任意の基底変更と可換する．
基底の被約性・正則性を仮定する必要はない．
\end{proposition}
\begin{proof}
前補題の有限局所射影解消を$A^e$で双対化する．
これにより$\Theta_A$は，有限局所射影$A^e$-加群の
有界複体で表される．$A$が$\OO_S$-有限局所自由なので，
各項は左右の$A$-加群としても局所射影であり，
$\OO_S$上も有限局所自由である．
したがって各固定$m$の導来tensor冪は，
この代表を$m$回tensorした有界複体で計算できる．
その各項は$\OO_S$上有限局所自由である．
有限局所射影加群の双対・tensorは任意基底変更と可換し，
先のbar解消も基底変更後の対角を解消するので
式\eqref{eq:qa-serre-tensor-all-base-change}が従う．
結合と単位との適合はtensor積の自然性から従う．
導来引戻しは直和を保つので，内部次数付きdg代数全体も
基底変更と可換する．

最後の主張はperfect複体のfiberによるTor-amplitude判定である．
ここでは局所的な有限自由複体を取り，各点の局所環で
微分行列に現れる可逆成分を収縮して極小複体にする．
剰余体へのtensorで微分が零となったこの複体に
次数零以外のcohomologyがないなら，
次数零以外の自由項はすべて零である．
使った可逆小行列はその点の開近傍で可逆であるから，
元の複体も近傍で次数零の有限自由加群に準同型となる．
これを$K_m(A)$へ各$m$ごとに適用する．
非被約底上の行列をfiber階数一定だけで分裂させたのではなく，
\emph{perfect複体の全導来fiberの次数集中}を使った点が重要である．
この結果から$H^0$と任意基底変更の可換性も従う．
全$m$を一度に制御する有限個の次数判定を仮定していない．
\end{proof}

\begin{proposition}[pointed Morita同値によるSerre tensor代数の移送]
\label{prop:qa-pointed-morita-serre-transport}
$A,B$を$S$上proper smoothな代数とし，
$P$を$(B,A)$-双加群，$Q$を$(A,B)$-双加群とする
可逆なperfect Morita datum
\[
P\otimes_A^{\mathbf L}Q\simeq B,
\qquad Q\otimes_B^{\mathbf L}P\simeq A
\]
を固定する．この同型には通常の単位・余単位の整合性を課す．
すると全$m\ge0$にわたり
\begin{equation}
Q\otimes_B^{\mathbf L}K_m(B)\otimes_B^{\mathbf L}P
\simeq K_m(A)
\label{eq:qa-pointed-morita-theta}
\end{equation}
が積と整合し，$\Pi_3(A)$と$\Pi_3(B)$は
内部次数を保って導来Morita同値となる．

さらに$fB$，$eA$を指定した右射影対象とし，
指定同型$fB\otimes_B^{\mathbf L}P\simeq eA$を仮定する．
このとき全次数のpointed orbit複体
\begin{equation}
\RHom_B(fB,fB\otimes_B^{\mathbf L}K_m(B))
\simeq
\RHom_A(eA,eA\otimes_A^{\mathbf L}K_m(A))
\label{eq:qa-pointed-morita-orbit}
\end{equation}
は積と整合する．両辺が次数零に集中する場合には
構造対象の反Serre軌道環の同型となる．
この構成は合成するMorita datumおよび任意基底変更と整合する．
従って，既にcoherentなpointed Morita降下datumがあるなら，
そのdatumはSerre tensor代数とpointed orbit代数へ移送される．
\end{proposition}
\begin{proof}
双加群圏の同値
$X\mapsto Q\otimes_B^{\mathbf L}X\otimes_B^{\mathbf L}P$
は対角$B$を$A$へ移し，tensor積を保つ．
有限perfect対角の双対化とも整合するので，
$\RHom_{B^e}(B,B^e)$は
$\RHom_{A^e}(A,A^e)$へ移る．
これは\cite[Proposition~3.10(e)]{KellerDeformedCY2011}
のinverse dualizing complexのMorita共変性である．
したがって$m=1$の式が成り立つ．
中間因子$P\otimes_A^{\mathbf L}Q\simeq B$の縮約により
全$m$の式と積への適合を得る．
tensor代数の導来Morita同値は
\cite[Proposition~4.2]{KellerDeformedCY2011}，
または同じ縮約を内部次数ごとに用いた双加群で得られる．
指定射影対象の同型とMorita全忠実性を適用すれば
式\eqref{eq:qa-pointed-morita-orbit}が従う．

右加群の軌道積は本稿の規約
$a*b=s^n(a)\circ b$（$\deg b=n$）であり，
同値と$s=-\otimes_A^{\mathbf L}\Theta_A$の交換同型を使うため
この順序も保存される．
合成・基底変更への適合は，選択したMorita datumの
単位・余単位の整合性と双対化の自然性から従う．
任意の二つの局所Morita同値が自動的にcoherentな
降下datumを成す，という主張はしていない．
\end{proof}

\begin{remark}[四頂点模型への適用と残る認識問題]
\label{rem:qa-relative-serre-application-scope}
階数五十四の$\mathsf A_{\mathrm{Mor}}$とその反対環には
$r=4$のbar解消を無条件に適用できる．
定理~\ref{thm:qa-unordered-rank45-morita-descent}のsquare模型では，
その非零な厳密増大chainは長さ二までなので
$J^{\otimes_D3}=0$となり，対角解消の長さは二に短縮する．
これは一般の四頂点有向代数に大域次元二を仮定することと異なる．

右Yonedaの比較には$\mathsf A_{\mathrm{Mor}}^{\mathrm{op}}$を使う。
幾何fiberにおけるfinite EndのMorita比較と
全整数helixのforward Ext消滅を証明済みの入力に取れば，
条件\eqref{eq:qa-relative-geometric-forward}を確認できる．
この場合，非被約底も含む通常のSerre tensor代数まで得られる．
しかしfinite Endの同型，そのdg completion，
全次数の次数零集中，Noetherian $\qgr$心と元の心の同値は
それぞれ別の段階である．Noether性とheartに必要な相対豊富性・
tails認識はMorita形式論だけからは出ず、後の
定理~\ref{thm:qa-symmetric-relative-pointed-heart}で元の族のtilting表示を使う。
同様に，対象や二根対を交換する\emph{新しい}Morita datumの存在は
本命題の仮定であって結論ではない．
\end{remark}


\begin{corollary}[pointed反Serre環の相対二次表示]
\label{cor:qa-morita-relative-antiserre-quadratic}
命題~\ref{prop:qa-relative-serre-tensor-base-change}の
全fiber集中条件を仮定し，指定冪等元$e$について
$R_m=eC_m(A)e$と置く．さらに全幾何fiberの次数付き環
$R_{\bar s}=\bigoplus_mR_m\otimes\kappa(\bar s)$が
論文AのH2対象の反Serre軌道環に同型であると仮定する．
このとき
\[
\operatorname{rk}R_m=(2m+1)^2,\qquad
Q=\ker(R_1\otimes R_1\longrightarrow R_2)
\]
はそれぞれ有限局所自由で，$\operatorname{rk}R_1=9$，
$\operatorname{rk}Q=56$，かつ
\[
R\simeq T_{\OO_S}(R_1)/(Q).
\]
この表示は任意の基底変更と可換する．
\end{corollary}
\begin{proof}
$e$によるcornerは有限局所自由加群の冪等直和因子である．
幾何fiberでのHilbert関数から指定rankを得る．
各$m\ge2$について乗法
$R_1\otimes R_{m-1}\to R_m$は全幾何fiberで全射であり，
余核は有限型なのでNakayamaにより全射である．
特に$m=2$のkernel $Q$は局所的に分裂した有限局所自由加群となり，
基底変更と可換する．
$\widetilde R=T(R_1)/(Q)$と置く．
各$\widetilde R_m$は有限表示$\OO_S$-加群であり，
その構成は任意基底変更と可換する．
自然な全射$\widetilde R_m\to R_m$は
$R_m$が局所射影なので局所的に分裂する．
そのkernelは有限型である．
各幾何fiberでは論文Aの反Serre環の二次性により
この全射は同型だから，kernelの全fiberは零である．
再びNakayamaによりkernelが零となる．
この議論は全$m$について行い，有限次数の検算だけで
全次数の二次性を推論してはいない．
\end{proof}

\begin{remark}[square模型の全次数数値]
\label{rem:qa-square-amplified-serre-ranks}
基本square foundationのEuler行列を
\[
G=\begin{pmatrix}1&2&2&4\\0&1&0&2\\0&0&1&2\\0&0&0&1\end{pmatrix}
\]
とし，増幅数を$r=(1,2,2,1)^t$とする．
Serre作用の数値行列は$S=G^{-1}G^t$で，
$(S-I)^3=0$である．$t=2m$と置くと
\[
GS^{-m}=
\begin{pmatrix}
(t+1)^2&(t+1)(t+2)&(t+1)(t+2)&(t+2)^2\\
t(t+1)&(t+1)^2&t(t+2)&(t+1)(t+2)\\
t(t+1)&t(t+2)&(t+1)^2&(t+1)(t+2)\\
t^2&t(t+1)&t(t+1)&(t+1)^2
\end{pmatrix}.
\]
これは$S^{-m}=I-m(S-I)+\binom{m+1}{2}(S-I)^2$
から全整数$m$について従う．
したがってforward集中を\emph{別に}証明した場合には，
増幅したcompletionの第$m$成分のrankは
\[
r^tGS^{-m}r=144m^2+144m+30,
\]
零番cornerのrankは$(2m+1)^2$である．
$r=(1,2,2,2)^t$という別の増幅では
$r^tGS^{-m}r=196m^2+196m+45$となる．
数値恒等式自体はforward集中やAS正則性の証明ではない．
ここで使用するSerre関手は\emph{有限ordinary directed End代数}
のperfect加群圏のSerre関手であり，
genus一曲線上の線束直和のderived Endとは異なる．
後者には対角$\Ext^1$があり，本計算の代数に置き換えてはならない．
\end{remark}

\begin{corollary}[対称模型の全次数Serre代数]
\label{cor:qa-rank45-all-degree-serre}
$\mathsf A=\mathcal A_\square^{45}$とする。
その$\Pi_3(\mathsf A)$は全ての内部次数でcohomological degree零に集中する。
$C_m=H^0(\Theta_{\mathsf A}^{\otimes_{\mathsf A}^{\mathbf L}m})$とおくと
\[
 \operatorname{rk}C_m=196m^2+196m+45,\qquad
 R_m=e_0C_me_0,\quad \operatorname{rk}R_m=(2m+1)^2 .
\]
全成分と積は任意の基底変更と可換し、
$R=\bigoplus_m R_m$は九生成・階数五十六の二次関係を持つ。
\end{corollary}
\begin{proof}
対角部分はrulingを順序付ける有限\'etale被覆上でAzumaya代数の積となり、
その下では中間二頂点間のHomは零である。
従って大域のseparable対角部分$D$とradical $J$について
$J^{\otimes_D3}=0$であり、補題~\ref{lem:qa-directed-azumaya-bar}は
長さ二の相対解消を与える。
ここで通常Endの右Yoneda規約では、同補題の冪等元を逆順に並べればよい。
各幾何fiberでは$\mathsf A$は系~\ref{cor:qa-symmetric-square-field-pointing}
のfull strong square foundationの増幅である。
$m\ge1$に対する全ての成分は
$\RHom(E_i,S^{-m}E_j[2m])$であり、
$j+4m>i$という全整数helixの前向きExt消滅により次数零に集中する。
$m=0$は元の有限代数である。
命題~\ref{prop:qa-relative-serre-tensor-base-change}から
全基底上の局所自由性と基底変更が従う。
rankは注意~\ref{rem:qa-square-amplified-serre-ranks}、
cornerの二次性は系~\ref{cor:qa-morita-relative-antiserre-quadratic}による。
\end{proof}

\begin{lemma}[交換指標を含む相対Serre線形化]
\label{lem:qa-symmetric-equivariant-serre-character}
$P$上の普遍二形式環$\mathcal B$について、
$\widetilde G=(\GL(W)\times\Gm^2)\rtimes S_2$と
$D\subset\widetilde G$を論文Bと同じ群とする。
$\varepsilon:\widetilde G\to\{\pm1\}$を二形式交換の符号指標、
$\mathcal D=\det W$と置く。
$\iota$は$c,d$の同時符号反転である。
canonical compact圏の\emph{同変}相対Serre関手は
\begin{equation}
 S=(-)_\iota(-2)\otimes\mathcal D\otimes\varepsilon[2],
 \qquad
 t=S^{-1}[2]=(-)_\iota(2)\otimes\mathcal D^{-1}\otimes\varepsilon.
 \label{eq:qa-equivariant-serre-sign}
\end{equation}
従って構造対象のcanonical反Serre環は
\begin{equation}
 \mathcal R^{\rm can}_n
 =\mathcal B_{2n}\otimes\mathcal D^{-n}\otimes\varepsilon^n,
 \qquad a\star b=a\,\iota^n(b)\quad(a\in\mathcal R^{\rm can}_n).
 \label{eq:qa-canonical-antiserre-sign}
\end{equation}
ここで$\varepsilon$は基底上の同変線束であり、
二形式の順序を固定した恒等成分上では自明である。
$c,d$の独立スカラー指標は式に残らない。
\end{lemma}
\begin{proof}
中心を$Z=\OO_P[u,v,h]$、$h=cd+dc$とすると
$\mathcal B=Z\oplus Zc\oplus Zd\oplus Zcd$である。
最後の係数を取る$Z$-線形写像を$\lambda:\mathcal B\to Z$と書く。
基底$(1,c,d,cd)$に対するpairing
$(x,y)\mapsto\lambda(xy)$の行列は
\[
 \begin{pmatrix}
 0&0&0&1\\
 0&0&1&0\\
 0&-1&0&0\\
 1&0&0&h
 \end{pmatrix}.
\]
その行列式は$-1$であり、全基底上でperfectである。
また$\lambda(xy)=\lambda(y\iota(x))$なので、
これはNakayama作用$\iota$を持つFrobenius pairingである。

恒等成分の$(A,a_c,a_d)$では$h$は$a_ca_d$倍される。
中心のcanonical生成元$du\wedge dv\wedge dh$の指標は
$\det(W)a_ca_d$、双対側の$\lambda$の指標は
$(a_ca_d)^{-1}$である。
両者は相殺して$\det W$だけが残る。
一方、交換$c\leftrightarrow d$は$h$を固定し、
$cd\mapsto dc=h-cd$とする。
従って$\lambda\circ\varepsilon=-\lambda$であり、
中心のcanonical生成元は交換で不変である。
このためFrobenius双対化には追加の指標$\varepsilon$が残る。
内部次数は中心のcanonical shift $-4$と
$\lambda$の次数$-2$から$-2$となる。
weighted projective中心上の相対双対化を行えば
式~\eqref{eq:qa-equivariant-serre-sign}を得る。
このpairingは係数環上の恒等式なので非平坦な基底変更にも保たれる。

式~\eqref{eq:qa-canonical-antiserre-sign}は$t$の反復と
右加群規約の軌道積から従う。
$D$は$\mathcal B_{2n}$に重み$2n$、
$\mathcal D^{-n}$に重み$-2n$で作用し、
$\varepsilon|_D=1$なので、全体は$G=\widetilde G/D$へ降下する。
\end{proof}

\begin{remark}[交換符号の検出と次数線束twist]
\label{rem:qa-antiserre-sign-test}
交換符号は単なる左右規約ではない。
$f_1=f_2=uv$で交換を考えると、
$H^2(\qgr B,\OO(-2))$の生成元はweighted \v{C}ech表示で
$cd/(uvh)$である。
交換はこれを$1/(uv)-cd/(uvh)$へ送り、
$1/(uv)$は\v{C}ech境界なので、この一次元空間への作用は$-1$である。
$H^0(\OO)=k$への作用は$+1$だから、同変Serre双対には
式~\eqref{eq:qa-equivariant-serre-sign}の符号が必要である。

なお、任意の基底線束$L$について
$R^{(L)}_n=R_n\otimes L^n$とすると、
$M_n\mapsto M_n\otimes L^n$はgraded加群圏の同値を与え、
有限生成性とtorsionを保つ。
従って$(\qgr R,\pi R)\simeq
(\qgr R^{(L)},\pi R^{(L)})$である。
符号を省いた環とcanonical反Serre環の違いは$L=\varepsilon$の
このtwistであり、構造対象付きのheartを変更しない。
ただしcanonical Serre関手の同変同定、またはgraded環そのものの
同変同型を述べる場合には符号を省略できない。
\end{remark}

\begin{theorem}[対称有限Morita表示の反Serre角と相対pointed heart]
\label{thm:qa-symmetric-relative-pointed-heart}
前述の二形式商$\mathfrak N^{\rm pres}=[P/G]$上へ降下した
対称有限Morita代数を$\mathsf A$とし、
構造対象を表す零頂点の射影を$P_0$とする。
局所的な基本表示は四項$(B,X_c,X_d,B(1))$であり、
rank-$45$表示では各spinorと最後の項を指定係数束で増幅する。
相対Serre関手から$t_{\mathsf A}=S_{\mathsf A}^{-1}[2]$を定め、
\[
 \mathcal R_{\mathsf A}
 =\bigoplus_{n\ge0}
 \Hom_{\Perf\mathsf A}(P_0,t_{\mathsf A}^{n}P_0)
\]
に軌道積を入れる。この環は
式~\eqref{eq:qa-canonical-antiserre-sign}の$\mathcal R^{\rm can}$に
自然に同型である。
各成分は階数$(2n+1)^2$の局所自由束であり、積と同型は
任意の局所有限型基底変更（非被約を含む）に適合する。
従って全ての有限型affine test $T\to\mathfrak N^{\rm pres}$について
\[
 (\qgr\mathcal R_{\mathsf A,T},\pi\mathcal R_{\mathsf A,T})
 \simeq (\mathcal C_T,\mathcal O_T)
\]
は内在的noncubic族とのpointed exact同値である。
これらの同値はmodule-object基底変更およびfppf降下と可換する。
これは\emph{点曲線$E_f$側の対称Morita avatarの比較}であり、
$C_q$や$\Sigma_2$由来の別avatarのheart、全QP stack、
全ての有限代数の同型groupoidの分類を述べない。
\end{theorem}
\begin{proof}
まず$f_i$の順序と分解を選んだfppf chartへ移る。
論文Bの\cite[\texttt{lem:nc-relative-normal-tilting}]{QuadricModuli2026}
を用いると、normal四項は元の二形式族のcanonical compact圏の
相対compact tilting generatorである。
そこで命題~\ref{prop:qa-symmetric-square-relative-mutation}を適用し、
対称四項へ置き換える。
Morita増幅は同じ加法閉包を生成するので、指定係数束を用いた
対称有限代数$\mathsf A$についても、
\[
 \Perf\mathsf A_T\simeq D(\operatorname{QGr}B_T)^c
\]
という相対Morita同値を得る。
零頂点は$\pi B_T$へ送られる。
ここで有限代数の同型だけを使ったのではなく、元の相対族における
\emph{compact tilting generator}を実際に用いている。

この入力の局所perfectnessは直接にも検出できる。
$X_c=\coker M_c$は$\ell_1$または$\ell_2$の可逆なchartでは
一生成自由な$B$-加群である。
$h$が可逆なchartでは$N=\operatorname{diag}(d,-d)$とすると
$M_cN+NM_c=hI$、$M_c^2=0$である。
従って$h^{-1}M_cN$は次数零冪等元で、
$\im M_c=\ker M_c$は直和因子となる。
これらのchartはweighted projective中心を覆う。
$X_d$も同様であり、以上の全式は任意基底変更に保たれる。
相対strongnessとfullnessはBの同補題の
perfect Hom集中とcompact generatorのcounit判定により成立する。

相対Morita同値は相対Serre双対とYoneda一意性によって自然に可換する。
零頂点の同型とこの自然同型を\emph{同じものの反復}で各次数へ拡張すれば、
すべての$\Hom(P_0,t^nP_0)$と積を同時に同定できる。
前補題により、右辺は正しい交換指標を含む
$\mathcal B_{2n}\mathcal D^{-n}\varepsilon^n$である。
この同定は元の行列表示、評価、Morita増幅から作られるので、
基底変換、因子スカラー、根交換、二形式交換の全てに自然である。
従って$G$への降下cocycleと整合し、$\mathfrak N^{\rm pres}$上へ降下する。

最後にheartを回収する。
局所的には$\mathcal D$を自明化し順序を選べば
$t=(-)_\iota(2)$であり、$(\pi B,t)$はAZ ampleである
\cite[\texttt{prop:nc-root-ample-orbit}]{QuadricModuli2026}。
相対反Serre再構成
\cite[\texttt{prop:relative-h2-antiserre-reconstruction}]{QuadricModuli2026}
により、この全軌道環の$\qgr$は元の構造対象付きheartに一致する。
符号線束が非自明でも前注意の次数線束twistがpointed heartを保つ。
再構成の逆関手は自由graded加群のシフトから軌道対象へ送り、
有限表示の余核へ延ばすものなので、非平坦な通常の基底変更にも可換する。
このため局所的な同値はmodule-object引戻しとfppf降下を保って貼り合う。
有限代数間の導来同値だけからAbel同値を推論する段階はない。
\end{proof}


% !TeX root = ../nc-quadric-avatar.tex
% 日本語 LuaLaTeX。2026-11-23 / noncubic構造頂点の全四位相。
\section{Noncubic avatarの別構造頂点と反Serre生成障害}
\label{sec:qa-spinor-pointing-obstruction}

本節は，非自明二捩れavatarに付随する同じAbel圏で
構造対象を変えた場合を扱う．
定理~\ref{thm:qa-point-all-noncubic-pointed}が認識した
最小十本矢型の一番頂点に対し，零番頂点を構造対象とする組は
強い$\HH$には属さない．
対応するrotated型の零番頂点についても同じ結論を得る．
障害は曲線上の線束の数値ではなく，
ある非零point moduleが構造対象の反Serre軌道から見えないことである．

\subsection{二周期point moduleとspinorの直交性}

$k$を標数零の代数閉体，$f_1,f_2\in k[u,v]_2$を
それぞれsquarefreeな二次形式とする．
二形式の共有根および根対一致を許す．
中央二重平方環を
\[
 B=B_f=k\langle u,v,c,d\rangle/
   (u,v\text{ は中心的},\ c^2=f_1,\ d^2=f_2)
\]
と書き，$h=cd+dc\in Z(B)_2$と置く．
$f_1=-\ell_1\ell_2$, $f_2=-m_1m_2$と因数分解し，
式~\eqref{eq:qa-point-mcm-matrices}の右$B$加群
\[
 X_c=\coker\begin{pmatrix}c&\ell_1\\\ell_2&-c\end{pmatrix},
 \qquad
 X_d=\coker\begin{pmatrix}d&m_1\\m_2&-d\end{pmatrix}
\]
を用いる．両余核は次数零に二つの生成元を持つ．

\begin{lemma}[二周期point moduleが与える直交対象]
\label{lem:qa-spinor-period-two-point-orthogonality}
次数付き右$B$加群$Q_c=\bigoplus_{j\ge0}ke_j$を
$\deg e_j=j$，$Q_cu=Q_cv=0$および
\[
 e_{2j}c=e_{2j+1},\quad e_{2j+1}c=0,
 \qquad
 e_{2j}d=0,\quad e_{2j+1}d=e_{2j+2}
 \qquad(j\ge0)
\]
で定義する．$Q_d$は$c,d$を交換した加群とする．
これらは一元生成のpoint moduleで，$\pi Q_c,\pi Q_d$は非零である．
すべての整数$n$について
\begin{equation}
 \Hom_{\qgr B}(\pi X_c(-2n),\pi Q_c)=0,
 \qquad
 \Hom_{\qgr B}(\pi X_d(-2n),\pi Q_d)=0
 \label{eq:qa-spinor-point-orthogonality}
\end{equation}
が成り立つ．
\end{lemma}
\begin{proof}
$u=v=0$上では$f_1=f_2=0$なので，上の作用は
$c^2=d^2=0$を満たし，$B$加群を定める．
$e_0$から$c,d$を交互に作用させて全基底を得る．
各次数の次元は一であるから$Q_c$は有限生成で無限次元，
したがって$\pi Q_c\ne0$である．また
\[
 e_jh=e_{j+2}\qquad(j\ge0)
\]
であり，$Q_c$は$h$-torsionを持たない．

中心元$h$を次数付きで可逆化し，$B_h=B[h^{-1}]$とする．
任意の有限生成次数付き加群$M$に対する自然写像
\begin{equation}
 \Hom_{\qgr B}(\pi M,\pi Q_c)
 \longrightarrow
 \Hom_{\mathrm{Gr}B_h}(M_h,(Q_c)_h)
 \label{eq:qa-spinor-localization-hom-injection}
\end{equation}
は単射である．実際，$Q_c$には非零有限次元部分加群がない．
従って$qgr$の射は，$M/M'$が有限次元となる次数付き部分加群
$M'\subset M$から$Q_c$への射で表せる．
有限次元加群は$h$の冪で殺されるので$M'_h=M_h$である．
この射が局所化後に零なら，その像は$Q_c$の$h$-torsion部分に
含まれるため，元の射も零である．
この議論はまた\eqref{eq:qa-spinor-localization-hom-injection}の
代表の選択への非依存性を示す．

$(Q_c)_h$には全整数$j$について$e_j$があり，各偶数次数で
\[
 (-)c:k e_{2j}\xrightarrow{\sim}k e_{2j+1}
\]
となる．$X_c(-2n)_h$の次数$2n$の二生成元の像を
$a e_{2n},b e_{2n}$と書く．
$u,v$は標的に零として作用するため，$X_c$の表示行列の関係は
この二つの像にそれぞれ$c$が零として作用することを要求する．
上の単射性から$a=b=0$であり，
局所化後の次数零Homは零である．
\eqref{eq:qa-spinor-localization-hom-injection}より
最初の消滅を得る．$c,d$を交換すれば第二の消滅を得る．
\end{proof}

\begin{theorem}[Noncubic spinor pointingは強い$\HH$に属さない]
\label{thm:qa-spinor-pointing-not-h2}
上の全三型のnoncubic環について
\[
 (\qgr B_f,\pi X_c)\notin\HH,
 \qquad
 (\qgr B_f,\pi X_d)\notin\HH.
\]
より強く，両構造対象の反Serre軌道は
Artin--Zhangの生成条件を満たさない．
\end{theorem}
\begin{proof}
$\iota$を$u,v$を固定して$c,d$を同時に反転する自己同型とする．
補題~\ref{lem:qa-symmetric-equivariant-serre-character}を
体上で用いると，反Serre自己同値は
\[
 t=S^{-1}[2]=(-)_\iota(2)
\]
である．$J=\operatorname{diag}(1,-1)$とすれば
\[
 \iota(M_c)=-J M_c J,
 \qquad \iota(M_d)=-J M_d J.
\]
従って表示の基底変換により$(X_c)_\iota\simeq X_c$および
$(X_d)_\iota\simeq X_d$を得る．特に全整数$n$について
\[
 t^{-n}\pi X_c\simeq\pi X_c(-2n),
 \qquad t^{-n}\pi X_d\simeq\pi X_d(-2n).
\]
前補題の非零対象$\pi Q_c$および$\pi Q_d$は，
対応する反Serre軌道の全対象から直交する．
よって有限個の負の軌道対象の直和からこれらへの全射は存在せず，
生成条件(AZ1)が破れる．
補題~\ref{lem:qa-antiserre-necessary}は強い$\HH$の対象に対して
反Serre軌道の豊富性を要求するので，結論を得る．
\end{proof}

\begin{remark}[障害の位置]
$Q_c,Q_d$は$u=v=0$上に支持される二周期point moduleである．
$cd-dc$は$Q_c$の偶数次数に$h$，奇数次数に$-h$として作用するので，
これらは可換化$B/(cd-dc)$上の加群ではない．
従って点曲線$E_f=\operatorname{Proj}(B/(cd-dc))$上の
線束とpencilの辞書だけでは，上の直交対象を検出できない．
この障害は$\operatorname{coh}\Sigma_2$の負断面による反例とは独立である．
\end{remark}

\subsection{Square，rotated，最小十本矢型の全四位相}

\begin{corollary}[Noncubic由来avatarの構造頂点の分類]
\label{cor:qa-noncubic-avatar-vertex-h2-status}
点曲線$E_f$の二pencilを$F=F_c$, $G=F_d$，
$P=GF^{-1}$とし，$P^2=\OO_{E_f}$，$P\ne\OO_{E_f}$とする．
このsquare avatarから作るsquare型代数を$A_{\rm sq}$，
対応するrotated型を$A_{\rm rsq}=A_{\rm sq}^{[1]}$，
指定mutationによる最小十本矢型を$A_{10}$とする．
各商圏の標準helixをそれぞれ$(E_i)$，$(R_i)$，$(H_i)$と書く．
既存のexact同値の下で，初期四項は次のように同定される：
\begin{equation}
\begin{array}{c|cccc|c}
 &0&1&2&3&\HH\text{に属する位相}\\ \hline
 (E_i)&\pi B&\pi X_d&\pi X_c&\pi B(1)&0,3\\
 (R_i)&\pi X_d&\pi X_c&\pi B(1)&\pi B(2)&2,3\\
 (H_i)&\pi X_d(-1)&\pi B&\pi X_c&\pi B(1)&1,3
\end{array}
\label{eq:qa-avatar-vertex-pointing-table}
\end{equation}
最右欄は添字を四で割った剰余を意味し，すべての整数頂点に適用する．
従ってrotated型の零番頂点と，最小十本矢型の零番頂点を
構造対象とする組はいずれも強い$\HH$に属さない．
\end{corollary}
\begin{proof}
定理~\ref{thm:qa-point-all-noncubic-pointed}が与える
\emph{頂点と全合成を保つ}finite End同型およびexact同値は
\[
 (H_1,H_2,H_3,H_4)
 \longmapsto(\pi B,\pi X_c,\pi B(1),\pi X_d(1))
\]
である．$H_4=tH_0$と，前定理の$(X_d)_\iota\simeq X_d$から
$H_0$の像は$\pi X_d(-1)$となる．

補題~\ref{lem:qa-as10-all-index-mutual-inverse}の逆mutationでは
$E_0=H_1$, $E_2=H_2$, $E_3=H_3$であり，$E_1$は余評価列
\[
 0\longrightarrow H_0\longrightarrow
 \Hom(H_0,H_1)^\vee\otimes H_1\longrightarrow E_1\longrightarrow0
\]
の余核である．$M_d^2=0$からの線形行列分解は
\[
 0\longrightarrow X_d(-1)\longrightarrow B^{\oplus2}
 \longrightarrow X_d\longrightarrow0
\]
を与える（式~\eqref{eq:qa-symmetric-square-spinor-evaluation}を一度shiftする）．
この包含の二成分は$\Hom(\pi X_d(-1),\pi B)$の基底なので，
前記余評価と基底変換を除き一致する．
従って$E_1$の像は$\pi X_d$であり，squareの行を得る．

定理~\ref{thm:qa-as-avatar-bijection}の証明における
添字shiftは$R_i=E_{i+1}$を与える．
また$E_4=tE_0$の像は$\pi B(2)$である．
これでrotatedの行も得る．
この同定は曲線上の共通線束tensorだけから結論したものではなく，
有限End同型，Serre完成，exact mutationおよび実際の行列分解を用いている．

$(\qgr B,\pi B)$は強い$\HH$の対象であり，
通常次数shift$(j)$はexact自己同値なので
$(\qgr B,\pi B(j))$も強い$\HH$に属する．
ここでは$(1)$自体を適合Serre平方根と主張する必要はなく，
既存の適合根をこのexact同値で共役すればよい．
一方，前定理と次数shiftにより$\pi X_c(j)$，$\pi X_d(j)$を
構造対象とする組は強い$\HH$に属さない．
反Serre関手$t$もexact自己同値なので，四周期を通してこの結論が保たれる．
\end{proof}

\begin{remark}[Spectral avatarへの帰結と主張範囲]
滑らかなspectral avatarの指定square実現にこの系を適用すると，
mutation後の一番頂点は
定理~\ref{thm:qa-spectral-avatar-pointed-involution}の
$\mathcal H(f^\vee)$を与えるが，零番頂点および対応する
rotated零番頂点は強い$\HH$の外にある．
従ってこれらを別のFricke像や元の$\mathcal H(f)$と
pointed exact同値にすることはできない．
本節はnoncubic root-pair由来の非自明二捩れavatarを扱い，
一般の$P$を持つ全rotated型avatarの構造頂点を分類するものではない．
\end{remark}

% !TeX root = ../nc-quadric-avatar.tex
% 2026-11-23 / scratch draft: vertex-ideal criterion and all classified avatars
\section{頂点idealによる反Serre生成判定と一般avatar}
\label{sec:qa-general-vertex-az}

前節の二捩れspinorによる障害を，一般のavatarへ無条件に
外挿することはできない．本節では先にAS helixの頂点idealに
よる生成判定を証明する．最小十本矢型では二本の追加矢が作る
無限pathを用いることで，二捩れを仮定せずに偶数頂点を排除できる．
その後，非被約曲線を許すsquare avatarについて，二つのmutation壁と
cubicのpointed逆構成を用いて全四位相を判定する．

\subsection{十分負の成分比較と頂点ideal}

$A$を本稿のsquare，rotated，cubic，または最小十本矢型の
左右Noetherian AS $\mathbb Z$-代数とする．
$P_i=e_iA$, $E_i=\pi P_i$，
$t=S^{-1}[2]$と書き，$tE_i\simeq E_{i+4}$の規約を用いる．
右加群の$j$成分は$M_j=Me_j$であり，
$(P_a)_j=A_{aj}$は$j\le a$に支持される．

\begin{lemma}[有限射影解消からの遠方成分比較]
\label{lem:qa-as-far-component-comparison}
任意の有限生成次数付き右$A$加群$M$に対し，$j\ll0$なら
自然な写像は同型
\[
 M_j\xrightarrow{\sim}\Hom_{\qgr A}(E_j,\pi M)
\]
を与え，さらに$\Ext^r_{\qgr A}(E_j,\pi M)=0$ ($r>0$)である．
必要な$j$の上限は$M$に依存してよい．
\end{lemma}
\begin{proof}
$A$は左右Noetherianであり，左右の単純加群は長さ三のAS解消を持つ．
従って$M$には有限個の標準射影$P_a$からなる長さ三以下の解消がある．
これは命題~\ref{prop:qa-as-avatar-coherent-morita}の証明と同じである：
右AS解消を双対化すると，AS--Gorenstein条件から全ての左単純の
長さ三の射影解消を得る．有限生成右加群は添字について上に有界で，
graded Nakayamaによる極小射影被覆を持つ．
有限生成syzygyの極小射影被覆を取ると，第四項が非零なら
いずれかの左単純との$\Tor_{>3}$で検出され，
左単純の長さ三の解消に矛盾する．
AS10についても定理~\ref{thm:qa-as10-global-bijection}および
命題~\ref{prop:qa-as10-noetherian-qgr}により同じ入力が成立する．

この有限解消を$P_\bullet\to M$とする．
補題~\ref{lem:qa-as-serre-unique-recovery}，
補題~\ref{lem:qa-as10-hilbert-foundation}および全整数mutationの
前向きExt消滅より，$j$を解消に現れる全$a$より十分小さく取れば
\[
 \RHom(E_j,\pi P_a)=A_{aj}[0]
\]
であり，同型は全写像と両立する．
従って$\RHom(E_j,\pi P_\bullet)$は元の射影解消の$j$成分そのもので，
そのコホモロジーは次数零の$M_j$だけである．
$\pi$は完全なので，これは$\RHom(E_j,\pi M)$を計算する．
ここでは任意の個別頂点の豊富性や，一般的な$\chi$定理を先に仮定していない．
\end{proof}

\begin{theorem}[頂点idealによる反Serre豊富性の必要十分条件]
\label{thm:qa-vertex-ideal-az-criterion}
$i\in\mathbb Z$を固定し，
\[
 I_i=\sum_{q\in\mathbb Z}Ae_{i+4q}A,
 \qquad C_i=A/I_i
\]
と置く．次は同値である：
\begin{enumerate}
\item $(E_i,t)$はArtin--Zhang ampleである．
\item すべての$j\in\mathbb Z$について右$C_i$加群$e_jC_i$は
      有限次元である．
\end{enumerate}
特に，削除後に無限次元の標準右射影が残る位相は
強い$\HH$の構造頂点にならない．
\end{theorem}
\begin{proof}
(2)を仮定する．任意の$\pi M\in\qgr A$と任意の下限$N$を取る．
$M$を十分負の添字まで切断しても$\pi M$は変わらないので，
$M_j=0$ ($j>i-4N$)としてよい．
この切断もNoether性により有限生成である．
選択したresidueの成分から生成される部分加群を
\[
 L=MI_i=\sum_q M_{i+4q}A
\]
と置く．$M/L$は有限生成$C_i$加群なので，(2)より有限次元である．
$L$はNoetherianな$M$の部分加群であるから有限生成であり，
有限個の$M_{i-4n}$ ($n\ge N$)の元で生成できる．
よって有限直和からの全射
\[
 \bigoplus_{\alpha}t^{-n_\alpha}E_i\twoheadrightarrow\pi M,
 \qquad n_\alpha\ge N
\]
を得る．これが任意に深い負の軌道による(AZ1)である．
任意の全射のkernelに前補題を適用すれば，十分大きい$n$で
$\Ext^1(t^{-n}E_i,\ker)=0$となり，(AZ2)も得る．

逆に$e_jC_i$が無限次元となる$j$が存在するとする．
$M=e_jC_i$は$P_j$の商なので有限生成で，$\pi M\ne0$である．
すべての$q$で$M_{i+4q}=0$だから，前補題によって
\[
 \Hom(t^{-n}E_i,\pi M)=0\qquad(n\gg0).
\]
十分大きい$N$を選べば，非零対象$t^N\pi M$への
全ての$\Hom(t^{-n}E_i,-)$ ($n\ge0$)が零になる．
これは(AZ1)に反する．最後の主張は
補題~\ref{lem:qa-antiserre-necessary}による．
\end{proof}

\subsection{最小十本矢型の追加矢が与える一般の障害}

\begin{proposition}[最小十本矢型では偶数位相が生成しない]
\label{prop:qa-as10-general-even-obstruction}
任意の$A\in\mathrm{AS}_{10}^{\mathrm{min}}$について，
$(E_i,t)$がAZ ampleであることと$i$が奇数であることは同値である．
特に全ての偶数$i$について
\[
 (\qgr A,E_i)\notin\HH.
\]
この主張に$P$の二捩れ性，avatar曲線の滑らかさは不要である．
\end{proposition}
\begin{proof}
定理~\ref{thm:qa-as10-global-bijection}と
式~\eqref{eq:qa-as10-reduced-potential}により，四周期表示のpotentialは
\[
 W=yq+px+H,
\]
ただし$x:1\to3$, $y:3\to1$は内部次数二，
$p,q$はcycle矢二本の線形結合，$H$はcycle矢四本の線形結合である．
全整数への展開では$x,y$が奇数頂点を二つずつ進める．

全偶数頂点を削除すると，八本のcycle矢はすべて零になり，
奇数頂点を結ぶ$x,y$だけが残る．
全てのcyclic derivativeは，削除されたcycle矢を少なくとも一本含む：
$\partial_xW=p$と$\partial_yW=q$は二本，
cycle矢による微分は一本または三本を含む．
従って残る$x,y$には関係がなく，商は奇数頂点上の一本の有向鎖の
道代数である．その標準右射影$M=e_1C$は，
すべての奇数$j\le1$で$M_j\simeq k$，偶数$j$で$M_j=0$となる
無限次元一元生成加群である．
各偶数位相だけを削除する商はさらに$C$へ全射するので，
前定理の有限次元条件を満たさない．

奇数の一つの位相を削除すると，追加矢$x,y$はともに消え，
残るcycle矢は各周期で長さ高々二の有向鎖になる．
従って削除後の各標準右射影は有限次元であり，前定理を適用できる．
\end{proof}

\begin{remark}[非二捩れでの障害と二周期point module]
この無限path加群はAS10の追加矢から構成され，
中央二重平方表示や中心元$cd+dc$を必要としない．
従って二捩れnoncubicの場合の二周期point moduleによる障害には，
一般$P$でも有効なAS10側の表現がある．
一方，一般avatarに中央二重平方環$B_f$や
通常の次数一生成point schemeが付随すると主張しているのではない．
\end{remark}

\subsection{一般曲線型avatarの全四位相}

以下，$(D,F,G,P)$を本稿のsquare曲線型avatarとする．
$D$は非被約でもよく，$F,G$は指定された$(2,2)$埋め込みの
相異なる二ruling，$P\in\Pic^0(D)\setminus\{\OO_D\}$である．
二つの壁を
\[
 D_F=\{G\simeq FP\},\qquad D_G=\{GP\simeq F\}
\]
と書く．ここで$D_F,D_G$は曲線そのものでなくパラメータ上の条件である．

\begin{theorem}[一般square曲線型avatarの構造頂点]
\label{thm:qa-general-square-pointings}
square AS helixを$(E_i)$とすると，各位相について
強い$\HH$への所属と反Serre軌道のAZ豊富性は同値であり，
次の表で与えられる：
\[
\begin{array}{c|c|c}
 i\bmod4& (E_i,t)\text{ がAZ ample}&(\qgr A,E_i)\in\HH\\ \hline
 0,3&\text{常に成立}&\text{常に成立}\\
 1&GP\not\simeq F&GP\not\simeq F\\
 2&G\not\simeq FP&G\not\simeq FP
\end{array}
\]
\end{theorem}
\begin{proof}
第一mutationの全整数exact比較は
\[
 H_0=\ker(E_0^{\oplus2}\to E_1),\quad
 H_1=E_0,\quad H_2=E_2,\quad H_3=E_3
\]
である．$D_F$上では最小十本矢型なので，前命題によって
$E_2=H_2$の反Serre生成条件が破れる．
$F,G$を交換して第二mutationを行えば，$D_G$上で
$E_1$について同じ結論を得る．

$D_F$の外では定理~\ref{thm:qa-as10-forward-as}により
第一mutationはcubic型AS代数である．この定理は曲線の被約性を要求しない．
命題~\ref{prop:qa-as10-noetherian-qgr}のexact同値は
上記の$H_1=E_0$, $H_2=E_2$, $H_3=E_3$を保つ．
論文A~\cite[\texttt{thm:all-admissible-inverse}]{QuadricIntrinsic2026}は，
\emph{任意の}三次三次元AS正則$\mathbb Z$-代数の標準零頂点を
構造対象とする組が強い$\HH$に属することを述べる．
これを得られたcubic代数および各添字shiftに適用すると，
全$H_i$を構造対象とする組が強い$\HH$に属する．
特に$E_0,E_2,E_3$の所属を得る．
ここで同じrootやcurve上のtensorがnc heartへ延びるとは仮定しない．

第二mutationを用いると，$D_G$の外で$E_0,E_1,E_3$の所属を得る．
二壁の交点では$P^2=\OO_D$, $P\ne\OO_D$, $G=FP$である．
補題~\ref{lem:qa-as10-double-wall-reduced-geometries}により，
この場合$D$はsmooth/nodal/$I_2$の三型に限られる．
全三型についての既存のnoncubic認識と
系~\ref{cor:qa-noncubic-avatar-vertex-h2-status}により
$E_0,E_3$が強い$\HH$に属する．
これで表の全ての肯定・否定を得た．
肯定側のAZ豊富性は補題~\ref{lem:qa-antiserre-necessary}から従い，
否定側はすでに生成障害そのものとして証明されている．
四周期はexact自己同値$t$なので，全整数へ結論が移る．
\end{proof}

\begin{corollary}[Rotated零頂点と一般最小十本矢型のpointing]
\label{cor:qa-general-rotated-as10-pointings}
曲線型rotated avatarを$(D,F_r,H_r,P)$，
$C=H_rF_r^{-2}$とし，標準helixを$(R_i)$とする．
強い$\HH$への所属は
\[
\begin{array}{c|c}
 i\bmod4&(\qgr A_{\rm rsq},R_i)\in\HH\\ \hline
 0&C\not\simeq\OO_D\\
 1&C\not\simeq P^2\\
 2,3&\text{常に成立}
\end{array}
\]
である．とくに零頂点はHilbert曲面が$\Sigma_0$の場合に属し，
$\Sigma_2$の場合には属さない．
また，全ての最小十本矢型avatarでは
構造頂点$H_i$が強い$\HH$に属することと$i$が奇数であることは同値である．
\end{corollary}
\begin{proof}
squareとの辞書は
$F_s=F_r$, $G_s=H_rF_r^{-1}P^{-1}$，
また$R_i=E_{i+1}$である．従って
\[
 G_sP\simeq F_s\iff C\simeq\OO_D,
 \qquad
 G_s\simeq F_sP\iff C\simeq P^2.
\]
前定理に代入すればrotatedの表を得る．
Hilbert曲面の判定は
定理~\ref{thm:qa-core-avatar-square-rotated}による．
この結論は$\coh\Sigma_2$との同一視から得たものではなく，
AS10の無限pathによる生成障害から得ている．

最小十本矢型では$D_F$上にあり，$H_1=E_0$，$H_3=E_3$なので
奇数位相は前定理によって強い$\HH$に属する．
偶数位相は命題~\ref{prop:qa-as10-general-even-obstruction}により除外される．
\end{proof}

\begin{proposition}[古典型の全位相]
\label{prop:qa-classical-all-vertex-pointings}
squareおよびrotatedの各一つの古典型では，
すべての標準頂点を構造対象とする組が強い$\HH$に属する．
\end{proposition}
\begin{proof}
命題~\ref{prop:qa-as-extraction}と
定理~\ref{thm:qa-as-avatar-bijection}の古典型は
$\PP^1\times\PP^1$であり，標準四項はそれぞれ
\[
 (\OO,\OO(1,0),\OO(0,1),\OO(1,1)),\qquad
 (\OO,\OO(-1,1),\OO(0,1),\OO(1,2))
\]
である．全整数helixも線束からなる．各線束とのtensorはexact自己同値で，
$(\coh(\PP^1\times\PP^1),\OO)$は通常のSegre表示による
強い$\HH$の対象だから，結論を得る．
古典型を曲線型の許容条件$P\ne\OO_D$へ形式的に代入したものではない．
\end{proof}

\begin{remark}[範囲と残る比較]
以上は本稿で分類した体上のsquare／rotated／最小十本矢型の
全標準構造頂点を扱う．任意の非標準exceptional対象を構造対象とする問題，
一般avatar族の相対pointed認識，任意基底上のAS/QP群oidの同値は別の課題である．
また，一般のpointed圏について「反Serre ampleなら強い$\HH$」と
主張したのではなく，表の肯定側では実際のcubic/noncubic認識を用いた．
\end{remark}

% !TeX root = ../nc-quadric-avatar.tex
% 2026-11-23 / v1.0
\section{相対avatarの構造頂点と二つの行列式開}
\label{sec:qa-relative-vertex-determinants}

本節は，体上の全構造頂点判定を，任意の非被約底変更と両立する
有限表示の開条件へ移す．まず相対曲線と実際の線束だけから
二つの乗法行列を作る．この幾何適合開の構成と，実際の相対
pointed $\HH$族の認識は区別し，後者に必要な入力を末尾に明記する．

\subsection{相対乗法束と全基底変更}

$T$を$k$-schemeとし，$f:\mathcal D\to T$をproper，flat，
finite presentationな射とする．$\mathcal F,\mathcal G,\mathcal P$を
$\mathcal D$上の可逆層とし，各幾何点$\bar t\to T$において
\[
 (\mathcal D_{\bar t},\mathcal F_{\bar t},
   \mathcal G_{\bar t},\mathcal P_{\bar t})
\]
が本稿のsquare曲線型avatarであると仮定する．
曲線の被約性も$T$の被約性も仮定しない．
とくに$\mathcal P_{\bar t}\in\Pic^0(\mathcal D_{\bar t})
\setminus\{\OO\}$である．

\begin{lemma}[二つのrank四乗法束]
\label{lem:qa-relative-pencil-multiplication-bundles}
$\mathcal V(L)=f_*L$と書く．
\[
\mathcal V(\mathcal F),\quad\mathcal V(\mathcal G),\quad
\mathcal V(\mathcal F\mathcal P^{-1}),\quad
\mathcal V(\mathcal G\mathcal P^{-1})
\]
は階数二，$\mathcal W=\mathcal V(\mathcal F\mathcal G\mathcal P^{-1})$
は階数四の有限局所自由層である．
これらの高次直像は零で，形成は任意の$T'\to T$と可換する．
従って切断の乗法はrank四束間の二つの写像
\begin{align}
 \mu_G:&\ \mathcal V(\mathcal F)\otimes
          \mathcal V(\mathcal G\mathcal P^{-1})\longrightarrow\mathcal W,
          \label{eq:qa-relative-mu-g}\\
 \mu_F:&\ \mathcal V(\mathcal G)\otimes
          \mathcal V(\mathcal F\mathcal P^{-1})\longrightarrow\mathcal W
          \label{eq:qa-relative-mu-f}
\end{align}
を定め，両写像も任意の底変更と可換する．
全幾何fiberにおける階数は三または四である．
\end{lemma}
\begin{proof}
補題~\ref{lem:qa-core-curve-vanishing}より，各fiberで対象線束の
$H^1$は零，$H^0$の次元は記載した二または四である．
各$Rf_*L$はperfectで，任意の導来底変更と可換する
（\href{https://stacks.math.columbia.edu/tag/0B91}
{Stacks Project, Tag 0B91, Lemma 36.30.4}）．
その全幾何fiberが次数零に集中するので，有限自由複体の可逆minorを
順に消去すれば，$Rf_*L$自体が記載した階数の局所自由層となる．
従って導来底変更は通常の底変更となる．
これは非被約底上でコホモロジー次元が一定という理由だけから
局所自由性を推論するものではない．
乗法の自然性は導来底変更およびprojection formulaと両立し，
階数三または四という主張は
補題~\ref{lem:qa-as10-rank-wall-pencils}から従う．
\end{proof}

\subsection{幾何適合開とschemeとしての悪いlocus}

\begin{theorem}[標準構造頂点の有限表示開部分関手]
\label{thm:qa-relative-vertex-determinant-open}
体上のAS/avatar対応で各幾何fiberから得られる$i$番標準構造頂点について，
$T$上のschemeの部分関手$\mathscr U_i$を，すべての幾何fiberの
指定pointingが強い$\HH$に属する底変更$T'\to T$の全体とする．
これは次の有限表示open immersion $U_i\hookrightarrow T$で表される：
\begin{equation}
 U_0=U_3=T,\qquad
 U_1=D(\det\mu_F),\qquad U_2=D(\det\mu_G).
 \label{eq:qa-relative-square-vertex-open}
\end{equation}
行列式はそれぞれ
$\det\mathcal W\otimes\det(\operatorname{dom}\mu)^{-1}$の切断である．
すべての$T'\to T$について
\[
 U_i(\mathcal D_{T'},\mathcal F_{T'},\mathcal G_{T'},\mathcal P_{T'})
 =U_i\times_T T'
\]
がschemeとして成立する．

悪い閉部分schemeを
\[
 Z_F=V\bigl(\operatorname{Fitt}_0\coker\mu_G\bigr),\qquad
 Z_G=V\bigl(\operatorname{Fitt}_0\coker\mu_F\bigr)
\]
と置けば，$U_1=T\setminus Z_G$，$U_2=T\setminus Z_F$である．
$Z_F,Z_G$も有限表示であり，任意の底変更と可換する．
各写像は局所的に
\begin{equation}
 \operatorname{diag}(1,1,1,\delta)
 \label{eq:qa-relative-wall-schur-normal-form}
\end{equation}
と同値であり，対応する悪い閉部分は$(\delta)$で定義される．
その余核は悪い閉部分上の線束のpushforwardである．
\end{theorem}
\begin{proof}
前補題により乗法束と行列式はすべての底変更と可換する．
定理~\ref{thm:qa-general-square-pointings}によって，幾何fiberの
一番頂点は$\mu_F$が可逆の場合に，二番頂点は$\mu_G$が可逆の場合に
ちょうど強い$\HH$へ属する．
線束の切断が全幾何点で非零であることは，その切断が可逆であることと
同値であり，非被約底でも成り立つ．
従って各部分関手は記載した行列式の可逆開で表される．
線束を局所自明化すればprincipal openであるため，
open immersionは有限表示である．

rank四束間の写像の第零Fitting idealは行列式で生成される．
Fitting idealの形成は任意の環準同型と可換する
（\href{https://stacks.math.columbia.edu/tag/07Z6}
{Stacks Project, Tag 07Z6, Lemma 15.8.4(3)--(4)}）．
全幾何fiberで階数が三以上なので，三次minorの可逆開が$T$を覆う．
その上で行列の消去を行えば
\eqref{eq:qa-relative-wall-schur-normal-form}を得る．
余核は局所的に$\OO_T/(\delta)$であり，これらの局所表示は
$V(\delta)$上の可逆層として貼り合う．
ここで$\delta$が非零因子であるとは仮定していない．
\end{proof}

\begin{corollary}[Rotated族と最小十本矢族の開部分]
\label{cor:qa-relative-rotated-as10-open}
rotated族から
$\mathcal F_s=\mathcal F_r$，
$\mathcal G_s=\mathcal H_r\mathcal F_r^{-1}\mathcal P^{-1}$
によって二乗法を作ると，rotatedの各位相の幾何適合開は
\[
 U^{\rm rsq}_0=D(\det\mu_F),\quad
 U^{\rm rsq}_1=D(\det\mu_G),\quad
 U^{\rm rsq}_2=U^{\rm rsq}_3=T.
\]
全幾何fiberが最小十本矢型である族では，偶数位相の幾何適合開は空，
奇数位相の幾何適合開は$T$である．
これらはいずれも任意の底変更と可換する．
\end{corollary}
\begin{proof}
前者は$R_i=E_{i+1}$と相対line bundleの式に，後者は
系~\ref{cor:qa-general-rotated-as10-pointings}に従う．
rotatedから得た線束の各fiberはsquare avatarなので，
前補題と前定理をそのまま適用できる．
\end{proof}

\subsection{普遍楕円パラメータ上のCartier壁}

\begin{proposition}[普遍壁はCartier，任意引戻しではFitting閉部分]
\label{prop:qa-relative-universal-cartier-walls}
$E\to S$を標数零の滑らかな楕円曲線族，$F$を相対次数二線束とし，
$J=\Pic^0(E/S)$と置く．
\[
 V=(J\setminus0)\times_S(J\setminus0),\qquad
 (P,\eta)\in V,\quad G=F\eta
\]
上で，必要ならPoincare束を持つfppf chartへ移って乗法束を作る．
このときscheme-theoreticに
\[
 Z_F=\{\eta=P\},\qquad Z_G=\{\eta=P^{-1}\}
\]
であり，両者は$V$上の有効Cartier因子である．
両者の交わりは$J[2]\setminus0$で，$S$上有限etale次数三である．
相対余次元二で横断的に交わる．

任意のavatar族への引戻しについては，悪い閉部分は上記Fitting理想の
引戻しであるが，有効Cartier因子であるとは限らない．
引き戻した局所行列式が非零因子である場合に限り，Cartierである．
\end{proposition}
\begin{proof}
次数二pencil $A$と$B$について，相対評価列は
\[
 0\to f^*\det\mathcal V(A)\otimes BA^{-1}
 \to f^*\mathcal V(A)\otimes B\to AB\to0
\]
を与える．正次数項の高次直像消滅から
\[
 Rf_*(BA^{-1})\otimes\det\mathcal V(A)
 \simeq
 [\mathcal V(A)\otimes\mathcal V(B)\xrightarrow{\mu}\mathcal V(AB)]
\]
を得る．右辺の二項は次数零，一に置く．
従って行列式の消滅は，次数零線束$BA^{-1}$のcohomology jumpを
そのscheme構造も含めて記録する．

普遍$J$の零断面上では対象線束は基底からの線束を除き自明であり，
kernelとcokernelは各一の線束となる．
ここでは$Rf_*\OO_E$のunitからの分裂と
$f_*\OO_E=\OO_S$を用いる
（\href{https://stacks.math.columbia.edu/tag/0E6B}
{Stacks Project, Tag 0E6B, Lemma 75.26.7}）．
したがって行列式は零断面上でscheme-theoreticに零となる．
各幾何fiberで零断面に垂直な非零接方向
$\xi\in H^1(E_s,\OO_{E_s})$を取ると，
対応する一次変形$Q_\epsilon$の接続写像は
\[
 H^0(\OO_{E_s})\longrightarrow H^1(\OO_{E_s}),\qquad1\longmapsto\xi
\]
である．従って残る一行一列の微分は非零で，行列式は零断面で一次の零点を持つ．
零断面の局所Cartier方程式を$z$とすれば，行列式は$zu$と書け，
このfiber上の計算から$u$は単元である．
ゆえにこの同定は非被約な$S$上でもscheme-theoreticに成立する．

$\mu_G$では$BA^{-1}=\eta P^{-1}$，
$\mu_F$では$BA^{-1}=\eta^{-1}P^{-1}$であるから，
記載した二つのgraphを得る．
滑らかな相対一次元群$J$におけるdiagonalおよび反転のgraphは
有効Cartierである．交点では$P^2=\OO$，かつ$P\ne\OO$であり，
$2$は可逆なので$J[2]\setminus0$は有限etale次数三である．
二つの相対法方向は$d\eta-dP$，$d\eta+dP$であり，
行列式$2$が可逆であることから横断性が従う．
Poincare束の基底からのline bundle twistは行列式の局所方程式を
単元倍するだけなので，これらの閉部分はchartから降下する．
最後の主張は局所正規形の$\delta$がCartier方程式となる条件そのものである．
\end{proof}

\begin{example}[非被約底での階数三条件]
普遍壁に横断する一次変形を
$T=\Spec k[\epsilon]/(\epsilon^2)$へ引き戻すと，局所乗法行列は
\[
 \operatorname{diag}(1,1,1,\epsilon)
\]
となる．全幾何fiberの階数は三だが，行列式は$T$上恒等的に零ではない．
悪いFitting閉部分は$V(\epsilon)$で，有効Cartier因子ではない．
一方，対応する幾何適合開は空である．
したがって幾何fiberの階数だけからscheme-theoreticな壁への因数分解，
kernel束の局所自由性，またはCartier性を推論しない．
\end{example}

\subsection{相対pointed認識へ進むための有限入力}

\begin{lemma}[既存H2開への相対因数分解]
\label{lem:qa-relative-foundation-open-factorization}
$T$を本稿の相対モジュライsite内の基底とする．
実際の有限局所自由foundation代数族と全積から
射$\lambda:T\to\mathfrak F_0$が与えられ，
その全幾何fiberが実現可能開$\mathfrak F_{\HH}$に入るとする．
このとき$\lambda$は一意に$\mathfrak F_{\HH}$を経由し，
定理~\ref{thm:qa-intrinsic-h2-open-in-foundation-stack}により
実際の相対pointed $\HH$族を定める．
この構成は任意のsite内の底変更と同型射に適合する．
\end{lemma}
\begin{proof}
開部分への因数分解は位相空間の像で判定できるので，
非被約底についても$\lambda$は$\mathfrak F_{\HH}$を経由する．
同定$\mathfrak F_{\HH}\simeq\mathfrak H_2^{\rm rul}$を合成すればよい．
開埋込とこの同値の自然性から底変更・同型射への適合性が従う．
\end{proof}

この補題で必要なのは，四頂点の全Hom束と全合成を持つ
\emph{実際の}相対foundation族，$\mathfrak F_0$の条件，
および幾何fiberにおけるnormal foundationとしての認識である．
行列式開の上にあるという事実だけからこれらを省略しない．
さらに，こうして得た$\HH$族と，別に与えた相対AS代数の$qgr$ heartを
同一視するには，構造対象を保つ相対tilting比較，Serre完成と
相対AZ再構成を確認する必要がある．
本節の二行列式だけで全AS/QP stackを$\HHstack$と同定したとは主張しない．

% !TeX root = ../nc-quadric-avatar.tex
% 2026-11-23 / v1.0

\section{一般square族の相対normal窓と二つの構造対象}
\label{sec:qa-relative-avatar-pointed-recognition}

本節では有限型基底上の実際の avatar 族から，相対 $\HHstack$ の対象を作る．
構成には，定理~\ref{thm:qa-intrinsic-h2-open-in-foundation-stack} の
\emph{相対的な}開埋込
\[
 \Xi:\mathfrak H_2^{\rm rul}\xrightarrow{\sim}\mathfrak F_{\HH}
 \subset\mathfrak F_0
\]
を既知の入力として使う．幾何 fiber の $H_2$ 性だけから
相対 Serre 平方根や相対 AS4 被覆の存在を推論するのではない．

\subsection{基底と実際の有限代数族}

$k$ は標数零の代数閉体とする．$S$ を $k$ 上局所有限型のスキームとし，
$p:D\to S$ を proper flat finite presentation な曲線族とする．
fppf 局所で $D\subset\PP^1_S\times_S\PP^1_S$ は
相対 $(2,2)$ Cartier 因子で，$F,G$ は二つの ruling の制限であるとする．
$P$ は $D$ 上の可逆層で，各幾何 fiber において
$(D_s,F_s,G_s,P_s)$ が本稿の square 曲線型 avatar であると仮定する．
特に $P_s\in\Pic^0(D_s)\setminus\{\OO_{D_s}\}$ である．
曲線の非被約性と可約性を許す．
古典型で反標準データが曲線から曲面へ変わる族は，このflat曲線族の範囲に含めない．
以下の構成はこのデータの同型に自然であるから，この定義の
avatar の fppf groupoid に対しても適用できる．
この groupoid を任意の AS/QP 表示 groupoid と同一視することはしない．

\[
 (L_0,L_1,L_2,L_3)=(P,F,G,FG),\qquad
 \mathcal H_{ij}=p_*(L_jL_i^{-1})\quad(i<j)
\]
と置く．補題~\ref{lem:qa-core-curve-vanishing} と
proper flat な射の perfect pushforward と任意基底変更
（\href{https://stacks.math.columbia.edu/tag/0B91}{Stacks Project, Tag 0B91}）
により，$\mathcal H_{ij}$ は所定階数の有限局所自由層であり，
$\mathcal H_{12}=0$，
\[
 (\operatorname{rk}\mathcal H_{01},\operatorname{rk}\mathcal H_{02},
  \operatorname{rk}\mathcal H_{03},\operatorname{rk}\mathcal H_{13},
  \operatorname{rk}\mathcal H_{23})=(2,2,4,2,2).
\]
各束と切断の積は非平坦なものを含む任意基底変更と可換する．
実際，導来 pushforward の各幾何 fiber が次数零に集中するので，
perfect 複体の fiber 判定を用いている．

右 Yoneda の規約で
\[
 e_j\Lambda e_i=\mathcal H_{ij}\ (i<j),\qquad
 e_i\Lambda e_i=\OO_S,\qquad e_i\Lambda e_j=0\ (i<j)
\]
と定め，積を切断乗法とする．これは階数十六の有限局所自由な
directed 代数である．この定義は $\Perf(D)$ の full derived End
ではない．特に対角の $H^1(\OO_D)$ を代数へ加えない．
$E_i=e_i\Lambda\in\Perf(\Lambda)$ とすれば
\[
 \RHom(E_i,E_j)=\mathcal H_{ij}[0]\quad(i<j)
\]
であり，逆向きの Hom は零である．
有限 normalized bar 解消によって $\Lambda$ は相対 smooth proper である．

\begin{proposition}[全基底上のnormal候補窓]
\label{prop:qa-relative-avatar-normal-window}
$V=\mathcal H_{13}=p_*G$ とし，実際の relative cone
\[
 Z_1=\operatorname{Cone}
  \bigl(E_1\xrightarrow{\rm coev}V^\vee\otimes E_3\bigr)
\]
を取る．すると
\[
 \mathcal N_F=(E_0,E_2,E_3,Z_1)
\]
は相対 full strong exceptional foundation であり，その通常の End
代数 $\Gamma_F$ は階数二十四である．この窓は
$\mathfrak F_0$ の点 $\lambda_F:S\to\mathfrak F_0$ を与える．
窓，End 代数，全ての合成，対応する compact Morita 同値と
$\lambda_F$ は任意の基底変更に可換する．

同様に $F,G$ と $E_1,E_2$ を交換して
\[
 Z_2=\operatorname{Cone}
 \bigl(E_2\to\mathcal H_{23}^\vee\otimes E_3\bigr),\qquad
 \mathcal N_G=(E_0,E_1,E_3,Z_2)
\]
を得る．論文Bの相対 ruling 交換 $\tau$ に対し
$\tau\lambda_F\simeq\lambda_G$ が coherent に成立する．
\end{proposition}
\begin{proof}
四項の間で新しく調べる必要がある Hom は
\[
 \RHom(E_0,Z_1)
 =\operatorname{Cone}\bigl(
   \alpha:\mathcal H_{01}\to V^\vee\otimes\mathcal H_{03}\bigr)
\]
である．$\alpha$ は $G$ の切断乗法の随伴である．
各幾何 fiber 上で，$a\in H^0(FP^{-1})$ がその kernel に入るなら，
$aH^0(G)=0$ である．$G$ は大域生成なので，局所的な生成切断で
縮約すれば $a=0$ となる．これは非被約な $D_s$ 上でも成立する．
従って $\alpha$ は基底上の部分束包含であり，
\[
 C=\operatorname{coker}\alpha
\]
は階数六の有限局所自由層である．
余核とその完全列は任意基底変更に保たれる．

残る Hom は cone の三角から直接
\[
 \RHom(E_2,Z_1)=V^\vee\otimes\mathcal H_{23},\qquad
 \RHom(E_3,Z_1)=V^\vee
\]
となり，いずれも次数零である．逆向きの Hom は零で，
$Z_1$ の自己 Hom は $\OO_S[0]$ である．
最後の二主張は相対 exceptional pair $(E_1,E_3)$ の mutation，
または同じ三角の Hom 列で確認できる．
$E_1$ はこの三角から回復できるので四項は full である．
従って通常 End の階数表は
\[
 \begin{pmatrix}1&2&4&6\\0&1&2&4\\0&0&1&2\\0&0&0&1\end{pmatrix}
\]
となる．

論文Bの二写像表示では
\[
 (W_1,W_2,W_3;P_{\rm old},Q_{\rm old})
 =(\mathcal H_{02},\mathcal H_{03},C;
   \mathcal H_{23},V^\vee).
\]
$G_{\rm old}:V^\vee\otimes\mathcal H_{03}\to C$ は定義上全射である．
$F_{\rm old}$ は $F$ の切断乗法であり，その随伴
\[
 \alpha_{F_{\rm old}}:
 \mathcal H_{02}\to\mathcal H_{23}^\vee\otimes\mathcal H_{03}
\]
も $F$ の大域生成性によって部分束包含である．
従って実際に $U_0$ の全射・部分束条件を満たす．
これは Euler 行列の一致だけから normal chart への所属を
推論したものではない．

四項は有界な有限 projective 複体とその cone から作られる．
通常 End の全 Hom が vector bundle であることを既に示したので，
tilting Morita 同値
$\Perf(\Gamma_F)\simeq\Perf(\Lambda)$ は全四対象を保ち，
任意基底変更に整合する．

最後にBの ruling 交換を $\mathcal N_F$ に適用すると，
新しい第二対象は
\[
 \operatorname{Cone}(V^\vee\otimes E_3\to Z_1)[-1]=E_1,
\]
第四対象は $\operatorname{Cone}(E_2\to\mathcal H_{23}^\vee\otimes E_3)=Z_2$
である．従って二つの窓が交換される．
この同定は同じ二つの cone 三角の交換から来るため，
二回適用と三重重なりの整合性も保たれる．
\end{proof}

\begin{remark}[相対determinant因子]
曲線へ制限した $Z_1$ は一般には $FG^2$ 自体ではない．
pencil の双対評価列から
\[
 0\to F\to(p_*G)^\vee\otimes FG
   \to FG^2\otimes\det(p_*G)^\vee\to0
\]
となるので，最後の determinant 因子を保持する．
$\mathfrak F_0$ では頂点ごとの基底線束 twist は剛体化に吸収されるが，
実際の相対 tilting 対象や End 代数の比較では省略しない．
上の cone による定義はこの因子を自動的に保持する．
\end{remark}

\begin{theorem}[一般square族の相対$E_0$・$E_3$認識]
\label{thm:qa-relative-square-e0-e3-h2}
上記の avatar 族に対して $\lambda_F$ は
$\mathfrak F_{\HH}$ を通り，したがって実際の射
\[
 \rho_0:S\longrightarrow\HHstack
\]
を与える．その相対 Abel 圏を $\mathcal C_0$，構造対象を
$\mathcal O_0$ と書くと，pointed compact Morita 同値
\[
 (\Perf\Lambda,E_0)\simeq(\mathcal C_0^{\rm perf},\mathcal O_0)
\]
があり，全ての相対 Serre 構造と任意基底変更に整合する．
ここで $\mathcal C_0^{\rm perf}$ は局所canonical compact chartと整合する相対perfect圏を表す．
今回の有限tiltingと内部共役による降下の範囲で扱い，任意affine上の
大域的 $D(\operatorname{Ind}\mathcal C_0)^c$ との同定や，
特異基底上での $D^b(\mathcal C_0)$ との同一視は主張しない．

四つの $E_i$ は同じ $\mathcal C_0$ の基底上平坦な compact 対象へ移る．
特に $E_3$ の像を $\mathcal M$ とすると
\[
 (\mathcal C_0,\mathcal M)\in\HHstack(S)
\]
であり，第二の実際の射 $\rho_3:S\to\HHstack$ を得る．
両構成は avatar の全同型と $F,G$ の交換に自然である．
後者の交換では内在的 ruling の選択だけが交換される．
\end{theorem}
\begin{proof}
まず $\lambda_F$ の各幾何 fiber が\emph{真正な normal foundation}
であることを確認する．$D_F=\{G=FP\}$，$D_G=\{GP=F\}$ と書く．

$D_F$ の外では第一mutationが cubic 型であり，
その標準四項 $(H_1,H_2,H_3,H_4)$ は $\mathcal N_F$ に一致する．
第四項の同定は full exceptional helix の Serre mutation から従う：
$H_0=\operatorname{Cone}(E_0\otimes\Hom(E_0,E_1)\to E_1)[-1]$
を右へ移し，$E_1$ と $E_2$ の直交性を用いると
$tH_0=\operatorname{Cone}(E_1\to\Hom(E_1,E_3)^\vee\otimes E_3)=Z_1$
となる．論文Aの全 cubic 逆構成と
\cite[\texttt{thm:all-admissible-inverse},\texttt{thm:cubic-coverage}]{QuadricIntrinsic2026} の必要性証明は，標準 cubic 窓の
第三対象を $s\mathcal O$，第二対象を実際の MCM ruling に同定する．
第四対象はBの \cite[\texttt{cor:two-intrinsic-foundations}]{QuadricModuli2026} により一意である．
従ってこれらの同定により，
これは $H_1=E_0$ を構造対象とする真正な normal foundation である．

$D_G$ の外では第二mutationによって $\mathcal N_G$ が同じ意味で
cubic normal foundation となる．前命題の相対 ruling 交換は
$\tau\mathcal N_F=\mathcal N_G$ を与え，Bの
\cite[\texttt{prop:ruling-involution-mutation}]{QuadricModuli2026} は真正な normal foundation を
もう一方の真正な normal foundation へ送る．
従って single-wall $D_F\setminus D_G$ でも $\mathcal N_F$ は normal である．
二壁の交点は非自明二捩れの全 noncubic locus であり，
既存の全三型の pointed 認識と実際の spinor 評価列が
$\mathcal N_F$ を normal foundation に同定する．
これで全幾何 fiber の $\mathfrak F_{\HH}$ への所属を得る．

$\mathfrak F_{\HH}\subset\mathfrak F_0$ は既に証明された\emph{開}埋込である．
従って $\lambda_F^{-1}(\mathfrak F_{\HH})$ は $S$ の開部分であり，
全幾何点を含むので $S$ 全体に等しい．
この開への因子化は非被約基底でも成立する．
$\Xi^{-1}\lambda_F$ と ruling 忘却を合成して $\rho_0$ を得る．
この段階で相対 $H_2$ の全公理は $\Xi^{-1}$ から供給されており，
fiber ごとの AS4 被覆を単に貼り合わせたのではない．

次に compact 比較を説明する．fppf 局所で $\Xi^{-1}\lambda_F$ の
normal 四対象と $\Gamma_F$ の頂点を比較する．
Bの相対 full normal tilting とその全 Hom の比較により，
相対 tilting Morita 同値 $\Perf\Gamma_F\simeq\mathcal C_0^{\rm perf}$ を得る．
上の $\Perf\Lambda\simeq\Perf\Gamma_F$ と合成する．
比較同型の変更は頂点スカラーによる内部共役である．
その実施単元は零頂点で一となるよう一意に正規化でき，
対応する Morita functor の自然同型は $E_0$ の標識上恒等となる．
三重重なりでも同じ内部単元の積により整合するので，
この pointed compact 比較は降下する．
全構成は有限 projective 複体，cone，dual と tensor から作られ，
任意の基底変更に整合する．

$\mathcal N_F$ の各項の像は相対 normal 四項なので
heart に属し，基底上平坦かつ compact である．
もう一方の ruling 窓の第二項が $E_1$ なので，$E_1$ にも同じことが成立する．
特に $E_3$ は normal foundation の第三対象 $M$ へ移る．
$H_2$ の normal helix では，fppf 局所に適合 Serre 根 $s$ を選ぶと
$M\simeq s\mathcal O$ である（基底線束 twist を含む場合も許す）．
exact 自己同値 $s$ と基底線束による tensor は $H_2$ 条件を移すから，
同じ heart の $(\mathcal C_0,\mathcal M)$ も相対 $H_2$ である．
大域的に一つの $s$ を選ぶ必要はない．
対象 $\mathcal M$ 自体は実際の $E_3$ の像として既に降下しており，
$H_2$ の性質は fppf 局所で課されるためである．

最後に $F,G$ の交換は $\lambda_F$ と $\lambda_G$ を交換する．
$\Xi$ の相対 ruling 対合への同変性と
$\HHstack=[\mathfrak H_2^{\rm rul}/S_2]$ により，
$\rho_0$ はこの交換と整合する．
Bの ruling 交換は第三対象 $M$ を保つので $\rho_3$ も同様である．
\end{proof}

\begin{corollary}[全Serre tensor次数の比較]
\label{cor:qa-relative-avatar-all-serre-degrees}
前定理の compact Morita 同値は，$E_0$ と $E_3$ の
全ての非負次数 $m\ge0$ の反Serre orbit 複体とその積を相対的に同定する．
各次数は finite locally free で，任意基底変更と可換する．
さらに $\Lambda$ 自体の非負内部次数付き Serre completion も，
命題~\ref{prop:qa-relative-serre-tensor-base-change} により
通常の次数付き代数であり，次数別に任意基底変更と可換する．
\end{corollary}
\begin{proof}
pointed Morita による inverse dualizing complex と全 tensor 冪の移送は
命題~\ref{prop:qa-pointed-morita-serre-transport} を用いる．
相対 $H_2$ の反Serre orbit は全非負次数で vector bundle である．
$\Lambda$ 全体については，各幾何 fiber の全整数 helix の
前向き Ext 消滅から同命題の geometric forward 条件が成立する．
finite diagonal 解消により各 tensor 冪が perfect なので，
導来 fiber の次数集中判定が全次数それぞれに適用できる．
\end{proof}

\begin{remark}[基底とheartの比較範囲]
有限代数，cone，vector bundle，compact Morita と Serre tensor の
構成そのものは，任意のスキームへの基底変更と可換する．
内在的な Abel 圏の全公理を持つ $\HHstack$ への射については，
論文Bで定義された局所有限型 $k$-scheme の site 上で述べている．
任意の Noetherian 基底に対する素朴な Abel 公理との一致や
strong Noether 性をこの構成から追加で結論しない．

また，ここで構成した heart は $\Xi^{-1}$ による実際の内在的 heart
$\mathcal C_0$ である．別に与えた全次数の相対 square AS
$\mathbb Z$-代数族の $\qgr$ heart と同じであることは，
その族の相対 Noether 性，torsion 局所化，または相対 AZ 比較を
確認した場合に述べる．finite End の同型や全 Serre tensor の一致だけで
この追加の $t$-exactness を結論しない．

$E_1,E_2$ の幾何的 $H_2$ 開部分では，別の cubic normal 窓から
同様の相対 $H_2$ 族と compact marking を構成する方法がある．
しかしその構成と $\mathcal C_0$ の\emph{同じ heart}内の
$(\mathcal C_0,E_i)$ との一致を主張するには，
cubic tails の相対認識または相対 AZ 比較が別に必要である．
本節の $E_3$ はこの問題を $M=s\mathcal O$ によって直接解決した場合である．
任意 AS/avatar/QP 表示 stack 全体の同値や，表示を忘れる全相対 fiber の
分類は本定理に含まれない．
\end{remark}

% !TeX root = ../nc-quadric-surface-avatar.tex
% 2026-11-23 / v1.0
\section{相対heartの比較とcubic代数族の回復}
\label{sec:qa-relative-avatar-heart-comparison}

\subsection{相互の生成対象による相対 heart の認識}
\label{subsec:qa-relative-mutual-heart-recognition}

次の議論は fiberwise な $t$-exactness の持上げではない．
両側で既に構成された相対 $H_2$ heart の中に，互いの構造対象が
実際に存在することを用いる．したがって非被約基底でも同じ証明が成立する．

\begin{lemma}[生成族に対する非正性の検出]
\label{lem:qa-generators-detect-right-texactness}
$\mathcal A,\mathcal B$ を局所 noetherian Grothendieck 圏とし，
$\mathcal A$ の noetherian 部分を $\mathcal A^{\rm noeth}$ とする．
$\{L_\alpha\}$ を $\mathcal A^{\rm noeth}$ 内の族であって，
全ての noetherian 対象がその有限直和の商であると仮定する．
標準の安定導来 enhancement の間の colimit を保つ完全関手
\[
 \mathsf F:D(\mathcal A)\longrightarrow D(\mathcal B)
\]
が全ての $L_\alpha$ を $D^{\leq0}(\mathcal B)$ に送るなら，
$\mathsf F$ は右 $t$-exact，すなわち
$\mathsf F(D^{\leq0}(\mathcal A))\subset D^{\leq0}(\mathcal B)$ である．
\end{lemma}
\begin{proof}
任意の noetherian 対象 $M$ に対し，生成族からの全射を取り，
その kernel にも同じ操作を繰り返す．kernel も noetherian なので，
\[
 \cdots\longrightarrow Q^{-2}\longrightarrow Q^{-1}
       \longrightarrow Q^0\longrightarrow M\longrightarrow0
\]
という完全列を得る．各 $Q^{-j}$ は $L_\alpha$ の有限直和である．
これは射影解消である必要はなく，通常の複体 $Q^\bullet$ が
$M$ とquasi-isomorphicであることだけを使う．

$Q^\bullet$ は，区間 $[-n,0]$ に制限した brutal truncation
$\sigma^{\geq-n}Q^\bullet$ の順方向 homotopy colimit である．
実際，これらは複体の包含写像による系であり，filtered colimit の
完全性からその homotopy colimit は項別 colimit とquasi-isomorphicである．
各有限複体の像は $D^{\leq0}(\mathcal B)$ に属する．
同部分圏は coproduct と extension に閉じているので，
sequential homotopy colimit にも閉じている．したがって
$\mathsf F(M)\in D^{\leq0}(\mathcal B)$ である．

任意の $\mathcal A$ の対象は noetherian 部分対象の filtered colimit である．
導来圏の標準 enhancement において，この colimit は homotopy colimit
とも一致し，$D^{\leq0}$ はこれらのfiltered homotopy colimitに閉じている．
これはGrothendieck圏のAB5による．$\mathsf F$ はこれを保つので，任意の heart 対象も
非正部分に送られる．
最後に，$D^{\leq0}(\mathcal A)$ の任意の対象を次数 $0$ 以下の複体で表し，
再び有限 brutal truncation の homotopy colimit として書けば結論を得る．
本稿で用いる $\mathsf F$ は perfect Morita 関手の Ind 延長であり，
この enhancement と全 colimit の保存を備えている．
\end{proof}

\begin{proposition}[相対 $H_2$ heart の相互標識判定]
\label{prop:qa-h2-mutual-markings-heart-test}
同じ基底 $S$ 上の二つの相対 $H_2$ 族を
$(\mathcal C,\mathcal O)$，$(\mathcal D,\mathcal P)$ とする．
共通の fppf 局所 chart 上の compact 圏の間に，
任意基底変更と整合する $S$-線形 Morita 同値
\[
 \mathsf F^c:\mathcal D^{\rm perf}\xrightarrow{\sim}
                  \mathcal C^{\rm perf}
\]
があるとする．さらに，実際の相対 heart に関する条件
\[
 \mathsf F^c(\mathcal P)\in\mathcal C,
 \qquad
 (\mathsf F^c)^{-1}(\mathcal O)\in\mathcal D
 \tag{MH}\label{eq:qa-mutual-heart-condition}
\]
がその chart の全ての affine 基底変更上で成立すると仮定する．
このとき $\mathsf F^c$ は唯一の colimit を保つ導来 Morita 同値へ延長され，
その延長は $t$-exact である．従って基底変更と整合する完全同値
\[
 \mathsf F:\mathcal D\xrightarrow{\sim}\mathcal C
\]
を得る．初めの compact 比較に coherent な降下データがある場合，
この完全同値にも同じ降下データがある．
\end{proposition}
\begin{proof}
相対 $H_2$ の定義により，各局所 affine chart 上で
$D(\operatorname{Ind}\mathcal C)$ と
$D(\operatorname{Ind}\mathcal D)$ は compactly generated である．
従って perfect Morita 同値はこれらの間の colimit を保つ同値に
一意に延長される（\cite[\S4.3,\S8.2]{Keller1994}の標準dg Morita）．

canonical な相対反 Serre 関手をそれぞれ $t_{\mathcal C},t_{\mathcal D}$
と書く．両者は heart 上で完全であり，Serre 双対の一意性により
\[
 \mathsf F^c t_{\mathcal D}\simeq
 t_{\mathcal C}\mathsf F^c
\]
が成り立つ．論文Bの
\cite[\texttt{prop:relative-h2-antiserre-reconstruction}]{QuadricModuli2026}
により，$\{t_{\mathcal C}^{-n}\mathcal O\}_{n\geq0}$ と
$\{t_{\mathcal D}^{-n}\mathcal P\}_{n\geq0}$ はそれぞれ AZ 生成族である．
条件~\eqref{eq:qa-mutual-heart-condition} と $t$ の完全性から，
これらの生成族は $\mathsf F$ とその逆によって相手の heart に送られる．
前補題を両向きに適用すると，$\mathsf F$ と $\mathsf F^{-1}$ は
ともに右 $t$-exact である．従って
$\mathsf F(D^{\leq0})=D^{\leq0}$ であり，直交性から非負部分も保つ．
これで $t$-exactness が従う．

得られる Ind heart の完全同値は noetherian 対象を保つので，
指定された coherent heart の完全同値に制限される．
任意の基底変更に対する導来比較は compact 生成対象上の比較から
一意に延長される．通常の module-object 引戻しはその次数零 cohomology
なので，$t$-exactness によって通常の引戻しとの coherent な比較を得る．
降下データの間の自然同型も compact 生成対象から一意に延長される．
したがって局所判定を貼り合わせられる．ここでは，任意の相対 compact
圏に独立の heart を与えたり，fiberwise な cohomology 一致から heart を
復元したりしていない．
\end{proof}

\begin{corollary}[中間頂点も同じ相対 heart 内で認識される]
\label{cor:qa-relative-middle-same-heart}
定理~\ref{thm:qa-relative-avatar-middle-h2-realizations} の記法で，
その compact Morita 比較は完全同値
\[
 \mathcal C_i\xrightarrow{\sim}\mathcal C_0|_{U_i},
 \qquad \mathcal O_i\longmapsto E_i\quad(i=1,2)
\]
へ延長される．従って，元の同じ相対 heart の対
$(\mathcal C_0|_{U_i},E_i)$ が実際に $\mathfrak H_2(U_i)$ の対象である．
この主張と比較は任意基底変更と整合する．
\end{corollary}
\begin{proof}
前節では全 $E_j$ が $\mathcal C_0$ の基底平坦な heart 対象であることを
既に示した．従って $\mathsf F^c(\mathcal O_i)=E_i$ は第一条件を満たす．
一方，$\mathcal N_i$ の第四対象は $t_\Lambda E_0$ である．
真正な normal foundation の第四対象は $\mathcal C_i$ の基底平坦な
heart 対象だから，canonical $t_i^{-1}$ を適用して
\[
 (\mathsf F^c)^{-1}(E_0)
   =t_i^{-1}(\text{$\mathcal N_i$ の第四対象})\in\mathcal C_i
\]
を得る．$t_i^{-1}$ は完全であり基底変更と可換するので，
この対象も基底上平坦であり，全ての affine 基底変更後に同じ議論が成立する．
これが第二条件である．前命題を適用すればよい．
頂点の相対枠を変えた際の基底線束との tensor は，heart の完全自己同値であり，
この二条件と基底平坦性を保つ．従って実際の cone に含まれる determinant
線束や，剛体化した foundation chart の頂点線束を自明化する必要もない．
\end{proof}

\begin{remark}[この判定の範囲]
前命題で必要なのは互いの構造対象が相手の実際の heart に属するという
二つの条件である．その検証を，幾何 fiber で heart に属することだけで
置き換えることはしない．また相手が未構成の一般 AS 商圏である場合には，
compact generation と相対 AZ 生成族の存在が別途必要である．
したがって本系は，任意の相対 AS/QP stack と $H_2$ stack の同一視ではない．
\end{remark}

\begin{theorem}[同じ相対heartの全標準構造頂点の開条件]
\label{thm:qa-relative-all-vertex-h2-open}
前節のflat曲線型square avatarから構成した相対heart $\mathcal C_0$ と
その四つの基底平坦な対象 $E_i$ を取る．
局所有限型siteの任意の底変更 $S'\to S$ について，
\[
 (\mathcal C_{0,S'},E_{i,S'})\in\HHstack(S')
 \quad\Longleftrightarrow\quad S'\to S\text{ が }U_i\text{ を経由する},
\]
ただし
\[
 U_0=U_3=S,\qquad U_1=D(\det\mu_F),\qquad U_2=D(\det\mu_G)
\]
である．したがってこの二行列式開は，幾何fiberの適合性だけでなく，
\emph{同じ実際の相対heart}の中の強い$\HH$構造対象条件を表す．
必要な根・AS4被覆は局所存在条件のままであり，標識へ加えない．
\end{theorem}
\begin{proof}
$U_0,U_3$ では定理~\ref{thm:qa-relative-square-e0-e3-h2}，
$U_1,U_2$ では前系を用いる．
いずれも任意のsite内の底変更と整合する．
逆に相対$\HH$の条件は全幾何fiberの$\HH$所属を含むので，
定理~\ref{thm:qa-relative-vertex-determinant-open}によって
底変更は該当する開を経由する．
\end{proof}


\subsection{有限窓からの実際のcubic族}

\begin{proposition}[avatarのnormal窓からの相対cubic代数族]
\label{prop:qa-relative-cubic-from-avatar-window}
flat な曲線型 square avatar の族に対し，
$U_2=D(\det\mu_G)$ 上の normal 窓
\[
 (E_0,E_2,E_3,Z_1),\qquad
 Z_1=\operatorname{Cone}
   (E_1\to\mathcal H_{13}^{\vee}\otimes E_3)
\]
を取る．その実際の有限 End 代数を $\Gamma_F$ とする．
すると非線型 geometric quintuple の相対族が基底変更と可換に復元され，
その cubic $\mathbb Z$-代数族の初期四項の End は
$\Gamma_F$ と全ての積を含めて同型になる．
全次数の成分は有限局所自由で，標準 cubic 解消と AS 双対は
任意基底変更後にも成立する．
局所 affine な Noetherian 底では全 projective row と column は
左右 Noetherian であり，任意の Noetherian 基底変更後も同様である．

この cubic 族の canonical tails は，実際の相対 $H_2$ 族である．
その指定零対象からの normal foundation は $\lambda_F$ を与えるので，
$\Xi^{-1}\lambda_F$ によって構成した $E_0$ の相対 $H_2$ 族と
pointed exact 同値である．
$U_1=D(\det\mu_F)$ 上にも $F,G$ を交換した同じ主張が成立する．
\end{proposition}
\begin{proof}
\emph{(1) 有限ベクトル束から実際のquintupleを作る．}
記号を
\[
 \mathcal U=\mathcal H_{02},\qquad
 \mathcal V=\mathcal H_{23},\qquad
 \mathcal W=\mathcal H_{13}^{\vee},\qquad
 C=\operatorname{coker}
   (\mathcal H_{01}\to\mathcal H_{13}^{\vee}\otimes\mathcal H_{03})
\]
とする．三つの束 $\mathcal U,\mathcal V,\mathcal W$ は階数二，
$C$ は階数六である．$U_2$ 上では第一の合成
\[
 \mathcal V\otimes\mathcal U\xrightarrow{\sim}\mathcal H_{03}
\]
が同型であり，最後の合成は全射なので
\[
 0\to R\longrightarrow
 \mathcal W\otimes\mathcal V\otimes\mathcal U
 \longrightarrow C\to0
\]
という階数二の核を持つ局所分裂完全列を得る．
従って包含をテンソルに読み替えると
\[
 w_R\in
 \mathcal U\otimes\mathcal V\otimes\mathcal W\otimes R^\vee
\]
となる．これは nowhere zero で，四つの階数二の束と
可逆なテンソル直線 $\OO\,w_R$ を与える．
全構成はベクトル束の核・双対・tensor と実際の積から来るので，
非平坦なものを含む任意基底変更と可換する．

これは論文B \cite[\texttt{prop:cubic-foundation-open}]{QuadricModuli2026} の構成そのものである．
各幾何 fiber は第一mutationで得た cubic AS 代数の有限窓だから，
$w_R$ は各 fiber で geometric である．
geometric 条件は双対射影束上の proper incidence の補集合という
実際の開条件なので，この族全体が geometric 開へ因数分解する．
枠を選ぶ局所化を除けば，この操作は avatar の全同型にも自然である．

\emph{(2) linear fiber がないこと．}
もしある幾何 fiber が linear cubic であれば，その標準四項を
\[
 (H_0,H_1,H_2,H_3)
 =(\OO,\OO(1,0),\OO(1,1),\OO(2,1))
\]
と取れる．逆mutationの余評価列は通常の pencil 評価列であり，
\[
 M=\operatorname{coker}(\OO\to\OO(1,0)^{\oplus2})
   =\OO(2,0).
\]
従って逆square窓 $(H_1,M,H_2,H_3)$ は，
$\OO(-1,0)$ を共通tensorすると
\[
 (\OO,\OO(1,0),\OO(0,1),\OO(1,1))
\]
という classical square になる．
補題~\ref{lem:qa-as10-all-index-mutual-inverse} の全整数での互逆性と
定理~\ref{thm:qa-as-avatar-bijection} の同型類分類により，
元の square avatar も classical でなければならない．
これは入力が Cartier 曲線型であることに反する．
有限 End から回復される一点表現も，曲線型では一次元の $D$，
classical では二次元の $Y$ なので両型は区別される．
従って全族は非線型 geometric 開へ因数分解する．
幾何 fiber に linear 点がないことからこの\emph{開}への因子化を
結論しているので，非被約基底でも問題はない．

\emph{(3) 全次数とcanonical tails．}
復元したquintupleから，三次関係を $w_R$ の縮約として定め，
四因子を巡回して全整数の cubic $\mathbb Z$-代数を作る．
初期四頂点では二段階合成に関係がなく，三段階合成の関係は
正に上の $R$ である．従ってその有限 End は $\Gamma_F$ と一致する．
この同定は単に Hom 階数の一致ではなく，全ての積を保持する．

有限型の非線型 geometric 普遍族の
論文B \cite[\texttt{lem:universal-flatness}]{QuadricModuli2026} を引き戻せば，
全次数の局所自由性，標準 cubic 解消，AS 双対，および全基底変更との
整合性を得る．非被約な入力基底の行列の fiber 階数一定性だけから
平坦性を推論する必要はない．
同論文 \cite[\texttt{thm:relative-geometric-cubic-noetherian}]{QuadricModuli2026} が
任意 Noetherian 基底変更後の row/column Noether 性を与え，
\cite[\texttt{prop:general-cubic-canonical-tails}]{QuadricModuli2026} が実際の module-object
基底変更と整合する canonical tails，四項 tilting と Serre 比較を与える．
右加群の添字反転は
定理~\ref{thm:qa-geometric-cubic-relative-forgetting} の規約に従う．

同定した初期四項を $Q_0,Q_1,Q_2,Q_3$ とし，
forward の添字規約で書けば
\[
 (Q_0,Q_1,Q_2,Q_3)=(E_0,E_2,E_3,Z_1)
\]
である．有限 End の相対 Morita 同値は相対 Serre functor を保つので，
canonical inverse dualizing complex とその各非負 tensor 冪は
$\Gamma_F$ のものと一致し，この比較は積と基底変更に整合する．
ここで得た実際の cubic 代数族に，四次 overlap の可逆束や
頂点線束を省略した厳密な四周期を追加で課すことはしない．
全整数添字付き代数と canonical Serre 完成の\emph{厳密な}表示同一視には，
これらの可逆束を含む vertex-line transport の明記が別に必要である．

\emph{(4) 既存の相対H2射とnormal窓の認識．}
定理~\ref{thm:qa-geometric-cubic-relative-forgetting} を適用すれば，
上の canonical tails 自体が $Q_0$ を構造対象とする
実際の $\HHstack$ 族となる．
非線型の特異・可約・非被約曲線もこの定理の範囲に入る．

その標準四項が相対的にも真正な normal 窓であることを確認する．
適合二周期を選ぶ fppf 局所では $Q_2=sQ_0$ が実際に成立する．
真の局所 normal ruling を $X,Y$ とすると，$Q_1$ は
各幾何 fiber で $X$ または $Y$ の一方である
（論文A \cite[\texttt{thm:cubic-coverage}]{QuadricIntrinsic2026} の必要性証明）．
双方は canonical compact なので $\RHom(Q_1,X)$ は
perfect で導来基底変更と可換し，fiber は $k[0]$ または零である．
従って matching locus は開閉で，該当側では Hom は可逆層となる．
評価
\[
 Q_1\otimes\Hom(Q_1,X)\longrightarrow X
\]
の perfect cone は全幾何 fiber で零なので，評価は同型である．
非matching側では $Y$ に同じ議論を適用する．
第四項も normal 第四項との perfect Hom と評価，または相対的な
左直交の一意性で同定される．
これで，基底線束 twist を除いて標準四項が実際の normal foundation
となることを確認できる．

従ってこの相対 $H_2$ 族の選択normal窓の有限表示は
実際の $\Gamma_F$，すなわち $\lambda_F$ に一致する．
$\Xi$ の全相対 groupoid 上での全忠実性により，
この canonical cubic tails は $\Xi^{-1}\lambda_F$ と
pointed exact 同値である．
これは fiberwise exact 同値を直接相対化したものではなく，
既に構成済みの二つの相対 $H_2$ 族を，実際の finite normal End を通して
比較している．$U_1$ は $F,G$ の交換で同じである．
\end{proof}

\begin{remark}[用途と範囲]
この命題は二つの開上で，有限 avatar 窓から実際の全次数 cubic 族を
回復する手続きを与える．別に与えた任意の相対 AS/QP 族とこの族が
同型であることは，その族の全整数積と Serre transport を比較して述べる．
また，この段階でも任意の表示 groupoid 全体の同値を主張しない．
\end{remark}

% !TEX root = ../nc-quadric-avatar.tex
% 2026-11-23 / 相対square完成と実際のtails認識。
\section{相対square完成のAS解消とtails認識}
\label{sec:qa-relative-square-completion}

\subsection{平坦compact対象の相対Ext次元}
\label{subsec:qa-flat-compact-ext-dimension}

以下では基底を局所有限型 $k$-scheme とし，必要に応じて共通の fppf
局所 affine chart $S=\operatorname{Spec}R$ を取る．特異な基底上で
coherent heart 全体の大域次元が二であるとは仮定しない．

\begin{lemma}[平坦compactな始域に対する相対Serre評価]
\label{lem:qa-flat-compact-ext-two}
$\mathcal C$ を相対 $H_2$ の heart とし，その局所 compact 圏が
$R$ 上有限 projective で相対 smooth な代数 $\Gamma$ によって
Morita 表示されているとする．$E\in\mathcal C$ が $R$ 上平坦かつ
compact なら，任意の $K\in\operatorname{Ind}\mathcal C$ に対して
\[
 \operatorname{Ext}^q_{\operatorname{Ind}\mathcal C}(E,K)=0
 \qquad(q>2)
 \tag{FE2}\label{eq:qa-flat-compact-ext-two}
\]
が成立する．この結論は同じ仮定を満たす任意基底変更後にも成立する．
\end{lemma}
\begin{proof}
$I$ を $R$-加群圏の injective cogenerator とする．
相対 Serre 双対を compact な第二変数だけに限定せず，次の随伴として使う：
\[
 \mathbf R\operatorname{Hom}_R
   (\mathbf R\operatorname{Hom}_{\mathcal C}(E,K),I)
 \simeq
 \mathbf R\operatorname{Hom}_{\mathcal C}
   (K,\mathsf S E\otimes_R^{\mathbf L}I).
 \tag{SA}\label{eq:qa-serre-injective-adjunction}
\]
この式は任意の $K\in D(\operatorname{Ind}\mathcal C)$ に対して成立する．
実際，Morita 表示で $E$ に対応する perfect 右 $\Gamma$-加群を $Q$，
\[
 Q^\vee=\mathbf R\operatorname{Hom}_\Gamma(Q,\Gamma)
\]
をその左双対とする．$\Gamma$ が $R$ 上 perfect なので $Q^\vee$ も
$R$ 上 perfect である．tensor--Hom 随伴と有限 perfect 双対によって
\[
 \begin{split}
 \mathbf R\operatorname{Hom}_R
  (K\otimes_\Gamma^{\mathbf L}Q^\vee,I)
 &\simeq \mathbf R\operatorname{Hom}_\Gamma
     (K,\mathbf R\operatorname{Hom}_R(Q^\vee,I))\\
 &\simeq \mathbf R\operatorname{Hom}_\Gamma
     (K,\mathbf R\operatorname{Hom}_R(Q^\vee,R)
                    \otimes_R^{\mathbf L}I)
 \end{split}
\]
を得る．最後の右加群は $\mathsf S Q\otimes_R^{\mathbf L}I$ であり，
これが式~\eqref{eq:qa-serre-injective-adjunction} である．
ここで $K$ の compactness や基底平坦性は使っていない．

相対 $H_2$ では $\mathsf S E=t^{-1}E[2]$ であり，canonical $t$ は
heart 上完全で $R$-線形である．従って $t^{-1}E$ も基底平坦で，
\[
 \mathsf S E\otimes_R^{\mathbf L}I
       =(t^{-1}E\otimes_R I)[2]
\]
となる．括弧内は Ind heart の対象である．$K$ も heart にあるので，
式~\eqref{eq:qa-serre-injective-adjunction} の右辺は cohomological
次数 $-2$ 未満で零である．一方 $I$ の injectivity により，左辺の
次数 $-q$ cohomology は
\[
 \operatorname{Hom}_R
       (\operatorname{Ext}^q_{\mathcal C}(E,K),I)
\]
である．$I$ が cogenerator であることから
式~\eqref{eq:qa-flat-compact-ext-two} が従う．
\end{proof}

\begin{proposition}[有限個の平坦対象を加えた相対AZ生成]
\label{prop:qa-relative-flat-foundation-az}
$\mathcal C,\mathcal O$ を相対 $H_2$ 族とし，
$E_0=\mathcal O,E_1,\ldots,E_r$ を基底平坦 compact な heart 対象とする．
各 affine 基底変更の共通 fppf 局所 chart 上で前補題の有限 Morita
表示があり，そこで各 $i$ に対し
\[
 \operatorname{Ext}^q(E_i,t^m\mathcal O)=0
 \qquad(q>0,\ m\gg0)
 \tag{FO}\label{eq:qa-flat-foundation-forward-o}
\]
を仮定する．このとき
\[
 T=\bigoplus_{i=0}^r E_i
\]
に対し $(T,t)$ は各 affine 基底変更上で AZ ample である．
\end{proposition}
\begin{proof}
$(\mathcal O,t)$ は論文Bの相対反 Serre 再構成により AZ ample であり，
$\mathcal O$ は $T$ の直和因子なので生成条件 AZ1 は直ちに従う．

まず上記の局所 Morita chart 上で議論する．AZ2 のため，任意の
coherent 対象 $K$ に対して
\[
 \operatorname{Ext}^1(t^{-n}E_i,K)=0\quad(n\gg0)
 \tag{FV}\label{eq:qa-flat-foundation-ext-one-vanish}
\]
を示す．AZ1 を二回使い，短完全列
\[
 0\longrightarrow L\longrightarrow P\longrightarrow K\longrightarrow0,
 \qquad
 0\longrightarrow L'\longrightarrow P'\longrightarrow L\longrightarrow0
\]
を取る．ここで $P,P'$ は有限個の $t^{-a}\mathcal O$ の直和である．
基底の正則性も $K,L,L'$ の基底平坦性も仮定しない．
条件~\eqref{eq:qa-flat-foundation-forward-o} により，十分大きい $n$ について
\[
 \operatorname{Ext}^1(t^{-n}E_i,P)=0,
 \qquad \operatorname{Ext}^2(t^{-n}E_i,P')=0
\]
となる．さらに $t^{-n}E_i$ は基底平坦 compact なので，前補題から
$\operatorname{Ext}^3(t^{-n}E_i,L')=0$ である．
二つの長完全列により
\[
 \operatorname{Ext}^1(t^{-n}E_i,K)
  \hookrightarrow\operatorname{Ext}^2(t^{-n}E_i,L)=0
\]
を得る．これが式~\eqref{eq:qa-flat-foundation-ext-one-vanish} である．
任意の全射の kernel にこの消滅を適用すれば AZ2 を得る．
各 $i$ は有限個なので，一つの十分大きい $n$ を選べる．

一般の affine noetherian 基底 $\operatorname{Spec}R$ 上の AZ2 を確認する．
その準compactnessにより，共通 Morita chart となる有限個の affine fppf
被覆を選べる．その非交和を $\operatorname{Spec}R'$ とすれば，
$R\to R'$ は faithfully flat である．与えられた全射を引き戻すと，
有限個の chart 上の消滅から一つの十分大きい $n$ の範囲で
\[
 \operatorname{Hom}(t^{-n}T,M)\otimes_R R'
 \longrightarrow
 \operatorname{Hom}(t^{-n}T,N)\otimes_R R'
\]
が全射となる．ここでは相対 $H_2$ の coherent Hom の flat base change
と $t$ の基底変更整合性を用いている．余核を faithfully flat に
引き戻すと零なので，もとの Hom 写像も全射である．
これにより AZ2 は基底全体で成立する．この最後の議論には大域の
$D(\operatorname{Ind}\mathcal C)$ に対する Morita 表示を使っていない．
\end{proof}

\begin{remark}[AZ生成における構造対象の位置と二段presentation]
\label{rem:qa-relative-az-noninitial-structure}
命題~\ref{prop:qa-relative-flat-foundation-az}では，
指定構造対象$\mathcal O$が$T$のいずれかの
直和因子であればよく，先頭の因子である必要はない．
実際，AZ1は$\mathcal O$の直和因子性のみを使い，
AZ2の証明は各$E_i$から十分前向きの
$t^m\mathcal O$への高次Ext消滅のみを使う．
二回の$\mathcal O$-軌道対象からの表示
$0\to L\to P\to K\to0$，
$0\to L'\to P'\to L\to0$は
$\operatorname{Ext}^1(E,P)=
 \operatorname{Ext}^2(E,P')=
 \operatorname{Ext}^3(E,L')=0$
から$\operatorname{Ext}^1(E,K)=0$を与える．
最後の消滅には平坦compact始域の
補題~\ref{lem:qa-flat-compact-ext-two}を使い，
基底の絶対大域次元二は仮定していない．
従ってこれは反対normalの第一頂点が
強い$H_2$ pointingでなくても，
他の頂点に適合構造対象を含む場合に適用できる．
ただし全ての候補対象が\emph{同じ相対heart}の
基底平坦compact対象であることは別途確認しなければならない．
\end{remark}

\subsection{相対square完成の標準tails心}
\label{subsec:qa-relative-square-tails-heart}

\begin{theorem}[square完成の相対標識付きtails認識]
\label{thm:qa-relative-square-completion-tails}
相対 curve-type square avatar の実際の有限代数を $\Lambda$，
その四つの標準対象を $E_0,\ldots,E_3$ とする．
定理~\ref{thm:qa-relative-square-e0-e3-h2} の相対 $H_2$ heart
$\mathcal C_0$ を用い，$T=\bigoplus_{i=0}^3E_i$ と置く．
その通常の相対 Serre 完成を
\[
 C(\Lambda)=\bigoplus_{n\geq0}
 H^0(\Theta_\Lambda^{\otimes_\Lambda^{\mathbf L}n}),
 \qquad
 \Theta_\Lambda=\mathbf R\operatorname{Hom}_{\Lambda^e}
                      (\Lambda,\Lambda^e)[2]
\]
とする．全ての本稿の affine 基底変更上で，
$C(\Lambda)$ は右 noetherian であり，実際の完全同値
\[
 \qgr C(\Lambda)\xrightarrow{\sim}\mathcal C_0
 \tag{SQH}\label{eq:qa-relative-square-completion-heart}
\]
が存在する．標準冪等元を $e_i$ とすると，この同値は
$\pi(e_iC(\Lambda))\mapsto E_i$ と次数移動 $(1)\mapsto t$ を保つ．
四頂点の周期的 unfolding による square $\mathbb Z$-代数を
$A(\Lambda)$ と書けば，同じ主張は $\qgr A(\Lambda)$ にも成立する．

同値~\eqref{eq:qa-relative-square-completion-heart} は任意の本稿の
基底変更に対する module-object 引戻しと coherent に整合する．
次数別代数 $C(\Lambda)$，$A(\Lambda)$ の形成自体は，前節の
Serre tensor 比較により任意の scheme 基底変更と可換する．
\end{theorem}
\begin{proof}
前節で全 $E_i$ は $\mathcal C_0$ の基底平坦 compact な heart 対象であり，
$T$ の有限 End は実際の $\Lambda$ と同定されている．
全 forward Hom 複体は基底上 perfect で，全幾何 fiber の geometric helix
上の前向き Ext 消滅から次数零に集中する．特に $m\geq1$ について
\[
 \mathbf R\operatorname{Hom}(E_i,t^mE_0)
\]
は次数零の有限局所自由層である．これは任意基底変更後にも成り立つ．
従って前命題の条件~\eqref{eq:qa-flat-foundation-forward-o} が成立し，
$(T,t)$ は相対 AZ ample である．

通常の軌道環
\[
 R(T,t)=\bigoplus_{n\geq0}\operatorname{Hom}(T,t^nT)
\]
に対し，$\operatorname{End}(T)=\Lambda$ は右 noetherian であり，
全 coherent Hom は有限 $R$-加群なので有限 $\Lambda$-加群でもある．
一般の End 環に対する Artin--Zhang 再構成を適用すると，
$R(T,t)$ の右 noetherian 性と
\[
 \mathcal C_0\simeq\qgr R(T,t)
\]
を得る．ここで用いる形は，論文B
\cite[\texttt{prop:relative-h2-antiserre-reconstruction}]{QuadricModuli2026}
が引用する \cite[Theorem 4.5(1)]{AZ1994}，
\cite[Theorem 7.3]{OACU2026} と同じである．

共通の fppf 局所 chart の実際の tilting Morita 同値の下で $t$ は
$-\otimes_\Lambda^{\mathbf L}\Theta_\Lambda$ に対応する．
従って
\[
 \mathbf R\operatorname{Hom}(T,t^nT)
  \simeq \Theta_\Lambda^{\otimes_\Lambda^{\mathbf L}n}
\]
は積を保つ同型である．全次数の次数零集中を取り，coherent Hom の
flat base change に沿って降下すると，基底全体で
$R(T,t)=C(\Lambda)$ を得る．四つの冪等元による成分分解は標準
$E_i$ を保つ．有限個の頂点ごとに graded 成分を並べ直す通常の
folding/unfolding は graded 加群圏の同値を与え，有限生成性と
torsion を保つので $\qgr A(\Lambda)$ についても同じ結論となる．

最後に，AZ 同値の逆向きは graded 自由加群のシフトを
$t^{-n}T$ に送り，有限表示の余核へ延長して作られる．
通常の module-object 引戻しは右完全なので，この有限表示の構成と
可換する．標準射影 $e_iC(\Lambda)$ についても同様である．
自由対象上の比較は合成に対して coherent であり，有限表示への
延長の一意性から全対象上の比較も coherent となる．
非平坦基底変更に対して全対象の低次数 $\Gamma_*$ が可換すると
仮定する必要はない．
非 affine な基底ではこの比較に沿って Abel 圏の descent を適用し，
実際の完全同値を貼り合わせる．
大域の $D(\operatorname{Ind}\mathcal C_0)$ と局所 compact 圏の descent
を同一視する追加命題は用いていない．
\end{proof}

\begin{remark}[結論の範囲]
これは実際の曲線 avatar 族から構成した有限 square 代数とその canonical
Serre 完成に対する認識である．任意の AS/QP 族の Noetherian 性，任意の
presentation の相対 $H_2$ 性，あるいは全 presentation stack と
$\mathfrak H_2$ の同一視を結論していない．また
式~\eqref{eq:qa-flat-compact-ext-two} は始域を基底平坦 compact とした
相対的な消滅であり，特異な基底上の heart 全体の大域次元二という
主張ではない．
\end{remark}

% !TEX root = ../nc-quadric-avatar.tex
% 2026-11-23 / 相対square完成と実際のtails認識。
\subsection{有限square族への限定的な拡張}
\label{subsec:qa-finite-square-family-opposite}

以下で基底は本稿の局所有限型 $k$-scheme の site に属するものとする。
次の補題は任意の四頂点代数を扱うのではなく、全幾何 fiber が既に分類された
square 型である有限局所自由族だけを扱う。

\begin{lemma}[分類済みfiberを持つ有限square族]
\label{lem:qa-classified-fiber-finite-square-recognition}
$S$ 上の有限局所自由代数 $B$ が四つの直交冪等元
$e_0,e_1,e_2,e_3$ を持ち、その右射影 $P_i=e_iB$ の Hom 束が
\[
 \left(\operatorname{rk}\operatorname{Hom}(P_i,P_j)\right)_{i,j}
 =\begin{pmatrix}1&2&2&4\\0&1&0&2\\0&0&1&2\\0&0&0&1\end{pmatrix}
\]
であるとする。対角成分は単位によって $\mathcal O_S$ と同定され、
逆向き成分は零とする。
さらに各幾何 fiber が、本稿で分類した square AS 型の初期四項の
有限 End 代数と、頂点と全積を保って同型であると仮定する。
古典型も許す。

このとき $B$ の実際の相対 perfect 圏から、$P_0$ を構造対象とする
相対 $\mathfrak H_2$ 族が構成される。四つの $P_i$ はこの同じ heart の
基底平坦対象となり、共通の fppf 局所 chart で compact である。
通常の canonical Serre 完成
\[
 C(B)=\bigoplus_{n\geq0}
 H^0\left(\Theta_B^{\otimes_B^{\mathbf L}n}\right),\qquad
 \Theta_B=\mathbf R\operatorname{Hom}_{B^e}(B,B^e)[2]
\]
は、各局所 affine chart とその許された基底変更上で右 Noetherian であり、
その tails はこの相対 heart と同値になる。
\end{lemma}
\begin{proof}
有限 directed 代数の normalized bar 解消により、$B$ は相対 smooth proper であり、
全 perfect Hom と相対 inverse Serre 関手は任意基底変更と可換する。
$H_{ij}=\operatorname{Hom}(P_i,P_j)$ とし、実際の cone
\[
 Z_1=\operatorname{Cone}(P_1\to H_{13}^{\vee}\otimes P_3),\qquad
 N_F=(P_0,P_2,P_3,Z_1)
\]
を取る。幾何 fiber の分類によって、二つの随伴
\[
 H_{01}\longrightarrow H_{13}^{\vee}\otimes H_{03},\qquad
 H_{02}\longrightarrow H_{23}^{\vee}\otimes H_{03}
\]
は各 fiber で単射である。曲線型では pencil の大域生成性、古典型では
通常の $\mathbb P^1\times\mathbb P^1$ の評価列による。
従って両者は基底上でも部分束の包含となる。

定理~\ref{thm:qa-relative-square-e0-e3-h2} の有限 cone の計算は、
これらの Hom 束と全積しか使わない。
したがって $N_F$ は相対 full strong exceptional 四項であり、
通常 End は階数二十四、$U_0$ の全射・部分束条件を満たす。
反対の二中間頂点を用いる窓 $N_G$ とは、B の相対 ruling 対合で交換される。

各幾何 fiber で $N_F$ が真正な normal foundation であることも、
同定理の体上の証明をそのまま適用できる。
cubic 開上では cubic 逆構成、一重壁では相対 ruling 交換の体上の形、
二重壁では既存の noncubic 認識を使う。
古典型では例えば
$(\mathcal O,\mathcal O(0,1),\mathcal O(1,1),\mathcal O(1,2))$
という normal 四項に一致する。
よって実際の射 $S\to\mathfrak F_0$ の全幾何 fiber が
$\mathfrak F_{\HH}$ に入り、既存の開埋込 $\Xi$ の逆を適用できる。
正規窓およびもう一方の ruling 窓から、四つの $P_i$ が
同じ相対 heart の基底平坦対象であることを得る。

全幾何 fiber の square helix の forward Ext 消滅と perfect-fiber 判定により、
$\mathbf R\operatorname{Hom}(P_i,t^mP_j)$ は $m\geq1$ に対して
次数零に集中する。とくに相対 AZ 命題の forward 条件が成立する。
従って $\bigoplus_iP_i$ は相対 AZ ample であり、
定理~\ref{thm:qa-relative-square-completion-tails} の証明を適用して
右 Noether 性と tails の認識を得る。
ここでは幾何 fiber の分類を入力とし、大域的な曲線 avatar の復元や
まだ得ていない左 Noether 性は使っていない。
\end{proof}

\subsection{反対代数の四頂点と添字の規約}

\begin{lemma}[square有限基礎の反対代数]
\label{lem:qa-square-foundation-opposite}
前補題の $B$ に対し $B^{\rm op}$ の標準冪等元を
\[
 f_i=e_{3-i}\qquad(0\leq i\leq3)
\]
と番号付け直す。このとき $B^{\rm op}$ も前補題の仮定を満たす。
\end{lemma}
\begin{proof}
$Q_i=f_iB^{\rm op}$ とすると
\[
 \operatorname{Hom}_{B^{\rm op}}(Q_i,Q_j)
 =f_jB^{\rm op}f_i
 =e_{3-i}Be_{3-j}
 =\operatorname{Hom}_B(P_{3-j},P_{3-i}).
\]
したがって Hom の階数表は再び square 型となり、積は元の積を逆転したものである。

幾何 fiber が曲線型 avatar $(D;P,F,G)$ に対応する場合を直接確認する。
\[
 P'=P^{-1},\qquad F'=FP^{-1},\qquad G'=GP^{-1}
\]
と置く。既存の曲線消滅補題により $F',G'$ は大域生成の次数二 pencil、
$F'G'=FGP^{-2}$ は very ample である。
また $G'(F')^{-1}=GF^{-1}$ は acyclic なので、二 pencil は同じ $D$ の
$(2,2)$ 埋込みを与える。$P'\ne\mathcal O_D$ も保たれる。
従って $(D;P',F',G')$ は再び square avatar である。

元の四線束を $(L_0,L_1,L_2,L_3)=(P,F,G,FG)$ とすれば、
新しい四線束は
\[
 (P',F',G',F'G')
 =(FGP^{-1})\otimes
 (L_3^{-1},L_2^{-1},L_1^{-1},L_0^{-1}).
\]
可逆層の双対は射の合成を逆転し、共通 tensor は Hom を保つ。
よって新しい directed End は、全積を含めて $B^{\rm op}$ と同型である。
本稿の avatar から AS 代数への逆構成により、これは分類済み square AS 型の
有限基礎となる。古典型でも、逆順双対 collection に
$\mathcal O(1,1)$ を共通 tensor すると標準 square collection を得る。
従って全幾何 fiber で所要の仮定が成立する。
\end{proof}

\begin{remark}[全整数AS代数で書く場合の反転]
本稿の正向きの規約が $A_{ji}=\operatorname{Hom}(E_i,E_j)$ であるとき、
上の有限基礎に対応する indexed opposite は
\[
 A^{\dagger}_{ji}=A_{3-i,\,3-j}
\]
であり、積を逆にする。
$i\mapsto-i$ だけを使うと初期位相がずれるので、同じ四位相表を主張するときは
定数 $3$ を保持する。
右単純の分解を双対化すると、元の添字 $i$ の右単純から
添字 $i-4$ の左単純の分解を得る。
新しい添字 $j$ に対して元の $i=7-j$ であり、距離 $k$ は $4-k$ に変わる。
従って $Q_{1,\rm new}$ は $Q_{2,\rm old}$、$Q_{2,\rm new}$ は
$Q_{1,\rm old}$ から得られる。
この変換は定義~\ref{def:qa-as-square-rotated-minimal} の square の四行を
それぞれ同じ四行へ送る。
ただし前補題の証明は、この全整数の双対性を未証明の入力にせず、
各体上の双対 avatar の有限 End を直接用いた。
\end{remark}

\subsection{inverse dualizing complexとopposite}

\begin{lemma}[Serre完成の反対代数への適合性]
\label{lem:qa-serre-completion-opposite}
$B$ を基底上有限局所自由で相対 smooth な代数とする。
$M^{\rm op}$ は $B$-双加群 $M$ の左右を交換した $B^{\rm op}$-双加群を表す。
有限対角解消に対して自然な同型
\[
 \Theta_{B^{\rm op}}\simeq(\Theta_B)^{\rm op}
\]
があり、非負 tensor 冪とその積を合わせると
\[
 \Pi_3(B^{\rm op})\simeq\Pi_3(B)^{\rm op}
\]
という内部次数付き dg 代数の自然な quasi-isomorphism を得る。
したがって通常の次数零部分についても
\[
 C(B^{\rm op})\simeq C(B)^{\rm op}
\]
が成立する。全ての比較は任意基底変更と可換する。
\end{lemma}
\begin{proof}
二つの enveloping 代数の間の因子交換
\[
 (B^{\rm op})^e=B^{\rm op}\otimes B
   \xrightarrow{\sim}B\otimes B^{\rm op}=B^e
\]
は対角双加群と有限局所射影対角解消を対応させる。
この解消を enveloping 代数で双対化し、標準の左右交換によって再び
双加群と解釈すれば、inverse dualizing complex の表示した同型を得る。
cohomological shift $[2]$ はこの比較と可換する。

左右交換は tensor 積について反モノイド的であり、
\[
 (M\otimes_B^{\mathbf L}N)^{\rm op}
 \xrightarrow{\sim}
 N^{\rm op}\otimes_{B^{\rm op}}^{\mathbf L}M^{\rm op}
\]
を与える。複体で書けば交換には cohomological Koszul 符号を付ける。
これを各 tensor 冪へ反復適用すると、積の順序を逆にした
内部次数付き dg 代数の quasi-isomorphism となる。
次数零 cohomology 上ではこの符号は通常の反対代数の積を与え、
内部次数に新たな符号を挿入しない。
有限局所射影解消、双対、tensor と因子交換はいずれも基底変更に自然なので、
全ての同型は任意基底変更と可換する。
\end{proof}

\begin{corollary}[square完成の左右Noether性]
\label{cor:qa-relative-square-completion-two-sided-noetherian}
実際の相対 curve-type square avatar から構成した $\Lambda$ について、
$C(\Lambda)$ は各局所 affine chart とその許された基底変更上で
左右 Noetherian である。
\end{corollary}
\begin{proof}
右 Noether 性は既存の tails 認識による。
反転した $\Lambda^{\rm op}$ は
補題~\ref{lem:qa-square-foundation-opposite} により
全幾何 fiber が分類済み square 型である有限局所自由族なので、
補題~\ref{lem:qa-classified-fiber-finite-square-recognition} により
$C(\Lambda^{\rm op})$ は右 Noetherian である。
前補題の同型によりこれは $C(\Lambda)^{\rm op}$ の右 Noether 性であり、
$C(\Lambda)$ の左 Noether 性を意味する。
必要な局所主張は中央の faithfully flat 基底変更によって降下する。
一般の任意 Noetherian $k$-基底への強 Noether 性や、任意 AS/QP 表示との
同一視は、この系から追加で主張しない。
\end{proof}

% !TEX root = ../nc-quadric-surface-avatar.tex
% 2026-11-23 / 相対square完成と実際のtails認識。
\subsection{四冪等元による全整数展開}
\label{subsec:qa-relative-square-unfolding}

以下で $S$ は前節と同じ flat な曲線型 square avatar の基底であり，
まず実際の線束と四標準対象を持つ chart 上で考える．
有限代数を $\Lambda$，標準冪等元を $e_0,\ldots,e_3$ とし，
通常の相対 Serre 完成を
\[
 B=C(\Lambda)=\bigoplus_{n\ge0}B_n,
 \qquad B_n=H^0(\Theta_\Lambda^{\otimes_\Lambda^{\mathbf L}n})
\]
と書く．特に $B_0=\Lambda$ であり，その零次数代数を
$\OO_S^{\oplus4}$ に置き換えることはしない．

\begin{proposition}[canonical square の全整数展開]
\label{prop:qa-relative-square-unfolding}
$a=4p+i$, $b=4q+j$，$0\le i,j\le3$ に対して
\[
 A_{ba}=
 \begin{cases}
 e_jB_{q-p}e_i,&q\ge p,\\
 0,&q<p
 \end{cases}
 \tag{UF}\label{eq:qa-square-unfolding-components}
\]
と定め，積を $B$ の積から移す．これは $A_{aa}=\OO_S$，
$A_{ba}=0$ ($b<a$) の positively directed な
$\mathbb Z$-代数であり，同時の四頂点移動
$A_{b+4,a+4}=A_{ba}$ を持つ．
全成分は有限局所自由で，積を含めて任意の scheme 基底変更と可換する．
相対 heart の標準対象を $E_i$ とし，
\[
 E_{4p+i}=t^pE_i,\qquad t=\mathsf S^{-1}[2]
\]
と置くと，$b\ge a$ で
\[
 A_{ba}=\operatorname{Hom}(E_a,E_b)
\]
が全ての積を含めて成立する．

右加群圏には，基底変更と可換な完全同値
\[
 \operatorname{Gr}_{\mathbb Z}(A)\xrightarrow{\sim}
 \operatorname{Gr}(B),
 \qquad (\operatorname{Pack}M)_n e_i=M_{i-4n}
 \tag{PK}\label{eq:qa-square-packing}
\]
がある．これは有限生成性，有限表示性，torsion を保ち，
\[
 P_{4p+i}=e_{4p+i}A\longmapsto e_iB(p),
 \qquad B(p)_n=B_{n+p}
\]
となる．従って前節の相対 AZ 再構成の範囲では
$A$ の全ての右 row は noetherian であり，
\[
 \qgr A\simeq\qgr B\simeq\mathcal C_0,
 \qquad \pi P_a\longmapsto E_a
 \tag{UQ}\label{eq:qa-square-unfolded-heart}
\]
という実際の完全同値を得る．この比較は同じ範囲の任意基底変更と
coherent に整合する．
\end{proposition}
\begin{proof}
$q=p$ のとき，$e_j\Lambda e_i=0$ ($j<i$) なので
式~\eqref{eq:qa-square-unfolding-components} は positively directed である．
結合則と単位は $B$ のものから従う．相対 Serre tensor の
基底変更命題により $B_n$ は全て有限局所自由であり，
冪等元 corner はその直和因子だから同じ性質を持つ．
compact Morita の Serre 比較は
\[
 e_jB_ne_i=\operatorname{Hom}(E_i,t^nE_j)\qquad(n\ge0)
\]
を，右規約の軌道積を含めて与える．$t$ が自己同値なので
$b\ge a$ の表示式が従う．

右 $A$-加群の作用は $M_b\otimes A_{ba}\to M_a$ である．
式~\eqref{eq:qa-square-packing} において
\[
 M_{j-4n}\otimes e_jB_me_i\longrightarrow M_{i-4(n+m)}
\]
は次数 $n$ と $m$ の通常の右 graded 作用になる．
逆は $M_{4p+i}=N_{-p}e_i$ であり，両構成は互いに厳密に逆である．
標準射影についての式は成分を代入して確かめられる．
従って有限個の標準射影による表示を両向きに移せる．

$A$ 側の torsion は各同次元が十分負の頂点への作用で消えること，
$B$ 側の torsion は各同次元が $B_{\ge N}$ により消えることである．
対応する頂点差は $4m+j-i$ であり，$|j-i|\le3$ なので，
一方の十分遠方という条件は他方と同値である．
よって全 graded 加群の torsion 局所化も同値になり，
noetherian 対象に制限すると $\qgr$ の同値となる．
$B$ の右 noetherian 性から各 $e_iB(p)$ は noetherian であり，
$A$ の全右 row に同じことが従う．
最後に，$B$ の AZ 同値と通常の module-object 基底変更との比較は
前節の有限表示の余核による構成である．packing 自体は成分の並べ替え
だけなので，その coherent な比較をそのまま移せる．
\end{proof}

\subsection{係数束を残した実際のAS解消}
\label{subsec:qa-relative-square-simple-resolutions}

上の代数の幾何 fiber が square AS 型であることと，
相対的な単純解消が実際に存在し全基底変更に保たれることを区別する．
後者は以下の有限な relation 表示により証明できる．

\begin{proposition}[相対squareの単純解消とAS双対]
\label{prop:qa-relative-square-as-resolutions}
式~\eqref{eq:qa-square-unfolding-components} の $A$ は，
任意の頂点 $i$ において，係数が有限局所自由な実際の最小 AS 解消を持つ．
その四位相の重複度は定義~\ref{def:qa-as-square-rotated-minimal} の
square 表の通りであり，最左項は canonical な規約で $P_{i-4}$ である．
これらの解消，その左向きの双対解消，および
\[
 \mathcal E xt_A^q(S_i,A)=0\quad(q\ne3),
 \qquad
 \mathcal E xt_A^3(S_i,A)\simeq {}_\ell S_{i-4}
 \tag{RAG}\label{eq:qa-square-relative-as-gorenstein}
\]
は任意の scheme 基底変更に保たれる．
ここで $S_i$ は頂点 $i$ に $\OO_S$ を置いた標準右加群である．
特異な係数環を含め，代数の絶対大域次元が三であるとは主張しない．
\end{proposition}
\begin{proof}
\emph{有限関係表示．}
$D_0=\bigoplus_{r=0}^3\OO_Se_r$ と置き，有限 square quiver の
四つの矢束を
\[
 V_{10}=\mathcal H_{01},\quad V_{20}=\mathcal H_{02},\quad
 V_{31}=\mathcal H_{13},\quad V_{32}=\mathcal H_{23}
\]
とする．いずれも階数二である．二段階合成は実際に全射
\[
 (V_{31}\otimes V_{10})\oplus(V_{32}\otimes V_{20})
       \longrightarrow\mathcal H_{03}
\]
であり，その核を $R$ とする．全射性は各幾何 fiber の square 有限表示と
Nakayama から従い，標的が局所自由なので $R$ は階数四の部分束である．
この完全列は任意基底変更に保たれる．
$V$ を四矢束の $D_0$-双加群直和とすると
\[
 \Lambda=T_{D_0}(V)/(R).
\]
有限 quiver の道長は高々二であり，$R$ は source $0$ から sink $3$ の
成分だけにある．従って関係の重なりはなく，通常の関係 ideal は
実際に $R$ によって生成される．

この表示には任意基底変更と可換な有限相対対角解消
\[
 0\longrightarrow\Lambda\otimes_{D_0}R\otimes_{D_0}\Lambda
 \longrightarrow\Lambda\otimes_{D_0}V\otimes_{D_0}\Lambda
 \longrightarrow\Lambda\otimes_{D_0}\Lambda
 \longrightarrow\Lambda\longrightarrow0
 \tag{DR}\label{eq:qa-square-short-diagonal}
\]
がある．最初の微分は，道 $ba$ に対する
$b\otimes a\otimes1+1\otimes b\otimes a$ の和，次は
$x\otimes v\otimes y\mapsto xv\otimes y-x\otimes vy$ である．
道長が高々二で，二段階道の積の核がちょうど $R$ であることから
各頂点対の成分で完全性を直接確認できる．
あるいは関係と矢の局所枠で通常の bar 解消の可縮な道部分を消去すれば，
同じ複体を得る．全項は $\Lambda^e$ 上有限局所射影である．

\emph{通常のcompletionとrelation potential．}
$R^\vee$ を戻り矢 $3\to0$ として加え，
$R\hookrightarrow V\otimes_{D_0}V$ と恒等 tensor による
tautological な巡回類
\[
 W=\sum_\alpha c_\alpha r_\alpha
 \tag{RP}\label{eq:qa-square-relative-relation-potential}
\]
を取る．この表示は $R$ の局所枠 $r_\alpha$ と双対枠 $c_\alpha$ に依存しない．
元の矢の block 次数を零，戻り矢の block 次数を一とする．
式~\eqref{eq:qa-square-short-diagonal} を双対化し $[2]$ だけ移すと，
$\Theta_\Lambda$ の次数零の生成束は $R^\vee$，
次数 $-1$ の生成束は $V^\vee$，次数 $-2$ は四頂点の単位束になる．
その tensor dg 代数の次数零部分では，元の関係 $R$ と，
$V^\vee$ の微分である $\partial_VW$ が課される．
$\partial_{R^\vee}W=R$ だから
\[
 B=H^0\Pi_3(\Lambda)
   \simeq \operatorname{Jac}(V\oplus R^\vee,W)
 \tag{RJ}\label{eq:qa-square-relative-jacobian}
\]
を得る．これは有限な解消の実際の双対計算であり，
一般の変形 parameter に対する CY 性や，arrow-adic 完備化を用いない．
体上では \cite[Theorem 6.10]{KellerDeformedCY2011} の relation QP 計算に
一致する．その訂正で要求される ordinary ideal の生成条件も，上の
道長二の表示で満たされる．訂正の一次資料は
\url{https://webusers.imj-prg.fr/~bernhard.keller/publ/KellerCorrCYCompletion.pdf}
で確認できる．

\emph{全整数の単純複体．}
戻り矢を展開すると，非零の矢束 $\mathcal V_{ba}$ は
\[
\begin{array}{c|c}
 (b,a)&\mathcal V_{ba}\\\hline
 (4q+1,4q)&V_{10}\\
 (4q+2,4q)&V_{20}\\
 (4q+3,4q+1)&V_{31}\\
 (4q+3,4q+2)&V_{32}\\
 (4q+4,4q+3)&R^\vee
\end{array}
\]
である．全矢は正の頂点差を持ち，$W$ の各巡回項の頂点差の和は四である．
通常の potential 複体は，基底を選ばず
\[
\begin{split}
0\longrightarrow P_{i-4}
 &\longrightarrow\bigoplus_j\mathcal V_{j,i-4}^{\vee}\otimes P_j
 \longrightarrow\bigoplus_j\mathcal V_{ij}\otimes P_j\\
 &\longrightarrow P_i\longrightarrow S_i\longrightarrow0
\end{split}
\tag{SR}\label{eq:qa-square-relative-simple-resolution}
\]
と書ける．右の微分は矢の評価，中間微分は cyclic derivative の
最後の矢による分解，左の微分は矢の余評価である．
合成が零であることは Jacobian 関係と巡回恒等式から従う．
係数束を枠で自明化すると，四位相はちょうど既存の square 最小分解表となる．

各幾何 fiber では，式~\eqref{eq:qa-square-relative-jacobian} は元の
square AS 代数の relation 表示であり，この potential 複体は完全である．
これは体上の全整数 AS 分類と，
補題~\ref{lem:qa-as10-ginzburg-concentration} の square の場合の
標準単純解消による．任意の固定した右頂点成分を取れば，
式~\eqref{eq:qa-square-relative-simple-resolution} の増大複体は
有限局所自由層の有界複体になる．全導来幾何 fiber が零なので複体自体が
零と準同型であり，局所的に分裂している．
従って式は実際に完全で，任意基底変更後にも完全である．
最小性は各微分の成分が正の頂点差を持つことから従う．

最後にこの有限射影解消を $A$ で双対化する．
各左頂点成分も有限局所自由層の有界複体であり，最終項から
標準左単純 ${}_\ell S_{i-4}$ への canonical augmentation を持つ．
幾何 fiber の AS 双対性により，この増大複体の全 fiber は完全である．
同じ perfect-fiber 判定で式~\eqref{eq:qa-square-relative-as-gorenstein} と
その任意基底変更を得る．左単純の解消はこの双対複体そのものである．
この議論に左 noetherian 性を先取りしていない．
\end{proof}

\begin{remark}[頂点線束と周期の標識]
\label{rem:qa-square-unfolding-vertex-lines}
実際の標準対象を $E_a'=E_a\otimes L_a$ に変更すると，
対応する全整数代数は
\[
 A'_{ba}=L_b\otimes A_{ba}\otimes L_a^{\vee}
 \tag{VG}\label{eq:qa-square-vertex-line-gauge}
\]
となる．右加群の同値は $M_a\mapsto M_a\otimes L_a^{\vee}$ である．
したがって線束を自明化せずに，有限表示，torsion，および tails の比較を
移せる．ただしこの操作は，選択なしの厳密な代数同型と同じではない．

有限四対象の基底線束変更を canonical $t$ で延長する場合は
$L_{4p+i}=L_i$ であり，四周期も保たれる．
一方，周期代表を $u=t\otimes N$ に変更すると
$L_{4p+i}=L_i\otimes N^p$ となり，block 次数 $m$ の成分は
さらに $N^m$ を tensor される．
この規約での単純解消の最左項は
\[
 P'_{i-4}\otimes L_i\otimes L_{i-4}^{\vee}
\]
であり，その係数線束を黙って消去してはならない．

本節で構成した $A$ は実際の canonical $t$ と有限 $\Lambda$ からの
unfolding である．別途与えた相対 AS $\mathbb Z$-代数族と同一視するには，
有限 End の同型に加えて，全整数対象を canonical Serre transport と
比較する同型，または式~\eqref{eq:qa-square-vertex-line-gauge} の線束を
含む比較を示す．この条件を省いた全 AS/QP presentation groupoid の
相対同値を主張していない．
\end{remark}

\begin{corollary}[展開後の左右Noether性]
\label{cor:qa-relative-square-unfolding-two-sided-noetherian}
前節と同じ局所 affine chart とその許された基底変更上で，
$A$ の全ての右 row および左 column は noetherian である．
\end{corollary}
\begin{proof}
右 row については命題~\ref{prop:qa-relative-square-unfolding} による．
系~\ref{cor:qa-relative-square-completion-two-sided-noetherian} によって
$B^{\mathrm{op}}$ も右 noetherian である．その冪等元を
$f_i=e_{3-i}$ と並べ直して同じ packing を用いると，
対応する全整数代数は
\[
 A^{\dagger}_{ba}=A_{3-a,\,3-b}
\]
であり，積は元の積を反転したものである．
その全右 row が noetherian なので，元の $A$ の全左 column も
noetherian となる．左側の有限性にはここで実際に反対代数の
相対 $H_2$ 認識と AZ 再構成を使っており，
単純 AS 双対解消だけから推論していない．
\end{proof}

% !TEX root = ../nc-quadric-surface-avatar.tex
% 2026-11-23 / 全基底上の有限表示とcanonicalな相対認識。
\section{有限square表示と全H2の相対foundation}
\label{sec:qa-relative-square-foundation-stack}

\subsection{二つの積による有限squareスタック}
\label{subsec:qa-square-normal-foundation-stack}

以下の有限代数の構成は任意の可換基底上で行える。
$\mathfrak H_2$ との比較には、本稿で既に固定した標数零の基礎体と
局所有限型基底の site を用いる。
四つの階数二の束 $A,B,C,D$ と階数四の束 $W$、および二つの積
\[
 f:C\otimes A\longrightarrow W,\qquad
 g:D\otimes B\longrightarrow W
\]
を取る。その随伴
\[
 \alpha_f:A\longrightarrow C^\vee\otimes W,\qquad
 \alpha_g:B\longrightarrow D^\vee\otimes W
\]
がともに部分束の包含であることを要求する。
対応する四頂点代数 $\Lambda_\square$ の標準射影を
$(E_0,E_1,E_2,E_3)$ と書き、非対角の Hom 束を
\[
 H_{01}=A,\quad H_{02}=B,\quad H_{13}=C,\quad
 H_{23}=D,\quad H_{03}=W
\]
とする。二つの中間頂点は直交し、積は $f,g$ だけで決まる。
長さ三の非定数な道はないので、追加の結合律方程式はない。
対角を単位によって構造層と同定すれば、この代数は階数十六である。

これらの束の基底空間を固定し、二つの随伴が階数二を持つ開部分を
\[
 V_\square\subset
 \operatorname{Hom}(C\otimes A,W)
 \oplus\operatorname{Hom}(D\otimes B,W)
 \simeq\mathbb A^{32}
\]
と書く。基底変更群と頂点スカラー部分群は
\[
 G_\square=\operatorname{GL}(A)\times\operatorname{GL}(B)
 \times\operatorname{GL}(C)\times\operatorname{GL}(D)
 \times\operatorname{GL}(W),
\]
\[
 K_\square=
 \left\{\left(\lambda_1,\lambda_2,
 \frac{\lambda_3}{\lambda_1},
 \frac{\lambda_3}{\lambda_2},\lambda_3\right)\right\}
 \simeq\mathbb G_m^3.
\]
各成分は対応する束のスカラー行列を表す。
$K_\square$ はパラメータ上自明に作用し、頂点零のスカラーを一とした
対角単元の内部共役に一致する。
\[
 \mathfrak S^{\rm lin}_\square=[V_\square/G_\square],\qquad
 \mathfrak S_\square=[V_\square/(G_\square/K_\square)]
\]
と置く。前者は頂点を保持した有限代数と全同型を、後者はその
頂点スカラーによる rigidification を表す。

\begin{theorem}[有限squareとnormal foundationの相対同値]
\label{thm:qa-square-normal-foundation-stack-equivalence}
論文Bの二写像表示
\[
 F:P\otimes W_1\to W_2,\qquad G:Q\otimes W_2\to W_3
\]
について、$U_0$ は $G$ が全射、$\alpha_F$ が部分束の包含である開部分とする。
階数は $\operatorname{rk}(W_1,W_2,W_3;P,Q)=(2,4,6;2,2)$ である。
任意基底変更と可換する同値
\[
 \Phi^{\rm lin}:\mathfrak S^{\rm lin}_\square
       \xrightarrow{\sim}[U_0/\mathcal G]
\]
が、次の式によって与えられる。
\[
 \begin{split}
 (W_1,W_2,W_3;P,Q)
   &=(B,W,\operatorname{coker}\alpha_f;D,C^\vee),\\
 F&=g,\qquad
 G:C^\vee\otimes W\twoheadrightarrow
             \operatorname{coker}\alpha_f.
 \end{split}
 \tag{SN}\label{eq:qa-square-to-normal-data}
\]
その逆は
\[
 (A,B,C,D,W)=(\ker G,W_1,Q^\vee,P,W_2),
 \tag{NS}\label{eq:qa-normal-to-square-data}
\]
および $g=F$、$f$ は包含 $\ker G\hookrightarrow Q\otimes W_2$
の随伴、という構成である。

この同値は $K_\square$ と論文Bの $K$ を同定し、
\[
 \Phi:\mathfrak S_\square\xrightarrow{\sim}\mathfrak F_0
\]
へ降下する。さらに中間頂点 $1,2$ の交換を $\sigma$ とすると
$\Phi\sigma\simeq\tau\Phi$ は coherent であり、
\[
 [\mathfrak S_\square/S_2]\xrightarrow{\sim}\mathfrak Q_{\rm fnd}.
\]
\end{theorem}
\begin{proof}
両随伴の部分束条件によって式~\eqref{eq:qa-square-to-normal-data} の
余核は階数六のベクトル束であり、商写像 $G$ は全射である。
$\alpha_F=\alpha_g$ なので、このデータは $U_0$ に属する。
逆に $U_0$ の全射 $G$ の核は階数二の部分束であり、
式~\eqref{eq:qa-normal-to-square-data} では
$\alpha_f$ がこの包含、$\alpha_g=\alpha_F$ となる。
従って逆データも $V_\square$ の束値の点である。

二つの構成を合成すると、$\ker(\operatorname{coker}\alpha_f)
=\operatorname{im}\alpha_f\simeq A$、
$(Q\otimes W_2)/\ker G\simeq W_3$ という標準同型により元のデータを得る。
縮約と随伴もこの同型の下で元の積に戻る。
核・余核の列はベクトル束の短完全列であり局所的に分裂するため、
全ての構成は任意の、非平坦なものを含む基底変更と可換する。
同型に関する自然性から、これは全群oidの相互逆同値である。

この変換は実際の相対 perfect mutation でもある。
元の四項に対して標準 dg 増強で
\[
 Z_1=\operatorname{Cone}(E_1\longrightarrow C^\vee\otimes E_3)
\]
を取り、$(E_0,E_2,E_3,Z_1)$ を考える。
随伴の部分束条件と二つの中間項の直交性によって、その forward Hom は
次数零に集中し、階数表は
\[
 \begin{pmatrix}1&2&4&6\\0&1&2&4\\0&0&1&2\\0&0&0&1\end{pmatrix}
\]
となる。特に
\[
 \operatorname{Hom}(E_0,Z_1)=\operatorname{coker}\alpha_f,
 \quad\operatorname{Hom}(E_3,Z_1)=C^\vee,
 \quad\operatorname{Hom}(E_2,Z_1)=C^\vee\otimes D.
\]
最後の同型は $E_3$ を経由する積そのものである。
逆向き Hom は零、各対角は構造層であり、三角から $E_1$ を回復できるので
この四項は相対 full strong exceptional foundation である。
その積は式~\eqref{eq:qa-square-to-normal-data} と一致する。
逆に normal 四項 $(\mathcal O,X,M,Z)$ から
\[
 Y=\operatorname{Cone}(Q\otimes M\longrightarrow Z)[-1]
\]
を取れば $(\mathcal O,Y,X,M)$ を得る。
論文Bの\cite[\texttt{prop:ruling-involution-mutation}]{QuadricModuli2026}
と同じ Hom 計算が式~\eqref{eq:qa-normal-to-square-data} を与える。
ここでは Euler 行列だけを一致させるのではなく、cone と全ての積を同定した。

square の頂点スカラーを $(1,\lambda_1,\lambda_2,\lambda_3)$ とすれば、
cone $Z_1$ 上のスカラーは $\lambda_1$ である。
従って normal 側の頂点スカラー $(1,\mu_1,\mu_2,\mu_3)$ との関係は
\[
 (\mu_1,\mu_2,\mu_3)=(\lambda_2,\lambda_3,\lambda_1).
\]
よって同値は二つの中心部分群を同定し、rigidification へ降下する。

middle swap は $(A,C,f)$ と $(B,D,g)$ を交換して $W$ を固定する。
その normal 像は
\[
 (A,W,\operatorname{coker}\alpha_g;C,D^\vee)
\]
と積 $f$ および商写像であり、これは論文Bの $\tau$ の式に一致する。
二つの短完全列
\[
 0\to A\xrightarrow{\alpha_f}C^\vee\otimes W
          \to\operatorname{coker}\alpha_f\to0,\qquad
 0\to B\xrightarrow{\alpha_g}D^\vee\otimes W
          \to\operatorname{coker}\alpha_g\to0
\]
の交換として記述すれば、二回交換と恒等の同型は自動的に整合する。
頂点スカラー上でも $\lambda_1\leftrightarrow\lambda_2$ は
$\mu_1\leftrightarrow\mu_3$ に一致する。
従って $S_2$ 商の同値を得る。
\end{proof}

\begin{corollary}[三十二次元の線形表示]
\label{cor:qa-square-linear-presentation-normal-stack}
四つの階数二の基底空間を同じ $k^2$ と同定すると、中間頂点の交換は
$V_\square\subset\mathbb A^{32}$ 上の線形作用であり、群の四つの
$\operatorname{GL}_2$ 因子を対応して交換する。
従って
\[
 \mathfrak Q_{\rm fnd}\simeq
 \left[V_\square\Big/
 \left(\left(\big((\operatorname{GL}_2)^4\times\operatorname{GL}_4\big)
 /\mathbb G_m^3\right)\rtimes S_2\right)\right].
\]
$\mathfrak S_\square$ とこの商は有限型・滑らかで次元三である。
\end{corollary}
\begin{proof}
上の $S_2$ 同変同値を、固定基底による線形表示で書き直す。
$V_\square$ は滑らかな三十二次元空間の開部分であり、
$\dim G_\square=4\cdot4+16=32$、$\dim K_\square=3$ である。
従って商の次元は三である。
\end{proof}

\subsection{二経路による生成と内在的$H_2$開部分}

\begin{proposition}[二積の同時全射性は追加の開条件]
\label{prop:qa-square-joint-generation-open}
二つの積の直和が全射である条件
\[
 V_\square^{\rm gen}
 =\left\{f\oplus g:(C\otimes A)\oplus(D\otimes B)
                 \twoheadrightarrow W\right\}
\]
は $G_\square\rtimes S_2$ 不変の開条件である。
この条件は $U_0$ の二条件からは従わない。
さらに各積の核が非零の分解可能テンソルを含まないと仮定しても、
この不足は解消されない。
\end{proposition}
\begin{proof}
全射性は階数四の最大階数条件なので開であり、基底変更と交換に不変である。
反例として $W=\langle x_0,x_1,x_2,x_3\rangle$ を取り、二つの積を同じ
行列
\[
 \begin{pmatrix}x_0&x_1\\x_1&x_2\end{pmatrix}
\]
で与える。例えば $f(c_0a_0)=x_0$、$f(c_0a_1)=f(c_1a_0)=x_1$、
$f(c_1a_1)=x_2$ と定め、$g$ も同様とする。
両随伴は単射であり、両積の階数は三である。
各核は $c_0\otimes a_1-c_1\otimes a_0$ 型のテンソルで生成され、
このテンソルの行列階数は二なので分解可能でない。
しかし二つの像はともに $\langle x_0,x_1,x_2\rangle$ であり、
直和写像の余核は $kx_3$ である。
定理の同値によって、これは実際の $U_0$ 点も与える。
\end{proof}

\begin{corollary}[内在的$H_2$のsquare表示]
\label{cor:qa-h2-square-foundation-open}
定理~\ref{thm:qa-intrinsic-h2-open-in-foundation-stack} の
$\tau$ 不変開部分 $\mathfrak F_{\HH}\subset\mathfrak F_0$ に対し
\[
 \mathfrak S_{\HH}=\Phi^{-1}(\mathfrak F_{\HH})
\]
と置く。これは $\sigma$ 不変の有限型開部分であり、
\[
 \mathfrak H_2^{\rm rul}\simeq\mathfrak S_{\HH},\qquad
 \mathfrak H_2\simeq[\mathfrak S_{\HH}/S_2].
\]
さらに
\[
 \mathfrak S_{\HH}\subset
 [V_\square^{\rm gen}/(G_\square/K_\square)].
\]
したがって全ての相対 $H_2$ 族で、二つの ruling を経由する積は合わせて
$\operatorname{Hom}(\mathcal O,M)$ を生成する。
この主張は非被約基底にも適用される。
\end{corollary}
\begin{proof}
開部分と全相対群oidの同値は、既存の $\Xi$ と前定理の同値を合成して得る。
残る同時全射性は幾何 fiber で確認すればよい。
各 $H_2$ 点で二つの積は二つの MCM 型 $X,Y$ を経由する積である。
論文Aの\cite[\texttt{cor:normal-cubic-defect}]{QuadricIntrinsic2026}
により、いずれの階数も三以上である。
一方が階数四なら全射性は明らかである。
双方が階数三なら、$p=\Omega(1)$ が二型を固定する場合は
\cite[\texttt{cor:w-distinct-supports}]{QuadricIntrinsic2026}、
交換する場合は
\cite[\texttt{cor:standard-supports-distinct}]{QuadricIntrinsic2026}
によって、二つの像の超平面が異なる。
従ってその和は四次元空間全体となる。
これは二元集合上の $p$ の全ての場合を尽くす。
古典型でも通常の二つの tensor 積はそれぞれ同型である。
よって任意の相対族の幾何 fiber で最大階数となり、
ベクトル束間の写像の全射開条件から基底上でも全射である。
\end{proof}

\begin{remark}[古典型とこの表示の射程]
古典的な四項
\[
 (\mathcal O,\mathcal O(1,0),\mathcal O(0,1),\mathcal O(1,1))
 \quad\text{on }\mathbb P^1\times\mathbb P^1
\]
も $\mathfrak S_{\HH}$ の点である。
二つの積はともに通常の二次元空間の tensor 積から四次元空間への同型であり、
その mutation は
$(\mathcal O,\mathcal O(0,1),\mathcal O(1,1),\mathcal O(1,2))$
という normal foundation になる。
従ってこの有限代数表示は、曲線 avatar の平坦族を先に指定せずに、
古典型と曲線型を同じ階数十六の有限スタックで扱う。
再構成された幾何データの次元が古典型で上がることを、有限 Hom 束の
階数変化と取り違える必要はない。

ここで明示したのは有限代数の全スタック同値と、既存の $H_2$ 開埋込の
その表示への移送である。
$V_\square^{\rm gen}$ の全点が $H_2$ であるとは主張しない。
その中の真正な開部分は定理~\ref{thm:qa-square-explicit-h2-open}で認識し、
系~\ref{cor:qa-square-h2-twenty-polynomials}で有限個の明示多項式により記述する。
一般の有限データで、mutation が特定の既知の Abel heart の中の短完全列となる
ことも主張していない。有限同値はまず perfect 圏で成立し、相対 heart は
$\mathfrak F_{\HH}$ 上で既存の $\Xi$ によって認識される。
線形 mutation と kernel--cokernel の一般原理自体の新規性は主張せず、
本稿では指定した二つの表示スタック、中心 rigidification、二 ruling の作用を
全相対群oidで明示比較した結果として用いる。
\end{remark}

% !TeX root = ../nc-quadric-avatar.tex
% 2026-11-23 / v4.18。A32 の二積から全 H2 開を有限式で判定する。

\subsection{随伴行列と交差縮約による明示的な$H_2$開}
\label{subsec:qa-square-explicit-h2-open}

本項では標数零の代数閉体 $k$ を固定する。
前項と同じく $A,B,C,D$ は階数二、$W$ は階数四とし、
\[
 U=C\otimes A,\qquad V=D\otimes B,\qquad
 f:U\to W,\quad g:V\to W
\]
を取る。固定基底で全二積の空間は $\mathbb A^{32}$ である。
ここでは $f,g$ の可逆性を最初に仮定しない。

\paragraph{線束を保持した随伴行列.}
\[
 L_f=\det W\otimes(\det U)^\vee,\qquad
 L_g=\det W\otimes(\det V)^\vee
\]
と置く。通常の随伴行列は束値では
\[
 \operatorname{adj}(f):W\to U\otimes L_f,
 \qquad \operatorname{adj}(f)f=1_U\otimes\det f
\]
という標準写像であり、$g$ も同様である。
従って四因子テンソル
\[
 T_f=\operatorname{adj}(f)g\in
 C\otimes A\otimes D^\vee\otimes B^\vee\otimes L_f,
 \qquad T_g=\operatorname{adj}(g)f
\]
が基底なしで定まる。
これを次の交差縮約として読む。
\[
 \begin{split}
 M_f &:C^\vee\otimes D\longrightarrow A\otimes B^\vee\otimes L_f,\\
 M_g &:D^\vee\otimes C\longrightarrow B\otimes A^\vee\otimes L_g.
 \end{split}
 \tag{CF}\label{eq:qa-square-cross-flattenings}
\]
例えば $T_f$ の係数を $(T_f)_{ca,db}$ とすれば
$(M_f)_{ab,cd}=(T_f)_{ca,db}$ とする。
添字は各因子の指定順で辞書式に並べる。
さらに線束を再び終域へ移した転置を
\[
 \begin{split}
 M_f^{\mathrm t}&:A^\vee\otimes B\longrightarrow C\otimes D^\vee\otimes L_f,\\
 M_g^{\mathrm t}&:B^\vee\otimes A\longrightarrow D\otimes C^\vee\otimes L_g
 \end{split}
\]
と書く。すなわち $M_f^{\mathrm t}=M_f^\vee\otimes L_f$ であり、$g$ も同様である。

各始域は二つの階数二束の tensor 積である。
写像 $M$ に対し、$\mathcal K(M)$ は、その kernel が各幾何 fiber で
非零の分解可能テンソルを含まないという条件を表す。
これは始域の二つの射影束の積上で、tautological な分解可能直線の
像が零になる incidence を作り、その proper 射影の像を除いた開条件である。

\begin{definition}[四つの交差縮約条件]
\label{def:qa-square-cross-geometric-open}
次の開部分を定める。
\[
 V_{\mathrm{cross}}=
 \mathcal K(M_f)\cap\mathcal K(M_f^{\mathrm t})
 \cap\mathcal K(M_g)\cap\mathcal K(M_g^{\mathrm t})
 \subset\mathbb A^{32}.
 \tag{CG}\label{eq:qa-square-four-cross-open}
\]
この条件は線束 $L_f,L_g$ の自明化にも、各 tensor 因子の基底にも依存しない。
従って $G_\square\rtimes S_2$ 不変であり、任意基底変更と整合する。
\end{definition}

\begin{lemma}[階数一テンソルの交差縮約]
\label{lem:qa-rank-one-cross-kronecker}
$T=\kappa\otimes\psi\in(C\otimes A)\otimes(D^\vee\otimes B^\vee)$ とする。
その交差縮約は
\[
 M_T=(\kappa:C^\vee\to A)\otimes(\psi:D\to B^\vee)
\]
である。$\mathcal K(M_T)$ が成立する必要十分条件は、
二つの二次正方行列 $\kappa,\psi$ がともに可逆であることである。
このとき $M_T$ と $M_T^{\mathrm t}$ はともに可逆である。
\end{lemma}
\begin{proof}
二写像が可逆なら tensor 積も可逆である。
一方、例えば $\kappa$ が特異なら、その非零 kernel 元 $c^\vee$ と
任意の非零 $d\in D$ の tensor が $M_T$ の kernel に入る。
$\psi$ が特異な場合も同様である。
\end{proof}

\begin{lemma}[可逆な二積領域における四因子条件]
\label{lem:qa-cross-invertible-cubic}
$f$ が可逆である幾何点で、$T=f^{-1}g$ を
\[
 w_T\in A\otimes C\otimes D^\vee\otimes B^\vee
\]
と読む。このとき式~\eqref{eq:qa-square-four-cross-open} は、$w_T$ が
geometric quintuple を与えることと同値である。
$g$ が可逆の場合も、$f,g$ を交換した同じ主張が成立する。
\end{lemma}
\begin{proof}
前項の交換した normal 窓では
$(W_1,P,Q,R)=(A,C,D^\vee,B)$ であり、
$W\simeq C\otimes A$ の下で $R\hookrightarrow Q\otimes W$ は
$g$ の随伴そのものである。従ってこの $w_T$ は論文Bの cubic chart の
関係 tensor に一致する。
geometric 条件は四つの隣接二因子の縮約が全ての非零分解可能元上で
非零であることを要求する
\cite[Definition 5.3.1]{VdB2011}。
この tensor の四条件は、$\ker T$、$\ker T^{\mathrm t}$、
交差縮約 $\operatorname{Cr}(T)$ の kernel とその転置の kernel が
それぞれ分解可能元を含まないことである。
$M_f=(\det f)\operatorname{Cr}(T)$ なので、後二条件は直接
$\mathcal K(M_f),\mathcal K(M_f^{\mathrm t})$ に一致する。

まず $g$ も可逆なら前二条件は空である。
残る交差条件が $T^{-1}$ の交差条件と同値であることを確認する。
$c^\vee\otimes d$ が $\operatorname{Cr}(T)$ の kernel に入ることは
\[
 T(kd\otimes B)\subset(\ker c^\vee)\otimes A
\]
と同値である。両辺は二次元であり $T$ が可逆なので、包含は等号となる。
逆写像を取れば、これは $T^{-1}$ の同じ種類の ruling 二平面条件である。
転置の分解可能 kernel についても
$T(D\otimes kb)\subset C\otimes\ker a^\vee$ を用いれば同様である。
従って $M_g$ とその転置の条件はこの場合に追加制約を与えない。

次に $\operatorname{rk}g=3$ とする。
\[
 \operatorname{adj}(g)=\kappa_g\otimes\lambda_g,
 \quad k\kappa_g=\ker g,\quad
 k\lambda_g=\operatorname{Ann}(\operatorname{im}g)
\]
と書くと
\[
 \operatorname{adj}(g)f=\kappa_g\otimes(\lambda_gf).
\]
前補題により、$M_g$ の条件は $\kappa_g$ と $\lambda_gf$ が
二次行列として可逆であることに等しい。
これらは正に $\ker T$ と $\ker T^{\mathrm t}$ を生成する tensor である。
同時に $M_g^{\mathrm t}$ の条件も満たされる。

$\operatorname{rk}g\le2$ なら $\operatorname{adj}(g)=0$ なので
式~\eqref{eq:qa-square-four-cross-open} は成立しない。
また $\mathbb P(\ker T)$ は少なくとも直線を含み、代数閉体上で
その直線は Segre 二次曲面と交わるため、$w_T$ も geometric でない。
以上で全場合を尽くす。
\end{proof}

\subsection{両階数三の支持conicと中央二重平方}

\begin{proposition}[両階数三における有限認識]
\label{prop:qa-double-rank-three-conic-recognition}
$\operatorname{rk}f=\operatorname{rk}g=3$ とする。
式~\eqref{eq:qa-square-four-cross-open} が成立する必要十分条件は、
以下の三条件である。
\begin{enumerate}
\item $\ker f$ と $\ker g$ の各生成 tensor は二次行列として可逆である。
\item $W_f=\operatorname{im}f$ と $W_g=\operatorname{im}g$ は異なる。
      従って $W_f+W_g=W$、$U_0=W_f\cap W_g$ は二次元である。
\item $f^\vee:W_f^\vee\hookrightarrow C^\vee\otimes A^\vee$ の像の行列式を
      $q_f$ とし、$q_g$ も同様に定める。
      これらは smooth な ternary conic を与え、
      \[
       q_f|_{\operatorname{Ann}_{W_f^\vee}(U_0)}\ne0,
       \qquad
       q_g|_{\operatorname{Ann}_{W_g^\vee}(U_0)}\ne0.
      \]
\end{enumerate}
このとき元の二積データは、全ての積を保つ基底変更の下で
二つの非退化 binary quadratic $h_f,h_g\in\operatorname{Sym}^2U_0$ による
中央二重平方環
\[
 B_h=k[u,v]\langle c,d\rangle/(c^2-h_f(u,v),d^2-h_g(u,v))
\]
の有限 square foundation と同型である。
特に真正な noncubic $H_2$ foundation である。
\end{proposition}
\begin{proof}
随伴行列を
$\operatorname{adj}(f)=\kappa_f\otimes\lambda_f$ と書く。
ここで $\lambda_f$ は $\operatorname{Ann}(W_f)$ を張る。
すると
\[
 T_f=\kappa_f\otimes(\lambda_fg).
\]
前補題により $M_f$ の条件は両因子の行列式の非消滅と同値である。
もし $W_f=W_g$ なら第二因子は零なので、この場合は除かれる。
異なる場合には $\lambda_f|_{W_g}$ は $U_0$ の annihilator 線を張り、
\[
 \det(\lambda_fg)=q_g(\lambda_f|_{W_g}).
\]
$f,g$ を交換した式も合わせれば、表示した三条件に一致する。
転置の条件は Kronecker 積が可逆となるため自動的に成立する。

kernel tensor の可逆性から conic の smooth 性が従うことも確認する。
可逆 kernel を二因子の基底変更で標準的な非退化 tensor に直せば、
その annihilator は二次行列の trace 零 hyperplane と同値である。
この三次元空間への行列式の制限は非退化二次形式である。
従って $q_f,q_g$ はともに非退化である。

有限代数を中央二重平方のものに直接同定する。
共通空間の基底を $u,v$ とし、$W_f=U_0\oplus kc$ と書く。
$f$ の積の行列を $M(c,u,v)$ とすると、$c$ の係数行列は
$q_f|_{\operatorname{Ann}(U_0)}\ne0$ によって可逆である。
左右基底変更によりこれを単位行列に直し、次に $c$ へ $u,v$ の一次式を
加えることで
\[
 M=cI_2+N(u,v),\qquad\operatorname{tr}N=0
\]
とできる。このとき
\[
 \det M=c^2+\det N(u,v)
\]
であり、ternary form の非退化性から binary form $\det N$ も非退化である。
$N$ の二次元係数空間は $\mathfrak{sl}_2$ の trace pairing について非退化なので、
その直交補は可逆な trace 零行列 $H$ により生成される。
$H$ の二固有値は異なるため、定数共役によって
$H=\operatorname{diag}(\lambda,-\lambda)$ とできる。
直交補条件から $N$ は off-diagonal となり、
\[
 M=\begin{pmatrix}c&\ell_1(u,v)\\\ell_2(u,v)&c\end{pmatrix},
 \qquad h_f=\ell_1\ell_2
\]
を得る。$h_f$ の非退化性から $\ell_1,\ell_2$ は一次独立である。
$W_g=U_0\oplus kd$ についても独立に同じ操作を行えるため、
共通の $u,v$ を変えずに $g$ を同じ形にすることができる。

得られた二行列は、中央二重平方環の二つの標準行列分解から得る
$\operatorname{Hom}(\mathcal O,X)\otimes\operatorname{Hom}(X,M)$ の積である。
例えば $\left(\begin{smallmatrix}c&\ell_1\\-\ell_2&-c\end{smallmatrix}\right)$
の片側の基底を符号変更すれば上の行列になる。
論文Bの\cite[\texttt{eq:nchh-plane} の直前の行列分解計算]{QuadricModuli2026}
は normal foundation の同じ積を計算しており、前項の逆 mutation が
対応する square を与える。
有限 square の全積はこの二行列だけで決まるため、これは単なる
同じ conic の認識ではなく、頂点と全積を保つ有限代数の同型である。
非退化二形式の中央二重平方環が noncubic $H_2$ を与えることには
\cite[\texttt{thm:ore-anticommuting-central},\texttt{thm:noncubic-classification}]
{QuadricIntrinsic2026} を用いる。
\end{proof}

\begin{corollary}[両階数三のscheme壁における相対正規形]
\label{cor:qa-double-rank-three-scheme-normal-form}
\[
 Z_{33}=V(\det f,\det g)\subset V_{\mathrm{cross}}
\]
を閉部分スキームとする。その任意の基底変更上で、kernel、image と共通部分
$U_0$ は所定階数のベクトル束となる。
\'etale 局所に前命題の中央二重平方正規形が成立し、
二つの binary form の判別式は可逆である。
基底の被約性は不要である。
\end{corollary}
\begin{proof}
四つの交差条件は各 $f,g$ の階数が三以上であることを強制する。
三次 minor を可逆化すると $f$ は $\operatorname{diag}(I_3,a)$ の形となる。
$\det f=0$ が基底環の元として成立するため $a=0$ であり、$g$ も同様である。
従って kernel と image は束となる。
前命題の Kronecker 積の可逆性は基底上でも成立し、二つの像の和は全体である。
その共通部分は階数二の部分束となる。

前証明の係数行列の正規化と trace の除去は一般の係数環上で行える。
非退化な二平面の直交補行列 $H$ は局所的に生成元を持ち、$\det H$ は単元である。
$\lambda^2=-\det H$ の有限 \'etale 被覆上で、
$(I\pm\lambda^{-1}H)/2$ による固有線束の分解を取り、各線束を局所自明化する。
これで同じ off-diagonal 正規形を得る。
両線形式の係数行列の行列式は単元なので、二形式は非退化である。
全操作は有限代数の全積を保持する。
\end{proof}

\begin{theorem}[全$H_2$開の有限明示判定]
\label{thm:qa-square-explicit-h2-open}
四交差条件の開部分は
\[
 V_{\mathrm{cross}}\subset V_\square^{\rm gen}
\]
を満たし、前項の同値を通じて
\[
 [V_{\mathrm{cross}}/(G_\square/K_\square)]
   =\mathfrak S_{\HH}\simeq\mathfrak H_2^{\rm rul}.
\]
従って
\[
 \mathfrak H_2\simeq
 \left[V_{\mathrm{cross}}\Big/
       \big((G_\square/K_\square)\rtimes S_2\big)\right].
\]
これは非被約有限型基底を含む全相対 groupoid の同値である。
\end{theorem}
\begin{proof}
まず幾何点で判定する。
$f$ または $g$ の階数が二以下なら、その随伴行列が零となるため交差条件に反する。
従って二つの階数は三以上である。
その場合は随伴 $A\to C^\vee W$ と $B\to D^\vee W$ が単射である。
実際、例えば最初の kernel 元 $a\ne0$ があれば二次元空間 $C\otimes ka$
が $\ker f$ に入ってしまう。
一方が可逆なら $f\oplus g$ は全射であり、両階数三の場合も
前命題によって二つの像の和は $W$ となる。
従って $V_{\mathrm{cross}}\subset V_\square^{\rm gen}$ である。

少なくとも一方が可逆なら、補題~\ref{lem:qa-cross-invertible-cubic} により
対応する normal 窓が geometric cubic quintuple を与える。
論文Bの\cite[\texttt{prop:cubic-foundation-open}]{QuadricModuli2026} と
既存の cubic 逆構成によって、元の有限 square は真正な $H_2$ foundation である。
この cubic 入力は linear な古典型も含む
\cite[Lemma 5.3.3, Theorem 5.5.10]{VdB2011}。
両階数三なら前命題の直接正規形が同じ結論を与える。

逆に真正な $H_2$ square を取る。
二つの合成階数は三または四である。
一方が四なら、その ruling の normal helix は cubic であり
\cite[\texttt{thm:cubic-coverage},\texttt{cor:normal-cubic-defect}]
{QuadricIntrinsic2026}、その四因子テンソルは geometric である。
補題~\ref{lem:qa-cross-invertible-cubic} により四交差条件を得る。
両階数三なら、noncubic の体上の分類により、構造対象と二つの ruling を保って
非退化二形式の中央二重平方表示に帰着する。
その二つの支持 conic は $c^2-h_f$ と $d^2-h_g$ であり、
共通 $\langle u,v\rangle$ の annihilator 上の値は単元である。
従って前命題により四交差条件が成立する。
ここで使う noncubic 分類は有限代数候補の domain 性を仮定するためではなく、
既に $H_2$ に属する点の必要性にだけ用いた。

以上で $V_{\mathrm{cross}}$ と、既存の $\mathfrak S_{\HH}$ を
$V_\square$ へ引き戻した開部分は、全幾何点の集合が同じである。
二つの開部分スキームは幾何点集合が同じなら同じ開部分である。
従って両者は scheme として一致する。
既存の全相対開埋込 $\Xi$ を用いると、この等式は非被約なものを含む
既存の相対 site で許容する任意の基底にも適用され、
表示した相対群oidの同値を与える。
閉境界とその開補集合に関する幾何点の判定だけから、相対 heart を直接貼り合わせた
わけではない。相対存在と全 Isom の同定は、既に証明された $\Xi$ による。
\end{proof}

\subsection{二十個以内の多項式による実装}

\begin{lemma}[Segreを避けるkernelの五生成開]
\label{lem:qa-segre-kernel-five-polynomials}
$M$ を、始域を $k^2\otimes k^2$ と同定した四次正方行列とする。
$Q(x)=x_{00}x_{11}-x_{01}x_{10}$ は分解可能 tensor の Segre 方程式とする。
基底列 $e_1,\ldots,e_4$ に対し
\[
 I_{\rm ker}(M)=
 \left(\det M,
       Q(\operatorname{adj}(M)e_1),\ldots,
       Q(\operatorname{adj}(M)e_4)\right)
 \tag{IK}\label{eq:qa-five-polynomial-kernel-ideal}
\]
と置く。すると $\mathcal K(M)=D(I_{\rm ker}(M))$ である。
ここで $D(I)$ は生成元の少なくとも一つが非零となる開部分を意味する。
\end{lemma}
\begin{proof}
階数四なら行列式が非零である。
階数三なら随伴行列は階数一で、その非零列が kernel 線を張る。
従って kernel 線が Segre の外にあることは、表示した四つの二次式の
少なくとも一つが非零であることと同値である。
階数二以下なら行列式も随伴行列も零である。
この場合、kernel の射影化は直線を含み、代数閉体上でその直線は
Segre 二次曲面と交わる。
これで幾何点における同値を得る。
両辺とも開部分なので、普遍行列空間上で同じ開部分スキームとなる。
\end{proof}

\begin{corollary}[四組の明示有限式]
\label{cor:qa-square-h2-twenty-polynomials}
全空間 $\mathbb A^{32}$ の固定基底表示で
\[
 V_{\mathrm{cross}}=
 D(I_{\rm ker}(M_f))\cap D(I_{\rm ker}(M_f^{\mathrm t}))
 \cap D(I_{\rm ker}(M_g))\cap D(I_{\rm ker}(M_g^{\mathrm t})).
\]
従って四組各五個、合計二十個以内の多項式で全 $H_2$ 開を判定できる。
各組で少なくとも一つが非零であることを要求し、二十個全ての非消滅を
要求するのではない。
$M_f,M_g$ の成分は $f,g$ の係数に関して次数四、各組の行列式は次数十六、
随伴列の Segre 二次式は次数二十四である。
転置の行列式は元の行列式と同じなので、表示した二十式には重複もある。

任意基底上では、四つの引戻し ideal がそれぞれ単位 ideal となることが
対応する条件である。
\end{corollary}

\begin{remark}[idealのscheme構造と閉境界を区別する]
式~\eqref{eq:qa-five-polynomial-kernel-ideal} の生の五生成 ideal が
全ての座標基底または局所枠の変更に対して同じ ideal として貼り合うとは主張しない。
列の線形結合では二次式の polar cross term も生じるため、
その scheme-theoretic な零集合を canonical な incidence 像と同定するには
別の検証が必要である。
ここで基底に依存せず決まるのは $D(I_{\rm ker})$、またはその閉補集合の
radical ideal であり、全 $H_2$ 開の scheme としての同定にはこれで十分である。

一方 $Z_{33}=V(\det f,\det g)$ は $H_2$ 開の中の実際の閉境界である。
この上の相対中央二形式正規形には、二つの行列式が基底環で零となることを使う。
全幾何 fiber で零という条件で置き換えてはならない。
例えば行列式が非零の nilpotent となる一次横断族は、$H_2$ 開には属しても
この閉境界へ scheme として因数分解しない。
この違いは論文Bの
\cite[\texttt{prop:lagrangian-relative-double-rank-three},
\texttt{cor:lagrangian-open-closed-recognition}]{QuadricModuli2026}
の閉境界と横断 chart の区別と一致する。
今回の交差条件は、そこで使った Lagrangian chart に限らず、
二積の全 $\mathbb A^{32}$ に書かれている。
\end{remark}

\begin{example}[二積の可逆性だけでは足りない]
四つの二次元空間を同じ基底で同定し、$f=g=I_4$ とする。
両積は同型であり、kernel 条件も同時全射性も満たす。
しかし $T_f=I_4$ の交差縮約は
\[
 c^\vee\otimes d\longmapsto c^\vee(d)\,I_2
\]
であって階数一であり、分解可能 kernel を持つ。
従ってこの aligned なデータは $H_2$ 開に属さない。
古典的 $\mathbb P^1\times\mathbb P^1$ の二経路では、相対写像は
二つの tensor 因子の交換である。その交差縮約は可逆であり、古典点は正しく残る。
\end{example}


\begin{example}[次数二十四の条件が必要となる点]
$f=I_4$ とし、
\[
 g=\begin{pmatrix}
 1&0&0&2\\0&-1&0&0\\0&0&-1&0\\2&0&0&1
 \end{pmatrix}
\]
と置く。$\det g=-3$ であるが、交差縮約は
\[
 M_f=\begin{pmatrix}
 1&0&0&-1\\0&2&0&0\\0&0&2&0\\-1&0&0&1
 \end{pmatrix}
\]
となり、階数三でkernelは非分解テンソル$(1,0,0,1)$で張られる。
その随伴列の四つのSegre値は$(16,0,0,16)$である。
$M_f$とその転置のkernel条件が成立し、$f,g$が可逆なので
補題~\ref{lem:qa-cross-invertible-cubic}により残る二条件も成立する。
従ってこれは$H_2$点であり、交差行列式の非消滅だけではこの点を失う。
随伴列の次数二十四の式を残す必要がある。
\end{example}

\begin{corollary}[明示atlas上の普遍相対square AS代数]
\label{cor:qa-explicit-h2-atlas-square-as}
$V_{\mathrm{cross}}$ 上の普遍階数十六square代数$\Lambda$は、
canonical Serre完成と四頂点展開によって実際の相対square AS代数を与える。
全次数の局所自由性、係数束付き単純解消とその双対、および各許容基底上の
左右Noether性と指定対象を保つtails認識が成立する。
曲線avatarが基底上で平坦な曲線族をなすことを入力にせず、古典fiberも含む。
\end{corollary}
\begin{proof}
定理~\ref{thm:qa-square-explicit-h2-open}によって全幾何fiberは真正な
$H_2$ squareである。補題~\ref{lem:qa-classified-fiber-finite-square-recognition}、
系~\ref{cor:qa-relative-square-completion-two-sided-noetherian}、
命題~\ref{prop:qa-relative-square-unfolding}と
命題~\ref{prop:qa-relative-square-as-resolutions}を適用する。
全$\HHstack$上で増幅前の四頂点が大域降下するとは主張しない。
そこでは命題~\ref{prop:qa-universal-h2-square-morita81}の階数八十一代数を使う。
\end{proof}

\begin{remark}[入力と本項の位置付け]
geometric quintuple の定義と cubic 分類は
\cite[Definition 5.3.1, Lemma 5.3.3, Theorem 5.5.10]{VdB2011}
の既知の入力である。一次資料 arXiv:0807.3753v3 の pp.~13--14,20 で
これらの記述を確認した。
noncubic 分類と全相対開埋込は論文A・Bおよび本稿の既存結果に依存する。
本項で追加したのは、二積の随伴行列からの四交差条件、両階数三の直接有限認識、
および二十式以内での全開部分の記述である。
この組合せの文献上の新規性は未確認であり、証明稿としての主張と
形式証明器による検証とを同一視しない。
\end{remark}

% !TeX root = ../nc-quadric-surface-avatar.tex
% 2026-11-23 / v4.18。全H2開上の階数低下因子。

\subsection{二つの階数低下因子の横断性}
\label{subsec:qa-square-rank-boundary-geometry}

前項の明示開 $V_{\mathrm{cross}}\subset\mathbb A^{32}$ を用いる。
$G=G_\square/K_\square$ と略記し，
\[
 D_f=V(\det f)\cap V_{\mathrm{cross}},\qquad
 D_g=V(\det g)\cap V_{\mathrm{cross}},\qquad
 Z=D_f\cap D_g
\]
と置く。交点と和集合にはそれぞれ
$(\det f,\det g)$ と $(\det f\det g)$ による scheme 構造を入れる。

\begin{proposition}[全域の横断的階数低下]
\label{prop:qa-square-global-transverse-rank-divisors}
$D_f,D_g$ は $V_{\mathrm{cross}}$ 上の滑らかな有効 Cartier 因子で，
$Z$ は滑らかな余次元二の閉部分スキームである。
$D_f$ と $D_g$ は横断的に交わる。

従って ordered ruling stack
$\mathfrak H_2^{\mathrm{rul}}=[V_{\mathrm{cross}}/G]$ 上に
二つの滑らかな有効 Cartier 因子
$\mathfrak D_f,\mathfrak D_g$ が定まり，
その交点 $\mathfrak Z$ は滑らかで余次元二，
和集合は単純正規交差因子である。
特に $\dim\mathfrak Z=1$，$\dim\mathfrak D_f=\dim\mathfrak D_g=2$ である。
unordered stack $\mathfrak H_2$ 上には，二枝の交換を許す
\'etale 局所の正規交差因子 $\mathfrak D$ が降下する。
\end{proposition}

\begin{proof}
明示開上では $\operatorname{rank}f,\operatorname{rank}g\geq3$ である。
四次正方行列の determinant hypersurface は階数三の点で滑らかである。
実際，行列式の微分の係数は三次小行列式であり，
階数三の点では少なくとも一つが非零である。
従って $D_f,D_g$ はともに滑らかな Cartier 因子である。

二積の全空間は
\[
 \mathbb A^{32}
  =\operatorname{Hom}(C\otimes A,W)
     \times\operatorname{Hom}(D\otimes B,W)
\]
である。$Z$ の任意の幾何点において
$d(\det f)$ と $d(\det g)$ は，この二つの異なる因子にそれぞれ属する
非零 cotangent vector であるため一次独立である。
よって $Z$ は滑らかな余次元二となり，交わりは横断的である。
同じ計算により，交点の \'etale 局所座標では和集合の方程式を $xy=0$ とできる。

$\det f,\det g$ は前節の線束 $L_f,L_g$ の $G$-equivariant section である。
従って Cartier 因子と交点は商 stack へ降下する。
$V_{\mathrm{cross}}\to[V_{\mathrm{cross}}/G]$ は滑らかな atlas であるため，
滑らかさ，余次元および横断性の主張も降下する。
$\dim[V_{\mathrm{cross}}/G]=3$ より表示した次元を得る。
二つの因子は ruling の交換によって互いに移るので，
和集合は $\mathfrak H_2=[V_{\mathrm{cross}}/(G\rtimes S_2)]$ へ降下する。
順序付き二重被覆の上で既に単純正規交差であることから最後の主張が従う。
\end{proof}

\begin{corollary}[unordered階数因子の正規化]
\label{cor:qa-square-rank-divisor-normalization}
有限射
\[
 \nu:\mathfrak D_f=[D_f/G]\longrightarrow
 \mathfrak D=[(D_f\cup D_g)/(G\rtimes S_2)]
\]
は $\mathfrak D$ の正規化である。
この射は noncubic 閉部分
$\mathfrak N=[Z/(G\rtimes S_2)]$ の外では同型であり，
$\mathfrak N$ 上への引戻しは
\[
 [Z/G]\longrightarrow[Z/(G\rtimes S_2)]
\]
という内在的な ruling の二重被覆である。
また $\mathfrak D$ 内の conductor 閉部分は $\mathfrak N$ である。
\end{corollary}

\begin{proof}
二つの滑らかな因子が横断的に交わるので，
$D_f\amalg D_g\to D_f\cup D_g$ は有限な正規化射である。
これは交点近傍の式 $xy=0$ とその二つの smooth branch によって確認できる。
この射は $G\rtimes S_2$-equivariant であり，
\[
 [(D_f\amalg D_g)/(G\rtimes S_2)]\simeq[D_f/G]
\]
である。従って上の正規化射を商へ降下すると $\nu$ を得る。
有限性・代表可能性・正規化であることは上の atlas への引戻しで確認できる。
交点を除けば正規化は同型であり，
交点上の二枝はちょうど二つの ruling を表す。
最後に，regular な局所係数環 $R$ 上の $R[x,y]/(xy)$ において，
二枝への正規化の conductor は $(x,y)$ であるため，
conductor の閉部分は scheme としても $Z$ と一致する。
これも商へ降下する。
\end{proof}

\begin{remark}[非被約横断族と相対サイト]
以上の Cartier 因子，交点，および正規化射は普遍的な代数 stack 上の構成である。
その任意の scheme 基底変更を取れるが，
任意に引き戻した因子や交点が再び滑らか・被約であるとは主張しない。
引き戻した有限射が再び正規化となることや，conductorの形成が任意の非平坦基底変更と
可換となることも主張しない。
また actual heart の相対 groupoid の同定には，
既存の $\mathfrak H_2$ の相対サイトの範囲を引き続き用いる。
$\det f$ または $\det g$ が非零の nilpotent となる族は，
全幾何 fiber が両階数三でも $Z$ へ因数分解しない。
\end{remark}

% !TeX root = ../nc-quadric-surface-avatar.tex
% 2026-11-23 / v4.19。全H2の飽和GIT開と良モジュライ。

\subsection{共通線形化による全$H_2$のgood moduli space}
\label{subsec:qa-global-h2-good-moduli}

本項では前項の代数閉体 $k$ の標数を零とする。
四つの二次元空間には固定した基底を取り，
$k^2\otimes k^2$ の基底を $(00,01,10,11)$ の順に並べる。
Segre 二次形式の対称行列を
\[
 J=\begin{pmatrix}
 0&0&0&1\\
 0&0&-1&0\\
 0&-1&0&0\\
 1&0&0&0
 \end{pmatrix}
 \quad\left(J^2=I_4,\quad\det J=1,\quad
 x^{\rm t}Jx=2(x_{00}x_{11}-x_{01}x_{10})\right)
\]
とする。以下の写像の始域・終域の両方にこの規約を用いる。

\begin{lemma}[左右kernel条件をまとめる二つの式]
\label{lem:qa-segre-kernel-two-semi-invariants}
四次正方行列 $M$ に対し
\[
 P(M)=\det M,\qquad
 Q(M)=e_3(JM^{\rm t}JM)
      =\operatorname{tr}\!\left(
           \operatorname{adj}(M)J\operatorname{adj}(M)^{\rm t}J
        \right)
\]
と置く。ここで $e_3(N)$ は
$\det(tI_4+N)$ の $t$ の係数である。
$P,Q$ の次数はそれぞれ四，六であり，
\[
 \mathcal K(M)\cap\mathcal K(M^{\rm t})=D(P(M),Q(M)).
 \tag{PQ}\label{eq:qa-cross-pair-two-polynomials}
\]
この等式は普遍行列空間の開部分として成立し，任意基底変更に保たれる。
\end{lemma}
\begin{proof}
$e_3(N)=\operatorname{tr}(\operatorname{adj}(N))$ と
$\operatorname{adj}(AB)=\operatorname{adj}(B)\operatorname{adj}(A)$，
$\operatorname{adj}(J)=J$ から二つの $Q$ の表示が一致する。
階数四なら $P\ne0$ である。
階数三の場合
$\operatorname{adj}(M)=\kappa\lambda^{\rm t}$ と書けば，
$\kappa,\lambda$ はそれぞれ $M,M^{\rm t}$ の kernel を張り，
\[
 Q(M)=(\kappa^{\rm t}J\kappa)(\lambda^{\rm t}J\lambda).
\]
従って $Q\ne0$ は二つの kernel 線がともに Segre を避けることと同値である。
階数二以下なら随伴行列が零なので $P=Q=0$ であり，
kernel の射影化は Segre と交わる。
以上で幾何点における同値を得る。
両辺が開部分であることから普遍空間上の等式と任意基底変更への整合性が従う。
\end{proof}

$X=\mathbb A^{32}$ とし，
\[
 \widetilde G=
 \operatorname{GL}(A)\times\operatorname{GL}(B)
 \times\operatorname{GL}(C)\times\operatorname{GL}(D)
 \times\operatorname{GL}(W),\qquad
 G=\widetilde G/K_\square
\]
と書く。文字 $a,b,c,d,w$ は各群要素の determinant を表すものとし，
次の指標を定める：
\[
 \chi_f=\frac{w}{a^2c^2},\qquad
 \chi_g=\frac{w}{b^2d^2},\qquad
 \chi_0=\frac{w}{abcd},\qquad
 \zeta_f=\chi_f\chi_0,\quad \zeta_g=\chi_g\chi_0.
 \tag{CH}\label{eq:qa-h2-git-characters}
\]
いずれも $K_\square$ 上で自明であり，
$\chi_f\chi_g=\chi_0^2$，$\zeta_f\zeta_g=\chi_0^4$ である。
指標 $\chi_0$ は ruling 交換で不変なので
$H=G\rtimes S_2$ へ，$S_2$ 上自明な指標として延びる。

\begin{proposition}[共通指標の半不変式による飽和被覆]
\label{prop:qa-h2-common-character-saturated-cover}
前項の交差行列 $M_f,M_g$ に対し
\[
 P_f=P(M_f),\quad Q_f=Q(M_f),\qquad
 P_g=P(M_g),\quad Q_g=Q(M_g)
\]
とする。その $G$-半不変指標は順に
$\zeta_f^2,\zeta_f^3,\zeta_g^2,\zeta_g^3$ である。
次の四式は全て指標 $\chi_0^{24}$ の半不変式である：
\[
 \begin{split}
 s_0&=(P_fP_g)^3,\qquad s_1=P_f^3Q_g^2,\\
 s_2&=Q_f^2P_g^3,\qquad s_3=(Q_fQ_g)^2.
 \end{split}
 \tag{SS}\label{eq:qa-h2-four-common-semi-invariants}
\]
これらは $X$ の座標の通常の総次数では九十六であり，
\[
 s_1s_2=s_0s_3,\qquad
 V_{\rm cross}=\bigcup_{i=0}^3D(s_i).
 \tag{SAT}\label{eq:qa-h2-saturated-semi-invariant-cover}
\]
また ruling 交換は $s_1,s_2$ を交換し $s_0,s_3$ を固定する。
従って $H$ についても，同じ開部分を
\[
 V_{\rm cross}=D(h_0)\cup D(h_1)\cup D(h_2),\qquad
 (h_0,h_1,h_2)=(s_0,s_3,s_1+s_2)
 \tag{UN}\label{eq:qa-h2-three-unordered-semi-invariants}
\]
という同一指標 $\chi_0^{24}$ の三半不変式で覆える。
\end{proposition}
\begin{proof}
$M_f:C^\vee\otimes D\to A\otimes B^\vee\otimes L_f$ の変換則は
\[
 M_f\longmapsto
  \chi_f\,(g_A\otimes g_B^{-{\rm t}})\,
  M_f\,(g_C^{-{\rm t}}\otimes g_D)^{-1}
\]
である。tensor 積上の determinant 二次形式に対し，
左側の similitude factor は $\chi_f^2a/b$，
右側の factor は $d/c$ である。
従ってその比は
\[
 \frac{\chi_f^2ac}{bd}=\chi_f\chi_0=\zeta_f
\]
となる。
$JM_f^{\rm t}JM_f$ は相似変換と $\zeta_f$ 倍によって変わるため，
$Q_f=e_3(JM_f^{\rm t}JM_f)$ の指標は $\zeta_f^3$ である。
また二次元 tensor 積の determinant はその similitude factor の二乗なので，
$P_f$ の指標は $\zeta_f^2$ である。$g$ 側も同様である。

補題~\ref{lem:qa-segre-kernel-two-semi-invariants} により
\[
 V_{\rm cross}=D(P_f,Q_f)\cap D(P_g,Q_g).
\]
この積の四通りの非消滅を正の冪で揃えたものが
式~\eqref{eq:qa-h2-saturated-semi-invariant-cover} である。
指標と恒等式は直接計算から従う。
ruling 交換は固定基底の係数式において $P_f,P_g$ および $Q_f,Q_g$ を交換する。
もし $s_1+s_2=0$ かつ $s_1$ または $s_2$ が非零なら両方とも非零であり，
$s_1s_2=s_0s_3$ から $s_0,s_3$ も非零となる。
これにより三式でも同じ開を覆うことが分かる。
\end{proof}

\begin{theorem}[全$H_2$の準射影的good moduli space]
\label{thm:qa-global-h2-quasiprojective-good-moduli}
$\Gamma=G$ または $H$ とし，半不変式の次数付き環を
\[
 R_\Gamma(\chi_0)=
   \bigoplus_{n\geq0}
   \{a\in k[X]\mid a(\gamma x)=\chi_0(\gamma)^n a(x)\}
\]
と定める。この環は有限生成で，次数零部分は $k$ である。
$\Gamma=G$ なら
\[
 M^{\rm rul}=\bigcup_{i=0}^3D_+(s_i)
             \subset\operatorname{Proj}R_G(\chi_0),
\]
$\Gamma=H$ なら
\[
 M=\bigcup_{i=0}^2D_+(h_i)
             \subset\operatorname{Proj}R_H(\chi_0)
\]
と置く。このとき
\[
 \pi^{\rm rul}:\mathfrak H_2^{\rm rul}\longrightarrow M^{\rm rul},
 \qquad
 \pi:\mathfrak H_2\longrightarrow M
\]
は good moduli space morphism である。
特に $M^{\rm rul},M$ は有限型で separated な準射影 $k$-scheme である。

ここで good moduli space morphism とは，
準コンパクトで，準連接層の直像が完全で，構造層の自然射が同型となる射をいう。
既存の定理~\ref{thm:qa-square-explicit-h2-open} による
全相対 groupoid の同定を用いており，
幾何点のみの粗い対応を主張しているのではない。
\end{theorem}
\begin{proof}
標数零で $G$ と $H$ は線型簡約である。
$X\times\mathbb A^1$ 上の点への作用を
$\gamma\cdot(x,z)=(\gamma x,\chi_0(\gamma)^{-1}z)$ とすると
\[
 R_\Gamma(\chi_0)\simeq(k[X][z])^\Gamma,\qquad
 a\text{ in degree }n\longleftrightarrow az^n.
\]
従って簡約群の不変式の有限生成性から $R_\Gamma(\chi_0)$ は有限生成である。
さらに $\operatorname{GL}(W)$ の scalar は $f,g$ の全座標を同じ正の重みで
変換するため，$k[X]^\Gamma=k$ である。
よって $\operatorname{Proj}R_\Gamma(\chi_0)$ は projective $k$-scheme であり，
表示した開部分は準射影である。

被覆に使った任意の半不変式 $h$（ordered の場合は $s_i$）は
$\chi_0^{24}$ の指標を持つ。
その localization について
\[
 \bigl(R_\Gamma(\chi_0)[h^{-1}]\bigr)_0
       =(k[X][h^{-1}])^\Gamma
 \tag{LC}\label{eq:qa-h2-localized-invariants}
\]
が成り立つ。実際，右辺の元を $a/h^m$ と書けば，
その不変性は $a$ が指標 $\chi_0^{24m}$ の半不変式であることと同値である。
従って各標準開の上では求める射は
\[
 [D(h)/\Gamma]\longrightarrow
 \operatorname{Spec}(k[X][h^{-1}])^\Gamma
       =D_+(h)
\]
という affine quotient である。
これは線型簡約群の affine quotient の good moduli space morphism である
\cite[Theorem 13.2]{AlperGoodModuli2013}。

二つの被覆式 $h,h'$ は同じ指標を持つので，
$D(h)\cap D(h')$ は $D(h)$ 上の不変正則関数 $h'/h$ の非消滅開である。
したがって重なりは $D_+(h)$ の対応する principal open の逆像であり，
飽和している。局所 quotient は同じ invariant localization によって貼り合う。
この Zariski 局所構成から，全体で準連接直像の完全性と
構造層の直像同型を得る。
飽和開の貼合せについては\cite[Proposition 7.9]{AlperGoodModuli2013}にも一致する。
これは GIT の semistable 開からの good moduli space の構成
\cite[Theorem 13.6]{AlperGoodModuli2013} の明示的な飽和開への制限である。

最後に
$\mathfrak H_2^{\rm rul}=[V_{\rm cross}/G]$ と
$\mathfrak H_2=[V_{\rm cross}/H]$
という定理~\ref{thm:qa-square-explicit-h2-open}の全相対同値を用いれば結論が従う。
\end{proof}

\begin{remark}[線形化の役割と三式の限界]
本定理の根拠は，単に $V_{\rm cross}$ が簡約群の不変開であることではない。
同一指標にそろえた半不変式によって affine 開を覆い，
重なりの飽和性を式~\eqref{eq:qa-h2-localized-invariants} で確認したことを使う。
三つの $h_i$ は被覆を与えるが，不変式環を生成するとは主張しない。
従ってこれらの三座標だけで同型類を分類することや，
$M$ の射影平面への閉埋込みを得たという主張はしない。

両階数三の閉部分では $M_f,M_g$ はともに可逆なので，
全 noncubic 閉部分は $D(s_0)=D(P_fP_g)$ の一つの affine quotient chart に入る。
構成したgood moduli射は標的$M$または$M^{\rm rul}$からの任意基底変更で
good moduli射のままである\cite[Proposition 4.7]{AlperGoodModuli2013}。
actual heartの群oidとの比較の範囲には、引き続き本稿の相対サイトを用いる。
この良モジュライは非同型対象の軌道閉包が交わる場合にそれらを同一視する。
従って論文B\cite[\texttt{cor:global-h2-no-coarse-space}]{QuadricModuli2026}
の粗モジュライ不存在と矛盾しない。
\end{remark}

% !TeX root = ../nc-quadric-avatar.tex
% 2026-11-23 / v4.19。Cubic開と既存四量子ビットGIT。

\subsection{Cubic開の飽和性と四量子ビットGIT}
\label{subsec:qa-global-cubic-git-comparison}

前項と同じく標数零の代数閉体 $k$ 上で考える。
$F=\det f$，$G_0=\det g$ と置き，
$\mathfrak U_f=[V_{\rm cross}\cap D(F)/G]$，
$\mathfrak U_g=[V_{\rm cross}\cap D(G_0)/G]$
を ordered $H_2$ の二つの cubic 開とする。
ここでいう cubic 開は linear quintuple に対応する古典型も含む。

\begin{proposition}[二つのcubic開は飽和している]
\label{prop:qa-global-cubic-opens-saturated}
$\mathfrak U_f,\mathfrak U_g$ は
$\pi^{\rm rul}:\mathfrak H_2^{\rm rul}\to M^{\rm rul}$
について飽和開である。具体的に
\[
 M_f^{\rm rul}
  =M^{\rm rul}\cap
       \bigl(D_+(F^2P_g)\cup D_+(F^3Q_g)\bigr)
\]
と置けば，
$\mathfrak U_f=(\pi^{\rm rul})^{-1}(M_f^{\rm rul})$ である。
$g$ 側は $f,g$ を交換した式で与えられる。
\end{proposition}
\begin{proof}
$F^2P_g$ と $F^3Q_g$ の指標はそれぞれ
\[
 \chi_f^2\zeta_g^2=\chi_0^6,\qquad
 \chi_f^3\zeta_g^3=\chi_0^9
\]
である。一方，$V_{\rm cross}$ 上では $P_g,Q_g$ の少なくとも一方が
非零なので，この開部分の内部で
\[
 V_{\rm cross}\cap D(F)
 =
 (V_{\rm cross}\cap D(F^2P_g))
 \cup
 (V_{\rm cross}\cap D(F^3Q_g))
\]
が成立する。右辺の各項は正の $\chi_0$ 冪の半不変式に対応する
$\operatorname{Proj}R_G(\chi_0)$ の標準開の逆像である。
これは $X=\mathbb A^{32}$ 全体で右辺が $V_{\rm cross}$ に入るという
主張を用いていない。
\end{proof}

\paragraph{可逆な一積による四因子表示.}
$f$ が可逆なら $W=C\otimes A$ と同一視して $f=1$ とし，
\[
 T=f^{-1}g:D\otimes B\longrightarrow C\otimes A,\qquad
 w_T\in A\otimes C\otimes D^\vee\otimes B^\vee
\]
と書ける。この gauge の残余群は
\[
 G_f=
 \bigl(\operatorname{GL}(A)\times\operatorname{GL}(B)
       \times\operatorname{GL}(C)\times\operatorname{GL}(D)\bigr)/K_f,
\]
\[
 K_f=
 \{(\lambda_A I,\lambda_B I,\lambda_C I,\lambda_D I)
          :\lambda_A\lambda_C=\lambda_B\lambda_D\}
       \simeq\mathbb G_m^3 .
\]
これは四因子 $A,C,D^\vee,B^\vee$ の tensor 表現の有効群である。
一次元の関係線を別に入れた raw quintuple 表示から
四次元の自明 scalar 核を消した群とも同型である
（命題~\ref{prop:qa-ou-scalar-gerbe}）。
群の点への作用は
\[
 T\longmapsto
 (g_C\otimes g_A)\,T\,(g_D\otimes g_B)^{-1}
\]
であり，この残余群上で
\[
 \chi_f=1,\qquad
 \chi_0=\chi=\frac{\det A\,\det C}{\det B\,\det D}.
 \tag{CQ}\label{eq:qa-cubic-residual-character}
\]
以下では $\operatorname{Cr}(T)$ を前項と同じ交差縮約とし，
$N(T)=\operatorname{Cr}(\operatorname{adj}T)$ と置く。
これは $f=1$ での $M_g$ である。

\begin{lemma}[Geometric条件の二半不変式による判定]
\label{lem:qa-cubic-adjugate-two-invariants}
$w_T$ が geometric quintuple を与える必要十分条件は
\[
 \mathcal K(N(T))\cap\mathcal K(N(T)^{\mathrm t})
\]
である。従って geometric 開 $\mathscr G$ は
\[
 \mathscr G
   =D\bigl(p_{\rm adj},q_{\rm adj}\bigr)
       \subset\operatorname{Hom}(D\otimes B,C\otimes A),
 \quad
 p_{\rm adj}=P(N(T)),\quad q_{\rm adj}=Q(N(T)).
 \tag{CA}\label{eq:qa-cubic-adjugate-open}
\]
$p_{\rm adj},q_{\rm adj}$ は通常の tensor 座標次数で $12,18$，
$G_f$ の指標で $\chi^6,\chi^9$ の半不変式である。
この開の等式は任意の係数スキームへの引戻しと整合する。
\end{lemma}
\begin{proof}
Geometric 条件は $T,T^{\mathrm t},
\operatorname{Cr}(T),\operatorname{Cr}(T)^{\mathrm t}$ の四つの kernel が
非零の分解可能 tensor を含まない条件である
\cite[Definition 5.3.1]{VdB2011}。
$\operatorname{rk}T\le2$ なら $\operatorname{adj}T=0$ なので本補題の条件は
成立しない。この場合，$\mathbb P(\ker T)$ は少なくとも直線を含み，
その直線が Segre 二次曲面と交わるため geometric 条件も成立しない。

$\operatorname{rk}T=4$ なら
$\operatorname{adj}T=(\det T)T^{-1}$ である。
可逆な $T$ について交差条件が $T^{-1}$ と同値になることは
補題~\ref{lem:qa-cross-invertible-cubic} の証明の ruling 二平面の議論による。
例えば $\operatorname{Cr}(T)(c^\vee\otimes d)=0$ は
\[
 T(kd\otimes B)=(\ker c^\vee)\otimes A
\]
と同値であり，逆写像を取れば $T^{-1}$ の対応する分解可能 kernel
条件を得る。転置も同様である。

残る $\operatorname{rk}T=3$ では
\[
 \operatorname{adj}T=\kappa\otimes\lambda,\qquad
 k\kappa=\ker T\subset D\otimes B,\qquad
 k\lambda=\operatorname{Ann}(\operatorname{im}T)
                    \subset C^\vee\otimes A^\vee
\]
と書く。補題~\ref{lem:qa-rank-one-cross-kronecker} により，
$N(T)$ の条件は $\kappa,\lambda$ がともに非分解可能であることと同値で，
そのとき $N(T),N(T)^{\mathrm t}$ は実際に可逆である。
Geometric 条件の最初の二条件は正にこの非分解可能性を要求する。
これから残る二条件も従うことを確認する。
仮に $\operatorname{Cr}(T)(c^\vee\otimes d)=0$ なら
\[
 T(kd\otimes B)\subset(\ker c^\vee)\otimes A .
\]
$\kappa$ が非分解可能なので，$T$ は左辺の入力二平面上で単射である。
従ってこの包含は等号となる。
しかし $\lambda$ が非分解可能なら，
$\ker\lambda=\operatorname{im}T$ は
$(\ker c^\vee)\otimes A$ のような ruling 二平面を含まない。
これは矛盾である。
転置の分解可能 kernel についても
$T(D\otimes kb)\subset C\otimes\ker a^\vee$ を使えば同じ証明となる。

これで全階数の場合を尽くした。
二式による表示は
補題~\ref{lem:qa-segre-kernel-two-semi-invariants} から従う。
$N(T)$ の成分は $T$ の三次式なので通常次数は $12,18$ となる。
また式~\eqref{eq:qa-cubic-residual-character} と
$\chi_g=\chi^2$ から $\zeta_g=\chi^3$ となり，
指標は $\zeta_g^2=\chi^6,\zeta_g^3=\chi^9$ である。
最後に，以上の幾何点による同値は普遍係数空間の二つの開部分の一致を示す。
従って非被約係数スキームを含めて任意の引戻し後にも同じ開を定める。
\end{proof}

\begin{proposition}[既存四量子ビット商との開部分比較]
\label{prop:qa-global-cubic-four-qubit-good-moduli}
四因子表現の $SL_2^4$ 不変式環に通常の多項式次数を入れ，
その $\operatorname{Proj}$ を $\overline M_{\rm qui}$ とする。
このとき
\[
 \mathfrak U_f\simeq[\mathscr G/G_f],
 \qquad
 M_f^{\rm rul}\simeq
 \Omega_{\rm geo}:=
       D_+(p_{\rm adj})\cup D_+(q_{\rm adj})
       \subset\overline M_{\rm qui}
\]
という同型があり，二つの good moduli 射と整合する。
$g$ 側も同様である。
既知の不変式環の計算により
\[
 \overline M_{\rm qui}\simeq\mathbb P(2,4,4,6)
                  \simeq\mathbb P(1,2,2,3)
\]
であり，右側の重みは $\chi$ による次数に対応する。
\end{proposition}
\begin{proof}
$f$ の可逆性によって $W=C\otimes A$ と同一視する操作は，
ベクトル束と全同型を保つ相対操作である。
従って，これは単なる体上の標準形ではなく，商スタックの同型
\[
 [\{(f,g)\in X:F\ne0\}/G]\simeq
 [\operatorname{Hom}(D\otimes B,C\otimes A)/G_f]
\]
を与える。明示的 $H_2$ 開の認識と前補題によって，その $H_2$ 開は
$[\mathscr G/G_f]$ に一致する。

四因子それぞれの $SL_2$ による不変式は偶数次数にのみ存在する。
通常次数 $2n$ の不変式は四つの scalar に対し
$(\lambda_A\lambda_C/\lambda_B\lambda_D)^{2n}$
で変換するので，$GL_2^4$ に対する指標は $\chi^n$ である。
逆も明らかで，従って
\[
 R_{G_f}(\chi)_n
   =k[A\otimes C\otimes D^\vee\otimes B^\vee]^{SL_2^4}_{2n}.
 \tag{CGIT}\label{eq:qa-cubic-half-degree-regrading}
\]
これは偶数次数を半分にした同じ環であり，$\operatorname{Proj}$ は変わらない。
前補題により $\mathscr G$ はこの GIT 商の飽和開で，
その good moduli は表示した $\Omega_{\rm geo}$ である。
一方，命題~\ref{prop:qa-global-cubic-opens-saturated} は
$\mathfrak U_f\to M_f^{\rm rul}$ も good moduli とする。
Good moduli space の一意性からこの二つの scheme の同型を得る。

最後の不変式環は既知の四量子ビット計算
\cite[\S3, (3.9)--(3.16), (3.19)]{OU2014} であり，
通常次数 $2,4,4,6$ の四つの代数的独立な生成元を持つ。
同計算の直交群と $W(D_4)$ の不変式による記述は，
標数零の代数閉体上で同様に成立する。
式~\eqref{eq:qa-cubic-half-degree-regrading} により
$\chi$ 次数の重みは $1,2,2,3$ となる。
\end{proof}

\begin{remark}[比較の範囲]
この命題は既存の重み付き射影商の計算を新たに主張するものではない。
本稿で追加した点は，全 $H_2$ の大域的線形化を cubic 開に制限したときの
指標の一致と，geometric 開の明示的な飽和性である。
全 $M^{\rm rul}$ または非順序 $M$ が
$\overline M_{\rm qui}$ と同型であるとは主張していない。
両階数三の noncubic 閉部分は両 cubic 開の外側にある。
また geometric 開には非閉軌道を含み得るため，
ここで比較するのは good moduli scheme であり，
pointed exact 同型類を一対一に分類する coarse moduli ではない。
\end{remark}

% !TeX root = ../nc-quadric-surface-avatar.tex
% 2026-11-23 / v4.20 / cubic geometric 開の不変式表示。

\subsection{四量子ビット不変式によるgeometric開の二本の補曲線}
\label{subsec:qa-cubic-geometric-open-invariants}

前項の $f=1$ chart と交差縮約の規約を保つ。
固定した Segre 双線型形式を $J$ とし，
\[
 K_T=JT^{\mathrm t}JT,\qquad
 e_i=e_i(K_T)\ (i=1,2,3),\qquad d=\det T
\]
と置く。右辺の $e_i$ は行列の固有値の第 $i$ 基本対称式である。
$e_1,e_2,d,e_3$ の通常の多項式次数は $2,4,4,6$，
前項の指標 $\chi$ による次数は $1,2,2,3$ である。
Okawa--Ueda の不変式環の計算
\cite[\S3, (3.9)--(3.16)]{OU2014} と Newton の恒等式により
\[
 k[T]^{SL_2^4}=k[e_1,e_2,d,e_3].
 \tag{CI}\label{eq:qa-cubic-elementary-invariant-ring}
\]
従って既存 GIT 商には同次座標
$[e_1:e_2:d:e_3]$ を持つ
$\mathbb P(1,2,2,3)$ を用いることができる。

\begin{proposition}[二半不変式の普遍恒等式]
\label{prop:qa-cubic-adjugate-explicit-invariants}
$N(T)=\operatorname{Cr}(\operatorname{adj}T)$，
$p_{\rm adj}=\det N(T)$，
$q_{\rm adj}=e_3(JN(T)^{\mathrm t}JN(T))$ とする。
任意の $4\times4$ 行列 $T$ に対し，次の多項式恒等式が成立する。
\[
 \begin{split}
 16p_{\rm adj}
   &=e_3^2-4d^2e_2+8d^3,\\
 16q_{\rm adj}
   &=e_3^3-4d^2e_2e_3+16d^4e_1-8d^3e_3.
 \end{split}
 \tag{CPQ}\label{eq:qa-cubic-adjugate-pq-formulas}
\]
特に
\[
 q_{\rm adj}-e_3p_{\rm adj}
     =d^3(de_1-e_3).
 \tag{CPQ-diff}\label{eq:qa-cubic-adjugate-pq-difference}
\]
\end{proposition}
\begin{proof}
基底を $e_{00},e_{01},e_{10},e_{11}$ の順とし，
$H$ の列を
\[
 e_{00}+e_{11},\quad e_{00}-e_{11},\quad
 e_{01}+e_{10},\quad e_{01}-e_{10}
\]
とする。
この各列を適当な定数で割れば $J$ に関する正規直交基底となる。
その定数対角行列は以下の対角行列と可換なので，
Cartan 部分空間を
\[
 T=H\operatorname{diag}(x_1,x_2,x_3,x_4)H^{-1}
\]
と書いてよい。この部分空間では
\[
 e_i=e_i(x_1^2,x_2^2,x_3^2,x_4^2),\qquad
 d=x_1x_2x_3x_4 .
\]
$\lambda_i=\prod_{j\ne i}x_j$ と置けば，
$\operatorname{adj}T$ の同じ基底での固有値は $\lambda_i$ である。
交差縮約の係数を直接並べ替えると，$N(T)$ も同じ $H$ で対角化され，
その固有値 $\mu_i$ は
\[
 \begin{pmatrix}\mu_1\\\mu_2\\\mu_3\\\mu_4\end{pmatrix}
 =
 \frac12
 \begin{pmatrix}
 1&1&1&1\\
 1&1&-1&-1\\
 1&-1&1&-1\\
 1&-1&-1&1
 \end{pmatrix}
 \begin{pmatrix}\lambda_1\\\lambda_2\\\lambda_3\\\lambda_4\end{pmatrix}.
 \tag{CH}\label{eq:qa-cross-adjugate-hadamard}
\]
この $N(T)$ は対称で $J$ と可換なので，
$JN(T)^{\mathrm t}JN(T)=N(T)^2$ である。

$E_i=e_i(\lambda_1^2,\ldots,\lambda_4^2)$，
$D=\lambda_1\lambda_2\lambda_3\lambda_4$ と置くと，
式~\eqref{eq:qa-cross-adjugate-hadamard} の四つの線型式の展開により
\[
 \begin{split}
 16\prod_i\mu_i&=E_1^2-4E_2+8D,\\
 16e_3(\mu_1^2,\ldots,\mu_4^2)
     &=E_1^3-4E_1E_2+16E_3-8DE_1.
 \end{split}
\]
一方，補積 $\lambda_i$ の定義から
\[
 E_1=e_3,\qquad E_2=d^2e_2,\qquad
 E_3=d^4e_1,\qquad D=d^3 .
\]
代入すると求める二式を得る。

以上は Cartan 部分空間上の計算だが，両辺は全て $SL_2^4$ 不変式である。
式~\eqref{eq:qa-cubic-elementary-invariant-ring} の直交群による記述では，
Cartan への制限は $W(D_4)$ 不変式環との同型，特に単射である。
従って二式は全行列空間上の普遍恒等式となる。
最後の式は二式を引き算すれば従う。
\end{proof}

\begin{corollary}[Geometric開の被約補集合]
\label{cor:qa-cubic-geometric-two-complement-curves}
$\Omega_{\rm geo}\subset\mathbb P(1,2,2,3)$ を前項の geometric 開の
良モジュライとする。その補集合を被約構造で考えると
\[
 \bigl(\mathbb P(1,2,2,3)\setminus\Omega_{\rm geo}\bigr)_{\rm red}
    =C_0\cup C_1,
\]
\[
 C_0=V(d,e_3),\qquad
 C_1=V(e_3-de_1,\ 4e_2-e_1^2-8d).
 \tag{CC}\label{eq:qa-cubic-two-bad-curves}
\]
両成分は $\mathbb P(1,2)$ と同型であり，
唯一の交点は $[2:1:0:0]$ である。
被約和の同次イデアルは
\[
 \sqrt{(p_{\rm adj},q_{\rm adj})}
 =
 (e_3-de_1,\ d(4e_2-e_1^2-8d)).
 \tag{CR}\label{eq:qa-cubic-geometric-radical}
\]
交点の近傍でこの和は通常の節点である。
\end{corollary}
\begin{proof}
$h=e_3-de_1$，$\ell=4e_2-e_1^2-8d$ と置く。
前命題から，二式のイデアルは単元の係数を除いて
\[
 I=(h^2+2de_1h-d^2\ell,\ d^3h)
\]
と書ける。
これを含む素イデアルは，$d$ を含む場合には $h$ も含み，
$d$ を含まない場合には $h,\ell$ を含む。
逆に $(d,h)$ と $(h,\ell)$ はともに $I$ を含む素イデアルである。
従って
\[
 \sqrt I=(d,h)\cap(h,\ell)=(h,d\ell).
\]
最後の等号は $h$ で割った多項式環において $d$ と $\ell$ が
互いに素であることによる。

$C_0$ の座標環は $k[e_1,e_2]$，$C_1$ の座標環は $k[e_1,d]$ なので，
両者は重み $(1,2)$ の射影直線となる。
例えば後者の閉埋込みは
\[
 [s:t]\longmapsto
 [s:(s^2+8t)/4:t:st]
\]
で与えられる。
共通点では $d=e_3=0$ と $4e_2=e_1^2$ が成り立つため，
射影点は $[2:1:0:0]$ に限られる。

この点は $e_1\ne0$ の滑らかな affine chart に入る。
その座標を
\[
 u=\frac{4e_2}{e_1^2}-1,\qquad
 v=\frac d{e_1^2},\qquad
 w=\frac{e_3}{e_1^3}
\]
とすると，被約和のイデアルは $(w-v,\ v(u-8v))$ である。
滑らかな面 $w=v$ の正則座標 $v,u-8v$ に関する二軸の和なので，
通常の節点である。
\end{proof}

\begin{remark}[被約補集合と検算の範囲]
式~\eqref{eq:qa-cubic-geometric-radical} は radical の等式であり，
二つの高次式が定める閉部分そのものを被約と主張してはいない。
ここでの二曲線は既存 cubic GIT の非 geometric 領域であり，
全 $H_2$ の大域的 good moduli に含まれる noncubic 閉曲線とは区別する。

再現コード \texttt{verify\_cubic\_adjugate\_invariants.py} は，
上の基底と交差縮約を係数から実装し，Cartan 恒等式，
階数四・三・二の非対角行列での直接式，
および radical 証明に使った代数恒等式を exact arithmetic で検算する。
全行列上での恒等式の証明は，これらの有限例ではなく，
上記の不変式制限の単射性によって与えている。
\end{remark}

% !TeX root = ../nc-quadric-surface-avatar.tex
% 2026-11-23 / v4.21 / cubic chart の貼合せと全域のGIT比較。

\subsection{二つのcubic chartの貼合せと重み付き逆行列写像}
\label{subsec:qa-cubic-inversion-gluing}

$\mathbb W=\mathbb P(1,2,2,3)$ の座標を
$[e_1:e_2:d:e_3]$ とし，前項の geometric 開を
$\Omega=\Omega_{\mathrm{geo}}=\mathbb W\setminus(C_0\cup C_1)$ と書く。
この式で用いる補集合は underlying closed set である。
二つの高次式が定める閉部分をその radical と取り違えない。
$\Omega^\times=\Omega\cap D_+(d)$ と置く。

\begin{proposition}[重み付きCremona対合とcubic chartの貼合せ]
\label{prop:qa-cubic-weighted-inversion-gluing}
逆行列操作 $T\mapsto T^{-1}$ が $\Omega^\times$ に誘導する対合は
\[
 \iota[e_1:e_2:d:e_3]
       =[e_3:d^2e_2:d^3:d^4e_1]
 \tag{INV}\label{eq:qa-cubic-weighted-cremona}
\]
である。右辺の四式はそれぞれ重み $3,6,6,9$ を持つ。
同じ式は $\mathbb W\setminus C_0$ から $\mathbb W$ への正則射を定め，
特に $\Omega$ 全体で定義される。しかし
\[
 \iota(\Omega\cap V(d))=\{q_0\},
 \qquad q_0=[1:0:0:0]\in C_0\setminus C_1
\]
であり，$\Omega^\times$ 上の対合は $\Omega\to\Omega$ の射へ延長できない。

全 ordered good moduli の cubic 開は，二つの $\Omega$ を
$\Omega^\times$ 上のこの対合で貼り合わせたスキームである：
\[
 M^{\mathrm{rul}}\setminus\overline N^{\mathrm{rul}}
   \simeq\Omega_f\mathop{\cup}_{\iota}\Omega_g.
 \tag{GL}\label{eq:qa-cubic-two-git-chart-gluing}
\]
ここで $\overline N^{\mathrm{rul}}\simeq\mathbb A^1_\chi$ は noncubic 閉像である。
ruling 交換は二つの chart を同じ同次座標のまま交換する。
全 $M^{\mathrm{rul}}$ 上のこの正則対合は
$\Omega^\times$ 上で式~\eqref{eq:qa-cubic-weighted-cremona} に一致し，
$\overline N^{\mathrm{rul}}$ 上では恒等である。
\end{proposition}
\begin{proof}
$K_T=JT^{\mathrm{t}}JT$ とすると，$K_{T^{-1}}$ は $K_T^{-1}$ と相似である。
$\det K_T=d^2$ だから，逆行列の基本対称式は
\[
 e_1(T^{-1})=\frac{e_3}{d^2},\qquad
 e_2(T^{-1})=\frac{e_2}{d^2},\qquad
 d(T^{-1})=\frac1d,\qquad
 e_3(T^{-1})=\frac{e_1}{d^2}.
\]
重み $(1,2,2,3)$ のスカラーとして $d^2$ を掛ければ
式~\eqref{eq:qa-cubic-weighted-cremona} を得る。
二回適用した四座標は
\[
 [d^4e_1:d^8e_2:d^9:d^{12}e_3]
\]
であり，これは重み付きスカラー $d^4$ によって元の点と同じである。
逆行列による geometric 条件の保存は
補題~\ref{lem:qa-cubic-adjugate-two-invariants} の ruling 二平面の議論による。

式の四座標が同時に零となる台は $d=e_3=0$，すなわち $C_0$ である。
従ってその補で正則射を与える。
$d=0$ では geometric 条件は $e_3\ne0$ なので，像は $q_0$ となる。
$\Omega$ は整であり，$\mathbb W$ は separated である。
仮に $\Omega\to\Omega$ への延長があれば，その $\mathbb W$ への合成は
稠密開 $\Omega^\times$ で上の正則射と一致するので，全体で一致する。
しかし $q_0\notin\Omega$ であり矛盾する。
同じ理由で $\Omega\cap V(d)$ のどの点の近傍にも $\Omega$ 値の延長はない。

一方，$f,g$ がともに可逆な square では
$f=1$ の chart の行列は $T=f^{-1}g$，
$g=1$ の chart の行列は $T'=g^{-1}f=T^{-1}$ である。
命題~\ref{prop:qa-global-cubic-four-qubit-good-moduli} の相対 chart 同型と，
命題~\ref{prop:qa-global-closed-boundary-good-images} の飽和性を用いると，
二つの good moduli chart の共通部分は正に $\Omega^\times$ となり，
その貼合せ写像が $\iota$ である。
二つの cubic 開の補は noncubic 閉像だから，表示した貼合せを得る。

ruling 交換は有限表示 square 上で $f,g$ と対応する四つの係数束を交換する。
従って二つの chart の間では同じ座標値を持つ点へ移り，
共通部分の一方の chart で見れば $\iota$ となる。
その作用は半不変式環上の作用，または good moduli の普遍性により
全 $M^{\mathrm{rul}}$ へ降下する。
補題~\ref{lem:qa-noncubic-ruling-sign-pushforward-zero} は
ordered および unordered の noncubic good moduli をともに
$\mathbb A^1_\chi$ とし，その間の射を恒等とする。
よって ruling 交換はこの閉部分の上でも scheme として恒等である。
\end{proof}

\subsection{重み二の慣性を含めた固定部分と粗二重射}

\begin{proposition}[可逆cubic領域の粗環と固定部分]
\label{prop:qa-cubic-inversion-weight-two-fixed-locus}
$D_+(d)\subset\mathbb W$ は
\[
 \operatorname{Spec}
 k[b,u,v,w]/(uw-v^2),
 \qquad
 b=\frac{e_2}{d},\quad
 u=\frac{e_1^2}{d},\quad
 v=\frac{e_1e_3}{d^2},\quad
 w=\frac{e_3^2}{d^3}
 \tag{AC}\label{eq:qa-cubic-invertible-affine-ring}
\]
であり，$\iota$ は $u,w$ を交換し $b,v$ を固定する。
その scheme-theoretic fixed locus は
\[
 u=w,\qquad (u-v)(u+v)=0
\]
で，被約な二成分である。
同次座標では $e_3=\pm de_1$ という二つの被約部分の和となる。
両成分の交差は $u=v=w=0$ であり，
$\Omega^\times$ 内では $b\ne2$ の直線である。
特に $e_1=e_3=0$ の重み二の点を除外する必要はない。

商の affine 環は
\[
 \bigl(k[b,u,v,w]/(uw-v^2)\bigr)^{\langle\iota\rangle}
       =k[b,s,v],\qquad s=u+w.
\]
$z=u-w$ と置けば粗二重射の方程式は
\[
 z^2=s^2-4v^2=(s-2v)(s+2v).
 \tag{DB}\label{eq:qa-cubic-inversion-double-cover}
\]
従って標的は smooth であり，二つの branch 因子は横断する Cartier 因子である。
源の二つの固定成分は，交差する重み二の locus ではそれぞれ Cartier ではない。
\end{proposition}
\begin{proof}
$d$ の平方根を用いて重み付き座標を $d=1$ に規格化すると，
残る $\mu_2$ は三座標 $(a,b,c)$ に
\[
 (a,b,c)\longmapsto(-a,b,-c)
\]
と作用する。ここで形式的には
$a=e_1/\sqrt d$，$b=e_2/d$，$c=e_3/d^{3/2}$ である。
この $\mu_2$ の不変式は $u=a^2,v=ac,w=c^2$ で，唯一の関係は $uw=v^2$。
これは degree-zero localization の計算でもあり，
式~\eqref{eq:qa-cubic-invertible-affine-ring} は粗スキーム上での同型である。
式~\eqref{eq:qa-cubic-weighted-cremona} はこの被覆上で $a,c$ の交換となる。
従って粗環では $u,w$ を交換し，固定イデアルは $(u-w)$ である。
商を取ると $(u-v)(u+v)=0$ となり，被約二成分の主張を得る。

$e_1=e_3=0$ では $u=v=w=0$ である。
この直線上で前項の二式は
$p_{\mathrm{adj}}=(2-b)/4$，$q_{\mathrm{adj}}=0$ となるから，
geometric 開に残る条件は $b\ne2$ である。

環を $s=u+w,z=u-w$ で書き直すと
\[
 k[b,u,v,w]/(uw-v^2)
   =k[b,s,v,z]/(z^2-s^2+4v^2).
\]
対合は $z\mapsto-z$ なので不変式環と二重射の式が従う。
$s-2v,s+2v$ の微分は交差上で一次独立だから，二つの branch 因子は横断する。
源の個々の固定成分のイデアルは
$(u-w,u-v)$ と $(u-w,u+v)$ である。
交差上ではこれらのイデアルの最小生成元は二つで，
height 一の主イデアルではない。従って個々の成分は Cartier でない。
\end{proof}

\begin{corollary}[両積可逆なunordered良モジュライの明示開]
\label{cor:qa-both-invertible-good-moduli-affine-open}
$M^{44}\subset M$ を両方の積が階数四となる飽和開の good moduli とする。
このとき
\[
 M^{44}\simeq
 \mathbb A^3_{b,s,v}\setminus
 V(s-2v,\ v-4b+8).
 \tag{M44}\label{eq:qa-both-invertible-a3-minus-line}
\]
右辺で除いているものは一本の affine 直線である。
\end{corollary}
\begin{proof}
$D_+(d)$ では $C_0$ は空であり，$C_1$ の粗方程式は
\[
 u=v=w=4b-8
\]
となる。これは $\iota$ によって点ごとに固定される曲線であり，
その商像は $s=2v,\ v=4b-8$ である。
有限群の商の fiber は幾何点の軌道なので，この曲線は位相的に飽和し，
その補の商は表示した直線の補となる。
ただし直線の scheme としての逆像は，上の二重射の式から $z^2=0$ の厚みを持つ．
ここで使った飽和性は台についてのものである．
全 $M^{\mathrm{rul}}$ から $M$ への射は ruling 交換による有限全射の商である。
実際，構成に用いた $H$-半不変式 $h_i$ の共通 chart 上で
\[
 \bigl(k[X][h_i^{-1}]^G\bigr)^{S_2}
     =k[X][h_i^{-1}]^H
\]
となり，有限群の不変式環への拡大は有限整拡大である。
この三つの affine chart は両 good moduli を覆うので，有限全射性が従う。
両積可逆な開の ordered chart は $\Omega^\times$ である。
よって上の quotient 計算が正に $M^{44}$ を与える。
\end{proof}

\subsection{随伴tensorによる全域の既存GIT商への射}

\begin{theorem}[二つの全域OU射]
\label{thm:qa-global-two-ou-morphisms}
全 ordered good moduli から二つの正則射
\[
 \rho_f,\rho_g:M^{\mathrm{rul}}\longrightarrow\mathbb W
\]
が存在し，次の性質を持つ。
\begin{enumerate}
\item $\rho_f$ は $f$ が可逆な cubic chart を $\Omega$ と同型に写し，
      $g$ が可逆な chart では式~\eqref{eq:qa-cubic-weighted-cremona} で与えられる。
      $\rho_g$ には $f,g$ を交換した記述が成立する。
\item 集合としての像はともに $\Omega\cup\{q_0\}$ であり，
      $\rho_f^{-1}(\Omega)=M_f^{\mathrm{rul}}$，
      $\rho_g^{-1}(\Omega)=M_g^{\mathrm{rul}}$ が開部分として成立する。
      $\rho_f$ は $\mathfrak D_f$ の全閉像を $q_0$ に送り，
      $\rho_g$ は $\mathfrak D_g$ の全閉像を $q_0$ に送る。
      特に noncubic 閉曲線 $\overline N^{\mathrm{rul}}\simeq\mathbb A^1_\chi$ は
      どちらの射でも $q_0$ へ写る。
\item ruling 交換 $\sigma$ に対し
      $\rho_g=\rho_f\circ\sigma$ である。
\end{enumerate}
これらの射の像は稠密な真部分であり，後の系で示す通り射は proper でない。
また二つの射の積は noncubic 閉曲線を復元できない。
\end{theorem}
\begin{proof}
全 $V_{\mathrm{cross}}$ 上で
\[
 T_f=\operatorname{adj}(f)g:
      D\otimes B\longrightarrow C\otimes A\otimes L_f
\]
を用いる。通常の四因子の二対縮約としてその Gram 不変式を取り，
\[
 E_i^{(f)}=e_i(JT_f^{\mathrm{t}}JT_f)\quad(i=1,2,3),
 \qquad D^{(f)}=\det T_f
\]
と書く。
基底変更による二つの Segre 形式の similitude の比は
\[
 \chi_f^2\frac{ac}{bd}=\zeta_f
\]
である。従って四式
$E_1^{(f)},E_2^{(f)},D^{(f)},E_3^{(f)}$ の指標は
\[
 \zeta_f,\quad\zeta_f^2,\quad\zeta_f^2,\quad\zeta_f^3
\]
で，通常の $f,g$ 座標の総次数は $8,16,16,24$ である。

この四式が同時に零にならないことを確認する。
$f$ が可逆なら
$T_f=(\det f)f^{-1}g$ は geometric tensor の非零スカラー倍であり，
前節の GIT 比較により semistable である。
従ってその四つの基本不変式は同時に零にならない。

$\operatorname{rk}f=3$ の場合には
$\operatorname{adj}(f)=\kappa\lambda^{\mathrm{t}}$ と書け，
$T_f=\kappa\psi^{\mathrm{t}}$，$\psi^{\mathrm{t}}=\lambda^{\mathrm{t}}g$ である。
$V_{\mathrm{cross}}$ の条件と階数一テンソルの交差縮約の補題から，
$\kappa,\psi$ はともに Segre 二次形式に関して非等方である。
従って
\[
 E_1^{(f)}=(\kappa^{\mathrm{t}}J\kappa)(\psi^{\mathrm{t}}J\psi)\ne0,
 \qquad E_2^{(f)}=D^{(f)}=E_3^{(f)}=0.
\]
全 $V_{\mathrm{cross}}$ では $\operatorname{rk}f\ge3$ なので，これで全場合を尽くす。

以上の四式は，実際の正則射
\[
 V_{\mathrm{cross}}\longrightarrow\mathbb W,
 \qquad
 (f,g)\longmapsto
 [E_1^{(f)}:E_2^{(f)}:D^{(f)}:E_3^{(f)}]
\]
を定める。指標が重み $(1,2,2,3)$ に従っているため，この射は $G$-不変であり，
$\mathfrak H_2^{\mathrm{rul}}$ から $\mathbb W$ への射となる。
ここで指標は $\zeta_f$ であって $\chi_0$ ではない。
$R_G(\chi_0)$ への同次数環準同型を仮定するのではなく，
good moduli space の代数空間を標的とする普遍性
\cite[Theorem 6.6]{AlperGoodModuli2013} によって
この射を $\rho_f:M^{\mathrm{rul}}\to\mathbb W$ へ降下させる。
$g$ 側にも同じ構成を行う。

$f=1$ chart では定義がそのまま $T=f^{-1}g$ の四不変式であり，
$\rho_f$ は既存の $\Omega$ との chart 同型である。
$g=1$ chart では $T_f=\operatorname{adj}(g^{-1}f)$ であるため，
$f$ も可逆な稠密開では逆行列写像に一致する。
両辺は $\Omega$ から separated な $\mathbb W$ への正則射なので，
全 $g=1$ chart で同じ式となる。
階数三の場合の計算は，$\mathfrak D_f$ 全体の像が $q_0$ であることも示す。
この閉部分は被約であり，その good moduli 閉像も被約であるから，
閉像上でも射は定値射である。

飽和性により，$M^{\mathrm{rul}}$ の各点は
$f$ 可逆な chart または $\mathfrak D_f$ の閉像に属する。
前者の像は $\Omega$，後者の像は $q_0\notin\Omega$ なので，
集合としての像と開逆像の主張を得る。
noncubic 閉部分は両階数三であり，二つの射でともに $q_0$ へ写る。
最後に，ruling 交換は $T_f,T_g$ の構成を交換するから
$\rho_g=\rho_f\circ\sigma$ である。
\end{proof}

\begin{corollary}[全good moduliの非proper性]
\label{cor:qa-global-good-moduli-nonproper}
$\rho_f,\rho_g$ は proper でなく，$M^{\mathrm{rul}}$ および $M$ は
$k$ 上 proper でない。
また $(\rho_f,\rho_g):M^{\mathrm{rul}}\to\mathbb W\times\mathbb W$ は
noncubic 閉曲線全体を $(q_0,q_0)$ へ送るため，閉埋込みではない。
\end{corollary}
\begin{proof}
$\Omega\cup\{q_0\}$ は $\mathbb W$ の稠密な真部分であり，閉じていない。
Proper 射の像は閉じているから，$\rho_f,\rho_g$ は proper でない。
$M^{\mathrm{rul}}$ と $M$ には noncubic good moduli
$\mathbb A^1_\chi$ が閉埋め込まれている。
Proper なスキームの閉部分は proper であるから，両者は $k$ 上 proper でない。
最後の主張は，閉埋込みは幾何点で単射であることと，
$\overline N^{\mathrm{rul}}\simeq\mathbb A^1_\chi$ 上で二つの射がともに定値であることから従う。
\end{proof}

\begin{remark}[既存分岐式との一致]
$\Omega^\times$ における二つの固定成分の式は
$e_3\pm de_1=0$ である。
Newton の恒等式によって，既存の OU 座標では
$D_\pm=6(e_3\pm de_1)$ となる
（系~\ref{cor:qa-ou-coarse-branch-invariant-equations}）。
上の計算はその式を $e_1=e_3=0$ の重み二の locus まで含めて
粗スキーム上で検証し，その点での源の二次錐と標的の横断 branch を区別する。
一方，$\Omega$ の補曲線 $C_0,C_1$ は非 geometric tensor の閉集合であり，
全 $H_2$ の noncubic 閉曲線ではない。
全域射 $\rho_f,\rho_g$ はこの区別を保ちながら，後者を一点 $q_0$ に送る。

再現コード \texttt{verify\_cubic\_inversion\_global\_maps.py} は，
重み付き逆写像の二乗，非対角行列の Gram 逆行列式，
階数一 tensor の普遍不変式，粗環の二重被覆式と指標を厳密算術で検算する。
射の存在，貼合せおよび降下の証明は，上記の環と相対商の議論による。
\end{remark}

% !TeX root = ../nc-quadric-avatar.tex
% 2026-11-23 / v4.19。全H2良モジュライの境界像。

\subsection{階数境界の良モジュライ像と正規化}
\label{subsec:qa-global-good-moduli-boundary-images}

前項の記号を用い，$\pi:\mathfrak H_2\to M$ を構成済みの
good moduli space morphism とする。
$F=\det f$，$G_0=\det g$ と置く。
$\mathfrak D$ は $FG_0=0$ の unordered 階数因子，
$\mathfrak N\subset\mathfrak D$ は $F=G_0=0$ の noncubic 閉部分である。

\begin{proposition}[閉境界の良モジュライと階数の飽和性]
\label{prop:qa-global-closed-boundary-good-images}
次の主張が成立する。
\begin{enumerate}
\item $\mathfrak N$ の良モジュライ射
      $\chi:\mathfrak N\to\mathbb A^1_\chi$ は，標準的な閉埋込み
      $\iota:\mathbb A^1_\chi\hookrightarrow M$ を通じて $\pi$ の制限に一致する。
      $\mathfrak D$ にも良モジュライ射
      $\pi_D:\mathfrak D\to\overline D$ があり，
      $\overline D\hookrightarrow M$ は被約閉部分スキームである。
\item underlying topological space について
      \[
       |\pi|^{-1}(|\iota(\mathbb A^1)|)=|\mathfrak N|,
       \qquad
       |\pi|^{-1}(|\overline D|)=|\mathfrak D|.
      \]
      従って cubic 点と noncubic 点が同じ良モジュライ fiber に入ることはない。
      ordered 商では二つの開 $D(F),D(G_0)$ がそれぞれ飽和している。
\end{enumerate}
これらは閉部分の scheme-theoretic inverse image が元の閉部分に一致する，
という主張ではない。
\end{proposition}
\begin{proof}
有限型 good moduli space の閉部分については，
各 affine quotient chart 上の不変式の完全性を用いる。
実際，$B$ をその affine atlas の環，$I$ を安定イデアルとすると
\[
 (B/I)^\Gamma=B^\Gamma/I^\Gamma.
\]
従って閉部分自身の良モジュライが $M$ に閉埋め込みとなる。
同じことは $\pi_*$ を
$0\to\mathcal I\to\mathcal O\to\mathcal O/\mathcal I\to0$
へ適用しても分かる
\cite[Lemma 4.14, Theorem 4.16(iii)]{AlperGoodModuli2013}。
既存の noncubic 表示と
\cite[\texttt{cor:intrinsic-boundary-good-moduli}]{QuadricModuli2026}
による $\mathfrak N\to\mathbb A^1_\chi$ は good moduli space なので，
その一意性により最初の同定を得る。
$\mathfrak D$ の atlas は被約であり，不変式環も被約なので
$\overline D$ は被約である。

飽和性を別に証明する。
命題~\ref{prop:qa-global-cubic-opens-saturated} により、
ordered 商では $D(F),D(G_0)$ はそれぞれ飽和している。
そこでは $F^2P_g,F^3Q_g$ と交換側の半不変式を用いた。
以下も全て $V_{\rm cross}$ の中での開集合の等式として読む。

unordered 商でも直接同じことを確認できる。
\[
 u_0=(F^2P_g)^3,\quad u_1=(F^3Q_g)^2,
 \qquad
 v_0=(G_0^2P_f)^3,\quad v_1=(G_0^3Q_f)^2
\]
は全て指標 $\chi_0^{18}$ を持ち，ruling 交換は $u_i,v_i$ を交換する。
従って $u_i+v_i$ と $u_iv_i$ は $H$ の正の $\chi_0$ 冪の半不変式である。
体上の等式
$D(u)\cup D(v)=D(u+v)\cup D(uv)$ によって，
\[
 D(F)\cup D(G_0)=
 \bigcup_{i=0}^1\bigl(D(u_i+v_i)\cup D(u_iv_i)\bigr).
\]
この開は unordered 良モジュライについても飽和している。
その補集合が $|\mathfrak N|$ である。
また $FG_0$ 自身は $H$ の指標 $\chi_0^2$ の半不変式であり，
$D(FG_0)$ は飽和している。その補集合が $|\mathfrak D|$ である。
閉像の良モジュライへの全射性と合わせて表示した逆像の等式を得る。
この議論は開集合の等式を用いており，閉イデアルの pullback の等式は用いていない。
\end{proof}

\begin{lemma}[ruling 符号線束の直像消滅]
\label{lem:qa-noncubic-ruling-sign-pushforward-zero}
$p:\mathfrak N^{\rm rul}\to\mathfrak N$ を ruling の二重被覆とし，
\[
 p_*\mathcal O_{\mathfrak N^{\rm rul}}
       =\mathcal O_{\mathfrak N}\oplus\mathcal L_-
\]
を deck 対合の固有分解とする。このとき
\[
 \chi_*\mathcal L_-=0.
\]
特に $\chi\circ p:\mathfrak N^{\rm rul}\to\mathbb A^1_\chi$ も
good moduli space morphism であり，ruling 二重被覆は
良モジュライ上では恒等射を誘導する。
\end{lemma}
\begin{proof}
全ての noncubic 幾何点は二形式の中央二重平方表示を持つ。
\cite[\texttt{lem:noncubic-local-exchange}]{QuadricModuli2026}
は二形式を交換する線型変換と，$c,d$ を互いに移す次数付き自己同型を与える。
この自己同型は構造対象を保ち，$c,d$ の行列分解から来る
二つの spinor 対象を交換する。従って内在的 ruling を交換する
pointed 自己同値となる。
その慣性元は $\mathcal L_-$ の fiber に $-1$ で作用する。

任意の開 $U\subset\mathbb A^1$ 上の
$\chi_*\mathcal L_-$ の切断を取る。
その各幾何点での値は慣性不変であるから，標数零の下で零である。
$\mathfrak N$ は smooth，従って被約であるため，
smooth atlas 上でこの切断は零となり，もとの切断も零である。
これで直像の消滅を得る。
$p$ は有限射であり $\chi$ は cohomologically affine だから合成もそうで，
構造層の直像は上の固有分解と消滅によって
$\mathcal O_{\mathbb A^1}$ となる。最後の主張が従う。
\end{proof}

\begin{corollary}[階数因子の良モジュライ像は正規]
\label{cor:qa-rank-divisor-good-image-normal}
$\nu:\widetilde{\mathfrak D}=\mathfrak D_f\to\mathfrak D$
を unordered 階数因子の正規化とする。
前命題の $\overline D$ は正規であり，
\[
 \widetilde{\mathfrak D}\xrightarrow{\nu}\mathfrak D
        \xrightarrow{\pi_D}\overline D
\]
の合成も good moduli space morphism である。
従って $\nu$ が良モジュライに誘導する射は同型である。
stack 上の conductor は $\mathfrak N$ であるが，
$\overline D$ の正規化の conductor 閉部分は空である。
\end{corollary}
\begin{proof}
$j:\mathfrak N\hookrightarrow\mathfrak D$ とする。
ordered atlas 上の二枝の貼合せ列は
\[
 0\longrightarrow\mathcal O_{D_f\cup D_g}
 \longrightarrow\mathcal O_{D_f}\oplus\mathcal O_{D_g}
 \xrightarrow{\ (a,b)\mapsto a|_Z-b|_Z\ }
 \mathcal O_Z\longrightarrow0.
\]
ruling 交換は右端に符号を掛けるため，unordered 商への降下は
\[
 0\longrightarrow\mathcal O_{\mathfrak D}
 \longrightarrow\nu_*\mathcal O_{\widetilde{\mathfrak D}}
 \longrightarrow j_*\mathcal L_-
 \longrightarrow0
 \tag{BD}\label{eq:qa-boundary-branch-sign-sequence}
\]
となる。
$\pi_{D*}$ の完全性と直前の補題を適用すれば
\[
 \mathcal O_{\overline D}
   \xrightarrow{\ \sim\ }
   (\pi_D\circ\nu)_*\mathcal O_{\widetilde{\mathfrak D}}
\]
を得る。$\nu$ は有限だから合成は cohomologically affine でもあり，
これが good moduli space の主張である。

正規性は同じ affine quotient chart 上で直接確認できる。
各 $H$-半不変式 $h$ の chart で，正規化の atlas は
\[
 (D_f\cap D(h))\amalg(D_g\cap D(h))
\]
であり，二つの滑らかな affine scheme の和である。
これを $H$ で割った不変式環は，枝交換によって
$D_f\cap D(h)$ の $G$-不変式環に同型である。
正規整域の不変式環は正規であり，前段の直像同型によって
この環が $\overline D\cap D_+(h)$ の環に一致する。
従って $\overline D$ は正規である。
正規化が同型なので，その conductor イデアルは単位イデアルとなる。
\end{proof}

\begin{proposition}[noncubic 閉像の引戻しは一致点で厚みを持つ]
\label{prop:qa-noncubic-good-image-noncartesian}
二根対一致の noncubic 点の近傍では，
\[
 \mathfrak N\subset
 \mathfrak H_2\times_M\mathbb A^1_\chi
\]
は真の閉包含である。
右辺の underlying reduced closed substack は $\mathfrak N$ であるが，
scheme-theoretic fiber product は一致点で被約ではない。
\end{proposition}
\begin{proof}
一致点を含む $\pi$ の affine quotient chart
$[\operatorname{Spec}B/H]\to\operatorname{Spec}B^H$ を取り，
その atlas の点を $x$ とする。
$I\subset B$ は $\mathfrak N$ のイデアル，$J=I^H\subset B^H$
はその閉像のイデアルである。
\cite[\texttt{prop:noncubic-coincident-full-normal-representation}]
{QuadricModuli2026} の符号自己同値 $e_1$ は，
境界の二次元法空間とその双対に $-I_2$ で作用する。
従って $x$ の stabilizer の作用に関し
\[
 e_1|_{I B_x/\mathfrak m_x I B_x}=-I_2.
\]
任意の $b\in J$ は stabilizer 不変なので，この商への像は零である。
ゆえに
\[
 J B_x\subset\mathfrak m_x I B_x\subsetneq I B_x.
\]
最後の真の包含は $I$ が rank 二の conormal を持つ正則閉埋込みの
イデアルであること，または Nakayama の補題から従う。
左辺は fiber product のイデアルであるため，主張する真の包含を得る。
前命題の位相的飽和性と $\mathfrak N$ の被約性により
$\sqrt{J B_x}=I B_x$ だから，この fiber product は被約でない。
具体的には，$e_1$ が $-1$ で作用する接ベクトルで，
法空間への像が非零なものを取る。このベクトルは
$\operatorname{Spec}k[\varepsilon]/(\varepsilon^2)$ 上の実際の
$H_2$ 族を与える。全 invariant 関数の微分はこのベクトルを消すため，
その良モジュライ像は一定である。一方，法方向が非零なので
この族は $\mathfrak N$ には因数分解しない。
従って上の厚みは，同じ幾何点を持つ一次変形の差を実際に検出している。
\end{proof}

\begin{remark}[閉像と引戻しの区別]
noncubic 境界は位相的には $\pi$ の fiber の和であり，
その良モジュライ $\mathbb A^1_\chi$ は $M$ に閉埋め込まれる。
それでも invariant ideal の拡大は $(F,G_0)$ そのものとは限らず，
直前の命題は一致点で実際に厚みが生じることを示す。
また stack 上の Cartier 性や正規交差性を，良モジュライ像の
Cartier 性や正規交差性にそのまま移してはいない。
上の正規性は，二枝の差が ruling の符号線束となり，
その直像が消滅することを用いた別の結論である。
\end{remark}

% !TeX root = ../nc-quadric-avatar.tex
% 2026-11-23 / v4.22. 二つの四変数二次形式のpencilによる全半不変環。

\subsection{二つのGram二次形式のpencilと全大域不変式環}
\label{subsec:qa-global-h2-invariant-hull}

標数零の代数閉体 $k$ 上で、全$H_2$の線形表示
$X=\operatorname{Hom}(C\otimes A,W)\oplus
\operatorname{Hom}(D\otimes B,W)=\mathbb A^{32}$
と $G,H=G\rtimes S_2$、線形化 $\chi_0$ は
第~\ref{subsec:qa-global-h2-good-moduli}節のものとする。
$A,B,C,D$ は二次元、$W$ は四次元である。
Segre双線型形式 $J$ を式~\eqref{eq:qa-cross-pair-two-polynomials} の直前と
同じ規約で固定する。$f,g$ を二つの四次正方行列として、
\[
 S_f=fJf^{\mathrm t},\qquad S_g=gJg^{\mathrm t},\qquad
 F=\det f,\quad G_0=\det g
\]
とする。次の四次pencilの係数を定義する：
\[
 \det(XS_f+YS_g)
 =F^2X^4+I_1X^3Y+I_2X^2Y^2+I_3XY^3+G_0^2Y^4.
 \tag{GP}\label{eq:qa-global-gram-pencil}
\]
この $X,Y$ はpencilの形式変数であり、全係数空間$X$とは区別する。

\begin{theorem}[全半不変式環の明示表示]
\label{thm:qa-global-h2-full-semiinvariant-ring}
同じ $\chi_0$ に関する半不変式環を
$R_G=R_G(\chi_0)$、$R_H=R_H(\chi_0)$ と書く。
\[
 A_2=FG_0,\quad B_2=I_2,\quad
 X_3=FI_3,\quad Y_3=G_0I_1,\quad D_4=I_1I_3
\]
と置くと、添字は$\chi_0$次数であり、自然な次数付き環の同型
\[
 \boxed{\quad
 R_G\simeq
 \frac{k[A_2,B_2,X_3,Y_3,D_4]}
 {(X_3Y_3-A_2D_4)},\qquad
 \deg(A_2,B_2,X_3,Y_3,D_4)=(2,2,3,3,4).
 \quad}
 \tag{RG}\label{eq:qa-global-h2-ordered-semiinvariant-ring}
\]
が成立する。ruling交換は $X_3\leftrightarrow Y_3$、
$A_2,B_2,D_4$ を固定する。従って
$C_3=X_3+Y_3$ と置けば
\[
 \boxed{\quad
 R_H=R_G^{S_2}\simeq k[A_2,B_2,C_3,D_4],
 \qquad \deg(A_2,B_2,C_3,D_4)=(2,2,3,4).
 \quad}
 \tag{RH}\label{eq:qa-global-h2-unordered-semiinvariant-ring}
\]
特に、構成済みの全good moduliは明示的な開埋め込み
\[
 M\hookrightarrow\mathbb P(2,2,3,4),\qquad
 M^{\rm rul}\hookrightarrow
 \operatorname{Proj}k[A_2,B_2,X_3,Y_3,D_4]/(X_3Y_3-A_2D_4)
 \tag{WP}\label{eq:qa-global-h2-weighted-projective-hull}
\]
を持つ。これは同型類の粗モジュライ空間の存在を主張するものではない。
\end{theorem}
\begin{proof}
まず $SL(A)\times SL(C)$ が $C\otimes A$ に及ぼす作用は
$SO(J)$ の標準作用に一致する（中心の有限核は不変式に影響しない）。
直交群の第一基本定理を四つの列ベクトルに適用すると、
$f$ についての不変式は対称行列 $S_f$ の十成分と
向きを表す $F$ で生成され、唯一の関係は
$\det S_f=F^2$ である。$g$ についても同じなので、
$SL(A)\times SL(B)\times SL(C)\times SL(D)$ の不変式環は
\[
 \frac{k[(S_f)_{ij},(S_g)_{ij},F,G_0]}
 {(\det S_f-F^2,\ \det S_g-G_0^2)}.
\]
続いて $SL(W)$ の不変式を取る。
二つの対称四次形式に対する $SL_4$ 不変式環は、
四次pencilの五係数で自由に生成される
\cite[\S4.2]{BhargavaGrossWang2015}。
標数零で $SL(W)$ は線型簡約なので、
上記二つの不変な関係による剰余と不変式の取得は交換する。
従って
\[
 k[X]^{SL(W)\times SL_2^4}
 =k[F,G_0,I_1,I_2,I_3].
 \tag{P0}\label{eq:qa-global-h2-five-pencil-invariants}
\]
ここで既知の二形式pencilの不変式定理自体の新規性は主張しない。

残る五つの$GL$の行列式文字を
$(a,b,c,d,w)$ とし、第~\eqref{eq:qa-h2-git-characters}式を用いる。
$F,G_0$ の指標はそれぞれ $w/(ac)^2,w/(bd)^2$、
$I_j$ の指標は $w^2/(ac)^{4-j}(bd)^j$ である。
ゆえに単項式 $F^uG_0^vI_1^pI_2^qI_3^r$ が
$\chi_0$ の正の冪を持つ必要十分条件は
\[
 u+p=v+r.
 \tag{bal}\label{eq:qa-global-h2-balanced-exponents}
\]
その場合の$\chi_0$次数は $u+v+2(p+q+r)$ である。
この非負整数半群は、$FG_0,I_2,FI_3,G_0I_1,I_1I_3$
の五元で生成され、その関係は
$(FI_3)(G_0I_1)=(FG_0)(I_1I_3)$ の一つだけである。
例えば重みが釣り合うように左側の組$(F,I_1)$と
右側の組$(G_0,I_3)$から積を組み、対角的な四つの組と
$I_2$を取り出せば半群の表示が直接得られる。
これで式~\eqref{eq:qa-global-h2-ordered-semiinvariant-ring} が従う。

二積交換は二次形式pencilの両端と中間の係数
$I_1,I_3$ を交換する。
$X_3Y_3=A_2D_4$ と $C_3=X_3+Y_3$ により
$R_G=k[A_2,B_2,C_3,D_4,X_3]/(X_3^2-C_3X_3+A_2D_4)$
は $k[A_2,B_2,C_3,D_4]$ 上の自由階数二代数である。
交換が $X_3\mapsto C_3-X_3$ なので、不変部分は
ちょうど基底環となる。最後の開埋め込みは
定理~\ref{thm:qa-global-h2-quasiprojective-good-moduli} と
その同一指標による飽和開の定義から従う。
\end{proof}

\begin{lemma}[交差縮約の多項式と中央判別式]
\label{lem:qa-global-h2-gram-cross-polynomials}
第~\ref{subsec:qa-global-h2-good-moduli}節の
$P_f,Q_f,P_g,Q_g$ は、全$X$上で
\[
 \begin{aligned}
 16P_f&=I_1^2-4F^2I_2+8F^3G_0,\\
 16Q_f&=I_1^3-4F^2I_1I_2+16F^4I_3-8F^3G_0I_1,\\
 16P_g&=I_3^2-4G_0^2I_2+8G_0^3F,\\
 16Q_g&=I_3^3-4G_0^2I_3I_2+16G_0^4I_1-8G_0^3FI_3.
 \end{aligned}
 \tag{PQGP}\label{eq:qa-global-h2-cross-pencil-identities}
\]
特に $L_2=B_2-2A_2$ と置けば、重み八の式
\[
 \boxed{\quad
 \Delta_8=(D_4+4A_2L_2)^2-4L_2C_3^2
             =256P_fP_g
 \quad}
 \tag{D8}\label{eq:qa-global-h2-central-discriminant}
\]
が成立する。
\end{lemma}
\begin{proof}
$f$ を可逆とする稠密開では $f$ を単位行列に規格化し、
$T=f^{-1}g$ と置ける。
このchartの四量子ビット不変式は
$I_1=e_1(JT^{\mathrm t}JT)$、
$I_2=e_2(JT^{\mathrm t}JT)$、
$I_3=e_3(JT^{\mathrm t}JT)$、$G_0=\det T$ である。
$M_f=\operatorname{Cr}(T)$ と
$M_g=\operatorname{Cr}(\operatorname{adj}T)$ に対する
Hadamardの四固有値恒等式
（命題~\ref{prop:qa-cubic-adjugate-explicit-invariants} の証明）
をそれぞれ適用すると、表示した四式を得る。
実際に各式の左右は同じ$GL$指標を持つので、
規格化に用いた$F$の冪を復元した式が上記となる。
稠密開で一致する多項式は全$X$で一致する。

最後に第一式と第三式を掛け、
$X_3Y_3=A_2D_4$ と $C_3=X_3+Y_3$ を使うと
\[
 \begin{split}
 (16)^2P_fP_g
 &= (D_4+4A_2(B_2-2A_2))^2\\
 &\hspace{1cm}-4(B_2-2A_2)(X_3+Y_3)^2
 \end{split}
\]
を得る。これは多項式恒等式なので階数三の壁も含む。
\end{proof}

\begin{proposition}[中央affine chartとnoncubic閉曲線の全座標]
\label{prop:qa-global-h2-central-affine-ring}
全$H_2$の非順序good moduli $M$ において、
二つの交差行列式がともに可逆となる飽和開は
\[
 \boxed{\quad M_0=D_+(\Delta_8)\subset\mathbb P(2,2,3,4),\qquad
 \Gamma(M_0,\mathcal O)
 =\bigl(k[A_2,B_2,C_3,D_4,\Delta_8^{-1}]\bigr)_0.\quad}
 \tag{AFF}\label{eq:qa-global-h2-central-affine-ring}
\]
この開は全noncubic閉像を含む。
その閉像はこの開の中で
\[
 \boxed{\quad \overline N=V_+(A_2,C_3)\cap M_0
  \simeq\operatorname{Spec}k[\chi],
  \qquad\chi=\frac{B_2^2}{4D_4}.\quad}
 \tag{NC}\label{eq:qa-global-h2-noncubic-in-pencils}
\]
すなわち閉像は被約schemeとしても二つの明示方程式で与えられる。
\end{proposition}
\begin{proof}
定理~\ref{thm:qa-global-h2-quasiprojective-good-moduli} の被覆では
$s_0=(P_fP_g)^3$ である。前補題の式~\eqref{eq:qa-global-h2-central-discriminant}
により、その標準開は $D_+(\Delta_8)$ と同じであり、
その環は表示した重み零の局所化である。

noncubicでは $F=G_0=0$ であるから $A_2=C_3=0$。
逆に $V_{\rm cross}$ 上で $A_2=C_3=0$ とする。
仮に $F=0,G_0\ne0$ なら第一の交差行列式は
$16P_f=I_1^2$ であり、$f$ が階数三のため
$P_f\ne0$、従って $I_1\ne0$ である。
しかし $C_3=G_0I_1\ne0$ となり矛盾する。
$f,g$ を交換した場合も同じなので、両者は階数三である。
この場合 $P_fP_g=(I_1I_3)^2/256$ なので
$D_4\ne0$。従って $A_2=C_3=0$ と
$D_4\ne0$ の閉部分は $D_+(\Delta_8)$ に含まれ、
\[
  V_+(A_2,C_3)\cap M_0=D_+(D_4)\subset\mathbb P(2,4)
    =\operatorname{Spec}k[B_2^2/D_4]
\]
となる。

$\chi$ の係数を確認する。両階数三の中央二重平方正規形
（命題~\ref{prop:qa-double-rank-three-conic-recognition}）で、
$h_f=uv$、$h_g=(au+bv)(cu+dv)$ と置く。
$W$の基底を$(c,u,v,d)$に取り、乗法の行列を
\[
 f=\begin{pmatrix}
 1&0&0&1\\0&1&0&0\\0&0&1&0\\0&0&0&0
 \end{pmatrix},\qquad
 g=\begin{pmatrix}
 0&0&0&0\\0&a&c&0\\0&b&d&0\\1&0&0&1
 \end{pmatrix}
\]
と選ぶと、式~\eqref{eq:qa-global-gram-pencil} の直接展開で
\[
 I_1=-4,\qquad I_2=-8(ad+bc),\qquad
 I_3=-4(ad-bc)^2.
\]
従って
\[
 \frac{B_2^2}{4D_4}
   =\frac{(ad+bc)^2}{(ad-bc)^2}
   =\frac{D(h_f,h_g)^2}{\Delta(h_f)\Delta(h_g)}
   =\chi
\]
となり、右辺は論文Bの
\cite[\texttt{prop:moduli-noncubic-good-moduli}]{QuadricModuli2026}
における不変式である。
既存の閉埋め込み
（命題~\ref{prop:qa-global-closed-boundary-good-images}）
とこの座標同型により、非被約な余分のscheme構造を
閉像そのものへ付け加えずに式~\eqref{eq:qa-global-h2-noncubic-in-pencils}
を得る。閉像の逆像の厚みは
命題~\ref{prop:qa-noncubic-good-image-noncartesian} と矛盾しない。
\end{proof}

\begin{corollary}[良モジュライの二つの孤立商特異点]
\label{cor:qa-global-h2-two-quotient-singularities}
標数零の代数閉体上で、全good moduli $M$ の
特異点は次の二点だけである：
\[
 p_3=[0:0:1:0],\qquad
 p_4=[0:0:0:1]\quad\in\mathbb P(2,2,3,4).
\]
それぞれの形式局所解析型は
\[
 p_3:\ \tfrac13(1,1,2),\qquad
 p_4:\ \tfrac12(1,1,1).
\]
$p_4$ はnoncubic閉曲線 $\chi=0$ 上にあり、
$p_3$ はcubic開にある。
それ以外の全点で $M$ は滑らかである。
\end{corollary}
\begin{proof}
$p_4$ では$\Delta_8=D_4^2\ne0$。
$p_3$ では $A_2=B_2=D_4=0$、
$C_3\ne0$ である。関係式 $X_3Y_3=A_2D_4$ により
一方の$X_3,Y_3$だけが非零である。
式~\eqref{eq:qa-global-h2-cross-pencil-identities} から
$Q_fQ_g$ は非零となるため
$s_3=(Q_fQ_g)^2\ne0$。
従って二点とも$M$に属する。

重み三の座標$C_3$のchartでは、原点近傍は
$\mathbb A^3/\mu_3$ であり、残り三座標
$(A_2,B_2,D_4)$への重みは $(2,2,1)$。
これは型$\tfrac13(1,1,2)$である。
重み四の座標$D_4$のchartでは、$\mu_4$ が
$(A_2,B_2,C_3)$へ重み$(2,2,3)$で作用する。
位数二の部分群は$C_3$の符号だけを変えるため、
まずその不変式を取れば座標$(A_2,B_2,C_3^2)$を得る。
残余$\mu_2$は三座標を同時に反転するので、
型$\tfrac12(1,1,1)$となる。
その他の重み二座標のchartでは、残る
$\mu_2$ の非自明元は$C_3$だけを反転する擬鏡映であり、
粗商は滑らかである。各chartの重なりにも
新たな特異点はない。
最後に前命題の$\chi$座標を適用すれば
$p_4\in\overline N$、$\chi(p_4)=0$を得る。
\end{proof}

\begin{remark}[新しい環の計算と先行研究の切り分け]
\label{rem:qa-global-h2-invariant-hull-status}
二対称形式の不変式環と直交群の第一基本定理は既知の入力である。
ここで新たに得たのは、内在的$H_2$の共通線形化$\chi_0$に対する
全半不変環の正確な抽出、ruling交換の不変式、
明示的な中央affine補完、noncubic曲線の方程式および
二商特異点の決定である。これらの組合せの先行研究に対する
包括的新規性は未確認である。
全相対AS/QP群oidの同型やright-equivalenceの比較は
この不変式環の計算だけでは得られない。
\end{remark}

% !TEX root = ../nc-quadric-avatar.tex
% 2026-11-23 / 全基底上の有限表示とcanonicalな相対認識。
\subsection{全H2上の通常のsquare Morita代数}
\label{subsec:qa-universal-h2-square-morita81}

本項では古典型を含む全ての $\HHstack$ を扱う．
normal foundation の選択は有限 \'etale 二重被覆
$\mathfrak H_2^{\rm rul}\to\HHstack$ 上で行い，
実際の対象の代表とその Hom 束はさらに fppf 局所で取る．
そのような chart で
\[
 (E_0,E_1,E_2,E_3)=(\mathcal O,Y,X,M),
 \qquad V_i=\operatorname{Hom}(\mathcal O,E_i),
 \qquad (r_0,r_1,r_2,r_3)=(1,2,2,4)
\]
と置く．ここで $Y$ は normal 第四項の evaluation kernel であり，
$M$ は内在的な第一軌道対象である．
四対象を大域に順序付けたり，基底線束を自明化したりする必要はない．

\begin{proposition}[canonical な階数八十一の通常代数]
\label{prop:qa-universal-h2-square-morita81}
局所対象
\[
 T^{\rm can}=\bigoplus_{i=0}^3 E_i\otimes V_i^\vee
 \tag{T81}\label{eq:qa-h2-canonical-amplified-generator}
\]
は，基底線束による代表変更と二つの ruling の交換に自然に整合し，
全 $\HHstack$ の相対 heart の対象へ降下する．
その End は通常の有限局所自由代数
\[
 \mathcal A_{81}=\operatorname{End}(T^{\rm can}),
 \qquad \operatorname{rk}\mathcal A_{81}=81
\]
を定める．各相対 $H_2$ 族への引戻しで $T^{\rm can}$ は基底平坦であり，
共通 fppf 局所 chart 上の canonical compact 圏の tilting generator となる．
従って，実際の局所 Morita 同値
\[
 \operatorname{Perf}\mathcal A_{81}\simeq\mathcal C^{\rm perf}
\]
が得られる．$\mathcal O$ の直和因子に対応する大域冪等元 $e_0$ を持ち，
$e_0\mathcal A_{81}e_0=\OO$，$e_0\mathcal A_{81}\mapsto\mathcal O$ である．

代数，Morita datum，Hom 束および全積は任意の scheme 基底変更と可換する．
ただし相対 heart の主張は本稿の局所有限型 $k$-scheme の site で述べ，
大域の unbounded 導来圏と局所 compact 圏の descent の同値を新たに仮定しない．
\end{proposition}
\begin{proof}
まず局所 normal foundation を $(\mathcal O,X,M,Z)$ とする．
論文Bの相対 ruling mutation によって
\[
 Y=\operatorname{Cone}
 \bigl(\operatorname{Hom}(M,Z)\otimes M\to Z\bigr)[-1]
\]
を取り，$(\mathcal O,Y,X,M)$ は相対 full strong square 四項となる．
これは一般の $U_0$ 上でも成立する有限 Hom の計算である．
$\mathfrak F_{\HH}$ 上では $Y$ は他方の実際の normalized ruling であるので，
四対象はいずれも基底平坦な heart 対象となる．
系~\ref{cor:qa-h2-square-foundation-open} により二積は全射である．

幾何点のこの square は既存の square AS 分類の有限基礎でもある．
実際，論文Aの \cite[\texttt{thm:normal-helix}]{QuadricIntrinsic2026}
の normal helix を一項前から始めると
\[
 (H_0,H_1,H_2,H_3,H_4)=(t^{-1}Z,\mathcal O,X,M,Z)
\]
となる．同論文の \cite[\texttt{cor:normal-cubic-defect}]{QuadricIntrinsic2026}
は最小単純解消を与え，合成階数が四なら cubic，三なら最小十本矢型である．
定理~\ref{thm:qa-as10-inverse-all-targets} と
補題~\ref{lem:qa-as10-all-index-mutual-inverse} の逆操作は
\[
 \operatorname{coker}
 (t^{-1}Z\to\operatorname{Hom}(t^{-1}Z,\mathcal O)^\vee\otimes\mathcal O)
   \simeq Y
\]
を与え，この square を回収する．古典型では通常の
$(\OO,\OO(1,0),\OO(0,1),\OO(1,1))$ である．
従って局所有限代数 $\Lambda=\operatorname{End}(\bigoplus_iE_i)$ は，
補題~\ref{lem:qa-classified-fiber-finite-square-recognition} の範囲に入る．
ここで大域的な曲線 avatar が存在するとは仮定していない．

\emph{頂点線束の消去．}
二つの局所代表の比較で $E_i'=E_i\otimes L_i$ なら，
pointed な $\mathcal O$ の比較の下で $V_i'=V_i\otimes L_i$ である．
従って自然な縮約により
\[
 E_i'\otimes(V_i')^\vee
 \simeq E_i\otimes L_i\otimes V_i^\vee\otimes L_i^\vee
 \simeq E_i\otimes V_i^\vee
\]
を得る．比較射を単元倍した際，対象上の単元と Hom 双対上の逆単元は
厳密に相殺される．縮約は tensor の結合と整合するので，三重重なりの
scalar cocycle も消える．これは有限 square 表示の頂点トーラス
$K\simeq\Gm^3$ の全作用を消すことに等しい．

pointed exact 同値は二つの ruling の非順序対と $M$ を保つ．
ruling の交換は上の直和の二つの中間項だけを入れ替える．
有限直和の自然な置換は二回で恒等となり，既存の ruling 交換の coherent
$S_2$ 作用と整合する．よって $T^{\rm can}$ は ruling を選ばない
$\HHstack$ 上へ降下する．$\mathcal O$ 項と $M\otimes V_3^\vee$ 項は
個別に降下し，中間二項についてはその\emph{和}が降下する．

同じ議論は代数を直接構成しても確認できる．局所の $(j,i)$ corner は
\[
 \operatorname{Hom}(E_i,E_j)\otimes V_j^\vee\otimes V_i
\]
であり，頂点スカラーの指標
$(\lambda_j/\lambda_i)\lambda_j^{-1}\lambda_i$ は一である．
合成は中間の $V_j\otimes V_j^\vee$ の縮約なので，この降下は全積を保つ．

\emph{階数とtilting．}
局所 square の Hom 階数行列を
\[
 G=\begin{pmatrix}1&2&2&4\\0&1&0&2\\0&0&1&2\\0&0&0&1\end{pmatrix},
 \qquad r=(1,2,2,4)^{\mathsf T}
\]
とすれば，End の階数は
\[
 r^{\mathsf T}Gr
 =(1+4+4+16)+(4+4+16+16+16)=81
\]
である．各 $V_i$ を局所自明化すれば $T^{\rm can}$ は $E_i$ を
$r_i$ 回ずつ増幅した四項の直和であり，full strong 性と tilting が従う．
$r_0=1$ と単位 $V_0=\operatorname{End}(\mathcal O)=\OO$ によって
指定冪等元の同定も得る．
局所的な compact Morita datum の比較においても，前述の scalar cancellation
と $S_2$ の置換が cocycle を与える．
全 Hom 束は perfect-fiber 判定による有限局所自由束であり，
この対角・tensor・縮約の構成は任意基底変更に保たれる．
\end{proof}

\begin{remark}[大域の対角代数]
\label{rem:qa-h2-rank81-diagonal}
この構成は，二つの ruling を大域にラベル付けした四つの分裂頂点を
要求しない．ruling 選択被覆を $\pi:\mathfrak H_2^{\rm rul}\to\HHstack$ とすると，
相対的な対角代数は
\[
 \mathcal D=\OO\oplus\pi_*\mathcal D_{\rm rul}\oplus\mathcal D_M
\]
と書ける．$\mathcal D_{\rm rul}$ は選択被覆上で局所的に
$\operatorname{End}(V_X^\vee)$ となる次数二の Azumaya 代数，
$\mathcal D_M$ は局所的に $\operatorname{End}(V_3^\vee)$ となる
次数四の Azumaya 代数である．それらを大域的な行列代数と仮定しない．
対角の階数は $1+8+16=25$ である．
厳密に前向きの corner の和を $\mathcal J$ とすると，
\[
 \mathcal A_{81}=\mathcal D\oplus\mathcal J,
 \qquad \mathcal J^{\otimes_{\mathcal D}3}=0.
\]
従って大域の separable 対角を持つ有限 bar 解消を適用できる．
\end{remark}

\begin{theorem}[全H2のcanonical完成とpointed tails]
\label{thm:qa-universal-h2-rank81-completion}
上の通常代数に対して
\[
 \mathcal C_{81}
 =\bigoplus_{n\ge0}
 H^0\left(\Theta_{\mathcal A_{81}}^{
                   \otimes_{\mathcal A_{81}}^{\mathbf L}n}\right)
\]
を取る．全 tensor 冪は cohomological degree 零に集中し，
\[
 \operatorname{rk}(\mathcal C_{81})_n=81(2n+1)^2
 \qquad(n\ge0).
\]
全次数と全積は任意の scheme 基底変更と可換する．
零頂点 corner は canonical な反 Serre 環であり，
\[
 e_0\mathcal C_{81}e_0
   =\bigoplus_{n\ge0}\operatorname{Hom}(\mathcal O,t^n\mathcal O),
 \qquad
 \operatorname{rk}e_0(\mathcal C_{81})_ne_0=(2n+1)^2.
\]

本稿の site の任意の affine test $S\to\HHstack$ 上で，
$(\mathcal C_{81})_S$ は左右 Noetherian であり，実際の pointed exact 同値
\[
 \bigl(\qgr(\mathcal C_{81})_S,\,
       \pi e_0(\mathcal C_{81})_S,\,(1)\bigr)
   \simeq (\mathcal C_S,\mathcal O_S,t)
 \tag{QT81}\label{eq:qa-h2-rank81-pointed-tails}
\]
が存在する．この比較は通常の module-object 基底変更と coherent に整合する．
\end{theorem}
\begin{proof}
前注意の finite separable diagonal 解消により
命題~\ref{prop:qa-relative-serre-tensor-base-change} が適用できる．
局所 square の Morita 表示では，その第 $n$ tensor 冪は
\[
 \mathbf R\operatorname{Hom}(T^{\rm can},t^nT^{\rm can})
\]
に対応する．幾何 fiber の全整数 square helix の前向き Ext 消滅によって，
$n=0$ を含む全 $n\ge0$ で次数零に集中する．
perfect-fiber 判定は全基底上の局所自由性と任意基底変更を与える．
積は相対 Serre transport と合成によるものなので，同じ比較を保つ．

数値 Serre 行列を $S_G=G^{-1}G^{\mathsf T}$ とすると，fiber の階数は
$r^{\mathsf T}GS_G^{-n}r$ である．$(S_G-I)^3=0$ を用いた直接計算は
\[
 r^{\mathsf T}GS_G^{-n}r=81(2n+1)^2
\]
を与える．これは前向き Ext 消滅を別に証明した後の階数計算である．
$e_0$ が $\mathcal O$ を保つため，corner は全積を含めて
表示した canonical 反 Serre 環に一致する．

\emph{左右Noether性．}
まず各 $V_i$ と四頂点が分裂する局所 chart を取る．
$B=C(\Lambda)$ と書くと，$(\mathcal C_{81})_n$ は
$e_jB_ne_i$ を $r_j\times r_i$ 個に増幅した block 行列である．
各頂点の $r_i$ コピーから一つずつを選ぶ零次数冪等元を $f$ とすれば
\[
 f\mathcal C_{81}f\simeq B.
\]
零次数の同一頂点内の matrix units によって
$\mathcal C_{81}f\mathcal C_{81}=\mathcal C_{81}$ である．
従ってこれは通常の degree-zero graded Morita 同値である．
局所 $\Lambda$ の全幾何 fiber は分類済み square 型であり，
補題~\ref{lem:qa-classified-fiber-finite-square-recognition}，
\ref{lem:qa-square-foundation-opposite}，
\ref{lem:qa-serre-completion-opposite} によって $B$ は左右 Noetherian である．
この性質は通常の Morita 同値で保たれ，中央の faithfully flat 基底変更で
降下する．よって全ての指定された affine test 上で左右 Noether 性を得る．
ここで用いた full idempotent $f$ は各増幅頂点からの一コピーの和であり，
$\mathcal O$ だけの冪等元 $e_0$ とは区別する．

\emph{実際のheart．}
$T^{\rm can}$ は基底平坦 compact な対象の有限直和で，
$\mathcal O$ を直和因子に含む．各項から十分前向きの
$t^m\mathcal O$ への高次 Ext は既に消えている．
従って命題~\ref{prop:qa-relative-flat-foundation-az} により
$(T^{\rm can},t)$ は相対 AZ ample である．
実際の軌道環は，全積を含めて
\[
 \bigoplus_{n\ge0}\operatorname{Hom}(T^{\rm can},t^nT^{\rm can})
    =\mathcal C_{81}
\]
である．一般の有限 End 環に対する AZ 再構成により
式~\eqref{eq:qa-h2-rank81-pointed-tails} を得る．
これは $e_0$ の corner だけから graded Morita 同値を導いたものではない．
指定対象は $T^{\rm can}$ の実際の直和因子として識別される．
基底変更との比較は，標準射影のシフトを $t$ の軌道対象へ送り，
有限表示の余核へ延長する構成である．したがって低次数 section の
非平坦基底変更を仮定せずに，通常の module-object 引戻しとの
coherent な比較を得る．
\end{proof}

\begin{remark}[この通常代数の役割]
この結果は，古典型を含む全ての内在的 $H_2$ 族に，
指定構造対象を保持する通常の有限 Morita 代数とその canonical 全 Serre 完成を
付随させる．四つの rank-one 対象の大域的な代表や，ruling の順序は不要である．
一方，任意の AS/QP presentation stack を同定するには，
その presentation の相対平坦性，実際の全整数積，頂点線束と Serre transport，
および全同型層を比較する別の証明が必要である．
また任意の Noetherian $k$-基底へ相対 heart の定義を拡張したものとの
同一視や，その意味での strong Noether 性は本定理に含めない．
\end{remark}

% !TEX root = ../nc-quadric-avatar.tex
% 2026-11-23 / 全基底上の有限表示とcanonicalな相対認識。
\section{回転・mutationと相対heart}
\label{sec:qa-relative-rotated-mutated}

\subsection{回転と第一mutationの相対tails比較}
\label{subsec:qa-relative-rotated-mutated-tails}

本節の基底は，引き続き本稿の局所有限型 $k$-scheme の site に属する．
実際の曲線型 square avatar から構成した同じ相対 heart を
$\mathcal C_0$，四標準対象を $E_0,\ldots,E_3$ とする．
基底線束の因子を省略しない実際の coevaluation cone を
\[
 Z_1=\operatorname{Cone}
   \bigl(E_1\longrightarrow\operatorname{Hom}(E_1,E_3)^\vee\otimes E_3\bigr)
\]
と書き，canonical な $t=\mathsf S^{-1}[2]$ を用いる．
次の二つの四項窓を固定する：
\[
 \begin{split}
 \mathcal T_{\rm rot}&=(E_1,E_2,E_3,tE_0),\\
 \mathcal T_{\rm mut}&=(t^{-1}Z_1,E_0,E_2,E_3).
 \end{split}
 \tag{RM}\label{eq:qa-relative-rot-mut-windows}
\]
これらの直和を $T_{\rm rot},T_{\rm mut}$ と書く．

\begin{proposition}[二つの実際の相対tilting窓]
\label{prop:qa-relative-rot-mut-finite-tilting}
式~\eqref{eq:qa-relative-rot-mut-windows} の全対象は
$\mathcal C_0$ の基底平坦な heart 対象であり，共通 fppf 局所 chart で
compact である．両窓は相対 full strong exceptional foundation であり，
通常 End 代数
\[
 \Lambda_{\rm rot}=\operatorname{End}(T_{\rm rot}),\qquad
 \Lambda_{\rm mut}=\operatorname{End}(T_{\rm mut})
\]
は階数二十四の有限局所自由 directed 代数となる．
Hom 階数表はそれぞれ
\[
 G_{\rm rot}=
 \begin{pmatrix}
 1&0&2&6\\0&1&2&6\\0&0&1&4\\0&0&0&1
 \end{pmatrix},\qquad
 G_{\rm mut}=
 \begin{pmatrix}
 1&2&4&6\\0&1&2&4\\0&0&1&2\\0&0&0&1
 \end{pmatrix}.
 \tag{RH}\label{eq:qa-relative-rot-mut-rank-tables}
\]
それぞれの perfect 圏は，四つの実際の対象を保って
$\mathcal C_0^{\rm perf}$ と Morita 同値である．
代数，Morita datum，および全ての積は任意基底変更と可換する．
\end{proposition}
\begin{proof}
全 $E_i$ と normal foundation の第四項 $Z_1$ は既に
$\mathcal C_0$ の基底平坦 compact 対象である．
$t$ は heart 上完全であり基底線形なので，$tE_0$ と $t^{-1}Z_1$ にも
同じ性質が成立する．

幾何 fiber 上で第一窓は square helix の一頂点回転である．
第二窓は第一 mutation の全整数 helix の初期四項であり，
\[
 (H_0,H_1,H_2,H_3,H_4)
   =(t^{-1}Z_1,E_0,E_2,E_3,Z_1)
\]
となる．既存の固定 mutation と逆 mutation の比較により，
各 fiber は cubic または最小 AS10 の full strong foundation である．
この二型の Hilbert 関数は同じであり，両型を同一の minimal quiver と
みなす必要はない．

階数表を確認するには square の Euler 行列
\[
 G=\begin{pmatrix}1&2&2&4\\0&1&0&2\\0&0&1&2\\0&0&0&1\end{pmatrix}
\]
と Serre の数値行列 $S=G^{-1}G^t$ を用いればよい．
各 fiber の $K_0$ では
$[Z_1]=2[E_3]-[E_1]$ であり，
$[t^{-1}Z_1]=2[E_0]-[E_1]$ となる．
二つの基底変更行列は
\[
 Q_{\rm rot}=
 \begin{pmatrix}0&0&0&1\\1&0&0&-2\\0&1&0&-2\\0&0&1&4\end{pmatrix},
 \qquad
 Q_{\rm mut}=
 \begin{pmatrix}2&1&0&0\\-1&0&0&0\\0&0&1&0\\0&0&0&1\end{pmatrix}
\]
であり，$Q^tGQ$ が式~\eqref{eq:qa-relative-rot-mut-rank-tables} を与える．
これは fiber の数値計算であり，実際の相対 cone の係数束を
自明な二次元束と同定しているわけではない．

全対象の相対 Hom は perfect であり，各導来幾何 fiber は次数零に
集中する．perfect-fiber 判定により Hom は指定階数の局所自由束となり，
任意基底変更と可換する．有限 directed 代数の相対 bar 解消と
補題~\ref{lem:qa-relative-avatar-perfect-tilting-detection} により，
幾何 fiber の full 性も実際の相対 tilting と Morita 同値へ移る．
\end{proof}

\begin{theorem}[回転・mutation完成の同じ相対heart]
\label{thm:qa-relative-rot-mut-completion-heart}
$\bullet\in\{\mathrm{rot},\mathrm{mut}\}$ に対して
\[
 C_\bullet=\bigoplus_{n\geq0}
 H^0\bigl(\Theta_{\Lambda_\bullet}^{\otimes_{\Lambda_\bullet}^{\mathbf L}n}\bigr)
\]
と置く．全 tensor 冪は次数零に集中し，各次数は有限局所自由である．
その形成と積は任意の scheme 基底変更と可換する．
全次数の階数は両場合とも
\[
 \operatorname{rk}(C_\bullet)_n=64n^2+64n+24\qquad(n\geq0)
 \tag{RHC}\label{eq:qa-relative-rot-mut-completion-ranks}
\]
である．

本稿の site の全 affine test 上で $C_\bullet$ は右 Noetherian であり，
実際の相対完全同値
\[
 \qgr C_{\rm rot}\xrightarrow{\sim}\mathcal C_0,
 \qquad
 \qgr C_{\rm mut}\xrightarrow{\sim}\mathcal C_0
 \tag{RMT}\label{eq:qa-relative-rot-mut-tails}
\]
が存在する．標準冪等元に対応する四対象はそれぞれ
$\mathcal T_{\rm rot},\mathcal T_{\rm mut}$ に移り，次数移動は canonical
$t$ に移る．四頂点の unfolding による全整数代数についても同じ結論となる．
通常の module-object 引戻しとの比較は coherent である．
\end{theorem}
\begin{proof}
共通 fppf chart 上の Morita 同値は相対 Serre 双対を保つ．
従って第 $n$ tensor 冪は
\[
 \mathbf R\operatorname{Hom}(T_\bullet,t^nT_\bullet)
\]
に対応する．各幾何 fiber では，$n=0$ の strong 性および $n\geq1$ の
全整数 geometric helix の forward Ext 消滅によって次数零に集中する．
各固定 $n$ ごとの perfect-fiber 判定により，全次数で局所自由性と
任意基底変更整合性を得る．cubic と AS10 の境界を横断する非被約基底にも
同じ議論が成立する．

階数は各 fiber で $Q_\bullet^tGS^{-n}Q_\bullet$ の成分和である．
$(S^{-1}-I)^3=0$ なので
\[
 S^{-n}=I+n(S^{-1}-I)+\frac{n(n-1)}2(S^{-1}-I)^2
\]
を代入すると式~\eqref{eq:qa-relative-rot-mut-completion-ranks} を得る．

$T_{\rm mut}$ は $E_0=\mathcal O$ を直和因子に含み，
$T_{\rm rot}$ は $tE_0$ を含む．$(tE_0,t)$ も AZ ample である．
各窓の対象から $t^mE_0$，または $t^m(tE_0)$ への十分前向きの
高次 Ext は，先ほどと同じ perfect-fiber 判定により零となる．
命題~\ref{prop:qa-relative-flat-foundation-az} を，第二の場合には
基準対象を $tE_0$ に換えて適用すると，$(T_\bullet,t)$ は AZ ample である．
この命題の有限個の affine fppf chart からなる被覆と Hom の flat base change による降下は，
大域の affine test 上で AZ2 を与える．

一般の有限 End 環に対する AZ 再構成を適用し，
\[
 \bigoplus_{n\geq0}\operatorname{Hom}(T_\bullet,t^nT_\bullet)
       =C_\bullet
\]
を用いれば右 Noether 性と式~\eqref{eq:qa-relative-rot-mut-tails} を得る．
この軌道環の同定は局所 compact Morita 比較を H0 の flat base change に
沿って降下したものである．
通常の基底変更との整合性は，標準射影のシフトを
$t^{-n}\mathcal T_{\bullet,i}$ に送り，有限表示の余核へ延長する構成から従う．
unfolding は四冪等元の成分を並べ替えるため，有限生成性と torsion を保つ．
\end{proof}

\begin{remark}[構造対象の指定]
式~\eqref{eq:qa-relative-rot-mut-tails} は同じ実際の Abel 圏との比較であり，
標準零対象まで $E_0$ に移すという主張ではない．
回転側の零対象は $E_1$，mutation 側の零対象は $t^{-1}Z_1$ に移る．
これらが強い $H_2$ pointing となる開部分は，既存の標準頂点判定で
個別に指定する．特に最小 AS10 fiber の零対象が $H_2$ でないことと，
その tails heart が $\mathcal C_0$ であることは両立する．
\end{remark}

\subsection{最小AS10表示とschemeとしての壁}
\label{subsec:qa-relative-minimal-as10-scheme-wall}

\begin{proposition}[最小矢束の平坦性はscheme壁を検出する]
\label{prop:qa-relative-as10-minimal-arrow-wall}
第一 mutation に対する乗法を
\[
 \mu_G:\operatorname{Hom}(E_2,E_3)\otimes\operatorname{Hom}(E_0,E_2)
             \longrightarrow\operatorname{Hom}(E_0,E_3)
\]
とし，$Z_F=V(\det\mu_G)$ と置く．局所的には
$\mu_G\sim\operatorname{diag}(I_3,\delta)$ である．
mutation の全整数代数で，最小矢の $1\to3$ 成分は
\[
 (J/J^2)_{3,1}=\operatorname{coker}\mu_G
                 \simeq\mathcal O_S/(\delta)
 \tag{MA}\label{eq:qa-relative-as10-extra-arrow-cokernel}
\]
となる．ここで $J$ は正の頂点差の射からなる ideal である．

したがって，指定の余分な最小矢が基底上階数一の線束となるための
必要十分条件は，基底の射が $Z_F$ を scheme として通ることである．
その壁上では，局所 Schur reduction により二次部分が零の reduced
最小十本矢表示を取れる．一方，壁の補集合では四組の消去により
八本矢 cubic 表示となる．
\end{proposition}
\begin{proof}
mutation 窓の $(H_1,H_2,H_3)=(E_0,E_2,E_3)$ において，
$1$ と $3$ の間の中間頂点は $2$ だけである．
従って分解可能射の像はちょうど $\operatorname{im}\mu_G$ であり，
式~\eqref{eq:qa-relative-as10-extra-arrow-cokernel} を得る．
局所正規形から cokernel は $\mathcal O_S/(\delta)$ である．
これが基底上の線束なら，局所的な全射
$\mathcal O_S\to\mathcal O_S/(\delta)$ は階数一束の同型であり，
$\delta=0$ となる．逆は明らかである．

相対 square 関係束の完全列を
\[
 0\longrightarrow R\longrightarrow T_F\oplus T_G
       \xrightarrow{\mu_F+\mu_G}W\longrightarrow0
\]
とする．$R\to T_F$ の cokernel は $\operatorname{coker}\mu_G$ と同型である．
これは $T_F$ の元を $W/\operatorname{im}\mu_G$ へ送る完全列から従う．
従って premutation の二次 pairing も局所的に
$\operatorname{diag}(I_3,u\delta)$，$u$ は単元，の形になる．
$\delta=0$ の壁上では kernel と cokernel は実際の線束であり，
三つの可逆な対を消去する．命題~\ref{prop:qa-as10-reduction} の
変数変換は可逆行列と有限回の多項式置換だけからなるので，この
scheme 壁上でも同じ計算が成立し，
\[
 W_{10}=yq+px+H
\]
という reduced 十本矢表示を得る．係数は各矢の線束・束に保持する．
$\delta$ が可逆なら第四の対も消去でき，cubic 八本矢表示となる．
代数族の平坦性と tails の同一性は，この表示の矢数から推論するのではなく，
前定理の実際の Hom と Serre 完成によって既に証明されている．
\end{proof}

\begin{example}[全fiberがAS10でも最小矢束が平坦でない場合]
普遍壁に横断する一次方向を
$S=\operatorname{Spec}k[\epsilon]/(\epsilon^2)$ に引き戻して
$\mu_G=\operatorname{diag}(1,1,1,\epsilon)$ とする．
唯一の幾何 fiber は階数三であり，mutation の fiber は最小 AS10 型である．
しかし $\epsilon\ne0$ なので $S$ は scheme として $Z_F$ を通らず，
\[
 (J/J^2)_{3,1}=k[\epsilon]/(\epsilon)
\]
は基底上平坦ではない．従って最小矢束が階数一の線束である
相対 reduced AS10 表示にはならない．

三つの可逆な対だけを消した局所表示には，単元倍を吸収して
\[
 W_\epsilon=\epsilon yx+yq+px+H
\]
という二次項が残る．これは十個の選択された生成矢を使う表示を
禁止するという意味ではない．その表示には nilpotent な線形関係
$\epsilon x+q$，$\epsilon y+p$ があり，最小矢の quotient が
線束になる reduced 十本矢族とは異なる．
それでも $\Lambda_{\rm mut}$，$C_{\rm mut}$ とその各 Hom 成分は
前定理によって平坦であり，その相対 tails は同じ $\mathcal C_0$ である．
\end{example}

\subsection{双対square心による回転・mutationの反対代数}
\label{subsec:qa-relative-rot-mut-opposite-heart}

引き続き実際の相対 square 有限代数を $\Lambda$ とする．
補題~\ref{lem:qa-square-foundation-opposite} と
\ref{lem:qa-classified-fiber-finite-square-recognition} により，
反対代数 $\Lambda^{\rm op}$ の逆順の標準射影に対応する相対 $H_2$ heart
$\mathcal C^\vee$ を取れる．この heart の四標準対象を
$Q_0,Q_1,Q_2,Q_3$ とし，$Q_0$ を構造対象とする．
その canonical な反 Serre 関手を $t^\vee$ と書く．

共通の fppf 局所 perfect chart 上で
\[
 \mathbb D=\mathbf R\operatorname{Hom}_\Lambda(-,\Lambda):
       \Perf\Lambda\xrightarrow{\sim}(\Perf\Lambda^{\rm op})^{\rm op}
\]
を用いれば，$Q_i=\mathbb D E_{3-i}$ である．
これは perfect 圏の反同値であり，元の heart の全対象に対する
反完全同値を仮定しているのではない．

\begin{lemma}[反同値とcanonical周期]
\label{lem:qa-relative-dual-reverses-antiserre}
任意基底変更と整合する自然な同型
\[
 \mathbb D t\simeq(t^\vee)^{-1}\mathbb D
\]
がある．従って perfect 対象の四項窓 $X_0,\ldots,X_3$ に対し，
$Y_i=\mathbb D X_{3-i}$ と置くと
\[
 \mathbf R\operatorname{Hom}(Y_i,(t^\vee)^nY_j)
 \simeq
 \mathbf R\operatorname{Hom}(X_{3-j},t^nX_{3-i})
 \tag{DF}\label{eq:qa-dual-window-forward-hom}
\]
が成立する．
\end{lemma}
\begin{proof}
相対 Serre 双対の表現対象の一意性から，反同値は Serre 関手を
逆 Serre 関手に送る．すなわち
$\mathbb D\mathsf S\simeq(\mathsf S^\vee)^{-1}\mathbb D$ である．
さらに $\mathbb D(E[2])=(\mathbb D E)[-2]$ なので
\[
 \mathbb D(\mathsf S^{-1}E[2])
     =\mathsf S^\vee\mathbb D(E)[-2]
     =(t^\vee)^{-1}\mathbb D(E)
\]
を得る．反同値の全忠実性とこの同型を使えば
式~\eqref{eq:qa-dual-window-forward-hom} が従う．
全ての比較は有限 perfect 双対と相対 Serre 比較から作られるため，
任意基底変更と可換する．
\end{proof}

\begin{lemma}[例外対象のfiberwise同型と相対線束]
\label{lem:qa-relative-exceptional-line-comparison}
相対 smooth proper な perfect 圏の二対象 $A,B$ が，各幾何 fiber で
同じ例外対象に同型であると仮定する．このとき
\[
 L=\mathbf R\operatorname{Hom}(A,B)
\]
は次数零の可逆層であり，評価は実際の同型
$L\otimes A\xrightarrow{\sim}B$ を与える．
この構成は任意基底変更と可換する．
\end{lemma}
\begin{proof}
相対 proper 性と perfectness により $L$ は基底上 perfect で，導来
基底変更と可換する．全幾何 fiber が $k[0]$ なので，perfect-fiber 判定で
$L$ は次数零の線束となる．評価の cone は相対 perfect であり，
各 fiber では同じ例外対象の一次元 Hom による評価だから零となる．
再び perfect-fiber 判定を適用すると cone は零である．
これは任意の二つの対象の fiberwise heart 条件を持ち上げる主張ではない．
cross-Hom が相対 perfect であることと，fiber の例外性が本質である．
また $L$ の自明性は主張していない．
\end{proof}

\begin{proposition}[二つの反対窓は同じ双対heartに入る]
\label{prop:qa-relative-rot-mut-dual-windows}
二つの逆順双対窓を
\[
 \begin{split}
 \mathcal Y_{\rm rot}
   &=((t^\vee)^{-1}Q_3,Q_0,Q_1,Q_2),\\
 \mathcal Y_{\rm mut}
   &=(Q_0,Q_1,Q_3,K),\qquad K=\mathbb D(t^{-1}Z_1)
 \end{split}
 \tag{DW}\label{eq:qa-relative-dual-windows}
\]
とする．その全対象は $\mathcal C^\vee$ の基底平坦な heart 対象である．
それぞれ相対 full strong tilting foundation であり，通常 End は
$\Lambda_{\rm rot}^{\rm op}$，$\Lambda_{\rm mut}^{\rm op}$ と同型となる．
\end{proposition}
\begin{proof}
式~\eqref{eq:qa-relative-dual-windows} は前補題から従う．
全対象の compactness，full 性，strong 性，および End の反対代数との
同定は perfect 反同値によって元の窓から移る．
回転側は全て $Q_i$ とその canonical 周期移動なので，基底平坦な
heart 対象であることも直ちに従う．

mutation 側の $K$ については，fiberwise な heart 条件をそのまま
相対 heart 条件に読み替えず，前補題の例外対象の比較を使う．
まず各幾何 fiber では，補題~\ref{lem:qa-as10-all-index-mutual-inverse}
の全整数 mutation の評価列
\[
 0\longrightarrow H_0\longrightarrow V\otimes E_0
       \longrightarrow E_1\longrightarrow0,
 \qquad V=\operatorname{Hom}(E_0,E_1),\quad H_0=t^{-1}Z_1
\]
がある．これに perfect 反同値を適用すると
\[
 Q_2\longrightarrow V^\vee\otimes Q_3
       \longrightarrow K\longrightarrow Q_2[1]
\]
を得る．$\operatorname{Hom}(Q_2,Q_3)=V$ であり，最初の射は coevaluation
なので，この fiber の $K$ は dual square のもう一方の normal 窓の
第四対象 $Z_2^\vee$ である．この実際の四項は
\[
 (Q_0,Q_1,Q_3,Z_2^\vee),\qquad
 Z_2^\vee=\operatorname{Cone}
  \bigl(Q_2\to\operatorname{Hom}(Q_2,Q_3)^\vee\otimes Q_3\bigr)
\]
であり，命題~\ref{prop:qa-relative-avatar-normal-window} の
$F,G$ を交換した窓である．その normal 性は
定理~\ref{thm:qa-relative-square-e0-e3-h2} の体上の比較，または有限
square 族に対する補題~\ref{lem:qa-classified-fiber-finite-square-recognition}
で既に確認されている．よって
$(Q_0,Q_1,Q_3,K)$ の各幾何 fiber は真正な normal foundation である．

実際の $Z_2^\vee$ は dual square の相対 normal cone として存在し，
基底平坦な heart 対象である．前補題をこの $Z_2^\vee$ と $K$ に適用すると
\[
 L=\mathbf R\operatorname{Hom}(K,Z_2^\vee),\qquad
 L\otimes K\xrightarrow{\sim}Z_2^\vee
\]
を得る．従って $K\simeq Z_2^\vee\otimes L^{-1}$ は実際に
$\mathcal C^\vee$ の基底平坦対象である．
この線束比較は任意基底変更と整合するが，$L$ を自明とした厳密な
第四対象の同定までは要求しない．
\end{proof}

\begin{theorem}[回転・mutation完成の左右Noether性]
\label{thm:qa-relative-rot-mut-two-sided-noetherian}
実際の相対曲線 avatar から構成した
$C_{\rm rot},C_{\rm mut}$ は，本稿の site の全 affine test 上で
左右 Noetherian である．さらに
\[
 \qgr C(\Lambda_{\rm rot}^{\rm op})\simeq\mathcal C^\vee,
 \qquad
 \qgr C(\Lambda_{\rm mut}^{\rm op})\simeq\mathcal C^\vee
\]
は，それぞれ式~\eqref{eq:qa-relative-dual-windows} の四対象と
canonical $t^\vee$ を保つ相対完全同値である．
\end{theorem}
\begin{proof}
前命題の二つの窓はいずれも $Q_0$ を直和因子に含む．
式~\eqref{eq:qa-dual-window-forward-hom} と元の全 forward Ext 消滅から，
各対象から $(t^\vee)^mQ_0$ への十分前向きの高次 Ext は零となる．
全対象が基底平坦 compact なので
命題~\ref{prop:qa-relative-flat-foundation-az} を適用できる．
有限個の affine fppf chart からなる被覆と coherent Hom の flat base change に沿った同命題の
降下を含めて，二つの窓の直和と $t^\vee$ は全 affine test 上で AZ ample
となる．通常の軌道環はそれぞれ $C(\Lambda_\bullet^{\rm op})$ であり，
一般 End 環版 AZ 再構成から右 Noether 性と表示した相対 tails 同値を得る．

補題~\ref{lem:qa-serre-completion-opposite} により
\[
 C(\Lambda_\bullet^{\rm op})\simeq C(\Lambda_\bullet)^{\rm op}
\]
である．従ってその右 Noether 性は $C_\bullet$ の左 Noether 性である．
右 Noether 性は定理~\ref{thm:qa-relative-rot-mut-completion-heart} で
既に得ているので結論が従う．
この証明は左 Noether 性を finite foundation の認識や最初の
relative tails 構成の入力に使っていない．
\end{proof}

\begin{remark}[非被約壁を横断する族]
この左 Noether 性の証明も，最小矢束の平坦性を仮定しない．
特に $\mu_G\sim\operatorname{diag}(I_3,\epsilon)$ の一次変形に対しても
成立する．そこで相対 reduced 最小 AS10 presentation を持たないことは，
通常の全 Hom 代数が左右 Noetherian であることを妨げない．
任意の site 外の Noetherian 基底に対する strong Noether 性や，
任意の AS/QP presentation stack の同一視は追加で主張しない．
\end{remark}

\subsection{回転の添字移動とmutationの相対非最小解消}
\label{subsec:qa-relative-mutation-nonminimal-as}

前節で構成した有限代数 $\Lambda_{\rm mut}$ とその通常の Serre 完成を
用いる．以下の表示は基底の局所 chart 上で行う．線束や高階数の係数束を
自明化しない場合には，tensor と双対を各矢束に施した同じ式と解釈する．

\begin{proposition}[rotated展開は実際の添字移動]
\label{prop:qa-relative-rotated-index-shift}
canonical な全整数展開について
\[
 (A_{\rm rot})_{ba}=(A_{\rm sq})_{b+1,a+1}
\]
という全積を保つ自然な同型がある．従って square の相対 AS 単純解消，
その双対，および任意基底変更との整合性は，添字を一つずらすだけで
rotated のものへ移る．
\end{proposition}
\begin{proof}
$R_i=E_{i+1}$ ($0\leq i\leq3$) であり，$R_3=tE_0=E_4$ である．
canonical $t$ で延長すると全整数で $R_a=E_{a+1}$ となる．
両展開はこれらの実際の対象の Hom と合成で定義されているので，
表示した全次数の代数同型が従う．
\end{proof}

第一 mutation の四項を $(H_0,H_1,H_2,H_3)$ とし，
$\Gamma=\Lambda_{\rm mut}$ と置く．
\[
 V_{10}=\operatorname{Hom}(H_0,H_1),\quad
 V_{21}=\operatorname{Hom}(H_1,H_2),\quad
 V_{32}=\operatorname{Hom}(H_2,H_3)
\]
はいずれも階数二である．$V_{21}\otimes V_{10}\to\operatorname{Hom}(H_0,H_2)$
は同型であり，
\[
 \mu:V_{32}\otimes V_{21}\longrightarrow\operatorname{Hom}(H_1,H_3)
\]
は既存の $\mu_G$ である．

\begin{lemma}[七生成矢・三関係による有限Endの表示]
\label{lem:qa-relative-mutation-seven-three-presentation}
三次 minor を可逆にした局所 chart 上で，$\Gamma$ は次の実際の表示を持つ．
六つの隣接生成矢に $x:1\to3$ を加え，$x$ の頂点差を二とする．
関係束は $1\to3$ の階数一と $0\to3$ の階数二であり，
局所枠を選ぶと
\[
 r=\delta x+q,
 \qquad r_\alpha\in
 (V_{32}\otimes V_{21}\otimes V_{10})
       \oplus (\mathcal O_Sx\otimes V_{10}),\quad \alpha=1,2
 \tag{73}\label{eq:qa-relative-mutation-seven-three}
\]
と書ける．$q\in V_{32}\otimes V_{21}$ は二段階道の部分束の基底の一つに
選べる．従って $q$ の少なくとも一係数は単元である．
$\delta$ は $\det\mu_G$ の局所方程式の単元倍であり，nilpotent でもよい．
\end{lemma}
\begin{proof}
三つの二段階道とその像を部分束の基底に選び，第四の道からその三つの
成分を引いて $q$ とする．標的の第四補元を $x$ として選ぶと，符号を
吸収して $\mu(q)=-\delta x$ である．従って
\[
 (V_{32}\otimes V_{21})\oplus\mathcal O_Sx
       \longrightarrow\operatorname{Hom}(H_1,H_3)
\]
は全射で，その kernel は $r=\delta x+q$ の張る線束である．
$q$ は二段階道の基底の一つなので，これは全道束の部分束である．

$0\to3$ の全道束は上の階数五束と $V_{10}$ の tensor，従って階数十である．
合成
\[
 \operatorname{Hom}(H_1,H_3)\otimes V_{10}
       \longrightarrow\operatorname{Hom}(H_0,H_3)
\]
は全幾何 fiber で全射である．cubic と AS10 の有限 quiver はともに，
頂点 $0$ から出る最小矢を $0\to1$ の二本だけに持つためである．
よって基底上でも全射であり，階数八から六へのこの写像の kernel は
階数二の束となる．

階数十の全道束から $\operatorname{Hom}(H_0,H_3)$ への kernel を $R_{30}$
と書くと，実際の束の完全列
\[
 0\longrightarrow\mathcal O_Sr\otimes V_{10}
       \longrightarrow R_{30}\longrightarrow\overline R_{30}\longrightarrow0,
 \qquad\operatorname{rk}R_{30}=4,\quad
 \operatorname{rk}\overline R_{30}=2
\]
がある．$r\otimes V_{10}$ は $q$ の単元係数により全道束の部分束であり，
上の商も実際に局所自由である．この列を局所的に分裂させ，
$\overline R_{30}$ の二つの基底の lift を $r_1,r_2$ と取る．

有限 quiver の他の頂点対には関係はなく，$0\to3$ の kernel は
$rV_{10}$ と $r_1,r_2$ で尽くされる．従って通常の path algebra の
relation ideal は正確に $(r,r_1,r_2)$ であり，表示した quotient は
全積を含めて実際の $\Gamma$ である．この証明は base の被約性を使わない．
\end{proof}

\begin{proposition}[有限対角解消とordinary Jacobian]
\label{prop:qa-relative-mutation-ordinary-jacobian}
前補題の七矢束を $V$，三関係束を $R$，頂点環を
$D_0=\bigoplus_{i=0}^3\mathcal O_Se_i$ とする．
この表示は任意基底変更に保たれる短い相対対角解消
\[
 0\longrightarrow\Gamma\otimes_{D_0}R\otimes_{D_0}\Gamma
 \longrightarrow\Gamma\otimes_{D_0}V\otimes_{D_0}\Gamma
 \longrightarrow\Gamma\otimes_{D_0}\Gamma
 \longrightarrow\Gamma\longrightarrow0
 \tag{MD}\label{eq:qa-relative-mutation-short-diagonal}
\]
を持つ．関係の双対戻り矢を加えた十矢 potential は
\[
 W_\delta=\delta yx+yq+px+H
 \tag{MP}\label{eq:qa-relative-mutation-delta-potential}
\]
と書け，通常の全次数代数として
\[
 C(\Gamma)\simeq\operatorname{Jac}(V\oplus R^\vee,W_\delta)
 \tag{MJ}\label{eq:qa-relative-mutation-ordinary-jacobian}
\]
が成立する．arrow-adic completion は用いない．
\end{proposition}
\begin{proof}
式~\eqref{eq:qa-relative-mutation-short-diagonal} の第一微分は
relation の普遍微分，次は矢の左右乗法の差である．
relation は実際の $\Gamma$ で零なので合成は零である．
各幾何 fiber で $\delta=0$ なら，これは既存の最小 AS10 finite foundation
の七矢三関係による対角解消である．
$\delta\ne0$ なら $r$ から $x$ を消去し，二つの $r_\alpha$ に代入する．
この操作は矢 $x$ と関係 $r$ の間の可逆な成分を消すことであり，
対角複体は cubic 六矢二関係の対角解消と一つの可縮部分に分かれる．
従って全幾何 fiber で増大複体は完全である．
有限個の頂点対成分では，全項は基底上有限局所自由な有界複体だから，
perfect-fiber 判定で実際の完全性と任意基底変更後の完全性を得る．
各項は $\Gamma^e$ 上有限局所射影である．

$r$ の双対を $y:3\to1$，$r_1,r_2$ の双対を $d_1,d_2:3\to0$ とする．
tautological relation potential
\[
 W_\delta=yr+d_1r_1+d_2r_2
\]
において，$r_\alpha$ の $xV_{10}$ 部分を $px$ にまとめ，
三段階道の部分を $H$ にまとめると式~\eqref{eq:qa-relative-mutation-delta-potential}
を得る．元の六隣接矢と $d_1,d_2$ は頂点差一，$x,y$ は頂点差二であり，
potential の全巡回項は頂点差四を持つ．

短対角解消を双対化して $[2]$ だけ移すと，$\Theta_\Gamma$ の
次数零の生成束は $R^\vee$，次数 $-1$ は $V^\vee$，次数 $-2$ は
頂点単位の束である．この非正モデルの tensor dg 代数の $H^0$ は，
元の関係 $R$ と，$V^\vee$ の微分である $\partial_VW_\delta$ を課す．
$\partial_{R^\vee}W_\delta=R$ だから
式~\eqref{eq:qa-relative-mutation-ordinary-jacobian} を得る．
この直接計算は relation に一次項 $\delta x$ がある場合にも成立する．
各 component の平坦性は，前節の相対 Hom による $C(\Gamma)$ の構成で
既に得ている．potential の存在だけから Jacobian の平坦性を推論していない．
\end{proof}

\begin{theorem}[壁を横断する実際の相対AS解消]
\label{thm:qa-relative-mutation-nonminimal-as-resolutions}
第一 mutation の全整数代数は，各局所 chart 上で，十生成矢
potential~\eqref{eq:qa-relative-mutation-delta-potential} による
実際の長さ三の右単純解消を持つ．係数束を自明化した重複度は
\[
 \begin{cases}
 0\to P_{i-4}\to P_{i-3}^{\oplus2}\to P_{i-1}^{\oplus2}
                  \to P_i\to S_i\to0,&i\text{ even},\\
 0\to P_{i-4}\to P_{i-3}^{\oplus2}\oplus P_{i-2}
       \to P_{i-1}^{\oplus2}\oplus P_{i-2}
                  \to P_i\to S_i\to0,&i\text{ odd}.
 \end{cases}
 \tag{MS}\label{eq:qa-relative-mutation-simple-complex}
\]
これらの解消，その左向きの双対解消，および
\[
 \mathcal E xt_A^q(S_i,A)=0\ (q\ne3),\qquad
 \mathcal E xt_A^3(S_i,A)\simeq {}_\ell S_{i-4}
\]
は任意の scheme 基底変更に保たれる．
$\delta$ が可逆なら奇数位相の $P_{i-2}$ の対を可縮部分として消せば
cubic の最小解消を得る．$\delta=0$ なら式~\eqref{eq:qa-relative-mutation-simple-complex}
自体が最小 AS10 解消である．
\end{theorem}
\begin{proof}
十矢束を $\mathcal V_{ba}$ と書くと，potential 複体は係数束を保持して
\[
 0\to P_{i-4}\to
 \bigoplus_j\mathcal V_{j,i-4}^\vee\otimes P_j\to
 \bigoplus_j\mathcal V_{ij}\otimes P_j\to P_i\to S_i\to0
\]
となる．微分は矢の余評価，cyclic derivative の分解，矢の評価である．
Jacobian 関係と巡回恒等式により合成は零である．
二次項 $\delta yx$ は，奇数位相の二つの $P_{i-2}$ 間に係数 $\delta$ の
成分を与える．他の成分は正の頂点差の矢を含む．

各幾何 fiber の $\delta=0$ では既存の最小 AS10 の potential 解消となり，
$\delta\ne0$ では二次対を除いた cubic の potential 解消に可縮対を加えた
ものとなる．両者は既存の体上 AS 分類により完全である．
任意の固定した右頂点 component は，前命題の平坦な ordinary Jacobian
の成分からなる有限局所自由層の有界増大複体である．
その全幾何 fiber が完全なので，perfect-fiber 判定により実際の完全性と
任意基底変更後の完全性を得る．

この有限射影解消を $A$ で双対化する．最終項から
${}_\ell S_{i-4}$ への canonical augmentation を付けた各左 component も
有限局所自由層の有界複体である．体上の AS 双対性と同じ fiber 判定から，
表示した Ext の消滅・同定と左単純解消を得る．
この議論に左 Noether 性を入力する必要はない．
\end{proof}

\begin{remark}[最小性と非被約横断方向]
$\delta=\epsilon$，$\epsilon^2=0$ のときにも上の実際の相対解消は存在する．
しかし中間微分に一次項の係数 $\epsilon$ が残り，
全ての微分が正の矢 ideal に入るという相対最小性は満たさない．
余分な最小矢の quotient は $\mathcal O_S/(\epsilon)$ であり，基底平坦ではない．
全幾何 fiber が最小 AS10 であること，最小矢束が平坦であること，
固定された局所自由項からなる相対 AS 解消が存在することは区別される．
係数環が特異な場合の代数の絶対大域次元三も主張していない．
頂点代表を線束で変更する場合の最左係数とAS双対は，
注意~\ref{rem:qa-square-unfolding-vertex-lines} と同じ規約で保持する．
\end{remark}

% !TeX root = ../nc-quadric-avatar.tex
% 2026-11-22 / C v4.8
\section{Spectral avatarのpointed実現と境界延長不能}
\label{sec:qa-spectral-pointed-atkin-lehner}

本節では基礎体$k$を標数零の代数閉体とし，
根を共有しない二つのsquarefree二次形式だけを扱う．
すなわち点曲線とspectral曲線はともに滑らかである．
avatarからのpointed対象とは，
定理~\ref{thm:qa-point-all-noncubic-pointed}と同じ
mutation後の一番頂点を指定したAS商圏を意味する．
この位相の指定を省略しない．
二同種の商公式とAtkin--Lehner対合自体は古典的な楕円曲線論である．
二同種の具体式は\cite[III, Example~4.5]{SilvermanAEC2009}と照合する．
本節で行うのは，それを本稿のpointed認識定理と組み合わせ，
二つのavatar分類が与えるH2対象の差を明示することである．
この適用の新規性について包括的な先行研究調査はまだ行っていない．

\subsection{二捩れ標識を保つ逆構成}

\begin{lemma}[smooth二捩れavatarの点曲線型逆構成]
\label{lem:qa-smooth-point-avatar-inverse}
$T$を滑らかな種数一曲線，$F$を次数二線束，
$0\ne P\in\Pic^0(T)[2]$とし，$G=FP$と置く．
このとき根を共有しないsquarefree二次形式$g_1,g_2$が存在して，
\[
T\simeq E_g=\{c^2=g_1(u,v),\ d^2=g_2(u,v)\}\subset\PP^3
\]
となり，二conic商から得るpencilはそれぞれ$F,G$となる．
根対の$\PGL_2$-同型類は$(T,F,G)$から一意に決まる．
またこの同型類は$(T,P)$だけに依存する．
従って点曲線型の対応
\[
\{\text{根を共有しない被約二根対}\}/\PGL_2
\longleftrightarrow
\{(T,P):T\text{ smooth genus one},\ 0\ne P\in\Pic^0(T)[2]\}/\simeq
\]
も全単射である．左辺の根対は二つの被約二点集合の非順序対を意味し，
各二次形式の非零スカラー倍は忘れている．
\end{lemma}
\begin{proof}
原点を選び，次数二線束$F$に対応する群の点を$a$，
$P$に対応する非零二捩れ点を$p$とする．
$|F|$と$|G|$のdeck involutionは
\[
\iota_F(x)=a-x,\qquad \iota_G(x)=a+p-x
\]
である．従って両者は可換で，積は自由な平行移動$t_p$となる．
生成する群$V=\langle\iota_F,\iota_G\rangle$は四元群である．
$Q_F=T/\langle\iota_F\rangle$と
$Q_G=T/\langle\iota_G\rangle$は射影直線であり，
さらに$Z=T/V$も射影直線となる．
双方の$Q_i\to Z$は次数二で，それぞれ二点で分岐する．

この二つの分岐因子は交わらない．実際，
滑らかな曲線の標数零有限群商の点の慣性群は巡回群である．
もし二つの中間二重被覆が同じ点で分岐すれば，
その点の慣性は両方の商で非自明でなければならず，
$V$内の可能な位数二部分群は$\langle t_p\rangle$だけとなる．
しかし$t_p$は固定点を持たない．
四元群全体は巡回群でないので慣性群にならない．
よって分岐点は四点相異なる．

$Z\simeq\PP^1_{[u:v]}$の座標を選べば二重被覆は
$Q_F=\{c^2=g_1(u,v)\}$，
$Q_G=\{d^2=g_2(u,v)\}$と書ける．
分岐の被約性から$g_1,g_2$はsquarefreeである．
四元群拡大の中間体の合成を取ると
$T\to Q_F\times_ZQ_G$は有限birationalであり，
右辺は分岐因子が交わらないため滑らかである．
したがって同型となる．
各$Q_i$上の次数一線束の引戻しは，
構成上$F,G$である．
この商の構成は根対を座標変換と交換を除いて一意に回収する．

$T$上で次数二線束を別の$F'$へ変えても，
$[2]:\Pic^0(T)\to\Pic^0(T)$が全射なので
ある平行移動$t$で$t^*F\simeq F'$となる．
平行移動は$\Pic^0(T)$に自明に作用するため
$t^*(FP)\simeq F'P$も成り立つ．
従って$(T,F,FP)$の同型類は$(T,P)$だけで決まる．
同じ議論で二pencilを交換できるから，根対の順序も
この同型類に追加の選択を与えない．
前向きの構成は定理~\ref{thm:qa-point-spectral-dual-two-isogeny}であり，
以上はその逆である．
\end{proof}

\begin{theorem}[spectral avatarのAS実現はAtkin--Lehner対合を与える]
\label{thm:qa-spectral-avatar-pointed-involution}
$f=(f_1,f_2)$を上記のsmooth根対とし，
$\mathcal H(f)=(\qgr B_f,\pi B_f)$と置く．
spectral二重被覆$p:C_f\to\PP^1$に対し
$F=p^*\OO(1)$，$P=M_C$，$G=FP$とする．
square avatar $(C_f,F,G,P)$から指定mutationを行い，
一番頂点を指定して得るpointed AS商圏を
$\mathcal H_{\mathrm{sp}}(f)$と書く．
この対象はnoncubic H2に属し，あるsmooth根対$f^\vee$について
\[
\mathcal H_{\mathrm{sp}}(f)\simeq\mathcal H(f^\vee)
\]
となる．幾何同型類上の$f\mapsto f^\vee$は対合である．
点曲線およびspectral曲線を二捩れ標識付きで書けば
\[
(E_{f^\vee},P_{E_{f^\vee}})\simeq(C_f,M_C),\qquad
(C_{f^\vee},M_{C_{f^\vee}})\simeq(E_f,P_{E_f}).
\]
すなわちこの対合は既存の二同種の双対化
$\mathcal W_2$である．

特に
\begin{equation}
\mathcal H_{\mathrm{sp}}(f)\simeq\mathcal H(f)
\quad\Longleftrightarrow\quad
(C_f,M_C)\simeq(E_f,P_{E_f}).
\label{eq:qa-spectral-pointed-selfdual-criterion}
\end{equation}
左辺は構造対象を保つexact同値を意味する．
一般にはこの条件は成立しない．
\end{theorem}
\begin{proof}
前補題を$(T,F,G)=(C_f,F,FP)$へ適用し，
$E_{f^\vee}\simeq C_f$と二pencilの同定を得る．
定理~\ref{thm:qa-point-all-noncubic-pointed}により，
このavatarの指定したpointed AS実現は
$\mathcal H(f^\vee)$とexact同値である．
また同定は二捩れの差も保つので
$(E_{f^\vee},P_{E_{f^\vee}})\simeq(C_f,M_C)$である．
定理~\ref{thm:qa-point-spectral-dual-two-isogeny}の双対性を
$f$と$f^\vee$へ適用すれば第二の同型も得られる．
二度双対を取ると元の二同種へ戻るので，
前補題の一意性から$f^{\vee\vee}\simeq f$である．

論文Aのnoncubic pointed分類は，
$\mathcal H(f)\simeq\mathcal H(g)$の必要十分条件が
根対の$\PGL_2$-同型であることを与える．
前補題はこれを点曲線型の二捩れ標識付き同型へ
言い換えるから，式\eqref{eq:qa-spectral-pointed-selfdual-criterion}を得る．
この主張は幾何同型類と指定したpointed実現の比較である．
相対射とその非忠実性は第~\ref{sec:qa-spectral-relative-nonfaithful-map}節で扱い，
粗対合を覆うstack自己同値は存在しないことを示す．
\end{proof}

\subsection{不変量と自己双対点}

\begin{proposition}[$\chi$座標での対合と二つの自己双対点]
\label{prop:qa-spectral-atkin-lehner-chi}
論文Bの不変式
$\chi=D(f_1,f_2)^2/(\Delta(f_1)\Delta(f_2))$を使う．
smooth根対の同型類は$\chi\in k\setminus\{1\}$で分類され，
前定理の対合は
\begin{equation}
\chi\longmapsto\chi^\vee=\frac{\chi}{\chi-1}
\label{eq:qa-spectral-atkin-lehner-chi}
\end{equation}
である．従ってspectral avatarの指定pointed AS実現が
元の$\mathcal H(f)$と同値となるsmooth幾何同型類は
正確に$\chi=0,2$の二つである．
その曲線の$j$値はそれぞれ$1728,8000$である．
一般の$j$不変式は
\begin{equation}
j(C_f)=\frac{64(\chi+3)^3}{(\chi-1)^2},\qquad
j(E_f)=\frac{64(4\chi-3)^3}{\chi-1}.
\label{eq:qa-two-j-functions-chi}
\end{equation}
\end{proposition}
\begin{proof}
二根を$0,\infty$へ送ると，スカラーを除いて
\[
f_1=uv,\qquad f_2=u^2+a uv+bv^2,
\qquad b\ne0,\quad a^2-4b\ne0
\]
と書ける．ここで$b\ne0$は根を共有しない条件である．
この正規形では
$\chi=a^2/(a^2-4b)$となる．
座標の対角変換と第二形式のスカラー変換で$b=1$にでき，
残る$a\leftrightarrow-a$も座標変換で実現する．
従って$a^2/b$，同じことだが$\chi$が完全不変量となる．
$\chi=1$は$b=0$を意味するのでsmooth域から除外される．

spectral曲線はWeierstrass表示
\[
C_f:\ y^2=x(x^2+ax+b)
\]
を持ち，指定二捩れは$(0,0)$である．
この点を核とする二同種の商は
\[
E':\ Y^2=X\bigl(X^2-2aX+(a^2-4b)\bigr).
\]
実際，
\[
X=x+a+\frac b x,\qquad
Y=y\left(1-\frac b{x^2}\right)
\]
を代入すれば右辺の方程式が成立する．
この射は有限次数二で，$O$と$(0,0)$を同じ無限遠点へ送り，
商の標識二捩れは再び$(0,0)$である．
定理~\ref{thm:qa-point-spectral-dual-two-isogeny}より
$(E',(0,0))\simeq(E_f,P_{E_f})$となる．
したがって$f^\vee$のspectral表示として
\[
f_1^\vee=uv,\qquad
f_2^\vee=u^2-2a uv+(a^2-4b)v^2
\]
を取れる．この二形式は判別式$1,16b$を持ち，
根を共有しない条件は$a^2-4b\ne0$だから，
再びsmooth根対である．不変式を代入して
\[
\chi^\vee=\frac{4a^2}{16b}
=\frac{\chi}{\chi-1}
\]
を得る．この変換の二乗は恒等で，固定点方程式は
$\chi(\chi-2)=0$である．

一般Weierstrass模型$y^2=x(x^2+ax+b)$の
$c_4=16(a^2-3b)$，$\Delta_E=16b^2(a^2-4b)$から
\[
j=\frac{256(a^2-3b)^3}{b^2(a^2-4b)}
\]
を得る．これと二同種先の係数$(-2a,a^2-4b)$を代入すると
式\eqref{eq:qa-two-j-functions-chi}となる．
最後に$\chi=0,2$を代入する．
\end{proof}

\begin{example}[元のpointed quadricと異なるspectral avatar]
\label{ex:qa-spectral-avatar-pointed-inequivalent}
既存の例~\ref{ex:qa-point-spectral-j-distinction}の
$f_1=x(x-1)$，$f_2=(x-2)(x-3)$では
$\Delta(f_1)=\Delta(f_2)=1$，$D(f_1,f_2)=-7$なので
$\chi=49$である．spectral avatarの指定pointed AS実現は
$\chi^\vee=49/48$のnoncubic H2対象となり，
$\chi=49$の元の対象とはpointed exact同値でない．
曲線の不一致も
$j(C_f)=35152/9$，$j(E_f)=28756228/3$
から直接見える．
\end{example}

\begin{example}[曲線が同型でもpointed対象は異なり得る]
\label{ex:qa-spectral-equal-j-not-pointed}
二つの$j$が等しい条件の分子は
\[
(4\chi-3)^3(\chi-1)-(\chi+3)^3
=\chi(\chi-2)(64\chi^2-81\chi+81)
\]
である．従って自己双対点$0,2$のほかにも
\[
\chi_\pm=\frac{81\pm45\sqrt{-7}}{128}
\]
で曲線は同型となる．しかし
$\chi_+^\vee=\chi_-$，$\chi_-^\vee=\chi_+$であり，
これらは相異なるpointed H2同型類である．
両点で$j(C_f)=j(E_f)=-3375$となることは
$64(\chi+3)^3+3375(\chi-1)^2$を
$64\chi^2-81\chi+81$で割った剰余が零であることから分かる．
この例は，式\eqref{eq:qa-spectral-pointed-selfdual-criterion}で
曲線だけでなく非零二捩れ標識も保存する必要を示す．
\end{example}

\subsection{二つの境界型での延長不能}

\begin{corollary}[二つのnoncubic境界型でのH2延長不能]
\label{cor:qa-spectral-involution-no-h2-extension}
論文Bの内在的noncubic境界$\mathfrak N\subset\mathfrak H_2$の
閉性と相対表示同値を入力とする．
上記のsmooth幾何点上のspectral-pointed対応を
全$\mathfrak N$上の射$\mathfrak N\to\mathfrak H_2$に延長することはできない．
より具体的には，次の二族の幾何generic fiberの対応先は，
いずれも$t=0$を越えるH2族に延長できない：
\[
\begin{array}{c|c|c|c}
f_1&f_2&\chi&\chi^\vee\\\hline
uv&u^2+uv+t v^2&(1-4t)^{-1}&(4t)^{-1}\\
uv&t u^2+uv+t v^2&(1-4t^2)^{-1}&(4t^2)^{-1}.
\end{array}
\]
第一族の特殊fiberは一根共有，第二族は根対一致である．
この延長不能は任意の有限分岐trait基底変更の後にも残る．
\end{corollary}
\begin{proof}
第一族は$\operatorname{Spec}k[t,(1-4t)^{-1}]$，
第二族は$\operatorname{Spec}k[t,(1-4t^2)^{-1}]$上で定義される．
各二次形式の判別式は全基底上可逆で，
論文Bの二形式表示は実際の$\mathfrak N$-族を与える．
$t\ne0$では四根相異なる．
表の$\chi$は判別式の極化の直接計算から，
$\chi^\vee$は式\eqref{eq:qa-spectral-atkin-lehner-chi}から得る．

以下では論文Bの代数的stackを全schemeへ通常に延長して扱い，
任意Noetherian基底上の素朴なAbel圏公理の同値を仮定しない．
いずれかの族の対応先に，$t=0$の局所trait上の
H2延長$h$が存在すると仮定する．
その幾何generic fiberはnoncubicだから，忠実平坦降下により
$h$のgeneric fiberは閉部分stack$\mathfrak N$に因数分解する．
ここで閉性は論文Bの
\cite[\texttt{thm:noncubic-boundary-smooth-codim-two}]{QuadricModuli2026}
を用いる．$h$による$\mathfrak N$の引戻しはtraitの閉部分であり，
そのidealはgeneric fiberで零となる．
traitの局所環は整域なのでこのideal自体が零であり，
延長$h$全体が$\mathfrak N$に因数分解する．

一方，論文Bの相対表示同値と
\cite[\texttt{prop:noncubic-good-moduli}]{QuadricModuli2026}により，
$\chi:\mathfrak N\to\mathbb A^1$は全域の正則射である．
従って$\chi\circ h$はtrait上の正則関数でなければならない．
そのgeneric値は表の$(4t)^{-1}$または$(4t^2)^{-1}$であり，
負の付値を持つので矛盾する．
有限分岐拡大後は$t$の付値が正の整数倍となるだけで，
この極は消えない．
全$\mathfrak N$からの延長射があれば
上の有限型二族へ制限できるので，その存在も否定される．
\end{proof}

\begin{remark}[$\Sigma_2$模型と相対範囲]
\label{rem:qa-sigma2-spectral-involution-scope}
式~\eqref{eq:qa-noncubic-sigma2-square-dictionary}によって
spectral $\Sigma_2$データをsquare avatarへ変換し，
本節の指定mutation・pointingでAS商圏を作る限り，
そのpointed同型類も上記の$\chi^\vee$である．
これにより，spectral $\Sigma_2$の集合論的分類対応を
元の$\mathcal H(f)$とのpointed同値と同一視できないことが分かる．
別に指定したrotated ASの構造頂点は未比較である．
smooth域の非被約基底を含む相対射は次節で構成するが，全忠実ではない．
特に$\chi=1$の二型では，直前の系によって
対応先そのものをH2族に延長できない．
これは固定spectral曲線上の通常Picard線束の非延長障害とは
別の，内在的pointed対象に対する障害である．
\end{remark}

% !TeX root = ../nc-quadric-surface-avatar.tex
% 2026-11-22 / C v4.8
\section{Spectral双対化の相対射と二次線束のgerbe}
\label{sec:qa-spectral-relative-nonfaithful-map}

本節の前半では標数零の体上で，二形式の四分岐根が相異なる開を扱う．
後半では二重根を許す平坦代数モデルと，その境界でのSerre障害を扱う．
$\mathfrak N^{\mathrm{sm}}=[P^{\mathrm{sm}}/G]$を非順序の二形式表示stack，
$\mathfrak N^{\mathrm{sm,ord}}=[P^{\mathrm{sm}}/G^\circ]$を
二形式の順序を保った表示stackとする．ここで
$G^\circ=(\GL_2\times\Gm^2)/D$であり，
$\mathfrak N^{\mathrm{sm,ord}}\to\mathfrak N^{\mathrm{sm}}$は
二形式の順序を選ぶ有限\'etale二重被覆である．
これは表示stackとして任意の基底へ引き戻せる構成である．
内在的H2族への解釈は，論文Bの相対表示同値が適用されるsiteで行う．

\subsection{二重被覆による表示とpencilのBrauer類}

\begin{lemma}[二conic被覆のgroupoid]
\label{lem:qa-rootpair-double-cover-groupoid}
$S$上の$\mathfrak N^{\mathrm{sm}}$の対象は，次のデータと同値である：
相対genus zero曲線$Z/S$と，順序を忘れた二つの有限平坦二重被覆
$Q_1,Q_2\to Z$で，各$Q_i/S$も相対genus zeroであり，
分岐因子は相対次数二の\'etale因子で互いに交わらないもの．
非順序の二被覆は二重\'etale被覆上で順序を付けて記述してよい．
順序を保ったデータは$\mathfrak N^{\mathrm{sm,ord}}$に対応する．
\end{lemma}
\begin{proof}
$Q_i\to Z$のtrace分解を
$\OO_Z\oplus L_i^{-1}$と書く．分岐次数二より$L_i$の相対次数は一である．
順序を選ぶ\'etale被覆上で$L_1$は実際の線束なので，
$Z\simeq\PP(q_*L_1)$となる．
$L_2L_1^{-1}$は基底から引き戻した線束である．
fppf局所的に$Z\simeq\PP^1$，$L_i\simeq\OO(1)$とすると，
被覆は$c_i^2=f_i(u,v)$になる．座標と二つのtracefree生成線の変更は
$\widetilde G$の作用そのもので，射影曲線と被覆への作用の核は
共通次数スカラー$D$である．したがって同型のgroupoidは
$[P^{\mathrm{sm}}/G]$と一致する．逆向きはこの局所式からの有効降下である．
\end{proof}

\begin{lemma}[spectral pencilのtwistと反不変線束]
\label{lem:qa-spectral-relative-pencil-brauer}
非順序の入力に対し，点曲線
$E=Q_1\times_ZQ_2$，同時符号$\delta$の商$C=E/\langle\delta\rangle$，
$p:C\to Z$を取る．$E\to C$は有限\'etale二重被覆であり，
\[
 r_*\OO_E=\OO_C\oplus M,\qquad
 \mu:M^{\otimes2}\xrightarrow{\sim}\OO_C
\]
は実際の線束と二次積である．$[M]$は全幾何fiberで非零二捩れである．
一方，$F=p^*\OO_Z(1)$という記法は一般に相対Picard類を表し，
実際の線束を無条件に表さない．その線束化障害は$Z$のBrauer類
$\alpha_Z\in\operatorname{Br}(S)[2]$である．
順序を選んだ被覆では$L_1$が存在するので$\alpha_Z$は消える．
\end{lemma}
\begin{proof}
四根が相異なるので$E/S,C/S$は滑らかなgenus one曲線であり，
$\delta$は固定点を持たない．trace分解と積は任意基底変更に可換である．
各幾何fiberで$E$が連結だから，その非自明な\'etale二重被覆を表す$M$は
非零二捩れである．これは射影曲線の構造層の固有成分から作るため，
共通次数スカラーは自明に作用する．
$\OO_Z(1)$は相対次数一のPicard類で，その境界類は$\alpha_Z$である．
相対Picardの境界写像は$p^*$と可換するので，$F$の障害も$\alpha_Z$となる．
従って$M$が実際の線束であることから$F$まで線束化してはならない．
\end{proof}

\begin{example}[pencilのBrauer障害が非零になる入力]
\label{ex:qa-spectral-relative-pencil-nonsplit}
$k$を標数零の代数閉体，$K=k(s,t)$とし，
\[
 f_1=u^2+v^2,\qquad f_2=u^2+t^2v^2
\]
を取る．四根は相異なる．$\widetilde G$の変換
\[
 h:(u,v,c,d)\longmapsto(-tv,u,d,-tc)
\]
はこの二形式を交換して保ち，$h^2$は全四変数への$-t$倍である．
従って$G$における$h$は位数二である．
$K(\sqrt{s})/K$と$h$によってこの対象をtwistする．
その基底conicは，$\PP^1$上の射影変換
$\left(\begin{smallmatrix}0&-t\\1&0\end{smallmatrix}\right)$による降下であり，
Brauer類は$(s,-t)=(s,t)\ne0$である．
非零性は$t$における剰余が$s\in k(s)^\times/k(s)^{\times2}$となることから従う．
この入力のspectral pencilは実際の線束にならないが，$(M,\mu)$は存在する．
以下の相対構成は$F$の線束化を要求しない．
\end{example}

\subsection{二つのinvolutionの商による相対射}

\begin{theorem}[spectral双対化の非忠実な相対射]
\label{thm:qa-spectral-relative-nonfaithful-map}
任意基底変更と両立する射
\[
 \mathscr D^{\mathrm{ord}}:\mathfrak N^{\mathrm{sm}}
       \longrightarrow\mathfrak N^{\mathrm{sm,ord}},\qquad
 \mathscr D=\operatorname{forget}\circ\mathscr D^{\mathrm{ord}}
       :\mathfrak N^{\mathrm{sm}}\longrightarrow\mathfrak N^{\mathrm{sm}}
\]
が存在し，幾何同型類上で$\chi\mapsto\chi/(\chi-1)$を誘導する．
具体的には，$\iota$を$p:C\to Z$のdeck involution，
$t_{[M]}$を$J(C/S)$の点$[M]$による$C$上の平行移動と置き，
\[
 \iota'=t_{[M]}\iota,\quad
 Z'=C/\langle\iota,t_{[M]}\rangle,\quad
 Q'_1=C/\langle\iota\rangle=Z,\quad
 Q'_2=C/\langle\iota'\rangle
\]
から得る\emph{順序付き}二被覆を$\mathscr D^{\mathrm{ord}}$とする．
特にこの構成はsmooth域上の相対射であり，stack自己同値ではない．
\end{theorem}
\begin{proof}
$C$は$J(C/S)$のtorsorであり，$[M]\in J(C/S)[2]$から
$t_{[M]}$が自然に定まる．$\iota$はJacobian上で$-1$を誘導するので
$t_{[M]}$と可換し，$\iota'$もinvolutionである．
$[M]$は全fiberで非零だから$t_{[M]}$は固定点を持たない．
fppf局所的に$C$に原点を選ぶと，$\iota$と$\iota'$は
$x\mapsto a-x$，$x\mapsto a+[M]-x$となる．
各固定点schemeは$[2]$のfiberで\'etale次数四であり，両者は交わらない．

有限群の位数は可逆なので商とその構造層の不変部分は任意の基底変更と可換する．
上の固定点の記述から$Q'_1,Q'_2,Z'$は相対genus zero曲線で，
$Q'_i\to Z'$の分岐は各相対次数二，両分岐は交わらない．
$C\to Q'_1\times_{Z'}Q'_2$は全幾何fiberで同型であり，
双方の$Z'$上の直像代数は階数四の局所自由加群だから同型である．
前補題のgroupoid記述によって，これが実際の表示stackの対象を与える．

入力の同型は$(C,p,M,\mu)$を自然に移し，$[M]$，二involution，
その商をすべて保つ．従って対象だけでなく射とその合成も保存される．
入力の二形式交換にも自然であり，共通次数スカラーはすべての商に自明に作用する．
この構成は非被約基底を含む任意基底変更に可換である．

二形式族を直接書くこともできる．$q:Z'\to S$とし，
$Q'_i$のtrace分解を$\OO_{Z'}\oplus (L'_i)^{-1}$と書く．
$L'_1$の相対次数は一なので$W'=q_*L'_1$は階数二で，
$Z'\simeq\PP(W')$となる．評価同型により
$L'_2\simeq L'_1\otimes q^*T$となる基底線束$T$が定まる．
二被覆の積は
\[
 q_1\in\operatorname{Sym}^2W',\qquad
 q_2\in\operatorname{Sym}^2W'\otimes T^{\otimes2}
\]
を与える．したがって一次生成束
$W'\oplus\OO_Sc'\oplus T^{-1}d'$に対し，$W'$を中心とし，
$c'^2=q_1$，$d'^2=q_2$を課した二形式環が実際に大域構成できる．
第二の式は$T^{-2}\to\operatorname{Sym}^2W'$という関係式を意味する．
基底線束$T$を捨てずに残すことがcocycleの適合を保証する．

正規形$f_1=uv$，$f_2=u^2+a uv+bv^2$では，
$C:y^2=x(x^2+ax+b)$上で
\[
 t_{[M]}(x,y)=(b/x,-by/x^2),\qquad
 X=x+b/x,\quad c'=x-b/x,\quad d'=y/x
\]
を取れる．その関係は
\[
 c'^2=X^2-4b,\qquad d'^2=X+a.
\]
従って出力二形式は
$g_1=U^2-4bV^2$，$g_2=UV+aV^2$である．
$U'=U+aV$と二形式交換によって
\[
 (g_2,g_1)=(U'V,\ U'^2-2aU'V+(a^2-4b)V^2)
\]
となる．これは既存の双対二同種模型なので，
$\chi\mapsto\chi/(\chi-1)$を与える．
非忠実性とfiberは次の命題で示す．
\end{proof}

\subsection{逆構成の二次線束gerbe}

\begin{proposition}[順序付き標的上の$\mu_2$-gerbe]
\label{prop:qa-spectral-relative-quadratic-gerbe}
$\mathscr D^{\mathrm{ord}}$は$\mu_2$でband付けされたgerbeである．
具体的に，順序付き標的$Q'_1,Q'_2\to Z'$の点曲線
$C=Q'_1\times_{Z'}Q'_2$上で，二つの次数二pencilの相対Picard類を
$F_1,F_2$，その差を$P=F_2F_1^{-1}\ne0$とする．
$\mathscr D^{\mathrm{ord}}$のfiberは，$P$を代表する実際の線束$M$と
$\mu:M^{\otimes2}\simeq\OO_C$のgroupoidである．
射はこの二次積を保つ線束同型とする．
このgerbeのBrauer像は
\[
 \beta=\alpha(Q'_2)-\alpha(Q'_1)\in\operatorname{Br}(S)[2]
\]
である．一般の基底では$\operatorname{Pic}(S)/2$の寄与もあり得るので，
$H^2(S,\mu_2)$のgerbe類そのものをBrauer像だけで同定しない．
\end{proposition}
\begin{proof}
Picard stackからPicard functorへの射がgerbeであることは
\cite[Lemma~108.9.1, Tag~0DME]{StacksRelativePicard}の一般論である．
以下ではその二次積付きのfiberを直接記述する．
線束と二乗自明化の$\mu_2$-torsorとの対応は
Kummer列と\cite[Tag~040Q]{StacksKummer}による．
相対Picard類$P$はfppf局所的に線束$M$で表される．
$P^2=1$なので$M^2$は基底から引き戻した線束で，
基底をさらに局所化すれば二次積$\mu$を選べる．
二つの選択は基底上の線束とその二乗自明化だけ異なり，
局所的に同型である．自己同型は$\mu_2$である．
ここでは$\pi_*\OO_C=\OO_S$を使っている．
従ってこの選択のgroupoidは$\mu_2$-gerbeをなす．
相対Picardの境界類の加法性から，$P$の線束化障害は
$\alpha(Q'_2)-\alpha(Q'_1)$である．
二次積を忘れるとこれがgerbe類のBrauer像となる．

このgerbeが実際に$\mathscr D^{\mathrm{ord}}$のfiberであることを確認する．
二次線束$(M,\mu)$から
$E=\operatorname{Spec}_C(\OO_C\oplus M)$を構成する．
これは幾何fiberで連結な\'etale二重被覆である．
$\iota_1$を$C\to Q'_1$のdeck involutionとする．
$\iota_1$はJacobian上$-1$として作用するため$[M]$を保つ．
従って二次線束$\iota_1^*(M,\mu)\otimes(M,\mu)^{-1}$の
相対Picard類は零であり，それは基底上の二次線束の引戻しである．
対応する基底上の$\mu_2$-torsorをfppf局所的に自明化すれば，
$\iota_1^*(M,\mu)\simeq(M,\mu)$を得る．
これが$\iota_1$の$E$へのliftを与え，その全体は$\mu_2$-torsorをなす．
二つのliftを$\widetilde\iota_1,\delta\widetilde\iota_1$と書く．
各liftの二乗は恒等である：$\iota_1$の固定点schemeは
\'etale全射で，その上の二点被覆では任意の置換の二乗が恒等である．
二乗は$E/C$のdeck変換なので，そこで恒等なら全体で恒等である．

二つのliftは可換し，積は$\delta$である．
liftのJacobian作用は$-1$である：norm isogeny
$r_*:J(E/S)\to J(C/S)$との可換性から，$-1$との差は
有限\'etale核に入るが，連結Jacobianからこの有限群への群準同型は零である．
従って各liftは幾何fiberで四つの固定点を持つ．
各商$E/\langle\widetilde\iota_1\rangle$，
$E/\langle\delta\widetilde\iota_1\rangle$は相対conicであり，
共通基底$Q'_1=C/\langle\iota_1\rangle$への二重被覆となる．
両分岐は交わらず，各分岐次数は二である．
liftを交換すると二被覆が交換されるだけだから，
その\emph{非順序}対は追加の選択なしに$S$上へ降下する．
これが$\mathfrak N^{\mathrm{sm}}$の対象である．

そのspectral商は構成した$C$，反不変二次線束は$(M,\mu)$，
spectral二重射は$C\to Q'_1$である．
$F_2F_1^{-1}=P$と二pencilのdeck involutionの関係から，
$\iota_2=t_P\iota_1$である．よって前定理の構成は
元の順序付き$Q'_1,Q'_2\to Z'$を回収する．
逆に入力からこの手続を行うと，元の$E\to C$と
二つのliftの非順序対を回収する．
線束同型と被覆同型の対応も自然なので，これはfiber groupoidの同値である．
\end{proof}

\begin{corollary}[慣性の核と幾何homotopy fiber]
\label{cor:qa-spectral-relative-two-mu2-gerbes}
$\mathscr D$の相対慣性の核は同時符号$\delta$による$\mu_2$である．
任意の幾何点におけるhomotopy fiberは
\[
 B\mu_2\ \amalg\ B\mu_2
\]
と同値である．一般点では慣性準同型の像は二conicを個別に保つ
位数八の部分群であり，全慣性の位数十六に対して指数二である．
したがって，幾何同型類上では対合でも，
$\mathscr D$は全忠実でもstack自己同値でもない．
\end{corollary}
\begin{proof}
$\mathscr D^{\mathrm{ord}}$は前命題の$\mu_2$-gerbeで，
$\mu_2$は$E\to C$のdeck変換そのものである．
順序忘却は表現可能な有限\'etale二重被覆だから，
幾何点への2-fiberは二点，その各点上のgerbeは$B\mu_2$となる．
一般点の慣性については，gerbeの慣性全射と順序忘却の慣性単射を合成する．
\end{proof}

\begin{remark}[一般点の逆構成に必要な拡大次数]
\label{rem:qa-spectral-relative-gerbe-index-four}
順序付き標的の独立二形式の普遍係数体では，既存の
定理~\ref{thm:qa-generic-point-period-two-index-four}により
$\beta=\alpha(Q'_2)-\alpha(Q'_1)$のindexは四である．
体上では$H^2(K,\mu_2)\to\operatorname{Br}(K)[2]$が同型なので，
この場合の二次線束gerbeは非neutralであり，
その対象を得るための有限拡大次数の最小値は四である．
実際、indexが拡大次数を割り、分岐根選択の四重被覆がこの下界を達成する．
従って幾何fiberが$B\mu_2$であることと，
元の体上に逆構成の対象が存在することは別である．
また$\mathscr D$は幾何同型類上の対合を持ち上げるが，
$\mathscr D^2\simeq\operatorname{Id}$という相対自然同型は主張しない．
\end{remark}

\subsection{Smooth spectral対合に対する慣性の中心bandの障害}
\label{subsec:qa-spectral-center-band-obstruction}

この小節では$k$を標数零の代数閉体とする。
以下で$\mathfrak N^{\mathrm{sm}}$は、二形式が根を共有しない開部分である。
その粗モジュライ座標は
$U=\mathbb A^1_\chi\setminus\{1\}$である。
幾何同型類の対合
$\varphi(\chi)=\chi/(\chi-1)$は$U$の自己同型であるが、
これを全pointed慣性を保つstack自己同値へ持ち上げることはできない。
各幾何点の抽象群だけの比較では、この障害は見えない。

\begin{proposition}[自己双対点と等$j$点の幾何慣性]
\label{prop:qa-spectral-geometric-inertia-comparison}
$D_8$を位数八の二面体群とする。
$\mathfrak N^{\mathrm{sm}}$の幾何慣性とspectral対合の対応は次の通りである。
\[
\begin{array}{c|c|c}
\chi&\varphi(\chi)&\text{pointed慣性}\\\hline
0&0&H_0\ (|H_0|=32)\\
2&2&D_8\times C_2\\
\chi_+=(81+45\sqrt{-7})/128&\chi_-&D_8\times C_2\\
\chi_-=(81-45\sqrt{-7})/128&\chi_+&D_8\times C_2\\
\chi\ne0,1&\varphi(\chi)\ne0,1&D_8\times C_2
\end{array}
\]
ここで$H_0$は論文Bの非分裂完全列
\[
1\longrightarrow\mu_2^2\longrightarrow H_0
\longrightarrow D_8\longrightarrow1
\]
で記述される群である。
特に二つの自己双対点の慣性は互いに異なるが、
対合は各点を固定するので、この事実自体は対合を排除しない。
等$j=-3375$の二点でも慣性の抽象群は同じである。
\end{proposition}
\begin{proof}
論文Bの\cite[\texttt{prop:noncubic-all-stabilizers}]{QuadricModuli2026}
はsmooth域で追加の根配置対称性が生じる条件を
$\chi=0$と同定する。
従って$\chi=2,\chi_\pm$には追加の慣性がない。
残りは式~\eqref{eq:qa-spectral-atkin-lehner-chi}と
例~\ref{ex:qa-spectral-equal-j-not-pointed}から従う。
\end{proof}

\begin{lemma}[一般慣性の中心とその二次分裂体]
\label{lem:qa-noncubic-generic-inertia-center}
$K=k(\chi)$、$b=(\chi-1)/(4\chi)$とし、$K$上の表示
\[
 f=(uv,\ u^2+uv+bv^2)
\]
の$G$内の安定化群を$H_f$とする。
その幾何群は$D_8\times C_2$であり、中心は
\begin{equation}
 Z(H_f)\simeq
 \operatorname{Spec}K\ \amalg\ \operatorname{Spec}K\
 \amalg\ \operatorname{Spec}K(\sqrt b)
 \label{eq:qa-noncubic-center-splitting}
\end{equation}
という基礎有限\'{e}taleスキームを持つ。
最初の二点は恒等元と$c,d$の同時符号反転である。
これは単なる位数の計算ではなく、中心へのGalois作用の計算である。
\end{lemma}
\begin{proof}
$r^2=b$を取り、$e_1,e_2$を$c,d$の個別符号反転とする。
列$(u,v,c,d)^{\mathsf t}$への置換行列を
\[
 R=\operatorname{diag}\left(
 \begin{pmatrix}0&b\\1&0\end{pmatrix},r,r\right)
\]
と置く。$R(uv)=buv$、
$R(u^2+uv+bv^2)=b(u^2+uv+bv^2)$なので、
$R$は$B_f$の次数付き自己同型である。
また$R^2=bI_4$だから、次数スカラーを除いた$G$内では位数二である。

幾何体上で$\rho_1,\rho_2$を$x^2+x+b$の二根、
$\theta^2=\rho_2$とし、もう一つの生成元を
\[
 T=\operatorname{diag}\left(
 \begin{pmatrix}\rho_2&-b\\1&-\rho_2\end{pmatrix},
 \begin{pmatrix}0&\theta\\
 \theta(\rho_2-\rho_1)&0\end{pmatrix}\right)
\]
と取れる。これは二形式を交換する。
直接計算すると、射影群で
\[
 R^2=T^2=e_1^2=e_2^2=1,\quad
 [R,e_i]=[e_1,e_2]=1,\quad
 Te_1T^{-1}=e_2,\quad [R,T]=e_1e_2.
\]
従って$z=Re_1$は中央対合である。
これらは一般根配置の全$V_4$を持ち上げ、核の$\mu_2^2$も含むため、
全安定化群を生成する。
その中心は
\[
 Z(H_f)(\overline K)=\{1,e_1e_2,z,z(e_1e_2)\}
 =\{1,e_1e_2,Re_1,Re_2\}.
\]

最初の二元は$K$上で定義される。
$r\mapsto-r$は$R\mapsto(e_1e_2)R$を与えるので、
残る二元を交換する。
$b$は$\chi=1$で奇数付値を持ち非平方だから、
その二元は正確に$K(\sqrt b)$で分裂する。
これで表示した有限\'{e}taleスキームとGalois作用を得る。
\end{proof}

\begin{theorem}[smooth spectral対合は全stack自己同値へ持ち上がらない]
\label{thm:qa-spectral-no-stack-autoequivalence}
幾何pointed同型類に
$\chi\mapsto\chi/(\chi-1)$を誘導する
$k$-stack自己同値
\[
 \mathfrak N^{\mathrm{sm}}\longrightarrow\mathfrak N^{\mathrm{sm}}
\]
は存在しない。
より強く、この粗モジュライ自己同型を覆う射は、
一般点でrepresentableになることができない。
したがって定理~\ref{thm:qa-spectral-relative-nonfaithful-map}の相対射に
現れる非自明な慣性核は、別の構成を選んでも一般点で除去できない。
\end{theorem}
\begin{proof}
まず一般coarse fiberを同じ体$K=k(\chi)$上で比較する。
前補題の表示は$K$-有理点を持つ。
$P\to\mathbb A^1_\chi$の一般fiberはsmoothである。
実際、上の正規形では$\chi=(1-4b)^{-1}$の微分が
$4(1-4b)^{-2}db\ne0$であり、群作用により全一般fiberで同じことが成立する。
また、
各幾何fiberの根対は$\chi$によって一意に決まる。
従って$G_K$はこのfiberに幾何的に推移的に作用する。
安定化群は有限\'{e}taleで、軌道写像はsmooth全射だから、
このfiberはfppf斉次空間$G_K/H_f$に同型である。
したがって$\mathfrak N^{\mathrm{sm}}$の一般coarse fiberは
neutral gerbe $BH_f$である。

対合で引き戻した標的の一般fiberにも同じ表示を使えるが、
第二形式の係数は
\[
 b'=\frac{\varphi(\chi)-1}{4\varphi(\chi)}
 =\frac1{4\chi}
\]
となる。その表示を$f'$と書くと、標的gerbeは$BH_{f'}$である。
前補題により、その中心の非有理二点の分裂体は$K(\sqrt{b'})$である。
しかし
\[
 \frac b{b'}=\chi-1\notin K^{\times2}.
\]
$b,b'$はともに非平方なので、二つの二次拡大
$K(\sqrt b)$と$K(\sqrt{b'})$は$K$上非同型である。
従って$Z(H_f)$と$Z(H_{f'})$は有限\'{e}tale群スキームとして非同型である。
ここでは標的を$\varphi$で引き戻してから$K$上で比較しており、
粗モジュライの体自己同型を二度適用していない。

gerbe同値がneutral objectをneutral objectへ送る必要はないが、
この自由度は中心の差を解消しない。
実際、$BH_f\simeq BH_{f'}$の下で自明$H_f$-torsorの像が
$H_{f'}$-torsor $Q$ならば、
\[
 H_f\simeq\operatorname{Aut}(Q)={}^{Q}H_{f'}
\]
となる。右辺は共役によるinner formである。
共役は中心に自明に作用するため
$Z({}^{Q}H_{f'})\simeq Z(H_{f'})$であり、先の非同型に矛盾する。
これでstack自己同値は排除される。

一般点でrepresentableな射では、\cite[Tag~04YY]{StacksModuli}により、
自明対象の慣性から
$\operatorname{Aut}(Q)$への写像がmonomorphismになる。
両方の幾何位数は十六であるから、有限\'{e}tale群のこの射は同型となる。
同じ中心の矛盾が成立する。
\end{proof}

\begin{remark}[幾何慣性とbandを区別する]
この障害は$\chi=0$での追加慣性や、境界$\chi=1$での極とは独立である。
すでに一般coarse fiberの上で現れる。
中心の二次分裂体は、点曲線$E_f$の非有理二捩れ二点の分裂体とも一致する。
二同種先の二捩れ群を元のものと同一視できないことが、
ここではpointed自己同値群の中心という内在的不変量に現れている。
幾何同型類の対合、慣性を非忠実に送る相対射、
全慣性を保つstack自己同値は別の主張である。
\end{remark}

\subsection{二重根を許す平坦表示と境界での無限高次 Ext}
\label{subsec:qa-spectral-singular-boundary}

smooth spectral同型類には、二重根を許す中央二重平方代数としての
明示的な平坦延長がある。しかし、その特殊fiberのtails圏は
$\mathcal H_2$には属さない。このことを、固定した二対象間の
全高次Extの計算によって示す。

\begin{proposition}[両境界traitのgeneric像に対する平坦な代数モデル]
\label{prop:qa-spectral-flat-singular-boundary}
$r\in\{1,2\}$とし、$A_r=k[t,(1-4t^r)^{-1}]$と置く。
入力をそれぞれ
\[
 \begin{array}{c|c|c}
 r&f_1&f_2\\\hline
 1&uv&u^2+uv+tv^2\\
 2&uv&tu^2+uv+tv^2
 \end{array}
\]
とし、両方に対して
\[
 g_1=U^2-4t^rV^2,\qquad g_2=UV+V^2,
 \qquad
 \mathcal B^{(r)}=A_r[U,V]\langle c,d\rangle/
       (c^2-g_1,d^2-g_2)
\]
と置く。$U,V$は中心にあり、四生成元は全て次数一とする。
次が成立する。
\begin{enumerate}
\item 各$\mathcal B^{(r)}_n$は階数$(n+1)^2$の自由$A_r$-加群である。
  $t\ne0$の幾何fiberでは二つの$g_i$がsquarefreeかつ根を共有せず、
  そのpointed同型類は入力のsmooth spectral同型類である。
  入力の$t=0$は$r=1$では一根共有、$r=2$では根対一致であり、
  二つの出力代数は共通の特殊fiber
  \[
   B_0=k[U,V]\langle c,d\rangle/
       (c^2-U^2,d^2-UV-V^2)
  \]
  を持つ。
\item
  $X=B_0/(c-U)B_0$とする。さらに$S=k[V]$に
  \[
   s\cdot U=s\cdot c=0,\qquad
   s\cdot V=Vs,\qquad s\cdot d=Vs
   \quad(s\in S)
  \]
  で次数付き右$B_0$-加群構造を入れる。
  このとき全ての$n\ge0$について
  \begin{equation}
    \operatorname{Ext}^n_{\qgr B_0}(\pi X,\pi S)\simeq k.
    \label{eq:qa-spectral-boundary-infinite-ext}
  \end{equation}
\item $D^b(\qgr B_0)$は右Serre関手を持たない。
  特に$(\qgr B_0,\pi B_0)$は$\mathcal H_2$の対象でなく、
  $\qgr B_0$の大域次元は無限である。
\end{enumerate}
第一項のgeneric幾何同型類の主張は、全慣性を保つ相対自己同値の
存在を意味しない。
\end{proposition}

\begin{proof}
\par\medskip\noindent\textbf{平坦性とgeneric同型類.}\par\nobreak\smallskip
二形式の判別式に仮定を置かない中央Clifford標準形
\cite[\texttt{lem:noncubic-relative-clifford-center}]{QuadricModuli2026}
により、$h=cd+dc$、$C_r=A_r[U,V,h]$とすれば
\[
 \mathcal B^{(r)}=C_r\oplus C_rc\oplus C_rd\oplus C_rcd,
 \qquad \deg h=2.
\]
従って各次数は自由であり、Hilbert級数は
\[
 \frac{1+2z+z^2}{(1-z)^2(1-z^2)}
 =\frac{1+z}{(1-z)^3}
 =\sum_{n\ge0}(n+1)^2z^n
\]
となる。これは二重根を持つ$t=0$でも変わらない。

入力の不変量は$\chi(f)=1/(1-4t^r)$である。
出力の判別式は$\operatorname{disc}(g_1)=16t^r$、
$\operatorname{disc}(g_2)=1$、混合判別式は$-2$なので、
\[
 \chi(g)=\frac1{4t^r}
 =\frac{\chi(f)}{\chi(f)-1}.
\]
$g_2$の二根$V=0,U+V=0$での$g_1$の値は、それぞれ
$U^2,(1-4t^r)V^2$である。従って$t(1-4t^r)\ne0$では
両形式はsquarefreeかつ根を共有しない。
smooth noncubicの幾何分類とspectral対合公式から第一項を得る。
generic fiberにおける表示を直接見るには、一般に
$f_2=pu^2+auv+bv^2$に対するspectral曲線
\[
 y^2=x(px^2+ax+b)
\]
で
\[
 \xi=px+\frac b x,\qquad
 c'=px-\frac b x,\qquad d'=\frac y x
\]
と置けばよい。すると
\[
 (c')^2=\xi^2-4pb,\qquad (d')^2=\xi+a
\]
となる。上の二族では$(p,a,b)=(1,1,t),(t,1,t)$なので、
$pb=t^r$であり、斉次化すると表示した$g_1,g_2$になる。
第二族では$p=t$が特殊点で可逆でないため、
ここではこの座標構成をgeneric fiber上でのみ用いている。
特殊fiberまでのcanonicalな相対spectral構成は主張しない。

\par\medskip\noindent\textbf{二周期自由解消の完全性.}\par\nobreak\smallskip
ここから$t=0$とする。
$R_c=k[U,V,c]/(c^2-U^2)$と置く。
同じClifford標準形により
\[
 B_0=R_c[h]\oplus R_c[h]d
\]
は左$R_c$-加群として自由である。
$R_c$では
\[
 \operatorname{Ann}_{R_c}(c-U)=(c+U)R_c,
 \qquad
 \operatorname{Ann}_{R_c}(c+U)=(c-U)R_c.
\]
例えば$a,b\in k[U,V]$に対し
$(c-U)(a+bc)=(a-Ub)c-U(a-Ub)$であるから、
第一の等式は基底$1,c$の係数比較だけで従う。
従って$R_c/(c-U)$の交互乗法による二周期自由解消を
$-\otimes_{R_c}B_0$で移すと、完全列
\begin{equation}
 \cdots\xrightarrow{\,c-U\,} B_0(-2)
 \xrightarrow{\,c+U\,}B_0(-1)
 \xrightarrow{\,c-U\,}B_0
 \longrightarrow X\longrightarrow0
 \label{eq:qa-spectral-boundary-periodic-resolution}
\end{equation}
を得る。各矢印は\emph{左乗法}による右加群準同型であり、
一般に第$i$微分は$i$が奇数なら$c-U$、偶数なら$c+U$である。
表示式の省略部分もこの規則に従う。

\par\medskip\noindent\textbf{tailsで自由項から点対象への高次Extが消えること.}\par\nobreak\smallskip
$C=k[U,V,h]$、
\[
 \mathcal Y=\mathbb P_k(1,1,2)
 =[(\operatorname{Spec}C\setminus V(C_+))/\mathbb G_m]
\]
を重み付き射影\emph{スタック}とし、
$\mathscr B$を$B_0$に対応する代数層とする。
中央Clifford標準形の基底次数は$0,1,1,2$なので、
$B_+$-torsionと$C_+$-torsionは一致する。
従って中心局所化と$\mathbb G_m$同変descentによって
\[
 \qgr B_0\simeq\operatorname{coh}(\mathcal Y,\mathscr B)
\]
を得る。この同値には二形式のsquarefree性を要さない。
これは\cite[\texttt{prop:nc-qgr-weighted}の証明]{QuadricModuli2026}
のsheafification部分と同じ議論である。

$S$上では$h=cd+dc$が零であり、$C$-加群として
$S=C/(U,h)$である。したがって$\mathscr S=\widetilde S$は
通常点$p=[0:1:0]\in\mathcal Y$に支持された
一次元のskyscraper加群である。
この点は次数一座標$V$が非零のaffine chartにあり、慣性を持たない。
各$\mathscr B(-i)$は局所的に自由な右$\mathscr B$-加群だから、
局所Extと大域cohomologyのスペクトル系列により
\[
 \operatorname{Ext}^q_{\qgr B_0}(\pi B_0(-i),\pi S)
 \simeq H^q(\mathcal Y,\mathscr S(i))
 =\begin{cases}k,&q=0,\\0,&q>0\end{cases}
 \quad(i\in\mathbb Z).
\]
ここでは点支持のcoherent層の高次cohomology消滅だけを用い、
$\mathscr B$の滑らかさや有限大域次元は仮定していない。

完全関手$\pi$で二周期解消を移し、
$\operatorname{Hom}_{\qgr B_0}(-,\pi S)$を適用する。
全自由項はこのHomに対して高次acyclicであるから、
得られたHom複体が$\operatorname{Ext}^*(\pi X,\pi S)$を計算する。
実際、任意の有限次数について解消を十分長く切り、
通常のdimension shiftingを用いればよい。
各項は$k$であり、$c-U,c+U$はいずれも$S$に零として作用するため、
全微分が零になる。従って
\eqref{eq:qa-spectral-boundary-infinite-ext}を得る。

\par\medskip\noindent\textbf{Serre関手の不存在.}\par\nobreak\smallskip
もし$D^b(\qgr B_0)$が右Serre関手$\mathbb S$を持てば、
\[
 \operatorname{Ext}^n(\pi X,\pi S)^*
 \simeq
 \operatorname{Hom}_{D^b(\qgr B_0)}
       (\pi S,\mathbb S(\pi X)[-n])
\]
となる。$\mathbb S(\pi X)$は有界複体なので、$n\gg0$では
右辺が標準$t$構造の直交性によって零となる。
これは左辺が全$n\ge0$で$k$であることに反する。
従って右Serre関手は存在せず、$\mathcal H_2$の最初の条件に反する。
また全高次Extの非消滅から、大域次元は無限である。
\end{proof}

\begin{corollary}[この表示族におけるsquarefree条件の必要十分性]
\label{cor:qa-double-square-degenerate-not-h2}
$k$を標数零の代数閉体、$f_1,f_2\in k[U,V]_2$を任意の二元二次形式とする。
零形式も許し、
\[
 B(f_1,f_2)=k[U,V]\langle c,d\rangle/(c^2-f_1,d^2-f_2)
\]
と置く。いずれかの$f_i$がsquarefreeでなければ、
有限生成次数付き右加群$X$と点加群$S_\delta$が存在し、
\[
 \operatorname{Ext}^n_{\qgr B(f_1,f_2)}(\pi X,\pi S_\delta)
 \simeq k\qquad(n\ge0)
\]
となる。従って$D^b(\qgr B(f_1,f_2))$は右Serre関手を持たない。

特に\emph{この中央二重平方表示族の中では}
\[
 (\qgr B(f_1,f_2),\pi B(f_1,f_2))\in\mathcal H_2
 \quad\Longleftrightarrow\quad
 \operatorname{disc}(f_1)\operatorname{disc}(f_2)\ne0.
\]
二つのsquarefree形式が根を共有したり、同じ根対を持ったりすることは許される。
\end{corollary}
\begin{proof}
必要なら二形式を交換する。
$f_1$が非零かつ非squarefreeならば、$\operatorname{GL}_2$による座標変更と
$c$の定数倍により$f_1=U^2$とできる。
$\delta\in k$を$\delta^2=f_2(0,1)$となるように選び、
\[
 S_\delta=k[V],\qquad
 s\cdot U=s\cdot c=0,\quad s\cdot V=Vs,\quad
 s\cdot d=\delta Vs
\]
とする。$f_2(0,V)=\delta^2V^2$なので、これは右$B(f_1,f_2)$-加群である。
$X=B(f_1,f_2)/(c-U)B(f_1,f_2)$には前命題と同じ交互$c\pm U$の
二周期解消がある。
この完全性には第二形式を用いていない。
中心$k[U,V,h]$上で$S_\delta=k[U,V,h]/(U,h)$であることも変わらないため、
同じskyscraperのExt消滅と零微分の計算から結論を得る。

$f_1=0$の場合には$R_c=k[U,V,c]/(c^2)$とし、
$X=B(f_1,f_2)/cB(f_1,f_2)$とする。
$\operatorname{Ann}_{R_c}(c)=cR_c$と左$R_c$-自由性により
\[
 \cdots\xrightarrow{\,c\,}B(f_1,f_2)(-2)
 \xrightarrow{\,c\,}B(f_1,f_2)(-1)
 \xrightarrow{\,c\,}B(f_1,f_2)
 \longrightarrow X\longrightarrow0
\]
は完全である。同じ$S_\delta$を用いれば、今度は全微分$c$が零として作用する。
従ってこの場合にも全次数のExtは$k$となる。
右Serre関手の不存在は前命題の最後の段落による。

判別式がともに非零の場合の十分性は、既存のnoncubic逆構成
\cite[\texttt{thm:noncubic-relative-effectivity}]{QuadricModuli2026}
を体上に適用して得られる。
そこで用いるZhang twistとVeronese比較は構造対象を保つので、
ここでの$(\qgr B(f_1,f_2),\pi B(f_1,f_2))$自身が$\mathcal H_2$に属する。
この十分性は先行する論文Bの結果であり、上の非squarefreeの場合の
Ext計算とは区別する。
\end{proof}

\begin{remark}[この境界計算の範囲]
$d$が$-V$として作用する点加群$S_-=k[V]$でも同じ計算が成立する。
一方、本命題は特定の平坦な代数延長候補についての計算であり、
全ての可能な$\mathcal H_2$への延長をそれだけで排除するものではない。
全ての延長を排除する主張には、noncubicの閉性と$\chi$の極を用いる
先のtraitの議論を適用する。
ここで得た無限高次Extは、自然な中央二重平方表示を境界まで保った際に
失われる圏論的条件を具体化している。
右Serre関手を排除した圏は有界coherent導来圏$D^b(\qgr B_0)$である。
perfect部分圏におけるClifford双加群のSerre公式とは矛盾しない。
実際、上で構成した$\pi X$はperfectではない。
\end{remark}

% !TeX root = ../nc-quadric-avatar.tex
% 2026-11-22 / smooth ruling interchange through a four-sheeted half-period torsor
\section{滑らかなnoncubicのruling交換と半周期torsor}
\label{sec:qa-half-period-ruling}

前節は一般係数体でruling交換が頂点固定のMorita
自己同値にならないことを示した。
一方、幾何的に四項を交換するための
\emph{有限\'{e}tale被覆}は、Jacobianの半周期から構成できる。
ここで点曲線が滑らかな開に限定し、
nodal/$I_2$への同じ次数四の延長は主張しない。

\begin{theorem}[半周期による二pencil交換の次数四torsor]
\label{thm:qa-half-period-ruling-torsor}
$U^\circ\subset U_2$を、二つの二次形式の四根が相異なる開とし、
$E\to U^\circ$を滑らかな種数一の点曲線、
$J=\Pic^0_{E/U^\circ}$をその相対Jacobianとする。
前節のBrauer gerbe上で定義した
$P=[GF^{-1}]$は、根の選択や線束の
基底捩れによらない非零な
$P\in J(U^\circ)[2]$へ降下する。
$[2]:J\to J$のこのsection上のfiber
\begin{equation}
 T_P=U^\circ\times_{P,J,[2]}J
 \xrightarrow{\quad}U^\circ
 \label{eq:qa-half-period-four-cover}
\end{equation}
は有限\'{e}tale次数四で、
$J[2]$のtorsorである。

$T_P$上のtautological sectionを$t$とする。
種数一torsor $E_{T_P}$に対する
Jacobianの平行移動$\tau_t$は、
相対Picard類について
\begin{equation}
 [\tau_t^*F]=[G],\quad
 [\tau_t^*G]=[F],\quad
 [\tau_t^*H]=[H],\quad
 [\tau_t^*(HG)]=[HF]
 \label{eq:qa-half-period-exchange}
\end{equation}
を満たす。
従って相対四項
$(\OO_E,F,H,HG)$と$(\OO_E,G,H,HF)$は、
$T_P$をBrauer gerbeへ引き戻した後、
必要な基底線束の捩れを除けば
$\tau_t^*$によって全Homと積を保って対応する。
実際の各頂点線束の同型を同時に選ぶには、
さらにfppf局所化してよい。
これは幾何的finite Endの比較であって、
相対AS mutationや相対pointed$\qgr$同値を証明しない。

一般係数体$K_{12}$上では$T_P$に有理点がない。
従って上記の半周期による四頂点交換を
一般係数体上で一点の選択として実現することはできない。
\end{theorem}
\begin{proof}
$E$は元の二conicの係数空間$U^\circ$上で定義される。
gerbe上の$F,G$の相対Picard類は
線束の基底からの捩れに不変であり、
各$\Gm$慣性も相対Picard類には自明に作用する。
よって$P=[GF^{-1}]$は$J$のsectionへ降下する。
各幾何fiberで非零二捩れであることは
補題~\ref{lem:qa-pointed-all-rootpair-geometry}による。
標数零の楕円曲線の$[2]$は有限\'{e}tale次数四であるため、
base changeした$T_P$も同じ性質を持ち、
各fiberは$J[2]$の自由かつ推移的なtorsorである。

滑らかな種数一曲線$E_s$を幾何点$s$上でとり、
一時的に原点を選ぶ。
次数$d$の線束$L$は次数$d$の因子の線型同値類であり、
平行移動によるpullbackは
\[
 [\tau_t^*L]-[L]=-d t\quad\text{in }\Pic^0(E_s)
 \tag{Translation-degree}\label{eq:qa-half-period-degree}
\]
となる。実際、$\sum_\nu x_\nu$のpullbackは
$\sum_\nu(x_\nu-t)$で、両者の差は$-dt$である。
この公式は原点の選択によらないため相対Picardに降下する。
$F,G$の次数はともに二、$H$の次数は四であり、
$2t=P$かつ$2P=0$だから
\[
 -2t=-P=P,\qquad -4t=-2P=0.
\]
従って$\tau_t^*$は$F$を$FP=G$に、
$G$を$GP=F$に、$H$を$H$に送る。
積についても同じであり、
式~\eqref{eq:qa-half-period-exchange}を得る。
曲線の自己同型のpullbackは射の合成を保つので、
各頂点の線束同型をfppf局所で選べば
有限Endの全合成が比較できる。
ただしその選択のdescent datumまではここで作っていない。

最後に$K_{12}$上で$2t=P$となる
$K_{12}$-有理点$t$が存在すると仮定する。
この$t$による平行移動は
$K_{12}$-自己同型$\tau_t$である。
幾何Picard類$F,G$のHochschild--Serre境界は
前節の命題~
\ref{prop:qa-universal-conic-brauer-obstruction}により
$\partial_E[F]=\alpha_c$、
$\partial_E[G]=\alpha_d$である。
この境界は基礎体上の自己同型pullbackと可換し、
右端の$\operatorname{Br}(K_{12})$には自明に作用する。
$\tau_t^*[F]=[G]$から
$\alpha_c=\alpha_d$が従うが、
同命題で$\alpha_c-\alpha_d\ne0$と証明している。
矛盾するので$T_P(K_{12})=\varnothing$である。
\end{proof}

\begin{remark}[異なる次数四被覆と特異境界]
\label{rem:qa-half-period-versus-root-cover}
半周期torsor $T_P$は$J[2]$-torsorであり、
前節の「各二次形式から分岐根を一つずつ選ぶ」
有限\'{e}tale次数四被覆とは定義が異なる。
両被覆の同型、OUの一般次数四射との比較、
およびそのGaloisモノドロミーの一致は未確認である。
nodal/$I_2$では相対$\Pic^0$の$[2]$の幾何次数が
smooth fiberの四から変わるため、
上記の同じ次数四の\'{e}tale torsorを
境界まで自動的に延長してはならない。
\end{remark}

% !TeX root = ../nc-quadric-avatar.tex
% 2026-11-22: 共有根 A1 の最小解消に残る相対 Pic^0 障害
\section{共有根の最小解消と相対 Picard の奇偶障害}
\label{sec:qa-nodal-resolution-parity}

定理~\ref{thm:qa-noncubic-rootpair-torsionfree-spectral-module}
の相対 spectral module は共有根で非局所自由となる．
この非自由性は，その点を最小解消して線束にすれば
自動的に通常の相対 $\Pic^0$ に戻る，という意味ではない．
以下は元の一径数基底を変えない最小解消での具体的な障害である．

\begin{proposition}[最小解消後にも残る奇数次数]
\label{prop:qa-nodal-resolution-odd-exceptional-degree}
$k$ は標数零の代数閉体とし，
$A=k[[a,b,y]]/(ab-y^2)$，
$s=a-b$ とおく．
符号商
\[
 A=k[[c,d]]^{\{\pm1\}},\qquad
 a=c^2,\quad b=d^2,\quad y=cd
\]
の反不変部分
$M=k[[c,d]]^- =Ac+Ad$
を考える．$\pi:\widetilde X\to\Spec A$ を
$A_1$ 特異点の最小解消，$E\simeq\PP^1$ を
例外曲線とする．
\[
 L=(\pi^*M)/(\text{torsion})
\]
は $\widetilde X$ 上可逆で，その例外曲線への制限は
\[
                    L|_E\simeq\OO_{\PP^1}(1).
\]
特に任意の $n\in\mathbb Z$ に対し
$\deg_E L(nE)=1-2n$ は奇数であり，
相対 $\Pic^0$ に属する多重次数零の線束にはならない．

さらに，$s$ を基底パラメータとする適切な大域曲線族で
特殊曲線が既約一節点であれば，最小解消後の特殊
fiber は二つの成分 $D'$ と $E$ を持ち，
$D'\cdot E=2$ である．
同じ generic fiber 上で $M$ と一致する
任意の可逆延長は，$L$ から垂直 Cartier 因子で
tensor したものであるから，
$E$ 上の次数の奇偶は変わらない．
従って\emph{元の一径数基底上で最小解消するだけでは}，
非自明二捩れの generic fiber を
多重次数零の可逆線束の相対族として延長できない．
\end{proposition}
\begin{proof}
最小解消の標準二枚のアフィン chart は
\[
\begin{array}{c|ccc}
 &a&y&b\\ \hline
 U_c&u&ut&ut^2\\
 U_d&vw^2&vw&v
\end{array},\qquad
v=ut^2,\quad w=t^{-1}.
\]
例外曲線は $u=0$ と $v=0$ を貼り合わせた
$\PP^1_{[c:d]}$ で，自己交点は $E^2=-2$ である．
$M$ の生成元を形式的に $e_c=c,\ e_d=d$ と書く．
$A$ 上の関係 $yc=ad$ と $yd=bc$ を引き戻して
torsion を除くと，
$U_c$ では $e_d=t e_c$，
$U_d$ では $e_c=w e_d$ となる．
従って各 chart では一つの自由生成元を持ち，
重なりでの遷移関数は
$e_d=t e_c$ である．これは
$\OO_{\PP^1}(1)$ の標準遷移である．
$\OO(E)|_E=\OO_E(-2)$ だから
$\deg_E L(nE)=1-2n$ が従う．

$s=a-b$ を引き戻すと
$U_c$ では $s=u(1-t^2)$，
$U_d$ では $s=v(w^2-1)$ となる．
従って特殊 fiber には例外成分 $E$ が重複度一で現れ，
元の節点の二枝の strict transform が
$E$ の $t=1,-1$ の二点に横断的に交わる．
大域的に元の特殊曲線が既約一節点なら，
この二枝は同じ既約 strict transform $D'$ に属する．
ゆえに $D'\cdot E=2$ である．
二つの可逆延長が generic fiber 上同型なら，
正則曲面上のその比は垂直 Cartier 因子で表される．
その成分は $E$ と $D'$ の整数係数の和なので，
$E$ 上の次数の変化は
$nE^2+mD'\cdot E=-2n+2m$ で偶数である．
よって多重次数零への延長は存在しない．
\end{proof}

\begin{corollary}[相対 avatar の修正に必要な追加操作]
\label{cor:qa-nodal-resolution-requires-more}
共有根の noncubic spectral avatar を
通常の相対 $\Pic^0$ へ移すには，
元の基底を変えずに全空間の $A_1$ を最小解消するだけでは
不十分である．
基底変更，非可逆 rank-one torsion-free 層を許す
compactified Picard，または安定モデルと偏極を
同時に変更する構成のいずれかを別途検討する必要がある．
この命題は，それらの追加操作を施した場合にも
相対射が存在しないと主張するものではない．
\end{corollary}
\begin{proof}
命題~\ref{prop:qa-nodal-resolution-odd-exceptional-degree}
で可逆延長の例外成分上の次数が奇数になる．
相対 $\Pic^0$ の多重次数零の線束にはなり得ない．
\end{proof}


\subsection{Spectral 二捩れ乗法の欠損と正規化上の極限}
\label{subsec:qa-spectral-compactified-limit}

前の奇偶障害は，線束の極限が存在しないという意味ではなく，
\emph{可逆線束に限った極限が存在しない}ことを意味する．
相対 torsion-free sheaf 自身には自然な二次積があり，
その余核によって境界を scheme-theoretic に測定できる．
以下では定理~\ref{thm:qa-noncubic-rootpair-torsionfree-spectral-module}
の記号を使い，二形式は各々二つの相異なる根を持つとする．

\begin{proposition}[乗法欠損の長さと自己双対性]
\label{prop:qa-spectral-product-defect}
$B=\PP^1_k$，$f_1,f_2\in H^0(B,\OO_B(2))$ は
非零で各々 squarefree，$k$ は標数零の代数閉体とする．
$C=\{y^2=f_1f_2\}$ と
$\mathcal M=(r_*\OO_{\widetilde C})^-$ を取る．
反不変元の積から生じる
\[
 \mu:\mathcal M\otimes_{\OO_C}\mathcal M\longrightarrow\OO_C
\]
について，次が成り立つ．
\begin{enumerate}
\item $\operatorname{im}\mu=\mathcal I=(f_1,f_2,y)$ であり，
  $T=V_B(f_1,f_2)$ を $y=0$ によって $C$ に埋め込むと
  \[
     \coker\mu\simeq\OO_T,\qquad
     \operatorname{length}(\coker\mu)
       =\deg\gcd(f_1,f_2)\in\{0,1,2\}.
     \tag{Sq-defect}\label{eq:qa-spectral-square-defect}
  \]
  この余核は二根対の共通零点の
  \emph{scheme-theoretic} な構造層である．
\item 各幾何 fiber において，積の随伴射
  \[
   \mu^\flat:\mathcal M
       \longrightarrow
       \mathcal H om_C(\mathcal M,\OO_C)
  \]
  は同型である．つまり $\mathcal M$ は自己双対だが，
  $\mu$ が全射となるのは二根対が交わらない場合に限る．
\end{enumerate}
特に $\mu^\flat$ の同型と
可逆線束としての $\mathcal M^{\otimes2}\simeq\OO_C$
とは，節点では異なる主張である．
\end{proposition}
\begin{proof}
相対射 $r:\widetilde C\to C$ に関する
同時符号の反不変元は局所的に $c,d$ で生成され，
$c^2=f_1,\ d^2=f_2,\ cd=y$ である．
従って積の像は記載のイデアルで，
\[
 \OO_C/(f_1,f_2,y)
 \simeq\OO_B/(f_1,f_2)
\]
が従う．$f_i$ は各々 squarefree なので，
その共通零点は $0,1,2$ 個の被約点である．

自己双対性は，二形式が共通根を持たない開集合では
通常の非分岐二重 torsor の二次積から従う．
共通根では完備局所座標と単元の平方根を選ぶと，
\[
 A_0=k[[x,y]]/(y^2-x^2),\quad
 M_0=\langle c,d\mid yc=xd,\ yd=xc\rangle,
 \qquad c^2=d^2=x,\quad cd=y .
\]
ここで $u=c+d,\ v=c-d$ とすれば
\[
  M_0\simeq A_0/(y-x)\,u
          \ \oplus\ A_0/(y+x)\,v,
  \qquad u^2=2(x+y),\quad
  v^2=2(x-y),\quad uv=0.                 \tag{Node-M}
\]
環 $A_0$ において
\[
 \begin{split}
  \Hom_{A_0}(A_0/(y-x),A_0)&=(x+y),\\
  \Hom_{A_0}(A_0/(y+x),A_0)&=(x-y)
 \end{split}
\]
であり，積からの $u\mapsto 2(x+y)$ と
$v\mapsto2(x-y)$ は各直和因子で同型となる．
異なる枝間の Hom は零なので，$\mu^\flat$ は同型である．
各共通根は同じ局所型であり，主張が従う．
\end{proof}

\begin{theorem}[noncubic 三型の compactified spectral limit]
\label{thm:qa-spectral-three-fiber-limits}
前命題の仮定を置き，$\nu:\widehat C\to C$ を
節点での正規化とする．
\begin{enumerate}
\item 二根対が交わらない smooth 型では
  $\mathcal M$ は非自明な $\Pic^0(C)[2]$ の線束であり，
  $\End_C(\mathcal M)=k$ である．
\item ちょうど一根共有する既約 nodal 型では
  $\widehat C\simeq\PP^1$ であり，
  \[
       \mathcal M\simeq\nu_*\OO_{\PP^1}(-1),
       \qquad \End_C(\mathcal M)=k .
       \tag{Node-limit}\label{eq:qa-spectral-nodal-limit}
  \]
  これは次数零の rank-one torsion-free \emph{単純}層だが，
  $\Pic^0(C)\simeq\Gm$ の非自明二捩れ線束
  （節点の貼合せ比 $-1$）とは同型でない．
\item 二根対が一致する $I_2$ 型では
  $\widehat C=C_+\amalg C_-$，$C_\pm\simeq\PP^1$ であり，
  \[
       \mathcal M\simeq
       \nu_*(\OO_{C_+}(-1)\oplus\OO_{C_-}(-1)),
       \qquad \End_C(\mathcal M)\simeq k\oplus k .
       \tag{I2-limit}\label{eq:qa-spectral-i2-limit}
  \]
  従って $\Aut_C(\mathcal M)\simeq\Gm^2$ で，
  この層は\emph{単純でない}．
\end{enumerate}
三型とも $\chi(\mathcal M)=0$ である．
従って，前定理の全相対族を受けるモジュライを
\emph{単純な} rank-one torsion-free sheaf に限定すると，
$I_2$ 型の点が脱落する．
ただし適切な偏極に関する半安定層の粗モジュライを
排除するものではない．
\end{theorem}
\begin{proof}
(1) は定理
~\ref{thm:qa-noncubic-rootpair-torsionfree-spectral-module}
により従う．
各節点での式 (Node-M) は
正規化の二枝の直像の局所模型そのものである．
従って非局所自由な全節点で正規化すると，
$\mathcal M=\nu_*Q$ となる線束 $Q$ が存在する．

(2) の正規化は種数零の滑らかな連結曲線であり，
$Q\simeq\OO_{\PP^1}(d)$ と書ける．
既知の
$p_*\mathcal M\simeq\OO_{\PP^1}(-1)^{\oplus2}$
によって $\chi(\mathcal M)=0$ なので，
$\chi(Q)=d+1=0$，従って $d=-1$ である．
また
$\mathcal E nd_C(\nu_*Q)\simeq\nu_*\OO_{\widehat C}$
であるから，大域 End は $k$ となる．
$\Pic^0(C)$ の元は節点でも局所自由なので，
非局所自由な $\mathcal M$ とは同型にならない．

(3) では比例する二形式をスカラーで規格化して
$f_1=f_2=f$ とできる．曲線は $y=\pm f$ の二成分で，
正規化後の各成分は $B\simeq\PP^1$ に同型である．
$y$ の $c,d$ 上の作用は
$\left(\begin{smallmatrix}0&f\\f&0\end{smallmatrix}\right)$
なので，$c+d$ と $c-d$ がそれぞれ
$y=f$ と $y=-f$ に支持される固有生成元となる．
$p_*\mathcal M=\OO_B(-1)c\oplus\OO_B(-1)d$ から
各固有成分は $\OO_B(-1)$ である．
異なる既約成分間には非零な sheaf Hom はなく，
各成分の自己準同型環は $H^0(\OO_{\PP^1})=k$ だから
End は $k\oplus k$ となる．
最後の留保は，simple/fine と
semistable/coarse の compactification を
混同しないためである．
\end{proof}

\begin{remark}[相対 avatar に必要なモジュライの選択]
\label{rem:qa-spectral-compactification-choice}
既約 nodal 型では $\nu_*\OO_{\PP^1}(-1)$ が
単純な compactified Picard の境界点となる一方，
$I_2$ 型では automorphism が $\Gm^2$ になる．
rank-one torsion-free sheaf の compactification と
安定性条件の一般論は
\cite{AltmanKleiman1980,ChristPayneShen2023} を参照する．
本稿の固有な入力は式~\eqref{eq:qa-spectral-square-defect}と
二根対の三型に沿った explicit な極限表示である．
今後の正確な比較対象は，
相対曲線 $C/B$ と $\mathcal M$ のみならず，
\emph{自己双対な二次積 $\mu^\flat$ と
長さ $0,1,2$ の積欠損}も保持する群oidである．
この群oidと avatar AS 心の同値は未証明である．
\end{remark}

% !TeX root = ../nc-quadric-avatar.tex
% 2026-11-22 / C v4.9
\section{節点spectral曲線上の二次自己双対層の代数性}
\label{sec:qa-tf-quad-artin}

本節では標数零の体$k$上で，幾何fiberがsmooth，既約nodal，
または$I_2$となるspectral二重被覆だけを扱う．
非被約曲線を許す元の$\mathfrak{TF}^{\rm quad}$全体の代数性・有界性を
ここから主張しない．またrank oneは\emph{各既約成分上でrank one}を意味する．
単純層であること，安定性，固定多重次数を定義に加えない．

\subsection{曲線の有限型基底と定義}

$Q_{\rm nod}\subset\operatorname{Sym}^4(k^2)$を，零でなく，
幾何的な根の重複度が高々二となる二元四次形式の開集合とする．
\[
 H=\bigl(\GL_2\times\Gm\bigr)
       /\{(\lambda I,\lambda^2)\},\qquad
 (A,a)\cdot h=a^{-2}\operatorname{Sym}^4(A)h
\]
とし，$\mathfrak D_{\rm nod}=[Q_{\rm nod}/H]$と置く．
この有限型Artin stackは，conic束$q:B\to S$と，
相対次数$-2$の線束$N$，射$h:N^{\otimes2}\to\OO_B$からなる
有限平坦二重被覆
\[
 p:C=\operatorname{Spec}_B(\OO_B\oplus N)\longrightarrow B
\]
で，上記三型の幾何fiberを持つものを分類する．
実際，fppf局所で$B=\PP^1$，$N=\OO(-2)$とすると$y^2=h(u,v)$となり，
表示変更の核は記載した共通重みスカラーである．
$C/S$は平坦・射影・有限表示で，幾何fiberは被約・連結・Gorensteinである．
二形式族からは$h=f_1f_2$によってこの基底への射を得る．

$\mathfrak T_{\rm nod}\to\mathfrak D_{\rm nod}$を次のstackとする．
対象は$S$-flatな有限表示$\OO_C$-module $M$で，
各幾何fiberでrank-one torsion-free，$\chi(M_s)=0$となるものと，
対称な$\OO_C$-双線型射
\[
 \mu:M\otimes_{\OO_C}M\longrightarrow\OO_C
\]
で，各幾何fiberの随伴射$\mu_s^\flat:M_s\to M_s^\vee$が同型となるもの．
射は被覆と二次積付き層の同型である．
これは既存の$\mathfrak{TF}^{\rm quad}$を上記の節点曲線基底へ制限したものである．

\subsection{Frobenius表示とtrace零条件}

\begin{lemma}[相対直像の局所自由性と双対]
\label{lem:qa-tf-frobenius-duality}
上記の$M$について$V=p_*M$は$B$上の階数二ベクトル束である．
$A=p_*\OO_C=\OO_B\oplus N$の積と射影$\lambda:A\to N$は完全pairing
\[
 A\otimes_{\OO_B}A\longrightarrow N,\qquad
 (a,b)\longmapsto\lambda(ab)
\]
を与える．従って
\begin{equation}
 p_*\mathcal H om_C(M,\OO_C)
   \simeq\mathcal H om_B(V,N)=V^\vee\otimes N
 \label{eq:qa-tf-frobenius-dual}
\end{equation}
は任意基底変更と可換である．
$\mu$の$N$成分$b=\lambda\circ p_*\mu$は
$\operatorname{Sym}^2V\to N$という非退化二次形式になる．
この非退化性は$\mu_s^\flat$の全幾何fiberでの同型と同値で，
その場合$\mu^\flat$自体も相対的に同型である．
\end{lemma}
\begin{proof}
各幾何fiberで$p_*M_s$は$\PP^1$上のtorsion-free層である．
実際，基底曲線の局所パラメータは$C_s$の全成分上で非零であり，
$M_s$に非零因子として作用する．従って$p_*M_s$は階数二のベクトル束となる．
$V$は$S$-flatで有限表示だから，fiberによる平坦性判定を適用すると
$B$上で平坦，したがって局所自由となる
\cite[Theorem~76.23.3, Tag~05X0]{StacksCoherentSheaves}．

$N$の局所生成元を$y$とすると，pairingの行列は
$\left(\begin{smallmatrix}0&1\\1&0\end{smallmatrix}\right)$である．
従って$A\simeq\mathcal H om_B(A,N)$であり，有限射の随伴から記載の双対同型を得る．
右辺はベクトル束の双対なので，非平坦基底変更にも可換する．
$\mu^\flat$はこの同型の下で$b^\flat:V\to V^\vee\otimes N$になる．
同階数ベクトル束の射が全幾何fiberで同型なら行列式は単元であり，
全体で同型である．
\end{proof}

\begin{lemma}[trace零は非被約基底でも成立する]
\label{lem:qa-tf-trace-zero}
$A$の$N$成分の作用を$\theta:N\otimes V\to V$と書くと，
\[
 \theta(1_N\otimes\theta)=h\otimes1_V,
 \qquad\operatorname{Tr}(\theta)=0
\]
が成立する．逆に$V$が階数二ベクトル束で，$\theta$がこの二条件を満たせば，
対応する$A$-moduleは全幾何fiberで各既約成分上rank oneかつtorsion-freeである．
\end{lemma}
\begin{proof}
第一式は$A$の作用の結合性である．第二式を幾何fiberで確認するだけでは不十分なので，
以下の議論を用いる．$h$は各fiberで零でない四次形式である．
fppf局所で$B=\PP^1_S$とし，$h$の零点でない無限遠点を選べば，
アフィンchart上の$h$はleading係数が単元の四次多項式となる．
従って$h$は任意基底変更後にも非零因子である．
$U=B\setminus V(h)$では$C_U\to U$は有限\'etaleで，$C_U/S$は滑らかである．
$M$は$S$-flatで，各fiberのこの開上では線束だから，
再びfiberによる平坦性判定により$M|_{C_U}$は実際の線束である．
その$N$作用のtraceは$A$自身のtraceと等しく零となる．
$\operatorname{Tr}(\theta)$は$N^\vee$のsectionであり，
$h$を可逆化した開で零，$h$は非零因子なので，全$B$上で零である．

逆に幾何fiberの$B_s$のgeneric pointで，$A_s$は二次\'etale代数である．
非分裂の場合，二次体上の$V$はrank oneとなる．
分裂する場合も，$\theta$のtrace零条件と$\theta^2=h$により，
異なる二固有値の固有空間はそれぞれ一次元である．
従って$M_s$は全成分上rank oneである．
$V_s$が$B_s$上局所自由なので$M_s$に零次元torsionはなく，
被約曲線上でtorsion-freeとなる．
\end{proof}

\begin{remark}[rank one条件を落としてはならない]
例えば$I_2$で$h=a^2$の場合，$\theta=aI$は$\theta^2=hI$を満たし，
非退化双曲型$b$に関して自己随伴であるが，一方の成分にrank two，
他方にrank zeroのmoduleを与える．これは本節の対象ではなく，
$\operatorname{Tr}(\theta)=0$がこの誤った成分を除外する．
\end{remark}

\subsection{一様有界性と有限型Artin stack}

\begin{lemma}[自己双対rank-one層の一様有界性]
\label{lem:qa-tf-uniform-boundedness}
任意の幾何fiberで
\[
 h^0(C_s,M_s)\le1,
\qquad
 p_*M_s\simeq
 \OO_{\PP^1}(-1)^{\oplus2}
 \quad\text{または}\quad
 \OO_{\PP^1}\oplus\OO_{\PP^1}(-2).
\]
後者は$M_s\simeq\OO_{C_s}$のときに限る．
$\mathcal L=p^*\omega_{B/S}^{-1}$と置くと，全幾何fiberで
$M_s\otimes\mathcal L_s$は大域生成され，$H^1=0$，$h^0=4$である．
$\mathcal L$に関する$M$のHilbert多項式は$4n$である．
\end{lemma}
\begin{proof}
$0\ne s\in H^0(C_s,M_s)$があれば，torsion-free性から
$s$は少なくとも一つの成分のgeneric pointで非零である．
そこでは$\mu$は一次元空間上の非退化対称形式なので$\mu(s,s)\ne0$である．
$C_s$は\emph{被約かつ連結}で$H^0(\OO_{C_s})=k(s)$だから，
$\mu(s,s)$は非零定数である．
よって$\OO_{C_s}\xrightarrow{s}M_s$は$\mu(s,-)/\mu(s,s)$により分裂する．
残る直和因子はrank zeroかつtorsion-freeなので零となり，$M_s\simeq\OO_{C_s}$である．
従って$h^0\le1$である．この論証を非被約曲線へ適用してはならない．

$\chi(p_*M_s)=0$より$p_*M_s$の次数は$-2$である．
$\PP^1$上の分解型と$h^0\le1$から，記載した二型だけが残る．
いずれも$(p_*M_s)(2)$は大域生成，$H^1=0$，$h^0=4$である．
有限射の評価$p^*p_*M_s\to M_s$は全射なので，$M_s\otimes\mathcal L_s$にも従う．
$\deg\mathcal L_s=4$と$\chi(M_s)=0$からHilbert多項式を得る．
\end{proof}

\begin{theorem}[節点曲線上の二次自己双対標的の代数性]
\label{thm:qa-tf-nodal-finite-type-artin}
$\mathfrak T_{\rm nod}$は$k$上有限型のArtin stackであり，
$\mathfrak D_{\rm nod}$に対して局所有限表示，対角はaffineかつ有限表示である．
$k$上の対角もaffineかつ有限表示である．
定義・構成は任意の基底変更と可換し，nonsimpleな$I_2$の層を含む．
\end{theorem}
\begin{proof}
まず$\mathfrak D_{\rm nod}$のatlas $Q_{\rm nod}$へ引き戻す．
この上では$B=\PP^1$，$N=\OO(-2)$である．
階数二・次数$-2$のベクトル束のstack $\operatorname{Bun}_{2,-2}(B)$を用いる．
これはcoherent sheaf stackの開部分で，Artin，局所有限表示，
affine有限表示対角を持つ
\cite[Theorem~99.6.1, Tag~09DS; Lemmas~108.4.1--108.4.2]{StacksCoherentSheaves}．

その普遍$V$上に$\theta:N\otimes V\to V$と
$b:\operatorname{Sym}^2V\to N$を加える．
これらのHom functorはaffine有限表示である
\cite[Proposition~99.3.10, Tag~08K6]{StacksCoherentSheaves}．対象側は有限表示・基底flat・proper supportである．
条件$\theta^2=h$，$\operatorname{Tr}\theta=0$，および
\[
 b(\theta(n\otimes v),w)=b(v,\theta(n\otimes w))
\]
はHom functor上の閉条件である．
$b$の非退化性は開条件である：行列式の零点の$B$上でのsupportは閉であり，
$B$が基底上properなのでその像も閉になる．

これらの閉条件と開条件で定めたstackは$\mathfrak T_{\rm nod}$と同値である．
前二補題が一方向を与える．逆向きには$\theta$で$A$-moduleを作ると，
trace零により全成分rank oneとtorsion-free性を得る．
$V$が$B$上局所自由なので対応する$M$は基底flatで有限表示である．
Frobenius双対により$b$は対称な$A$-双線型$\mu$を一意に与える．
実際，$N$の局所生成元$y$で$b=b_0y$，$\theta(y,-)=Y$と書けば
\[
 \mu(v,w)=b_0(Yv,w)+y\,b_0(v,w).
\]
自己随伴条件と$Y^2=h$から$A$-双線型性と対称性が直接確認できる．
$b$が非退化なので$\mu^\flat$も同型である．
これはすべての射と合成について自然で，任意基底変更とも両立する．

この構成からArtin性と局所有限表示を得る．
二対象のIsomはベクトル束のaffine有限表示Isomの中に
$\theta$と$b$を保存する閉条件を課したものなので，対角もaffine有限表示である．
$\mathfrak D_{\rm nod}=[Q_{\rm nod}/H]$自身もaffine有限表示対角を持つため，
$k$上の対角にも従う．

最後にquasi-compactnessを確認する．前補題により全対象で
$V(2)$は大域生成，$R^1q_*V(2)=0$，$q_*V(2)$は階数四で基底変更と可換である．
従ってそのframeを取ると，全対象は一つの固定Hilbert多項式
$2n+4$のQuot空間
\[
 \operatorname{Quot}^{2n+4}_{\OO_{\PP^1}^{\oplus4}/\PP^1/Q_{\rm nod}}
\]
で被覆される．具体的にはquotientが階数二ベクトル束で，
大域切断への写像$\OO^4\to q_*V(2)$が同型となる開部分を取る．
固定Hilbert多項式のQuotは有限表示であり，底$Q_{\rm nod}$はNoetherianなので，
この開もquasi-compactで有限型である
\cite[Lemma~108.6.1, Tag~0DP9]{StacksCoherentSheaves}．
その上に上記二Homを加え閉条件・開条件を課すと有限型空間を得る．
frameを忘れる射は$\GL_4$-torsorで滑らか全射である．
従って$\mathfrak T_{\rm nod}$は有限型となる．
\end{proof}

\begin{remark}[固定Hilbert多項式だけでは足りない]
自己双対性を外すと，同じ$I_2$曲線上の多重次数$(m,-m)$の線束は
すべてHilbert多項式$4n$を持つが，$m$は非有界である．
上の有限型性は固定多項式だけによらず，
rank one・被約連結性・非退化二次積から得た一様生成性に依存する．
また本節は非被約quartic曲線上の標的や，その全体の有限型性を扱わない．
\end{remark}

\subsection{Balanced開と等方線の相対降下}

\begin{lemma}[balanced開と基底上の等方線二重被覆]
\label{lem:qa-tf-balanced-isotropic-descent}
$\mathfrak T_{\rm nod}$内の条件$H^0(C_s,M_s)=0$は開であり，
$V_s=\OO_{\PP^1}(-1)^{\oplus2}$と同値である．
この開上で$b$の等方線を分類する$Y\to B$は有限\'etale二重被覆となる．
ある有限\'etale二重被覆$S'\to S$が自然に定まり，
$Y\simeq B\times_SS'$となる．
$S'$上で$V=K_1\oplus K_2$と等方線に分裂し，
$K_i$の相対次数は$-1$，$K_1\otimes K_2\simeq N$である．
\end{lemma}
\begin{proof}
$H^0$の上半連続性と分解型の補題から開性と同値を得る．
非退化な二元対称形式の等方線schemeは，標数二でないので有限\'etale次数二である．
各幾何fiberの$B_s=\PP^1$上ではこの被覆は自明な二重被覆である．
$\pi_Y:Y\to S$はproper smoothで，幾何fiberは二つの$\PP^1$の非交和，
$H^0=2$，$H^1=0$である．従って$\pi_{Y*}\OO_Y$は任意基底変更と可換な
階数二の局所自由代数で，全幾何fiberが\'etaleなので有限\'etale代数となる．
$S'=\operatorname{Spec}_S(\pi_{Y*}\OO_Y)$と置く．
$Y\to B\times_SS'$に対応する$B$上の階数二局所自由代数間の射は
全幾何fiberで同型だから，全体でも同型である．
二つの異なる等方線の和は$V$全体であり，交差pairingは完全である．
相対次数はbalanced分解型から従う．
\end{proof}

等方線上の平方は$\mu(K_i^{\otimes2})\subset\OO_B$となり，
$K_i^{-2}$の相対二次sectionを与える．
このsectionがsquarefreeである条件は開であり，二重被覆上の両条件を課して
基底へ降下できる．既存の二形式族はこの開に入り，
$I_2$の$M=\nu_*(\OO_{\PP^1}(-1)\oplus\OO_{\PP^1}(-1))$も含まれる．
この層は$\End(M)=k^2$でnonsimpleだが，二次積を保つ慣性は$\mu_2^2$である．
したがってsimple sheafだけのcompactified Picardに置換してはいけない．
等方線と平方sectionから二形式を回収するgroupoid同値は，
相対認識の別命題で述べる．

\subsection{Balanced二次積の相互逆な有限表示}

\begin{proposition}[二形式の再構成と結合律の自動性]
\label{prop:qa-tf-balanced-quadratic-inverse}
補題~\ref{lem:qa-tf-trace-zero}の状況で，$V=p_*M$はfiberwiseに
$\OO_{\PP^1}(-1)^{\oplus2}$であるとする．
また対称な$\OO_C$-双線型写像
$\mu:M\otimes_{\OO_C}M\to\OO_C$が各幾何fiberで自己双対性を誘導するとする．
このときfppf局所的に二元二次形式$f_1,f_2$と生成元$c,d,y$を選んで
\[
 \begin{gathered}
 C=\{y^2=f_1f_2\},\qquad
 p_*M=\OO_B(-1)c\oplus\OO_B(-1)d,\\
 yc=f_1d,\qquad yd=f_2c,\\
 \mu(c,c)=f_1,\qquad\mu(d,d)=f_2,\qquad\mu(c,d)=y
 \end{gathered}
\]
と書ける．
従って$\OO_C\oplus M$に
\[
 (a,m)(b,n)=(ab+\mu(m,n),\ an+bm)
\]
で定めた積は可換かつ結合的であり，
\begin{equation}
 \operatorname{Spec}_C(\OO_C\oplus M)
 \simeq
 \operatorname{Spec}_B
 \OO_B[c,d]/(c^2-f_1,d^2-f_2)
 \label{eq:qa-tf-reconstructed-cover}
\end{equation}
となる．右辺では$c,d$は次数一のtwisted生成元である．

同型の変更は$B$の座標変換，$c,d$の個別スカラー，二項の交換であり，
共通次数スカラーを除けば二形式表示の同型に一致する．
特に追加の結合律条件は，このbalancedかつgenerically reducedな範囲では不要である．
\end{proposition}
\begin{proof}
Balanced条件とcohomology and base changeにより，$S$上fppf局所的に
$V=\OO_B(-1)^{\oplus2}$と書ける．
tracezero線も$L=\OO_B(-2)$と書ける．
補題~\ref{lem:qa-tf-frobenius-duality}により$q=\lambda\mu$は非退化である．
その行列の成分は$H^0(B,\OO_B)=\OO_S$に属するので，
rank two非退化二次形式を$S$上fppf局所的にsplitすれば，
\[
 q(c,c)=q(d,d)=0,\qquad q(c,d)=y
\]
とできる．従って二元二次形式$f_1,f_2,a$を用いて
\[
 \mu(c,c)=f_1,\quad\mu(d,d)=f_2,
 \quad\mu(c,d)=a+y
\]
と書ける．

まず結合律を仮定せず，$\OO_C$-双線型性だけを使う．
$Yc=r c+s d$と置き，
$\mu(Yc,c)=y\mu(c,c)$および
$\mu(Yc,d)=y\mu(c,d)$の$1,y$係数を比較すると
\[
 s=f_1,\qquad r=a,\qquad 2af_1=0,
 \qquad h=a^2+f_1f_2
\]
を得る．$Yd$についても同様に
\[
 Yd=f_2c+a d,\qquad2af_2=0.
\]
従って$Y$の行列は
\[
                 Y=\begin{pmatrix}a&f_2\\f_1&a\end{pmatrix}.
\]
補題~\ref{lem:qa-tf-trace-zero}の$\operatorname{Tr}(Y)=0$と$2$の可逆性より$a=0$．
以上の全表示を得る．
なお，少なくとも一方の$f_i$がfiberwise squarefreeである開では，
その$f_i$の相対Cartier性と$2af_i=0$だけからも$a=0$が直接従う．

$\OO_C\oplus M$は$B$上で$1,c,d,y$を基底とするrank four vector bundleであり，
得られた式は$c^2=f_1,d^2=f_2,cd=y$の積表に一致する．
右辺の二つのmonic二次関係は$B$上同じ四元基底を持つため，
式\eqref{eq:qa-tf-reconstructed-cover}は代数同型である．
特に結合律はこの実際の代数同型から従う．

最後に同型の形を確認する．tracezero線は有限平坦二重射のtraceで内在的に定まり，
したがって全同型は$q$のsimilitudeを誘導する．
$q(xc+zd)=2xz\,y$の二つの等方線は$c,d$の非順序対である．
任意の可換環上で$\left(\begin{smallmatrix}\alpha&\beta\\\gamma&\delta\end{smallmatrix}\right)$が
この形式のsimilitudeなら，
\[
 \alpha\gamma=\beta\delta=0,
 \qquad\alpha\delta+\beta\gamma\in R^\times.
\]
$e=\alpha\delta/(\alpha\delta+\beta\gamma)$は冪等元であり，
$e=1$側では行列は対角，$e=0$側では反対角となる．
従って冪零元による隠れた混合はなく，等方線の非順序対は任意基底で回復される．
二形式の二乗とcross積が全残余の係数を決める．
射影曲線とその被覆に自明に作用する共通スカラーは，表示群の$D$に一致する．
\end{proof}

\subsection{等方線による二重平方表示の相対認識}
\label{subsec:qa-tf-relative-recognition}

$\mathfrak T_{\rm bal}\subset\mathfrak T_{\rm nod}$を
$H^0(C_s,M_s)=0$で定まるbalanced開とする．
$W=k^2$に対し
\[
 \overline P_{\rm nod}
 =\{(f_1,f_2)\in\operatorname{Sym}^2W\times\operatorname{Sym}^2W
       \mid f_1f_2\ne0,\ \text{全ての幾何的根の重複度が高々二}\}
\]
と置く．これは六次元アフィン空間の開集合である．
各$f_i$が二重根を持つ場合を許すが，その積の三重根は許さない．
表示群を
\[
 G=((\operatorname{GL}(W)\times\Gm^2)\rtimes S_2)/D,
 \qquad D=\{(\lambda I,\lambda,\lambda)\}
\]
とする．両形式がsquarefreeとなる開は従来の$P$である．

\begin{theorem}[balanced標的の群商表示と元のnoncubic開]
\label{thm:qa-tf-balanced-quotient-recognition}
二重平方表示からの構成は，任意の$k$-schemeを底とするfppf stackの同値
\[
 [\overline P_{\rm nod}/G]\xrightarrow{\sim}\mathfrak T_{\rm bal}
\]
を与える．非被約底も含み，任意の基底変更と可換する．

この同値の下で，元の表示stack$\mathfrak N^{\rm pres}=[P/G]$は，
二つの等方線上の平方sectionが両方ともsquarefreeとなる開部分
$\mathfrak T_{\rm adm}\subset\mathfrak T_{\rm bal}$に一致する．
特に既存の射
\[
 \mathfrak N^{\rm pres}\longrightarrow\mathfrak{TF}^{\rm quad}
\]
は全忠実であり，節点曲線上の標的$\mathfrak T_{\rm nod}$への射としては
representableな開埋込みである．この主張は内在的$\mathcal H_2$への
spectral双対化射の全忠実性を意味しない．
\end{theorem}
\begin{proof}
\paragraph{二形式からの構成.}
$(f_1,f_2)\in\overline P_{\rm nod}(S)$について
\[
 A=\OO_B\oplus\OO_B(-2)y,\quad y^2=f_1f_2,
 \qquad V=\OO_B(-1)c\oplus\OO_B(-1)d
\]
と置き，$y$の$V$上の作用と対称積を
\[
 Y=\begin{pmatrix}0&f_2\\f_1&0\end{pmatrix},\qquad
 \mu(c,c)=f_1,\quad\mu(d,d)=f_2,\quad\mu(c,d)=y
\]
とする．これは従来の同時符号の反不変部分と同じ構成である．
$Y^2=f_1f_2I$，$\operatorname{Tr}Y=0$だから，
補題~\ref{lem:qa-tf-trace-zero}により対応する$C$上の層は
全幾何fiberで各成分上rank oneかつtorsion-freeである．
$V$が$B$上局所自由なので底上平坦であり，Euler特性は零である．
積のtrace-zero成分は行列
$\left(\begin{smallmatrix}0&1\\1&0\end{smallmatrix}\right)$を持つので，
補題~\ref{lem:qa-tf-frobenius-duality}により$\mu^\flat$は同型である．
従ってこれは$\mathfrak T_{\rm bal}$の対象となる．
二形式のsquarefree性はこの構成には必要ない．

\paragraph{逆構成と全忠実性.}
補題~\ref{lem:qa-tf-balanced-isotropic-descent}により等方線の順序は
底上の有限\'etale二重被覆で選べる．
命題~\ref{prop:qa-tf-balanced-quadratic-inverse}は，任意の非被約底でも
$Y=\left(\begin{smallmatrix}a&f_2\\f_1&a\end{smallmatrix}\right)$の
trace零から$a=0$，$h=f_1f_2$を与え，二形式と全積を回収する．
もとの$C$が幾何的に被約節点曲線だから，
$f_1f_2$の根の重複度は高々二で，局所表示は$\overline P_{\rm nod}$に入る．
同命題のsimilitude行列の冪等元分解は，各同型が
対角または反対角の二生成元変換と底の$\operatorname{GL}_2$変換に由来し，
不定性が正確に共通次数スカラー$D$であることを示す．
この一意な回収はfppf局所で自然なので降下し，Isom sheafの同型を与える．
従って局所再構成と合わせてstack同値を得る．

全操作は有限平坦直像，trace，二次積，等方線とその平方sectionで与えられ，
任意の基底変更と可換する．
平方sectionのsquarefree性は相対二次形式の判別式が単元という開条件であり，
二等方線を交換しても不変である．従って$P$に対応する開は
正確に$\mathfrak T_{\rm adm}$である．
\end{proof}

\begin{corollary}[smoothなbalanced標的とH2開の判別式]
\label{cor:qa-tf-h2-admissible-open}
$\mathfrak T_{\rm bal}$はsmooth既約・有限型の純次元一のArtin stackである．
その内在的に定めた開$\mathfrak T_{\rm adm}$は，中央二重平方表示族が
H2を与える領域と正確に一致する．内在的相対H2族との解釈は
論文Bの相対表示同値が適用される局所有限型siteで行う．
\end{corollary}
\begin{proof}
$\overline P_{\rm nod}$は六次元アフィン空間の非空開で，
$G$は次元五のsmooth群なので，前定理の商表示から最初の主張を得る．
二形式のH2条件は系~\ref{cor:qa-double-square-degenerate-not-h2}により
$\Delta_1\Delta_2\ne0$と同値であり，これは$\mathfrak T_{\rm adm}$の定義である．
相対解釈には論文Bの
\cite[\texttt{thm:noncubic-intrinsic-presentation-equivalence}]{QuadricModuli2026}
を用いる．
\end{proof}

\begin{corollary}[順序付けと表示射の幾何fiber]
\label{cor:qa-tf-ordered-double-cover}
$G^{\rm ord}=(\operatorname{GL}(W)\times\Gm^2)/D$とする．
等方線の順序を保持したbalanced標的には
\[
 [\overline P_{\rm nod}/G^{\rm ord}]
       \simeq\mathfrak T_{\rm bal}^{\rm ord}
\]
という同値があり，順序忘却はrepresentable有限\'etale次数二である．
元のnoncubic開にも同じ主張が成立する．
従って二次積付き層とpencilを全て保持した標的に対し，
非順序付き表示射の像上の幾何2-fiberは一点，
順序付き表示射の幾何2-fiberは二点である．
これは$B\mu_2$というfiberではない．
\end{corollary}
\begin{proof}
二つの等方線は，底上の有限\'etale二重被覆として回収される．
その順序を選ぶことはこの被覆を自明化して一方のsheetを選ぶことと同じである．
前定理の全忠実性を順序を保つ射へ制限すると主張が従う．
\end{proof}


\begin{corollary}[全noncubic境界での元の符号被覆の回収]
\label{cor:qa-tf-original-point-cover-reconstruction}
二つの$f_i$が各fiberでsquarefreeなら，共有根・一致根対を含め，
$(C\to B,M,\mu)$から元の点曲線
$E_f=\{c^2=f_1,d^2=f_2\}$と同時符号$\delta$を任意の基底変更と可換に回収できる．
$\delta$は$\OO_C\oplus M$の二項に$+1,-1$として作用する．
その不変部分は$\OO_C$であり，固定点閉部分は
\[
 \operatorname{Spec}_C\bigl(\OO_C/\operatorname{im}\mu\bigr)
 =V_C(f_1,f_2,y).
\]
従って共有根で$\mu$が全射でなくても，符号被覆の代数は完全に回復される．
この被覆は$B$上有限平坦rank fourであるが，共有根のあるfiberでは
$C$上の有限平坦二重被覆ではない．
\end{corollary}
\begin{proof}
前命題を普遍二形式の表示に適用する．
同時符号の固定点idealは反不変部分$M$で生成され，
その偶数部分には$\mu(M,M)$が入るため，固定点環は記載の商である．
$C$上平坦なら$\OO_C\oplus M$はrank two vector bundleであり，
その直和因子$M$はline bundleとなる．共有根の節点で$M$は非局所自由なので
この平坦性は成立しない．
\end{proof}

\subsection{二つのnoncubic境界における慣性と忘却の核}
\label{subsec:qa-tf-boundary-inertia}

以下では$k$を標数零の代数閉体とする．
$\delta$は$c,d$の同時符号，$e_1,e_2$は個別符号とする．
一致型では$\tau$を$c,d$の交換とする．
各群は自然同型類をさらに割った群ではなく，
曲線・層・二次積の実際の自己同型群スキームである．

\begin{proposition}[全二次データの慣性と曲線への射]
\label{prop:qa-tf-boundary-inertia-maps}
元のnoncubic表示の二つの境界では，二次積付きTF標的の全慣性は
入力の慣性と同型である．具体的に，一根共有のnodal型では
\[
 H_{\rm nod}\simeq D_8,
 \qquad H_{\rm nod}^{\rm ord}\simeq\mu_2^2,
\]
一致する$I_2$型では，$O(uv)=\Gm\rtimes C_2$として
\[
 H_{\rm eq}\simeq(O(uv)\times D_8)/\langle(-I,\delta)\rangle,
 \qquad
 H_{\rm eq}^{\rm ord}
   \simeq(O(uv)\times\mu_2^2)/\langle(-I,\delta)\rangle.
\]
従って表示射自身の慣性準同型の核は，両境界で自明である．

一方，二次積付き層を忘れて被覆$(C\xrightarrow p B)$だけを残すと，
非順序付き入力の慣性準同型は全射となり，
\[
 \begin{array}{c|c|c}
 \text{型}&\text{核}&\operatorname{Aut}(C\xrightarrow pB)\\\hline
 \text{既約nodal}&\langle\delta\rangle\simeq\mu_2&C_2\times C_2\\
 I_2&\langle\delta,\tau\rangle\simeq\mu_2^2
          &(\Gm\rtimes C_2)\times C_2
 \end{array}
\]
となる．順序付き入力では核は両型とも$\langle\delta\rangle$である．
nodal型の順序付き像は被覆のdeck involutionからなる$C_2$で指数二，
$I_2$型の順序付き像は標的の全群である．
\end{proposition}
\begin{proof}
全TF慣性については前定理の全忠実性と，論文Bの全慣性計算
\cite[\texttt{prop:noncubic-all-stabilizers}]{QuadricModuli2026}を使う．
順序付き部分群は$c,d$の交換を行わない部分群である．

被覆を固定する核は
$\operatorname{Aut}_C(M,\mu)$であり，
系~\ref{cor:qa-spectral-pair-finite-inertia}から
nodal型では$\mu_2$，$I_2$型では$\mu_2^2$である．
nodal型では$-1_M$が$\delta$である．
$I_2$型の正規化成分上の生成元$c+d,c-d$に対し，
$\delta$は両方の符号，$\tau$は後者だけの符号を反転する．
従って核は記載した部分群であり，順序を保持すると$\tau$は除かれる．

nodal型では$f_1=v(u-v)$，$f_2=v(u+v)$と規格化できる．
被覆は$y^2=v^2(u^2-v^2)$であり，底上では唯一の二重根を固定し，
残りの二根を交換してよい．従って底の群は$C_2$，
そのliftと$y\mapsto-y$は可換な対合だから，被覆の群は$C_2^2$である．
個別符号はdeck involutionを与える．
$i^2=-1$を選ぶと
\[
 (u,v,c,d)\longmapsto(-u,v,id,-ic)
\]
は底の残りの対合を持ち上げる．よって非順序付きの像は全群である．
順序付きの底自己同型は三つの相異なる根を個別に固定するので恒等であり，
像はdeck群だけとなる．

$I_2$型では$f_1=f_2=uv$とできる．
二成分$y=uv,y=-uv$は各々$B$に同型であり，被覆の自己同型は
二節点を保つ底の自己同型と，独立な成分交換からなる．
従って$\operatorname{Aut}(C\to B)=(\Gm\rtimes C_2)\times C_2$である．
$H_{\rm eq}$からの射は，$O(uv)$の射影化と
\[
 D_8\longrightarrow C_2,\qquad e_1,e_2\longmapsto-1,
       \quad\tau\longmapsto1
\]
の積から誘導される．これらはともにfppf全射であり，
$D_8$を$\mu_2^2$へ制限しても後者は全射である．
従って$I_2$では順序付きの場合も全群を得る．
\end{proof}

\begin{remark}[smooth二次線束gerbeとの相違]
smooth域のspectral双対化で現れた$\mu_2$核は，
実際の二次線束$(M,\mu)$からそのPicard類へ移り，
層の$-1$自己同型を忘れる段階に由来する．
本節の標的は$(M,\mu)$を実際の層として保持するので，
$\delta=-1_M$を消さない．
境界で単に$\operatorname{Aut}_C(M,\mu)$を計算して得る
$\mu_2$または$\mu_2^2$は，表示射の核ではなく，
\emph{そこからさらに層を忘れる射の核}である．
\end{remark}

\begin{remark}[一般ribbonを加える際に残る仮定]
\label{rem:qa-tf-associativity-ribbon-warning}
上のtraceの議論は，分岐四次形式が各fiberで非零であることを使う．
この条件とgeneric rank one条件を落とすと，自己双対な対称積だけでは結合律は従わない．
例えば$R=k[\epsilon]/(\epsilon^2)$，$A=R[y]/(y^2)$，
$M=Rc\oplus Rd$に$y\cdot c=\epsilon c$，$y\cdot d=\epsilon d$と作用させ，
\[
 \mu(c,c)=\mu(d,d)=0,\qquad\mu(c,d)=\epsilon+y
\]
と置く．これは対称かつ$A$-双線型で，
$\operatorname{Ann}_A(y-\epsilon)=R(y+\epsilon)$から自己双対である．
しかし
\[
 (c\cdot c)\cdot d=0,
 \qquad c\cdot(c\cdot d)=2\epsilon c\ne0.
\]
ここで特殊fiberの$M$はribbonのgeneric点で自由rank oneでない．
従ってこれは本節のnodal TFの反例ではなく，より広い非被約曲線を含める場合に
rankの定義と結合律条件を省略できないことを示す例である．
\end{remark}

\subsection{Fricke出力の平坦点曲線との比較}

\begin{proposition}[二つのtraitでのgeneric同型と特殊曲線の相違]
\label{prop:qa-tf-spectral-flat-model-point-curve}
$k$を標数零の代数閉体とし，$r=1,2$について
\[
 p_1=1,\quad p_2=t,\qquad b_1=b_2=t,
 \qquad
 C_r=\{y^2=uv(p_ru^2+uv+b_rv^2)\}
\]
と置く．平坦出力代数$\mathcal B^{(r)}$の可換点曲線は
\[
 E'_r=\{c^2=U^2-4t^rV^2,\quad d^2=V(U+V)\}\subset\PP^3.
\]
そのgeneric fiberと$C_r$の間の同型は，有理式
\[
 \Phi_r:[u:v:y]\dashrightarrow
 [p_ru^2+b_rv^2:uv:p_ru^2-b_rv^2:y]
\]
によって与えられる．
しかしこの指定generic同型は$t=0$を越える曲線族の同型には延長できない．
これは有限trait基底変更後にも成立する．
特に$r=2$では，別のgeneric自己同型を合成した全ての同型について
延長不能を主張しているのではない．
\end{proposition}
\begin{proof}
四座標は同じ重み二であり，
\[
 (p_ru^2-b_rv^2)^2=(p_ru^2+b_rv^2)^2-4p_rb_r(uv)^2,
 \qquad y^2=(uv)(p_ru^2+b_rv^2+uv)
\]
なので，$p_rb_r=t^r$から像は$E'_r$に入る．
$t\ne0$，$uv\ne0$の開では
$x=u/v$を$(U+c)/(2p_rV)$で回復できるのでbirationalである．
generic fiberは双方smooth proper genus oneだから，このbirational mapは同型へ延びる．

共通の特殊fiber$E'_0$は$c=U$と$c=-U$の二つのsmooth conic
$d^2=V(U+V)$からなり，$[0:1:0:1]$と$[0:1:0:-1]$で横断的に交わる．
従って$I_2$型である．
一方，$C_{1,0}=\{y^2=u^2v(u+v)\}$は既約nodal曲線なので，
$r=1$では特殊fiber自体が同型でない．

$r=2$では$C_{2,0}=\{y^2=u^2v^2\}$も$I_2$型だが，
$uv\ne0$上の$\Phi_2$は
\[
                     [u:v:y]\longmapsto[0:1:0:\pm1]
\]
となり，各成分の稠密開を一点へ収縮する．
もしgeneric同型を延長する族同型が存在すれば，分離性により
$\Phi_2$が定義される開上でこの式と一致するため矛盾する．
有限trait基底変更は特殊fiberの型やこの式の収縮を変えない．
\end{proof}

\begin{remark}[再構成と境界双対化の区別]
二次積付きTFデータの代数$\OO_C\oplus M$は，入力の元の点曲線と符号を復元する．
Fricke出力の点曲線はsmooth域で$C$と同型になるが，前命題の固定モデルでは
その同型が境界を越えない．従って元の符号被覆の再構成だけから，
Fricke射のcompactificationが自動的に定まるわけではない．
節点で$M$が非局所自由なら$[M]$は通常の$\Pic^0(C)$の点ではなく，
smooth構成に用いた$t_{[M]}$もそのまま定義できない．
次の課題は，必要な曲線修正とpointingを指定し，二つのtraitで異なる収縮を
同型・慣性・基底変更と両立させる境界対応を構成することである．
\end{remark}

% !TeX root = ../nc-quadric-avatar.tex
% 日本語 LuaLaTeX。2026-11-23 / 二境界のgraph・二次線束・四捩れ補正。
\section{Fricke境界の共通半安定模型と二次線束}
\label{sec:qa-spectral-boundary-models}

命題~\ref{prop:qa-tf-spectral-flat-model-point-curve}の指定generic同型は
元の二模型の同型へ延長しない．本節ではそのgraph閉包を実際に構成し，
符号被覆と二次積を共通正則模型へ延長する．
さらに根対一致型には，四捩れtranslationで補正した別の模型同型が存在する．
曲線の構造層を保つ比較，conicの順序，H2の構造対象は区別する．

\subsection{指定した一般同型のgraph閉包と共通半安定模型}
\label{subsec:qa-boundary-graph-semistable}

$k$を標数零の代数閉体，$R=k[[t]]$，$K=k((t))$とする．
$r=1,2$に対して
\[
 (p_1,b_1)=(1,t),\qquad(p_2,b_2)=(t,t),\qquad p_rb_r=t^r
\]
と置く．曲線模型を
\[
 C_r:\ y^2=uv(p_ru^2+uv+b_rv^2)
       \subset\mathbb P_R(1,1,2),
\]
\[
 E'_r:\ c^2=U^2-4t^rV^2,\quad d^2=V(U+V)
       \subset\mathbb P_R^3
\]
で定める．一般fiber上の同型は
\[
 \varphi_r:[u:v:y]\longmapsto
 [p_ru^2+b_rv^2:uv:p_ru^2-b_rv^2:y].
 \tag{G-phi}\label{eq:qa-boundary-graph-phi}
\]
$A=(U+c)/2$，$B=(U-c)/2$とすると
\[
 E'_r:\ AB=t^rV^2,\quad d^2=V(A+B+V),\qquad
 \varphi_r=[p_ru^2:b_rv^2:uv:y]
 \tag{G-AB}\label{eq:qa-boundary-graph-ab}
\]
となる．以下でいうgraph閉包は，この\emph{指定された}一般同型の
scheme-theoreticな閉包である．一般同型を変更した比較とは区別する．

\begin{theorem}[二traitのgraph閉包]
\label{thm:qa-boundary-graph-common-regular-model}
$C_r\times_R E'_r$内の$\varphi_{r,K}$のgraph閉包を$X_r$とする．
このとき
\[
 X_r\simeq\operatorname{Bl}_{\mathcal I_r}C_r,
 \qquad
 \mathcal I_r=(p_ru^2,b_rv^2,uv,y)\otimes\mathcal O_{C_r}(-2)
\]
であり，$X_r$は既に正則な射影的半安定$R$-模型である．
追加の正規化または底の分岐被覆は必要ない．
二つの射を$a_r:X_r\to C_r$，$b_r:X_r\to E'_r$と書く．
\begin{enumerate}
 \item $C_1$の全空間は$[u:v:y]=[0:1:0]$に一つの$A_1$特異点を持つ．
 $a_1$はその極大idealのblowupによる最小解消で，$X_{1,0}$は
 二つの有理曲線が二点で横断的に交わる$I_2$型である．
 $E'_1$の全空間は正則で，$b_1$は同型である．
 \item $C_2$の全空間は$[0:1:0]$，$[1:0:0]$に二つの$A_1$特異点を持つ．
 $a_2$は両点の極大idealをblowupする最小解消であり，
 $X_{2,0}$は四つの有理曲線からなる$I_4$型である．
 $D_+,D_-$を$C_{2,0}$の成分$y=\pm uv$のstrict transform，
 $L_0,L_\infty$を$a_2$の例外曲線とすると，
 各$D_\pm$は各$L_0,L_\infty$と一度ずつ交わり，同種の二成分は交わらない．
 $E'_2$は$P_\pm=[A:B:V:d]=[0:0:1:\pm1]$に二つの$A_1$特異点を持つ．
 $b_2$は$D_\pm$をそれぞれ$P_\pm$へ収縮する最小解消であり，
 $L_0,L_\infty$は$E'_{2,0}$の二成分へ同型に写る．
\end{enumerate}
従って$X_r$は共通の最小正則模型であり，特殊fiberは$I_{2r}$型である．
\end{theorem}

\begin{proof}
四つのsectionから得るRees代数の次数一生成系は
$\operatorname{Bl}_{\mathcal I_r}C_r\hookrightarrow C_r\times_R\mathbb P^3_R$
を与える．blowupは整で一般fiberは元のgraphであるから，その像は
scheme-theoreticなgraph閉包である．二つの方程式
\eqref{eq:qa-boundary-graph-ab}はこの閉包上で成立するので，像は
$C_r\times_R E'_r$内にある．この議論はRees代数によるblowupの定義\cite[Lemmas 31.33.2--31.33.3]{StacksBlowups}から従い，
単に二項の比例式だけを課して余分な成分を加えたincidence schemeではない．

$v=1$のchartで$x=u$と置くと，$C_r$は
\[
 y^2=x^2+tx+p_rx^3
\]
である．$N_0=(x,y,t)=(0,0,0)$における$\mathcal I_r$は
$(x,y,t)$に等しい．$z_0=t+x+p_rx^2$と置けば
$y^2=xz_0$であり，$\partial z_0/\partial t$は$r=1$で$1$，
$r=2$で$1+x^2$である．従って全空間の完備局所環は$A_1$特異点
$k[[x,y,z_0]]/(y^2-xz_0)$である．
$r=2$の$N_\infty$でも$u,v$を交換して同じ計算ができる．
その他の特殊fiber上の点ではJacobian判定から全空間は正則である．

極大idealのblowupの$x$-chartで
\[
 y=xz,\qquad t=xs
\]
と書く．strict transformの方程式と$E'_r$への射は，それぞれ
\[
 \begin{array}{c|c|c}
 r&\text{方程式}&[A:B:V:d]\\\hline
 1&z^2=1+x+s&[x:s:1:z]\\
 2&z^2=1+s(1+x^2)&[sx^2:s:1:z]
 \end{array}
 \tag{G-chart}\label{eq:qa-boundary-graph-chart}
\]
である．例外曲線は$x=0$，元の特殊fiberのstrict transformは$s=0$であり，
交点は$z=\pm1$となる．その近傍では$x,s$が正則パラメータで，
基底のパラメータは$t=xs$である．
例外曲線上でこのchartから外れる点は$t$-chart
$x=t\alpha$，$y=t\beta$で覆われ，方程式は
\[
 \beta^2=\alpha^2+\alpha+p_rt\alpha^3
\]
となる．$(t,\alpha,\beta)=(0,0,0)$では$\alpha$方向の偏微分が$-1$なので正則である．
これら二chartは例外曲線全体を覆う：例外conicの式
$y^2=x^2+tx$で$x=t=0$なら$y=0$となり，射影点を与えないためである．
従って全blowupは正則で，特殊fiberは被約な単純正規交差曲線である．

$r=1$では$C_{1,0}$の正規化は有理曲線であり，その節点の二枝に
例外$\mathbb P^1$が接続するので$I_2$を得る．
式\eqref{eq:qa-boundary-graph-chart}により，二成分は$E'_{1,0}$の二成分へ
それぞれ次数一で写る．$E'_1$の二節点の全空間は$AB=t$で正則であり，
それ以外も正則である．従って$b_1$はproperでquasi-finiteなbirational射で，
正規な標的への有限birational射だから同型である．

$r=2$では二つの$A_1$点の解消が，各々$D_+,D_-$の間に一つの
例外曲線を挿入するので，記載した四成分cycleを得る．
$E'_r$の$P_\pm$近傍では$V=1$とでき，$2d$が単元なので，
完備局所環は$AB=t^r$に一致する．特に$r=2$では$A_1$である．
式\eqref{eq:qa-boundary-graph-chart}の第二行は
\[
 B=s,\qquad t=xs,\qquad A=x^2s
\]
という$AB=t^2$の標準解消chartそのものである．
$N_\infty$のchartは$A,B$を交換したchartとなる．
従って$b_2$は正に$D_+,D_-$を収縮し，例外成分$L_0,L_\infty$は
それぞれ$A=0,B=0$のconicへ同型に写る．

全特殊fiberが被約であり，各成分は他成分と合計二回横断的に交わる．
fiberとの交点数が零であることから各成分の自己交点は$-2$である．
従って相対的な$(-1)$曲線はなく，これらは最小正則模型である．
\end{proof}

\begin{corollary}[符号作用とその半安定商]
\label{cor:qa-boundary-graph-sign-quotient}
$E'_r$上の同時符号作用は
\[
 \delta':[A:B:V:d]\longmapsto[B:A:V:-d]
\]
である．これは$X_r$へ持ち上がり，全空間上固定点を持たない．
特殊fiberの$I_{2r}$-cycleでは半回転として作用する．
従って
\[
 Y_r=X_r/\langle\delta'\rangle
\]
は正則な射影的半安定$R$-曲線で，特殊fiberは$I_r$型である．
$C'_r=E'_r/\langle\delta'\rangle$と書くと，$Y_r\to C'_r$は
$r=1$で同型，$r=2$で唯一の$A_1$点の最小解消となる．
\end{corollary}
\begin{proof}
$r=1$では$X_1\simeq E'_1$である．$r=2$では
\[
 [u:v:y]\longmapsto[v:u:-y]
\]
が$C_2$上の自己同型であり，二つのblowup中心を交換するので
$X_2$へ持ち上がる．一般fiberで$\varphi_2$により$\delta'$と一致し，
従って全体でも$b_2$は同変である．作用は$D_+,D_-$を交換し，
$L_0,L_\infty$を交換するので四成分cycleの半回転である．

$E'_r$で固定点があると仮定する．$V\ne0$なら$A=B$，$d=0$となり，
$d^2=V(2A+V)$より$A=-V/2$，従って$4t^r=1$となる．
これは$R$上では不可能である．$V=0$の場合も$AB=0,d^2=0$を用いると
射影座標が全て零となるため不可能である．従って$\delta'$は自由で，
その同変な$X_r$上の作用も自由である．
射影的scheme上の有限自由作用の商はschemeとして存在し\cite[Lemma 66.14.2]{StacksFreeQuotients}，$X_r\to Y_r$は有限\'{e}tale次数二である．
正則性は\'{e}tale局所的なので$Y_r$は正則である．
成分と節点が半回転で二つずつ同一視されるので特殊fiberは$I_r$となる．
特に$r=1$では一つの有理成分に一つの自己節点ができる．
$r=2$では$b_2$が収縮する$D_+,D_-$が商で一つの$(-2)$曲線となり，
それを$C'_2$の唯一の節点へ収縮する．
商曲線の式
\[
 C'_r:\ Y^2=(U^2-4t^rV^2)V(U+V)
\]
からも，その節点の全空間が$r=1$で正則，$r=2$で$A_1$であると確認できる．
\end{proof}

\begin{proposition}[構造層と完全複体の比較]
\label{prop:qa-boundary-graph-structure-sheaf}
二つの射影$a_r,b_r$について
\[
 \mathbf R a_{r*}\mathcal O_{X_r}\simeq\mathcal O_{C_r},
 \qquad
 \mathbf R b_{r*}\mathcal O_{X_r}\simeq\mathcal O_{E'_r}.
\]
従って$\mathbf L a_r^*$と$\mathbf L b_r^*$は各曲線模型の
$\operatorname{Perf}$から$\operatorname{Perf}(X_r)$への全忠実関手であり，
構造層を$\mathcal O_{X_r}$へ送る．
この命題は曲線模型の完全複体についての比較であり，
特異なAS heartのH2所属またはその間の同値を主張するものではない．
\end{proposition}
\begin{proof}
射は正規な標的の孤立$A_1$点の解消，または同型である．
正規性から零次順像は構造層である．
一つの例外曲線を$L\simeq\mathbb P^1$とすると，
$L^2=-2$より，その$n$次無限小近傍の構造層は
$\mathcal O_{\mathbb P^1}(2j)$，$0\le j<n$を逐次商に持つ．
これらの$H^1$は零であり，$A_1$の極大idealの逆像は$\mathcal O(-L)$だから，
形式関数定理\cite[Theorem 30.20.5]{StacksFormalFunctions}より一階高次順像も零となる．
二階以上も同じ近傍のcohomology消滅から零となる．
最後の主張は完全複体に対するprojection formulaから従う．
\end{proof}

\begin{proposition}[有限分岐基底変更]
\label{prop:qa-boundary-graph-ramified-base-change}
$t=s^e$で$R'=k[[s]]$へ基底変更する．
指定した一般同型のgraph閉包は$X_r\times_RR'$であり，
その特異点は元の各節点に現れる$xy=s^e$型だけである．
それぞれを最小解消すると$I_{2re}$型の最小正則模型を得る．
符号作用は自由な半回転として持ち上がり，商は$I_{re}$型となる．
\end{proposition}
\begin{proof}
有限DVR拡大$R\to R'$は平坦であり，Rees代数とblowupはこの基底変更と可換である．
元の正則半安定模型の節点では$t=xy$なので，基底変更後は$xy=s^e$となる．
この特異点の解消は$j=0,\ldots,e-1$の正則chart
\[
 s=\alpha\beta,\qquad
 x=\alpha^{j+1}\beta^j,\qquad
 y=\alpha^{e-j-1}\beta^{e-j}
\]
で記述できる．隣接chartの貼合せは
$\alpha'=\beta^{-1}$，$\beta'=\alpha\beta^2$である．
従って各辺に$e-1$本の有理成分が挿入され，cycleの長さは$e$倍となる．
符号作用は元の節点を二つずつ交換し，対応する解消chainを交換するため自由である．
\end{proof}

\begin{remark}[基底変更の範囲]
ここで確認したのはflat base changeと有限DVR拡大に対する挙動である．
graph閉包や最小正則解消が任意の非平坦・非被約基底変更と可換であるとは主張しない．
また，$r=2$では別途構成する補正同型$\Psi$が元の二つの模型の$R$同型を与える．
指定した$\varphi_2$のgraphが非自明であることは，この別同型の存在と両立する．
\end{remark}

\subsection{二つのtraitにおけるFricke二次線束の延長}

記号は前節を引き継ぎ，$E'_r$上の個別符号を$e_c,e_d$，
同時符号を$\delta'=e_ce_d$とする．
入力の元の点曲線$E_{f_r}$と，出力の点曲線$E'_r$は区別する．

\begin{proposition}[二捩れ平行移動と順序付きconic]
\label{prop:qa-fricke-trait-sign-conjugacy}
指定したgeneric同型
\[
 \Phi_r:[u:v:y]\longmapsto
 [p_ru^2+b_rv^2:uv:p_ru^2-b_rv^2:y]
\]
による$\delta'$の共役は
\[
 T_r:[u:v:y]\longmapsto[b_rv:p_ru:-p_rb_ry].
 \tag{F-sign}\label{eq:qa-fricke-trait-sign}
\]
これは入力の二次線束$[M]$による平行移動$t_{[M]}$である．
入力のhyperelliptic involutionを$\iota:y\mapsto-y$とすると，
\[
 \Phi_r\iota=e_d\Phi_r,\qquad
 \Phi_r(T_r\iota)=e_c\Phi_r.
\]
従って$E'_r$の\emph{順序付き}二conic
\[
 Q'_1=E'_r/\langle e_d\rangle=\{c^2=U^2-4t^rV^2\},\qquad
 Q'_2=E'_r/\langle e_c\rangle=\{d^2=V(U+V)\}
\]
は，generic fiber上で自然な射$\mathscr D^{\rm ord}$の
$C_r/\langle\iota\rangle$，$C_r/\langle t_{[M]}\iota\rangle$
にこの順序で一致する．
\end{proposition}
\begin{proof}
$v=1$上で$x=u/v$と書く．$C_{r,K}$の原点を無限遠点$O$とし，
$P=(0,0)$とすると，$\operatorname{div}(x)=2P-2O$である．
入力の符号被覆$E_{f_r,K}\to C_{r,K}$は$x$の平方根を添加して得られるので，
その反不変線束の類は$\OO(P-O)$である．
曲線$y^2=x(p_rx^2+x+b_r)$上で$P$を通る直線の交点を計算すると，
\[
 t_P(x,y)=\left(\frac{b_r}{p_rx},
                   -\frac{b_ry}{p_rx^2}\right).
\]
これは\eqref{eq:qa-fricke-trait-sign}のaffine表示である．
同次式へ$T_r$を代入すると$U,V,c,d$は
$p_rb_r(U,V,-c,-d)$となる．従って$\Phi_rT_r=\delta'\Phi_r$である．
$\iota$は$\Phi_r$の$d$だけを反転するので，二つの個別符号に関する式も従う．
これらの商は単に同型なconicというだけでなく，
同じ$Z'=C_r/\langle\iota,T_r\rangle$への被覆として一致する．
\end{proof}

\begin{theorem}[二つのtrait上でのbalanced標的への延長]
\label{thm:qa-fricke-trait-balanced-extension}
\[
 g_{1,r}=U^2-4t^rV^2,\qquad g_{2,r}=V(U+V),\qquad
 C'_r=\{Y^2=g_{1,r}g_{2,r}\}
\]
と置く．順序付き二形式$(g_{1,r},g_{2,r})$は，
$\mathscr D^{\rm ord}(f_{r,K})$の
$\mathfrak T_{\rm bal}^{\rm ord}$への$R$上の延長を与える．
この延長に付随する反不変層と二次積は
\[
 q'_*\OO_{E'_r}=\OO_{C'_r}\oplus M'_r,
 \qquad \mu'_r:(M'_r)^{\otimes2}\xrightarrow{\sim}\OO_{C'_r}
\]
であり，$M'_r$は\emph{全族上の線束}，
$q':E'_r\to C'_r$は有限\'etale二重被覆である．
この代数，二次線束，被覆，trace分解の構成は，
非平坦・非被約なものを含む任意の$R$-schemeへの基底変更と可換である．

特殊fiberでは$C'_0$は既約nodal曲線，$E'_0$は$I_2$曲線であり，
$M'_0$は非自明な二捩れ線束である．
従って固定被覆上のbalanced fiberにおける追加の被約$B\mu_2$成分に属し，
元のnoncubic開$\mathfrak N^{\rm pres}$には属さない．
特にこの延長はH2への境界延長を与えない．
\end{theorem}
\begin{proof}
$g_{1,r}g_{2,r}$は各幾何fiberで非零であり，
generic fiberでは四根が相異なり，特殊fiberでは一つの二重根と二つの単根を持つ．
従ってこの二形式は$\overline P_{\rm nod}$上の$R$-点である．
balanced商表示の認識により，実際の二次積付きTF対象が得られる．
前命題はgeneric fiber上で自然な\emph{順序付き}Fricke射との同型を指定している．

線束性と基底変更は直接確認できる．
$Z'=\PP^1_R$上で$A=\OO_{Z'}\oplus\OO_{Z'}(-2)Y$，$Y^2=g_{1,r}g_{2,r}$と書く．
反不変部分の$Z'$上の直像は
\[
 p'_*M'_r=\OO_{Z'}(-1)c\oplus\OO_{Z'}(-1)d,
 \quad Yc=g_{1,r}d,\quad Yd=g_{2,r}c,
\]
であり，二次積は
\[
 c^2=g_{1,r},\qquad d^2=g_{2,r},\qquad cd=Y
\]
である．$g_{1,r}$と$g_{2,r}$は共通零点を持たない．
実際，$g_{2,r}=0$の二つの点$V=0$，$U+V=0$における$g_{1,r}$の値は，
局所的な線束の自明化を除き，それぞれ$1$と$1-4t^r$である．
これらは$R$の単元であり，任意の基底変更後にも単元である．

$g_{1,r}$が可逆な開では$M'_r$は$c$で生成され，
$d=(Y/g_{1,r})c$，$c^2=g_{1,r}$となる．
$g_{2,r}$が可逆な開では同様に$d$で生成される．
従って$M'_r$は線束で，$\mu'_r$は同型である．
$2$が可逆なので$\operatorname{Spec}_{C'_r}(\OO\oplus M'_r)$は
有限\'etale二重被覆であり，局所式からこれは$E'_r$と一致する．
全ての式は局所自由加群の代数として与えられており，
固有部分への射影$(1\pm\delta')/2$も基底変更と可換する．
よって主張した任意基底変更との両立が得られる．

$p'_*M'_r=\OO(-1)^{\oplus2}$からbalanced条件と$\chi(M'_s)=0$が従う．
特殊fiberは
\[
 C'_0:\ Y^2=U^2V(U+V),\qquad
 E'_0:\ c=\pm U,\quad d^2=V(U+V)
\]
である．$E'_0$の二つのsmooth conicは
$[U:V:c:d]=[0:1:0:\pm1]$で交わり，連結な$I_2$曲線となる．
したがって$E'_0\to C'_0$は非自明な\'etale二重被覆であり，
$M'_0$は非自明な二捩れ線束である．
幾何的にも，$C'_0$の正規化$d^2=V(U+V)$の二節点逆像において，
被覆の二つの成分は交差して貼り合わされる．
これは$\Pic^0(C'_0)\simeq\Gm$の元$-1$に対応する．

最後に$g_{1,r}|_{t=0}=U^2$は重根を持ち，
$g_{2,r}|_{t=0}=V(U+V)$はsquarefreeである．
従って特殊対象はbalanced開には属するが，
両平方がsquarefreeである元のnoncubic開には属さない．
H2所属判定$\Delta_1\Delta_2\ne0$とも整合する．
\end{proof}

\begin{proposition}[共通正則モデル上の符号と二次線束]
\label{prop:qa-fricke-regular-quotient-line}
$r=1$では$X_1=E'_1$，$Y_1=C'_1$とする．
$r=2$では$Y_2\to C'_2$を唯一の特殊節点のblowupとし，
\[
 X_2=E'_2\times_{C'_2}Y_2
\]
と置く．このとき$X_r,Y_r$は正則な半安定曲線族で，
特殊fiberはそれぞれ$I_{2r}$，$I_r$型である．
$X_r\to E'_r$は$r=2$の場合には二つの特殊節点の同時blowupである．
同時符号$\delta'$は$X_r$に自由に作用し，
\[
 X_r/\langle\delta'\rangle=Y_r,
 \qquad X_r\to Y_r\text{ は有限\'etale二重被覆}
\]
となる．$\beta_r:Y_r\to C'_r$とすると，その二次線束は
\[
 \widetilde M_r=\beta_r^*M'_r,\qquad
 \widetilde\mu_r=\beta_r^*\mu'_r
\]
である．$X_r$を指定generic同型の共通正則モデルと同定したとき，
この作用は入力の$t_{[M]}$の延長となる．
\end{proposition}
\begin{proof}
$C'_r$の特殊節点の近傍で$V=1$とする．
$U+1$の平方根はその完備局所環の単元として取れるので，
$z=Y/\sqrt{U+1}$と書けば
\[
 (U-z)(U+z)=4t^r.
\]
従って$r=1$では全空間が正則，$r=2$では$A_1$特異点であり，
そのblowupは一つの例外$\PP^1$を挿入する正則半安定解消である．
同じ局所式は$E'_r$の二つの節点で，$z$を$c$として成立する．
$E'_r\to C'_r$が有限\'etaleであるため，節点のidealの引戻しは
二節点のidealであり，blowupの平坦基底変更から
$X_2$は二節点の同時blowupとなる．
各節点に一つの例外成分が挿入されるので，特殊fiberの型は記載の通りである．

被覆と二次線束の式は前定理の有限\'etale代数を$Y_r$へ引き戻したものである．
特に作用は自由であり，商は$Y_r$である．
$\delta'$は元の二つのconic成分と二つの節点をそれぞれ交換するので，
$I_{2r}$の双対グラフ上では半回転として作用する．
generic fiber上の作用との一致は
\eqref{eq:qa-fricke-trait-sign}と共通モデルのgeneric同型による．
\end{proof}

\begin{remark}[基底変更と標識の範囲]
ここで任意基底変更と可換としたのは，構成した有限\'etale被覆と
その二次線束，およびそれらの引戻しである．
graph閉包，正規化，最小正則モデルを作る操作そのものが，
任意の非平坦基底変更と可換すると主張してはいない．
また二つの指定traitへの$\mathscr D^{\rm ord}$の延長を構成したのであり，
全noncubicスタック上の大域的境界射や，H2 heartの延長を構成したわけではない．
\end{remark}


\subsection{一致型の四捩れ補正とconic順序}

\begin{proposition}[一致型の別同型はtrait全体へ延長する]
\label{prop:qa-coincident-alternative-model-isomorphism}
$a\in R^\times$を$a^2=1+2t,\ a\equiv1\pmod t$で定める．
次の射は$R$上の同型である：
\[
 \begin{split}
 \Psi:C_2&\xrightarrow{\sim}E'_2,\\
 [u:v:y]&\longmapsto
 \bigl[a^2(u+v)^2-(v-u)^2:(v-u)^2:4ay:a(u+v)(v-u)\bigr].
 \end{split}
\]
従って$C_2$と$E'_2$の間で延長しないのは，指定された同型$\phi_2$であり，
抽象的な$R$-曲線としての同型の存在は妨げられない．
一方$C_1$と$E'_1$は，特殊fiberの既約成分数が異なるので$R$上同型ではない．
\end{proposition}
\begin{proof}
中間の重み付きquarticを
\[
 D:\quad w^2=(s^2-z^2)^2-4t^2z^4
\]
とする．線形変換および重み二の単元倍
\[
 [s:z:w]=[a(u+v):v-u:4ay]
\]
は$R$上可逆である．また
\[
 \begin{split}
 &(a^2(u+v)^2-(1+2t)(v-u)^2)\\
 &\hspace{12mm}\cdot(a^2(u+v)^2-(1-2t)(v-u)^2)\\
 &=16a^2uv\bigl(tu^2+uv+tv^2\bigr)
 \end{split}
\]
であるから，これは$C_2\simeq D$を与える．
$D$の二次Veroneseは$s^2,sz,z^2,w$で生成され，
\[
 [s:z:w]\longmapsto[s^2-z^2:z^2:w:sz]
\]
により$E'_2$と同型になる．この二つを合成したものが$\Psi$である．

$t=0$で$C_2$の二成分$y=\pm uv$は，$\Psi$により
$E'_2$の二成分$c=\pm U$へそれぞれ同型に移る．
これは$\phi_2$が二成分の稠密開を二つの節点へ収縮することとも整合する．
\end{proof}

\begin{proposition}[四捩れtranslationとpencilの交換]
\label{prop:qa-coincident-quarter-period-correction}
一般fiber上で$F=\Psi^{-1}\phi_2$と置くと
\[
 F([u:v:y])=[y-auv:y+auv:at uv(u^2-v^2)]
\]
であり，
\[
 F^2([u:v:y])=[v:u:-y],\qquad F^4=\mathrm{id}.
\]
$C_{2,K}$に$[1:0:0]$を原点とする群構造を入れると，$F$は位数四の点
$Q=[1:1:a]$によるtranslationである．従って
\[
 \Psi=\phi_2\circ T_{-Q}.
\]

$\iota:[u:v:y]\mapsto[u:v:-y]$，
$j:[u:v:y]\mapsto[v:u:-y]$，
$e_c,e_d$を$E'_2$の$c,d$の個別符号とする．このとき
\[
 \phi_2\iota=e_d\phi_2,
 \qquad\Psi\iota=e_c\Psi,
 \qquad\Psi j=(e_ce_d)\Psi.
\]
従って$\Psi$は同時符号には同変であるが，二つのconic quotientの順序を交換する．
\end{proposition}
\begin{proof}
$D\dashrightarrow E'_2$の逆を一般fiberの開上で
$[s:z:w]=[d:V:cV]$と書けば，$F$の公式を得る．
$u_1=y-auv,v_1=y+auv,y_1=at uv(u^2-v^2)$とすると
\[
 u_1v_1=tuv(u-v)^2.
\]
$L=2at uv(u-v)$と置いて再度$F$を作用させると
\[
 u_2=Lv,\qquad v_2=Lu,\qquad y_2=-L^2y
\]
となるため，重み付き射影座標で記載の$F^2$を得る．

$X=t u/v,Y=t y/v^2$によりWeierstrass方程式
\[
 Y^2=X^3+X^2+t^2X
\]
を得る．$Q=(t,ta)$における接線の傾きは$a$であり，倍点公式から
$2Q=(0,0)$を得る．さらに通常の加法公式に
$m=(y/v^2-a)/(u/v-1)$を代入すると，$T_Q$は上記$F$に一致する．
従って四捩れtranslationという記述は自己同型群の分類に依存しない．
最後の三つの同変等式は，表示された座標式から直ちに従う．
\end{proof}

\begin{remark}[選択と主張の範囲]
$a$の符号を変えると$\Psi$は$(e_ce_d)\Psi$に，$F$は$F^{-1}$に変わる．
代数的なparameter空間で同じ式を使う場合は，$a^2=1+2t$という
有限\'etale被覆上で構成する．上の結果はこのtraitの曲線模型と，
指定された二つのquotient pencilの比較である．全moduli上の標準的な補正，
指定AS10の構造対象の保存，あるいはH2への境界延長を主張しない．
実際，$E'_2$の平方形式の一つは$t=0$で二重根を持ち，
既存の無限ExtによるH2延長障害はそのまま残る．
\end{remark}

% !TeX root = ../nc-quadric-avatar.tex
% 日本語 LuaLaTeX。2026-11-23 / v4.11。
\section{自然なFricke射の全noncubic境界への延長}
\label{sec:qa-global-fricke-boundary}

本節では二trait上の構成を，全表示groupoid上の射へ延長する．
標的には二次自己双対線束を持つ節点対象を加える．
一致点の全慣性による局所表示から始め，補助平方根の選択を変える
全同型を比較して，任意基底変更と両立する大域射を構成する．

\subsection{一致根対に沿うFricke射の全慣性付き形式延長}

以下では$k$を標数零の代数閉体とする．
入力の一致点を$(uv,uv)$，出力の線束境界点を
$\ell_0=(U^2,V(U+V))$で表す．
$H=H_{\rm eq}$は入力の全慣性，$\Gamma$は$\ell_0$の全慣性を表す．
入力の次数一生成元を$c_0,d_0$，二形式の変形parameterを$a,b$とし，
出力の次数一生成元を$C,D$とする．

\begin{proposition}[線束境界点の慣性と形式変形]
\label{prop:qa-fricke-line-formal-target}
$\ell_0$の全慣性は$\Gamma\simeq\mu_2^3$であり，
一次変形空間は一次元，$\Gamma$の接作用は自明である．
この点のresidual gerbeに沿うbalanced標的の形式完備化は
\[
 \widehat{\mathfrak T}_{\mathrm{bal},\ell_0}
 \simeq[\operatorname{Spf}k[[q]]/\Gamma]
\]
であり，$\Gamma$は$q$に自明に作用する．
第一形式が二重根を持つ順序付きの点を選んだ
$\mathfrak T_{\mathrm{bal}}^{\rm ord}$の形式完備化も同じ表示である．
ここではこの順序付きの持上げを用いる．普遍一変数族として
\[
 g_{1,q}=U^2-4qV^2,\qquad g_2=V(U+V)
\]
を取れる．この近傍では$1/\chi=4q$である．
\end{proposition}
\begin{proof}
$u=U,v=U+2V$と第二生成元の二倍によって，閉点を
$(u^2,v^2-u^2)$へ移す．第一形式の二重根と第二形式の二単根を保つ
$\PGL_2$部分群は$v\mapsto-v$による$C_2$だけである．
二つの形式の型が異なるため，形式交換は安定化群に入らない．
$v\mapsto-v$は二形式を正確に保存する位数二の行列で持ち上がり，
二つの生成元の独立な符号と可換する．従って全群は$\mu_2^3$となる．

行列の一次変分を$\left(\begin{smallmatrix}\alpha&\beta\\\gamma&\delta\end{smallmatrix}\right)$，
生成元のスカラー変分を$r,s$とすると，六係数への軌道微分は
\[
 \bigl(2(\alpha-r),2\beta,0,
       2(s-\alpha),2(\gamma-\beta),2(\delta-s)\bigr).
\]
その階数は五，核は共通次数スカラーだけである．
従って第一形式の$v^2$係数が横断parameterとなる．
上の三つの符号対称性はこの係数を保つので接作用は自明である．

形式的陰関数定理により，閉点で恒等に合同な群変換を使って
任意のArtin変形をこの一変数族へ正規化できる．
閉点の各慣性元には，下に与える族全体の明示的な持上げがある．
また$1/\chi=4q$はこの近傍で正則なので，二つの正規形が同型なら
nilpotent parameterも含めて$q$は等しい．
軌道微分の核を共通スカラーで割ると零になるため，
閉点の同型を指定した正規化の持上げは一意である．
よって全Artin同型を含む形式商表示が得られる．
\end{proof}

\begin{theorem}[自然なFricke射の形式延長]
\label{thm:qa-fricke-formal-equivariant-extension}
$A=k[[a,b]]$，$q=ab$とし，$A$内の一意な平方根
\[
 \rho^2=1-4ab,\quad\rho\equiv1,\qquad
 \ell^2=(1+\rho)/2,\quad\ell\equiv1
\]
を取る．一致根対の全慣性
\[
 H=\bigl((\Gm\rtimes C_2)\times D_8\bigr)
       /\langle(-1,e_ce_d)\rangle
\]
を用いた形式表示上で，$\mathscr D^{\rm ord}$は
\[
 \boxed{
 [\operatorname{Spf}A/H]\longrightarrow
 [\operatorname{Spf}k[[q]]/\Gamma],\qquad
 q=ab
 }
\]
へ延びる．慣性の準同型は
\[
 H\longrightarrow H/H^\circ\xrightarrow{\sim}\Gamma,
 \qquad H^\circ=\Gm
\]
である．この射は，全群作用と全Artin同型を含む形式射である．
平方根を補助変数とする下記の\'etale chartへ代数化すると，
その$ab\ne0$の開上で指定した自然な順序付きFricke射と一致する．
\end{theorem}
\begin{proof}
\paragraph{入力の全群作用.}
入力族を
\[
 c_0^2=uv,\qquad d_0^2=au^2+uv+bv^2
\]
とする．$H$の生成元$T_z,r,e_c,e_d,\tau$のparameterへの作用は
\[
 T_z(a,b)=(z^{-2}a,z^2b),\quad r(a,b)=(b,a),\quad
 e_c(a,b)=e_d(a,b)=(a,b),\quad\tau(a,b)=(-a,-b).
\]
生成元の持上げを次で定める：
\[
 \begin{array}{c|cccc}
 &u'&v'&c'_0&d'_0\\\hline
 T_z&zu&z^{-1}v&c_0&d_0\\
 r&v&u&c_0&d_0\\
 e_c&u&v&-c_0&d_0\\
 e_d&u&v&c_0&-d_0\\
 \tau&\ell u+(b/\ell)v&(a/\ell)u+\ell v&d_0&\rho c_0
 \end{array}
\]
全生成元は$q,\rho,\ell$を保つ．$\tau$の行列を$L(a,b)$とすると
\[
 \ell^2+ab/\ell^2=1,\qquad
 \ell^2-ab/\ell^2=\rho,
\]
より
\[
 u'v'=au^2+uv+bv^2,
 \qquad -a(u')^2+u'v'-b(v')^2=\rho^2uv
\]
が成り立つ．従って$\tau$はparameter$(-a,-b)$の入力族への同型である．

群関係も実際の同型として確認できる．
\[
 L(-a,-b)L(a,b)=\rho I_2
\]
であり，$\tau^2$は四つの次数一生成元全ての$\rho$倍となるので，
共通次数スカラーを割った群では恒等である．
$r^2=1$，$rT_zr=T_{z^{-1}}$，$\tau e_c\tau=e_d$が成立し，
$\tau$は$r,T_z$と可換する．例えば
\[
 L(z^{-2}a,z^2b)\operatorname{diag}(z,z^{-1})
 =\operatorname{diag}(z,z^{-1})L(a,b)
\]
が最後の可換性を与える．さらに$T_{-1}$と$e_ce_d$の差は
全次数一生成元への$-1$倍なので，中央商の関係も満たされる．
従ってこれは単なる閉点の接表現でなく，作用groupoid全体の有効なcocycleである．

この族が一致点に沿う全形式変形を表すことも直接確認できる．
Artin局所基底上で第一の非退化二形式を$uv$にし，
第二形式の可逆な$uv$係数を生成元の単元倍で一にすれば上の形になる．
順序を保つ正規形間の同型は$T_z,r,e_c,e_d$で尽くされる．
順序を交換する同型は，上の$\tau$を合成して順序を保つ同型へ帰着する．
よって本表示は全同型を含み，論文Bの一致点の形式表示と一致する．

\paragraph{出力の全群作用.}
出力族を
\[
 C^2=U^2-4qV^2,\qquad D^2=V(U+V)
\]
とし，入力の生成元の像を
\[
 \begin{array}{c|cccc}
 &U'&V'&C'&D'\\\hline
 T_z&U&V&C&D\\
 r&U&V&-C&D\\
 e_c,e_d&U&V&C&-D\\
 \tau&-U-4qV&U+V&-\rho C&\rho D
 \end{array}
\]
で与える．$\tau$の二次元行列の二乗は$\rho^2I_2$であり，
二形式はいずれも$\rho^2$倍される．従ってこれは実際の出力同型で，
その二乗は共通次数スカラーである．
三つの非自明な生成元$r,e_c,\tau$の像は可換な独立の対合であり，
閉点では出力の全慣性$\mu_2^3$を生成する．
入力の連結トーラスが核で，その商は$\mu_2^3$であるから，
表示した慣性準同型が得られる．

\paragraph{自然なgeneric射との一致.}
入力のspectral曲線上の$y=c_0d_0$に対し
\[
 (U,V,C,D)=(au^2+bv^2,uv,au^2-bv^2,y)
\]
を用いる．平方根を補助変数として添加した代数的chartの$ab\ne0$上では，
これは指定したFrickeのgeneric同型である．
ここで$\operatorname{Spf}A$自体に$ab\ne0$という開があると主張してはいない．
$T_z,r,e_c,e_d$に関する同変性は直ちに従う．
$\tau$では$y'=\rho y$であり，直接計算により
\[
 \begin{split}
 -a(u')^2-b(v')^2&=-U-4qV,\\
 u'v'&=U+V,\\
 -a(u')^2+b(v')^2&=-\rho C
 \end{split}
\]
となる．従って出力の表は自然なFricke射の全入力同型への作用そのものである．
特にこの射は，粗式だけが同じで一般慣性核が位数八になる弱い共変射とは異なる．
\end{proof}

\begin{corollary}[特殊慣性と形式fiber]
\label{cor:qa-fricke-formal-boundary-fiber}
前定理の射の一致点での相対慣性は$\Gm$であり，
特殊慣性写像は$H_{\rm eq}\twoheadrightarrow\mu_2^3$である．
一方，代数化したchartの$ab=0$の各穿孔枝では入力慣性は$D_8$，核は
$\mu_2=\langle e_ce_d\rangle$，像は$\mu_2^2$で指数二である．
同じ代数的chartの$ab\ne0$の一般点でも核は$\mu_2$で，
順序付き出力の全慣性への像の位数は八である．
非順序付き出力の一般慣性は位数十六なので，順序忘却後の像は指数二となる．

出力の閉点とその同型を固定した形式fiberは
\[
 [\operatorname{Spf}k[[a,b]]/(ab)\,/\Gm],
 \qquad z\cdot(a,b)=(z^{-2}a,z^2b)
\]
である．これは$[\Spec k[a,b]/(ab)\,/\Gm]$の原点に沿う形式完備化であり，
その代数的模型の二つの穿孔枝は各々$B\mu_2$，原点は$B\Gm$となる．
\end{corollary}
\begin{proof}
群準同型$H\twoheadrightarrow\Gamma$の核が$\Gm$であるため，
出力のframeを固定した2-fiberではその核だけで割る．
基底のfiber方程式は$ab=0$である．
軸上の非零点を保つトーラス元は$z^2=1$だけで，
$r$は二つの軸を交換するためその点の安定化群に入らない．
$e_c,e_d$と，必要なトーラス元を合成した$\tau$により
入力慣性$D_8$を得て，像は$r$の像を含まない$\mu_2^2$となる．
$ab\ne0$では$r$にもトーラス元を合成して点を固定できるので，
像は全$\Gamma$となり，核は引き続き$\mu_2$である．
\end{proof}


\subsection{全noncubic境界へのFricke射の降下}

本節では前節の形式計算を，平方根を補助変数とする\'etale chartで
全同型の計算へ拡張する．出力の二形式には順序を保持する．
\[
 \mathfrak U=
 \{(g_1,g_2)\in\mathfrak T_{\rm bal}^{\rm ord}
       \mid \Delta(g_2)\ne0,\ \operatorname{Res}(g_1,g_2)\ne0\}
\]
と置く．この条件は表示の変更に不変である．
非退化な$g_1$を持つ部分は$\mathfrak N^{\rm sm,ord}$であり，
残る幾何点は第一形式が二重根，第二形式がsquarefreeで，
両者の根が交わらない線束境界点である．

\begin{theorem}[自然な順序付きFricke射の大域的境界延長]
\label{thm:qa-global-fricke-balanced-extension}
自然なsmooth射$\mathscr D^{\rm ord}$を延長する射
\[
 \overline{\mathscr D}^{\rm ord}:\mathfrak N^{\rm pres}
       \longrightarrow\mathfrak U
\]
が存在する．構成は表示groupoidの全同型を保存し，
任意基底変更と両立する．
その一根共有および根対一致の特殊像は，ともに
$(U^2,V(U+V))$が表す線束境界点である．
この出力は二次自己双対線束を持つが，元のnoncubic開には属さない．
\end{theorem}
\begin{proof}
\paragraph{一致境界を含む開と\'etale補助変数.}
$\mathfrak N_*=\{\chi\ne0\}\subset\mathfrak N^{\rm pres}$とし，
\[
 q=\frac{\chi-1}{4\chi},\qquad
 T=\Spec k[q,(1-4q)^{-1}]
\]
と置く．$\mathfrak N_*$から$T$への射は正則である．
\[
 \widetilde T=
 \Spec k[q,\rho,\ell,(\rho\ell)^{-1}]/
 (\rho^2-1+4q,\ 2\ell^2-1-\rho)
\]
は$T$の全射\'etale被覆である．
実際，$2\rho$と$4\ell$は可逆であり，
幾何点では$\rho$の二つの選択の少なくとも一つについて$1+\rho\ne0$となる．
特に$q=0$上では$\rho=1$，$\ell=\pm1$である．
\[
 S=\Spec k[a,b,(1-4ab)^{-1}]\times_{q=ab,T}\widetilde T
\]
とし，前節と同じ$H=H_{\rm eq}$を作用させる．
$q,\rho,\ell$は$H$-不変とする．
前節の全生成元と全群関係の式は，この代数的chart上でも定義される．

\paragraph{全groupoidの正規化.}
表示族$(uv,au^2+uv+bv^2)$は同値
\[
 [S/H]\xrightarrow{\sim}
 \mathfrak N_*\times_T\widetilde T
 \tag{F-chart}\label{eq:qa-global-fricke-chart}
\]
を与える．これを全基底上の同型について確認する．
まずfppf局所的に第一の非退化二形式を$uv$へ移す．
$\chi\ne0$より第二形式の$uv$係数は可逆であるため，
第二生成元を単元倍してその係数を一にできる．
このとき$\chi=(1-4ab)^{-1}$なので，補助変数の条件も保存される．
従って表示族は対象について局所的に全射である．

順序を保つ二つの正規形の同型を取る．
第一生成元のスカラーを共通次数スカラーで一にすると，
座標行列は$O(uv)$に属する．第二形式の$uv$係数が一なので，
第二生成元の残るスカラーは$\pm1$である．
従って全ordered同型は$T_z,r,e_c,e_d$で尽くされる．
これは任意環上でも成り立つ．$uv$を保つ行列の対角型・反対角型は，
対応する冪等元による開閉分解上で決まるからである．
非順序の同型は，明示的な$\tau$を一度合成してordered同型へ帰着する．
前節で全群関係を確認したので，\eqref{eq:qa-global-fricke-chart}は
幾何点だけでなく全groupoidの同値である．

\paragraph{chart上の出力.}
$S$上で
\[
 (g_1,g_2)=(U^2-4abV^2,V(U+V))
\]
を取り，前節の出力作用表を用いる．
$\Delta(g_2)=1$かつ$\operatorname{Res}(g_1,g_2)=1-4ab=\rho^2$は可逆なので，
これは$\mathfrak U$の対象である．全群関係が成立するため，
\eqref{eq:qa-global-fricke-chart}上の射が得られる．
出力の対象自体は$\rho,\ell$の選択に依存しない．
以下で矢についても補助変数から降下することを確認する．

\paragraph{平方根の選択を変えた際の全Isom.}
二重overlapで二つの選択を$(\rho,\ell)$，$(\rho',\ell')$と書く．
$2,\rho,\ell$が可逆なので，底を開閉分解して符号関係を固定できる．
ordered矢の持上げと出力は補助変数に依存しないため，
$\tau$の持上げだけを比較すれば十分である．

$\rho'=\rho,\ell'=-\ell$の場合，入力の$\tau$の持上げは
$T_{-1}\tau=(e_ce_d)\tau$へ変わる．
出力は$T_{-1}$を恒等へ写すため，出力の矢は変わらない．

$\rho'=-\rho$の場合は，$\ell,\ell'$の可逆性より$q$も可逆である．
$m=\ell\ell'$と置くと$m^2=q=ab$であり，
\[
 R_{a,b,m}=\begin{pmatrix}0&b/m\\a/m&0\end{pmatrix}\in O(uv)
\]
が定義される．入力の座標行列は正確に
\[
 L(\rho',\ell')L(\rho,\ell)^{-1}=R_{a,b,m}
\]
を満たす．第一生成元の像は同じで，第二生成元の像は反対符号となるので，
入力の二つの$\tau$の差は，parameter$(-a,-b)$の対象の
ordered自己同型$R_{a,b,m}e_d$である．
$R_{a,b,m}=T_{b/m}r$だから，その出力は$C$の符号であり，
$R_{a,b,m}e_d$の出力は$\delta'=(C,D)\mapsto(-C,-D)$となる．
一方，出力の$\tau$の式は
\[
 (U,V,C,D)\longmapsto
 (-U-4qV,U+V,-\rho C,\rho D)
\]
なので，$\rho$の符号を変えると正確に$\delta'$だけ変わる．
従って入力の同じ矢は，どちらの補助変数を使っても同じ出力の矢を与える．

以上は対象上の恒等同型を比較として用いた等式であるため，
三重overlapのcocycleも成立する．全てのunordered矢は
ordered矢と$\tau$の合成なので，これで全Isomの降下が確認できた．
有限局所自由束・二次積およびその同型の有効fppf降下により，
$\mathfrak N_*$上の射を得る．

\paragraph{smooth射との貼合せ.}
$\mathfrak N_*\cap\mathfrak N^{\rm sm}$では$ab$が可逆である．
前節の$\Phi$に関する全生成元の同変計算により，
構成した射は自然な$\mathscr D^{\rm ord}$と全groupoid上で一致する．
$\chi=0$の入力はresultantが非零なのでsmooth域に属する．
従って$\mathfrak N_*\cup\mathfrak N^{\rm sm}=\mathfrak N^{\rm pres}$であり，
二つの射を貼り合わせて主張の大域射を得る．
構成は局所自由な二次代数とその同型の式および有効降下によるので，
任意基底変更と両立する．
\end{proof}

\begin{corollary}[線束境界上の全fiber]
\label{cor:qa-global-fricke-boundary-fiber}
第一形式が二重根となる順序付き線束境界点$\ell_0$とそのframeを固定する．
前定理の射のscheme-theoretic $2$-fiberは
\[
 \boxed{
 (\overline{\mathscr D}^{\rm ord})^{-1}(\ell_0)
 \simeq[\Spec k[a,b]/(ab)\,/\Gm],
 \qquad z\cdot(a,b)=(z^{-2}a,z^2b).
 }
\]
二つの穿孔枝は各々$B\mu_2$であり，原点は$B\Gm$である．
この原点が一致根対，穿孔枝が一根共有に対応する．
一致点の慣性写像は$H_{\rm eq}\twoheadrightarrow\mu_2^3$で核$\Gm$，
一根共有点では$D_8\to\mu_2^3$で核$\mu_2$，像は位数四となる．
\end{corollary}
\begin{proof}
像が$\ell_0$となる入力は$\chi=1$である．
前定理のchartを$q=0,\rho=\ell=1$へ制限すると，
入力は$[\Spec k[a,b]/(ab)/H]$，出力のresidual gerbeは$B\Gamma$となる．
$H\twoheadrightarrow\Gamma$の核は$\Gm$であるから，
出力のframeを固定した2-fiberは核だけで割った記載の商となる．
非零軸上のトーラス安定化群は$\mu_2$，原点の安定化群は$\Gm$であり，
慣性写像については前節の全生成元表から従う．
\end{proof}

\begin{remark}[延長の標的とH2障害]
この定理は自然なspectral射の全groupoid上の境界延長を与えるが，
標的は二次積付きTF標的の開部分$\mathfrak U$である．
特殊出力の第一形式は重根を持つためH2には属さず，
既存のH2延長不能・無限Extの結論は変わらない．
また表示chart上で出力conicがsplitであることから，
降下後の全conicもsplitであると推論しない．
出力の全同型作用には$r,e_c,\tau$の独立な三対合が残り，
一般慣性核は$\mu_2$，非順序付き標的への一般像は位数八である．
この点で弱い共変式の一般慣性核八・像二とは異なる．
\end{remark}

% !TEX root = ../nc-quadric-surface-avatar.tex
% 2026-11-23（canonical square完成の相対tails認識を含む）
\section{内在的なモジュライへの忘却}
\label{sec:qa-intrinsic-forgetting}

有限 foundation、周期的 AS 代数、avatar、および構造対象だけを指定した
Abel 圏は、保持する構造も同型の意味も異なる。
本節では、前節までの開部分と内在的モジュライとの関係を、
既に構成した射と、相対的な認識を要する射とに分けて述べる。

\subsection{標的となる強い内在的クラス}

本稿の \(\HH\)、\(\HHstack\) は
\cite[\texttt{def:h2}]{QuadricIntrinsic2026} と
\cite[\texttt{def:relative-h2}]{QuadricModuli2026} の意味で用いる。
特に、任意の三次元 AS 正則 \(\mathbb Z\)-代数の tails 圏を
無条件に \(\HH\) と呼ぶわけではない。

\begin{definition}[体上の強い内在的クラス]
\label{def:qa-h2}
\(\HH\) の対象は \((\mathcal C,\OO)\) であり、次を満たす。
\begin{enumerate}
\item \(\mathcal C\) は \(\kk\)-線形 noetherian Abel 圏で、
\(D^b(\mathcal C)\) は Hom 有限で Serre 関手 \(S\) を持つ。
\item \(\OO\) は rank-one torsion-free exceptional 対象である。
ここで $\mathcal C_0$ は有限長対象の部分圏、$\mathcal C_1$ は
$\mathcal C/\mathcal C_0$ で有限長となる対象の逆像とし、
$q:\mathcal C\to\mathcal C/\mathcal C_1$ とする。
条件は $q\OO$ が非零単純かつ $\Hom(\mathcal C_1,\OO)=0$ である。
また exceptional とは $\End(\OO)=\kk$、$\Ext^{i}(\OO,\OO)=0$ ($i>0$) をいう。
\item 完全自己同値 \(s\) が存在し、
\(s^2\simeq S^{-1}[2]\) かつ \((\OO,s)\) は Artin--Zhang ample である。
\item 同じ \(s\) に対して、全軌道環
\[
 B(\mathcal C,\OO,s)
   =\bigoplus_{n\ge0}\Hom(\OO,s^n\OO)
\]
が \(A/(g)\) と同型である。
ここで \(A\) は左右 noetherian、二次 Koszul AS 正則、
型 \((1,4,6,4,1)\) の環、\(g\in A_2\) は左右正則正規元である。
\end{enumerate}
\(s\)、根同型、\(A\)、\(g\)、商表示は指定構造ではない。
射は \(\OO\) の同型を伴う完全同値であり、
その同型を保つ自然同型によって一射を同一視する。
\end{definition}

ここで型 \((1,4,6,4,1)\) は、自明加群の最小解消の階数列が
この列で、global dimension と Gorenstein parameter がともに四であることを意味する。
\(g\) の中心性は要求しない。
また、ある商圏との同値だけでなく、指定した \(\OO\) と同じ \(s\) の
\emph{全軌道環} に対する表示を要求する。

相対版 \(\HHstack\) は、標数零の代数閉体上の局所有限型スキームの
fppf site 上で定義する。
圏の基底変更は Ind 圏における加群対象の圏によるものであり、
Hom を tensor するだけの基底変更ではない。
対象と完全列の fppf descent、有限 Hom と flat base change、
平坦な \(\OO\) を要求する。
局所 canonical compact 導来圏は相対 smooth proper で、
相対 Serre 双対性と導来基底変更に適合する。
全幾何ファイバーは定義~\ref{def:qa-h2} を満たす。
さらに、\(s\)、AZ 条件、全軌道環の AS4 被覆は
\emph{非被約基底も含めて fppf 局所に存在する}ことを要求する。
根や被覆を重なり上で同一視することは要求しない。
これらは標識ではなく存在条件だからである。

\begin{remark}[相対存在は幾何点条件より強い]
\label{rem:qa-relative-root}
全幾何ファイバーが \(\HH\) に入っても、
相対平方根、AZ ampleness、相対 AS4 被覆が存在することは別問題である。
例外対象の無障害性や有限代数の \(\operatorname{HH}^3\) の消滅だけでは、
これらの条件は得られない。
\cite[\texttt{prop:chosen-root-does-not-lift}]{QuadricModuli2026} は、可換 quadric 上の
ある適合根が、特定の一次 Abel 変形へ持ち上がらない例を与える。
これは別の根の存在まで否定する例ではなく、
\(\HHstack\) の非代数性や許容部分の非開性を証明するものでもない。
\end{remark}

\subsection{canonical square 完成で確立した定義域}

定理~\ref{thm:qa-relative-square-completion-tails} と
命題~\ref{prop:qa-relative-square-unfolding} により，実際の曲線 avatar
から構成した有限 square 代数の canonical 完成と全整数展開には，
標準対象を保つ相対 tails 認識がある．これは幾何点の同値だけではなく，
本稿の site 上の module-object 基底変更と coherent に整合する完全同値である．
左右 Noether 性と係数束付きの実際の AS 解消も証明済みである．
補題~\ref{lem:qa-classified-fiber-finite-square-recognition} は，
古典型を含む分類済み square fiber を持つ有限局所自由族まで扱う．
従ってこの canonical 部分の相対 heart 回復を未解決のままにしない．
一方，任意の別途与えた AS/QP 族との全整数同一視には
注意~\ref{rem:qa-square-unfolding-vertex-lines} の頂点線束・周期比較が必要であり，
以下の一般 presentation stack 全体の認識と区別する．

さらに，定理~\ref{thm:qa-square-normal-foundation-stack-equivalence} は
有限 square と normal foundation の全スタックを同定し，
系~\ref{cor:qa-h2-square-foundation-open} は既存の内在的開をこの表示へ移す．
全 $\HHstack$ については
定理~\ref{thm:qa-universal-h2-rank81-completion} が，ruling の選択を要しない
階数八十一の通常代数と pointed tails を与える．
ここまでの有限表示・canonical 完成の認識は確立済みである．
定理~\ref{thm:qa-square-explicit-h2-open}と
系~\ref{cor:qa-square-h2-twenty-polynomials}は，この全$H_2$開を
三十二次元線形表示上の四組・計二十式以内で記述する．
任意の別途与えた AS/QP presentation 群oidとの比較は後続課題である．

\subsection{忘却射の正確な定義域}

\begin{definition}[内在的条件に適合する族]
\label{def:qa-compatible-chart}
\(\mathfrak U\) を、周期的代数または avatar から
基底変更と同型に関して自然に \((\mathcal C,\OO)\) を構成する
fppf パラメータスタックとする。
\(\mathfrak U_{\HH}\) は、その構成した族が
上記の \(\HHstack\) の相対条件を満たす対象の全部分スタックとする。
この定義において、\(\mathfrak U_{\HH}\) が代数的または開であることは
仮定も結論もしない。
\end{definition}

\begin{proposition}[適合族からの忘却]
\label{prop:qa-compatible-forgetting}
定義~\ref{def:qa-compatible-chart} の構成は射
\[
 \pi_{\mathfrak U}:\mathfrak U_{\HH}\longrightarrow\HHstack,
 \qquad u\longmapsto(\mathcal C_u,\OO_u)
\]
を与える。
二つの chart 間の変換が、基底変更と両立する pointed exact 同値を
与えるなら、二つの忘却射はその変換に沿って自然に二同型となる。
\end{proposition}
\begin{proof}
\(\mathfrak U_{\HH}\) の条件は fppf 局所的で、
pointed exact 同値と基底変更に安定である。
従って同型で飽和した全部分スタックを定める。
構成した圏の系、構造対象、その比較同型を保持して
残る標識を忘れれば、定義された標的への functor を得る。
二つの構成を結ぶ pointed exact 同値が任意の基底変更と整合する場合には、
それらの同値がそのまま二つの functor の二同型を与える。
\end{proof}

この命題は忘却の意味を明示するものであり、
未知の条件を有限個の方程式で認識したという定理ではない。
特に \(\mathfrak F_{\le12}\) は有限代数だけを分類するので、
そこから直接 \(\mathcal C\) を取ることはできない。
\(\operatorname{Perf}(A)\) の構成だけでは、元の曲面の heart は選ばれていない。
有限代数の標準的な加群 heart と、tilting によって移された曲面の heart を
同一視してはならない。

「十二本矢の適当な開部分から \(\HHstack\) への自然な射」という主張を
有限な表示で実現するためには、少なくとも次を示す必要がある。
\begin{enumerate}
\item 使用する periodic algebra、grading、foundation と heart を
相対的に構成し、全次数の平坦性、noetherian 性、基底変更を確認する。
\item 指定した \(\OO\) の相対例外性・rank 条件、相対 Serre 公式、
適合する完全平方根と AZ ampleness を確認する。
\item 同じ根の全軌道環が、相対 AS4 環の正則正規二次切断となることを確認する。
\item これらの条件を持つ部分関手が、考えている表示スタックの
不変開部分で表されることを証明し、表示変更の全ての射を降下させる。
\end{enumerate}
根と全軌道環が既に与えられた場合、固定軌道環への AS4 被覆を選ぶ空間は
\cite[\texttt{thm:relative-finite-as-cover}]{QuadricModuli2026} により、
\(\operatorname{Gr}(6,K)\) 内の有限表示局所閉 incidence として記述できる。
しかし、その被覆を忘れた fppf 像の開性・代数性や、根そのものの存在は
この有限判定の結論ではない。

\subsection{十二本矢の全体を送ることへの具体的な障害}

\begin{lemma}[反 Serre 軌道の生成は必要である]
\label{lem:qa-antiserre-necessary}
\((\mathcal C,\OO)\in\HH\) ならば、
\(t=S^{-1}[2]\) に対する \((\OO,t)\) も AZ ample である。
特に、非零対象 \(M\) に対して
\(\Hom(t^{-n}\OO,M)=0\) が全ての \(n\ge0\) で成立することはない。
\end{lemma}
\begin{proof}
定義の根 \(s\) を選び、\(E=\OO\) と置く。
\(B=A/(g)\) は次数一で生成される。
AZ の生成条件を \(sE\) に適用し、得られた全射に
\(s^{-1}\) を作用させると、\(E\) は
\(s^{-m}E\)（\(m>0\)）の有限直和の商になる。
次数一生成により各成分 \(s^{-m}E\to E\) は、
最後の段が \(s^{-1}E\to E\) である合成の線形結合である。
従って評価 \(\Hom(s^{-1}E,E)\otimes s^{-1}E\to E\) は全射となる。
Hom 有限性により、この評価の源は有限直和である。
これは \cite[\texttt{lem:ample-square-from-generation}]{QuadricIntrinsic2026}
の評価全射である。

任意の対象への AZ 生成全射で奇数指数 \(m\) の項があれば、
この評価全射を \(s^{-m}\) で移して \(s^{-(m+1)}E\) の
有限直和に置き換える。従って偶数指数の項だけで生成できる。
指数に任意の下限を課す場合は、先に対象を正に shift して
生成全射を作り、負の shift で戻す。
全射に対する Hom の最終的な全射性は
元の AZ 条件を偶数部分列に制限すれば従う。
\(s^2\simeq t\) なので \((E,t)\) は AZ ample であり、
最後の主張はその生成条件から従う。
\end{proof}

\begin{example}[重み付き十二本矢と residual gerbe]
\label{ex:qa-weighted-not-h2}
\(S=\kk[x,y,z]\)、\(\deg x=\deg y=1\)、\(\deg z=2\) とし、
\(\mathcal X=\PP(1,1,2)\) を重み付き射影 \emph{stack} とする。
この環は global dimension 三、Gorenstein parameter 四の
noetherian AS 正則環である。
しかし
\[
 (\coh\mathcal X,L)\notin\HH\qquad(L\in\Pic(\mathcal X)).
\]
\end{example}
\begin{proof}
通常の三変数 Koszul 解消に各変数の重みを入れると
\[
 0\longrightarrow S(-4)
 \longrightarrow S(-3)^2\oplus S(-2)
 \longrightarrow S(-1)^2\oplus S(-2)
 \longrightarrow S\longrightarrow\kk\longrightarrow0
\]
を得る。これは AS 条件と Gorenstein parameter を与える。

\((\OO,\OO(1),\OO(2),\OO(3))\) の有限 End 代数を考える。
隣接する \(x_i,y_i\) が六本、距離二の \(z_0,z_1\) が二本であり、
最小関係は
\[
 \begin{aligned}
 x_1y_0-y_1x_0&=0,&x_2y_1-y_2x_1&=0,\\
 z_1x_0-x_2z_0&=0,&z_1y_0-y_2z_0&=0.
 \end{aligned}
\]
従って relation completion は十二本矢である。
前節の四写像表示では \(a,b\) はともに階数三で、核は交代テンソルの直線、
\(c,d\) は全射である。よってこの点は実際に
\(\mathfrak F_{\le12}\) の最小十二本矢 locus に属する。
この四対象が full strong であることも直接確認できる。
重み付き多項式環の局所コホモロジーから
\(H^1(\mathcal X,\OO(n))=0\) であり、
Serre 双対により \(n\ge-3\) では \(H^2(\mathcal X,\OO(n))=0\) である。
これで strong 性が従う。
三変数の sheaf Koszul 完全列は四つの連続する twist から
全ての twist を生成し、有限生成 graded \(S\)-加群の有限自由解消により
これらは \(D^b(\coh\mathcal X)\) 全体を生成する。

一方、重み付き Euler 列から
\(\omega_{\mathcal X}=\OO(-4)\) であり、\(t=-\otimes\OO(4)\) である。
点 \([0:0:1]\) の residual gerbe を
\(i:B\mu_2\hookrightarrow\mathcal X\) とし、
\(\chi\) を \(\mu_2\) の非自明な一次元表現とする。
任意の可逆層 \(L\) に対し、非零 coherent 対象
\(M_L=i_*(i^*L\otimes\chi)\) を取る。
\(\OO(4n)\) の gerbe 上の指標は自明なので、引戻しと直像の随伴により
\[
 \Hom(L(-4n),M_L)=
 H^0(B\mu_2,i^*\OO(4n)\otimes\chi)=0\qquad(n\ge0).
\]
特に標準 helix の四位相すべてで生成条件が破れる。
補題~\ref{lem:qa-antiserre-necessary} に反するので、主張を得る。
\end{proof}

\begin{example}[\(\Sigma_2\) 自身の構造対象]
\label{ex:qa-hirzebruch-not-h2}
\(Y=\Sigma_2\) の負断面を \(C\)、ファイバーを \(f\) とする。
\(C^2=-2\)、\(-K_Y=2C+4f\) なので、
\(\OO_Y(-K_Y)|_C\simeq\OO_C\) である。
任意の可逆層 \(L\) に対し、
\(M_L=i_*(i^*L\otimes\OO_C(-1))\ne0\) とおくと
\[
 \Hom\bigl(L\otimes\OO_Y(nK_Y),M_L\bigr)
   =H^0(C,\OO_C(-1))=0\qquad(n\ge0).
\]
同じ必要条件から \((\coh\Sigma_2,L)\notin\HH\) となる。
これは \emph{avatar の一点 Hilbert 曲面が \(\Sigma_2\) である}ことと、
元の Abel 圏が \(\coh\Sigma_2\) であることを混同してよいという意味ではない。
Hilbert 曲面の型だけから、対応する非可換圏の \(\HH\) 適合性は決まらない。
\end{example}

この二例は、十二本矢の全 locus、AS 正則性、あるいは
full strong foundation の存在だけで忘却の標的が \(\HHstack\) になる、
という主張を否定する。
適切な開部分が存在する可能性を否定するものではなく、
その開部分には構造対象の反 Serre 軌道に対する生成条件も必要であることを示す。

\subsection{標識、射、既存の表示スタック}

\(\HHstack\) への忘却では \(\OO\) を保持する。
従って、foundation の mutation が構造対象を別の対象へ移す場合には、
そのまま同じ pointed 対象の表示変更と扱えない。
QP の mutation が dg 導来同値を与えても、
指定した heart の完全同値と \(\OO\) の同型は別途必要である。
\cite[\texttt{cor:pointed-derived-exact}]{QuadricIntrinsic2026} は体上の \(\HH\) 対象の間の判定であり、
任意の相対 periodic family について同じ結論を自動的に与えるものではない。

同様に、消去可能な二本の矢 \(x,y\) と項 \(yx\) を付加した生の QP には、
\((x,y)\mapsto(\lambda x,\lambda^{-1}y)\) という自己同型がある。
消去後の表示ではこの自己同型は見えなくなる。
従って、生の安定化 QP groupoid の消去を、
射を補正しないままスタック同値とすることはできない。
有限 foundation の頂点スカラーの剛体化、
周期的環の次数スカラーの除去、
構造対象を保つ自己同値全体での同一視は、それぞれ異なる操作である。

現在の内在的研究との関係をまとめると、
\cite[\texttt{prop:foundation-stack}]{QuadricModuli2026} の有限二写像表示 \(\mathfrak F_0\) は、
一方の長さ二の積を可逆にした foundation chart である。

\begin{proposition}[幾何点における十本矢 chart の被覆]
\label{prop:qa-h2-normal-ten-chart}
体上の \(\HH\) の任意の対象は、
\(\mathfrak F_{\le10}^{(a)}\) に属する normal foundation を持つ。
これは幾何点での主張であり、任意の \(\HHstack\) の族における
foundation の相対的な存在・降下を結論するものではない。
\end{proposition}
\begin{proof}
\cite[\texttt{thm:normal-helix}]{QuadricIntrinsic2026} と
\cite[\texttt{prop:ruling-involution-mutation}]{QuadricModuli2026} により、
\((E_0,E_1,E_2,E_3)=(\OO,X,s\OO,Z)\) という foundation があり、
本節の \(b\) は同型、\(c\) は全射である。
\(E_0,E_1,E_2\) は rank-one torsion-free 対象である。
このような対象間の非零射は単射である。
実際、像が低い Gabriel--Krull 段階に入れば標的の torsion-free 性に反し、
従って rank-one 商では非零射、すなわち同型を誘導する。
その核は低い段階に入るが、始域の torsion-free 性によって零である。
従って二つの隣接する非零射の合成は零にならず、
\(\ker a\) は非零の分解可能テンソルを含まない。
\(b\) は同型なので他方の核は零であり、
\(c\) の全射性から定義~\ref{def:qa-foundation-open} の全条件を満たす。
\end{proof}

相対的な ruling 交換を \(\tau\) とすると、
\[
 \mathfrak Q_{\mathrm{fnd}}=[\mathfrak F_0/S_2]
\]
は有限型 Artin スタックであり、体上では
\(\HH(\kk)\) をその full subgroupoid として実現する
（\cite[\texttt{cor:h2-point-groupoid-foundation}]{QuadricModuli2026}）。
この幾何点の結果だけから、\(\mathfrak Q_{\mathrm{fnd}}\) 全体または
その未特定の開部分が \(\HHstack\) と同値であるとはいえない。

相対的には、同論文の反 Serre 認識によって全忠実な射
\[
 \HHstack\hookrightarrow
 \mathfrak P=[X_3/\PGL_9]
\]
がある。
ここで \(X_3\subset\operatorname{Gr}(56,81)\) は、
九生成元・五十六二次関係の環の三次成分が階数四十九となる
scheme-theoretic な局所閉 locus である。
\[
 X_{\mathrm{as}}=X_3\times_{\mathfrak P}\HHstack
\]
は不変 fppf 部分層として定義され、
\(\HHstack\simeq[X_{\mathrm{as}}/\PGL_9]\) が成立する。
\(X_{\mathrm{as}}\) の有限型・滑らか・既約な純次元83のスキームとしての表現可能性は
相対モジュライ論文の定理~\cite[\texttt{thm:global-h2-smooth-finite-type}]{QuadricModuli2026} と
系~\cite[\texttt{cor:global-xas-smooth-scheme}]{QuadricModuli2026} で証明されている。
ただし $X_{\mathrm{as}}\to X_3$ は一致根対で開ではなく、
局所閉性と明示的な偏極認識方程式は未解決である。
従って本稿の八・十・十二本矢の明示的な開部分は、
この相対認識問題への有限表示の候補を与えるが、
それ自体で内在的モジュライの代数性を解決するものではない。

% !TeX root = ../quiver-avatar-moduli.tex
% 2026-11-22 / v0.5（OU z2座標とPicard根の明示対応、Kummer障害を追加）
\section{既存GITと内在的モジュライの比較}
\label{sec:qa-git-comparison}

本節は Okawa--Ueda \cite[Theorem 1.1]{OU2014} の
quintuple と conifold potential の GIT 商を、
本稿の有限 foundation と内在的 $\HHstack$ に対する比較問題へ接続する。
二つの既存 GIT 商自体を本稿の新しい成果とはしない。
また、表示を忘れた $\HHstack$ への射は
定義~\ref{def:qa-compatible-chart} の適合する領域でのみ構成し、
GIT compactification 全体で自動的に存在するとは主張しない。

\subsection{既存の次数四射と三つの忘却}

$V_0,V_1,V_2,V_3$ を二次元、$W$ を一次元の $k$-ベクトル空間とし、
\[
 T=(V_0\otimes V_1\otimes V_2\otimes V_3)\otimes W^\vee,\qquad
 \widetilde G=\prod_{i=0}^3\GL(V_i)\times\GL(W)
\]
と置く。Okawa--Uedaの生のquintuple表示スタックを
$\mathfrak M_{\mathrm{qui}}^{\mathrm{raw}}=[T/\widetilde G]$ と書く。
$T$ の零元を含む表示スタックの記法であり、
$\HHstack$ を定義するものではない。
同論文のpotential表示を $\mathfrak M_{\mathrm{pot}}$ とし、
その compact GIT 商を $\overline M_{\mathrm{pot}}$、
quintuple側の商を $\overline M_{\mathrm{qui}}$ とする。
既知結果は
\[
 \overline M_{\mathrm{pot}}\simeq\PP(1,2,3,4),\qquad
 \overline M_{\mathrm{qui}}\simeq\PP(2,4,4,6),\qquad
 \overline M_{\mathrm{pot}}\longrightarrow\overline M_{\mathrm{qui}}
 \text{ が有限次数 }4
\]
である。右端の射は潜在的なSerre平方根の忘却射ではなく、
\emph{potentialからquintupleへの表示の射}である。

本研究では区別すべき射が少なくとも三つある。
potentialからquintupleへの射、quintupleから四項foundationへの
表示上の射、そして適合foundationから pointed exact 対象
$(\mathcal C,\OO)$ への忘却射である。
最後の射では、quintupleの標識、ruling、Serre平方根とその同型、
ambientを保持しない。これらの射のどの組合せが
相対スタック射として存在するかは、対象と自然同型を個別に調べる必要がある。

\subsection{スカラー慣性の核}

\begin{proposition}[生のquintuple表示の四次元スカラー核]
\label{prop:qa-ou-scalar-gerbe}
次の中央部分群スキームは $\widetilde G$ に含まれる。
\[
 K=\left\{
 \bigl(\lambda_0\mathrm{id}_{V_0},\ldots,
       \lambda_3\mathrm{id}_{V_3},
       (\lambda_0\lambda_1\lambda_2\lambda_3)\mathrm{id}_{W}\bigr)
 :\lambda_i\in\Gm
 \right\}\simeq\Gm^4.
 \tag{GIT-K}\label{eq:qa-ou-scalar-kernel}
\]
$K$ は $T$ 全体に自明に作用する。従って
\[
 [T/\widetilde G]\longrightarrow[T/(\widetilde G/K)]
 \tag{GIT-rig}\label{eq:qa-ou-rigidification}
\]
は $K$ を相対慣性とする fppf gerbe である。

さらに、quintupleから対応する $\mathbb Z$-代数のtails圏と
添字零の構造対象を作る通常の構成が実際の相対族を与え、
その対象が $\HHstack$ の公理を満たす領域では、
$K$ の全元は pointed exact 自己同値の\emph{恒等射}へ写る。
従って、この領域の生の表示スタックから $\HHstack$ への
標準的な忘却射は表現可能ではない。
少なくともこのスカラー慣性を消さずに
全忠実な相対比較射を得ることはできない。
\end{proposition}
\begin{proof}
$\widetilde G$ のスカラー部分の作用は
$T$ 上で
\[
 (\lambda_0,\ldots,\lambda_3,\mu)\cdot t
   =(\lambda_0\lambda_1\lambda_2\lambda_3)\mu^{-1}t
\]
である。従って式~\eqref{eq:qa-ou-scalar-kernel} は
$T$ への作用の核に含まれる中央な四次元トーラスである。
$\widetilde G/K$ は群スキームとして存在し、
$\widetilde G\to\widetilde G/K$ は忠実平坦なので、
$K$ の作用が自明である場合の商スタックの一般的な構成から
式~\eqref{eq:qa-ou-rigidification} はgerbeになる。

後半は表現を変えるスカラーが圏に与える作用を直接調べる。
quintupleを使って構成した $\mathbb Z$-代数を $D$ とし、
その隣接Hom空間を $V_i$（または採用する双対規約）で表す。
スカラー変更は、各頂点の単元を変えない
\[
 D_{ij}\longrightarrow D_{ij},\qquad
 a\longmapsto\gamma_i^{-1}\gamma_j\,a
\]
という次数ごとの共役に帰着する。
ここで隣接空間での所定のスカラーに合うよう
$\gamma_i\in k^\times$ を漸化的に選ぶ。
四つの隣接スカラーの積が一次関係線 $W$ のスカラーと一致する
ことが、quintupleの関係を保つ条件である。
右graded $D$-加群の次数 $i$ の成分に
$\gamma_i^{-1}$ を掛ける自然変換は、この共役による
twist functorと恒等 functorの自然同型を与える。
したがってtailsへの降下でも恒等と自然同型であり、
添字零の構造対象の同型も同じスカラーで規格化できる。
相対族でも単元を用いて同じ計算ができる。

一方、元の表示スタックでは任意の $K$ の元は
各点の自己同型として残る。
忘却による相対慣性の核が $K$ を含むため、
射が表現可能であるための必要条件（相対慣性の自明性）に反する。
\end{proof}

\begin{remark}[GITの次数と慣性を混同しない]
\label{rem:qa-ou-degree-four-not-root}
Okawa--Ueda \cite[Theorem 1.1(4)]{OU2014} の次数四は、
compact GIT 商の有限射の一般ファイバーの次数である。
式~\eqref{eq:qa-ou-scalar-kernel} の群は
quintuple表示の自己同型に既に存在し、
pointed忘却では消える。
したがってその次数四を、標識を忘れた
$\HHstack$ の一般ファイバーの位数、
あるいは可換点でのSerre平方根の個数と同一視できない。
適合する族の範囲、表示側の剛体化、全pointed同型の像と核、
および閉軌道の選択が別途必要である。
\end{remark}

\subsection{相対ファイバーの群oidとしての公式}

\begin{proposition}[忘却射の幾何ファイバーと慣性核]
\label{prop:qa-forgetful-homotopy-fiber}
$F:\mathfrak X\to\mathfrak Y$ を $k$-上の groupoid stack の射とし、
$y\in\mathfrak Y(k)$ を取る。
幾何ファイバーの groupoid
$\mathfrak X_y=\mathfrak X\times_{\mathfrak Y}\operatorname{Spec}k$
の対象は、$x\in\mathfrak X(k)$ と同型
$\alpha:F(x)\xrightarrow{\sim}y$ の組である。
その点 $(x,\alpha)$ の自己同型群は
\[
 \operatorname{Aut}_{\mathfrak X_y}(x,\alpha)
 =\ker\bigl(\operatorname{Aut}_{\mathfrak X}(x)
       \longrightarrow\operatorname{Aut}_{\mathfrak Y}(F(x))\bigr).
 \tag{Fib-I}\label{eq:qa-fiber-inertia}
\]
さらに、$F(x)\simeq y$ となる $\mathfrak X(k)$ の同型類を
一つずつ選ぶと、$\mathfrak X_y(k)$ の同型類の集合は
\[
 \coprod_{[x]:\,F(x)\simeq y}
 \operatorname{Aut}_{\mathfrak Y}(y)\Big/
  \alpha\,F(\operatorname{Aut}_{\mathfrak X}(x))\,\alpha^{-1}
 \tag{Fib-C}\label{eq:qa-fiber-cosets}
\]
と同一視される。右辺は群の商ではなく、一般には剰余\emph{集合}である。
従って同型類のファイバーの濃度、スタックの相対慣性、
GITの有限射の次数はそれぞれ異なる量である。
\end{proposition}
\begin{proof}
二つの対象 $(x,\alpha),(x',\alpha')$ の間の射は
$f:x\to x'$ で $\alpha'F(f)=\alpha$ を満たすものである。
$x=x'$、$\alpha=\alpha'$ とすれば
式~\eqref{eq:qa-fiber-inertia} が直ちに従う。
固定した $x$ について同型 $\alpha:F(x)\to y$ を一つ選ぶと、
他の全ての同型は $g\alpha$、
$g\in\operatorname{Aut}_{\mathfrak Y}(y)$ の形である。
二つの選択は $\operatorname{Aut}_{\mathfrak X}(x)$ の像による作用で
同型になるため、式~\eqref{eq:qa-fiber-cosets} が得られる。
\end{proof}


\subsection{次数四射の $S_4$-monodromy と Picard 平方根}

\begin{proposition}[OU次数四射の一般Galois閉包]
\label{prop:qa-ou-fourfold-s4-monodromy}
標数零の代数閉体上で、Okawa--Ueda \cite[\S5, (5.8)]{OU2014}
のGIT有限射を、対角化による稠密な一般開で考える。
同論文の表示は
\[
 \PP^3/S_3\longrightarrow
 \PP^3/\bigl(S_3\ltimes(\mathbb Z/2)^3\bigr)
 \tag{OU-4}\label{eq:qa-ou-fourfold-diagonal}
\]
であり、$\PP^3$ では同時符号反転が自明に作用する。
従って有効な符号変更群と全群はそれぞれ
\[
 \Gamma=(\mathbb F_2)^3/
   \langle(1,1,1)\rangle\simeq(\mathbb F_2)^2,\qquad
 G=\Gamma\rtimes S_3\simeq\operatorname{AGL}_2(\mathbb F_2)
 \simeq S_4.
\]
$S_3$ は $\Gamma$ の三つの非零元を置換する標準的な
$\operatorname{GL}_2(\mathbb F_2)$-作用を持つ。
一般の射~\eqref{eq:qa-ou-fourfold-diagonal} は次数四で、
\emph{Galois射ではなく}、その関数体拡大のGalois閉包の群は $S_4$ である。
$\Gamma$ を自明化する表示の方へ引き戻した場合にのみ、
$\Gamma$-torsorとしての四点の符号変更を直接見ることができる。
\end{proposition}
\begin{proof}
$\mathbb F_2^3$ の標準基底の剰余類を $e_1,e_2,e_3$ とする。
商 $\Gamma$ では $e_1+e_2+e_3=0$ なので、
三つはまさに非零元である。$S_3$ がこれらを置換する作用は忠実で、
$\operatorname{Aut}_{\mathbb F_2}(\Gamma)
=\operatorname{GL}_2(\mathbb F_2)$ は位数六なので同型となる。
$\Gamma$ の四元上への平行移動とこの線型作用を合わせた
$G$ は四元上の忠実で推移的な位数 $4\cdot6=24$ の置換群、
従って $S_4$ である。
$H=S_3$ は原点の安定化群で、$[G:H]=4$。
標数零の一般対角座標はこの有限群作用に対し自明な
幾何安定化群を持つので、
関数体拡大 $k(\PP^3)^H/k(\PP^3)^G$ の次数は四である。
$H$ の $G$ における core（全共役の共通部分）は、
四元上の推移作用の核と一致し、忠実性から自明である。
従ってこの非正規拡大のGalois閉包は $k(\PP^3)$ であり、
Galois群は $G\simeq S_4$ となる。
最終のtorsorに関する文は、
$H$ の選択を消した一般の基底では
$\Gamma$ が一定群として作用する四重Galois被覆を
無条件に得られない、という意味である。
置換群の全24元とcoreは
\path{computations/verify_ou_fourfold.py} で再現できる。
\end{proof}

\begin{proposition}[適合Picard平方根の幾何fiberの濃度]
\label{prop:qa-picard-roots-degeneration-count}
論文Bの非線型geometric cubic quintupleの相対曲線
$C\to S$ と $N\in\Pic^0_{C/S}(S)$ を考え、
同論文 \cite[\texttt{lem:picard-halves}]{QuadricModuli2026}
の平方根空間
\[
 T=S\times_{\Pic^0_{C/S},[-2]}\Pic^0_{C/S}
 \tag{PR-4}\label{eq:qa-picard-root-torsor}
\]
を取る。これは $S$ 上 étale 全射である。
滑らかな楕円曲線の locus $S_{\mathrm{sm}}$ 上では
$T_{S_{\mathrm{sm}}}\to S_{\mathrm{sm}}$ は
$\Pic^0_{C/S}[2]$-torsorであり、有限 étale 次数四である。
各代数閉幾何点での平方根の個数は、曲線が
滑らかな楕円曲線、既約一節点の種数一曲線、
既約一尖点の種数一曲線である場合、それぞれ
\[
 \begin{array}{c|ccc}
 C_{\bar s}&\text{smooth elliptic}&\text{nodal}&\text{cuspidal}\\ \hline
 \Pic^0(C_{\bar s})&E&\mathbb G_m&\mathbb G_a\\
 \#T_{\bar s}&4&2&1
 \end{array}
 \tag{PR-count}\label{eq:qa-picard-root-count}
\]
となる。従って滑らかなfiberと節点fiberが同じ連結基底族に
現れるとき、$T\to S$ はその族上で有限 étale
\emph{次数四のまま}延長することはできない。
\end{proposition}
\begin{proof}
$\Pic^0_{C/S}$ は滑らかな相対一次元可換群代数空間で、
$[-2]$ は étale 全射であることを論文Bの
\cite[\texttt{lem:picard-halves}]{QuadricModuli2026} が示す。
滑らかな種数一のfiberのPicard零成分は楕円曲線の
相対Jacobianであり、$[-2]$ は有限 étale
次数 $2^{2}=4$ のisogenyである。
そのpullbackは四点の$E[2]$-torsorとなる。

一節点曲線の正規化は $\PP^1$ で、節点の二つの原像の
fiberの貼合せスカラーにより
$\Pic^0\simeq\mathbb G_m$ となる。
$[-2]$ は $z\mapsto z^{-2}$ なので、
代数閉体の各非零値の逆像は二点である。
一尖点曲線では正規化との貼合せの一次jetの加法群から
$\Pic^0\simeq\mathbb G_a$ となる。
$[-2]$ は $z\mapsto-2z$、従って一意の平方根を持つ。
いずれも標数零では étale である。
最後の主張は有限 étale 射の次数が連結基底上で
局所定数となることから従う。
\end{proof}

\begin{remark}[四点の抽象的な一致と比較すべきtorsor class]
\label{rem:qa-ou-picard-fourfold-not-identical}
OU側の $\Gamma\simeq\mathbb F_2^2$ と
楕円曲線上の $E[2](\bar k)\simeq\mathbb F_2^2$ は、
$S_3=\operatorname{GL}_2(\mathbb F_2)$ の
三非零元に対する作用も含めて同じ抽象モデルを持つ。
これは次数四の一致を説明する\emph{候補}であって、
OUの有限射とPicard平方根空間の同型を証明するものではない。
特にOU側は一般に$S_4$-monodromyの非Galois四重射であり、
Picard側は変動する群スキーム$E[2]$のtorsorである。
二つが同じ基底上に引き戻され、
$E[2]$の局所系とOUの符号変更局所系が同定できた後に、
両torsorの相対同型を確認する必要がある。
標識、自己同型群、さらにpointed同値による商は
その後に初めて比較できる。
節点・尖点における式~\eqref{eq:qa-picard-root-count} は
Picard根空間そのもののfiberの個数であり、
OU射の分岐fiberや$\HHstack$への忘却fiberを
数えたものではない。
\end{remark}

\begin{proposition}[OU楕円四つ組の $z_2$ と適合Picard根]
\label{prop:qa-ou-elliptic-half-coordinate}
滑らかな楕円曲線 $E$ に原点 $o$ を取り、
Okawa--Ueda \cite[\S4, (4.12)--(4.13)]{OU2014}
の楕円四つ組を
\[
 (E,\mathcal L_0,\mathcal L_1,\mathcal L_2)
 =
 (E,\OO_E(2o),\OO_E(2z_1),\OO_E(2z_2))
\]
と書く。論文Bの
$N=\mathcal L_2\mathcal L_0^{-1}\in\Pic^0(E)$ に対する
$a^{-2}=N$ の解集合は
\[
 \left\{
 \OO_E(o-z_2)\otimes\xi:\xi\in E[2]
 \right\}
 \tag{OU-PR}\label{eq:qa-ou-elliptic-half-coordinate}
\]
である。従ってOUの楕円四つ組の birational 表示において、
$\mathcal L_0,\mathcal L_1,\mathcal L_2$ を固定したまま
$z_2$ に2-torsionを加える操作は、
\emph{対応するPicard平方根の四つの選択}に一致する。
これは式~(4.13)の $(\mathbb Z/2)^4$ のうち
$z_2$ に作用する一つの $(\mathbb Z/2)^2$ の因子である。

ただし、式~\eqref{eq:qa-ou-fourfold-diagonal} の
potential GITからquintuple GITへの次数四の
\emph{四つのsheet}がこの $z_2$ 因子のtorsorと同型であることは、
この命題の結論ではない。
\end{proposition}
\begin{proof}
$N=\OO_E(2z_2-2o)$ であるから
$a_0=\OO_E(o-z_2)$ は
$a_0^{\otimes(-2)}=N$ を満たす。
二つの根$a,a_0$の比は二乗が自明な
$\Pic^0(E)[2]=E[2]$ の元であり、逆にその全元は解を与える。
$z_2$ の取り換え $z_2\mapsto z_2+t$（$2t=0$）は
$\OO_E(2z_2)$ を変えず、$a_0$ を
$\OO_E(-t)$ 倍する。従って四つの選択と自由かつ推移的に対応する。
OUの式~(4.11)--(4.13)は $z_1,z_2$ の独立した
2-torsion平行移動を含むので、最後の因子の同定も従う。
\end{proof}

\begin{lemma}[Picard平方根のKummer class]
\label{lem:qa-picard-kummer-torsor}
滑らかな楕円曲線の相対族 $E\to S$ と
$N\in E(S)$ に対し、平方根空間
$T_N=S\times_{E,[-2]}E$ の
$E[2]$-torsor class は、étale siteのKummer列
\[
 0\longrightarrow E[2]\longrightarrow E
 \xrightarrow{[-2]}E\longrightarrow0
\]
から得られる連結写像による
$\delta_{[-2]}(N)\in H^1_{\mathrm{\acute et}}(S,E[2])$
である。
従って、同じ基底上でOUの四重射に
$E[2]$-torsor構造を与える局所系の同型を
別途構成できた場合に限り、そのtorsorとPicard根空間が同型となる
必要十分条件は二つの $H^1$ class の一致である。
\end{lemma}
\begin{proof}
$[-2]$ は滑らかな楕円群スキーム上の有限étale isogenyで
核は$E[2]$なので、示した短完全列はétale層として完全。
$N$ の局所持上げの差から得るČech 1-cocycleこそ
$T_N$ のtorsor classである。
比較先にも同じ構造群スキームを同定した後は、
torsorの同型類の分類が$H^1$で与えられることから従う。
\end{proof}

\subsection{一般非線型 cubic の実際の相対比較射}

\begin{theorem}[geometric quintuple から pointed 相対スタックへの射]
\label{thm:qa-geometric-cubic-relative-forgetting}
$\mathscr G^{\mathrm{nl}}$ を、式~\eqref{eq:qa-ou-scalar-kernel}
のテンソル表示空間 $T\setminus\{0\}$ のうち
Van den Bergh の\emph{非線型 geometric quintuple}
を与える $\widetilde G$-不変開部分とする。
このとき局所有限型 $k$-スキーム上の fppf site において
\[
 \mathcal F_{\mathrm{geo}}:
 [\mathscr G^{\mathrm{nl}}/\widetilde G]
 \longrightarrow\HHstack,\qquad
 (V_0,V_1,V_2,V_3,W\hookrightarrow
       V_0\otimes\cdots\otimes V_3)
 \longmapsto
 (\operatorname{qgr}C,\pi e_0C)
 \tag{GIT-H2}\label{eq:qa-geometric-cubic-comparison-map}
\]
という\emph{実際の相対群oid射}が定義される。
ここで $D$ はquintupleのcubic $\mathbb Z$-代数、
右加群・向きの規約を論文Bに合わせて
$C_{ij}=D_{-j,-i}$ とする。
$\operatorname{qgr}C$ の相対heartは
module-object 基底変更・fppf descent を伴って構成される。

この射は $K\simeq\Gm^4$ の慣性を自明にする自然な
二同型を持ち、その剛体化を通じて
\[
 \overline{\mathcal F}_{\mathrm{geo}}:
 [\mathscr G^{\mathrm{nl}}/(\widetilde G/K)]
 \longrightarrow\HHstack
 \tag{GIT-H2-rig}\label{eq:qa-geometric-cubic-rigid-map}
\]
と因数分解する。

さらに、論文Bの相対Picard二等分空間
$T_{\mathrm{root}}\to\mathscr G^{\mathrm{nl}}$ で選んだ
二周期 $\beta$ と Serre 平方同型を使って得る
$\HHstack$-対象は、選択にかかわらず
式~\eqref{eq:qa-geometric-cubic-comparison-map}
の像の引戻しである。
特に滑らかな楕円曲線上の四つのPicard根は、
\emph{固定quintupleの四つの異なるpointed対象ではない}。

この定理は $\overline{\mathcal F}_{\mathrm{geo}}$ の
表現可能性・全忠実性・幾何ファイバーの有限性、
およびGIT compactification全体への延長を主張しない。
\end{theorem}
\begin{proof}
$\mathscr G^{\mathrm{nl}}$ 上の普遍quintupleは
相対的なcubic $\mathbb Z$-代数 $D$ を定める。
双対・添字反転により $C_{ij}=D_{-j,-i}$ とするのは
論文B \cite[\texttt{thm:geometric-cubic-relative-effectivity}]{QuadricModuli2026}
の右tails規約である。
同論文 \cite[\texttt{prop:general-cubic-canonical-tails}]{QuadricModuli2026}
により、二周期を選ぶ\emph{前}から
$\operatorname{qgr}C$ と $\pi e_0C$ は
全係数基底変更に整合する相対Abel族であり、
compact Morita理論・Serre双加群も基底上に構成される。

一方、同論文
\cite[\texttt{lem:picard-halves},
\texttt{thm:relative-compatible-root}]{QuadricModuli2026}
は普遍基底を被覆するétale局所で
実際の二周期 $\beta$ とそのSerre平方比較を構成し、
\cite[\texttt{thm:geometric-cubic-relative-effectivity}]{QuadricModuli2026}
が同じ相対圏と構造対象に対して
AZ性・全軌道AS4被覆を含む $\HHstack$ の条件を証明する。
これらは\emph{存在条件}で、追加標識ではなく
fppf局所で判定される。
従って基底上にもともとある
$(\operatorname{qgr}C,\pi e_0C)$ が
$\HHstack$ の実際の対象を定める。
ここでは選んだ $\beta$ を降下させたのではなく、
既に存在する圏と構造対象に対して、
必要な局所witnessの存在を確認している。
同じ議論は任意の局所有限型の非被約基底変更で成立する。

quintupleの同型は $\mathbb Z$-代数の同型を誘導し、
右tailsの相対heartおよび $\pi e_0C$ の同型を保つ。
この操作は恒等、合成、任意の基底変更と両立する。
従って $\widetilde G$-equivariantな族として
式~\eqref{eq:qa-geometric-cubic-comparison-map} が
fppf descentにより構成される。

命題~\ref{prop:qa-ou-scalar-gerbe} の頂点ごとの
共役による自然同型は、添字零でスカラーを1に規格化すれば
pointedである。各 $\gamma_i$ は四つのスカラーの
積による指標として選べるので、
$K$ の積に関する二同型の整合性も成立する。
従って $K$ の作用を剛体化した表示スタックへ
この射は降下する。
最後に根空間上で$\beta$を替えても
$\operatorname{qgr}C$と$\pi e_0C$は元の基底から
引き戻した\emph{同じ}対象であり、結論が従う。
\end{proof}

\begin{corollary}[quintuple の剛体化後の相対次元]
\label{cor:qa-quintuple-rigid-dimension}
前定理の剛体化された表示スタックは
$k$ 上smoothで、純次元三である。
その $\HHstack$ への射~\eqref{eq:qa-geometric-cubic-rigid-map}
は同じ局所次元三のsmooth Artinスタック間の射である。
この次元の一致は相対射の étale 性や
pointed同値群の全忠実性を単独で意味しない。
\end{corollary}
\begin{proof}
$\mathscr G^{\mathrm{nl}}$ は16次元テンソル空間の
smoothな開部分である。
$\dim\widetilde G=4\cdot\dim\GL_2+\dim\GL_1=17$、
$\dim K=4$ なので
$\dim(\widetilde G/K)=13$ である。
従ってそのquotient stackの次元は $16-13=3$、
smooth性は群と表示開のsmooth性から従う。
一方、論文B
\cite[\texttt{thm:global-h2-smooth-finite-type}]{QuadricModuli2026}
が $\HHstack$ のsmooth純次元三を与える。
両者の間の相対射の微分・慣性核の計算は
この数値だけからは従わない。
\end{proof}

\begin{remark}[Cで次に必要な具体的計算]
\label{rem:qa-git-next-fiber}
実際の相対射は定理~\ref{thm:qa-geometric-cubic-relative-forgetting} で構成した。
次の段階では、式~\eqref{eq:qa-ou-fourfold-diagonal} の
$S_4$-monodromyと式~\eqref{eq:qa-picard-root-torsor} の
相対$E[2]$-torsorを同じsmooth elliptic quintuple基底へ引き戻し、
二つの四重射のモノドロミーとtorsor classを比較する。
そのうえで一般geometric cubic quintupleの
適合する有限型開において
$\operatorname{Aut}(\mathcal C,\OO)$ と
標識を保存する自己同型群の像を計算する。
次に上の二式を用いて、rulingと四窓を忘れる
pointed相対ファイバーを求める。
その後に、potentialからquintupleへの既存次数四射と照合する。
楕円型の構造対象の格子・特殊な自己同型、
noncubic境界への特殊化は別の段階とする。
$\HHstack$ 全体のgood moduliの存在、
標識適合領域の開性・全表示との相対同値は
現時点では未証明として保持する。
\end{remark}


\section{全相対AS・幾何avatar・QPの比較}
\label{sec:qa-remaining}
\label{prob:qa-all-relative-as-avatar-qp-stacks}

基礎体は標数零の代数閉体とし、局所有限型schemeのfppf siteで
指定のsquare、rotated square、cubic、最小十本矢型を扱う。
相対ASには各成分の有限局所自由性を課し、
QPの射は頂点と重みを保つ全right-equivalenceとする。
最小十本矢の幾何fiber条件と、その最小矢束が線束となる
scheme壁は、非被約基底では異なる。

成分局所自由な全指定型AS族では、命題~
\ref{prop:qa-relative-as-perfectness-automatic}により
有限対角と単純の非最小相対解消が自動的に存在する。
系~\ref{cor:qa-relative-as-line-decorated-foundation}は
有限foundationと双対線束から全整数成分、合成、
全Isomを復元する。系~\ref{cor:qa-curve-local-period-criterion}により、
滑らかな連結射影曲線の冪零肥厚上で局所四周期化できる必要十分条件は、
被約曲線上の双対線束の類が有限生成部分群を生成することである。
raw QPについては、命題~
\ref{prop:qa-raw-qp-isom-equations}が全重み付き同型を、
命題~\ref{prop:qa-qp-as-matrix-recognition}が
全頂点差を使ったAS適合性の判定を与える。
定理~\ref{thm:qa-raw-qp-as-finite-type}は有限窓を補助変数に
AS適合QPの有限型表示を与える。
定理~\ref{thm:qa-all-relative-square-as-heart}と
系~\ref{cor:qa-all-relative-as-type-hearts}は
全指定AS族の相対heartを復元する。
曲線型の元の標識と全Isomは定理~
\ref{thm:qa-flat-curve-avatar-cross-equivalence}で回収し、
平坦な点表示が失敗する変形には向き付き二次式を用いる。
全標準頂点の相対$H_2$開とファイバー、QP比較像と余剰慣性は
定理~\ref{thm:qa-all-as-relative-h2-pointing-fibers}、
\ref{thm:qa-open-qp-canonical-as-image}および
系~\ref{cor:qa-open-qp-h2-pointing-fibers}に記述する。
独立した数学監査と文献照合の範囲は研究ノートに記録する。


点曲線側の全noncubic相対族については、定理~\ref{thm:qa-symmetric-relative-pointed-heart}により
構造対象を保つ比較が成立した。
smooth spectral側の指定pointingについても、
定理~\ref{thm:qa-spectral-avatar-pointed-involution}と
命題~\ref{prop:qa-spectral-atkin-lehner-chi}で像の幾何同型類は決定した。
このsmooth対応は定理~\ref{thm:qa-spectral-relative-nonfaithful-map}と
命題~\ref{prop:qa-spectral-relative-quadratic-gerbe}により、
非忠実な相対射とその二次線束gerbeとして記述できた。
定理~\ref{thm:qa-spectral-no-stack-autoequivalence}は、別の構成でも
一般点で表現可能なliftを得られないことを示す。
第~\ref{sec:qa-spinor-pointing-obstruction}節はnoncubic由来の全四位相を判定し，
rotated零番頂点と最小十本矢型零番頂点を反Serre生成障害で排除する．
第~\ref{sec:qa-general-vertex-az}節は，非被約曲線・古典型を含む
全分類AS型の標準頂点へ体上判定を拡張する．
第~\ref{sec:qa-relative-vertex-determinants}節はこの判定を任意底変更に適合する二行列式開へ移し，
第~\ref{sec:qa-relative-avatar-pointed-recognition}節はflat曲線型square族の零番・三番頂点の実際の相対$\HH$族を構成する．
第~\ref{sec:qa-relative-avatar-heart-comparison}節は中間二頂点も同じheartで認識し，
適合開上の実際のcubic族を復元する．全指定AS族の同じ相対heartとの比較は系~\ref{cor:qa-all-relative-as-type-hearts}で与える。
系~\ref{cor:qa-spectral-involution-no-h2-extension}により、
spectral対応を全noncubic境界からH2へ延長する目標は成立しない。
境界を含む比較にはH2の外側に標的を広げ、
二次積付きtorsion-free spectral対象などの適切なcompactificationを選ぶ必要がある。
命題~\ref{prop:qa-spectral-flat-singular-boundary}は二つの平坦代数モデルを与えるが、
その特殊fiberは周期Extを持つため、smoothなH2とは異なる境界対象である。
第~\ref{sec:qa-tf-quad-artin}節では節点TF標的の有限型Artin性と
balanced部分の群商表示、元のnoncubicの全忠実な開埋込みまで証明した。
第~\ref{sec:qa-spectral-boundary-models}節では，二境界traitのgraph閉包を
正則な$I_2/I_4$模型として構成し，有限\'etaleな符号被覆と二次線束を延長した．
根対一致型には四捩れtranslationで補正した別の模型同型も存在するが，
二conicの順序を交換する．入力二捩れの平方根障害は，任意の有限分岐後の
最小正則模型にも残る．第~\ref{sec:qa-global-fricke-boundary}節では，
全表示groupoidの正規化と平方根選択の比較から，自然な順序付きFricke射を
全noncubicからbalanced標的へ延長し，線束境界上の全2-fiberと慣性を決定した．
この標的の拡張はH2延長不能の結論と両立する．
全整数ASの線束標識には有限型Artin表示の障害がある。
命題~\ref{prop:qa-as-infinite-lines-no-common-atlas}は、
重み付き表示が全整数ASをfppf局所にも被覆できない例を与える。
一方、周期標識付きの許容窓とAS適合QPは有限型に表示される。
raw QPの余剰慣性を除いた後のruling交換は相対mutationの比較である。
個々のraw potential間の右同値は
命題~\ref{prop:qa-raw-qp-isom-equations}の方程式で判定する。

\bibliographystyle{alpha}
\bibliography{references}
\end{document}
